%============================================================
%  Bd. 27 - Endliche Wörter und Klammerungsbäume
%============================================================

\documentclass[main.tex]{subfiles}

\ifSubfilesClassLoaded{
  \BandSetupStandalone{B27}
}{
}

\title{Bd. 27 - Endliche Wörter und Klammerungsbäume}
\author{Martin Kunze}
\date{}
\setcounter{file}{27}

\hypersetup{
  pdftitle={Bd. 27 - Endliche Wörter und Klammerungsbäume},
  pdfauthor={Martin Kunze}
}

\begin{document}

\title{Bd. 27 - Endliche Wörter und Klammerungsbäume}
\author{Martin Kunze}
\date{}
\hypersetup{pageanchor=false}
\maketitle
\hypersetup{pageanchor=true}
\setcounter{tocdepth}{3}
\tableofcontents
\setcounter{file}{27}
\renewcommand{\FormulaBandID}{27}

% Die gemeinsame Tabellendarstellung setzt Formel und Begründung in feste
% Absatzspalten. Für lange Argumentlisten bleiben die lokalen Umbruchhilfen
% dieses Bands aktiv.
\let\BandTwentySevenOriginalRuleRefWithArg\RuleRefWithArg
\ExplSyntaxOn
\cs_set_eq:NN
  \mx_band_twentyseven_original_formula_ref_suffix:n
  \mx_formula_ref_suffix:n
\NewDocumentCommand{\BandTwentySevenBreakArgs}{m}
  {
    \tl_set:Nn \l_tmpa_tl {#1}
    \tl_replace_all:Nnn \l_tmpa_tl {,} {,\allowbreak}
    \tl_use:N \l_tmpa_tl
  }
\ExplSyntaxOff
\makeatletter
\newcol@{G}[0]{>{\raggedleft\arraybackslash\scriptsize}p{.10\linewidth}}
\makeatother
\ExplSyntaxOn
\cs_set_protected:Npn \mx_formula_ref_suffix:n #1
  {
    \IfNoValueF{#1}
      {\tl_if_blank:nF {#1}
        {\allowbreak (\BandTwentySevenBreakArgs{#1})}}
  }
\ExplSyntaxOff
\renewcommand{\RuleRefWithArg}[3]{%
  \ensuremath{\hyperref[rule:#1]{#2}\allowbreak
    (\BandTwentySevenBreakArgs{#3})}}

% Fallbeweise innerhalb eines bereits nummerierten Resultats erhalten eine
% untergeordnete Nummerierung; die Zuordnung der Theorempunkte bleibt erhalten.
\newenvironment{BandTwentySevenProofCases}{%
  \begingroup
  \edef\BandTwentySevenSavedProofPartCounter{\arabic{proofpartnr}}%
  \edef\BandTwentySevenCaseParent{%
    \mxProofPartParent(\romannumeral\numexpr\value{proofpartnr}+1\relax)}%
  \expandafter\ProofPartSetParent\expandafter{\BandTwentySevenCaseParent}%
  \setcounter{proofpartnr}{0}%
}{%
  \setcounter{proofpartnr}{\BandTwentySevenSavedProofPartCounter}%
  \endgroup
}

% Lokale Notation für die in Teil II konstruierten Positionsbäume. Die
% Symbole stehen bereits für die Übersicht am Bandanfang zur Verfügung;
% ihre Definitionen folgen im Kapitel über Klammerungsbäume.
\newcommand{\BinaryAddresses}{\mathcal W_2}
\newcommand{\AddressPrefixMap}[1]{\operatorname{vor}_{#1}}
\newcommand{\AddressPrefix}[2]{\AddressPrefixMap{#1}(#2)}
\newcommand{\AddressGraft}[2]{\Gamma(#1,#2)}
\newcommand{\GoodAddressSet}[1]{\mathsf{VollPraef}(#1)}
\newcommand{\TreePositions}[2]{\operatorname{Pos}_{#1}(#2)}
\newcommand{\AddressEdges}[1]{E_{#1}^{\mathrm{adr}}}
\newcommand{\AddressInner}[1]{I_{#1}^{\mathrm{adr}}}
\newcommand{\AddressLeft}[1]{\lambda_{#1}^{\mathrm{adr}}}
\newcommand{\AddressRight}[1]{\rho_{#1}^{\mathrm{adr}}}
\newcommand{\TreePositionGraph}[2]{\mathcal G_{#1}^{\mathrm{adr}}(#2)}

\chapter{Einführung}

In einem verwurzelten Baum wird jeder Knoten entlang seines eindeutigen
Wurzelwegs erreicht. Beschriftet man die Schritte dieses Weges, so liest
man eine endliche Folge von Buchstaben. Solche Folgen sind aus Band~21
bekannt. Wir untersuchen nun ihren Aufbau: Man beginnt mit dem leeren
Wort und fügt jeweils einen Buchstaben an. Danach betrachten wir, wie
sich dieselbe Buchstabenfolge durch verschiedene binäre Baumformen
klammern lässt.

Der Aufbau greift zwei weitere vertraute Gedanken auf. Band~10 gewinnt
die natürlichen Zahlen als kleinste Menge, die die Null enthält und
unter der Nachfolgerbildung abgeschlossen ist. Band~20 behandelt die Adjunktion an endliche
Mengen und ihre Kardinalzahlen. Für Wörter verbinden wir diese beiden
Gedanken: Eine feste Anfügungsregel erzeugt aus der leeren Menge genau
die gesuchten Gegenstände.

\section{Wörter durch Anfügen aufbauen}

Ein Alphabet \(A\) ist eine beliebige Menge. Ein Wort über \(A\) soll
die Reihenfolge seiner Buchstaben und auch Wiederholungen festhalten.
Die Menge \(\{a,b\}\) allein unterscheidet weder \(ab\) von \(ba\)
noch von \(aba\). Wir speichern deshalb zu jedem angefügten Buchstaben
die Anzahl der bereits gespeicherten Buchstaben. So entsteht zum Beispiel
\[
 \varnothing\ \longmapsto\ \{(0,a)\}
 \ \longmapsto\ \{(0,a),(1,b)\}
 \ \longmapsto\ \{(0,a),(1,b),(2,a)\}.
\]
Die letzte Menge stellt das Wort \(aba\) dar. Die beiden Buchstaben
\(a\) sind an verschiedenen Stellen hinzugekommen.

Für jede endliche Teilmenge \(w\) von \(\mathbb N\times A\) und
jedes \(a\in A\) ist der Anfügungsschritt
\[
 J(w,a)=w\cup\{(\FiniteCard(w),a)\}
\]
bereits mithilfe der Kardinalzahl aus Band~20 beschrieben. Die
Wortmenge ist die kleinste Familie solcher endlichen Teilmengen, welche
die leere Menge enthält und unter allen diesen Schritten abgeschlossen
ist. Sie wird als Durchschnitt abgeschlossener Familien konstruiert,
wie die kleinste induktive Menge in Band~10.

Aus der Konstruktion folgt anschließend, dass jedes erzeugte Wort
genau den Anfangsabschnitt \(\NatSeg{\FiniteCard(w)}\) als Positionsbereich
besitzt, mit einem Buchstaben an jeder Position. Beim Leerwort ist
dieser Bereich leer. Damit erhalten wir die endlichen Folgen aus Band~21
zurück. Die Darstellung als Funktion wird so zu einem bewiesenen
Hilfsmittel für das Lesen einzelner Stellen. Die Wortlänge ist die
endliche Kardinalzahl der aufgebauten Menge.

\section{Von Wurzelwegen zu Klammerungen}

Der Anfügungsprozess lässt sich als Baum aller möglichen Fortsetzungen
veranschaulichen. Das Leerwort steht an der Wurzel; für jeden Buchstaben
führt von einem Wort ein Schritt zu seiner Fortsetzung. Ein einzelner
Wurzelweg wird etwa durch
\[
 \varepsilon\xrightarrow{a}a\xrightarrow{b}ab\xrightarrow{a}aba
\]
veranschaulicht. Das ganze System kann unendlich sein, während jedes
einzelne Wort nach endlich vielen Schritten erreicht wird. Diese Sicht
knüpft an die endlichen Wurzelwege in den möglicherweise unendlichen
Bäumen aus Band~26 an. Die Buchstaben beschriften die Schritte;
verschiedene Schritte dürfen denselben Buchstaben tragen.

Bei einer Klammerung wird zusätzlich festgehalten, wie zwei schon
gebildete Teile zusammengefügt werden. Dafür bleibt der Graphenbegriff
der Bände~22 bis~26 maßgeblich: Jeder innere Knoten hat ein linkes und
ein rechtes Kind, und die Blätter tragen Buchstaben. Verschiedene
Baumformen können dieselbe von links nach rechts gelesene Blattfolge
besitzen. Der Band konstruiert diese Klammerungsbäume, verbindet sie
mit den geordneten vollen Binärbäumen aus Band~26 und untersucht ihre
Blattwörter und Auswertungen.

Die Strukturaxiome werden jeweils eingeführt, sobald ihre Eigenschaften
am konkreten Modell bewiesen sind. Bei Wörtern geschieht dies nach der
Konstruktion der kleinsten abgeschlossenen Familie und den Grundgesetzen
der Anfügung, vor der allgemeinen Konkatenation. Bei Bäumen geschieht
es nach dem Nachweis der Konstruktor- und Trägereigenschaften. Wie die
Peano-Axiome beschreiben diese Axiome den Aufbau der jeweiligen Struktur.
Bei Wörtern steht die Induktion als Axiom W3 bereit; die Rekursion wird
daraus bewiesen. Die verschiedenen Formen der Induktion werden als
gleichwertig begründet. Eine Gegenüberstellung macht die Übersetzung
von Null und Nachfolger in Anfangswort und Buchstabenanfügung sichtbar.
Für den Rekursionsbeweis verwenden wir die Zahlenrekursion aus Band~10:
Sie erzeugt Stufen von Wertzuordnungen. Zahleninduktion sichert deren
Verträglichkeit, Wortinduktion ihre Vollständigkeit auf der Wortstruktur.
Die Modellnachweise begründen die Erfüllbarkeit der Strukturaxiome;
weitere Mengenexistenzaxiome werden nicht benötigt.

Band~28 verwendet diese Rekursionen für Produkte und Umklammerungen;
Band~48 verwendet sie für den Aufbau von Formeln und Herleitungen.

\chapter{Endliche Wörter}

\section{Die Wortmenge aus Erzeugungsregeln}

\subsection{Erzeugung durch Adjunktion}

Ein Alphabet \(A\) ist eine beliebige Menge. Wir beginnen mit der leeren
Menge und fügen einen Buchstaben zusammen mit der Anzahl der bisher
gebildeten Einträge an. So entstehen beispielsweise
\[
  \varnothing,\qquad \{(0,a)\},\qquad
  \{(0,a),(1,b)\},\qquad \{(0,a),(1,b),(2,a)\}.
\]
Die Paare halten Reihenfolge und Wiederholungen fest. Der Aufbau verwendet
nur endliche Teilmengen, ihre Kardinalzahl und die Adjunktion aus Band~20.
Wie bei der Konstruktion der natürlichen Zahlen in Band~10 nehmen wir
anschließend die kleinste unter den angegebenen Regeln abgeschlossene
Menge.

\newcommand{\FiniteOccurrenceSets}[1]{\mathcal E_{#1}}
\newcommand{\OccurrenceAdjunction}[2]{J(#1,#2)}
\newcommand{\WordGenerationFamilies}[1]{\mathfrak C_{#1}}

\FormulaDefDeltaK[Umgebungsbereich der endlichen Teilmengen]{%
  \FiniteOccurrenceSets A\coloneqq\FinSubsets{\mathbb N\times A}
}{FiniteOccurrenceSetUniverseDef}{%
  \DeltaRow{Mengen}{A}
  \DeltaRow{\textbf{Neues Symbol}}{}
  \DeltaPrem{Umgebungsbereich}{\FiniteOccurrenceSets A}
}

Endlich sind die Elemente von \(\FiniteOccurrenceSets A\); der
Umgebungsbereich selbst braucht nicht endlich zu sein.

\Needspace{16\baselineskip}
\FormulaThmDeltaK[Elemente des Umgebungsbereichs]{%
  w\in\FiniteOccurrenceSets A\eqvdash
    w\subseteq\mathbb N\times A\land\Finite w
}{FiniteOccurrenceSetUniverseMembership}{\DeltaRow{Mengen}{A\dsep w}}
\begin{tabproofsplitwide}
  \Needspace{9\baselineskip}
  \proofpartwide{Vom Umgebungsbereich zur endlichen Teilmenge}
    \proofstepwidestar[1]{w\in\FiniteOccurrenceSets A}{\rA}
    \proofstepwidestar[]{\FiniteOccurrenceSets A=\FinSubsets{\mathbb N\times A}}{%
      \FormulaRefAuto{FiniteOccurrenceSetUniverseDef}[def]}
    \proofstepwidestar[1]{w\in\FinSubsets{\mathbb N\times A}}{\rIE{2,1}}
    \proofstepwidestar[1]{w\subseteq\mathbb N\times A\land\Finite w}{%
      \FormulaRefAuto{F\in\FinSubsets{M}\eqvdash F\subseteq M\land\Finite{F}}[thm]{3}}
  \closeproofpartwide
  \Needspace{9\baselineskip}
  \proofpartwide{Von der endlichen Teilmenge zum Umgebungsbereich}
    \proofstepwidestar[1]{w\subseteq\mathbb N\times A\land\Finite w}{\rA}
    \proofstepwidestar[1]{w\in\FinSubsets{\mathbb N\times A}}{%
      \FormulaRefAuto{F\in\FinSubsets{M}\eqvdash F\subseteq M\land\Finite{F}}[thm]{1}}
    \proofstepwidestar[]{\FiniteOccurrenceSets A=\FinSubsets{\mathbb N\times A}}{%
      \FormulaRefAuto{FiniteOccurrenceSetUniverseDef}[def]}
    \proofstepwidestar[1]{w\in\FiniteOccurrenceSets A}{\rIE{3,2}}
  \closeproofpartwide
\end{tabproofsplitwide}

\FormulaDefDeltaK[Adjunktion an der Kardinalzahlposition]{%
  w\in\FiniteOccurrenceSets A\dsep a\in A\vdash
    \OccurrenceAdjunction w a\coloneqq w\cup\{(\FiniteCard(w),a)\}
}{OccurrenceCardinalAdjunctionDef}{%
  \DeltaRow{Mengen}{A\dsep w\dsep a}
  \DeltaRow{\textbf{Neues Symbol}}{}
  \DeltaPrem{Erzeugungsschritt}{\OccurrenceAdjunction w a}
}

Die Kardinalzahl \(\FiniteCard(w)\) ist bereits aus Band~20 bekannt. Sie zählt
hier die Elemente einer endlichen Menge; eine Wortlänge wird für den
Erzeugungsschritt noch nicht vorausgesetzt. Auf einer beliebigen endlichen
Teilmenge muss das angefügte Paar nicht neu sein. Für die Abgeschlossenheit
des Umgebungsbereichs genügt deshalb die Adjunktion eines beliebigen
Elements.

\Needspace{16\baselineskip}
\FormulaThmDeltaK[Die leere Menge liegt im Umgebungsbereich]{%
  \varnothing\in\FiniteOccurrenceSets A
}{FiniteOccurrenceUniverseEmpty}{\DeltaRow{Mengen}{A}}
\begin{tabproofwide}
  \proofstepwidestar[]{\varnothing\subseteq\mathbb N\times A}{%
    \FormulaRefAuto{\varnothing\subseteq A}[thm]}
  \proofstepwidestar[]{\Finite\varnothing}{\FormulaRefAuto{FiniteEmptyAxiom}[ax]}
  \proofstepwidestar[]{\varnothing\subseteq\mathbb N\times A\land\Finite\varnothing}{\rAI{1,2}}
  \proofstepwidestar[]{\varnothing\in\FiniteOccurrenceSets A}{%
    \FormulaRefAuto{FiniteOccurrenceSetUniverseMembership}[thm]{3}}
\end{tabproofwide}

\Needspace{16\baselineskip}
\FormulaThmDeltaK[Der Erzeugungsschritt bleibt im Umgebungsbereich]{%
  w\in\FiniteOccurrenceSets A\dsep a\in A\vdash
    \OccurrenceAdjunction w a\in\FiniteOccurrenceSets A
}{FiniteOccurrenceUniverseClosed}{\DeltaRow{Mengen}{A\dsep w\dsep a}}
\begin{tabproofwide}
  \proofstepwidestar[1]{w\in\FiniteOccurrenceSets A}{\rA}
  \proofstepwidestar[2]{a\in A}{\rA}
  \proofstepwidestar[1]{w\subseteq\mathbb N\times A\land\Finite w}{%
    \FormulaRefAuto{FiniteOccurrenceSetUniverseMembership}[thm]{1}}
  \proofstepwidestar[1]{w\subseteq\mathbb N\times A}{\rAEa{3}}
  \proofstepwidestar[1]{\Finite w}{\rAEb{3}}
  \proofstepwidestar[1]{\FiniteCard(w)\in\mathbb N}{\FormulaRefAuto{FiniteCardNatural}[thm]{5}}
  \proofstepwidestar[1,2]{(\FiniteCard(w),a)\in\mathbb N\times A}{%
    \FormulaRefAuto{(a,b)\in A\times B\eqvdash a\in A\land b\in B}[thm]{\rAI{6,2}}}
  \proofstepwidestar[1,2]{\{(\FiniteCard(w),a)\}\subseteq\mathbb N\times A}{%
    \FormulaRefAuto{a\in A\vdash\{a\}\subseteq A}[thm]{7}}
  \proofstepwidestar[1,2]{w\cup\{(\FiniteCard(w),a)\}\subseteq\mathbb N\times A}{%
    \FormulaRefAuto{A\subseteq C\dsep B\subseteq C\vdash A\cup B\subseteq C}[thm]{4,8}}
  \proofstepwidestar[1]{\Finite{w\cup\{(\FiniteCard(w),a)\}}}{%
    \FormulaRefAuto{FiniteInsert}[thm]{5}}
  \proofstepwidestar[1,2]{w\cup\{(\FiniteCard(w),a)\}\in\FiniteOccurrenceSets A}{%
    \FormulaRefAuto{FiniteOccurrenceSetUniverseMembership}[thm]{\rAI{9,10}}}
  \proofstepwidestar[1,2]{\OccurrenceAdjunction w a=w\cup\{(\FiniteCard(w),a)\}}{%
    \FormulaRefAuto{OccurrenceCardinalAdjunctionDef}[def]{1,2}}
  \proofstepwidestar[1,2]{\OccurrenceAdjunction w a\in\FiniteOccurrenceSets A}{\rIE{12,11}}
\end{tabproofwide}

\subsection{Die kleinste abgeschlossene Menge}

\noindent\begin{minipage}{\linewidth}
\FormulaDefDeltaK[Familie der abgeschlossenen Mengen]{%
  \begin{aligned}[t]
    \WordGenerationFamilies A\coloneqq\bigl\{C\in\powerset(\FiniteOccurrenceSets A)\bigm|{}&
      \varnothing\in C\ \land{}\\[-2pt]
    &\forall w\in C\,\forall a\in A\,
      (\OccurrenceAdjunction w a\in C)\bigr\}
  \end{aligned}
}{WordGenerationFamilyDef}{%
  \DeltaRow{Mengen}{A\dsep C\dsep w\dsep a}
  \DeltaRow{\textbf{Neues Symbol}}{}
  \DeltaPrem{Abgeschlossene Mengen}{\WordGenerationFamilies A}
}
\end{minipage}\par

\Needspace{18\baselineskip}
\noindent\begin{minipage}{\linewidth}
\FormulaThmDeltaK[Elementkriterium der abgeschlossenen Mengen]{%
  \begin{aligned}[t]
    C\in\WordGenerationFamilies A\eqvdash{}&
    C\subseteq\FiniteOccurrenceSets A\ \land{}\\[-2pt]
    &\varnothing\in C\ \land
      \forall w\in C\,\forall a\in A\,(\OccurrenceAdjunction w a\in C)
  \end{aligned}
}{WordGenerationFamilyMembership}{\DeltaRow{Mengen}{A\dsep C\dsep w\dsep a}}
\end{minipage}\par
\begin{tabproofsplitwide}
  \Needspace{9\baselineskip}
  \proofpartwide{Eigenschaften eines Familienglieds}
    \proofstepwidestar[1]{C\in\WordGenerationFamilies A}{\rA}
    \proofstepwidestar[1]{\begin{aligned}[t]
      &C\in\powerset(\FiniteOccurrenceSets A)\land{}\\[-2pt]
      &\varnothing\in C\land
        \forall w\in C\,\forall a\in A\,(\OccurrenceAdjunction w a\in C)
    \end{aligned}}{%
      \FormulaRefAuto{x\in\{u\in A\mid P(u)\}\eqvdash x\in A\land P(x)}[thm]{%
        \rIE{\FormulaRefAuto{WordGenerationFamilyDef}[def],1}}}
    \proofstepwidestar[1]{C\subseteq\FiniteOccurrenceSets A}{%
      \FormulaRefAuto{x\in\powerset(A)\eqvdash x\subseteq A}[thm]{\rAEa{2}}}
    \proofstepwidestar[1]{\begin{aligned}[t]
      &C\subseteq\FiniteOccurrenceSets A\land{}\\[-2pt]
      &\varnothing\in C\land
        \forall w\in C\,\forall a\in A\,(\OccurrenceAdjunction w a\in C)
    \end{aligned}}{\rAI{3,\rAEb{2}}}
  \closeproofpartwide
  \Needspace{9\baselineskip}
  \proofpartwide{Aufnahme einer abgeschlossenen Menge}
    \proofstepwidestar[1]{\begin{aligned}[t]
      &C\subseteq\FiniteOccurrenceSets A\land{}\\[-2pt]
      &\varnothing\in C\land
        \forall w\in C\,\forall a\in A\,(\OccurrenceAdjunction w a\in C)
    \end{aligned}}{\rA}
    \proofstepwidestar[1]{C\in\powerset(\FiniteOccurrenceSets A)}{%
      \FormulaRefAuto{x\in\powerset(A)\eqvdash x\subseteq A}[thm]{\rAEa{1}}}
    \proofstepwidestar[1]{\begin{aligned}[t]
      &C\in\powerset(\FiniteOccurrenceSets A)\land{}\\[-2pt]
      &\varnothing\in C\land
        \forall w\in C\,\forall a\in A\,(\OccurrenceAdjunction w a\in C)
    \end{aligned}}{\rAI{2,\rAEb{1}}}
    \proofstepwidestar[1]{C\in\WordGenerationFamilies A}{%
      \rIE{\FormulaRefAuto{WordGenerationFamilyDef}[def],%
        \FormulaRefAuto{x\in\{u\in A\mid P(u)\}\eqvdash x\in A\land P(x)}[thm]{3}}}
  \closeproofpartwide
\end{tabproofsplitwide}

\Needspace{16\baselineskip}
\FormulaThmDeltaK[Der Umgebungsbereich ist eine abgeschlossene Menge]{%
  \FiniteOccurrenceSets A\in\WordGenerationFamilies A
}{WordGenerationUniverseMember}{\DeltaRow{Mengen}{A\dsep w\dsep a}}
\begin{tabproofwide}
  \proofstepwidestar[]{\FiniteOccurrenceSets A\subseteq\FiniteOccurrenceSets A}{%
    \FormulaRefAuto{A\subseteq A}[thm]}
  \proofstepwidestar[]{\varnothing\in\FiniteOccurrenceSets A}{%
    \FormulaRefAuto{FiniteOccurrenceUniverseEmpty}[thm]}
  \proofstepwidestar[3]{w\in\FiniteOccurrenceSets A}{\rA}
  \proofstepwidestar[4]{a\in A}{\rA}
  \proofstepwidestar[3,4]{\OccurrenceAdjunction w a\in\FiniteOccurrenceSets A}{%
    \FormulaRefAuto{FiniteOccurrenceUniverseClosed}[thm]{3,4}}
  \proofstepwidestar[]{\forall w\in\FiniteOccurrenceSets A\,\forall a\in A\,
    (\OccurrenceAdjunction w a\in\FiniteOccurrenceSets A)}{\rUI{\rRI{3,\rUI{\rRI{4,5}}}}}
  \proofstepwidestar[]{\FiniteOccurrenceSets A\in\WordGenerationFamilies A}{%
    \FormulaRefAuto{WordGenerationFamilyMembership}[thm]{\rAI{1,\rAI{2,6}}}}
\end{tabproofwide}

\Needspace{16\baselineskip}
\FormulaThmDeltaK[Die Familie abgeschlossener Mengen ist nicht leer]{%
  \WordGenerationFamilies A\neq\varnothing
}{WordGenerationFamilyNonempty}{\DeltaRow{Mengen}{A}}
\begin{tabproofwide}
  \proofstepwidestar[]{\FiniteOccurrenceSets A\in\WordGenerationFamilies A}{%
    \FormulaRefAuto{WordGenerationUniverseMember}[thm]}
  \proofstepwidestar[]{\exists C\,(C\in\WordGenerationFamilies A)}{\rEI{1}}
  \proofstepwidestar[]{\WordGenerationFamilies A\neq\varnothing}{%
    \FormulaRefAuto{A\neq\varnothing\eqvdash\exists x\,(x\in A)}[thm]{2}}
\end{tabproofwide}

Damit ist der folgende Durchschnitt nach Band~03 definiert.

\FormulaDefDeltaK[Menge der endlichen Wörter]{%
  \FiniteWords A\coloneqq\bigcap\WordGenerationFamilies A
}{FiniteWordSetDef}{%
  \DeltaRow{Mengen}{A}
  \DeltaRow{\textbf{Neues Symbol}}{}
  \DeltaPrem{Endliche Wörter}{\FiniteWords A}
}

Die Elemente von \(\FiniteWords A\) heißen \textbf{endliche Wörter über
\(A\)}. Die Definition fordert ausschließlich den leeren Ausgangsgegenstand
und den Erzeugungsschritt. Die folgenden Sätze zeigen, dass ihr
Durchschnitt genau die kleinste abgeschlossene Menge ist.

\Needspace{16\baselineskip}
\FormulaThmDeltaK[Elementkriterium des erzeugten Wortbereichs]{%
  w\in\FiniteWords A\leftrightarrow
    \forall C\in\WordGenerationFamilies A\,(w\in C)
}{GeneratedWordMembership}{\DeltaRow{Mengen}{A\dsep w\dsep C}}
\begin{tabproofwide}
  \proofstepwidestar[]{\WordGenerationFamilies A\neq\varnothing}{%
    \FormulaRefAuto{WordGenerationFamilyNonempty}[thm]}
  \proofstepwidestar[]{\begin{aligned}[t]
    &w\in\bigcap\WordGenerationFamilies A\leftrightarrow{}\\[-2pt]
    &\forall C\in\WordGenerationFamilies A\,(w\in C)
  \end{aligned}}{%
    \FormulaRefAuto{M\neq\varnothing\vdash x\in\bigcap M\leftrightarrow\forall X\in M\,(x\in X)}[thm]{1}}
  \proofstepwidestar[]{\FiniteWords A=\bigcap\WordGenerationFamilies A}{%
    \FormulaRefAuto{FiniteWordSetDef}[def]}
  \proofstepwidestar[]{\begin{aligned}[t]
    &w\in\FiniteWords A\leftrightarrow{}\\[-2pt]
    &\forall C\in\WordGenerationFamilies A\,(w\in C)
  \end{aligned}}{\rIE{3,2}}
\end{tabproofwide}

\Needspace{16\baselineskip}
\FormulaThmDeltaK[Jede abgeschlossene Menge enthält alle Wörter]{%
  C\in\WordGenerationFamilies A\vdash\FiniteWords A\subseteq C
}{FiniteWordSetGenerationLeast}{\DeltaRow{Mengen}{A\dsep C\dsep w\dsep D}}
\begin{tabproofwide}
  \proofstepwidestar[1]{C\in\WordGenerationFamilies A}{\rA}
  \proofstepwidestar[2]{w\in\FiniteWords A}{\rA}
  \proofstepwidestar[]{w\in\FiniteWords A\leftrightarrow
    \forall D\in\WordGenerationFamilies A\,(w\in D)}{%
    \FormulaRefAuto{GeneratedWordMembership}[thm]}
  \proofstepwidestar[2]{\forall D\in\WordGenerationFamilies A\,(w\in D)}{\rRE{2,\rLREa{3}}}
  \proofstepwidestar[1,2]{w\in C}{\rRE{1,\rUE{4}}}
  \proofstepwidestar[1]{\FiniteWords A\subseteq C}{%
    \FormulaRefAuto{SubsetIntroductionRule}[thm]{[2]\vdots 5}}
\end{tabproofwide}

\Needspace{16\baselineskip}
\FormulaThmDeltaK[Die Wortmenge liegt im Umgebungsbereich]{%
  \FiniteWords A\subseteq\FiniteOccurrenceSets A
}{FiniteWordSetSubsetUniverse}{\DeltaRow{Mengen}{A}}
\begin{tabproofwide}
  \proofstepwidestar[]{\FiniteOccurrenceSets A\in\WordGenerationFamilies A}{%
    \FormulaRefAuto{WordGenerationUniverseMember}[thm]}
  \proofstepwidestar[]{\FiniteWords A\subseteq\FiniteOccurrenceSets A}{%
    \FormulaRefAuto{FiniteWordSetGenerationLeast}[thm]{1}}
\end{tabproofwide}

\Needspace{16\baselineskip}
\FormulaThmDeltaK[Jedes Wort liegt im Umgebungsbereich]{%
  w\in\FiniteWords A\vdash w\in\FiniteOccurrenceSets A
}{FiniteWordInUniverse}{\DeltaRow{Mengen}{A\dsep w}}
\begin{tabproofwide}
  \proofstepwidestar[1]{w\in\FiniteWords A}{\rA}
  \proofstepwidestar[]{\FiniteWords A\subseteq\FiniteOccurrenceSets A}{%
    \FormulaRefAuto{FiniteWordSetSubsetUniverse}[thm]}
  \proofstepwidestar[1]{w\in\FiniteOccurrenceSets A}{%
    \FormulaRefAuto{A\subseteq B\dsep x\in A\vdash x\in B}[thm]{2,1}}
\end{tabproofwide}

\Needspace{16\baselineskip}
\FormulaThmDeltaK[Jedes Wort ist eine endliche Menge]{%
  w\in\FiniteWords A\vdash\Finite w
}{FiniteWordUnderlyingFinite}{\DeltaRow{Mengen}{A\dsep w}}
\begin{tabproofwide}
  \proofstepwidestar[1]{w\in\FiniteWords A}{\rA}
  \proofstepwidestar[1]{w\in\FiniteOccurrenceSets A}{\FormulaRefAuto{FiniteWordInUniverse}[thm]{1}}
  \proofstepwidestar[1]{w\subseteq\mathbb N\times A\land\Finite w}{%
    \FormulaRefAuto{FiniteOccurrenceSetUniverseMembership}[thm]{2}}
  \proofstepwidestar[1]{\Finite w}{\rAEb{3}}
\end{tabproofwide}

\Needspace{18\baselineskip}
\FormulaThmDeltaK[Das leere Wort]{%
  \varnothing\in\FiniteWords A
}{GeneratedWordEmpty}{\DeltaRow{Mengen}{A\dsep C\dsep w\dsep a}}
\begin{tabproofwide}
  \proofstepwidestar[1]{C\in\WordGenerationFamilies A}{\rA}
  \proofstepwidestar[1]{\begin{aligned}[t]
    &C\subseteq\FiniteOccurrenceSets A\land\varnothing\in C\land{}\\[-2pt]
    &\forall w\in C\,\forall a\in A\,(\OccurrenceAdjunction w a\in C)
  \end{aligned}}{\FormulaRefAuto{WordGenerationFamilyMembership}[thm]{1}}
  \proofstepwidestar[1]{\varnothing\in C}{\rAEa{\rAEb{2}}}
  \proofstepwidestar[]{\forall C\in\WordGenerationFamilies A\,(\varnothing\in C)}{\rUI{\rRI{1,3}}}
  \proofstepwidestar[]{\varnothing\in\FiniteWords A\leftrightarrow
    \forall C\in\WordGenerationFamilies A\,(\varnothing\in C)}{%
    \FormulaRefAuto{GeneratedWordMembership}[thm]}
  \proofstepwidestar[]{\varnothing\in\FiniteWords A}{\rRE{4,\rLREb{5}}}
\end{tabproofwide}

\Needspace{16\baselineskip}
\FormulaThmDeltaK[Wörter bleiben unter dem Erzeugungsschritt abgeschlossen]{%
  w\in\FiniteWords A\dsep a\in A\vdash\OccurrenceAdjunction w a\in\FiniteWords A
}{GeneratedWordClosed}{\DeltaRow{Mengen}{A\dsep w\dsep a\dsep C\dsep u\dsep b}}
\begin{tabproofwide}
  \proofstepwidestar[1]{w\in\FiniteWords A}{\rA}
  \proofstepwidestar[2]{a\in A}{\rA}
  \proofstepwidestar[1]{w\in\FiniteOccurrenceSets A}{\FormulaRefAuto{FiniteWordInUniverse}[thm]{1}}
  \proofstepwidestar[4]{C\in\WordGenerationFamilies A}{\rA}
  \proofstepwidestar[4]{\FiniteWords A\subseteq C}{\FormulaRefAuto{FiniteWordSetGenerationLeast}[thm]{4}}
  \proofstepwidestar[1,4]{w\in C}{%
    \FormulaRefAuto{A\subseteq B\dsep x\in A\vdash x\in B}[thm]{5,1}}
  \proofstepwidestar[4]{\begin{aligned}[t]
    &C\subseteq\FiniteOccurrenceSets A\land\varnothing\in C\land{}\\[-2pt]
    &\forall u\in C\,\forall b\in A\,(\OccurrenceAdjunction u b\in C)
  \end{aligned}}{\FormulaRefAuto{WordGenerationFamilyMembership}[thm]{4}}
  \proofstepwidestar[4]{\forall u\in C\,\forall b\in A\,(\OccurrenceAdjunction u b\in C)}{\rAEb{\rAEb{7}}}
  \proofstepwidestar[1,4]{\forall b\in A\,(\OccurrenceAdjunction w b\in C)}{\rRE{6,\rUE{8}}}
  \proofstepwidestar[1,2,4]{\OccurrenceAdjunction w a\in C}{\rRE{2,\rUE{9}}}
  \proofstepwidestar[1,2]{\forall C\in\WordGenerationFamilies A\,(\OccurrenceAdjunction w a\in C)}{\rUI{\rRI{4,10}}}
  \proofstepwidestar[1,2]{\begin{aligned}[t]
    &\OccurrenceAdjunction w a\in\FiniteWords A\leftrightarrow{}\\[-2pt]
    &\forall C\in\WordGenerationFamilies A\,(\OccurrenceAdjunction w a\in C)
  \end{aligned}}{\FormulaRefAuto{GeneratedWordMembership}[thm]}
  \proofstepwidestar[1,2]{\OccurrenceAdjunction w a\in\FiniteWords A}{\rRE{11,\rLREb{12}}}
\end{tabproofwide}

\Needspace{16\baselineskip}
\FormulaThmDeltaK[Die Wortmenge ist eine abgeschlossene Menge]{%
  \FiniteWords A\in\WordGenerationFamilies A
}{FiniteWordSetGenerationClosed}{\DeltaRow{Mengen}{A\dsep w\dsep a}}
\begin{tabproofwide}
  \proofstepwidestar[]{\FiniteWords A\subseteq\FiniteOccurrenceSets A}{\FormulaRefAuto{FiniteWordSetSubsetUniverse}[thm]}
  \proofstepwidestar[]{\varnothing\in\FiniteWords A}{\FormulaRefAuto{GeneratedWordEmpty}[thm]}
  \proofstepwidestar[3]{w\in\FiniteWords A}{\rA}
  \proofstepwidestar[4]{a\in A}{\rA}
  \proofstepwidestar[3,4]{\OccurrenceAdjunction w a\in\FiniteWords A}{\FormulaRefAuto{GeneratedWordClosed}[thm]{3,4}}
  \proofstepwidestar[]{\forall w\in\FiniteWords A\,\forall a\in A\,(\OccurrenceAdjunction w a\in\FiniteWords A)}{\rUI{\rRI{3,\rUI{\rRI{4,5}}}}}
  \proofstepwidestar[]{\FiniteWords A\in\WordGenerationFamilies A}{%
    \FormulaRefAuto{WordGenerationFamilyMembership}[thm]{\rAI{1,\rAI{2,6}}}}
\end{tabproofwide}

\Needspace{16\baselineskip}
\noindent\begin{minipage}{\linewidth}
\FormulaThmDeltaK[Minimalität der erzeugten Wortmenge]{%
  \begin{aligned}[t]
    &C\subseteq\FiniteOccurrenceSets A\dsep\varnothing\in C\dsep{}\\[-2pt]
    &\forall w\in C\,\forall a\in A\,(\OccurrenceAdjunction w a\in C)
      \vdash\FiniteWords A\subseteq C
  \end{aligned}
}{GeneratedWordMinimality}{\DeltaRow{Mengen}{A\dsep C\dsep w\dsep a}}
\end{minipage}\par
\begin{tabproofwide}
  \proofstepwidestar[1]{C\subseteq\FiniteOccurrenceSets A}{\rA}
  \proofstepwidestar[2]{\varnothing\in C}{\rA}
  \proofstepwidestar[3]{\forall w\in C\,\forall a\in A\,(\OccurrenceAdjunction w a\in C)}{\rA}
  \proofstepwidestar[1,2,3]{C\in\WordGenerationFamilies A}{%
    \FormulaRefAuto{WordGenerationFamilyMembership}[thm]{\rAI{1,\rAI{2,3}}}}
  \proofstepwidestar[1,2,3]{\FiniteWords A\subseteq C}{\FormulaRefAuto{FiniteWordSetGenerationLeast}[thm]{4}}
\end{tabproofwide}

\subsection{Induktion aus der Konstruktion}

Die Minimalität liefert bereits ein Induktionsprinzip für das konkrete
Modell. Es wird hier aus der Schnittkonstruktion bewiesen. Die späteren
Strukturaxiome werden dafür nicht verwendet.

\Needspace{16\baselineskip}
\noindent\begin{minipage}{\linewidth}
\FormulaThmDeltaK[Induktion für erzeugte Wörter]{%
  \begin{aligned}[t]
    &P(\varnothing)\dsep{}\\[-2pt]
    &\forall w\in\FiniteWords A\,\forall a\in A\,
      \bigl(P(w)\rightarrow P(\OccurrenceAdjunction w a)\bigr)\dsep{}\\[-2pt]
    &v\in\FiniteWords A\vdash P(v)
  \end{aligned}
}{GeneratedWordInduction}{%
  \DeltaRow{Mengen}{A\dsep v\dsep w\dsep a}
  \DeltaRow{Prädikat}{P}
  \DeltaRow{Abkürzung}{U=\{w\in\FiniteWords A\mid P(w)\}}
}
\end{minipage}\par
\begin{tabproofwide}
  \proofstepwidestar[1]{P(\varnothing)}{\rA}
  \proofstepwidestar[2]{\forall w\in\FiniteWords A\,\forall a\in A\,
    \bigl(P(w)\rightarrow P(\OccurrenceAdjunction w a)\bigr)}{\rA}
  \proofstepwidestar[3]{v\in\FiniteWords A}{\rA}
  \proofstepwidestar[4]{w\in U}{\rA}
  \proofstepwidestar[4]{w\in\FiniteWords A\land P(w)}{%
    \FormulaRefAuto{x\in\{u\in A\mid P(u)\}\eqvdash x\in A\land P(x)}[thm]{4}}
  \proofstepwidestar[4]{w\in\FiniteWords A}{\rAEa{5}}
  \proofstepwidestar[4]{P(w)}{\rAEb{5}}
  \proofstepwidestar[4]{w\in\FiniteOccurrenceSets A}{\FormulaRefAuto{FiniteWordInUniverse}[thm]{6}}
  \proofstepwidestar[]{U\subseteq\FiniteOccurrenceSets A}{%
    \FormulaRefAuto{SubsetIntroductionRule}[thm]{[4]\vdots 8}}
  \proofstepwidestar[]{\varnothing\in\FiniteWords A}{\FormulaRefAuto{GeneratedWordEmpty}[thm]}
  \proofstepwidestar[1]{\varnothing\in U}{%
    \FormulaRefAuto{x\in\{u\in A\mid P(u)\}\eqvdash x\in A\land P(x)}[thm]{\rAI{10,1}}}
  \proofstepwidestar[12]{a\in A}{\rA}
  \proofstepwidestar[2,4]{\forall a\in A\,
    \bigl(P(w)\rightarrow P(\OccurrenceAdjunction w a)\bigr)}{\rRE{6,\rUE{2}}}
  \proofstepwidestar[2,4,12]{P(w)\rightarrow P(\OccurrenceAdjunction w a)}{\rRE{12,\rUE{13}}}
  \proofstepwidestar[2,4,12]{P(\OccurrenceAdjunction w a)}{\rRE{7,14}}
  \proofstepwidestar[4,12]{\OccurrenceAdjunction w a\in\FiniteWords A}{%
    \FormulaRefAuto{GeneratedWordClosed}[thm]{6,12}}
  \proofstepwidestar[2,4,12]{\OccurrenceAdjunction w a\in U}{%
    \FormulaRefAuto{x\in\{u\in A\mid P(u)\}\eqvdash x\in A\land P(x)}[thm]{\rAI{16,15}}}
  \proofstepwidestar[2]{\forall w\in U\,\forall a\in A\,(\OccurrenceAdjunction w a\in U)}{\rUI{\rRI{4,\rUI{\rRI{12,17}}}}}
  \proofstepwidestar[1,2]{\FiniteWords A\subseteq U}{%
    \FormulaRefAuto{GeneratedWordMinimality}[thm]{9,11,18}}
  \proofstepwidestar[1,2,3]{v\in U}{%
    \FormulaRefAuto{A\subseteq B\dsep x\in A\vdash x\in B}[thm]{19,3}}
  \proofstepwidestar[1,2,3]{v\in\FiniteWords A\land P(v)}{%
    \FormulaRefAuto{x\in\{u\in A\mid P(u)\}\eqvdash x\in A\land P(x)}[thm]{20}}
  \proofstepwidestar[1,2,3]{P(v)}{\rAEb{21}}
\end{tabproofwide}


\Needspace{22\baselineskip}
\noindent\begin{minipage}{\linewidth}
\FormulaThmDeltaK[Induktionsregel für erzeugte Wörter]{%
  \begin{aligned}[t]
    &P(\varnothing)\dsep{}\\[-2pt]
    &[u\in\FiniteWords A,a\in A,P(u)]\vdots
      P(\OccurrenceAdjunction u a)\dsep{}\\[-2pt]
    &v\in\FiniteWords A\vdash P(v)
  \end{aligned}
}{GeneratedWordInductionRule}{%
  \DeltaRow{Mengen}{A\dsep u\dsep a\dsep v}
  \DeltaRow{Einstelliges Prädikat}{P}
}
\end{minipage}\par
Die Regel fasst Anfang, lokalen Anfügungsschritt und Induktionsschluss
zusammen, wie die Induktionsregel aus Band~10. Die Schrittvariablen
\(u,a\) sind beliebig; sie dürfen in den übrigen offenen Annahmen nicht
frei vorkommen. Zusätzliche Parameter des Prädikats bleiben fest.
Eine lokale Annahme darf im Schritt unbenutzt bleiben.
\Needspace{9\baselineskip}
\begin{tabproofwide}
  \proofstepwidestar[1]{P(\varnothing)}{\rA}
  \proofstepwidestar[2]{v\in\FiniteWords A}{\rA}
  \proofstepwidestar[3]{u\in\FiniteWords A}{\rA}
  \proofstepwidestar[4]{a\in A}{\rA}
  \proofstepwidestar[5]{P(u)}{\rA}
  \proofstepwidestar[3,4,5]{P(\OccurrenceAdjunction u a)}{%
    \begin{aligned}[t]
      &\text{Schrittprämisse}\\[-2pt]
      &(3,4,5)
    \end{aligned}}
  \proofstepwidestar[]{\forall u\in\FiniteWords A\,\forall a\in A\,
    (P(u)\rightarrow P(\OccurrenceAdjunction u a))}{%
    \rUI{\rRI{3,\rUI{\rRI{4,\rRI{5,6}}}}}}
  \proofstepwidestar[1,2]{P(v)}{%
    \FormulaRefAuto{GeneratedWordInduction}[thm]{1,7,2}}
\end{tabproofwide}

In Anwendungen genügen die Verweise auf Induktionsanfang (IA),
Induktionsschritt (IS) und Wortzugehörigkeit. Bei einem Schritt innerhalb
derselben Tabelle bezeichnet \([i,j,k]\vdots l\) die Teilableitung mit
den lokalen Annahmen \(i,j,k\) und der Endzeile \(l\). Diese Annahmen
werden beim Induktionsschluss entlassen. Für einen natürlichen Index
\(n\) steht bereits \FormulaRefAuto{PeanoInductionRuleAddOne}[thm] zur
Verfügung; dabei wird \(P(n)\) über Zahlen statt über Wörter bewiesen.


\Needspace{30\baselineskip}
\section{Der Anschluss an endliche Folgen}

% Charakterisierung der durch Leerwort und Adjunktion erzeugten Menge.
% Vorausgesetzt: generated-set sowie die früheren Mengen-, Graphen- und Folgenresultate.
% Keine spätere Wortaxiomatik.
\subsection{Die erzeugten Mengen sind genau die endlichen Folgen}

Die Definition der Wortmenge verwendet nur das Leerwort und einen
Erzeugungsschritt. Jetzt zeigen wir, welche Mengen dabei entstehen:
Genau die endlichen Folgen aus Band~21 gehören zur erzeugten Menge.
Die lückenlose Reihenfolge ist somit eine Folgerung aus den Regeln.

\Needspace{16\baselineskip}
\noindent\begin{minipage}{\linewidth}
\FormulaThmDeltaK[Die Erzeugungsregeln liefern lückenlose Funktionsgraphen]{%
  w\in\FiniteWords A\vdash w\colon\NatSeg{\FiniteCard(w)}\to A
}{GeneratedWordGraphTyping}{%
  \DeltaRow{Mengen}{A\dsep w\dsep u\dsep a}
  \DeltaRow{Abkürzung}{P(u)\coloneqq
    u\colon\NatSeg{\FiniteCard(u)}\to A}}
\end{minipage}\par
Für Wortargumente \(u\in\FiniteWords A\) ist die Endlichkeit bereits
durch \FormulaRefAuto{FiniteWordUnderlyingFinite}[thm] gesichert.
Die Kardinalzahl in \(P(u)\) ist damit definiert. Im Anfangsfall verwenden
wir die bereits bestimmte Kardinalzahl der leeren Menge.
\Needspace{11\baselineskip}
\begin{tabproofsplitwide}
  \proofpartwideindR[Anfang der Erzeugung]{P(\varnothing)}{%
    \varnothing\colon\NatSeg{\FiniteCard(\varnothing)}\to A
  }[GeneratedWordGraphTypingBase]
    \proofstepwidestar[]{\varnothing\colon\varnothing\to A}{%
      \FormulaRefAuto{\varnothing\colon \varnothing\to M}[thm]}
    \proofstepwidestar[]{\NatSeg0=\varnothing}{\FormulaRefAuto{PeanoNatSegZero}[thm]}
    \proofstepwidestar[]{\FiniteCard(\varnothing)=0}{\FormulaRefAuto{B27FiniteCardEmpty}[thm]}
    \proofstepwidestar[]{P(\varnothing)}{\rIE{2,3,1}}
  \closeproofpartwideindR
  \Needspace{11\baselineskip}
  \proofpartwideindR[Ein Erzeugungsschritt]{%
    u\in\FiniteWords A\dsep a\in A\dsep P(u)\vdash P(\OccurrenceAdjunction u a)}{%
    \begin{aligned}[t]
      &u\in\FiniteWords A\dsep a\in A\dsep
        u\colon\NatSeg{\FiniteCard(u)}\to A\vdash{}\\[-2pt]
      &\OccurrenceAdjunction u a\colon\NatSeg{\FiniteCard(\OccurrenceAdjunction u a)}\to A
    \end{aligned}
  }[GeneratedWordGraphTypingStep]
    \proofstepwidestar[1]{u\in\FiniteWords A}{\rA}
    \proofstepwidestar[2]{a\in A}{\rA}
    \proofstepwidestar[3]{P(u)}{\rA}
    \proofstepwidestar[1]{\Finite u}{\FormulaRefAuto{FiniteWordUnderlyingFinite}[thm]{1}}
    \proofstepwidestar[1]{\FiniteCard(u)\in\mathbb N}{\FormulaRefAuto{FiniteCardNatural}[thm]{4}}
    \proofstepwidestar[1,2,3]{u\cup\{(\FiniteCard(u),a)\}\colon\NatSeg{\FiniteCard(u)+1}\to A}{%
      \FormulaRefAuto{EarlyWordAdjunctionGraphTyping}[thm]{5,3,2}}
    \proofstepwidestar[1]{u\in\FiniteOccurrenceSets A}{\FormulaRefAuto{FiniteWordInUniverse}[thm]{1}}
    \proofstepwidestar[1,2]{\OccurrenceAdjunction u a=u\cup\{(\FiniteCard(u),a)\}}{%
      \FormulaRefAuto{OccurrenceCardinalAdjunctionDef}[def]{7,2}}
    \proofstepwidestar[1,2,3]{\OccurrenceAdjunction u a\colon\NatSeg{\FiniteCard(u)+1}\to A}{\rIE{8,6}}
    \proofstepwidestar[1]{\FiniteCard(u)+1\in\mathbb N}{\FormulaRefAuto{PeanoAddOneClosure}[thm]{5}}
    \proofstepwidestar[1,2,3]{\FiniteCard(\OccurrenceAdjunction u a)=\FiniteCard(u)+1}{%
      \FormulaRefAuto{FiniteWordCardinalityFromTyping}[thm]{10,9}}
    \proofstepwidestar[1,2,3]{P(\OccurrenceAdjunction u a)}{\rIE{11,9}}
  \closeproofpartwideindR
  \Needspace{7\baselineskip}
  \proofpartwide{Induktionsschluss}
    \proofstepwidestar[1]{w\in\FiniteWords A}{\rA}
    \proofstepwidestar[1]{w\colon\NatSeg{\FiniteCard(w)}\to A}{\FormulaRefAuto{GeneratedWordInductionRule}[thm]{IA,IS,1}}
  \closeproofpartwide
\end{tabproofsplitwide}

\noindent\begin{minipage}{\linewidth}
\FormulaThmDeltaK[Jede endliche Folge wird durch die Wortregeln erzeugt]{%
  n\in\mathbb N\dsep w\colon\NatSeg n\to A\vdash w\in\FiniteWords A
}{B27FiniteGraphGenerated}{%
  \DeltaRow{Mengen}{A\dsep n\dsep k\dsep w\dsep u}
  \DeltaRow{Abkürzungen}{Q(k)\coloneqq
    \forall u\,(u\colon\NatSeg k\to A\rightarrow u\in\FiniteWords A)
    \dsep r\coloneqq u\restriction_{\NatSeg k}}}
\end{minipage}\par
Hier wird die Zahleninduktion aus Band~10 benutzt: Ein Graph über
\(\NatSeg{k+1}\) entsteht aus seiner Einschränkung auf \(\NatSeg k\)
durch genau den zugelassenen Erzeugungsschritt.
\Needspace{11\baselineskip}
\begin{tabproofsplitwide}
  \proofpartwideindR[Induktionsanfang]{Q(0)}{%
    \forall u\,(u\colon\NatSeg0\to A\rightarrow u\in\FiniteWords A)
  }[B27FiniteGraphGeneratedBase]
    \proofstepwidestar[1]{u\colon\NatSeg0\to A}{\rA}
    \proofstepwidestar[]{\NatSeg0=\varnothing}{\FormulaRefAuto{PeanoNatSegZero}[thm]}
    \proofstepwidestar[1]{u\colon\varnothing\to A}{\rIE{2,1}}
    \proofstepwidestar[1]{u=\varnothing}{\FormulaRefAuto{f\colon \varnothing\to M \vdash f=\varnothing}[thm]{3}}
    \proofstepwidestar[]{\varnothing\in\FiniteWords A}{\FormulaRefAuto{GeneratedWordEmpty}[thm]}
    \proofstepwidestar[1]{u\in\FiniteWords A}{\rIE{4,5}}
    \proofstepwidestar[]{Q(0)}{\rUI{\rRI{1,6}}}
  \closeproofpartwideindR
  \Needspace{11\baselineskip}
  \proofpartwideindR[Induktionsschritt]{k\in\mathbb N\dsep Q(k)\vdash Q(k+1)}{%
    \begin{aligned}[t]
      &k\in\mathbb N\dsep
        \forall u\,(u\colon\NatSeg k\to A\rightarrow u\in\FiniteWords A)\vdash{}\\[-2pt]
      &\forall u\,(u\colon\NatSeg{k+1}\to A\rightarrow u\in\FiniteWords A)
    \end{aligned}
  }[B27FiniteGraphGeneratedStep]
    \proofstepwidestar[1]{k\in\mathbb N}{\rA}
    \proofstepwidestar[2]{Q(k)}{\rA}
    \proofstepwidestar[3]{u\colon\NatSeg{k+1}\to A}{\rA}
    \proofstepwidestar[1]{\NatSeg k\subseteq\NatSeg{k+1}}{\FormulaRefAuto{PeanoNatSegSubsetAddOne}[thm]{1}}
    \proofstepwidestar[1,3]{r\colon\NatSeg k\to A}{\FormulaRefAuto{RestrictionDef}[def]{3,4}}
    \proofstepwidestar[1,2,3]{r\in\FiniteWords A}{\rRE{5,\rUE{2}}}
    \proofstepwidestar[1,3]{\FiniteCard(r)=k}{\FormulaRefAuto{FiniteWordCardinalityFromTyping}[thm]{1,5}}
    \proofstepwidestar[1]{k\in\NatSeg{k+1}}{\FormulaRefAuto{PeanoNatSegEndpointInAddOne}[thm]{1}}
    \proofstepwidestar[1,3]{u(k)\in A}{\FormulaRefAuto{F\colon A\to B\dsep x\in A\vdash F(x)\in B}[thm]{3,8}}
    \proofstepwidestar[1,2,3]{\OccurrenceAdjunction r{u(k)}\in\FiniteWords A}{\FormulaRefAuto{GeneratedWordClosed}[thm]{6,9}}
    \proofstepwidestar[1,2,3]{r\in\FiniteOccurrenceSets A}{\FormulaRefAuto{FiniteWordInUniverse}[thm]{6}}
    \proofstepwidestar[1,2,3]{\OccurrenceAdjunction r{u(k)}=r\cup\{(\FiniteCard(r),u(k))\}}{\FormulaRefAuto{OccurrenceCardinalAdjunctionDef}[def]{11,9}}
    \proofstepwidestar[1,3]{u=r\cup\{(k,u(k))\}}{\FormulaRefAuto{B27FiniteGraphLastAdjunction}[thm]{1,3}}
    \proofstepwidestar[1,2,3]{u=\OccurrenceAdjunction r{u(k)}}{\rIE{7,12,13}}
    \proofstepwidestar[1,2,3]{u\in\FiniteWords A}{\rIE{14,10}}
    \proofstepwidestar[1,2]{Q(k+1)}{\rUI{\rRI{3,15}}}
  \closeproofpartwideindR
  \Needspace{7\baselineskip}
  \proofpartwide{Schluss durch Zahleninduktion}
    \proofstepwidestar[1]{n\in\mathbb N}{\rA}
    \proofstepwidestar[2]{w\colon\NatSeg n\to A}{\rA}
    \proofstepwidestar[1]{Q(n)}{\FormulaRefAuto{PeanoInductionRuleAddOne}[thm]{IA,IS,1}}
    \proofstepwidestar[1,2]{w\in\FiniteWords A}{\rRE{2,\rUE{3}}}
  \closeproofpartwide
\end{tabproofsplitwide}

\Needspace{17\baselineskip}
\noindent\begin{minipage}{\linewidth}
\FormulaThmDeltaK[Funktionenkriterium der endlichen Wortmenge]{%
  w\in\FiniteWords A\eqvdash
  \Bigl(w\in\powerset(\mathbb N\times A)\land
    \exists n\in\mathbb N\,(w\colon\NatSeg n\to A)\Bigr)
}{FiniteWordSetMembership}{\DeltaRow{Mengen}{A\dsep w\dsep n}}
\end{minipage}\par
\Needspace{7\baselineskip}
\begin{tabproofsplitwide}
  \proofpartwide{Von der Erzeugung zum endlichen Funktionsgraphen}
    \proofstepwidestar[1]{w\in\FiniteWords A}{\rA}
    \proofstepwidestar[1]{w\in\FiniteOccurrenceSets A}{\FormulaRefAuto{FiniteWordInUniverse}[thm]{1}}
    \proofstepwidestar[1]{w\subseteq\mathbb N\times A\land\Finite w}{\FormulaRefAuto{FiniteOccurrenceSetUniverseMembership}[thm]{2}}
    \proofstepwidestar[1]{w\in\powerset(\mathbb N\times A)}{\FormulaRefAuto{x\in\powerset(A)\eqvdash x\subseteq A}[thm]{\rAEa{3}}}
    \proofstepwidestar[1]{\FiniteCard(w)\in\mathbb N}{\FormulaRefAuto{FiniteCardNatural}[thm]{\rAEb{3}}}
    \proofstepwidestar[1]{w\colon\NatSeg{\FiniteCard(w)}\to A}{\FormulaRefAuto{GeneratedWordGraphTyping}[thm]{1}}
    \proofstepwidestar[1]{\exists n\in\mathbb N\,(w\colon\NatSeg n\to A)}{\rEI{\rAI{5,6}}}
    \proofstepwidestar[1]{w\in\powerset(\mathbb N\times A)\land
      \exists n\in\mathbb N\,(w\colon\NatSeg n\to A)}{\rAI{4,7}}
  \closeproofpartwide
  \Needspace{7\baselineskip}
  \proofpartwide{Vom endlichen Funktionsgraphen zur Erzeugung}
    \proofstepwidestar[1]{w\in\powerset(\mathbb N\times A)\land
      \exists n\in\mathbb N\,(w\colon\NatSeg n\to A)}{\rA}
    \proofstepwidestar[2]{n\in\mathbb N\land w\colon\NatSeg n\to A}{\rA}
    \proofstepwidestar[2]{w\in\FiniteWords A}{\FormulaRefAuto{B27FiniteGraphGenerated}[thm]{\rAEa{2},\rAEb{2}}}
    \proofstepwidestar[1]{w\in\FiniteWords A}{\rEE{\rAEb{1},2,3}}
  \closeproofpartwide
\end{tabproofsplitwide}

Damit darf ein Wort jetzt als endliche Folge gelesen werden. Die
Folgenbeschreibung ersetzt die Erzeugungsdefinition nicht, sondern
beschreibt genau ihre Elemente und stellt die Positionsnotation bereit.




\noindent\begin{minipage}{\linewidth}
\FormulaThmDeltaK[Einführung eines endlichen Wortes]{%
  w\in\powerset(\mathbb N\times A)\land
    \exists n\in\mathbb N\,(w\colon\NatSeg n\to A)
  \vdash w\in\FiniteWords A%
}{FiniteWordSetIntroduction}{%
  \DeltaRow{Mengen}{A\dsep w\dsep n}
}
\end{minipage}\par
\begin{tabproofwide}
  \proofstepwidestar[1]{%
    w\in\powerset(\mathbb N\times A)\land
      \exists n\in\mathbb N\,(w\colon\NatSeg n\to A)}{\rA}
  \proofstepwidestar[1]{w\in\FiniteWords A}{%
    \FormulaRefAuto{FiniteWordSetMembership}[thm]{1}}
\end{tabproofwide}

\FormulaThmDeltaK[Typisierte endliche Folgen sind endliche Wörter]{%
  n\in\mathbb N\dsep w\colon\NatSeg{n}\to A
  \vdash w\in\FiniteWords{A}%
}{FiniteWordFromTyping}{%
  \DeltaRow{Mengen}{A\dsep n\dsep k\dsep w}
}
\begin{tabproofwide}
  \proofstepwidestar[1]{n\in\mathbb N}{\rA}
  \proofstepwidestar[2]{w\colon\NatSeg{n}\to A}{\rA}
  \proofstepwidestar[1,2]{w\in\FiniteWords{A}}{%
    \FormulaRefAuto{B27FiniteGraphGenerated}[thm]{1,2}}
\end{tabproofwide}

\FormulaThmDeltaK[Eindeutige Länge eines endlichen Wortes]{%
  w\in\FiniteWords{A}\vdash
  \exists!n\in\mathbb N\,
    (w\colon\NatSeg{n}\to A)%
}{FiniteWordLengthUnique}{%
  \DeltaRow{Mengen}{A\dsep w\dsep m\dsep n}
}
\begin{tabproofsplitwide}
    \proofpartwideindR[Existenz einer Wortlänge]{%
    \begin{aligned}[t]
      &w\in\FiniteWords{A}\vdash{}\\[-2pt]
      &\exists n\in\mathbb N\,(w\colon\NatSeg{n}\to A)
    \end{aligned}
  }{\begin{aligned}[t]
      &w\in\FiniteWords{A}\vdash{}\\[-2pt]
      &\exists n\in\mathbb N\,(w\colon\NatSeg{n}\to A)
    \end{aligned}}[FiniteWordLengthExistencePart]
    \proofstepwidestar[1]{%
      w\in\FiniteWords{A}}{%
      \rA}
    \proofstepwidestar[1]{%
      w\in\powerset(\mathbb N\times A)\land
      \exists n\in\mathbb N\,(w\colon\NatSeg{n}\to A)}{%
      \FormulaRefAuto{FiniteWordSetMembership}[thm]{1}}
    \proofstepwidestar[1]{%
      \exists n\in\mathbb N\,(w\colon\NatSeg{n}\to A)}{%
      \rAEb{2}}
  \closeproofpartwideindR

  \proofpartwide{Zusammenführung der Fälle}
    \proofstepwidestar[1]{%
      w\in\FiniteWords{A}}{%
      \rA}
    \proofstepwidestar[1]{%
      \exists n\in\mathbb N\,(w\colon\NatSeg{n}\to A)}{%
      \FormulaRefAuto{FiniteWordLengthExistencePart}[thm]{1}}
    \proofstepwidestar[3]{%
      m\in\mathbb N\land
    w\colon\NatSeg{m}\to A}{%
      \rA}
    \proofstepwidestar[4]{%
      n\in\mathbb N\land
    w\colon\NatSeg{n}\to A}{%
      \rA}
    \proofstepwidestar[3,4]{%
      m=n}{%
      \FormulaRefAuto{FiniteWordLengthUniquenessPart}[thm]{3,4}}
    \proofstepwidestar[1]{%
      \exists!n\in\mathbb N\,(w\colon\NatSeg{n}\to A)}{%
      \UEI{2,3,4,5}}
  \closeproofpartwide

\end{tabproofsplitwide}

\section{Wortlänge}

Die Länge zählt die Stellen eines Wortes; gleiche Buchstaben an
verschiedenen Stellen werden mehrfach gezählt. Wir übernehmen dafür
die endliche Kardinalzahl aus Band~20. Nach
\FormulaRefAuto{FiniteWordCardinalityFromTyping}[thm] stimmt sie mit
der Länge der zugehörigen endlichen Folge überein.

\FormulaDefDeltaK[Länge eines endlichen Wortes]{%
  w\in\FiniteWords A\vdash\WordLength w\coloneqq\FiniteCard(w)
}{FiniteWordLengthDef}{%
  \DeltaRow{Mengen}{A\dsep w}
  \DeltaRow{\textbf{Neues Symbol}}{}
  \DeltaPrem{Wortlänge}{\WordLength w}
}

\FormulaThmDeltaK[Wortlänge ist die Kardinalzahl der aufgebauten Menge]{%
  w\in\FiniteWords A\vdash\WordLength w=\FiniteCard(w)
}{FiniteWordLengthCardinality}{\DeltaRow{Mengen}{A\dsep w}}
\begin{tabproofwide}
  \proofstepwidestar[1]{w\in\FiniteWords A}{\rA}
  \proofstepwidestar[1]{\Finite w}{%
    \FormulaRefAuto{FiniteWordUnderlyingFinite}[thm]{1}}
  \proofstepwidestar[1]{\WordLength w=\FiniteCard(w)}{%
    \FormulaRefAuto{FiniteWordLengthDef}[def]{1}}
\end{tabproofwide}

\Needspace{18\baselineskip}
\noindent\begin{minipage}{\linewidth}
\FormulaThmDeltaK[Typisierung durch die Wortlänge]{%
  w\in\FiniteWords A\vdash
  \begin{aligned}[t]
    &\text{(i)}\;\WordLength w\in\mathbb N\\[-2pt]
    &\text{(ii)}\;w\colon\NatSeg{\WordLength w}\to A
  \end{aligned}
}{FiniteWordTyping}{\DeltaRow{Mengen}{A\dsep w\dsep n}}
\end{minipage}\par
\Needspace{11\baselineskip}
\begin{tabproofsplitwide}
  \proofpartwideindR[Natürlichkeit der Wortlänge]{%
    w\in\FiniteWords A\vdash\WordLength w\in\mathbb N
  }{w\in\FiniteWords A\vdash\WordLength w\in\mathbb N}[FiniteWordLengthNatural]
    \proofstepwidestar[1]{w\in\FiniteWords A}{\rA}
    \proofstepwidestar[1]{\Finite w}{%
      \FormulaRefAuto{FiniteWordUnderlyingFinite}[thm]{1}}
    \proofstepwidestar[1]{\FiniteCard(w)\in\mathbb N}{%
      \FormulaRefAuto{FiniteCardNatural}[thm]{2}}
    \proofstepwidestar[1]{\WordLength w=\FiniteCard(w)}{%
      \FormulaRefAuto{FiniteWordLengthCardinality}[thm]{1}}
    \proofstepwidestar[1]{\WordLength w\in\mathbb N}{\rIE{4,3}}
  \closeproofpartwideindR
  \Needspace{11\baselineskip}
  \proofpartwideindR[Typisierung des Wortes]{%
    w\in\FiniteWords A\vdash w\colon\NatSeg{\WordLength w}\to A
  }{w\in\FiniteWords A\vdash w\colon\NatSeg{\WordLength w}\to A}[FiniteWordTypingAtLength]
    \proofstepwidestar[1]{w\in\FiniteWords A}{\rA}
    \proofstepwidestar[1]{\exists n\in\mathbb N\,(w\colon\NatSeg n\to A)}{%
      \FormulaRefAuto{FiniteWordLengthExistencePart}[thm]{1}}
    \proofstepwidestar[3]{n\in\mathbb N\land w\colon\NatSeg n\to A}{\rA}
    \proofstepwidestar[3]{\FiniteCard(w)=n}{%
      \FormulaRefAuto{FiniteWordCardinalityFromTyping}[thm]{\rAEa{3},\rAEb{3}}}
    \proofstepwidestar[1]{\WordLength w=\FiniteCard(w)}{%
      \FormulaRefAuto{FiniteWordLengthCardinality}[thm]{1}}
    \proofstepwidestar[1,3]{w\colon\NatSeg{\WordLength w}\to A}{\rIE{5,4,\rAEb{3}}}
    \proofstepwidestar[1]{w\colon\NatSeg{\WordLength w}\to A}{\rEE{2,3,6}}
  \closeproofpartwideindR
\end{tabproofsplitwide}


\FormulaThmDeltaK[Länge aus einer gegebenen Worttypisierung]{%
  w\in\FiniteWords{A}\dsep n\in\mathbb N\dsep
  w\colon\NatSeg{n}\to A
  \vdash \WordLength{w}=n%
}{FiniteWordLengthFromTyping}{%
  \DeltaRow{Mengen}{A\dsep w\dsep n\dsep k}
}
\begin{tabproofwide}
  \proofstepwidestar[1]{w\in\FiniteWords A}{\rA}
  \proofstepwidestar[2]{n\in\mathbb N}{\rA}
  \proofstepwidestar[3]{w\colon\NatSeg n\to A}{\rA}
  \proofstepwidestar[2,3]{\FiniteCard(w)=n}{%
    \FormulaRefAuto{FiniteWordCardinalityFromTyping}[thm]{2,3}}
  \proofstepwidestar[1]{\WordLength w=\FiniteCard(w)}{%
    \FormulaRefAuto{FiniteWordLengthCardinality}[thm]{1}}
  \proofstepwidestar[1,2,3]{\WordLength w=n}{\rChain{5,4}}
\end{tabproofwide}

\Needspace{16\baselineskip}
\section{Leerwort, Buchstabenwörter und Nichtleerwörter}

\noindent\begin{minipage}{\linewidth}
\FormulaDefDeltaK[Leerwort]{%
  \EmptyWord{A}\coloneqq\varnothing%
}{EmptyWordDef}{%
  \DeltaRow{Mengen}{A\dsep\EmptyWord{A}}
  \DeltaRow{\textbf{Neues Symbol}}{}
  \DeltaPrem{Leerwort}{\EmptyWord{A}}
}
\end{minipage}\par


\FormulaThmDeltaK[Typisierung des Leerwortes]{%
  \EmptyWord A\colon\NatSeg 0\to A%
}{EmptyWordTyping}{%
  \DeltaRow{Mengen}{A}
}
\begin{tabproofwide}
  \proofstepwidestar[]{\NatSeg 0=\varnothing}{%
    \FormulaRefAuto{PeanoNatSegZero}[thm]}
  \proofstepwidestar[]{\EmptyWord A=\varnothing}{%
    \FormulaRefAuto{EmptyWordDef}[def]}
  \proofstepwidestar[]{\varnothing\colon\varnothing\to A}{%
    \FormulaRefAuto{\varnothing\colon\varnothing\to M}[thm]}
  \proofstepwidestar[]{\EmptyWord A\colon\NatSeg 0\to A}{\rIE{1,2,3}}
\end{tabproofwide}

\Needspace{18\baselineskip}
\noindent\begin{minipage}{\linewidth}
\FormulaThmDeltaK[Das Leerwort hat Länge null]{%
  \begin{aligned}[t]
    &\text{(i)}\;\EmptyWord{A}\in\FiniteWords{A}\\[-2pt]
    &\text{(ii)}\;\WordLength{\EmptyWord{A}}=0
  \end{aligned}%
}{EmptyWordLength}{%
  \DeltaRow{Mengen}{A}
}
\end{minipage}\par
\Needspace{11\baselineskip}
\begin{tabproofsplitwide}
  \proofpartwideindR[Das Leerwort ist ein endliches Wort]{%
    \EmptyWord A\in\FiniteWords A
  }{\EmptyWord A\in\FiniteWords A}[EmptyWordFinite]
    \proofstepwidestar[]{0\in\mathbb N}{\FormulaRefAuto{PeanoZero}[ax]}
    \proofstepwidestar[]{\EmptyWord A\colon\NatSeg 0\to A}{%
      \FormulaRefAuto{EmptyWordTyping}[thm]}
    \proofstepwidestar[]{\EmptyWord A\in\FiniteWords A}{%
      \FormulaRefAuto{FiniteWordFromTyping}[thm]{1,2}}
  \closeproofpartwideindR
  \Needspace{11\baselineskip}
  \proofpartwideindR[Die Länge ist null]{%
    \WordLength{\EmptyWord A}=0
  }{\WordLength{\EmptyWord A}=0}[EmptyWordLengthZero]
    \proofstepwidestar[]{\EmptyWord A\in\FiniteWords A}{%
      \FormulaRefAuto{EmptyWordFinite}[thm]}
    \proofstepwidestar[]{0\in\mathbb N}{\FormulaRefAuto{PeanoZero}[ax]}
    \proofstepwidestar[]{\EmptyWord A\colon\NatSeg 0\to A}{%
      \FormulaRefAuto{EmptyWordTyping}[thm]}
    \proofstepwidestar[]{\WordLength{\EmptyWord A}=0}{%
      \FormulaRefAuto{FiniteWordLengthFromTyping}[thm]{1,2,3}}
  \closeproofpartwideindR
\end{tabproofsplitwide}

\Needspace{18\baselineskip}
\FormulaThmDeltaK[Länge null kennzeichnet das Leerwort]{%
  w\in\FiniteWords A\vdash
    \bigl(\WordLength w=0\leftrightarrow w=\EmptyWord A\bigr)%
}{FiniteWordLengthZeroIffEmpty}{%
  \DeltaRow{Mengen}{A\dsep w}
}
\Needspace{11\baselineskip}
\begin{tabproofsplitwide}
  \proofpartwideindR[Von Länge null zum Leerwort]{%
    w\in\FiniteWords A\dsep\WordLength w=0\vdash w=\EmptyWord A
  }{w\in\FiniteWords A\dsep\WordLength w=0\vdash w=\EmptyWord A}%
    [FiniteWordLengthZeroImpliesEmpty]
    \proofstepwidestar[1]{w\in\FiniteWords A}{\rA}
    \proofstepwidestar[2]{\WordLength w=0}{\rA}
    \proofstepwidestar[1]{w\colon\NatSeg{\WordLength w}\to A}{%
      \FormulaRefAuto{FiniteWordTypingAtLength}[thm]{1}}
    \proofstepwidestar[]{\NatSeg 0=\varnothing}{%
      \FormulaRefAuto{PeanoNatSegZero}[thm]}
    \proofstepwidestar[1,2]{w\colon\varnothing\to A}{\rIE{2,3,4}}
    \proofstepwidestar[1,2]{w=\varnothing}{%
      \FormulaRefAuto{f\colon\varnothing\to M\vdash f=\varnothing}[thm]{5}}
    \proofstepwidestar[]{\EmptyWord A=\varnothing}{%
      \FormulaRefAuto{EmptyWordDef}[def]}
    \proofstepwidestar[1,2]{w=\EmptyWord A}{\rIE{7,6}}
  \closeproofpartwideindR
  \Needspace{11\baselineskip}
  \proofpartwideindR[Vom Leerwort zur Länge null]{%
    w\in\FiniteWords A\dsep w=\EmptyWord A\vdash\WordLength w=0
  }{w\in\FiniteWords A\dsep w=\EmptyWord A\vdash\WordLength w=0}%
    [FiniteWordEmptyImpliesLengthZero]
    \proofstepwidestar[1]{w=\EmptyWord A}{\rA}
    \proofstepwidestar[]{\WordLength{\EmptyWord A}=0}{%
      \FormulaRefAuto{EmptyWordLengthZero}[thm]}
    \proofstepwidestar[1]{\WordLength w=0}{\rIE{1,2}}
  \closeproofpartwideindR
\end{tabproofsplitwide}

\Needspace{12\baselineskip}
\FormulaThmDeltaK[Nichtleeres Wort und von null verschiedene Länge]{%
  w\in\FiniteWords A\vdash
    \bigl(w\neq\EmptyWord A\leftrightarrow\WordLength w\neq0\bigr)%
}{FiniteWordLengthNonzeroIffNonempty}{%
  \DeltaRow{Mengen}{A\dsep w}
}
\begin{tabproofwide}
  \proofstepwidestar[1]{w\in\FiniteWords A}{\rA}
  \proofstepwidestar[1]{\WordLength w=0\leftrightarrow w=\EmptyWord A}{%
    \FormulaRefAuto{FiniteWordLengthZeroIffEmpty}[thm]{1}}
  \proofstepwidestar[1]{\WordLength w\neq0\leftrightarrow w\neq\EmptyWord A}{%
    \FormulaRefAuto{P \leftrightarrow Q \eqvdash
      \neg P \leftrightarrow \neg Q}[thm]{2}}
  \proofstepwidestar[1]{w\neq\EmptyWord A\leftrightarrow\WordLength w\neq0}{%
    \FormulaRefAuto{P \leftrightarrow Q \vdash Q \leftrightarrow P}[thm]{3}}
\end{tabproofwide}


\Needspace{9\baselineskip}
\FormulaDefDeltaK[Menge der nichtleeren Wörter]{%
  \NonemptyWords{A}\coloneqq
  \FiniteWords{A}\setminus\{\EmptyWord{A}\}%
}{NonemptyWordSetDef}{%
  \DeltaRow{Mengen}{A\dsep\FiniteWords{A}\dsep\NonemptyWords{A}}
  \DeltaRow{\textbf{Neues Symbol}}{}
  \DeltaPrem{Nichtleere Wörter}{\NonemptyWords{A}}
}


\Needspace{18\baselineskip}
\FormulaThmDeltaK[Elementkriterium für nichtleere Wörter]{%
  w\in\NonemptyWords A\eqvdash
  w\in\FiniteWords A\land\WordLength w\neq0%
}{NonemptyWordMembership}{%
  \DeltaRow{Mengen}{A\dsep w}
}
\Needspace{11\baselineskip}
\begin{tabproofsplit}
  \proofpartindR[Elimination]{%
    w\in\NonemptyWords A\vdash w\in\FiniteWords A\land\WordLength w\neq0
  }{w\in\NonemptyWords A\vdash w\in\FiniteWords A\land\WordLength w\neq0}%
    [NonemptyWordMembershipElimination]
    \proofstep{1}{w\in\NonemptyWords A}{\rA}
    \proofstep{1}{w\in\FiniteWords A\setminus\{\EmptyWord A\}}{%
      \FormulaRefAuto{NonemptyWordSetDef}[def]{1}}
    \proofstep{1}{w\in\FiniteWords A\land w\notin\{\EmptyWord A\}}{%
      \FormulaRefAuto{x\in A\setminus B\eqvdash x\in A\land x\notin B}[thm]{2}}
    \proofstep{1}{w\in\FiniteWords A}{\rAEa{3}}
    \proofstep{1}{w\neq\EmptyWord A}{%
      \FormulaRefAuto{x\notin\{a\}\eqvdash x\neq a}[thm]{\rAEb{3}}}
    \proofstep{1}{w\neq\EmptyWord{A}\leftrightarrow\WordLength{w}\neq0}{%
      \FormulaRefAuto{FiniteWordLengthNonzeroIffNonempty}[thm]{4}}
    \proofstep{1}{\WordLength w\neq0}{%
      \rRE{5,\rLREa{6}}}
    \proofstep{1}{w\in\FiniteWords A\land\WordLength w\neq0}{\rAI{4,7}}
  \closeproofpartindR
  \Needspace{11\baselineskip}
  \proofpartindR[Einführung]{%
    w\in\FiniteWords A\land\WordLength w\neq0\vdash w\in\NonemptyWords A
  }{w\in\FiniteWords A\land\WordLength w\neq0\vdash w\in\NonemptyWords A}%
    [NonemptyWordMembershipIntroduction]
    \proofstep{1}{w\in\FiniteWords A\land\WordLength w\neq0}{\rA}
    \proofstep{1}{w\in\FiniteWords A}{\rAEa{1}}
    \proofstep{1}{w\neq\EmptyWord{A}\leftrightarrow\WordLength{w}\neq0}{%
      \FormulaRefAuto{FiniteWordLengthNonzeroIffNonempty}[thm]{2}}
    \proofstep{1}{w\neq\EmptyWord A}{%
      \rRE{\rAEb{1},\rLREb{3}}}
    \proofstep{1}{w\notin\{\EmptyWord A\}}{%
      \FormulaRefAuto{x\notin\{a\}\eqvdash x\neq a}[thm]{4}}
    \proofstep{1}{w\in\FiniteWords A\setminus\{\EmptyWord A\}}{%
      \FormulaRefAuto{x\in A\setminus B\eqvdash x\in A\land x\notin B}[thm]{\rAI{2,5}}}
    \proofstep{1}{w\in\NonemptyWords A}{%
      \FormulaRefAuto{NonemptyWordSetDef}[def]{6}}
  \closeproofpartindR
\end{tabproofsplit}

\Needspace{19\baselineskip}
\FormulaThmDeltaK[Typisierung nichtleerer Wörter]{%
  w\in\NonemptyWords{A}\vdash
  \begin{aligned}[t]
    &\text{(i)}\;w\in\FiniteWords{A}\\[-2pt]
    &\text{(ii)}\;\WordLength{w}\in\mathbb N\\[-2pt]
    &\text{(iii)}\;w\colon\NatSeg{\WordLength{w}}\to A
  \end{aligned}%
}{NonemptyWordTyping}{%
  \DeltaRow{Mengen}{A\dsep w}
}
\Needspace{11\baselineskip}
\begin{tabproofsplitwide}
  \proofpartwideindR[Endlichkeit]{%
    w\in\NonemptyWords{A}\vdash w\in\FiniteWords{A}
  }{w\in\NonemptyWords{A}\vdash w\in\FiniteWords{A}}%
    [NonemptyWordFinite]
    \proofstepwidestar[1]{w\in\NonemptyWords{A}}{\rA}
    \proofstepwidestar[1]{%
      w\in\FiniteWords{A}\land\WordLength w\neq0}{%
      \FormulaRefAuto{NonemptyWordMembership}[thm]{1}}
    \proofstepwidestar[1]{w\in\FiniteWords{A}}{\rAEa{2}}
  \closeproofpartwideindR

  \Needspace{11\baselineskip}
  \proofpartwideindR[Natürlichkeit der Länge]{%
    w\in\NonemptyWords{A}\vdash\WordLength{w}\in\mathbb N
  }{w\in\NonemptyWords{A}\vdash\WordLength{w}\in\mathbb N}%
    [NonemptyWordLengthNatural]
    \proofstepwidestar[1]{w\in\NonemptyWords{A}}{\rA}
    \proofstepwidestar[1]{%
      w\in\FiniteWords{A}\land\WordLength w\neq0}{%
      \FormulaRefAuto{NonemptyWordMembership}[thm]{1}}
    \proofstepwidestar[1]{w\in\FiniteWords{A}}{\rAEa{2}}
    \proofstepwidestar[1]{\WordLength{w}\in\mathbb N}{%
      \FormulaRefAuto{FiniteWordLengthNatural}[thm]{3}}
  \closeproofpartwideindR

  \Needspace{11\baselineskip}
  \proofpartwideindR[Worttypisierung]{%
    w\in\NonemptyWords{A}\vdash
      w\colon\NatSeg{\WordLength{w}}\to A
  }{w\in\NonemptyWords{A}\vdash
      w\colon\NatSeg{\WordLength{w}}\to A}%
    [NonemptyWordTypingAtLength]
    \proofstepwidestar[1]{w\in\NonemptyWords{A}}{\rA}
    \proofstepwidestar[1]{%
      w\in\FiniteWords{A}\land\WordLength w\neq0}{%
      \FormulaRefAuto{NonemptyWordMembership}[thm]{1}}
    \proofstepwidestar[1]{w\in\FiniteWords{A}}{\rAEa{2}}
    \proofstepwidestar[1]{%
      w\colon\NatSeg{\WordLength{w}}\to A}{%
      \FormulaRefAuto{FiniteWordTypingAtLength}[thm]{3}}
  \closeproofpartwideindR
\end{tabproofsplitwide}

\Needspace{9\baselineskip}
\FormulaDefDeltaK[Buchstabenwort]{%
  a\in A\vdash
  \LetterWord{a}\coloneqq\{(0,a)\}%
}{LetterWordDef}{%
  \DeltaRow{Mengen}{A\dsep a\dsep\LetterWord{a}}
  \DeltaRow{\textbf{Neues Symbol}}{}
  \DeltaPrem{Buchstabenwort}{\LetterWord{a}}
}


\Needspace{20\baselineskip}
\FormulaThmDeltaK[Grundgesetze der Buchstabenwörter]{%
  a\in A\vdash
  \begin{aligned}[t]
    &\text{(i)}\;\LetterWord a\colon\NatSeg 1\to A\\[-2pt]
    &\text{(ii)}\;\WordLength{\LetterWord a}=1\\[-2pt]
    &\text{(iii)}\;\LetterWord a\in\NonemptyWords A\\[-2pt]
    &\text{(iv)}\;\LetterWord a(0)=a
  \end{aligned}%
}{LetterWordFacts}{%
  \DeltaRow{Mengen}{A\dsep a}
}
\Needspace{11\baselineskip}
\begin{tabproofsplitwide}
  \proofpartwideindR[Typisierung]{%
    a\in A\vdash\LetterWord a\colon\NatSeg 1\to A
  }{a\in A\vdash\LetterWord a\colon\NatSeg 1\to A}[LetterWordTyping]
    \proofstepwidestar[1]{a\in A}{\rA}
    \proofstepwidestar[1]{\{(0,a)\}\colon\{0\}\to A}{%
      \FormulaRefAuto{SingletonGraphFunction}[thm]{1}}
    \proofstepwidestar[]{\NatSeg 1=\{0\}}{%
      \FormulaRefAuto{PeanoNatSegOneEquality}[thm]}
    \proofstepwidestar[1]{\LetterWord a=\{(0,a)\}}{%
      \FormulaRefAuto{LetterWordDef}[def]{1}}
    \proofstepwidestar[1]{\LetterWord a\colon\NatSeg 1\to A}{\rIE{3,4,2}}
  \closeproofpartwideindR
  \Needspace{11\baselineskip}
  \proofpartwideindR[Länge eins]{%
    a\in A\vdash\WordLength{\LetterWord a}=1
  }{a\in A\vdash\WordLength{\LetterWord a}=1}[LetterWordLengthOne]
    \proofstepwidestar[1]{a\in A}{\rA}
    \proofstepwidestar[1]{\LetterWord a\colon\NatSeg 1\to A}{%
      \FormulaRefAuto{LetterWordTyping}[thm]{1}}
    \proofstepwidestar[]{1\in\mathbb N}{\FormulaRefAuto{PeanoOneInN}[thm]}
    \proofstepwidestar[1]{\LetterWord a\in\FiniteWords A}{%
      \FormulaRefAuto{FiniteWordFromTyping}[thm]{3,2}}
    \proofstepwidestar[1]{\WordLength{\LetterWord a}=1}{%
      \FormulaRefAuto{FiniteWordLengthFromTyping}[thm]{4,3,2}}
  \closeproofpartwideindR
  \Needspace{11\baselineskip}
  \proofpartwideindR[Nichtleeres Wort]{%
    a\in A\vdash\LetterWord a\in\NonemptyWords A
  }{a\in A\vdash\LetterWord a\in\NonemptyWords A}[LetterWordNonempty]
    \proofstepwidestar[1]{a\in A}{\rA}
    \proofstepwidestar[1]{\LetterWord a\colon\NatSeg 1\to A}{%
      \FormulaRefAuto{LetterWordTyping}[thm]{1}}
    \proofstepwidestar[]{1\in\mathbb N}{\FormulaRefAuto{PeanoOneInN}[thm]}
    \proofstepwidestar[1]{\LetterWord a\in\FiniteWords A}{%
      \FormulaRefAuto{FiniteWordFromTyping}[thm]{3,2}}
    \proofstepwidestar[1]{\WordLength{\LetterWord a}=1}{%
      \FormulaRefAuto{LetterWordLengthOne}[thm]{1}}
    \proofstepwidestar[]{1\neq0}{\FormulaRefAuto{PeanoOneNeZero}[thm]}
    \proofstepwidestar[1]{\WordLength{\LetterWord a}\neq0}{\rIE{5,6}}
    \proofstepwidestar[1]{\LetterWord a\in\NonemptyWords A}{%
      \FormulaRefAuto{NonemptyWordMembershipIntroduction}[thm]{\rAI{4,7}}}
  \closeproofpartwideindR
  \Needspace{11\baselineskip}
  \proofpartwideindR[Auswertung an null]{%
    a\in A\vdash\LetterWord a(0)=a
  }{a\in A\vdash\LetterWord a(0)=a}[LetterWordValueZero]
    \proofstepwidestar[1]{a\in A}{\rA}
    \proofstepwidestar[1]{\{(0,a)\}(0)=a}{%
      \FormulaRefAuto{SingletonGraphValue}[thm]{1}}
    \proofstepwidestar[1]{\LetterWord a=\{(0,a)\}}{%
      \FormulaRefAuto{LetterWordDef}[def]{1}}
    \proofstepwidestar[1]{\LetterWord a(0)=a}{\rIE{3,2}}
  \closeproofpartwideindR
\end{tabproofsplitwide}



% Früh einfügen: nach Leerwort/Buchstabenwörtern, vor Konkatenation.
% Benötigt nur die Wortmenge, ihre Länge und die Bände 03--20.
\section{Die primitive Anfügung und ihre Grundgesetze}

Bevor zwei beliebige Wörter miteinander verknüpft werden, betrachten wir
die bereits durch \(J(u,a)\) beschriebene Anfügung eines einzelnen
Buchstabens. Für die anschließende Axiomatik benötigen wir diesen
Erzeugungsschritt als Funktion auf der Wortmenge. Wir fassen ihn deshalb
zu einer Abbildung \(\sigma_A:\FiniteWords A\times A\to\FiniteWords A\)
zusammen und weisen ihre Grundgesetze nach. Sie übernimmt im konkreten
Wortmodell die Rolle der späteren abstrakten Anfügungsfunktion \(s\).
Die Minimalität ist durch die Anfangskonstruktion bereits gesichert.

\Needspace{18\baselineskip}
\noindent\begin{minipage}{\linewidth}
\FormulaThmDeltaK[Die primitive Anfügung bleibt in der Wortmenge]{%
  u\in\FiniteWords A\dsep a\in A\vdash
    u\cup\{(\WordLength u,a)\}\in\FiniteWords A
}{EarlyWordAdjunctionClosed}{\DeltaRow{Mengen}{A\dsep u\dsep a}}
\end{minipage}\par
\begin{tabproofwide}
  \proofstepwidestar[1]{u\in\FiniteWords A}{\rA}
  \proofstepwidestar[2]{a\in A}{\rA}
  \proofstepwidestar[1]{u\in\FiniteOccurrenceSets A}{\FormulaRefAuto{FiniteWordInUniverse}[thm]{1}}
  \proofstepwidestar[1,2]{\OccurrenceAdjunction u a\in\FiniteWords A}{\FormulaRefAuto{GeneratedWordClosed}[thm]{1,2}}
  \proofstepwidestar[1,2]{\OccurrenceAdjunction u a=u\cup\{(\FiniteCard(u),a)\}}{\FormulaRefAuto{OccurrenceCardinalAdjunctionDef}[def]{3,2}}
  \proofstepwidestar[1]{\WordLength u=\FiniteCard(u)}{\FormulaRefAuto{FiniteWordLengthCardinality}[thm]{1}}
  \proofstepwidestar[1,2]{u\cup\{(\WordLength u,a)\}\in\FiniteWords A}{\rIE{5,6,4}}
\end{tabproofwide}

\noindent\begin{minipage}{\linewidth}
\FormulaDefDeltaK[Anfügungsfunktion des Wortmodells]{%
  \begin{aligned}[t]
    &\sigma_A\coloneqq\bigl\{(p,z)\in
      (\FiniteWords A\times A)\times\FiniteWords A\bigm|{}\\[-2pt]
    &\qquad z=\pi_1(p)\cup\{(\WordLength{\pi_1(p)},\pi_2(p))\}\bigr\}
  \end{aligned}
}{WordModelSuccessorDef}{%
  \DeltaRow{Mengen}{A\dsep p\dsep z}
  \DeltaRow{\textbf{Neues Symbol}}{}
  \DeltaPrem{Anfügungsfunktion}{\sigma_A}
}
\end{minipage}\par
Die Projektionen beziehen sich hier auf \(\FiniteWords A\times A\).
Wir schreiben \(\sigma_A(u,a)\) für den Wert am geordneten Paar \((u,a)\).

\noindent\begin{minipage}{\linewidth}
\FormulaThmDeltaK[W0: Typisierung des konkreten Wortmodells]{%
  \EmptyWord A\in\FiniteWords A\land
    \sigma_A\colon\FiniteWords A\times A\to\FiniteWords A
}{EarlyWordSuccessorTyping}{\DeltaRow{Mengen}{A\dsep p\dsep u\dsep a}}
\end{minipage}\par
\begin{tabproofwide}
  \proofstepwidestar[]{\EmptyWord A\in\FiniteWords A}{\FormulaRefAuto{EmptyWordFinite}[thm]}
  \proofstepwidestar[2]{p\in\FiniteWords A\times A}{\rA}
  \proofstepwidestar[2]{\pi_1(p)\in\FiniteWords A\land\pi_2(p)\in A}{%
    \rAI{\FormulaRefAuto{F\colon A\to B\dsep x\in A\vdash F(x)\in B}[thm]{\FormulaRefAuto{FirstProjectionDef}[def],2},\FormulaRefAuto{F\colon A\to B\dsep x\in A\vdash F(x)\in B}[thm]{\FormulaRefAuto{SecondProjectionDef}[def],2}}}
  \proofstepwidestar[2]{\pi_1(p)\cup\{(\WordLength{\pi_1(p)},\pi_2(p))\}\in\FiniteWords A}{\FormulaRefAuto{EarlyWordAdjunctionClosed}[thm]{\rAEa{3},\rAEb{3}}}
  \proofstepwidestar[]{\begin{aligned}[t]
    &\forall p\in\FiniteWords A\times A\\[-2pt]
    &\quad\pi_1(p)\cup\{(\WordLength{\pi_1(p)},\pi_2(p))\}\in\FiniteWords A
  \end{aligned}}{\rUI{\rRI{2,4}}}
  \proofstepwidestar[]{\sigma_A\colon\FiniteWords A\times A\to\FiniteWords A}{%
    \rIE{\FormulaRefAuto{WordModelSuccessorDef}[def],\FormulaRefAuto{TypedTermGraphFunction}[thm]{5}}}
  \proofstepwidestar[]{\EmptyWord A\in\FiniteWords A\land
    \sigma_A\colon\FiniteWords A\times A\to\FiniteWords A}{\rAI{1,6}}
\end{tabproofwide}

\Needspace{16\baselineskip}
\noindent\begin{minipage}{\linewidth}
\FormulaThmDeltaK[Grundgleichung der primitiven Anfügung]{%
  u\in\FiniteWords A\dsep a\in A\vdash
    \sigma_A(u,a)=u\cup\{(\WordLength u,a)\}
}{EarlyWordSuccessorEquation}{\DeltaRow{Mengen}{A\dsep u\dsep a}}
\end{minipage}\par
\begin{tabproofwide}
  \proofstepwidestar[1]{u\in\FiniteWords A}{\rA}
  \proofstepwidestar[2]{a\in A}{\rA}
  \proofstepwidestar[1,2]{(u,a)\in\FiniteWords A\times A}{\FormulaRefAuto{(a,b)\in A\times B\eqvdash a\in A\land b\in B}[thm]{\rAI{1,2}}}
  \proofstepwidestar[1,2]{u\cup\{(\WordLength u,a)\}\in\FiniteWords A}{\FormulaRefAuto{EarlyWordAdjunctionClosed}[thm]{1,2}}
  \proofstepwidestar[1,2]{\begin{aligned}[t]
    &((u,a),u\cup\{(\WordLength u,a)\})\\[-2pt]
    &\quad\in(\FiniteWords A\times A)\times\FiniteWords A
  \end{aligned}}{\FormulaRefAuto{(a,b)\in A\times B\eqvdash a\in A\land b\in B}[thm]{\rAI{3,4}}}
  \proofstepwidestar[1,2]{\pi_1(u,a)=u}{\FormulaRefAuto{(x,y)\in A\times B\vdash\pi_1(x,y)=x}[thm]{3}}
  \proofstepwidestar[1,2]{\pi_2(u,a)=a}{\FormulaRefAuto{(x,y)\in A\times B\vdash\pi_2(x,y)=y}[thm]{3}}
  \proofstepwidestar[1,2]{\begin{aligned}[t]
    &u\cup\{(\WordLength u,a)\}=\pi_1(u,a)\cup{}\\[-2pt]
    &\qquad\{(\WordLength{\pi_1(u,a)},\pi_2(u,a))\}
  \end{aligned}}{\rIE{6,7,\rII}}
  \proofstepwidestar[1,2]{((u,a),u\cup\{(\WordLength u,a)\})\in\sigma_A}{%
    \rIE{\FormulaRefAuto{WordModelSuccessorDef}[def],\FormulaRefAuto{x\in\{u\in A\mid P(u)\}\eqvdash x\in A\land P(x)}[thm]{\rAI{5,8}}}}
  \proofstepwidestar[]{\sigma_A\colon\FiniteWords A\times A\to\FiniteWords A}{\rAEb{\FormulaRefAuto{EarlyWordSuccessorTyping}[thm]}}
  \proofstepwidestar[1,2]{\sigma_A(u,a)=u\cup\{(\WordLength u,a)\}}{\FormulaRefAuto{F\colon A\to B\dsep(x,y)\in F\vdash F(x)=y}[thm]{10,9}}
\end{tabproofwide}

\noindent\begin{minipage}{\linewidth}
\FormulaThmDeltaK[Die Anfügungsfunktion führt die Erzeugungsregel aus]{%
  u\in\FiniteWords A\dsep a\in A\vdash
    \sigma_A(u,a)=\OccurrenceAdjunction u a
}{EarlyWordSuccessorGenerationEquation}{\DeltaRow{Mengen}{A\dsep u\dsep a}}
\begin{tabproofwide}
  \proofstepwidestar[1]{u\in\FiniteWords A}{\rA}
  \proofstepwidestar[2]{a\in A}{\rA}
  \proofstepwidestar[1]{u\in\FiniteOccurrenceSets A}{\FormulaRefAuto{FiniteWordInUniverse}[thm]{1}}
  \proofstepwidestar[1,2]{\sigma_A(u,a)=u\cup\{(\WordLength u,a)\}}{\FormulaRefAuto{EarlyWordSuccessorEquation}[thm]{1,2}}
  \proofstepwidestar[1]{\WordLength u=\FiniteCard(u)}{\FormulaRefAuto{FiniteWordLengthCardinality}[thm]{1}}
  \proofstepwidestar[1,2]{\OccurrenceAdjunction u a=u\cup\{(\FiniteCard(u),a)\}}{\FormulaRefAuto{OccurrenceCardinalAdjunctionDef}[def]{3,2}}
  \proofstepwidestar[1,2]{\sigma_A(u,a)=\OccurrenceAdjunction u a}{\rIE{5,6,4}}
\end{tabproofwide}

\end{minipage}\par
\Needspace{16\baselineskip}
\noindent\begin{minipage}{\linewidth}
\FormulaThmDeltaK[Wortzugehörigkeit und Länge eines primitiven Nachfolgers]{%
  u\in\FiniteWords A\dsep a\in A\vdash
  \begin{aligned}[t]
    &\sigma_A(u,a)\in\FiniteWords A\land{}\\[-2pt]
    &\WordLength{\sigma_A(u,a)}=\WordLength u+1
  \end{aligned}
}{EarlyWordSuccessorLength}{\DeltaRow{Mengen}{A\dsep u\dsep a}}
\end{minipage}\par
\begin{tabproofwide}
  \proofstepwidestar[1]{u\in\FiniteWords A}{\rA}
  \proofstepwidestar[2]{a\in A}{\rA}
  \proofstepwidestar[1]{\WordLength u\in\mathbb N}{\FormulaRefAuto{FiniteWordLengthNatural}[thm]{1}}
  \proofstepwidestar[1]{u\colon\NatSeg{\WordLength u}\to A}{\FormulaRefAuto{FiniteWordTypingAtLength}[thm]{1}}
  \proofstepwidestar[1,2]{\sigma_A(u,a)=u\cup\{(\WordLength u,a)\}}{\FormulaRefAuto{EarlyWordSuccessorEquation}[thm]{1,2}}
  \proofstepwidestar[1,2]{u\cup\{(\WordLength u,a)\}\colon\NatSeg{\WordLength u+1}\to A}{\FormulaRefAuto{EarlyWordAdjunctionGraphTyping}[thm]{3,4,2}}
  \proofstepwidestar[1,2]{\sigma_A(u,a)\colon\NatSeg{\WordLength u+1}\to A}{\rIE{5,6}}
  \proofstepwidestar[1]{\WordLength u+1\in\mathbb N}{\FormulaRefAuto{PeanoAddOneClosure}[thm]{3}}
  \proofstepwidestar[1,2]{\sigma_A(u,a)\in\FiniteWords A}{\FormulaRefAuto{FiniteWordFromTyping}[thm]{8,7}}
  \proofstepwidestar[1,2]{\WordLength{\sigma_A(u,a)}=\WordLength u+1}{\FormulaRefAuto{FiniteWordLengthFromTyping}[thm]{9,8,7}}
  \proofstepwidestar[1,2]{\sigma_A(u,a)\in\FiniteWords A\land\WordLength{\sigma_A(u,a)}=\WordLength u+1}{\rAI{9,10}}
\end{tabproofwide}

\noindent\begin{minipage}{\linewidth}
\FormulaThmDeltaK[W1: Das Leerwort ist kein primitiver Nachfolger]{%
  u\in\FiniteWords A\dsep a\in A\vdash\sigma_A(u,a)\ne\EmptyWord A
}{EarlyWordSuccessorNonempty}{\DeltaRow{Mengen}{A\dsep u\dsep a}}
\end{minipage}\par
\begin{tabproofwide}
  \proofstepwidestar[1]{u\in\FiniteWords A}{\rA}
  \proofstepwidestar[2]{a\in A}{\rA}
  \proofstepwidestar[1,2]{\sigma_A(u,a)\in\FiniteWords A\land\WordLength{\sigma_A(u,a)}=\WordLength u+1}{\FormulaRefAuto{EarlyWordSuccessorLength}[thm]{1,2}}
  \proofstepwidestar[1]{\WordLength u\in\mathbb N}{\FormulaRefAuto{FiniteWordLengthNatural}[thm]{1}}
  \proofstepwidestar[1]{\WordLength u+1\ne0}{\FormulaRefAuto{PeanoAddOneNonzero}[thm]{4}}
  \proofstepwidestar[1,2]{\WordLength{\sigma_A(u,a)}\ne0}{\rIE{\rAEb{3},5}}
  \proofstepwidestar[1,2]{\sigma_A(u,a)\ne\EmptyWord A}{%
    \FormulaRefAuto{FiniteWordLengthNonzeroIffNonempty}[thm]{\rAEa{3},6}}
\end{tabproofwide}

\noindent\begin{minipage}{\linewidth}
\FormulaThmDeltaK[Werte vor und an der neuen letzten Position]{%
  u\in\FiniteWords A\dsep a\in A\vdash
  \begin{aligned}[t]
    &(\forall i\in\NatSeg{\WordLength u}\,
      \sigma_A(u,a)(i)=u(i))\land{}\\[-2pt]
    &\sigma_A(u,a)(\WordLength u)=a
  \end{aligned}
}{EarlyWordSuccessorValues}{\DeltaRow{Mengen}{A\dsep u\dsep a\dsep i}}
\end{minipage}\par
\begin{tabproofwide}
  \proofstepwidestar[1]{u\in\FiniteWords A}{\rA}
  \proofstepwidestar[2]{a\in A}{\rA}
  \proofstepwidestar[1,2]{\sigma_A(u,a)\in\FiniteWords A}{\rAEa{\FormulaRefAuto{EarlyWordSuccessorLength}[thm]{1,2}}}
  \proofstepwidestar[1,2]{\sigma_A(u,a)\colon\NatSeg{\WordLength{\sigma_A(u,a)}}\to A}{\FormulaRefAuto{FiniteWordTypingAtLength}[thm]{3}}
  \proofstepwidestar[1,2]{\sigma_A(u,a)=u\cup\{(\WordLength u,a)\}}{\FormulaRefAuto{EarlyWordSuccessorEquation}[thm]{1,2}}
  \proofstepwidestar[6]{i\in\NatSeg{\WordLength u}}{\rA}
  \proofstepwidestar[1,6]{(i,u(i))\in u}{%
    \FormulaRefAuto{F\colon A\to B\dsep x\in A\vdash(x,F(x))\in F}[thm]{\FormulaRefAuto{FiniteWordTypingAtLength}[thm]{1},6}}
  \proofstepwidestar[1,2,6]{(i,u(i))\in\sigma_A(u,a)}{\rIE{5,\FormulaRefAuto{z\in A\vdash z\in A\cup B}[thm]{7}}}
  \proofstepwidestar[1,2,6]{\sigma_A(u,a)(i)=u(i)}{\FormulaRefAuto{F\colon A\to B\dsep(x,y)\in F\vdash F(x)=y}[thm]{4,8}}
  \proofstepwidestar[1,2]{\forall i\in\NatSeg{\WordLength u}\,\sigma_A(u,a)(i)=u(i)}{\rUI{\rRI{6,9}}}
  \proofstepwidestar[1,2]{(\WordLength u,a)\in\sigma_A(u,a)}{\rIE{5,\FormulaRefAuto{z\in B\vdash z\in A\cup B}[thm]{\FormulaRefAuto{a\in\{a\}}[thm]}}}
  \proofstepwidestar[1,2]{\sigma_A(u,a)(\WordLength u)=a}{\FormulaRefAuto{F\colon A\to B\dsep(x,y)\in F\vdash F(x)=y}[thm]{4,11}}
  \proofstepwidestar[1,2]{\begin{aligned}[t]
    &(\forall i\in\NatSeg{\WordLength u}\,\sigma_A(u,a)(i)=u(i))\land{}\\[-2pt]
    &\sigma_A(u,a)(\WordLength u)=a
  \end{aligned}}{\rAI{10,12}}
\end{tabproofwide}

\noindent\begin{minipage}{\linewidth}
\FormulaThmDeltaK[W2: Gemeinsame Injektivität der primitiven Anfügung]{%
  u,v\in\FiniteWords A\dsep a,b\in A\dsep
    \sigma_A(u,a)=\sigma_A(v,b)\vdash u=v\land a=b
}{EarlyWordSuccessorInjective}{\DeltaRow{Mengen}{A\dsep u\dsep v\dsep a\dsep b\dsep i}}
\end{minipage}\par
\begin{tabproofwide}
  \proofstepwidestar[1]{u,v\in\FiniteWords A}{\rA}
  \proofstepwidestar[2]{a,b\in A}{\rA}
  \proofstepwidestar[3]{\sigma_A(u,a)=\sigma_A(v,b)}{\rA}
  \proofstepwidestar[1]{\WordLength u,\WordLength v\in\mathbb N}{\rAI{\FormulaRefAuto{FiniteWordLengthNatural}[thm]{\rAEa{1}},\FormulaRefAuto{FiniteWordLengthNatural}[thm]{\rAEb{1}}}}
  \proofstepwidestar[1,2]{\WordLength{\sigma_A(u,a)}=\WordLength u+1}{\rAEb{\FormulaRefAuto{EarlyWordSuccessorLength}[thm]{\rAEa{1},\rAEa{2}}}}
  \proofstepwidestar[1,2]{\WordLength{\sigma_A(v,b)}=\WordLength v+1}{\rAEb{\FormulaRefAuto{EarlyWordSuccessorLength}[thm]{\rAEb{1},\rAEb{2}}}}
  \proofstepwidestar[1,2,3]{\WordLength u+1=\WordLength v+1}{\rIE{3,5,6}}
  \proofstepwidestar[1,2,3]{\WordLength u=\WordLength v}{\FormulaRefAuto{PeanoAddOneInjective}[thm]{\rAEa{4},\rAEb{4},7}}
  \proofstepwidestar[1,2]{\sigma_A(u,a)(\WordLength u)=a}{\rAEb{\FormulaRefAuto{EarlyWordSuccessorValues}[thm]{\rAEa{1},\rAEa{2}}}}
  \proofstepwidestar[1,2]{\sigma_A(v,b)(\WordLength v)=b}{\rAEb{\FormulaRefAuto{EarlyWordSuccessorValues}[thm]{\rAEb{1},\rAEb{2}}}}
  \proofstepwidestar[1,2,3]{a=b}{\rIE{3,8,9,10}}
  \proofstepwidestar[12]{i\in\NatSeg{\WordLength u}}{\rA}
  \proofstepwidestar[1,2,12]{\sigma_A(u,a)(i)=u(i)}{\rRE{12,\rUE{\rAEa{\FormulaRefAuto{EarlyWordSuccessorValues}[thm]{\rAEa{1},\rAEa{2}}}}}}
  \proofstepwidestar[1,2,3,12]{i\in\NatSeg{\WordLength v}}{\rIE{8,12}}
  \proofstepwidestar[1,2,3,12]{\sigma_A(v,b)(i)=v(i)}{\rRE{14,\rUE{\rAEa{\FormulaRefAuto{EarlyWordSuccessorValues}[thm]{\rAEb{1},\rAEb{2}}}}}}
  \proofstepwidestar[1,2,3,12]{u(i)=v(i)}{\rIE{3,13,15}}
  \proofstepwidestar[1,2,3]{\forall i\in\NatSeg{\WordLength u}\,u(i)=v(i)}{\rUI{\rRI{12,16}}}
  \proofstepwidestar[1,2,3]{u\colon\NatSeg{\WordLength u}\to A\land v\colon\NatSeg{\WordLength u}\to A}{%
    \rAI{\FormulaRefAuto{FiniteWordTypingAtLength}[thm]{\rAEa{1}},\rIE{8,\FormulaRefAuto{FiniteWordTypingAtLength}[thm]{\rAEb{1}}}}}
  \proofstepwidestar[1,2,3]{u=v}{\FormulaRefAuto{FunctionExtensionality}[thm]{\rAEa{18},\rAEb{18},17}}
  \proofstepwidestar[1,2,3]{u=v\land a=b}{\rAI{19,11}}
\end{tabproofwide}

\subsection{Die Minimalität des Anfügungsmodells}

Die Wortmenge wurde als kleinste unter der Erzeugungsregel abgeschlossene
Familie konstruiert. Nach
\FormulaRefAuto{EarlyWordSuccessorGenerationEquation}[thm] führt
\(\sigma_A\) genau diese Regel aus. Daher gilt die bereits bewiesene
Minimalität auch in der Sprache der Anfügungsfunktion.

\noindent\begin{minipage}{\linewidth}
\FormulaThmDeltaK[W3: Minimalität des konkreten Wortmodells]{%
  \begin{aligned}[t]
    &U\subseteq\FiniteWords A\dsep\EmptyWord A\in U\dsep{}\\[-2pt]
    &\forall u\in U\,\forall a\in A\,(\sigma_A(u,a)\in U)
      \vdash U=\FiniteWords A
  \end{aligned}
}{EarlyWordModelMinimality}{%
  \DeltaRow{Mengen}{A\dsep U\dsep u\dsep a}
}
\end{minipage}\par
\begin{tabproofwide}
  \proofstepwidestar[1]{U\subseteq\FiniteWords A}{\rA}
  \proofstepwidestar[2]{\EmptyWord A\in U}{\rA}
  \proofstepwidestar[3]{\forall u\in U\,\forall a\in A\,(\sigma_A(u,a)\in U)}{\rA}
  \proofstepwidestar[]{\FiniteWords A\subseteq\FiniteOccurrenceSets A}{\FormulaRefAuto{FiniteWordSetSubsetUniverse}[thm]}
  \proofstepwidestar[1]{U\subseteq\FiniteOccurrenceSets A}{\FormulaRefAuto{A\subseteq B\dsep B\subseteq C\vdash A\subseteq C}[thm]{1,4}}
  \proofstepwidestar[]{\EmptyWord A=\varnothing}{\FormulaRefAuto{EmptyWordDef}[def]}
  \proofstepwidestar[2]{\varnothing\in U}{\rIE{6,2}}
  \proofstepwidestar[8]{u\in U}{\rA}
  \proofstepwidestar[9]{a\in A}{\rA}
  \proofstepwidestar[1,8]{u\in\FiniteWords A}{\FormulaRefAuto{A\subseteq B\dsep x\in A\vdash x\in B}[thm]{1,8}}
  \proofstepwidestar[3,8]{\forall a\in A\,(\sigma_A(u,a)\in U)}{\rRE{8,\rUE{3}}}
  \proofstepwidestar[3,8,9]{\sigma_A(u,a)\in U}{\rRE{9,\rUE{11}}}
  \proofstepwidestar[1,8,9]{\sigma_A(u,a)=\OccurrenceAdjunction u a}{\FormulaRefAuto{EarlyWordSuccessorGenerationEquation}[thm]{10,9}}
  \proofstepwidestar[1,3,8,9]{\OccurrenceAdjunction u a\in U}{\rIE{13,12}}
  \proofstepwidestar[1,3]{\forall u\in U\,\forall a\in A\,(\OccurrenceAdjunction u a\in U)}{\rUI{\rRI{8,\rUI{\rRI{9,14}}}}}
  \proofstepwidestar[1,2,3]{\FiniteWords A\subseteq U}{\FormulaRefAuto{GeneratedWordMinimality}[thm]{5,7,15}}
  \proofstepwidestar[1,2,3]{U=\FiniteWords A}{\FormulaRefAuto{A\subseteq B\dsep B\subseteq A\vdash A=B}[thm]{1,16}}
\end{tabproofwide}

Damit sind die vier konkreten Grundgesetze bewiesen. Wir halten nun
genau diese Gesetze als Axiome einer abstrakten Wortstruktur fest
und leiten daraus die allgemeine Induktion und Rekursion her.

\section{Das Axiomsystem der Wortstrukturen}

Die Axiomatik beschreibt, welche Eigenschaften eine Menge von Wörtern
mit einem Anfangswort und einer Anfügungsfunktion besitzt. Dabei haben
die Symbole folgende Bedeutung:

\begingroup
\renewcommand{\arraystretch}{1.15}
\noindent\begin{tabular}{@{}p{.14\linewidth}p{.81\linewidth}@{}}
\textbf{Symbol} & \textbf{Bedeutung} \\
\(A\) & Das Alphabet, also die Menge der zugelassenen Buchstaben. \\
\(W\) & Die Trägermenge der Struktur; ihre Elemente werden als ganze Wörter aufgefasst. \\
\(e\) & Das ausgezeichnete Anfangswort in \(W\); es übernimmt die Rolle des leeren Wortes. \\
\(s\) & Die Anfügungsfunktion: \(s(u,a)\) entsteht aus dem Wort \(u\), indem rechts der Buchstabe \(a\) angefügt wird. \\
\(u,v\) & Beliebige Wörter, also Elemente von \(W\). \\
\(a,b\) & Beliebige Buchstaben, also Elemente von \(A\). \\
\(U\) & Eine beliebige Teilmenge von \(W\), an der die Minimalitätsbedingung geprüft wird.
\end{tabular}\par
\endgroup

Da \(s\) auf \(W\times A\) definiert ist, bedeutet \(s(u,a)\) den
Funktionswert am geordneten Paar \((u,a)\). Die Funktion hat somit zwei
Eingaben: ein bisheriges Wort und einen einzelnen Buchstaben.
Beispielsweise stellt \(s(e,a)\) das aus \(a\) bestehende Wort dar;
\(s(s(e,a),b)\) stellt das Wort mit den aufeinanderfolgenden Buchstaben
\(a,b\) dar. Die Buchstaben dürfen dabei gleich sein.

Im zuvor aufgebauten Anfügungsmodell werden die Symbole durch
\[
 W=\FiniteWords A,\qquad e=\EmptyWord A,\qquad s=\sigma_A
\]
belegt. Nach \FormulaRefAuto{EarlyWordSuccessorEquation}[thm] ist dort
für \(u\in\FiniteWords A\) und \(a\in A\)
\[
 \sigma_A(u,a)=u\cup\{(\WordLength u,a)\}.
\]
Der neue Buchstabe erhält also die nächste freie Position. In einer
beliebigen Wortstruktur sind \(W,e,s\) dagegen zunächst die gewählten
Strukturdaten; ihre konkrete Darstellung ist durch die Axiomatik nicht
festgelegt. Insbesondere muss \(e\) dort nicht die leere Menge sein.

Für eine Menge \(A\) heißt \((W,e,s)\) eine \emph{freie Wortstruktur über
\(A\)}, wenn die folgenden Bedingungen gelten.

\Needspace{19\baselineskip}
\FormulaDefDeltaK[Freie Wortstruktur]{%
  \mathsf{WortStr}_A(W,e,s)
}{WordStructureDef}{%
  \DeltaRow{Mengen}{A\dsep W\dsep e\dsep U\dsep u\dsep v\dsep a\dsep b}
  \DeltaRow{Funktionen}{s}
  \DeltaRow{\textbf{Strukturaxiome}}{}
  \DeltaRow{W0: Typisierung}{e\in W\land s\colon W\times A\to W}
  \DeltaRow{W1: Anfang}{\forall u\in W\,\forall a\in A\;s(u,a)\ne e}
  \DeltaRow{W2: Eindeutigkeit}{%
    \begin{aligned}[t]
      &\forall u,v\in W\,\forall a,b\in A\\[-2pt]
      &\quad(s(u,a)=s(v,b)\rightarrow u=v\land a=b)
    \end{aligned}}
  \DeltaRow{W3: Minimalität}{%
    \begin{aligned}[t]
      &\forall U\subseteq W\,\Bigl(\bigl(e\in U\land{}\\[-2pt]
      &\quad(\forall u\in U\,\forall a\in A\;s(u,a)\in U)\bigr)
        \rightarrow U=W\Bigr)
    \end{aligned}}
  \DeltaRow{\textbf{Neues Symbol}}{}
  \DeltaPrem{Freie Wortstruktur}{\mathsf{WortStr}_A(W,e,s)}
}

Die bei den einzelnen Axiomen genannten Herkunftssätze beweisen die
jeweilige Eigenschaft für das Anfügungsmodell. Für beliebige Daten
\((W,e,s)\) werden diese Eigenschaften in
\FormulaRefAuto{WordStructureDef}[def] gefordert. Unter der Annahme
\(\mathsf{WortStr}_A(W,e,s)\) geben die folgenden nummerierten Axiome
jeweils die entsprechende Bedingung dieser Definition wieder.

\noindent\begin{minipage}{\linewidth}
\FormulaAxiomDeltaK[Typisierung einer Wortstruktur]{%
  \mathsf{WortStr}_A(W,e,s)\vdash e\in W\land s\colon W\times A\to W
}{WordStructureTypingAxiom}{%
  \DeltaRow{Mengen}{A\dsep W\dsep e}\DeltaRow{Funktionen}{s}}

\smallskip\noindent\textit{Zu W0.} Das Anfangswort gehört zu \(W\), und
die Anfügung ist für jedes Paar aus einem Wort und einem Buchstaben
definiert und liefert wieder ein Wort. Für
\((\FiniteWords A,\EmptyWord A,\sigma_A)\) ist dies genau
\FormulaRefAuto{EarlyWordSuccessorTyping}[thm].
\end{minipage}\par

\noindent\begin{minipage}{\linewidth}
\FormulaAxiomDeltaK[Das Anfangswort ist kein Nachfolger]{%
  \mathsf{WortStr}_A(W,e,s)\dsep u\in W\dsep a\in A\vdash s(u,a)\ne e
}{WordStructureZeroAxiom}{%
  \DeltaRow{Mengen}{A\dsep W\dsep e\dsep u\dsep a}\DeltaRow{Funktionen}{s}}

\smallskip\noindent\textit{Zu W1.} Durch das Anfügen eines Buchstabens
entsteht niemals das Anfangswort. Im Anfügungsmodell ist diese Aussage
durch \FormulaRefAuto{EarlyWordSuccessorNonempty}[thm] bewiesen.
\end{minipage}\par

\noindent\begin{minipage}{\linewidth}
\FormulaAxiomDeltaK[Gemeinsame Injektivität der Wortanfügung]{%
  \begin{aligned}[t]
    &\mathsf{WortStr}_A(W,e,s)\dsep u,v\in W\dsep a,b\in A\dsep{}\\[-2pt]
    &s(u,a)=s(v,b)\vdash u=v\land a=b
  \end{aligned}
}{WordStructureInjectivityAxiom}{%
  \DeltaRow{Mengen}{A\dsep W\dsep e\dsep u\dsep v\dsep a\dsep b}
  \DeltaRow{Funktionen}{s}}

\smallskip\noindent\textit{Zu W2.} Gleiche angefügte Wörter haben sowohl
denselben bisherigen Wortteil als auch denselben letzten Buchstaben.
Die beiden Eingaben lassen sich also eindeutig zurückgewinnen.
Injektivität nur bei festgehaltenem Buchstaben würde dafür nicht
genügen. Der konkrete Nachweis ist
\FormulaRefAuto{EarlyWordSuccessorInjective}[thm].
\end{minipage}\par

\noindent\begin{minipage}{\linewidth}
\FormulaAxiomDeltaK[Minimalität einer Wortstruktur]{%
  \begin{aligned}[t]
    &\mathsf{WortStr}_A(W,e,s)\dsep U\subseteq W\dsep e\in U\dsep{}\\[-2pt]
    &\forall u\in U\,\forall a\in A\;s(u,a)\in U\vdash U=W
  \end{aligned}
}{WordStructureMinimalityAxiom}{%
  \DeltaRow{Mengen}{A\dsep W\dsep e\dsep U\dsep u\dsep a}
  \DeltaRow{Funktionen}{s}}

\smallskip\noindent\textit{Zu W3.} Enthält eine Teilmenge \(U\) das
Anfangswort und mit jedem ihrer Wörter auch alle Anfügungen, so umfasst
sie bereits ganz \(W\). Damit gibt es keine echte Teilmenge von \(W\),
die unter diesen Erzeugungsregeln abgeschlossen ist und den Anfang
enthält. Im Anfügungsmodell liefert
\FormulaRefAuto{EarlyWordModelMinimality}[thm] genau diese Aussage.
\end{minipage}\par

\Needspace{40\baselineskip}
\FormulaThmDeltaK[Existenz einer freien Wortstruktur]{%
  \mathsf{WortStr}_A(\FiniteWords A,\EmptyWord A,\sigma_A)
}{WordStructureConcreteModel}{\DeltaRow{Mengen}{A\dsep U\dsep u\dsep v\dsep a\dsep b}}
\begin{tabproofwide}
  \proofstepwidestar[]{\EmptyWord A\in\FiniteWords A\land
    \sigma_A\colon\FiniteWords A\times A\to\FiniteWords A}{%
    \FormulaRefAuto{EarlyWordSuccessorTyping}[thm]}
  \proofstepwidestar[2]{u\in\FiniteWords A}{\rA}
  \proofstepwidestar[3]{a\in A}{\rA}
  \proofstepwidestar[2,3]{\sigma_A(u,a)\ne\EmptyWord A}{%
    \FormulaRefAuto{EarlyWordSuccessorNonempty}[thm]{2,3}}
  \proofstepwidestar[2]{\forall a\in A\;\sigma_A(u,a)\ne\EmptyWord A}{%
    \rUI{\rRI{3,4}}}
  \proofstepwidestar[]{\forall u\in\FiniteWords A\,\forall a\in A\;
    \sigma_A(u,a)\ne\EmptyWord A}{\rUI{\rRI{2,5}}}
  \proofstepwidestar[7]{u\in\FiniteWords A}{\rA}
  \proofstepwidestar[8]{v\in\FiniteWords A}{\rA}
  \proofstepwidestar[9]{a\in A}{\rA}
  \proofstepwidestar[10]{b\in A}{\rA}
  \proofstepwidestar[11]{\sigma_A(u,a)=\sigma_A(v,b)}{\rA}
  \proofstepwidestar[7,8,9,10,11]{u=v\land a=b}{%
    \FormulaRefAuto{EarlyWordSuccessorInjective}[thm]{\rAI{7,8},\rAI{9,10},11}}
  \proofstepwidestar[7,8,9,10]{\sigma_A(u,a)=\sigma_A(v,b)\rightarrow
    u=v\land a=b}{\rRI{11,12}}
  \proofstepwidestar[7,8,9]{\forall b\in A\;
    (\sigma_A(u,a)=\sigma_A(v,b)\rightarrow u=v\land a=b)}{%
    \rUI{\rRI{10,13}}}
  \proofstepwidestar[7,8]{\begin{aligned}[t]
    &\forall a,b\in A\\[-2pt]
    &\quad(\sigma_A(u,a)=\sigma_A(v,b)\rightarrow u=v\land a=b)
  \end{aligned}}{\rUI{\rRI{9,14}}}
  \proofstepwidestar[7]{\begin{aligned}[t]
    &\forall v\in\FiniteWords A\,\forall a,b\in A\\[-2pt]
    &\quad(\sigma_A(u,a)=\sigma_A(v,b)\rightarrow u=v\land a=b)
  \end{aligned}}{\rUI{\rRI{8,15}}}
  \proofstepwidestar[]{\begin{aligned}[t]
    &\forall u,v\in\FiniteWords A\,\forall a,b\in A\\[-2pt]
    &\quad(\sigma_A(u,a)=\sigma_A(v,b)\rightarrow u=v\land a=b)
  \end{aligned}}{\rUI{\rRI{7,16}}}
  \proofstepwidestar[18]{U\subseteq\FiniteWords A}{\rA}
  \proofstepwidestar[19]{\begin{aligned}[t]
    &\EmptyWord A\in U\land{}\\[-2pt]
    &(\forall u\in U\,\forall a\in A\;\sigma_A(u,a)\in U)
  \end{aligned}}{\rA}
  \proofstepwidestar[18,19]{U=\FiniteWords A}{%
    \FormulaRefAuto{EarlyWordModelMinimality}[thm]{18,\rAEa{19},\rAEb{19}}}
  \proofstepwidestar[18]{\begin{aligned}[t]
    &\bigl(\EmptyWord A\in U\land{}\\[-2pt]
    &\quad(\forall u\in U\,\forall a\in A\;\sigma_A(u,a)\in U)\bigr)
      \rightarrow U=\FiniteWords A
  \end{aligned}}{\rRI{19,20}}
  \proofstepwidestar[]{\begin{aligned}[t]
    &\forall U\subseteq\FiniteWords A\;\Bigl(\bigl(\EmptyWord A\in U\land{}\\[-2pt]
    &\quad(\forall u\in U\,\forall a\in A\;\sigma_A(u,a)\in U)\bigr)
       \rightarrow U=\FiniteWords A\Bigr)
  \end{aligned}}{\rUI{\rRI{18,21}}}
  \proofstepwidestar[]{\mathsf{WortStr}_A(\FiniteWords A,\EmptyWord A,\sigma_A)}{%
    \FormulaRefAuto{WordStructureDef}[def]{1,6,17,22}}
\end{tabproofwide}

Damit ist die Erfüllbarkeit der Axiome nachgewiesen. Die folgenden
Struktursätze werden aus W0--W3 hergeleitet; für das Anfügungsmodell
werden sie anschließend durch Einsetzen dieses Modells angewendet.

\subsection{Induktion und Zerlegung aus den Wortaxiomen}

Die folgenden Hilfssätze machen die Typisierung und die Behandlung eines geordneten Paars als Funktionsargument ausdrücklich verfügbar.

\noindent\begin{minipage}{\linewidth}
\FormulaThmDeltaK[Typisierung eines Anfügungsschritts]{%
\mathsf{WortStr}_A(W,e,s)\dsep u\in W\dsep a\in A\vdash s(u,a)\in W
}{WordStructureStepTyping}{\DeltaRow{Mengen}{A\dsep W\dsep e\dsep u\dsep a}\DeltaRow{Funktionen}{s}}
\end{minipage}\par
\begin{tabproofwide}
  \proofstepwidestar[1]{\mathsf{WortStr}_A(W,e,s)}{\rA}
  \proofstepwidestar[2]{u\in W}{\rA}
  \proofstepwidestar[3]{a\in A}{\rA}
  \proofstepwidestar[1]{e\in W\land s\colon W\times A\to W}{\FormulaRefAuto{WordStructureTypingAxiom}[ax]{1}}
  \proofstepwidestar[2,3]{(u,a)\in W\times A}{\FormulaRefAuto{(a,b)\in A\times B\eqvdash a\in A\land b\in B}[thm]{\rAI{2,3}}}
  \proofstepwidestar[1,2,3]{s(u,a)\in W}{\FormulaRefAuto{F\colon A\to B\dsep x\in A\vdash F(x)\in B}[thm]{\rAEb{4},5}}
\end{tabproofwide}


\noindent\begin{minipage}{\linewidth}
\FormulaThmDeltaK[Injektivität auf geordneten Paaren]{%
\mathsf{WortStr}_A(W,e,s)\dsep p,q\in W\times A\dsep s(p)=s(q)\vdash p=q
}{WordStructurePairInjective}{\DeltaRow{Mengen}{A\dsep W\dsep e\dsep p\dsep q}\DeltaRow{Funktionen}{s}\DeltaRow{Projektionen}{\pi_1\dsep\pi_2}}
\end{minipage}\par
\begin{tabproofwide}
  \proofstepwidestar[1]{\mathsf{WortStr}_A(W,e,s)}{\rA}
  \proofstepwidestar[2]{p,q\in W\times A}{\rA}
  \proofstepwidestar[3]{s(p)=s(q)}{\rA}
  \proofstepwidestar[2]{\pi_1(p)\in W}{\FormulaRefAuto{z\in A\times B\vdash\pi_1(z)\in A}[thm]{\rAEa{2}}}
  \proofstepwidestar[2]{\pi_2(p)\in A}{\FormulaRefAuto{z\in A\times B\vdash\pi_2(z)\in B}[thm]{\rAEa{2}}}
  \proofstepwidestar[2]{\pi_1(q)\in W}{\FormulaRefAuto{z\in A\times B\vdash\pi_1(z)\in A}[thm]{\rAEb{2}}}
  \proofstepwidestar[2]{\pi_2(q)\in A}{\FormulaRefAuto{z\in A\times B\vdash\pi_2(z)\in B}[thm]{\rAEb{2}}}
  \proofstepwidestar[2]{p=(\pi_1(p),\pi_2(p))}{\FormulaRefAuto{z\in A\times B\vdash z=\bigl(\pi_1(z),\pi_2(z)\bigr)}[thm]{\rAEa{2}}}
  \proofstepwidestar[2]{q=(\pi_1(q),\pi_2(q))}{\FormulaRefAuto{z\in A\times B\vdash z=\bigl(\pi_1(z),\pi_2(z)\bigr)}[thm]{\rAEb{2}}}
  \proofstepwidestar[2,3]{s(\pi_1(p),\pi_2(p))=s(\pi_1(q),\pi_2(q))}{\rIE{8,9,3}}
  \proofstepwidestar[1,2,3]{\pi_1(p)=\pi_1(q)\land\pi_2(p)=\pi_2(q)}{\FormulaRefAuto{WordStructureInjectivityAxiom}[ax]{1,4,6,5,7,10}}
  \proofstepwidestar[1,2,3]{p=q}{\rIE{\FormulaRefAuto{a=b\vdash b=a}[thm]{9},\rIE{\rAEa{11},\rAEb{11},8}}}
\end{tabproofwide}





\noindent\begin{minipage}{\linewidth}
\FormulaThmDeltaK[Induktion in einer Wortstruktur]{%
\begin{aligned}[t]&\mathsf{WortStr}_A(W,e,s)\dsep P(e)\dsep{}\\[-2pt]&\forall u\in W\,\forall a\in A\,(P(u)\rightarrow P(s(u,a)))\vdash\forall w\in W\,P(w)\end{aligned}
}{WordStructureInduction}{\DeltaRow{Mengen}{A\dsep W\dsep e\dsep u\dsep a\dsep w}\DeltaRow{Funktionen}{s}\DeltaRow{Prädikat}{P}\DeltaRow{Abkürzung}{U=\{w\in W\mid P(w)\}}}
\end{minipage}\par
\begin{tabproofwide}
  \proofstepwidestar[1]{\mathsf{WortStr}_A(W,e,s)}{\rA}
  \proofstepwidestar[2]{P(e)}{\rA}
  \proofstepwidestar[3]{\forall u\in W\,\forall a\in A\,(P(u)\rightarrow P(s(u,a)))}{\rA}
  \proofstepwidestar[]{U\subseteq W}{\FormulaRefAuto{\{x\in A\mid P(x)\}\subseteq A}[thm]}
  \proofstepwidestar[1]{e\in W\land s\colon W\times A\to W}{\FormulaRefAuto{WordStructureTypingAxiom}[ax]{1}}
  \proofstepwidestar[1,2]{e\in U}{\FormulaRefAuto{x\in\{u\in A\mid P(u)\}\eqvdash x\in A\land P(x)}[thm]{\rAI{\rAEa{5},2}}}
  \proofstepwidestar[7]{u\in U}{\rA}
  \proofstepwidestar[8]{a\in A}{\rA}
  \proofstepwidestar[7]{u\in W\land P(u)}{\FormulaRefAuto{x\in\{u\in A\mid P(u)\}\eqvdash x\in A\land P(x)}[thm]{7}}
  \proofstepwidestar[3,7,8]{P(u)\rightarrow P(s(u,a))}{\rRE{8,\rUE{\rRE{\rAEa{9},\rUE{3}}}}}
  \proofstepwidestar[3,7,8]{P(s(u,a))}{\rRE{\rAEb{9},10}}
  \proofstepwidestar[1,7,8]{s(u,a)\in W}{\FormulaRefAuto{WordStructureStepTyping}[thm]{1,\rAEa{9},8}}
  \proofstepwidestar[1,3,7,8]{s(u,a)\in U}{\FormulaRefAuto{x\in\{u\in A\mid P(u)\}\eqvdash x\in A\land P(x)}[thm]{\rAI{12,11}}}
  \proofstepwidestar[1,3]{\forall u\in U\,\forall a\in A\,s(u,a)\in U}{\rUI{\rRI{7,\rUI{\rRI{8,13}}}}}
  \proofstepwidestar[1,2,3]{U=W}{\FormulaRefAuto{WordStructureMinimalityAxiom}[ax]{1,4,6,14}}
  \proofstepwidestar[16]{w\in W}{\rA}
  \proofstepwidestar[1,2,3,16]{w\in U}{\rIE{\FormulaRefAuto{a=b\vdash b=a}[thm]{15},16}}
  \proofstepwidestar[1,2,3,16]{w\in W\land P(w)}{\FormulaRefAuto{x\in\{u\in A\mid P(u)\}\eqvdash x\in A\land P(x)}[thm]{17}}
  \proofstepwidestar[1,2,3]{\forall w\in W\,P(w)}{\rUI{\rRI{16,\rAEb{18}}}}
\end{tabproofwide}


Die Mengenform und die Prädikatsform der Induktion sind gleichwertig: In der Rückrichtung setzt man \(P(w)\coloneqq w\in U\) und benutzt die Teilmengendefinition.


\Needspace{32\baselineskip}
\noindent\begin{minipage}{\linewidth}
\FormulaThmDeltaK[Induktionsregel in einer Wortstruktur]{%
  \begin{aligned}[t]
    &\mathsf{WortStr}_A(W,e,s)\dsep P(e)\dsep{}\\[-2pt]
    &[u\in W,a\in A,P(u)]\vdots P(s(u,a))\dsep{}\\[-2pt]
    &v\in W\vdash P(v)
  \end{aligned}
}{WordStructureInductionRule}{%
  \DeltaRow{Mengen}{A\dsep W\dsep e\dsep u\dsep a\dsep v}
  \DeltaRow{Funktionen}{s}
  \DeltaRow{Einstelliges Prädikat}{P}
}
\end{minipage}\par
Diese Regel wird aus der durch W3 begründeten Induktion abgeleitet und
gilt für jede Wortstruktur. Wie bei der konkreten Wortregel sind
\(u,a\) beliebig und kommen in den übrigen offenen Annahmen nicht frei
vor. Die Notation \([i,j,k]\vdots l\) in einem Regelverweis nennt die
entlassenen lokalen Annahmen und die Endzeile des Schrittbeweises.
\Needspace{9\baselineskip}
\begin{tabproofwide}
  \proofstepwidestar[1]{\mathsf{WortStr}_A(W,e,s)}{\rA}
  \proofstepwidestar[2]{P(e)}{\rA}
  \proofstepwidestar[3]{v\in W}{\rA}
  \proofstepwidestar[4]{u\in W}{\rA}
  \proofstepwidestar[5]{a\in A}{\rA}
  \proofstepwidestar[6]{P(u)}{\rA}
  \proofstepwidestar[4,5,6]{P(s(u,a))}{%
    \begin{aligned}[t]
      &\text{Schrittprämisse}\\[-2pt]
      &(4,5,6)
    \end{aligned}}
  \proofstepwidestar[]{\forall u\in W\,\forall a\in A\,
    (P(u)\rightarrow P(s(u,a)))}{%
    \rUI{\rRI{4,\rUI{\rRI{5,\rRI{6,7}}}}}}
  \proofstepwidestar[1,2,3]{P(v)}{%
    \rRE{3,\rUE{\FormulaRefAuto{WordStructureInduction}[thm]{1,2,8}}}}
\end{tabproofwide}

\Needspace{16\baselineskip}
\noindent\begin{minipage}{\linewidth}
\FormulaThmDeltaK[Eindeutige Endzerlegung in einer Wortstruktur]{%
\mathsf{WortStr}_A(W,e,s)\dsep w\in W\vdash w=e\lor\exists!p\in W\times A\;w=s(p)
}{WordStructureDecomposition}{\DeltaRow{Mengen}{A\dsep W\dsep e\dsep w\dsep u\dsep a\dsep p\dsep q}\DeltaRow{Funktionen}{s}\DeltaRow{Abkürzung}{P(w)\coloneqq w=e\lor\exists p\in W\times A\;w=s(p)}}
\end{minipage}\par
\begin{tabproofwide}
  \proofstepwidestar[1]{\mathsf{WortStr}_A(W,e,s)}{\rA}
  \proofstepwidestar[2]{w\in W}{\rA}
  \proofstepwidestar[]{e=e}{\rII}
  \proofstepwidestar[]{P(e)}{\rOIa{3}}
  \proofstepwidestar[5]{u\in W}{\rA}
  \proofstepwidestar[6]{a\in A}{\rA}
  \proofstepwidestar[5,6]{(u,a)\in W\times A}{\FormulaRefAuto{(a,b)\in A\times B\eqvdash a\in A\land b\in B}[thm]{\rAI{5,6}}}
  \proofstepwidestar[]{s(u,a)=s(u,a)}{\rII}
  \proofstepwidestar[5,6]{\exists p\in W\times A\;s(u,a)=s(p)}{\rEI{\rAI{7,8}}}
  \proofstepwidestar[5,6]{P(s(u,a))}{\rOIb{9}}
  \proofstepwidestar[1,2]{w=e\lor\exists p\in W\times A\;w=s(p)}{%
    \FormulaRefAuto{WordStructureInductionRule}[thm]{1,4,[5,6]\vdots10,2}}
  \proofstepwidestar[12]{w=e}{\rA}
  \proofstepwidestar[12]{w=e\lor\exists!p\in W\times A\;w=s(p)}{\rOIa{12}}
  \proofstepwidestar[14]{\exists p\in W\times A\;w=s(p)}{\rA}
  \proofstepwidestar[15]{p\in W\times A\land w=s(p)}{\rA}
  \proofstepwidestar[16]{q\in W\times A\land w=s(q)}{\rA}
  \proofstepwidestar[15,16]{s(q)=s(p)}{\FormulaRefAuto{a=b\dsep a=c\vdash b=c}[thm]{\rAEb{16},\rAEb{15}}}
  \proofstepwidestar[1,15,16]{q=p}{\FormulaRefAuto{WordStructurePairInjective}[thm]{1,\rAI{\rAEa{16},\rAEa{15}},17}}
  \proofstepwidestar[1,15]{\forall q\,((q\in W\times A\land w=s(q))\rightarrow q=p)}{\rUI{\rRI{16,18}}}
  \proofstepwidestar[1,15]{\exists!p\in W\times A\;w=s(p)}{\FormulaRefAuto{WordUniqueExistenceFromWitness}[thm]{15,19}}
  \proofstepwidestar[1,15]{w=e\lor\exists!p\in W\times A\;w=s(p)}{\rOIb{20}}
  \proofstepwidestar[1,14]{w=e\lor\exists!p\in W\times A\;w=s(p)}{\rEE{14,15,21}}
  \proofstepwidestar[1,2]{w=e\lor\exists!p\in W\times A\;w=s(p)}{\rOE{11,12,13,14,22}}
\end{tabproofwide}


Die beiden Fälle schließen einander nach \FormulaRefAuto{WordStructureZeroAxiom}[ax] aus. Für ein einelementiges Alphabet ist die Anfügung ein einzelner Nachfolger; die Bedingungen W0 bis W3 entsprechen damit Anfangselement, Nachfolgerabbildung, Injektivität und Induktion der Peano-Struktur.

\subsection{Rekursion unmittelbar aus den Wortaxiomen}
Die Wortaxiome liefern Rekursion, indem aus allen abgeschlossenen Relationsgraphen ein kleinster ausgewählt wird. Die folgenden Definitionen halten die dafür benutzten Mengen und ihre Elementkriterien ausdrücklich fest. Ihre Parameter \(A,W,e,s,X,x_0,r\) bleiben in diesem Abschnitt fest.

\FormulaDefDeltaK[Rekursionsgleichungen einer Wortstruktur]{%
\begin{aligned}[t]&\mathsf{WRec}_A(F;W,e,s;X,x_0,r)\coloneqq{}\\[-2pt]
&F\colon W\to X\land F(e)=x_0\land{}\\[-2pt]
&\forall u\in W\,\forall a\in A\;F(s(u,a))=r(F(u),a)
\end{aligned}}{WordRecursionSpecificationDef}{%
\DeltaRow{Mengen}{A\dsep W\dsep e\dsep X\dsep x_0\dsep u\dsep a}
\DeltaRow{Funktionen}{F\dsep s\dsep r}
\DeltaRow{Neues Symbol}{\mathsf{WRec}_A(F;W,e,s;X,x_0,r)}}

\FormulaDefDeltaK[Umgebungsmenge der Rekursionsgraphen]{B\coloneqq W\times X}{WordRecursionAmbientDef}{\DeltaRow{Mengen}{W\dsep X}\DeltaRow{Neues Symbol}{B}}
\FormulaDefDeltaK[Abgeschlossenheit eines Rekursionsgraphen]{%
\begin{aligned}[t]&C(H)\coloneqq H\subseteq B\land(e,x_0)\in H\land{}\\[-2pt]
&\forall u\in W\,\forall x\in X\,\forall a\in A\;\\[-2pt]
&\qquad((u,x)\in H\rightarrow(s(u,a),r(x,a))\in H)
\end{aligned}}{WordRecursionClosedGraphDef}{\DeltaRow{Mengen}{A\dsep W\dsep X\dsep e\dsep x_0\dsep H\dsep u\dsep x\dsep a}\DeltaRow{Funktionen}{s\dsep r}\DeltaRow{Neues Symbol}{C(H)}}
\FormulaDefDeltaK[Familie der abgeschlossenen Graphen]{\mathcal C\coloneqq\{H\in\powerset B\mid C(H)\}}{WordRecursionGraphFamilyDef}{\DeltaRow{Mengen}{H}\DeltaRow{Neues Symbol}{\mathcal C}}
\FormulaDefDeltaK[Der gemeinsame Teil aller abgeschlossenen Graphen]{G\coloneqq\{p\in B\mid\forall H\in\mathcal C\;p\in H\}}{WordRecursionLeastGraphDef}{\DeltaRow{Mengen}{p\dsep H}\DeltaRow{Neues Symbol}{G}}


Die Mengen in diesen Definitionen sind Aussonderungen aus vorhandenen Mengen. Ihre Existenz und Eindeutigkeit liegen in
\FormulaRefAuto{\{x\in A\,\mid\,P(x)\}\coloneq\iota B\bigl(\forall x\,(x\in B\leftrightarrow(x\in A\land P(x)))\bigr)}[def]. Ein Graph muss hier noch keine Funktion sein.
Für \(B=W\times X\) verwenden wir den allgemeinen Faserfilter
\(\Phi^{W,X}_{t,c}(K)\) aus
\FormulaRefAuto{WordRecursionFiberFilterDef}[def].


\noindent\begin{minipage}{\linewidth}
\FormulaThmDeltaK[Mitgliedschaft in der Familie abgeschlossener Graphen]{%
H\in\mathcal C\eqvdash C(H)
}{WordRecursionFamilyMembership}{\DeltaRow{Mengen}{A\dsep W\dsep e\dsep X\dsep x_0\dsep H\dsep K\dsep u\dsep v\dsep x\dsep y\dsep a\dsep b\dsep p\dsep z\dsep t\dsep c}\DeltaRow{Funktionen}{s\dsep r}}
\end{minipage}\par
\begin{tabproofwide}
  \proofstepwidestar[]{\mathcal C=\{H\in\powerset B\mid C(H)\}}{\FormulaRefAuto{WordRecursionGraphFamilyDef}[def]}
  \proofstepwidestar[]{H\in\mathcal C\leftrightarrow H\in\powerset B\land C(H)}{\rIE{\FormulaRefAuto{a=b\vdash b=a}[thm]{1},\FormulaRefAuto{x\in\{u\in A\mid P(u)\}\eqvdash x\in A\land P(x)}[thm]}}
  \proofstepwidestar[3]{C(H)}{\rA}
  \proofstepwidestar[3]{H\subseteq B}{\rAEa{\FormulaRefAuto{WordRecursionClosedGraphDef}[def]{3}}}
  \proofstepwidestar[3]{H\in\powerset B}{\FormulaRefAuto{x\in\powerset(A)\eqvdash x\subseteq A}[thm]{4}}
  \proofstepwidestar[]{C(H)\rightarrow H\in\powerset B}{\rRI{3,5}}
  \proofstepwidestar[]{H\in\mathcal C\leftrightarrow C(H)}{%
    \rLRS{\FormulaRefAuto{P\rightarrow Q\vdash
      (Q\land P)\leftrightarrow P}[thm]{6},2}}
\end{tabproofwide}


\noindent\begin{minipage}{\linewidth}
\FormulaThmDeltaK[Elementkriterium des kleinsten Rekursionsgraphen]{%
p\in G\eqvdash p\in B\land\forall H\in\mathcal C\;p\in H
}{WordRecursionLeastMembership}{\DeltaRow{Mengen}{A\dsep W\dsep e\dsep X\dsep x_0\dsep H\dsep K\dsep u\dsep v\dsep x\dsep y\dsep a\dsep b\dsep p\dsep z\dsep t\dsep c}\DeltaRow{Funktionen}{s\dsep r}}
\end{minipage}\par
\begin{tabproofwide}
  \proofstepwidestar[]{G=\{p\in B\mid\forall H\in\mathcal C\;p\in H\}}{\FormulaRefAuto{WordRecursionLeastGraphDef}[def]}
  \proofstepwidestar[]{p\in G\leftrightarrow p\in B\land\forall H\in\mathcal C\;p\in H}{\rIE{\FormulaRefAuto{a=b\vdash b=a}[thm]{1},\FormulaRefAuto{x\in\{u\in A\mid P(u)\}\eqvdash x\in A\land P(x)}[thm]}}
\end{tabproofwide}


\noindent\begin{minipage}{\linewidth}
\FormulaThmDeltaK[Trägermenge eines abgeschlossenen Graphen]{%
C(H)\vdash H\subseteq B
}{WordRecursionClosedGraphSubset}{\DeltaRow{Mengen}{A\dsep W\dsep e\dsep X\dsep x_0\dsep H\dsep K\dsep u\dsep v\dsep x\dsep y\dsep a\dsep b\dsep p\dsep z\dsep t\dsep c}\DeltaRow{Funktionen}{s\dsep r}}
\end{minipage}\par
\begin{tabproofwide}
  \proofstepwidestar[1]{C(H)}{\rA}
  \proofstepwidestar[1]{\begin{aligned}[t]&H\subseteq B\land(e,x_0)\in H\land{}\\&\begin{aligned}[t]&\forall u\in W\,\forall x\in X\,\forall a\in A\\[-2pt]&\quad((u,x)\in H\rightarrow(s(u,a),r(x,a))\in H)\end{aligned}\end{aligned}}{\FormulaRefAuto{WordRecursionClosedGraphDef}[def]{1}}
  \proofstepwidestar[1]{ H\subseteq B}{\rAEa{2}}
\end{tabproofwide}


\noindent\begin{minipage}{\linewidth}
\FormulaThmDeltaK[Anfangspunkt eines abgeschlossenen Graphen]{%
C(H)\vdash(e,x_0)\in H
}{WordRecursionClosedGraphBase}{\DeltaRow{Mengen}{A\dsep W\dsep e\dsep X\dsep x_0\dsep H\dsep K\dsep u\dsep v\dsep x\dsep y\dsep a\dsep b\dsep p\dsep z\dsep t\dsep c}\DeltaRow{Funktionen}{s\dsep r}}
\end{minipage}\par
\begin{tabproofwide}
  \proofstepwidestar[1]{C(H)}{\rA}
  \proofstepwidestar[1]{\begin{aligned}[t]&H\subseteq B\land(e,x_0)\in H\land{}\\&\begin{aligned}[t]&\forall u\in W\,\forall x\in X\,\forall a\in A\\[-2pt]&\quad((u,x)\in H\rightarrow(s(u,a),r(x,a))\in H)\end{aligned}\end{aligned}}{\FormulaRefAuto{WordRecursionClosedGraphDef}[def]{1}}
  \proofstepwidestar[1]{(e,x_0)\in H}{\rAEa{\rAEb{2}}}
\end{tabproofwide}


\noindent\begin{minipage}{\linewidth}
\FormulaThmDeltaK[Einzelner Schritt in einem abgeschlossenen Graphen]{%
\begin{aligned}[t]&C(H)\dsep u\in W\dsep x\in X\dsep a\in A\dsep{}\\[-2pt]&(u,x)\in H\vdash(s(u,a),r(x,a))\in H\end{aligned}
}{WordRecursionClosedGraphStep}{\DeltaRow{Mengen}{A\dsep W\dsep e\dsep X\dsep x_0\dsep H\dsep K\dsep u\dsep v\dsep x\dsep y\dsep a\dsep b\dsep p\dsep z\dsep t\dsep c}\DeltaRow{Funktionen}{s\dsep r}}
\end{minipage}\par
\begin{tabproofwide}
  \proofstepwidestar[1]{C(H)}{\rA}
  \proofstepwidestar[2]{u\in W}{\rA}
  \proofstepwidestar[3]{x\in X}{\rA}
  \proofstepwidestar[4]{a\in A}{\rA}
  \proofstepwidestar[5]{(u,x)\in H}{\rA}
  \proofstepwidestar[1]{\begin{aligned}[t]&\forall u\in W\,\forall x\in X\,\forall a\in A\\[-2pt]&\quad((u,x)\in H\rightarrow(s(u,a),r(x,a))\in H)\end{aligned}}{\rAEb{\rAEb{\FormulaRefAuto{WordRecursionClosedGraphDef}[def]{1}}}}
  \proofstepwidestar[1,2,3,4]{(u,x)\in H\rightarrow(s(u,a),r(x,a))\in H}{\rRE{4,\rUE{\rRE{3,\rUE{\rRE{2,\rUE{6}}}}}}}
  \proofstepwidestar[1,2,3,4,5]{(s(u,a),r(x,a))\in H}{\rRE{5,7}}
\end{tabproofwide}


\noindent\begin{minipage}{\linewidth}
\FormulaThmDeltaK[Kleinster abgeschlossener Rekursionsgraph]{%
\begin{aligned}[t]&\mathsf{WortStr}_A(W,e,s)\dsep x_0\in X\dsep r\colon X\times A\to X\vdash{}\\[-2pt]&C(G)\land\forall H\,(C(H)\rightarrow G\subseteq H)\end{aligned}
}{WordStructureRecursionLeastGraph}{\DeltaRow{Mengen}{A\dsep W\dsep e\dsep X\dsep x_0\dsep H\dsep K\dsep u\dsep v\dsep x\dsep y\dsep a\dsep b\dsep p\dsep z\dsep t\dsep c}\DeltaRow{Funktionen}{s\dsep r}}
\end{minipage}\par
\begin{tabproofwide}
  \proofstepwidestar[1]{\mathsf{WortStr}_A(W,e,s)}{\rA}
  \proofstepwidestar[2]{x_0\in X}{\rA}
  \proofstepwidestar[3]{r\colon X\times A\to X}{\rA}
  \proofstepwidestar[1]{e\in W\land s\colon W\times A\to W}{\FormulaRefAuto{WordStructureTypingAxiom}[ax]{1}}
  \proofstepwidestar[]{G\subseteq B}{\rIE{\FormulaRefAuto{a=b\vdash b=a}[thm]{\FormulaRefAuto{WordRecursionLeastGraphDef}[def]},\FormulaRefAuto{\{x\in A\mid P(x)\}\subseteq A}[thm]}}
  \proofstepwidestar[1,2]{(e,x_0)\in B}{\rIE{\FormulaRefAuto{a=b\vdash b=a}[thm]{\FormulaRefAuto{WordRecursionAmbientDef}[def]},\FormulaRefAuto{(a,b)\in A\times B\eqvdash a\in A\land b\in B}[thm]{\rAI{\rAEa{4},2}}}}
  \proofstepwidestar[7]{H\in\mathcal C}{\rA}
  \proofstepwidestar[7]{C(H)}{\FormulaRefAuto{WordRecursionFamilyMembership}[thm]{7}}
  \proofstepwidestar[7]{(e,x_0)\in H}{\FormulaRefAuto{WordRecursionClosedGraphBase}[thm]{8}}
  \proofstepwidestar[]{\forall H\in\mathcal C\;(e,x_0)\in H}{\rUI{\rRI{7,9}}}
  \proofstepwidestar[1,2]{(e,x_0)\in G}{\FormulaRefAuto{WordRecursionLeastMembership}[thm]{\rAI{6,10}}}
  \proofstepwidestar[12]{u\in W}{\rA}
  \proofstepwidestar[13]{y\in X}{\rA}
  \proofstepwidestar[14]{a\in A}{\rA}
  \proofstepwidestar[15]{(u,y)\in G}{\rA}
  \proofstepwidestar[1,12,14]{s(u,a)\in W}{\FormulaRefAuto{WordStructureStepTyping}[thm]{1,12,14}}
  \proofstepwidestar[13,14]{(y,a)\in X\times A}{\FormulaRefAuto{(a,b)\in A\times B\eqvdash a\in A\land b\in B}[thm]{\rAI{13,14}}}
  \proofstepwidestar[3,13,14]{r(y,a)\in X}{\FormulaRefAuto{F\colon A\to B\dsep x\in A\vdash F(x)\in B}[thm]{3,17}}
  \proofstepwidestar[1,3,12,13,14]{(s(u,a),r(y,a))\in B}{\rIE{\FormulaRefAuto{a=b\vdash b=a}[thm]{\FormulaRefAuto{WordRecursionAmbientDef}[def]},\FormulaRefAuto{(a,b)\in A\times B\eqvdash a\in A\land b\in B}[thm]{\rAI{16,18}}}}
  \proofstepwidestar[15]{(u,y)\in B\land\forall H\in\mathcal C\;(u,y)\in H}{\FormulaRefAuto{WordRecursionLeastMembership}[thm]{15}}
  \proofstepwidestar[21]{H\in\mathcal C}{\rA}
  \proofstepwidestar[21]{C(H)}{\FormulaRefAuto{WordRecursionFamilyMembership}[thm]{21}}
  \proofstepwidestar[15,21]{(u,y)\in H}{\rRE{21,\rUE{\rAEb{20}}}}
  \proofstepwidestar[12,13,14,15,21]{(s(u,a),r(y,a))\in H}{\FormulaRefAuto{WordRecursionClosedGraphStep}[thm]{22,12,13,14,23}}
  \proofstepwidestar[12,13,14,15]{\forall H\in\mathcal C\;(s(u,a),r(y,a))\in H}{\rUI{\rRI{21,24}}}
  \proofstepwidestar[1,3,12,13,14,15]{(s(u,a),r(y,a))\in G}{\FormulaRefAuto{WordRecursionLeastMembership}[thm]{\rAI{19,25}}}
  \proofstepwidestar[1,3]{\begin{aligned}[t]&\forall u\in W\,\forall y\in X\,\forall a\in A\\[-2pt]&\quad((u,y)\in G\rightarrow(s(u,a),r(y,a))\in G)\end{aligned}}{\rUI{\rRI{12,\rUI{\rRI{13,\rUI{\rRI{14,\rRI{15,26}}}}}}}}
  \proofstepwidestar[1,2,3]{C(G)}{\FormulaRefAuto{WordRecursionClosedGraphDef}[def]{\rAI{5,\rAI{11,27}}}}
  \proofstepwidestar[29]{C(H)}{\rA}
  \proofstepwidestar[29]{H\in\mathcal C}{\FormulaRefAuto{WordRecursionFamilyMembership}[thm]{29}}
  \proofstepwidestar[31]{p\in G}{\rA}
  \proofstepwidestar[31]{p\in B\land\forall H\in\mathcal C\;p\in H}{\FormulaRefAuto{WordRecursionLeastMembership}[thm]{31}}
  \proofstepwidestar[29,31]{p\in H}{\rRE{30,\rUE{\rAEb{32}}}}
  \proofstepwidestar[29]{G\subseteq H}{\FormulaRefAuto{SubsetIntroductionRule}[thm]{[31]\vdots 33}}
  \proofstepwidestar[]{\forall H\,(C(H)\rightarrow G\subseteq H)}{\rUI{\rRI{29,34}}}
  \proofstepwidestar[1,2,3]{C(G)\land\forall H\,(C(H)\rightarrow G\subseteq H)}{\rAI{28,35}}
\end{tabproofwide}


\noindent\begin{minipage}{\linewidth}
\FormulaThmDeltaK[Eindeutiger Anfangswert des kleinsten Graphen]{%
\begin{aligned}[t]&\mathsf{WortStr}_A(W,e,s)\dsep x_0\in X\dsep r\colon X\times A\to X\vdash{}\\[-2pt]&\exists!x\in X\;(e,x)\in G\end{aligned}
}{WordStructureRecursionBaseFiber}{\DeltaRow{Mengen}{A\dsep W\dsep e\dsep X\dsep x_0\dsep H\dsep K\dsep u\dsep v\dsep x\dsep y\dsep a\dsep b\dsep p\dsep z\dsep t\dsep c}\DeltaRow{Funktionen}{s\dsep r}\DeltaRow{Abkürzung}{H_0=\Phi^{W,X}_{e,x_0}(B)}}
\end{minipage}\par
\begin{tabproofwide}
  \proofstepwidestar[1]{\mathsf{WortStr}_A(W,e,s)}{\rA}
  \proofstepwidestar[2]{x_0\in X}{\rA}
  \proofstepwidestar[3]{r\colon X\times A\to X}{\rA}
  \proofstepwidestar[1,2,3]{C(G)\land\forall H\,(C(H)\rightarrow G\subseteq H)}{\FormulaRefAuto{WordStructureRecursionLeastGraph}[thm]{1,2,3}}
  \proofstepwidestar[1]{e\in W\land s\colon W\times A\to W}{\FormulaRefAuto{WordStructureTypingAxiom}[ax]{1}}
  \proofstepwidestar[]{H_0\subseteq B}{\FormulaRefAuto{WordRecursionFilterSubset}[thm]}
  \proofstepwidestar[]{B\subseteq B}{\FormulaRefAuto{A\subseteq A}[thm]}
  \proofstepwidestar[1,2]{(e,x_0)\in B}{\rIE{\FormulaRefAuto{a=b\vdash b=a}[thm]{\FormulaRefAuto{WordRecursionAmbientDef}[def]},\FormulaRefAuto{(a,b)\in A\times B\eqvdash a\in A\land b\in B}[thm]{\rAI{\rAEa{5},2}}}}
  \proofstepwidestar[]{x_0=x_0}{\rII}
  \proofstepwidestar[]{e\ne e\lor x_0=x_0}{\rOIb{9}}
  \proofstepwidestar[1,2]{(e,x_0)\in H_0}{\FormulaRefAuto{WordRecursionFilterMembership}[thm]{\rIE{\FormulaRefAuto{WordRecursionAmbientDef}[def],7},\rAI{8,10}}}
  \proofstepwidestar[12]{u\in W}{\rA}
  \proofstepwidestar[13]{y\in X}{\rA}
  \proofstepwidestar[14]{a\in A}{\rA}
  \proofstepwidestar[1,12,14]{s(u,a)\in W}{\FormulaRefAuto{WordStructureStepTyping}[thm]{1,12,14}}
  \proofstepwidestar[13,14]{(y,a)\in X\times A}{\FormulaRefAuto{(a,b)\in A\times B\eqvdash a\in A\land b\in B}[thm]{\rAI{13,14}}}
  \proofstepwidestar[3,13,14]{r(y,a)\in X}{\FormulaRefAuto{F\colon A\to B\dsep x\in A\vdash F(x)\in B}[thm]{3,16}}
  \proofstepwidestar[1,3,12,13,14]{(s(u,a),r(y,a))\in B}{\rIE{\FormulaRefAuto{a=b\vdash b=a}[thm]{\FormulaRefAuto{WordRecursionAmbientDef}[def]},\FormulaRefAuto{(a,b)\in A\times B\eqvdash a\in A\land b\in B}[thm]{\rAI{15,17}}}}
  \proofstepwidestar[1,12,14]{s(u,a)\ne e}{\FormulaRefAuto{WordStructureZeroAxiom}[ax]{1,12,14}}
  \proofstepwidestar[1,12,14]{s(u,a)\ne e\lor r(y,a)=x_0}{\rOIa{19}}
  \proofstepwidestar[1,3,12,13,14]{(s(u,a),r(y,a))\in H_0}{\FormulaRefAuto{WordRecursionFilterMembership}[thm]{\rIE{\FormulaRefAuto{WordRecursionAmbientDef}[def],7},\rAI{18,20}}}
  \proofstepwidestar[1,3,12,13,14]{(u,y)\in H_0\rightarrow(s(u,a),r(y,a))\in H_0}{\FormulaRefAuto{Q\vdash P\rightarrow Q}[thm]{21}}
  \proofstepwidestar[1,3]{\begin{aligned}[t]&\forall u\in W\,\forall y\in X\,\forall a\in A\\[-2pt]&\quad((u,y)\in H_0\rightarrow(s(u,a),r(y,a))\in H_0)\end{aligned}}{\rUI{\rRI{12,\rUI{\rRI{13,\rUI{\rRI{14,22}}}}}}}
  \proofstepwidestar[1,2,3]{C(H_0)}{\FormulaRefAuto{WordRecursionClosedGraphDef}[def]{\rAI{6,\rAI{11,23}}}}
  \proofstepwidestar[1,2,3]{G\subseteq H_0}{\rRE{24,\rUE{\rAEb{4}}}}
  \proofstepwidestar[1,2,3]{(e,x_0)\in G}{\FormulaRefAuto{WordRecursionClosedGraphBase}[thm]{\rAEa{4}}}
  \proofstepwidestar[27]{z\in X\land(e,z)\in G}{\rA}
  \proofstepwidestar[1,2,3,27]{(e,z)\in H_0}{\FormulaRefAuto{A\subseteq B\dsep x\in A\vdash x\in B}[thm]{25,\rAEb{27}}}
  \proofstepwidestar[1,2,3,27]{z=x_0}{\FormulaRefAuto{WordRecursionFilterOwnFiber}[thm]{\rIE{\FormulaRefAuto{WordRecursionAmbientDef}[def],7},28}}
  \proofstepwidestar[1,2,3]{\forall z\,((z\in X\land(e,z)\in G)\rightarrow z=x_0)}{\rUI{\rRI{27,29}}}
  \proofstepwidestar[1,2,3]{\exists!x\in X\;(e,x)\in G}{\FormulaRefAuto{WordUniqueExistenceFromWitness}[thm]{\rAI{2,26},30}}
\end{tabproofwide}


Für den Induktionsschritt wird die Faser über \(s(u,a)\) gefiltert. Der folgende Beweis führt die dabei benötigte Injektivität und Faser-Eindeutigkeit getrennt aus.

\noindent\begin{minipage}{\linewidth}
\FormulaThmDeltaK[Eindeutige Werte bleiben beim Anfügen eindeutig]{%
\begin{aligned}[t]&\mathsf{WortStr}_A(W,e,s)\dsep x_0\in X\dsep r\colon X\times A\to X\dsep{}\\[-2pt]&u\in W\dsep a\in A\dsep\exists!x\in X\;(u,x)\in G\vdash{}\\[-2pt]&\exists!z\in X\;(s(u,a),z)\in G\end{aligned}
}{WordStructureRecursionStepFiber}{\DeltaRow{Mengen}{A\dsep W\dsep e\dsep X\dsep x_0\dsep H\dsep K\dsep u\dsep v\dsep x\dsep y\dsep a\dsep b\dsep p\dsep z\dsep t\dsep c}\DeltaRow{Funktionen}{s\dsep r}\DeltaRow{Abkürzungen}{x=\iota y\in X\;(u,y)\in G\dsep t=s(u,a)\dsep c=r(x,a)}\DeltaRow{Abkürzung}{H=\Phi^{W,X}_{t,c}(G)}}
\end{minipage}\par
\begin{tabproofwide}
  \proofstepwidestar[1]{\mathsf{WortStr}_A(W,e,s)}{\rA}
  \proofstepwidestar[2]{x_0\in X}{\rA}
  \proofstepwidestar[3]{r\colon X\times A\to X}{\rA}
  \proofstepwidestar[4]{u\in W}{\rA}
  \proofstepwidestar[5]{a\in A}{\rA}
  \proofstepwidestar[6]{\exists!x\in X\;(u,x)\in G}{\rA}
  \proofstepwidestar[6]{x\in X\land(u,x)\in G}{\FormulaRefAuto{IotaSymbolDef}[def]{6}}
  \proofstepwidestar[1,2,3]{C(G)\land\forall K\,(C(K)\rightarrow G\subseteq K)}{\FormulaRefAuto{WordStructureRecursionLeastGraph}[thm]{1,2,3}}
  \proofstepwidestar[1,2,3]{G\subseteq B}{\FormulaRefAuto{WordRecursionClosedGraphSubset}[thm]{\rAEa{8}}}
  \proofstepwidestar[]{H\subseteq G}{\FormulaRefAuto{WordRecursionFilterSubset}[thm]}
  \proofstepwidestar[1,2,3]{H\subseteq B}{\FormulaRefAuto{A\subseteq B,B\subseteq C\vdash A\subseteq C}[thm]{10,9}}
  \proofstepwidestar[1,4,5]{t\in W}{\FormulaRefAuto{WordStructureStepTyping}[thm]{1,4,5}}
  \proofstepwidestar[6,5]{(x,a)\in X\times A}{\FormulaRefAuto{(a,b)\in A\times B\eqvdash a\in A\land b\in B}[thm]{\rAI{\rAEa{7},5}}}
  \proofstepwidestar[3,6,5]{c\in X}{\FormulaRefAuto{F\colon A\to B\dsep x\in A\vdash F(x)\in B}[thm]{3,13}}
  \proofstepwidestar[1,2,3]{(e,x_0)\in G}{\FormulaRefAuto{WordRecursionClosedGraphBase}[thm]{\rAEa{8}}}
  \proofstepwidestar[1,4,5]{t\ne e}{\FormulaRefAuto{WordStructureZeroAxiom}[ax]{1,4,5}}
  \proofstepwidestar[1,4,5]{e\ne t}{\FormulaRefAuto{a\neq b\vdash b\neq a}[thm]{16}}
  \proofstepwidestar[1,4,5]{e\ne t\lor x_0=c}{\rOIa{17}}
  \proofstepwidestar[1,2,3,4,5]{(e,x_0)\in H}{\FormulaRefAuto{WordRecursionFilterMembership}[thm]{\rIE{\FormulaRefAuto{WordRecursionAmbientDef}[def],9},\rAI{15,18}}}
  \proofstepwidestar[20]{v\in W}{\rA}
  \proofstepwidestar[21]{y\in X}{\rA}
  \proofstepwidestar[22]{b\in A}{\rA}
  \proofstepwidestar[23]{(v,y)\in H}{\rA}
  \proofstepwidestar[23]{(v,y)\in G}{\FormulaRefAuto{A\subseteq B\dsep x\in A\vdash x\in B}[thm]{10,23}}
  \proofstepwidestar[1,2,3,20,21,22,23]{(s(v,b),r(y,b))\in G}{\FormulaRefAuto{WordRecursionClosedGraphStep}[thm]{\rAEa{8},20,21,22,24}}
  \proofstepwidestar[26]{s(v,b)=t}{\rA}
  \proofstepwidestar[1,4,5,20,22,26]{v=u\land b=a}{\FormulaRefAuto{WordStructureInjectivityAxiom}[ax]{1,20,4,22,5,26}}
  \proofstepwidestar[1,4,5,20,22,26,23]{(u,y)\in G}{\rIE{\rAEa{27},24}}
  \proofstepwidestar[1,4,5,20,21,22,26,23,6]{y=x}{\FormulaRefAuto{\exists! x\,P(x),\; P(a),\; P(b) \vdash a = b}[thm]{6,\rAI{21,28},7}}
  \proofstepwidestar[]{r(x,a)=c}{\rII}
  \proofstepwidestar[1,4,5,20,21,22,26,23,6]{r(y,b)=c}{\rIE{\FormulaRefAuto{a=b\vdash b=a}[thm]{29},\FormulaRefAuto{a=b\vdash b=a}[thm]{\rAEb{27}},30}}
  \proofstepwidestar[1,4,5,20,21,22,23,6]{s(v,b)=t\rightarrow r(y,b)=c}{\rRI{26,31}}
  \proofstepwidestar[1,4,5,20,21,22,23,6]{s(v,b)\ne t\lor r(y,b)=c}{\FormulaRefAuto{P\rightarrow Q\eqvdash\neg P\lor Q}[thm]{32}}
  \proofstepwidestar[1,2,3,4,5,20,21,22,23,6]{(s(v,b),r(y,b))\in H}{\FormulaRefAuto{WordRecursionFilterMembership}[thm]{\rIE{\FormulaRefAuto{WordRecursionAmbientDef}[def],9},\rAI{25,33}}}
  \proofstepwidestar[1,2,3,4,5,6]{\begin{aligned}[t]&\forall v\in W\,\forall y\in X\,\forall b\in A\\&\quad((v,y)\in H\rightarrow(s(v,b),r(y,b))\in H)\end{aligned}}{\rUI{\rRI{20,\rUI{\rRI{21,\rUI{\rRI{22,\rRI{23,34}}}}}}}}
  \proofstepwidestar[1,2,3,4,5,6]{C(H)}{\FormulaRefAuto{WordRecursionClosedGraphDef}[def]{\rAI{11,\rAI{19,35}}}}
  \proofstepwidestar[1,2,3,4,5,6]{G\subseteq H}{\rRE{36,\rUE{\rAEb{8}}}}
  \proofstepwidestar[1,2,3,4,5,6]{(t,c)\in G}{\FormulaRefAuto{WordRecursionClosedGraphStep}[thm]{\rAEa{8},4,\rAEa{7},5,\rAEb{7}}}
  \proofstepwidestar[39]{z\in X\land(t,z)\in G}{\rA}
  \proofstepwidestar[1,2,3,4,5,6,39]{(t,z)\in H}{\FormulaRefAuto{A\subseteq B\dsep x\in A\vdash x\in B}[thm]{37,\rAEb{39}}}
  \proofstepwidestar[1,2,3,4,5,6,39]{z=c}{\FormulaRefAuto{WordRecursionFilterOwnFiber}[thm]{\rIE{\FormulaRefAuto{WordRecursionAmbientDef}[def],9},40}}
  \proofstepwidestar[1,2,3,4,5,6]{\forall z\,((z\in X\land(t,z)\in G)\rightarrow z=c)}{\rUI{\rRI{39,41}}}
  \proofstepwidestar[1,2,3,4,5,6]{\exists!z\in X\;(s(u,a),z)\in G}{\FormulaRefAuto{WordUniqueExistenceFromWitness}[thm]{\rAI{14,38},42}}
\end{tabproofwide}


\noindent\begin{minipage}{\linewidth}
\FormulaThmDeltaK[Der kleinste Rekursionsgraph ist eine Funktion]{%
\begin{aligned}[t]&\mathsf{WortStr}_A(W,e,s)\dsep x_0\in X\dsep r\colon X\times A\to X\vdash{}\\[-2pt]&G\colon W\to X\end{aligned}
}{WordStructureRecursionGraphFunction}{\DeltaRow{Mengen}{A\dsep W\dsep e\dsep X\dsep x_0\dsep H\dsep K\dsep u\dsep v\dsep x\dsep y\dsep a\dsep b\dsep p\dsep z\dsep t\dsep c}\DeltaRow{Funktionen}{s\dsep r}\DeltaRow{Abkürzung}{P(u)\coloneqq\exists!x\in X\;(u,x)\in G}}
\end{minipage}\par
\begin{tabproofwide}
  \proofstepwidestar[1]{\mathsf{WortStr}_A(W,e,s)}{\rA}
  \proofstepwidestar[2]{x_0\in X}{\rA}
  \proofstepwidestar[3]{r\colon X\times A\to X}{\rA}
  \proofstepwidestar[1,2,3]{P(e)}{\FormulaRefAuto{WordStructureRecursionBaseFiber}[thm]{1,2,3}}
  \proofstepwidestar[5]{u\in W}{\rA}
  \proofstepwidestar[6]{a\in A}{\rA}
  \proofstepwidestar[7]{P(u)}{\rA}
  \proofstepwidestar[1,2,3,5,6,7]{P(s(u,a))}{\FormulaRefAuto{WordStructureRecursionStepFiber}[thm]{1,2,3,5,6,7}}
  \proofstepwidestar[1,2,3]{\forall u\in W\,\forall a\in A\;(P(u)\rightarrow P(s(u,a)))}{\rUI{\rRI{5,\rUI{\rRI{6,\rRI{7,8}}}}}}
  \proofstepwidestar[1,2,3]{\forall u\in W\,\exists!x\in X\;(u,x)\in G}{\FormulaRefAuto{WordStructureInduction}[thm]{1,4,9}}
  \proofstepwidestar[1,2,3]{C(G)\land\forall H\,(C(H)\rightarrow G\subseteq H)}{\FormulaRefAuto{WordStructureRecursionLeastGraph}[thm]{1,2,3}}
  \proofstepwidestar[1,2,3]{G\subseteq W\times X}{\rIE{\FormulaRefAuto{WordRecursionAmbientDef}[def],\FormulaRefAuto{WordRecursionClosedGraphSubset}[thm]{\rAEa{11}}}}
  \proofstepwidestar[1,2,3]{G\colon W\to X}{\FormulaRefAuto{UniqueValuedGraphFunction}[thm]{12,10}}
\end{tabproofwide}


\noindent\begin{minipage}{\linewidth}
\FormulaThmDeltaK[Rekursionsgleichungen des kleinsten Graphen]{%
\begin{aligned}[t]&\mathsf{WortStr}_A(W,e,s)\dsep x_0\in X\dsep r\colon X\times A\to X\vdash{}\\[-2pt]&\mathsf{WRec}_A(G;W,e,s;X,x_0,r)\end{aligned}
}{WordStructureRecursionGraphEquations}{\DeltaRow{Mengen}{A\dsep W\dsep e\dsep X\dsep x_0\dsep H\dsep K\dsep u\dsep v\dsep x\dsep y\dsep a\dsep b\dsep p\dsep z\dsep t\dsep c}\DeltaRow{Funktionen}{s\dsep r}}
\end{minipage}\par
\begin{tabproofwide}
  \proofstepwidestar[1]{\mathsf{WortStr}_A(W,e,s)}{\rA}
  \proofstepwidestar[2]{x_0\in X}{\rA}
  \proofstepwidestar[3]{r\colon X\times A\to X}{\rA}
  \proofstepwidestar[1,2,3]{G\colon W\to X}{\FormulaRefAuto{WordStructureRecursionGraphFunction}[thm]{1,2,3}}
  \proofstepwidestar[1,2,3]{C(G)\land\forall H\,(C(H)\rightarrow G\subseteq H)}{\FormulaRefAuto{WordStructureRecursionLeastGraph}[thm]{1,2,3}}
  \proofstepwidestar[1,2,3]{(e,x_0)\in G}{\FormulaRefAuto{WordRecursionClosedGraphBase}[thm]{\rAEa{5}}}
  \proofstepwidestar[1,2,3]{G(e)=x_0}{\FormulaRefAuto{F\colon A\to B\dsep(x,y)\in F\vdash F(x)=y}[thm]{4,6}}
  \proofstepwidestar[8]{u\in W}{\rA}
  \proofstepwidestar[9]{a\in A}{\rA}
  \proofstepwidestar[1,2,3,8]{G(u)\in X}{\FormulaRefAuto{F\colon A\to B\dsep x\in A\vdash F(x)\in B}[thm]{4,8}}
  \proofstepwidestar[1,2,3,8]{(u,G(u))\in G}{\FormulaRefAuto{F\colon A\to B\dsep x\in A\vdash(x,F(x))\in F}[thm]{4,8}}
  \proofstepwidestar[1,2,3,8,9]{(s(u,a),r(G(u),a))\in G}{\FormulaRefAuto{WordRecursionClosedGraphStep}[thm]{\rAEa{5},8,10,9,11}}
  \proofstepwidestar[1,2,3,8,9]{G(s(u,a))=r(G(u),a)}{\FormulaRefAuto{F\colon A\to B\dsep(x,y)\in F\vdash F(x)=y}[thm]{4,12}}
  \proofstepwidestar[1,2,3]{\forall u\in W\,\forall a\in A\;G(s(u,a))=r(G(u),a)}{\rUI{\rRI{8,\rUI{\rRI{9,13}}}}}
  \proofstepwidestar[1,2,3]{\mathsf{WRec}_A(G;W,e,s;X,x_0,r)}{\FormulaRefAuto{WordRecursionSpecificationDef}[def]{\rAI{4,\rAI{7,14}}}}
\end{tabproofwide}


Die Rekursionsgleichungen bestimmen denselben Anfangswert. Stimmen zwei
Lösungen an einer Stelle überein, so stimmen sie auch nach jedem
Anfügungsschritt überein. Die Wortinduktion überträgt diese Gleichheit
auf den ganzen Strukturträger.

\Needspace{40\baselineskip}
\noindent\begin{minipage}{\linewidth}
\FormulaThmDeltaK[Eindeutigkeit der Wortrekursion]{%
\begin{aligned}[t]&\mathsf{WortStr}_A(W,e,s)\dsep x_0\in X\dsep r\colon X\times A\to X\dsep{}\\[-2pt]&\mathsf{WRec}_A(F;W,e,s;X,x_0,r)\vdash F=G\end{aligned}
}{WordStructureRecursionUniqueness}{%
  \DeltaRow{Mengen}{A\dsep W\dsep e\dsep X\dsep x_0\dsep u\dsep v\dsep a}
  \DeltaRow{Funktionen}{s\dsep r}
  \DeltaRow{Funktion}{F}
  \DeltaRow{Prädikatsabkürzung}{P(u)\coloneqq F(u)=G(u)}
}
\end{minipage}\par
\begin{tabproofwide}
  \proofstepwidestar[1]{\mathsf{WortStr}_A(W,e,s)}{\rA}
  \proofstepwidestar[2]{x_0\in X}{\rA}
  \proofstepwidestar[3]{r\colon X\times A\to X}{\rA}
  \proofstepwidestar[4]{\mathsf{WRec}_A(F;W,e,s;X,x_0,r)}{\rA}
  \proofstepwidestar[4]{\begin{aligned}[t]&F\colon W\to X\land F(e)=x_0\land{}\\&\forall u\in W\,\forall a\in A\;F(s(u,a))=r(F(u),a)\end{aligned}}{\FormulaRefAuto{WordRecursionSpecificationDef}[def]{4}}
  \proofstepwidestar[1,2,3]{\begin{aligned}[t]&G\colon W\to X\land G(e)=x_0\land{}\\&\forall u\in W\,\forall a\in A\;G(s(u,a))=r(G(u),a)\end{aligned}}{\FormulaRefAuto{WordRecursionSpecificationDef}[def]{\FormulaRefAuto{WordStructureRecursionGraphEquations}[thm]{1,2,3}}}
  \proofstepwidestar[1,2,3,4]{P(e)}{\FormulaRefAuto{a=b\dsep c=b\vdash a=c}[thm]{\rAEa{\rAEb{5}},\rAEa{\rAEb{6}}}}
  \proofstepwidestar[8]{u\in W}{\rA}
  \proofstepwidestar[9]{a\in A}{\rA}
  \proofstepwidestar[10]{P(u)}{\rA}
  \proofstepwidestar[4,8,9]{F(s(u,a))=r(F(u),a)}{\rRE{9,\rUE{\rRE{8,\rUE{\rAEb{\rAEb{5}}}}}}}
  \proofstepwidestar[1,2,3,8,9]{G(s(u,a))=r(G(u),a)}{\rRE{9,\rUE{\rRE{8,\rUE{\rAEb{\rAEb{6}}}}}}}
  \proofstepwidestar[4,8,9,10]{F(s(u,a))=r(G(u),a)}{\rIE{10,11}}
  \proofstepwidestar[1,2,3,4,8,9,10]{P(s(u,a))}{\FormulaRefAuto{a=b\dsep c=b\vdash a=c}[thm]{13,12}}
  \proofstepwidestar[15]{v\in W}{\rA}
  \proofstepwidestar[1,2,3,4,15]{P(v)}{\FormulaRefAuto{WordStructureInductionRule}[thm]{1,7,[8,9,10]\vdots14,15}}
  \proofstepwidestar[1,2,3,4]{\forall v\in W\;F(v)=G(v)}{\rUI{\rRI{15,16}}}
  \proofstepwidestar[1,2,3,4]{F=G}{\FormulaRefAuto{FunctionExtensionality}[thm]{\rAEa{5},\rAEa{6},17}}
\end{tabproofwide}


\noindent\begin{minipage}{\linewidth}
\FormulaThmDeltaK[Rekursion in jeder freien Wortstruktur]{%
\begin{aligned}[t]&\mathsf{WortStr}_A(W,e,s)\dsep x_0\in X\dsep r\colon X\times A\to X\vdash{}\\[-2pt]&\exists!F\;\mathsf{WRec}_A(F;W,e,s;X,x_0,r)\end{aligned}
}{WordStructureRecursion}{\DeltaRow{Mengen}{A\dsep W\dsep e\dsep X\dsep x_0\dsep H\dsep K\dsep u\dsep v\dsep x\dsep y\dsep a\dsep b\dsep p\dsep z\dsep t\dsep c}\DeltaRow{Funktionen}{s\dsep r}\DeltaRow{Funktion}{F}}
\end{minipage}\par
\begin{tabproofwide}
  \proofstepwidestar[1]{\mathsf{WortStr}_A(W,e,s)}{\rA}
  \proofstepwidestar[2]{x_0\in X}{\rA}
  \proofstepwidestar[3]{r\colon X\times A\to X}{\rA}
  \proofstepwidestar[1,2,3]{\mathsf{WRec}_A(G;W,e,s;X,x_0,r)}{\FormulaRefAuto{WordStructureRecursionGraphEquations}[thm]{1,2,3}}
  \proofstepwidestar[5]{\mathsf{WRec}_A(F;W,e,s;X,x_0,r)}{\rA}
  \proofstepwidestar[1,2,3,5]{F=G}{\FormulaRefAuto{WordStructureRecursionUniqueness}[thm]{1,2,3,5}}
  \proofstepwidestar[1,2,3]{\forall F\;(\mathsf{WRec}_A(F;W,e,s;X,x_0,r)\rightarrow F=G)}{\rUI{\rRI{5,6}}}
  \proofstepwidestar[1,2,3]{\exists!F\;\mathsf{WRec}_A(F;W,e,s;X,x_0,r)}{\FormulaRefAuto{WordUniqueExistenceFromWitness}[thm]{4,7}}
\end{tabproofwide}


Dieser Beweis verwendet die Wortaxiome sowie die Mengen- und Funktionentheorie der früheren Bände. Für ein leeres Alphabet sind alle Schrittbedingungen leer; derselbe Beweis liefert die eindeutige Funktion mit Anfangswert \(x_0\).

\noindent\begin{minipage}{\linewidth}
\FormulaThmDeltaK[Rekursion vom Leerwort aus]{%
\begin{aligned}[t]&x_0\in X\dsep r\colon X\times A\to X\vdash{}\\[-2pt]&\exists!F\;\mathsf{WRec}_A(F;\FiniteWords A,\EmptyWord A,\sigma_A;X,x_0,r)\end{aligned}
}{FiniteWordRecursion}{\DeltaRow{Mengen}{A\dsep X\dsep x_0}\DeltaRow{Funktionen}{r\dsep F}}
\end{minipage}\par
\begin{tabproofwide}
  \proofstepwidestar[1]{x_0\in X}{\rA}
  \proofstepwidestar[2]{r\colon X\times A\to X}{\rA}
  \proofstepwidestar[]{\mathsf{WortStr}_A(\FiniteWords A,\EmptyWord A,\sigma_A)}{\FormulaRefAuto{WordStructureConcreteModel}[thm]}
  \proofstepwidestar[1,2]{\exists!F\;\mathsf{WRec}_A(F;\FiniteWords A,\EmptyWord A,\sigma_A;X,x_0,r)}{\FormulaRefAuto{WordStructureRecursion}[thm]{3,1,2}}
\end{tabproofwide}

\subsection{Kanonische Darstellung der Wortstruktur}

Wir konstruieren die Abbildungen in beiden Richtungen durch Rekursion. Die Rekursionseindeutigkeit zeigt, dass ihre Kompositionen die Identitäten sind. Damit folgt die Bijektivität aus Band~08.

\noindent\begin{minipage}{\linewidth}
\FormulaThmDeltaK[Zwei Lösungen derselben Wortrekursion sind gleich]{%
\begin{aligned}[t]&\mathsf{WortStr}_A(W,e,s)\dsep x_0\in X\dsep r\colon X\times A\to X\dsep{}\\[-2pt]&\mathsf{WRec}_A(F;W,e,s;X,x_0,r)\dsep \mathsf{WRec}_A(H;W,e,s;X,x_0,r)\vdash F=H\end{aligned}
}{WordRecursionTwoSolutionsEqual}{\DeltaRow{Mengen}{A\dsep W\dsep e\dsep V\dsep d\dsep X\dsep x_0\dsep u\dsep a}\DeltaRow{Funktionen}{s\dsep t\dsep r\dsep F\dsep H}}
\end{minipage}\par
\begin{tabproofwide}
  \proofstepwidestar[1]{\mathsf{WortStr}_A(W,e,s)}{\rA}
  \proofstepwidestar[2]{x_0\in X}{\rA}
  \proofstepwidestar[3]{r\colon X\times A\to X}{\rA}
  \proofstepwidestar[4]{\mathsf{WRec}_A(F;W,e,s;X,x_0,r)}{\rA}
  \proofstepwidestar[5]{\mathsf{WRec}_A(H;W,e,s;X,x_0,r)}{\rA}
  \proofstepwidestar[1,2,3]{\exists!K\;\mathsf{WRec}_A(K;W,e,s;X,x_0,r)}{\FormulaRefAuto{WordStructureRecursion}[thm]{1,2,3}}
  \proofstepwidestar[1,2,3,4,5]{F=H}{\FormulaRefAuto{\exists! x\,P(x),\; P(a),\; P(b) \vdash a = b}[thm]{6,4,5}}
\end{tabproofwide}


\noindent\begin{minipage}{\linewidth}
\FormulaThmDeltaK[Die Identität erfüllt die Wortrekursion]{%
\mathsf{WortStr}_A(W,e,s)\vdash \mathsf{WRec}_A(\Id_W;W,e,s;W,e,s)
}{WordRecursionIdentity}{\DeltaRow{Mengen}{A\dsep W\dsep e\dsep V\dsep d\dsep X\dsep x_0\dsep u\dsep a}\DeltaRow{Funktionen}{s\dsep t\dsep r\dsep F\dsep H}}
\end{minipage}\par
\begin{tabproofwide}
  \proofstepwidestar[1]{\mathsf{WortStr}_A(W,e,s)}{\rA}
  \proofstepwidestar[1]{e\in W\land s\colon W\times A\to W}{\FormulaRefAuto{WordStructureTypingAxiom}[ax]{1}}
  \proofstepwidestar[]{\Id_W\colon W\to W}{\FormulaRefAuto{IdA}[thm]}
  \proofstepwidestar[1]{\Id_W(e)=e}{\FormulaRefAuto{x\in A\vdash\Id_A(x)=x}[thm]{\rAEa{2}}}
  \proofstepwidestar[5]{u\in W}{\rA}
  \proofstepwidestar[6]{a\in A}{\rA}
  \proofstepwidestar[1,5,6]{s(u,a)\in W}{\FormulaRefAuto{WordStructureStepTyping}[thm]{1,5,6}}
  \proofstepwidestar[1,5,6]{\Id_W(s(u,a))=s(u,a)}{\FormulaRefAuto{x\in A\vdash\Id_A(x)=x}[thm]{7}}
  \proofstepwidestar[5]{\Id_W(u)=u}{\FormulaRefAuto{x\in A\vdash\Id_A(x)=x}[thm]{5}}
  \proofstepwidestar[1,5,6]{\Id_W(s(u,a))=s(\Id_W(u),a)}{\rIE{\FormulaRefAuto{a=b\vdash b=a}[thm]{9},8}}
  \proofstepwidestar[1]{\forall u\in W\,\forall a\in A\;\Id_W(s(u,a))=s(\Id_W(u),a)}{\rUI{\rRI{5,\rUI{\rRI{6,10}}}}}
  \proofstepwidestar[1]{\mathsf{WRec}_A(\Id_W;W,e,s;W,e,s)}{\FormulaRefAuto{WordRecursionSpecificationDef}[def]{\rAI{3,\rAI{4,11}}}}
\end{tabproofwide}


\noindent\begin{minipage}{\linewidth}
\FormulaThmDeltaK[Komposition von Rekursionsabbildungen]{%
\begin{aligned}[t]&\mathsf{WortStr}_A(W,e,s)\dsep \mathsf{WortStr}_A(V,d,t)\dsep{}\\[-2pt]&\mathsf{WRec}_A(F;W,e,s;V,d,t)\dsep \mathsf{WRec}_A(H;V,d,t;X,x_0,r)\vdash{}\\[-2pt]&\mathsf{WRec}_A(H\circ F;W,e,s;X,x_0,r)\end{aligned}
}{WordRecursionComposition}{\DeltaRow{Mengen}{A\dsep W\dsep e\dsep V\dsep d\dsep X\dsep x_0\dsep u\dsep a}\DeltaRow{Funktionen}{s\dsep t\dsep r\dsep F\dsep H}}
\end{minipage}\par
\begin{tabproofwide}
  \proofstepwidestar[1]{\mathsf{WortStr}_A(W,e,s)}{\rA}
  \proofstepwidestar[2]{\mathsf{WortStr}_A(V,d,t)}{\rA}
  \proofstepwidestar[3]{\mathsf{WRec}_A(F;W,e,s;V,d,t)}{\rA}
  \proofstepwidestar[4]{\mathsf{WRec}_A(H;V,d,t;X,x_0,r)}{\rA}
  \proofstepwidestar[3]{\begin{aligned}[t]&F\colon W\to V\land F(e)=d\land{}\\&\forall u\in W\,\forall a\in A\;F(s(u,a))=t(F(u),a)\end{aligned}}{\FormulaRefAuto{WordRecursionSpecificationDef}[def]{3}}
  \proofstepwidestar[4]{\begin{aligned}[t]&H\colon V\to X\land H(d)=x_0\land{}\\&\forall v\in V\,\forall a\in A\;H(t(v,a))=r(H(v),a)\end{aligned}}{\FormulaRefAuto{WordRecursionSpecificationDef}[def]{4}}
  \proofstepwidestar[3,4]{H\circ F\colon W\to X}{\FormulaRefAuto{CompositionDef}[def]{\rAEa{5},\rAEa{6}}}
  \proofstepwidestar[1]{e\in W\land s\colon W\times A\to W}{\FormulaRefAuto{WordStructureTypingAxiom}[ax]{1}}
  \proofstepwidestar[1,3,4]{(H\circ F)(e)=H(F(e))}{\FormulaRefAuto{x\in A\vdash(F\circ G)(x)=F(G(x))}[thm]{\rAEa{5},\rAEa{6},\rAEa{8}}}
  \proofstepwidestar[1,3,4]{(H\circ F)(e)=x_0}{\rIE{\rAEa{\rAEb{5}},\rAEa{\rAEb{6}},9}}
  \proofstepwidestar[11]{u\in W}{\rA}
  \proofstepwidestar[12]{a\in A}{\rA}
  \proofstepwidestar[1,11,12]{s(u,a)\in W}{\FormulaRefAuto{WordStructureStepTyping}[thm]{1,11,12}}
  \proofstepwidestar[3,11]{F(u)\in V}{\FormulaRefAuto{F\colon A\to B\dsep x\in A\vdash F(x)\in B}[thm]{\rAEa{5},11}}
  \proofstepwidestar[1,3,4,11,12]{(H\circ F)(s(u,a))=H(F(s(u,a)))}{\FormulaRefAuto{x\in A\vdash(F\circ G)(x)=F(G(x))}[thm]{\rAEa{5},\rAEa{6},13}}
  \proofstepwidestar[3,11,12]{F(s(u,a))=t(F(u),a)}{\rRE{12,\rUE{\rRE{11,\rUE{\rAEb{\rAEb{5}}}}}}}
  \proofstepwidestar[3,4,11,12]{H(t(F(u),a))=r(H(F(u)),a)}{\rRE{12,\rUE{\rRE{14,\rUE{\rAEb{\rAEb{6}}}}}}}
  \proofstepwidestar[3,4,11]{(H\circ F)(u)=H(F(u))}{\FormulaRefAuto{x\in A\vdash(F\circ G)(x)=F(G(x))}[thm]{\rAEa{5},\rAEa{6},11}}
  \proofstepwidestar[1,3,4,11,12]{(H\circ F)(s(u,a))=r((H\circ F)(u),a)}{\rIE{\FormulaRefAuto{a=b\vdash b=a}[thm]{18},\rIE{17,\rIE{16,15}}}}
  \proofstepwidestar[1,3,4]{\forall u\in W\,\forall a\in A\;(H\circ F)(s(u,a))=r((H\circ F)(u),a)}{\rUI{\rRI{11,\rUI{\rRI{12,19}}}}}
  \proofstepwidestar[1,3,4]{\mathsf{WRec}_A(H\circ F;W,e,s;X,x_0,r)}{\FormulaRefAuto{WordRecursionSpecificationDef}[def]{\rAI{7,\rAI{10,20}}}}
\end{tabproofwide}


\noindent\begin{minipage}{\linewidth}
\FormulaThmDeltaK[Rekursionsabbildungen zwischen Wortstrukturen sind Isomorphismen]{%
\begin{aligned}[t]&\mathsf{WortStr}_A(W,e,s)\dsep \mathsf{WortStr}_A(V,d,t)\dsep{}\\[-2pt]&\mathsf{WRec}_A(F;W,e,s;V,d,t)\vdash F\colon W\bij V\end{aligned}
}{WordStructureRecursorIsomorphism}{\DeltaRow{Mengen}{A\dsep W\dsep e\dsep V\dsep d\dsep X\dsep x_0\dsep u\dsep a}\DeltaRow{Funktionen}{s\dsep t\dsep r\dsep F\dsep H}}
\end{minipage}\par
\begin{tabproofwide}
  \proofstepwidestar[1]{\mathsf{WortStr}_A(W,e,s)}{\rA}
  \proofstepwidestar[2]{\mathsf{WortStr}_A(V,d,t)}{\rA}
  \proofstepwidestar[3]{\mathsf{WRec}_A(F;W,e,s;V,d,t)}{\rA}
  \proofstepwidestar[1]{e\in W\land s\colon W\times A\to W}{\FormulaRefAuto{WordStructureTypingAxiom}[ax]{1}}
  \proofstepwidestar[2]{d\in V\land t\colon V\times A\to V}{\FormulaRefAuto{WordStructureTypingAxiom}[ax]{2}}
  \proofstepwidestar[1,2]{\exists!H\;\mathsf{WRec}_A(H;V,d,t;W,e,s)}{\FormulaRefAuto{WordStructureRecursion}[thm]{2,\rAEa{4},\rAEb{4}}}
  \proofstepwidestar[1,2]{\exists H\;\mathsf{WRec}_A(H;V,d,t;W,e,s)}{\FormulaRefAuto{\exists! x\,P(x)\vdash \exists x\,P(x)}[thm]{6}}
  \proofstepwidestar[8]{\mathsf{WRec}_A(H;V,d,t;W,e,s)}{\rA}
  \proofstepwidestar[3]{F\colon W\to V}{\rAEa{\FormulaRefAuto{WordRecursionSpecificationDef}[def]{3}}}
  \proofstepwidestar[8]{H\colon V\to W}{\rAEa{\FormulaRefAuto{WordRecursionSpecificationDef}[def]{8}}}
  \proofstepwidestar[1,2,3,8]{\mathsf{WRec}_A(H\circ F;W,e,s;W,e,s)}{\FormulaRefAuto{WordRecursionComposition}[thm]{1,2,3,8}}
  \proofstepwidestar[1,2,3,8]{\mathsf{WRec}_A(F\circ H;V,d,t;V,d,t)}{\FormulaRefAuto{WordRecursionComposition}[thm]{2,1,8,3}}
  \proofstepwidestar[1]{\mathsf{WRec}_A(\Id_W;W,e,s;W,e,s)}{\FormulaRefAuto{WordRecursionIdentity}[thm]{1}}
  \proofstepwidestar[2]{\mathsf{WRec}_A(\Id_V;V,d,t;V,d,t)}{\FormulaRefAuto{WordRecursionIdentity}[thm]{2}}
  \proofstepwidestar[1,2,3,8]{H\circ F=\Id_W}{\FormulaRefAuto{WordRecursionTwoSolutionsEqual}[thm]{1,\rAEa{4},\rAEb{4},11,13}}
  \proofstepwidestar[1,2,3,8]{F\circ H=\Id_V}{\FormulaRefAuto{WordRecursionTwoSolutionsEqual}[thm]{2,\rAEa{5},\rAEb{5},12,14}}
  \proofstepwidestar[1,2,3,8]{F\colon W\bij V}{\FormulaRefAuto{MutuallyInverseFunctionBijection}[thm]{10,9,16,15}}
  \proofstepwidestar[1,2,3]{F\colon W\bij V}{\rEE{7,8,17}}
\end{tabproofwide}


\noindent\begin{minipage}{\linewidth}
\FormulaThmDeltaK[Kanonische Darstellung jeder Wortstruktur]{%
\begin{aligned}[t]&\mathsf{WortStr}_A(W,e,s)\vdash\exists!h\;\Bigl(h\colon\FiniteWords A\bij W\land{}\\[-2pt]&\qquad\mathsf{WRec}_A(h;\FiniteWords A,\EmptyWord A,\sigma_A;W,e,s)\Bigr)\end{aligned}
}{WordStructureCanonicalIsomorphism}{\DeltaRow{Mengen}{A\dsep W\dsep e}\DeltaRow{Funktionen}{s\dsep h\dsep k}}
\end{minipage}\par
\begin{tabproofwide}
  \proofstepwidestar[1]{\mathsf{WortStr}_A(W,e,s)}{\rA}
  \proofstepwidestar[1]{e\in W\land s\colon W\times A\to W}{\FormulaRefAuto{WordStructureTypingAxiom}[ax]{1}}
  \proofstepwidestar[]{\mathsf{WortStr}_A(\FiniteWords A,\EmptyWord A,\sigma_A)}{\FormulaRefAuto{WordStructureConcreteModel}[thm]}
  \proofstepwidestar[1]{\exists!h\;\mathsf{WRec}_A(h;\FiniteWords A,\EmptyWord A,\sigma_A;W,e,s)}{\FormulaRefAuto{FiniteWordRecursion}[thm]{\rAEa{2},\rAEb{2}}}
  \proofstepwidestar[1]{\exists h\;\mathsf{WRec}_A(h;\FiniteWords A,\EmptyWord A,\sigma_A;W,e,s)}{\FormulaRefAuto{\exists! x\,P(x)\vdash \exists x\,P(x)}[thm]{4}}
  \proofstepwidestar[6]{\mathsf{WRec}_A(h;\FiniteWords A,\EmptyWord A,\sigma_A;W,e,s)}{\rA}
  \proofstepwidestar[1,6]{h\colon\FiniteWords A\bij W}{\FormulaRefAuto{WordStructureRecursorIsomorphism}[thm]{3,1,6}}
  \proofstepwidestar[1,6]{\begin{aligned}[t]&h\colon\FiniteWords A\bij W\land{}\\[-2pt]&\quad\mathsf{WRec}_A(h;\FiniteWords A,\EmptyWord A,\sigma_A;W,e,s)\end{aligned}}{\rAI{7,6}}
  \proofstepwidestar[9]{\begin{aligned}[t]&k\colon\FiniteWords A\bij W\land{}\\[-2pt]&\quad\mathsf{WRec}_A(k;\FiniteWords A,\EmptyWord A,\sigma_A;W,e,s)\end{aligned}}{\rA}
  \proofstepwidestar[1,6,9]{k=h}{\FormulaRefAuto{\exists! x\,P(x),\; P(a),\; P(b) \vdash a = b}[thm]{4,\rAEb{9},6}}
  \proofstepwidestar[1,6]{\begin{aligned}[t]&\forall k\;\bigl((k\colon\FiniteWords A\bij W\land{}\\[-2pt]&\quad\mathsf{WRec}_A(k;\FiniteWords A,\EmptyWord A,\sigma_A;W,e,s))\rightarrow k=h\bigr)\end{aligned}}{\rUI{\rRI{9,10}}}
  \proofstepwidestar[1,6]{\begin{aligned}[t]&\exists!h\;\bigl(h\colon\FiniteWords A\bij W\land{}\\[-2pt]&\quad\mathsf{WRec}_A(h;\FiniteWords A,\EmptyWord A,\sigma_A;W,e,s)\bigr)\end{aligned}}{\FormulaRefAuto{WordUniqueExistenceFromWitness}[thm]{8,11}}
  \proofstepwidestar[1]{\begin{aligned}[t]&\exists!h\;\bigl(h\colon\FiniteWords A\bij W\land{}\\[-2pt]&\quad\mathsf{WRec}_A(h;\FiniteWords A,\EmptyWord A,\sigma_A;W,e,s)\bigr)\end{aligned}}{\rEE{5,6,12}}
\end{tabproofwide}



\section{Konkatenation}



\FormulaDefDeltaK[Konkatenation endlicher Wörter]{%
  u,v\in\FiniteWords{A}\vdash
  \WordConcat{u}{v}\coloneqq
  u\cup
  \bigl\{(\WordLength{u}+j,v(j))\bigm|
    j\in\NatSeg{\WordLength{v}}\bigr\}%
}{WordConcatenationDef}{%
  \DeltaRow{Mengen}{A\dsep u\dsep v\dsep j\dsep p}
  \DeltaRow{\textbf{Neues Symbol}}{}
  \DeltaPrem{Wortkonkatenation}{\WordConcat{u}{v}}
}

\subsection{Die primitive Anfügung als Spezialfall}

Das Anhängen eines Buchstabens ist auf der Mengenebene genau die Adjunktion
eines neuen Elements. Neu ist das markierte Vorkommen; der Buchstabe selbst
darf im bisherigen Wort bereits auftreten.

\noindent\begin{minipage}{\linewidth}
\FormulaThmDeltaK[Anfügung ist Adjunktion des letzten Vorkommens]{%
  u\in\FiniteWords A\dsep a\in A\vdash
  \WordConcat u{\LetterWord a}=u\cup\{(\WordLength u,a)\}
}{WordAppendOccurrenceEquation}{\DeltaRow{Mengen}{A\dsep u\dsep a}}
\end{minipage}\par

\begin{tabproofwide}
  \proofstepwidestar[1]{u\in\FiniteWords A}{\rA}
  \proofstepwidestar[2]{a\in A}{\rA}
  \proofstepwidestar[2]{\WordLength{\LetterWord a}=1}{%
    \FormulaRefAuto{LetterWordLengthOne}[thm]{2}}
  \proofstepwidestar[]{\NatSeg1=\{0\}}{\FormulaRefAuto{PeanoNatSegOneEquality}[thm]}
  \proofstepwidestar[2]{\LetterWord a(0)=a}{\FormulaRefAuto{LetterWordValueZero}[thm]{2}}
  \proofstepwidestar[1]{\WordLength u\in\mathbb N}{\FormulaRefAuto{FiniteWordLengthNatural}[thm]{1}}
  \proofstepwidestar[1]{\WordLength u+0=\WordLength u}{\FormulaRefAuto{PeanoAddRightZero}[thm]{6}}
  \proofstepwidestar[1,2]{\begin{aligned}[t]
    &\{(\WordLength u+j,\LetterWord a(j))\mid
      j\in\NatSeg{\WordLength{\LetterWord a}}\}\\[-2pt]
    &\qquad=\{(\WordLength u,a)\}
  \end{aligned}}{\rIE{4,3,5,7,\FormulaRefAuto{B27TermImageSingleton}[thm]}}
  \proofstepwidestar[2]{\LetterWord a\in\NonemptyWords A}{%
    \FormulaRefAuto{LetterWordNonempty}[thm]{2}}
  \proofstepwidestar[2]{\LetterWord a\in\FiniteWords A}{%
    \FormulaRefAuto{NonemptyWordFinite}[thm]{9}}
  \proofstepwidestar[1,2]{\WordConcat u{\LetterWord a}=u\cup\{(\WordLength u,a)\}}{%
    \rIE{8,\FormulaRefAuto{WordConcatenationDef}[def]{1,10}}}
\end{tabproofwide}

\noindent\begin{minipage}{\linewidth}
\FormulaThmDeltaK[Das angefügte Vorkommen ist neu]{%
  u\in\FiniteWords A\vdash(\WordLength u,a)\notin u
}{WordAppendOccurrenceFresh}{\DeltaRow{Mengen}{A\dsep u\dsep a}}
\end{minipage}\par
\begin{tabproofwide}
  \proofstepwidestar[1]{u\in\FiniteWords A}{\rA}
  \proofstepwidestar[1]{u\colon\NatSeg{\WordLength u}\to A}{%
    \FormulaRefAuto{FiniteWordTypingAtLength}[thm]{1}}
  \proofstepwidestar[1]{\WordLength u\in\mathbb N}{%
    \FormulaRefAuto{FiniteWordLengthNatural}[thm]{1}}
  \proofstepwidestar[1]{\WordLength u\notin\NatSeg{\WordLength u}}{%
    \FormulaRefAuto{PeanoNatSegEndpointNotIn}[thm]{3}}
  \proofstepwidestar[5]{(\WordLength u,a)\in u}{\rA}
  \proofstepwidestar[1,5]{\WordLength u\in\NatSeg{\WordLength u}}{%
    \FormulaRefAuto{F\colon A\to B\dsep(x,y)\in F\vdash x\in A}[thm]{2,5}}
  \proofstepwidestar[1,5]{\bot}{\rBI{6,4}}
  \proofstepwidestar[1]{(\WordLength u,a)\notin u}{\rCI{5,7}}
\end{tabproofwide}

\FormulaThmDeltaK[Endlichkeit der Anfügung aus dem Adjunktionsaxiom]{%
  u\in\FiniteWords A\dsep a\in A\vdash\Finite{\WordConcat u{\LetterWord a}}
}{WordAppendFiniteByAdjunction}{\DeltaRow{Mengen}{A\dsep u\dsep a}}
\begin{tabproofwide}
  \proofstepwidestar[1]{u\in\FiniteWords A}{\rA}
  \proofstepwidestar[2]{a\in A}{\rA}
  \proofstepwidestar[1]{\Finite u}{\FormulaRefAuto{FiniteWordUnderlyingFinite}[thm]{1}}
  \proofstepwidestar[1]{(\WordLength u,a)\notin u}{\FormulaRefAuto{WordAppendOccurrenceFresh}[thm]{1}}
  \proofstepwidestar[1]{\Finite{u\cup\{(\WordLength u,a)\}}}{%
    \FormulaRefAuto{FiniteAdjunctionAxiom}[ax]{4,3}}
  \proofstepwidestar[1,2]{\WordConcat u{\LetterWord a}=u\cup\{(\WordLength u,a)\}}{%
    \FormulaRefAuto{WordAppendOccurrenceEquation}[thm]{1,2}}
  \proofstepwidestar[1,2]{\Finite{\WordConcat u{\LetterWord a}}}{\rIE{6,5}}
\end{tabproofwide}


\noindent\begin{minipage}{\linewidth}
\FormulaThmDeltaK[Die primitive Anfügung ist Konkatenation mit einem Buchstabenwort]{%
  u\in\FiniteWords A\dsep a\in A\vdash
    \sigma_A(u,a)=\WordConcat u{\LetterWord a}
}{WordModelSuccessorConcatenation}{\DeltaRow{Mengen}{A\dsep u\dsep a}}
\end{minipage}\par

\begin{tabproofwide}
  \proofstepwidestar[1]{u\in\FiniteWords A}{\rA}
  \proofstepwidestar[2]{a\in A}{\rA}
  \proofstepwidestar[1,2]{\sigma_A(u,a)=u\cup\{(\WordLength u,a)\}}{%
    \FormulaRefAuto{EarlyWordSuccessorEquation}[thm]{1,2}}
  \proofstepwidestar[1,2]{\WordConcat u{\LetterWord a}=u\cup\{(\WordLength u,a)\}}{%
    \FormulaRefAuto{WordAppendOccurrenceEquation}[thm]{1,2}}
  \proofstepwidestar[1,2]{\sigma_A(u,a)=\WordConcat u{\LetterWord a}}{%
    \FormulaRefAuto{a=b\dsep c=b\vdash a=c}[thm]{3,4}}
\end{tabproofwide}


\subsection{Koordinaten und Typisierung der Konkatenation}

\FormulaThmDeltaK[Elementkriterium der Wortkonkatenation]{%
  \begin{aligned}[t]
    &u,v\in\FiniteWords{A}\vdash{}\\[-2pt]
    &\Bigl(p\in\WordConcat{u}{v}\leftrightarrow
      \Bigl(p\in u\lor
        \exists j\in\NatSeg{\WordLength v}\,
          p=(\WordLength u+j,v(j))\Bigr)\Bigr)
  \end{aligned}%
}{WordConcatenationMembership}{%
  \DeltaRow{Mengen}{A\dsep u\dsep v\dsep j\dsep p}
}
\begin{tabproofwide}
  \proofstepwidestar[1]{u,v\in\FiniteWords{A}}{\rA}
  \proofstepwidestar[1]{\WordConcat{u}{v}=u\cup\{(\WordLength{u}+j,v(j))\mid j\in\NatSeg{\WordLength{v}}\}}{%
      \FormulaRefAuto{WordConcatenationDef}[def]{1}}
  \proofstepwidestar[1]{%
    \begin{aligned}[t]
      &u\cup\{(\WordLength u+j,v(j))\mid
        j\in\NatSeg{\WordLength v}\}\\[-2pt]
      &\qquad=\WordConcat{u}{v}
    \end{aligned}}{%
    \FormulaRefAuto{a=b\vdash b=a}[thm]{%
      2}}
  \proofstepwidestar[]{%
    \begin{aligned}[t]
      &p\in u\cup\{(\WordLength u+j,v(j))\mid
        j\in\NatSeg{\WordLength v}\}\\[-2pt]
      &\quad\leftrightarrow p\in u\lor
        p\in\{(\WordLength u+j,v(j))\mid
          j\in\NatSeg{\WordLength v}\}
    \end{aligned}}{%
    \FormulaRefAuto{z\in A\cup B\eqvdash z\in A\lor z\in B}[thm]}
  \proofstepwidestar[]{%
    \begin{aligned}[t]
      &p\in\{(\WordLength u+j,v(j))\mid
        j\in\NatSeg{\WordLength v}\}\\[-2pt]
      &\quad\leftrightarrow
        \exists j\in\NatSeg{\WordLength v}\,
          p=(\WordLength u+j,v(j))
    \end{aligned}}{%
    \FormulaRefAuto{ShiftedWordGraphDef}[def]}
  \proofstepwidestar[]{%
    \begin{aligned}[t]
      &p\in u\cup\{(\WordLength u+j,v(j))\mid
        j\in\NatSeg{\WordLength v}\}\\[-2pt]
      &\quad\leftrightarrow p\in u\lor
        \exists j\in\NatSeg{\WordLength v}\,
          p=(\WordLength u+j,v(j))
    \end{aligned}}{\rLRS{4,5}}
  \proofstepwidestar[1]{%
    \begin{aligned}[t]
      &p\in\WordConcat{u}{v}\leftrightarrow p\in u\lor{}\\[-2pt]
      &\qquad\exists j\in\NatSeg{\WordLength v}\,
        p=(\WordLength u+j,v(j))
    \end{aligned}}{\rIE{3,6}}
\end{tabproofwide}



\FormulaThmDeltaK[Disjunktheit der Indexblöcke einer Konkatenation]{%
  u,v\in\FiniteWords{A}\dsep
    j\in\NatSeg{\WordLength v}\vdash
    \WordLength u+j\notin\NatSeg{\WordLength u}%
}{WordConcatenationIndexBlocksDisjoint}{%
  \DeltaRow{Mengen}{A\dsep u\dsep v\dsep j}
}
\begin{tabproofwide}
  \proofstepwidestar[1]{u,v\in\FiniteWords{A}}{\rA}
  \proofstepwidestar[2]{j\in\NatSeg{\WordLength v}}{\rA}
  \proofstepwidestar[1]{u\in\FiniteWords{A}}{\rAEa{1}}
  \proofstepwidestar[1]{v\in\FiniteWords{A}}{\rAEb{1}}
  \proofstepwidestar[1]{\WordLength u\in\mathbb N}{%
    \FormulaRefAuto{FiniteWordLengthNatural}[thm]{3}}
  \proofstepwidestar[1]{\WordLength v\in\mathbb N}{%
    \FormulaRefAuto{FiniteWordLengthNatural}[thm]{4}}
  \proofstepwidestar[1,2]{j\in\mathbb N}{%
    \FormulaRefAuto{PeanoNatSegElementNatural}[thm]{6,2}}
  \proofstepwidestar[1,2]{%
    \WordLength u+j\notin\NatSeg{\WordLength u}}{%
    \FormulaRefAuto{PeanoTailAvoidsInitialSegmentLeft}[thm]{5,7}}
\end{tabproofwide}

\FormulaThmDeltaK[Typisierung des Anfangsblocks einer Konkatenation]{%
  u,v\in\FiniteWords{A}\vdash
  u\subseteq\NatSeg{\WordLength u+\WordLength v}\times A%
}{WordConcatenationInitialGraphSubset}{%
  \DeltaRow{Mengen}{A\dsep a\dsep i\dsep p\dsep u\dsep v}
}
\begin{tabproofwide}
  \proofstepwidestar[1]{u,v\in\FiniteWords{A}}{\rA}
  \proofstepwidestar[1]{u\in\FiniteWords{A}}{\rAEa{1}}
  \proofstepwidestar[1]{v\in\FiniteWords{A}}{\rAEb{1}}
  \proofstepwidestar[1]{%
    u\colon\NatSeg{\WordLength u}\to A}{%
    \FormulaRefAuto{FiniteWordTypingAtLength}[thm]{2}}
  \proofstepwidestar[1]{%
    u\subseteq\NatSeg{\WordLength u}\times A}{%
    \FormulaRefAuto{FunctionGraphSubsetProduct}[thm]{4}}
  \proofstepwidestar[1]{\WordLength u\in\mathbb N}{%
    \FormulaRefAuto{FiniteWordLengthNatural}[thm]{2}}
  \proofstepwidestar[1]{\WordLength v\in\mathbb N}{%
    \FormulaRefAuto{FiniteWordLengthNatural}[thm]{3}}
  \proofstepwidestar[1]{%
    \NatSeg{\WordLength u}\subseteq
      \NatSeg{\WordLength u+\WordLength v}}{%
    \FormulaRefAuto{PeanoNatSegSubsetAddRight}[thm]{6,7}}
  \proofstepwidestar[1]{%
    \NatSeg{\WordLength u}\times A\subseteq
      \NatSeg{\WordLength u+\WordLength v}\times A}{%
    \FormulaRefAuto{CartesianProductMonotoneLeft}[thm]{8}}
  \proofstepwidestar[1]{%
    u\subseteq\NatSeg{\WordLength u+\WordLength v}\times A}{%
    \FormulaRefAuto{A\subseteq B\dsep B\subseteq C
      \vdash A\subseteq C}[thm]{5,9}}
\end{tabproofwide}

\noindent\begin{minipage}{\linewidth}
\FormulaThmDeltaK[Typisierung des verschobenen Konkatenationsblocks]{%
  \begin{aligned}[t]
    &\text{(i)}\;u,v\in\FiniteWords A\dsep
      j\in\NatSeg{\WordLength v}\vdash{}\\[-2pt]
    &\qquad(\WordLength u+j,v(j))\in
      \NatSeg{\WordLength u+\WordLength v}\times A\\[-2pt]
    &\text{(ii)}\;u,v\in\FiniteWords A\vdash{}\\[-2pt]
    &\qquad\{(\WordLength u+j,v(j))\mid
      j\in\NatSeg{\WordLength v}\}
      \subseteq\NatSeg{\WordLength u+\WordLength v}\times A
  \end{aligned}%
}{WordConcatenationShiftedGraphSubset}{%
  \DeltaRow{Mengen}{A\dsep j\dsep p\dsep u\dsep v}
}
\end{minipage}\par
\begin{tabproofsplitwide}
  \proofpartwideindR[Typisierung eines verschobenen Graphpunktes]{%
    \begin{aligned}[t]
      &u,v\in\FiniteWords A\dsep j\in\NatSeg{\WordLength v}\vdash{}\\[-2pt]
      &(\WordLength u+j,v(j))\in
        \NatSeg{\WordLength u+\WordLength v}\times A
    \end{aligned}
  }{u,v\in\FiniteWords A\dsep j\in\NatSeg{\WordLength v}\vdash
    (\WordLength u+j,v(j))\in
      \NatSeg{\WordLength u+\WordLength v}\times A}%
    [WordConcatenationShiftedGraphPoint]
    \proofstepwidestar[1]{u,v\in\FiniteWords A}{\rA}
    \proofstepwidestar[2]{j\in\NatSeg{\WordLength v}}{\rA}
    \proofstepwidestar[1]{u\in\FiniteWords A}{\rAEa{1}}
    \proofstepwidestar[1]{v\in\FiniteWords A}{\rAEb{1}}
    \proofstepwidestar[1]{\WordLength u\in\mathbb N}{%
      \FormulaRefAuto{FiniteWordLengthNatural}[thm]{3}}
    \proofstepwidestar[1]{\WordLength v\in\mathbb N}{%
      \FormulaRefAuto{FiniteWordLengthNatural}[thm]{4}}
    \proofstepwidestar[1]{v\colon\NatSeg{\WordLength v}\to A}{%
      \FormulaRefAuto{FiniteWordTypingAtLength}[thm]{4}}
    \proofstepwidestar[1,2]{v(j)\in A}{%
      \FormulaRefAuto{F\colon A\to B\dsep x\in A\vdash F(x)\in B}[thm]{7,2}}
    \proofstepwidestar[1,2]{%
      \WordLength u+j\in\NatSeg{\WordLength u+\WordLength v}}{%
      \FormulaRefAuto{PeanoNatSegShiftMembership}[thm]{5,6,2}}
    \proofstepwidestar[1,2]{%
      (\WordLength u+j,v(j))\in
        \NatSeg{\WordLength u+\WordLength v}\times A}{%
      \FormulaRefAuto{a\in A\dsep b\in B\vdash(a,b)\in A\times B}[thm]{9,8}}
  \closeproofpartwideindR

  \proofpartwideindR[Typisierung des gesamten verschobenen Blocks]{%
    \begin{aligned}[t]
      &u,v\in\FiniteWords A\vdash{}\\[-2pt]
      &\{(\WordLength u+j,v(j))\mid j\in\NatSeg{\WordLength v}\}
        \subseteq\NatSeg{\WordLength u+\WordLength v}\times A
    \end{aligned}
  }{u,v\in\FiniteWords A\vdash
    \{(\WordLength u+j,v(j))\mid j\in\NatSeg{\WordLength v}\}
      \subseteq\NatSeg{\WordLength u+\WordLength v}\times A}%
    [WordConcatenationShiftedGraphSubsetPart]
    \proofstepwidestar[1]{u,v\in\FiniteWords A}{\rA}
    \proofstepwidestar[2]{%
      p\in\{(\WordLength u+j,v(j))\mid j\in\NatSeg{\WordLength v}\}}{\rA}
    \proofstepwidestar[2]{%
      \exists j\in\NatSeg{\WordLength v}\,p=(\WordLength u+j,v(j))}{%
      \FormulaRefAuto{ShiftedWordGraphDef}[def]{2}}
    \proofstepwidestar[4]{%
      j\in\NatSeg{\WordLength v}\land p=(\WordLength u+j,v(j))}{\rA}
    \proofstepwidestar[4]{j\in\NatSeg{\WordLength v}}{\rAEa{4}}
    \proofstepwidestar[4]{p=(\WordLength u+j,v(j))}{\rAEb{4}}
    \proofstepwidestar[1,4]{%
      (\WordLength u+j,v(j))\in
        \NatSeg{\WordLength u+\WordLength v}\times A}{%
      \FormulaRefAuto{WordConcatenationShiftedGraphPoint}[thm]{1,5}}
    \proofstepwidestar[1,4]{%
      p\in\NatSeg{\WordLength u+\WordLength v}\times A}{\rIE{6,7}}
    \proofstepwidestar[1,2]{%
      p\in\NatSeg{\WordLength u+\WordLength v}\times A}{\rEE{3,4,8}}
    \proofstepwidestar[1]{%
      \{(\WordLength u+j,v(j))\mid j\in\NatSeg{\WordLength v}\}
        \subseteq\NatSeg{\WordLength u+\WordLength v}\times A}{%
      \FormulaRefAuto{SubsetIntroductionRule}[thm]{[2]\vdots 9}}
  \closeproofpartwideindR
\end{tabproofsplitwide}


\FormulaThmDeltaK[Totalität des Konkatenationsgraphen]{%
  \begin{aligned}[t]
    &u,v\in\FiniteWords{A}\dsep
      q\in\NatSeg{\WordLength u+\WordLength v}\vdash{}\\[-2pt]
    &\exists a\in A\,((q,a)\in\WordConcat{u}{v})
  \end{aligned}%
}{WordConcatenationGraphTotal}{%
  \DeltaRow{Mengen}{A\dsep a\dsep j\dsep q\dsep u\dsep v}
}
\begin{tabproofsplitwide}
  \proofpartwideindR[Fall 1: Index im Anfangsblock]{%
    \begin{aligned}[t]
      &u,v\in\FiniteWords A\dsep q\in\NatSeg{\WordLength u}\vdash{}\\[-2pt]
      &\exists a\in A\,((q,a)\in\WordConcat uv)
    \end{aligned}
  }{u,v\in\FiniteWords A\dsep q\in\NatSeg{\WordLength u}\vdash
    \exists a\in A\,((q,a)\in\WordConcat uv)}%
    [WordConcatenationGraphTotalInitial]
    \proofstepwidestar[1]{u,v\in\FiniteWords A}{\rA}
    \proofstepwidestar[2]{q\in\NatSeg{\WordLength u}}{\rA}
    \proofstepwidestar[1]{u\colon\NatSeg{\WordLength u}\to A}{%
      \FormulaRefAuto{FiniteWordTypingAtLength}[thm]{\rAEa{1}}}
    \proofstepwidestar[1,2]{u(q)\in A}{%
      \FormulaRefAuto{F\colon A\to B\dsep x\in A\vdash F(x)\in B}[thm]{3,2}}
    \proofstepwidestar[1,2]{(q,u(q))\in u}{%
      \FormulaRefAuto{F\colon A\to B\dsep x\in A\vdash
        (x,F(x))\in F}[thm]{3,2}}
    \proofstepwidestar[1,2]{%
      (q,u(q))\in u\lor
      \exists j\in\NatSeg{\WordLength v}\,
        (q,u(q))=(\WordLength u+j,v(j))}{\rOIa{5}}
    \proofstepwidestar[1,2]{(q,u(q))\in\WordConcat uv}{%
      \FormulaRefAuto{WordConcatenationMembership}[thm]{1,6}}
    \proofstepwidestar[1,2]{%
      \exists a\in A\,((q,a)\in\WordConcat uv)}{\rEI{\rAI{4,7}}}
  \closeproofpartwideindR

  \proofpartwideindR[Fall 2: Index im verschobenen Block]{%
    \begin{aligned}[t]
      &u,v\in\FiniteWords A\dsep
        \exists j\in\NatSeg{\WordLength v}\,(q=\WordLength u+j)
        \vdash{}\\[-2pt]
      &\exists a\in A\,((q,a)\in\WordConcat uv)
    \end{aligned}
  }{u,v\in\FiniteWords A\dsep
    \exists j\in\NatSeg{\WordLength v}\,(q=\WordLength u+j)\vdash
    \exists a\in A\,((q,a)\in\WordConcat uv)}%
    [WordConcatenationGraphTotalShifted]
    \proofstepwidestar[1]{u,v\in\FiniteWords A}{\rA}
    \proofstepwidestar[2]{%
      \exists j\in\NatSeg{\WordLength v}\,(q=\WordLength u+j)}{\rA}
    \proofstepwidestar[1]{v\colon\NatSeg{\WordLength v}\to A}{%
      \FormulaRefAuto{FiniteWordTypingAtLength}[thm]{\rAEb{1}}}
    \proofstepwidestar[4]{%
      j\in\NatSeg{\WordLength v}\land q=\WordLength u+j}{\rA}
    \proofstepwidestar[4]{j\in\NatSeg{\WordLength v}}{\rAEa{4}}
    \proofstepwidestar[4]{q=\WordLength u+j}{\rAEb{4}}
    \proofstepwidestar[1,4]{v(j)\in A}{%
      \FormulaRefAuto{F\colon A\to B\dsep x\in A\vdash F(x)\in B}[thm]{3,5}}
    \proofstepwidestar[4]{%
      (q,v(j))=(\WordLength u+j,v(j))}{\rIE{6,\rII}}
    \proofstepwidestar[4]{%
      \exists k\in\NatSeg{\WordLength v}\,
        (q,v(j))=(\WordLength u+k,v(k))}{\rEI{\rAI{5,8}}}
    \proofstepwidestar[4]{%
      (q,v(j))\in u\lor
      \exists k\in\NatSeg{\WordLength v}\,
        (q,v(j))=(\WordLength u+k,v(k))}{\rOIb{9}}
    \proofstepwidestar[1,4]{(q,v(j))\in\WordConcat uv}{%
      \FormulaRefAuto{WordConcatenationMembership}[thm]{1,10}}
    \proofstepwidestar[1,4]{%
      \exists a\in A\,((q,a)\in\WordConcat uv)}{\rEI{\rAI{7,11}}}
    \proofstepwidestar[1,2]{%
      \exists a\in A\,((q,a)\in\WordConcat uv)}{\rEE{2,4,12}}
  \closeproofpartwideindR

  \proofpartwide{Zusammenführung der beiden Fälle}
    \proofstepwidestar[1]{u,v\in\FiniteWords A}{\rA}
    \proofstepwidestar[2]{q\in\NatSeg{\WordLength u+\WordLength v}}{\rA}
    \proofstepwidestar[1]{\WordLength u\in\mathbb N}{%
      \FormulaRefAuto{FiniteWordLengthNatural}[thm]{\rAEa{1}}}
    \proofstepwidestar[1]{\WordLength v\in\mathbb N}{%
      \FormulaRefAuto{FiniteWordLengthNatural}[thm]{\rAEb{1}}}
    \proofstepwidestar[1,2]{%
      q\in\NatSeg{\WordLength u}\lor
      \exists j\in\NatSeg{\WordLength v}\,(q=\WordLength u+j)}{%
      \FormulaRefAuto{PeanoNatSegAddMembership}[thm]{3,4,2}}
    \proofstepwidestar[6]{q\in\NatSeg{\WordLength u}}{\rA}
    \proofstepwidestar[1,6]{%
      \exists a\in A\,((q,a)\in\WordConcat uv)}{%
      \FormulaRefAuto{WordConcatenationGraphTotalInitial}[thm]{1,6}}
    \proofstepwidestar[8]{%
      \exists j\in\NatSeg{\WordLength v}\,(q=\WordLength u+j)}{\rA}
    \proofstepwidestar[1,8]{%
      \exists a\in A\,((q,a)\in\WordConcat uv)}{%
      \FormulaRefAuto{WordConcatenationGraphTotalShifted}[thm]{1,8}}
    \proofstepwidestar[1,2]{%
      \exists a\in A\,((q,a)\in\WordConcat uv)}{\rOE{5,6,7,8,9}}
  \closeproofpartwide
\end{tabproofsplitwide}


\FormulaThmDeltaK[Funktionalität des Konkatenationsgraphen]{%
  \begin{aligned}[t]
    &u,v\in\FiniteWords{A}\dsep
      q\in\NatSeg{\WordLength u+\WordLength v}\dsep a,b\in A\dsep{}\\[-2pt]
    &(q,a)\in\WordConcat{u}{v}\dsep
      (q,b)\in\WordConcat{u}{v}\vdash a=b
  \end{aligned}%
}{WordConcatenationGraphFunctional}{%
  \DeltaRow{Mengen}{A\dsep a\dsep b\dsep j\dsep k\dsep q\dsep u\dsep v}
}
\begin{tabproofsplitwide}
  \proofpartwideindR[Beide Paare im Anfangsblock]{%
    \begin{aligned}[t]
      &u,v\in\FiniteWords{A}\dsep
        (q,a)\in u\dsep(q,b)\in u\vdash a=b
    \end{aligned}
  }{u,v\in\FiniteWords{A}\dsep
    (q,a)\in u\dsep(q,b)\in u\vdash a=b}%
    [WordConcatenationInitialBlockFunctional]
    \proofstepwidestar[1]{u,v\in\FiniteWords{A}}{\rA}
    \proofstepwidestar[2]{(q,a)\in u}{\rA}
    \proofstepwidestar[3]{(q,b)\in u}{\rA}
    \proofstepwidestar[1]{u\colon\NatSeg{\WordLength u}\to A}{%
      \FormulaRefAuto{FiniteWordTypingAtLength}[thm]{\rAEa{1}}}
    \proofstepwidestar[1,2]{u(q)=a}{%
      \FormulaRefAuto{F\colon A\to B\dsep (x,y)\in F
        \vdash F(x)=y}[thm]{4,2}}
    \proofstepwidestar[1,3]{u(q)=b}{%
      \FormulaRefAuto{F\colon A\to B\dsep (x,y)\in F
        \vdash F(x)=y}[thm]{4,3}}
    \proofstepwidestar[1,2,3]{a=b}{%
      \FormulaRefAuto{a=b\dsep a=c\vdash b=c}[thm]{5,6}}
  \closeproofpartwideindR

  \proofpartwideindR[Gemischte Blöcke]{%
    \begin{aligned}[t]
      &u,v\in\FiniteWords{A}\dsep(q,a)\in u\dsep{}\\[-2pt]
      &\exists k\in\NatSeg{\WordLength v}\,
        (q,b)=(\WordLength u+k,v(k))\vdash a=b
    \end{aligned}
  }{u,v\in\FiniteWords{A}\dsep(q,a)\in u\dsep
    \exists k\in\NatSeg{\WordLength v}\,
      (q,b)=(\WordLength u+k,v(k))\vdash a=b}%
    [WordConcatenationMixedBlocksImpossible]
    \proofstepwidestar[1]{u,v\in\FiniteWords{A}}{\rA}
    \proofstepwidestar[2]{(q,a)\in u}{\rA}
    \proofstepwidestar[3]{%
      \exists k\in\NatSeg{\WordLength v}\,
        (q,b)=(\WordLength u+k,v(k))}{\rA}
    \proofstepwidestar[1]{u\in\FiniteWords{A}}{\rAEa{1}}
    \proofstepwidestar[1]{u\colon\NatSeg{\WordLength u}\to A}{%
      \FormulaRefAuto{FiniteWordTypingAtLength}[thm]{4}}
    \proofstepwidestar[1,2]{q\in\NatSeg{\WordLength u}}{%
      \FormulaRefAuto{F\colon A\to B\dsep (x,y)\in F
        \vdash x\in A}[thm]{5,2}}
    \proofstepwidestar[7]{%
      k\in\NatSeg{\WordLength v}\land
        (q,b)=(\WordLength u+k,v(k))}{\rA}
    \proofstepwidestar[7]{k\in\NatSeg{\WordLength v}}{\rAEa{7}}
    \proofstepwidestar[7]{%
      (q,b)=(\WordLength u+k,v(k))}{\rAEb{7}}
    \proofstepwidestar[7]{q=\WordLength u+k\land b=v(k)}{%
      \FormulaRefAuto{(a,b)=(c,d)\eqvdash
        a=c\land b=d}[thm]{9}}
    \proofstepwidestar[7]{q=\WordLength u+k}{\rAEa{10}}
    \proofstepwidestar[1,7]{%
      \WordLength u+k\notin\NatSeg{\WordLength u}}{%
      \FormulaRefAuto{WordConcatenationIndexBlocksDisjoint}[thm]{1,8}}
    \proofstepwidestar[1,2,7]{%
      \WordLength u+k\in\NatSeg{\WordLength u}}{\rIE{11,6}}
    \proofstepwidestar[1,2,7]{\bot}{\rBI{12,13}}
    \proofstepwidestar[1,2,7]{a=b}{%
      \FormulaRefAuto{\bot\vdash P}[thm]{14}}
    \proofstepwidestar[1,2,3]{a=b}{\rEE{3,7,15}}
  \closeproofpartwideindR

  \proofpartwideindR[Beide Paare im verschobenen Block]{%
    \begin{aligned}[t]
      &u,v\in\FiniteWords{A}\dsep
        \exists j\in\NatSeg{\WordLength v}\,
          (q,a)=(\WordLength u+j,v(j))\dsep{}\\[-2pt]
      &\exists k\in\NatSeg{\WordLength v}\,
          (q,b)=(\WordLength u+k,v(k))\vdash a=b
    \end{aligned}
  }{u,v\in\FiniteWords{A}\dsep
    \exists j\in\NatSeg{\WordLength v}\,
      (q,a)=(\WordLength u+j,v(j))\dsep
    \exists k\in\NatSeg{\WordLength v}\,
      (q,b)=(\WordLength u+k,v(k))\vdash a=b}%
    [WordConcatenationShiftedBlockFunctional]
    \proofstepwidestar[1]{u,v\in\FiniteWords{A}}{\rA}
    \proofstepwidestar[2]{%
      \exists j\in\NatSeg{\WordLength v}\,
        (q,a)=(\WordLength u+j,v(j))}{\rA}
    \proofstepwidestar[3]{%
      \exists k\in\NatSeg{\WordLength v}\,
        (q,b)=(\WordLength u+k,v(k))}{\rA}
    \proofstepwidestar[1]{u\in\FiniteWords{A}}{\rAEa{1}}
    \proofstepwidestar[1]{v\in\FiniteWords{A}}{\rAEb{1}}
    \proofstepwidestar[1]{\WordLength u\in\mathbb N}{%
      \FormulaRefAuto{FiniteWordLengthNatural}[thm]{4}}
    \proofstepwidestar[1]{\WordLength v\in\mathbb N}{%
      \FormulaRefAuto{FiniteWordLengthNatural}[thm]{5}}
    \proofstepwidestar[1]{v\colon\NatSeg{\WordLength v}\to A}{%
      \FormulaRefAuto{FiniteWordTypingAtLength}[thm]{5}}
    \proofstepwidestar[2]{%
      (q,a)\in\{(\WordLength u+j,v(j))\mid
        j\in\NatSeg{\WordLength v}\}}{%
      \FormulaRefAuto{ShiftedWordGraphDef}[def]{2}}
    \proofstepwidestar[3]{%
      (q,b)\in\{(\WordLength u+k,v(k))\mid
        k\in\NatSeg{\WordLength v}\}}{%
      \FormulaRefAuto{ShiftedWordGraphDef}[def]{3}}
    \proofstepwidestar[1,2,3]{a=b}{%
      \FormulaRefAuto{ShiftedWordGraphFunctional}[thm]{6,7,8,9,10}}
  \closeproofpartwideindR

  \proofpartwide{Fallunterscheidung}
    \proofstepwidestar[1]{%
      u,v\in\FiniteWords{A}}{%
      \rA}
    \proofstepwidestar[2]{%
      q\in\NatSeg{\WordLength u+\WordLength v}}{%
      \rA}
    \proofstepwidestar[3]{%
      a\in A}{%
      \rA}
    \proofstepwidestar[4]{%
      b\in A}{%
      \rA}
    \proofstepwidestar[5]{%
      (q,a)\in\WordConcat{u}{v}}{%
      \rA}
    \proofstepwidestar[6]{%
      (q,b)\in\WordConcat{u}{v}}{%
      \rA}
    \proofstepwidestar[1,5]{%
      \begin{aligned}[t]&(q,a)\in u\\[-2pt]&\quad\lor\exists j\in\NatSeg{\WordLength v}\,((q,a)=(\WordLength u+j,v(j)))\end{aligned}}{%
      \FormulaRefAuto{WordConcatenationMembership}[thm]{1,5}}
    \proofstepwidestar[1,6]{%
      \begin{aligned}[t]&(q,b)\in u\\[-2pt]&\quad\lor\exists k\in\NatSeg{\WordLength v}\,((q,b)=(\WordLength u+k,v(k)))\end{aligned}}{%
      \FormulaRefAuto{WordConcatenationMembership}[thm]{1,6}}
    \proofstepwidestar[9]{%
      (q,a)\in u}{%
      \rA}
    \proofstepwidestar[10]{%
      (q,b)\in u}{%
      \rA}
    \proofstepwidestar[1,9,10]{%
      a=b}{%
      \FormulaRefAuto{WordConcatenationInitialBlockFunctional}[thm]{1,9,10}}
    \proofstepwidestar[12]{%
      \exists k\in\NatSeg{\WordLength v}\,((q,b)=(\WordLength u+k,v(k)))}{%
      \rA}
    \proofstepwidestar[1,9,12]{%
      a=b}{%
      \FormulaRefAuto{WordConcatenationMixedBlocksImpossible}[thm]{1,9,12}}
    \proofstepwidestar[1,6,9]{%
      a=b}{%
      \rOE{8,10,11,12,13}}
    \proofstepwidestar[15]{%
      \exists j\in\NatSeg{\WordLength v}\,((q,a)=(\WordLength u+j,v(j)))}{%
      \rA}
    \proofstepwidestar[16]{%
      (q,b)\in u}{%
      \rA}
    \proofstepwidestar[1,15,16]{%
      b=a}{%
      \FormulaRefAuto{WordConcatenationMixedBlocksImpossible}[thm]{1,16,15}}
    \proofstepwidestar[1,15,16]{%
      a=b}{%
      \FormulaRefAuto{a=b\vdash b=a}[thm]{17}}
    \proofstepwidestar[19]{%
      \exists k\in\NatSeg{\WordLength v}\,((q,b)=(\WordLength u+k,v(k)))}{%
      \rA}
    \proofstepwidestar[1,15,19]{%
      a=b}{%
      \FormulaRefAuto{WordConcatenationShiftedBlockFunctional}[thm]{1,15,19}}
    \proofstepwidestar[1,6,15]{%
      a=b}{%
      \rOE{8,16,18,19,20}}
    \proofstepwidestar[1,5,6]{%
      a=b}{%
      \rOE{7,9,14,15,21}}
  \closeproofpartwide
\end{tabproofsplitwide}

\FormulaThmDeltaK[Typisierung und Koordinaten der Konkatenation]{%
  \begin{aligned}[t]
    &\text{(i)}\;u,v\in\FiniteWords{A}\vdash
      \WordConcat{u}{v}\colon
      \NatSeg{\WordLength{u}+\WordLength{v}}\to A\\[-2pt]
    &\text{(ii)}\;u,v\in\FiniteWords{A}\dsep
      i\in\NatSeg{\WordLength{u}}\vdash
      \WordConcat{u}{v}(i)=u(i)\\[-2pt]
    &\text{(iii)}\;u,v\in\FiniteWords{A}\dsep
      j\in\NatSeg{\WordLength{v}}\vdash
      \WordConcat{u}{v}(\WordLength{u}+j)=v(j)
  \end{aligned}%
}{WordConcatenationTyping}{%
  \DeltaRow{Mengen}{A\dsep u\dsep v\dsep i\dsep j}
}
\begin{tabproofsplitwide}
  \proofpartwideindR[Typisierung]{%
    u,v\in\FiniteWords{A}\vdash
    \WordConcat{u}{v}\colon
      \NatSeg{\WordLength{u}+\WordLength{v}}\to A
  }{%
    u,v\in\FiniteWords{A}\vdash
    \WordConcat{u}{v}\colon
      \NatSeg{\WordLength{u}+\WordLength{v}}\to A
  }[WordConcatenationTypingPart]
    \proofstepwidestar[1]{%
      u\in\FiniteWords A}{%
      \rA}
    \proofstepwidestar[2]{%
      v\in\FiniteWords A}{%
      \rA}
    \proofstepwidestar[1,2]{%
      u\subseteq\NatSeg{\WordLength u+\WordLength v}\times A}{%
      \FormulaRefAuto{WordConcatenationInitialGraphSubset}[thm]{1,2}}
    \proofstepwidestar[1,2]{%
      \{(\WordLength u+j,v(j))\mid j\in\NatSeg{\WordLength v}\}\subseteq\NatSeg{\WordLength u+\WordLength v}\times A}{%
      \FormulaRefAuto{WordConcatenationShiftedGraphSubsetPart}[thm]{1,2}}
    \proofstepwidestar[1,2]{%
      u\cup\{(\WordLength u+j,v(j))\mid j\in\NatSeg{\WordLength v}\}\subseteq\NatSeg{\WordLength u+\WordLength v}\times A}{%
      \FormulaRefAuto{UnionSubsetCommonSuperset}[thm]{3,4}}
    \proofstepwidestar[1,2]{\WordConcat{u}{v}=u\cup\{(\WordLength{u}+j,v(j))\mid j\in\NatSeg{\WordLength{v}}\}}{%
      \FormulaRefAuto{WordConcatenationDef}[def]{1,2}}
    \proofstepwidestar[1,2]{%
      u\cup\{(\WordLength u+j,v(j))\mid j\in\NatSeg{\WordLength v}\}=\WordConcat uv}{%
      \FormulaRefAuto{a=b\vdash b=a}[thm]{6}}
    \proofstepwidestar[1,2]{%
      \WordConcat uv\subseteq\NatSeg{\WordLength u+\WordLength v}\times A}{%
      \rIE{7,5}}
    \proofstepwidestar[9]{%
      q\in\NatSeg{\WordLength u+\WordLength v}}{%
      \rA}
    \proofstepwidestar[1,2,9]{%
      \exists a\in A\,((q,a)\in\WordConcat uv)}{%
      \FormulaRefAuto{WordConcatenationGraphTotal}[thm]{1,2,9}}
    \proofstepwidestar[11]{%
      a\in A\land(q,a)\in\WordConcat uv}{%
      \rA}
    \proofstepwidestar[12]{%
      b\in A\land(q,b)\in\WordConcat uv}{%
      \rA}
    \proofstepwidestar[11]{%
      a\in A}{%
      \rAEa{11}}
    \proofstepwidestar[12]{%
      b\in A}{%
      \rAEa{12}}
    \proofstepwidestar[11]{%
      (q,a)\in\WordConcat uv}{%
      \rAEb{11}}
    \proofstepwidestar[12]{%
      (q,b)\in\WordConcat uv}{%
      \rAEb{12}}
    \proofstepwidestar[1,2,9,11,12]{%
      a=b}{%
      \FormulaRefAuto{WordConcatenationGraphFunctional}[thm]{\rAI{1,2},9,13,14,15,16}}
    \proofstepwidestar[1,2,9]{%
      \exists!a\in A\,((q,a)\in\WordConcat uv)}{%
      \UEI{10,11,12,17}}
    \proofstepwidestar[1,2]{%
      \forall q\in\NatSeg{\WordLength u+\WordLength v}\,\exists!a\in A\,((q,a)\in\WordConcat uv)}{%
      \rUI{\rRI{9,18}}}
    \proofstepwidestar[1,2]{%
      \WordConcat uv\colon\NatSeg{\WordLength u+\WordLength v}\to A}{%
      \FormulaRefAuto{UniqueValuedGraphFunction}[thm]{8,19}}
  \closeproofpartwideindR

  \proofpartwideindR[Anfangsblock]{%
    u,v\in\FiniteWords{A}\dsep
      i\in\NatSeg{\WordLength{u}}\vdash
      \WordConcat{u}{v}(i)=u(i)
  }{u,v\in\FiniteWords{A}\dsep
      i\in\NatSeg{\WordLength{u}}\vdash
      \WordConcat{u}{v}(i)=u(i)}[WordConcatenationInitialCoordinate]
    \proofstepwidestar[1]{u,v\in\FiniteWords{A}}{\rA}
    \proofstepwidestar[2]{i\in\NatSeg{\WordLength{u}}}{\rA}
    \proofstepwidestar[1]{%
      \WordConcat{u}{v}\colon
        \NatSeg{\WordLength{u}+\WordLength{v}}\to A}{%
      \FormulaRefAuto{WordConcatenationTypingPart}[thm]{1}}
    \proofstepwidestar[1]{%
      u\colon\NatSeg{\WordLength{u}}\to A}{%
      \FormulaRefAuto{FiniteWordTypingAtLength}[thm]{1}}
    \proofstepwidestar[1,2]{(i,u(i))\in u}{%
      \FormulaRefAuto{F\colon A\to B\dsep x\in A\vdash
        (x,F(x))\in F}[thm]{4,2}}
    \proofstepwidestar[1]{%
      \WordConcat{u}{v}=u\cup
      \{(\WordLength{u}+j,v(j))\mid
        j\in\NatSeg{\WordLength{v}}\}}{%
      \FormulaRefAuto{WordConcatenationDef}[def]{1}}
    \proofstepwidestar[1,2]{(i,u(i))\in\WordConcat{u}{v}}{%
      \FormulaRefAuto{z\in A\vdash z\in A\cup B}[thm]{5,6}}
    \proofstepwidestar[1,2]{\WordConcat{u}{v}(i)=u(i)}{%
      \FormulaRefAuto{F\colon A\to B\dsep (x,y)\in F
        \vdash F(x)=y}[thm]{3,7}}
  \closeproofpartwideindR

  \proofpartwideindR[Verschobener Block]{%
    u,v\in\FiniteWords{A}\dsep
      j\in\NatSeg{\WordLength{v}}\vdash
      \WordConcat{u}{v}(\WordLength{u}+j)=v(j)
  }{u,v\in\FiniteWords{A}\dsep
      j\in\NatSeg{\WordLength{v}}\vdash
      \WordConcat{u}{v}(\WordLength{u}+j)=v(j)}%
    [WordConcatenationShiftedCoordinate]
    \proofstepwidestar[1]{u,v\in\FiniteWords{A}}{\rA}
    \proofstepwidestar[2]{j\in\NatSeg{\WordLength{v}}}{\rA}
    \proofstepwidestar[1]{%
      \WordConcat{u}{v}\colon
        \NatSeg{\WordLength{u}+\WordLength{v}}\to A}{%
      \FormulaRefAuto{WordConcatenationTypingPart}[thm]{1}}
    \proofstepwidestar[1,2]{%
      (\WordLength{u}+j,v(j))
      \in\{(\WordLength{u}+k,v(k))\mid
        k\in\NatSeg{\WordLength{v}}\}}{%
      \FormulaRefAuto{ShiftedWordGraphDef}[def]{%
        \rEI{\rAI{2,\rII}}}}
    \proofstepwidestar[1]{%
      \WordConcat{u}{v}=u\cup
      \{(\WordLength{u}+k,v(k))\mid
        k\in\NatSeg{\WordLength{v}}\}}{%
      \FormulaRefAuto{WordConcatenationDef}[def]{1}}
    \proofstepwidestar[1,2]{%
      (\WordLength{u}+j,v(j))\in\WordConcat{u}{v}}{%
      \FormulaRefAuto{z\in B\vdash z\in A\cup B}[thm]{4,5}}
    \proofstepwidestar[1,2]{%
      \WordConcat{u}{v}(\WordLength{u}+j)=v(j)}{%
      \FormulaRefAuto{F\colon A\to B\dsep (x,y)\in F
        \vdash F(x)=y}[thm]{3,6}}
  \closeproofpartwideindR
\end{tabproofsplitwide}

\FormulaThmDeltaK[Abgeschlossenheit der endlichen Wortmenge unter Konkatenation]{%
  u,v\in\FiniteWords A\vdash\WordConcat uv\in\FiniteWords A%
}{WordConcatenationClosure}{%
  \DeltaRow{Mengen}{A\dsep u\dsep v}
}
\begin{tabproofwide}
  \proofstepwidestar[1]{u\in\FiniteWords A}{\rA}
  \proofstepwidestar[2]{v\in\FiniteWords A}{\rA}
  \proofstepwidestar[1]{\WordLength u\in\mathbb N}{%
    \FormulaRefAuto{FiniteWordLengthNatural}[thm]{1}}
  \proofstepwidestar[2]{\WordLength v\in\mathbb N}{%
    \FormulaRefAuto{FiniteWordLengthNatural}[thm]{2}}
  \proofstepwidestar[1,2]{\WordLength u+\WordLength v\in\mathbb N}{%
    \FormulaRefAuto{PeanoAddClosure}[thm]{3,4}}
  \proofstepwidestar[1,2]{\WordConcat uv\colon\NatSeg{\WordLength u+\WordLength v}\to A}{%
    \FormulaRefAuto{WordConcatenationTypingPart}[thm]{1,2}}
  \proofstepwidestar[1,2]{\WordConcat uv\in\FiniteWords A}{%
    \FormulaRefAuto{FiniteWordFromTyping}[thm]{5,6}}
\end{tabproofwide}

\FormulaThmDeltaK[Länge einer Konkatenation]{%
  u,v\in\FiniteWords{A}\vdash
  \WordLength{\WordConcat{u}{v}}
    =\WordLength{u}+\WordLength{v}%
}{WordConcatenationLength}{%
  \DeltaRow{Mengen}{A\dsep u\dsep v}
}
\begin{tabproofwide}
  \proofstepwidestar[1]{u\in\FiniteWords{A}}{\rA}
  \proofstepwidestar[2]{v\in\FiniteWords{A}}{\rA}
  \proofstepwidestar[1,2]{%
    \WordConcat{u}{v}\colon
    \NatSeg{\WordLength{u}+\WordLength{v}}\to A}{%
    \FormulaRefAuto{WordConcatenationTypingPart}[thm]{1,2}}
  \proofstepwidestar[1]{\WordLength{u}\in\mathbb N}{%
    \FormulaRefAuto{FiniteWordLengthNatural}[thm]{1}}
  \proofstepwidestar[2]{\WordLength{v}\in\mathbb N}{%
    \FormulaRefAuto{FiniteWordLengthNatural}[thm]{2}}
  \proofstepwidestar[1,2]{%
    \WordLength{u}+\WordLength{v}\in\mathbb N}{%
    \FormulaRefAuto{PeanoAddClosure}[thm]{4,5}}
  \proofstepwidestar[1,2]{\WordConcat{u}{v}\in\FiniteWords{A}}{%
    \FormulaRefAuto{FiniteWordFromTyping}[thm]{6,3}}
  \proofstepwidestar[1,2]{%
    \WordLength{\WordConcat{u}{v}}
      =\WordLength{u}+\WordLength{v}}{%
    \FormulaRefAuto{FiniteWordLengthFromTyping}[thm]{7,6,3}}
\end{tabproofwide}

\noindent\begin{minipage}{\linewidth}
\FormulaThmDeltaK[Kürzung bei festem Konkatenationsschnitt]{%
  \begin{aligned}[t]
    &u,v,u',v'\in\FiniteWords{A}\dsep
      \WordLength{u}=\WordLength{u'}\dsep
      \WordConcat{u}{v}=\WordConcat{u'}{v'}\\[-2pt]
    &\vdash
      \begin{aligned}[t]
        &\text{(i)}\;u=u'\\[-2pt]
        &\text{(ii)}\;v=v'
      \end{aligned}
  \end{aligned}%
}{WordConcatenationFixedCutCancellation}{%
  \DeltaRow{Mengen}{A\dsep u\dsep v\dsep u'\dsep v'\dsep i\dsep j}
}
\end{minipage}\par
\begin{tabproofsplitwide}
  \proofpartwideindR[Linker Faktor]{%
    \begin{aligned}[t]
      &u,v,u',v'\in\FiniteWords{A}\dsep
        \WordLength{u}=\WordLength{u'}\dsep{}\\[-2pt]
      &\WordConcat{u}{v}=\WordConcat{u'}{v'}\vdash u=u'
    \end{aligned}
  }{u,v,u',v'\in\FiniteWords{A}\dsep
    \WordLength{u}=\WordLength{u'}\dsep
    \WordConcat{u}{v}=\WordConcat{u'}{v'}\vdash u=u'}%
    [WordConcatenationFixedCutLeft]
    \proofstepwidestar[1]{u,v,u',v'\in\FiniteWords{A}}{\rA}
    \proofstepwidestar[2]{\WordLength{u}=\WordLength{u'}}{\rA}
    \proofstepwidestar[3]{%
      \WordConcat{u}{v}=\WordConcat{u'}{v'}}{\rA}
    \proofstepwidestar[1]{%
      u\colon\NatSeg{\WordLength u}\to A}{%
      \FormulaRefAuto{FiniteWordTypingAtLength}[thm]{1}}
    \proofstepwidestar[1]{%
      u'\colon\NatSeg{\WordLength{u'}}\to A}{%
      \FormulaRefAuto{FiniteWordTypingAtLength}[thm]{1}}
    \proofstepwidestar[1,2]{%
      u'\colon\NatSeg{\WordLength u}\to A}{\rIE{2,5}}
    \proofstepwidestar[7]{i\in\NatSeg{\WordLength u}}{\rA}
    \proofstepwidestar[1,7]{\WordConcat{u}{v}(i)=u(i)}{%
      \FormulaRefAuto{WordConcatenationInitialCoordinate}[thm]{1,7}}
    \proofstepwidestar[1,2,7]{\WordConcat{u'}{v'}(i)=u'(i)}{%
      \FormulaRefAuto{WordConcatenationInitialCoordinate}[thm]{1,2,7}}
    \proofstepwidestar[1,2,3,7]{u(i)=u'(i)}{\rIE{3,8,9}}
    \proofstepwidestar[1,2,3]{%
      \forall i\in\NatSeg{\WordLength u}\,u(i)=u'(i)}{%
      \rUI{\rRI{7,10}}}
    \proofstepwidestar[1,2,3]{u=u'}{%
      \FormulaRefAuto{FunctionExtensionality}[thm]{4,6,11}}
  \closeproofpartwideindR

  \proofpartwideindR[Rechter Faktor]{%
    \begin{aligned}[t]
      &u,v,u',v'\in\FiniteWords{A}\dsep
        \WordLength{u}=\WordLength{u'}\dsep{}\\[-2pt]
      &\WordConcat{u}{v}=\WordConcat{u'}{v'}\vdash v=v'
    \end{aligned}
  }{u,v,u',v'\in\FiniteWords{A}\dsep
    \WordLength{u}=\WordLength{u'}\dsep
    \WordConcat{u}{v}=\WordConcat{u'}{v'}\vdash v=v'}%
    [WordConcatenationFixedCutRight]
    \proofstepwidestar[1]{u,v,u',v'\in\FiniteWords{A}}{\rA}
    \proofstepwidestar[2]{\WordLength{u}=\WordLength{u'}}{\rA}
    \proofstepwidestar[3]{%
      \WordConcat{u}{v}=\WordConcat{u'}{v'}}{\rA}
    \proofstepwidestar[1]{%
      \WordLength{\WordConcat{u}{v}}
        =\WordLength u+\WordLength v}{%
      \FormulaRefAuto{WordConcatenationLength}[thm]{1}}
    \proofstepwidestar[1]{%
      \WordLength{\WordConcat{u'}{v'}}
        =\WordLength{u'}+\WordLength{v'}}{%
      \FormulaRefAuto{WordConcatenationLength}[thm]{1}}
    \proofstepwidestar[1,3]{%
      \WordLength u+\WordLength v
        =\WordLength{u'}+\WordLength{v'}}{%
      \rIE{3,4,5}}
    \proofstepwidestar[1,2,3]{%
      \WordLength u+\WordLength v
        =\WordLength u+\WordLength{v'}}{\rIE{2,6}}
    \proofstepwidestar[1]{%
      \WordLength u,\WordLength v,\WordLength{v'}\in\mathbb N}{%
      \FormulaRefAuto{FiniteWordLengthNatural}[thm]{1}}
    \proofstepwidestar[1,2,3]{%
      \WordLength v=\WordLength{v'}}{%
      \FormulaRefAuto{PeanoAddLeftCancellation}[thm]{8,7}}
    \proofstepwidestar[1]{%
      v\colon\NatSeg{\WordLength v}\to A}{%
      \FormulaRefAuto{FiniteWordTypingAtLength}[thm]{1}}
    \proofstepwidestar[1]{%
      v'\colon\NatSeg{\WordLength{v'}}\to A}{%
      \FormulaRefAuto{FiniteWordTypingAtLength}[thm]{1}}
    \proofstepwidestar[1,2,3]{%
      v'\colon\NatSeg{\WordLength v}\to A}{\rIE{9,11}}
    \proofstepwidestar[13]{j\in\NatSeg{\WordLength v}}{\rA}
    \proofstepwidestar[1,13]{%
      \WordConcat{u}{v}(\WordLength u+j)=v(j)}{%
      \FormulaRefAuto{WordConcatenationShiftedCoordinate}[thm]{1,13}}
    \proofstepwidestar[1,2,3,13]{%
      j\in\NatSeg{\WordLength{v'}}}{\rIE{9,13}}
    \proofstepwidestar[1,2,3,13]{%
      \WordConcat{u'}{v'}(\WordLength{u'}+j)=v'(j)}{%
      \FormulaRefAuto{WordConcatenationShiftedCoordinate}[thm]{1,15}}
    \proofstepwidestar[1,2,3,13]{%
      \WordConcat{u'}{v'}(\WordLength u+j)=v'(j)}{\rIE{2,16}}
    \proofstepwidestar[1,2,3,13]{v(j)=v'(j)}{\rIE{3,14,17}}
    \proofstepwidestar[1,2,3]{%
      \forall j\in\NatSeg{\WordLength v}\,v(j)=v'(j)}{%
      \rUI{\rRI{13,18}}}
    \proofstepwidestar[1,2,3]{v=v'}{%
      \FormulaRefAuto{FunctionExtensionality}[thm]{10,12,19}}
  \closeproofpartwideindR
\end{tabproofsplitwide}

\FormulaThmDeltaK[Leerwortgesetze der Konkatenation]{%
  w\in\FiniteWords{A}\vdash
  \begin{aligned}[t]
    &\text{(i)}\;\WordConcat{\EmptyWord{A}}{w}=w\\[-2pt]
    &\text{(ii)}\;\WordConcat{w}{\EmptyWord{A}}=w
  \end{aligned}%
}{WordConcatenationIdentity}{%
  \DeltaRow{Mengen}{A\dsep w}
}
\begin{tabproofsplitwide}
  \proofpartwideindR[Linkes Leerwortgesetz]{%
    w\in\FiniteWords{A}\vdash
      \WordConcat{\EmptyWord{A}}{w}=w
  }{w\in\FiniteWords{A}\vdash
      \WordConcat{\EmptyWord{A}}{w}=w}%
    [WordConcatenationLeftIdentity]
    \proofstepwidestar[1]{w\in\FiniteWords{A}}{\rA}
    \proofstepwidestar{\EmptyWord{A}\in\FiniteWords{A}}{%
      \FormulaRefAuto{EmptyWordFinite}[thm]}
    \proofstepwidestar{\WordLength{\EmptyWord{A}}=0}{%
      \FormulaRefAuto{EmptyWordLengthZero}[thm]}
    \proofstepwidestar[1]{\WordLength w\in\mathbb N}{%
      \FormulaRefAuto{FiniteWordLengthNatural}[thm]{1}}
    \proofstepwidestar[1]{%
      \WordConcat{\EmptyWord{A}}{w}\colon
      \NatSeg{\WordLength{\EmptyWord{A}}+\WordLength{w}}\to A}{%
      \FormulaRefAuto{WordConcatenationTypingPart}[thm]{2,1}}
    \proofstepwidestar[1]{0+\WordLength w=\WordLength w}{%
      \FormulaRefAuto{PeanoAddLeftZero}[thm]{4}}
    \proofstepwidestar[1]{%
      \WordConcat{\EmptyWord{A}}{w}\colon
      \NatSeg{\WordLength{w}}\to A}{%
      \rIE{3,6,5}}
    \proofstepwidestar[1]{%
      w\colon\NatSeg{\WordLength{w}}\to A}{%
      \FormulaRefAuto{FiniteWordTypingAtLength}[thm]{1}}
    \proofstepwidestar[9]{i\in\NatSeg{\WordLength{w}}}{\rA}
    \proofstepwidestar[1,9]{i\in\mathbb N}{%
      \FormulaRefAuto{PeanoNatSegElementNatural}[thm]{4,9}}
    \proofstepwidestar[1,9]{%
      \WordConcat{\EmptyWord{A}}{w}
        (\WordLength{\EmptyWord{A}}+i)=w(i)}{%
      \FormulaRefAuto{WordConcatenationShiftedCoordinate}[thm]{2,1,9}}
    \proofstepwidestar[1,9]{0+i=i}{%
      \FormulaRefAuto{PeanoAddLeftZero}[thm]{10}}
    \proofstepwidestar[1,9]{%
      \WordConcat{\EmptyWord{A}}{w}(i)=w(i)}{%
      \rIE{3,12,11}}
    \proofstepwidestar[1]{%
      \forall i\in\NatSeg{\WordLength{w}}\,
      \WordConcat{\EmptyWord{A}}{w}(i)=w(i)}{\rUI{\rRI{9,13}}}
    \proofstepwidestar[1]{\WordConcat{\EmptyWord{A}}{w}=w}{%
      \FormulaRefAuto{FunctionExtensionality}[thm]{7,8,14}}
  \closeproofpartwideindR

  \proofpartwideindR[Rechtes Leerwortgesetz]{%
    w\in\FiniteWords{A}\vdash
      \WordConcat{w}{\EmptyWord{A}}=w
  }{w\in\FiniteWords{A}\vdash
      \WordConcat{w}{\EmptyWord{A}}=w}%
    [WordConcatenationRightIdentity]
    \proofstepwidestar[1]{w\in\FiniteWords{A}}{\rA}
    \proofstepwidestar{\EmptyWord{A}\in\FiniteWords{A}}{%
      \FormulaRefAuto{EmptyWordFinite}[thm]}
    \proofstepwidestar{\WordLength{\EmptyWord{A}}=0}{%
      \FormulaRefAuto{EmptyWordLengthZero}[thm]}
    \proofstepwidestar[1]{\WordLength w\in\mathbb N}{%
      \FormulaRefAuto{FiniteWordLengthNatural}[thm]{1}}
    \proofstepwidestar[1]{%
      \WordConcat{w}{\EmptyWord{A}}\colon
      \NatSeg{\WordLength{w}+\WordLength{\EmptyWord{A}}}\to A}{%
      \FormulaRefAuto{WordConcatenationTypingPart}[thm]{1,2}}
    \proofstepwidestar[1]{\WordLength w+0=\WordLength w}{%
      \FormulaRefAuto{PeanoAddRightZero}[thm]{4}}
    \proofstepwidestar[1]{%
      \WordConcat{w}{\EmptyWord{A}}\colon
      \NatSeg{\WordLength{w}}\to A}{%
      \rIE{3,6,5}}
    \proofstepwidestar[1]{%
      w\colon\NatSeg{\WordLength{w}}\to A}{%
      \FormulaRefAuto{FiniteWordTypingAtLength}[thm]{1}}
    \proofstepwidestar[9]{i\in\NatSeg{\WordLength{w}}}{\rA}
    \proofstepwidestar[1,9]{%
      \WordConcat{w}{\EmptyWord{A}}(i)=w(i)}{%
      \FormulaRefAuto{WordConcatenationInitialCoordinate}[thm]{1,2,9}}
    \proofstepwidestar[1]{%
      \forall i\in\NatSeg{\WordLength{w}}\,
      \WordConcat{w}{\EmptyWord{A}}(i)=w(i)}{\rUI{\rRI{9,10}}}
    \proofstepwidestar[1]{\WordConcat{w}{\EmptyWord{A}}=w}{%
      \FormulaRefAuto{FunctionExtensionality}[thm]{7,8,11}}
  \closeproofpartwideindR
\end{tabproofsplitwide}

\subsection{Konkatenationsgesetze aus der Wortstruktur}

Die konkrete Formel für die Konkatenation wird nun mit der primitiven
Anfügung verbunden. Danach genügt für das Assoziativitätsgesetz die
axiomatische Wortinduktion W3.

\FormulaThmDeltaK[Konkatenation verträgt die Anfügung rechts]{%
  u,v\in\FiniteWords A\dsep a\in A\vdash
    \WordConcat u{\sigma_A(v,a)}=\sigma_A(\WordConcat u v,a)
}{WordConcatenationSuccessor}{\DeltaRow{Mengen}{A\dsep u\dsep v\dsep a\dsep j}}
Für die Tabelle seien \(m=\WordLength u\), \(n=\WordLength v\) und
\(q=\sigma_A(v,a)\). Die beiden verschobenen Mengen werden durch
\(D=\{(m+j,q(j))\mid j\in\NatSeg n\}\) und
\(E=\{(m+j,v(j))\mid j\in\NatSeg n\}\) abgekürzt.
\begin{tabproofwide}
  \proofstepwidestar[1]{u,v\in\FiniteWords A}{\rA}
  \proofstepwidestar[2]{a\in A}{\rA}
  \proofstepwidestar[1,2]{q\in\FiniteWords A\land\WordLength q=n+1}{\FormulaRefAuto{EarlyWordSuccessorLength}[thm]{\rAEb{1},2}}
  \proofstepwidestar[1]{m,n\in\mathbb N}{\rAI{\FormulaRefAuto{FiniteWordLengthNatural}[thm]{\rAEa{1}},\FormulaRefAuto{FiniteWordLengthNatural}[thm]{\rAEb{1}}}}
  \proofstepwidestar[1]{\NatSeg{n+1}=\NatSeg n\cup\{n\}}{\FormulaRefAuto{PeanoNatSegAddOne}[thm]{\rAEb{4}}}
  \proofstepwidestar[1,2]{\begin{aligned}[t]
    &\{(m+j,q(j))\mid j\in\NatSeg{\WordLength q}\}\\[-2pt]
    &\quad=D\cup\{(m+j,q(j))\mid j\in\{n\}\}
  \end{aligned}}{\rIE{5,\rAEb{3},\FormulaRefAuto{B27TermImageUnion}[thm]}}
  \proofstepwidestar[1,2]{\forall j\in\NatSeg n\;q(j)=v(j)}{\rAEa{\FormulaRefAuto{EarlyWordSuccessorValues}[thm]{\rAEb{1},2}}}
  \proofstepwidestar[8]{j\in\NatSeg n}{\rA}
  \proofstepwidestar[1,2,8]{(m+j,q(j))=(m+j,v(j))}{\rIE{\rRE{8,\rUE{7}},\rII}}
  \proofstepwidestar[1,2]{D=E}{\FormulaRefAuto{B27TermImagePointwiseEquality}[thm]{\rUI{\rRI{8,9}}}}
  \proofstepwidestar[1,2]{q(n)=a}{\rAEb{\FormulaRefAuto{EarlyWordSuccessorValues}[thm]{\rAEb{1},2}}}
  \proofstepwidestar[1,2]{\{(m+j,q(j))\mid j\in\{n\}\}=\{(m+n,a)\}}{\rIE{11,\FormulaRefAuto{B27TermImageSingleton}[thm]}}
  \proofstepwidestar[1,2]{\WordConcat u q=u\cup(E\cup\{(m+n,a)\})}{\rIE{6,10,12,\FormulaRefAuto{WordConcatenationDef}[def]{\rAEa{1},\rAEa{3}}}}
  \proofstepwidestar[1]{\WordConcat u v=u\cup E}{\FormulaRefAuto{WordConcatenationDef}[def]{1}}
  \proofstepwidestar[1]{\WordConcat u v\in\FiniteWords A}{\FormulaRefAuto{WordConcatenationClosure}[thm]{1}}
  \proofstepwidestar[1]{\WordLength{\WordConcat u v}=m+n}{\FormulaRefAuto{WordConcatenationLength}[thm]{1}}
  \proofstepwidestar[1,2]{\WordConcat u q=\WordConcat u v\cup\{(\WordLength{\WordConcat u v},a)\}}{\rIE{\FormulaRefAuto{(A\cup B)\cup C=A\cup(B\cup C)}[thm],14,16,13}}
  \proofstepwidestar[1,2]{\sigma_A(\WordConcat u v,a)=\WordConcat u v\cup\{(\WordLength{\WordConcat u v},a)\}}{\FormulaRefAuto{EarlyWordSuccessorEquation}[thm]{15,2}}
  \proofstepwidestar[1,2]{\WordConcat u{\sigma_A(v,a)}=\sigma_A(\WordConcat u v,a)}{\FormulaRefAuto{a=b\dsep c=b\vdash a=c}[thm]{17,18}}
\end{tabproofwide}


\Needspace{20\baselineskip}
\noindent\begin{minipage}{\linewidth}
\FormulaThmDeltaK[Assoziativität der Wortkonkatenation]{%
  u,v,w\in\FiniteWords{A}\vdash
  \WordConcat{\WordConcat{u}{v}}{w}
  =\WordConcat{u}{\WordConcat{v}{w}}%
}{WordConcatenationAssociative}{%
  \DeltaRow{Mengen}{A\dsep u\dsep v\dsep w}
}
\end{minipage}\par

Für feste \(u,v\in\FiniteWords A\) setzen wir
\[
 P(t)\coloneqq
   \WordConcat{\WordConcat u v}t=\WordConcat u{\WordConcat v t}.
\]
\begin{tabproofwide}
  \proofstepwidestar[1]{u,v,w\in\FiniteWords A}{\rA}
  \proofstepwidestar[1]{\WordConcat u v\in\FiniteWords A}{\FormulaRefAuto{WordConcatenationClosure}[thm]{\rAEa{1},\rAEa{\rAEb{1}}}}
  \proofstepwidestar[1]{\WordConcat{\WordConcat u v}{\EmptyWord A}=\WordConcat u v}{\FormulaRefAuto{WordConcatenationRightIdentity}[thm]{2}}
  \proofstepwidestar[1]{\WordConcat v{\EmptyWord A}=v}{\FormulaRefAuto{WordConcatenationRightIdentity}[thm]{\rAEa{\rAEb{1}}}}
  \proofstepwidestar[1]{P(\EmptyWord A)}{\rIE{\FormulaRefAuto{a=b\vdash b=a}[thm]{4},3}}
  \proofstepwidestar[6]{t\in\FiniteWords A}{\rA}
  \proofstepwidestar[7]{a\in A}{\rA}
  \proofstepwidestar[8]{P(t)}{\rA}
  \proofstepwidestar[1,6]{\WordConcat v t\in\FiniteWords A}{\FormulaRefAuto{WordConcatenationClosure}[thm]{\rAEa{\rAEb{1}},6}}
  \proofstepwidestar[1,6,7]{\WordConcat{\WordConcat u v}{\sigma_A(t,a)}=\sigma_A(\WordConcat{\WordConcat u v}t,a)}{\FormulaRefAuto{WordConcatenationSuccessor}[thm]{2,6,7}}
  \proofstepwidestar[1,6,7,8]{\WordConcat{\WordConcat u v}{\sigma_A(t,a)}=\sigma_A(\WordConcat u{\WordConcat v t},a)}{\rIE{8,10}}
  \proofstepwidestar[1,6,7]{\WordConcat u{\sigma_A(\WordConcat v t,a)}=\sigma_A(\WordConcat u{\WordConcat v t},a)}{\FormulaRefAuto{WordConcatenationSuccessor}[thm]{\rAEa{1},9,7}}
  \proofstepwidestar[1,6,7]{\WordConcat v{\sigma_A(t,a)}=\sigma_A(\WordConcat v t,a)}{\FormulaRefAuto{WordConcatenationSuccessor}[thm]{\rAEa{\rAEb{1}},6,7}}
  \proofstepwidestar[1,6,7,8]{P(\sigma_A(t,a))}{\rIE{\FormulaRefAuto{a=b\vdash b=a}[thm]{13},\FormulaRefAuto{a=b\dsep c=b\vdash a=c}[thm]{11,12}}}
  \proofstepwidestar[]{\mathsf{WortStr}_A(\FiniteWords A,\EmptyWord A,\sigma_A)}{\FormulaRefAuto{WordStructureConcreteModel}[thm]}
  \proofstepwidestar[1]{\WordConcat{\WordConcat u v}w=\WordConcat u{\WordConcat v w}}{\FormulaRefAuto{WordStructureInductionRule}[ax]{15,5,[6,7,8]\vdots14,\rAEb{\rAEb{1}}}}
\end{tabproofwide}


\Needspace{16\baselineskip}
\FormulaThmDeltaK[Abgeschlossenheit nichtleerer Konkatenationen]{%
  u,v\in\NonemptyWords{A}\vdash
  \WordConcat{u}{v}\in\NonemptyWords{A}%
}{NonemptyWordConcatenationClosure}{%
  \DeltaRow{Mengen}{A\dsep u\dsep v}
}
\begin{tabproofwide}
  \proofstepwidestar[1]{u,v\in\NonemptyWords A}{\rA}
  \proofstepwidestar[1]{u,v\in\FiniteWords A}{%
    \FormulaRefAuto{NonemptyWordFinite}[thm]{1}}
  \proofstepwidestar[1]{\WordLength u\in\mathbb N}{%
    \FormulaRefAuto{NonemptyWordLengthNatural}[thm]{1}}
  \proofstepwidestar[1]{\WordLength v\in\mathbb N}{%
    \FormulaRefAuto{NonemptyWordLengthNatural}[thm]{1}}
  \proofstepwidestar[1]{u\in\FiniteWords{A}\land\WordLength{u}\neq0}{%
      \FormulaRefAuto{NonemptyWordMembershipElimination}[thm]{1}}
  \proofstepwidestar[1]{\WordLength u\neq0}{%
    \rAEb{5}}
  \proofstepwidestar[1]{\WordLength u+\WordLength v\neq0}{%
    \FormulaRefAuto{PeanoAddNonzeroLeft}[thm]{3,4,6}}
  \proofstepwidestar[1]{\WordConcat uv\in\FiniteWords A}{%
    \FormulaRefAuto{WordConcatenationClosure}[thm]{2}}
  \proofstepwidestar[1]{\WordLength{\WordConcat uv}=\WordLength u+\WordLength v}{%
    \FormulaRefAuto{WordConcatenationLength}[thm]{2}}
  \proofstepwidestar[1]{\WordLength{\WordConcat uv}\neq0}{\rIE{9,7}}
  \proofstepwidestar[1]{\WordConcat uv\in\NonemptyWords A}{%
    \FormulaRefAuto{NonemptyWordMembershipIntroduction}[thm]{\rAI{8,10}}}
\end{tabproofwide}

\section{Endzerlegung, Induktion und Rekursion}

\FormulaThmDeltaK[Injektivität der Buchstabenwörter]{%
  a,b\in A\vdash
    \bigl(\LetterWord a=\LetterWord b\leftrightarrow a=b\bigr)%
}{LetterWordInjectivity}{%
  \DeltaRow{Mengen}{A\dsep a\dsep b}
}
\begin{tabproofwide}
  \proofstepwidestar[1]{a,b\in A}{\rA}
  \proofstepwidestar[2]{\LetterWord a=\LetterWord b}{\rA}
  \proofstepwidestar[1]{\LetterWord a(0)=a}{%
    \FormulaRefAuto{LetterWordValueZero}[thm]{1}}
  \proofstepwidestar[1]{\LetterWord b(0)=b}{%
    \FormulaRefAuto{LetterWordValueZero}[thm]{1}}
  \proofstepwidestar[1,2]{a=b}{\rIE{2,3,4}}
  \proofstepwidestar[1]{\LetterWord a=\LetterWord b\rightarrow a=b}{%
    \rRI{2,5}}
  \proofstepwidestar[7]{a=b}{\rA}
  \proofstepwidestar[1,7]{\LetterWord a=\LetterWord b}{\rIE{7}}
  \proofstepwidestar[1]{a=b\rightarrow\LetterWord a=\LetterWord b}{%
    \rRI{7,8}}
  \proofstepwidestar[1]{\LetterWord a=\LetterWord b\leftrightarrow a=b}{%
    \rLRI{6,9}}
\end{tabproofwide}

\Needspace{22\baselineskip}
\FormulaThmDeltaK[Wörter der Länge eins sind Buchstabenwörter]{%
  w\in\FiniteWords A\dsep\WordLength w=1\vdash
  \begin{aligned}[t]
    &\text{(i)}\;w(0)\in A\\[-2pt]
    &\text{(ii)}\;w=\LetterWord{w(0)}
  \end{aligned}%
}{FiniteWordLengthOneCharacterization}{%
  \DeltaRow{Mengen}{A\dsep w\dsep i}
}
\begin{tabproofsplitwide}
  \proofpartwideindR[Typisierung des einzigen Buchstabens]{%
    w\in\FiniteWords A\dsep\WordLength w=1\vdash w(0)\in A
  }{w\in\FiniteWords A\dsep\WordLength w=1\vdash w(0)\in A}[FiniteWordLengthOneLetter]
    \proofstepwidestar[1]{%
      w\in\FiniteWords A}{%
      \rA}
    \proofstepwidestar[2]{%
      \WordLength w=1}{%
      \rA}
    \proofstepwidestar[1]{%
      w\colon\NatSeg{\WordLength w}\to A}{%
      \FormulaRefAuto{FiniteWordTypingAtLength}[thm]{1}}
    \proofstepwidestar[1,2]{%
      w\colon\NatSeg 1\to A}{%
      \rIE{2,3}}
    \proofstepwidestar[]{%
      0\in\NatSeg 1}{%
      \FormulaRefAuto{PeanoNatSegOneMembership}[thm]}
    \proofstepwidestar[1,2]{%
      w(0)\in A}{%
      \FormulaRefAuto{F\colon A\to B\dsep x\in A\vdash F(x)\in B}[thm]{4,5}}
  \closeproofpartwideindR

  \proofpartwideindR[Darstellung als Buchstabenwort]{%
    w\in\FiniteWords A\dsep\WordLength w=1\vdash w=\LetterWord{w(0)}
  }{w\in\FiniteWords A\dsep\WordLength w=1\vdash w=\LetterWord{w(0)}}[FiniteWordLengthOneWord]
    \proofstepwidestar[1]{%
      w\in\FiniteWords A}{%
      \rA}
    \proofstepwidestar[2]{%
      \WordLength w=1}{%
      \rA}
    \proofstepwidestar[1]{%
      w\colon\NatSeg{\WordLength w}\to A}{%
      \FormulaRefAuto{FiniteWordTypingAtLength}[thm]{1}}
    \proofstepwidestar[1,2]{%
      w\colon\NatSeg 1\to A}{%
      \rIE{2,3}}
    \proofstepwidestar[]{%
      0\in\NatSeg 1}{%
      \FormulaRefAuto{PeanoNatSegOneMembership}[thm]}
    \proofstepwidestar[1,2]{%
      w(0)\in A}{%
      \FormulaRefAuto{FiniteWordLengthOneLetter}[thm]{1,2}}
    \proofstepwidestar[1,2]{%
      \LetterWord{w(0)}\colon\NatSeg 1\to A}{%
      \FormulaRefAuto{LetterWordTyping}[thm]{6}}
    \proofstepwidestar[1,2]{%
      \LetterWord{w(0)}(0)=w(0)}{%
      \FormulaRefAuto{LetterWordValueZero}[thm]{6}}
    \proofstepwidestar[9]{%
      i\in\NatSeg 1}{%
      \rA}
    \proofstepwidestar[9]{%
      i=0}{%
      \FormulaRefAuto{PeanoNatSegOneMembership}[thm]{9}}
    \proofstepwidestar[1,2,9]{%
      w(i)=\LetterWord{w(0)}(i)}{%
      \FormulaRefAuto{a=b\vdash b=a}[thm]{\rIE{10,8}}}
    \proofstepwidestar[1,2]{%
      \forall i\in\NatSeg 1\,w(i)=\LetterWord{w(0)}(i)}{%
      \rUI{\rRI{9,11}}}
    \proofstepwidestar[1,2]{%
      w=\LetterWord{w(0)}}{%
      \FormulaRefAuto{FunctionExtensionality}[thm]{4,7,12}}
  \closeproofpartwideindR
\end{tabproofsplitwide}


\Needspace{16\baselineskip}
\FormulaThmDeltaK[Länge nach Anhängen eines Buchstabens]{%
  u\in\FiniteWords A\dsep a\in A\vdash
    \WordLength{\WordConcat u{\LetterWord a}}=\WordLength u+1%
}{WordAppendLetterLength}{%
  \DeltaRow{Mengen}{A\dsep u\dsep a}
}
\begin{tabproofwide}
  \proofstepwidestar[1]{u\in\FiniteWords A}{\rA}
  \proofstepwidestar[2]{a\in A}{\rA}
  \proofstepwidestar[2]{\LetterWord a\in\NonemptyWords A}{%
    \FormulaRefAuto{LetterWordNonempty}[thm]{2}}
  \proofstepwidestar[2]{\LetterWord a\in\FiniteWords A}{%
    \FormulaRefAuto{NonemptyWordFinite}[thm]{3}}
  \proofstepwidestar[2]{\WordLength{\LetterWord a}=1}{%
    \FormulaRefAuto{LetterWordLengthOne}[thm]{2}}
  \proofstepwidestar[1,2]{%
    \WordLength{\WordConcat u{\LetterWord a}}
      =\WordLength u+\WordLength{\LetterWord a}}{%
    \FormulaRefAuto{WordConcatenationLength}[thm]{1,4}}
  \proofstepwidestar[1,2]{%
    \WordLength{\WordConcat u{\LetterWord a}}=\WordLength u+1}{\rIE{5,6}}
\end{tabproofwide}

\noindent\begin{minipage}{\linewidth}
\FormulaThmDeltaK[Letzter Wert nach Anhängen eines Buchstabens]{%
  u\in\FiniteWords A\dsep a\in A\vdash
    \WordConcat u{\LetterWord a}(\WordLength u)=a%
}{WordAppendLetterLastCoordinate}{%
  \DeltaRow{Mengen}{A\dsep u\dsep a}
}
\end{minipage}\par

\begin{tabproofwide}
  \proofstepwidestar[1]{u\in\FiniteWords A}{\rA}
  \proofstepwidestar[2]{a\in A}{\rA}
  \proofstepwidestar[2]{\LetterWord a\in\NonemptyWords A}{%
    \FormulaRefAuto{LetterWordNonempty}[thm]{2}}
  \proofstepwidestar[2]{\LetterWord a\in\FiniteWords A}{%
    \FormulaRefAuto{NonemptyWordFinite}[thm]{3}}
  \proofstepwidestar[2]{\WordLength{\LetterWord a}=1}{%
    \FormulaRefAuto{LetterWordLengthOne}[thm]{2}}
  \proofstepwidestar[]{0\in\NatSeg 1}{%
    \FormulaRefAuto{PeanoNatSegOneMembership}[thm]}
  \proofstepwidestar[2]{0\in\NatSeg{\WordLength{\LetterWord a}}}{\rIE{5,6}}
  \proofstepwidestar[1,2]{%
    \WordConcat u{\LetterWord a}(\WordLength u+0)=\LetterWord a(0)}{%
    \FormulaRefAuto{WordConcatenationShiftedCoordinate}[thm]{1,4,7}}
  \proofstepwidestar[2]{\LetterWord a(0)=a}{%
    \FormulaRefAuto{LetterWordValueZero}[thm]{2}}
  \proofstepwidestar[1]{\WordLength u\in\mathbb N}{%
    \FormulaRefAuto{FiniteWordLengthNatural}[thm]{1}}
  \proofstepwidestar[1]{\WordLength u+0=\WordLength u}{%
    \FormulaRefAuto{PeanoAddRightZero}[thm]{10}}
  \proofstepwidestar[1,2]{\WordConcat u{\LetterWord a}(\WordLength u)=a}{%
    \rIE{11,8,9}}
\end{tabproofwide}

\Needspace{22\baselineskip}
\noindent\begin{minipage}{\linewidth}
\FormulaThmDeltaK[Eindeutige Rechtszerlegung eines Nichtleerwortes]{%
  w\in\NonemptyWords{A}\vdash
  \exists!u\in\FiniteWords{A}\,
  \exists!a\in A\,
    \bigl(w=\WordConcat{u}{\LetterWord{a}}\bigr)%
}{NonemptyWordRightDecomposition}{%
  \DeltaRow{Mengen}{A\dsep w\dsep u\dsep a}
}
\end{minipage}\par

Die eindeutige Endzerlegung folgt nun aus dem abstrakten
Zerlegungssatz und W2. Die zuvor für den Modellnachweis verwendete
Restriktion eines Vorkommensgraphen wird dafür nicht erneut benötigt.

\begin{tabproofsplitwide}
  \proofpartwideindR[Existenz]{%
    w\in\NonemptyWords A\vdash
    \exists u\in\FiniteWords A\,\exists a\in A\;
      w=\WordConcat u{\LetterWord a}
  }{%
    w\in\NonemptyWords A\vdash
    \exists u\in\FiniteWords A\,\exists a\in A\;
      w=\WordConcat u{\LetterWord a}
  }[NonemptyWordRightDecompositionExistence]
    \proofstepwidestar[1]{w\in\NonemptyWords A}{\rA}
    \proofstepwidestar[1]{w\in\FiniteWords A\setminus\{\EmptyWord A\}}{\FormulaRefAuto{NonemptyWordSetDef}[def]{1}}
    \proofstepwidestar[1]{w\in\FiniteWords A\land w\notin\{\EmptyWord A\}}{\FormulaRefAuto{x\in A\setminus B\eqvdash x\in A\land x\notin B}[thm]{2}}
    \proofstepwidestar[1]{w\ne\EmptyWord A}{\FormulaRefAuto{x \notin \{a\} \eqvdash x \neq a}[thm]{\rAEb{3}}}
    \proofstepwidestar[]{\mathsf{WortStr}_A(\FiniteWords A,\EmptyWord A,\sigma_A)}{\FormulaRefAuto{WordStructureConcreteModel}[thm]}
    \proofstepwidestar[1]{w=\EmptyWord A\lor\exists!p\in\FiniteWords A\times A\;w=\sigma_A(p)}{\FormulaRefAuto{WordStructureDecomposition}[thm]{5,\rAEa{3}}}
    \proofstepwidestar[1]{\exists!p\in\FiniteWords A\times A\;w=\sigma_A(p)}{\FormulaRefAuto{P\lor Q,\neg P\vdash Q}[thm]{6,4}}
    \proofstepwidestar[1]{\exists p\in\FiniteWords A\times A\;w=\sigma_A(p)}{\FormulaRefAuto{\exists! x\,P(x)\vdash \exists x\,P(x)}[thm]{7}}
    \proofstepwidestar[9]{p\in\FiniteWords A\times A\land w=\sigma_A(p)}{\rA}
    \proofstepwidestar[9]{\exists u\in\FiniteWords A\,\exists a\in A\;p=(u,a)}{\FormulaRefAuto{x\in A\times B\eqvdash\exists a\in A\,\exists b\in B\,x=(a,b)}[thm]{\rAEa{9}}}
    \proofstepwidestar[11]{u\in\FiniteWords A\land\exists a\in A\;p=(u,a)}{\rA}
    \proofstepwidestar[12]{a\in A\land p=(u,a)}{\rA}
    \proofstepwidestar[9,12]{w=\sigma_A(u,a)}{\rIE{\rAEb{12},\rAEb{9}}}
    \proofstepwidestar[11,12]{\sigma_A(u,a)=\WordConcat u{\LetterWord a}}{\FormulaRefAuto{WordModelSuccessorConcatenation}[thm]{\rAEa{11},\rAEa{12}}}
    \proofstepwidestar[9,11,12]{w=\WordConcat u{\LetterWord a}}{\rIE{14,13}}
    \proofstepwidestar[9,11,12]{\exists a\in A\;w=\WordConcat u{\LetterWord a}}{\rEI{\rAI{\rAEa{12},15}}}
    \proofstepwidestar[9,11]{\exists a\in A\;w=\WordConcat u{\LetterWord a}}{\rEE{\rAEb{11},12,16}}
    \proofstepwidestar[9,11]{\exists u\in\FiniteWords A\,\exists a\in A\;w=\WordConcat u{\LetterWord a}}{\rEI{\rAI{\rAEa{11},17}}}
    \proofstepwidestar[9]{\exists u\in\FiniteWords A\,\exists a\in A\;w=\WordConcat u{\LetterWord a}}{\rEE{10,11,18}}
    \proofstepwidestar[1]{\exists u\in\FiniteWords A\,\exists a\in A\;w=\WordConcat u{\LetterWord a}}{\rEE{8,9,19}}
  \closeproofpartwideindR

  \proofpartwideindR[Eindeutigkeit]{%
    \begin{aligned}[t]
      w\in\NonemptyWords A\vdash
      &\forall u,v\in\FiniteWords A\,\forall a,b\in A\\[-2pt]
      &\quad\bigl((w=\WordConcat u{\LetterWord a}\land
        w=\WordConcat v{\LetterWord b})\\[-2pt]
      &\qquad\rightarrow u=v\land a=b\bigr)
    \end{aligned}
  }{%
    \begin{aligned}[t]
      w\in\NonemptyWords A\vdash
      &\forall u,v\in\FiniteWords A\,\forall a,b\in A\\[-2pt]
      &\quad\bigl((w=\WordConcat u{\LetterWord a}\land
        w=\WordConcat v{\LetterWord b})\\[-2pt]
      &\qquad\rightarrow u=v\land a=b\bigr)
    \end{aligned}
  }[NonemptyWordRightDecompositionJointUniqueness]
    \proofstepwidestar[1]{w\in\NonemptyWords A}{\rA}
    \proofstepwidestar[2]{u\in\FiniteWords A}{\rA}
    \proofstepwidestar[3]{v\in\FiniteWords A}{\rA}
    \proofstepwidestar[4]{a\in A}{\rA}
    \proofstepwidestar[5]{b\in A}{\rA}
    \proofstepwidestar[6]{w=\WordConcat u{\LetterWord a}\land w=\WordConcat v{\LetterWord b}}{\rA}
    \proofstepwidestar[2,4]{\sigma_A(u,a)=\WordConcat u{\LetterWord a}}{\FormulaRefAuto{WordModelSuccessorConcatenation}[thm]{2,4}}
    \proofstepwidestar[3,5]{\sigma_A(v,b)=\WordConcat v{\LetterWord b}}{\FormulaRefAuto{WordModelSuccessorConcatenation}[thm]{3,5}}
    \proofstepwidestar[2,4,6]{w=\sigma_A(u,a)}{\rIE{7,\rAEa{6}}}
    \proofstepwidestar[3,5,6]{w=\sigma_A(v,b)}{\rIE{8,\rAEb{6}}}
    \proofstepwidestar[2,3,4,5,6]{\sigma_A(u,a)=\sigma_A(v,b)}{\FormulaRefAuto{a=b\dsep a=c\vdash b=c}[thm]{9,10}}
    \proofstepwidestar[]{\mathsf{WortStr}_A(\FiniteWords A,\EmptyWord A,\sigma_A)}{\FormulaRefAuto{WordStructureConcreteModel}[thm]}
    \proofstepwidestar[2,3,4,5,6]{u=v\land a=b}{\FormulaRefAuto{WordStructureInjectivityAxiom}[ax]{12,\rAI{2,3},\rAI{4,5},11}}
    \proofstepwidestar[]{\begin{aligned}[t]
      &\forall u,v\in\FiniteWords A\,\forall a,b\in A\\[-2pt]
      &\quad((w=\WordConcat u{\LetterWord a}\land w=\WordConcat v{\LetterWord b})\\[-2pt]
      &\qquad\rightarrow u=v\land a=b)
    \end{aligned}}{\rUI{\rRI{2,\rUI{\rRI{3,\rUI{\rRI{4,\rUI{\rRI{5,\rRI{6,13}}}}}}}}}}
  \closeproofpartwideindR

  \proofpartwide{Eindeutige Wahl des Rumpfs und des Endbuchstabens}
    \proofstepwidestar[1]{w\in\NonemptyWords A}{\rA}
    \proofstepwidestar[1]{\exists u\in\FiniteWords A\,\exists a\in A\;
      w=\WordConcat u{\LetterWord a}}{%
      \FormulaRefAuto{NonemptyWordRightDecompositionExistence}[thm]{1}}
    \proofstepwidestar[1]{\begin{aligned}[t]
      &\forall u,v\in\FiniteWords A\,\forall a,b\in A\\[-2pt]
      &\quad\bigl((w=\WordConcat u{\LetterWord a}\land
        w=\WordConcat v{\LetterWord b})\\[-2pt]
      &\qquad\rightarrow u=v\land a=b\bigr)
    \end{aligned}}{\FormulaRefAuto{NonemptyWordRightDecompositionJointUniqueness}[thm]{1}}
    \proofstepwidestar[1]{\exists!u\in\FiniteWords A\,\exists!a\in A\;
      w=\WordConcat u{\LetterWord a}}{\FormulaRefAuto{NestedUniqueExistenceFromJointUniqueness}[thm]{2,3}}
  \closeproofpartwide
\end{tabproofsplitwide}


% Definierte Werkzeuge für die bereits bewiesene eindeutige Endzerlegung.
\subsection{Rumpf und Endbuchstabe}

Der Rumpf entsteht durch Entfernen des letzten Buchstabens. Er darf leer
sein; der Endbuchstabe liegt im Alphabet. Die eindeutige Endzerlegung
rechtfertigt beide Funktionen. Ihre Rechenregeln machen die Zerlegung
anschließend unmittelbar benutzbar.

\Needspace{18\baselineskip}
\FormulaDefDeltaK[Rumpffunktion eines nichtleeren Wortes]{%
  \begin{aligned}[t]
    &\WordRumpfMap A\colon\NonemptyWords A\to\FiniteWords A,\\[-2pt]
    &\forall w\in\NonemptyWords A\,\Bigl(\WordRumpf A w\coloneqq
      \iota u\in\FiniteWords A\,\exists!a\in A\,
        w=\WordConcat u{\LetterWord a}\Bigr)
  \end{aligned}
}{WordRumpfDef}{%
  \DeltaRow{Mengen}{A\dsep w\dsep u\dsep a}
  \DeltaRow{Funktionen}{\WordRumpfMap A}
  \DeltaRow{\textbf{Neues Symbol}}{}
  \DeltaPrem{Rumpf}{\WordRumpf A w}
}
\begin{tabproofwide}
  \proofstepwidestar[1]{w\in\NonemptyWords A}{\rA}
  \proofstepwidestar[1]{\exists!u\in\FiniteWords A\,\exists!a\in A\,
    w=\WordConcat u{\LetterWord a}}{%
    \FormulaRefAuto{NonemptyWordRightDecomposition}[thm]{1}}
  \proofstepwidestar[1]{\Bigl(\iota u\in\FiniteWords A\,\exists!a\in A\,
    w=\WordConcat u{\LetterWord a}\Bigr)\in\FiniteWords A}{%
    \rAEa{\FormulaRefAuto{IotaSymbolDef}[def]{2}}}
  \proofstepwidestar[]{\begin{aligned}[t]
    &\forall w\in\NonemptyWords A\,\Bigl(\\[-2pt]
    &\quad\iota u\in\FiniteWords A\,\exists!a\in A\,
      w=\WordConcat u{\LetterWord a}\Bigr)\in\FiniteWords A
  \end{aligned}}{%
    \rUI{\rRI{1,3}}}
  \proofstepwidestar[]{\begin{aligned}[t]
    &\exists!F\,\Bigl(F\colon\NonemptyWords A\to\FiniteWords A\land{}\\[-2pt]
    &\quad\forall w\in\NonemptyWords A\,F(w)=\\[-2pt]
    &\qquad\iota u\in\FiniteWords A\,\exists!a\in A\,
        w=\WordConcat u{\LetterWord a}\Bigr)
  \end{aligned}}{\FormulaRefAuto{TypedFunctionDefinitionPrinciple}[thm]{4}}
\end{tabproofwide}

\Needspace{10\baselineskip}
\FormulaThmDeltaK[Typisierung der Rumpffunktion]{%
  \WordRumpfMap A\colon\NonemptyWords A\to\FiniteWords A
}{WordRumpfFunction}{\DeltaRow{Mengen}{A}}
\begin{tabproofwide}
  \proofstepwidestar[]{\WordRumpfMap A\colon\NonemptyWords A\to\FiniteWords A}{%
    \FormulaRefAuto{WordRumpfDef}[def]}
\end{tabproofwide}

\FormulaThmDeltaK[Der Rumpf ist ein endliches Wort]{%
  w\in\NonemptyWords A\vdash\WordRumpf A w\in\FiniteWords A
}{WordRumpfTyping}{\DeltaRow{Mengen}{A\dsep w}}
\begin{tabproofwide}
  \proofstepwidestar[1]{w\in\NonemptyWords A}{\rA}
  \proofstepwidestar[]{\WordRumpfMap A\colon\NonemptyWords A\to\FiniteWords A}{%
    \FormulaRefAuto{WordRumpfFunction}[thm]}
  \proofstepwidestar[1]{\WordRumpf A w\in\FiniteWords A}{%
    \FormulaRefAuto{F\colon A\to B\dsep x\in A\vdash F(x)\in B}[thm]{2,1}}
\end{tabproofwide}

\FormulaThmDeltaK[Zum Rumpf gehört genau ein Endbuchstabe]{%
  w\in\NonemptyWords A\vdash
    \exists!a\in A\,w=\WordConcat{\WordRumpf A w}{\LetterWord a}
}{WordRumpfCharacterization}{\DeltaRow{Mengen}{A\dsep w\dsep u\dsep a}}
\begin{tabproofwide}
  \proofstepwidestar[1]{w\in\NonemptyWords A}{\rA}
  \proofstepwidestar[1]{\exists!u\in\FiniteWords A\,\exists!a\in A\,
    w=\WordConcat u{\LetterWord a}}{%
    \FormulaRefAuto{NonemptyWordRightDecomposition}[thm]{1}}
  \proofstepwidestar[1]{\WordRumpf A w=
    \iota u\in\FiniteWords A\,\exists!a\in A\,
      w=\WordConcat u{\LetterWord a}}{\FormulaRefAuto{WordRumpfDef}[def]{1}}
  \proofstepwidestar[1]{\exists!a\in A\,
    w=\WordConcat{\WordRumpf A w}{\LetterWord a}}{%
    \rIE{\FormulaRefAuto{x=y\vdash y=x}[thm]{3},
      \rAEb{\FormulaRefAuto{IotaSymbolDef}[def]{2}}}}
\end{tabproofwide}

\FormulaDefDeltaK[Endbuchstabenfunktion]{%
  \begin{aligned}[t]
    &\WordEndMap A\colon\NonemptyWords A\to A,\\[-2pt]
    &\forall w\in\NonemptyWords A\,\Bigl(\WordEnd A w\coloneqq
      \iota a\in A\,w=\WordConcat{\WordRumpf A w}{\LetterWord a}\Bigr)
  \end{aligned}
}{WordEndDef}{%
  \DeltaRow{Mengen}{A\dsep w\dsep a}
  \DeltaRow{Funktionen}{\WordEndMap A}
  \DeltaRow{\textbf{Neues Symbol}}{}
  \DeltaPrem{Endbuchstabe}{\WordEnd A w}
}
\begin{tabproofwide}
  \proofstepwidestar[1]{w\in\NonemptyWords A}{\rA}
  \proofstepwidestar[1]{\exists!a\in A\,
    w=\WordConcat{\WordRumpf A w}{\LetterWord a}}{%
    \FormulaRefAuto{WordRumpfCharacterization}[thm]{1}}
  \proofstepwidestar[1]{\Bigl(\iota a\in A\,
    w=\WordConcat{\WordRumpf A w}{\LetterWord a}\Bigr)\in A}{%
    \rAEa{\FormulaRefAuto{IotaSymbolDef}[def]{2}}}
  \proofstepwidestar[]{\begin{aligned}[t]
    &\forall w\in\NonemptyWords A\,\Bigl(\\[-2pt]
    &\quad\iota a\in A\,
      w=\WordConcat{\WordRumpf A w}{\LetterWord a}\Bigr)\in A
  \end{aligned}}{%
    \rUI{\rRI{1,3}}}
  \proofstepwidestar[]{\begin{aligned}[t]
    &\exists!F\,\Bigl(F\colon\NonemptyWords A\to A\land{}\\[-2pt]
    &\quad\forall w\in\NonemptyWords A\,F(w)=\\[-2pt]
    &\qquad\iota a\in A\,w=\WordConcat{\WordRumpf A w}{\LetterWord a}\Bigr)
  \end{aligned}}{\FormulaRefAuto{TypedFunctionDefinitionPrinciple}[thm]{4}}
\end{tabproofwide}

\FormulaThmDeltaK[Typisierung der Endbuchstabenfunktion]{%
  \WordEndMap A\colon\NonemptyWords A\to A
}{WordEndFunction}{\DeltaRow{Mengen}{A}}
\begin{tabproofwide}
  \proofstepwidestar[]{\WordEndMap A\colon\NonemptyWords A\to A}{%
    \FormulaRefAuto{WordEndDef}[def]}
\end{tabproofwide}

\FormulaThmDeltaK[Der Endbuchstabe liegt im Alphabet]{%
  w\in\NonemptyWords A\vdash\WordEnd A w\in A
}{WordEndTyping}{\DeltaRow{Mengen}{A\dsep w}}
\begin{tabproofwide}
  \proofstepwidestar[1]{w\in\NonemptyWords A}{\rA}
  \proofstepwidestar[]{\WordEndMap A\colon\NonemptyWords A\to A}{%
    \FormulaRefAuto{WordEndFunction}[thm]}
  \proofstepwidestar[1]{\WordEnd A w\in A}{%
    \FormulaRefAuto{F\colon A\to B\dsep x\in A\vdash F(x)\in B}[thm]{2,1}}
\end{tabproofwide}

\FormulaThmDeltaK[Rumpf und Endbuchstabe setzen das Wort wieder zusammen]{%
  w\in\NonemptyWords A\vdash
    w=\WordConcat{\WordRumpf A w}{\LetterWord{\WordEnd A w}}
}{WordRightDecompositionEquation}{\DeltaRow{Mengen}{A\dsep w\dsep a}}
\begin{tabproofwide}
  \proofstepwidestar[1]{w\in\NonemptyWords A}{\rA}
  \proofstepwidestar[1]{\exists!a\in A\,
    w=\WordConcat{\WordRumpf A w}{\LetterWord a}}{%
    \FormulaRefAuto{WordRumpfCharacterization}[thm]{1}}
  \proofstepwidestar[1]{\WordEnd A w=
    \iota a\in A\,w=\WordConcat{\WordRumpf A w}{\LetterWord a}}{%
    \FormulaRefAuto{WordEndDef}[def]{1}}
  \proofstepwidestar[1]{w=\WordConcat{\WordRumpf A w}{\LetterWord{\WordEnd A w}}}{%
    \rIE{\FormulaRefAuto{x=y\vdash y=x}[thm]{3},
      \rAEb{\FormulaRefAuto{IotaSymbolDef}[def]{2}}}}
\end{tabproofwide}

\Needspace{10\baselineskip}
\FormulaThmDeltaK[Eine gegebene Endzerlegung bestimmt Rumpf und Endbuchstabe]{%
  \begin{aligned}[t]
    &w\in\NonemptyWords A\dsep u\in\FiniteWords A\dsep a\in A\dsep
      w=\WordConcat u{\LetterWord a}\vdash{}\\[-2pt]
    &\WordRumpf A w=u\land\WordEnd A w=a
  \end{aligned}
}{WordRightDecompositionIdentification}{\DeltaRow{Mengen}{A\dsep w\dsep u\dsep v\dsep r\dsep a\dsep b\dsep c}}
\begin{tabproofwide}
  \proofstepwidestar[1]{w\in\NonemptyWords A}{\rA}
  \proofstepwidestar[2]{u\in\FiniteWords A}{\rA}
  \proofstepwidestar[3]{a\in A}{\rA}
  \proofstepwidestar[4]{w=\WordConcat u{\LetterWord a}}{\rA}
  \proofstepwidestar[1]{\WordRumpf A w\in\FiniteWords A}{%
    \FormulaRefAuto{WordRumpfTyping}[thm]{1}}
  \proofstepwidestar[1]{\WordEnd A w\in A}{%
    \FormulaRefAuto{WordEndTyping}[thm]{1}}
  \proofstepwidestar[1]{w=\WordConcat{\WordRumpf A w}{\LetterWord{\WordEnd A w}}}{%
    \FormulaRefAuto{WordRightDecompositionEquation}[thm]{1}}
  \proofstepwidestar[1]{\begin{aligned}[t]
    &\forall r\in\FiniteWords A\,\forall v\in\FiniteWords A\,
      \forall c\in A\,\forall b\in A\,\Bigl(\\[-2pt]
    &\quad\bigl(w=\WordConcat r{\LetterWord c}\land
      w=\WordConcat v{\LetterWord b}\bigr)\\[-2pt]
    &\quad\rightarrow(r=v\land c=b)\Bigr)
  \end{aligned}}{\FormulaRefAuto{NonemptyWordRightDecompositionJointUniqueness}[thm]{1}}
  \proofstepwidestar[1]{\begin{aligned}[t]
    &\forall v\in\FiniteWords A\,\forall c\in A\,\forall b\in A\,\Bigl(\\[-2pt]
    &\quad\bigl(w=\WordConcat{\WordRumpf A w}{\LetterWord c}\land
      w=\WordConcat v{\LetterWord b}\bigr)\\[-2pt]
    &\quad\rightarrow(\WordRumpf A w=v\land c=b)\Bigr)
  \end{aligned}}{\rRE{5,\rUE{8}}}
  \proofstepwidestar[1,2]{\begin{aligned}[t]
    &\forall c\in A\,\forall b\in A\,\Bigl(\\[-2pt]
    &\quad\bigl(w=\WordConcat{\WordRumpf A w}{\LetterWord c}\land
      w=\WordConcat u{\LetterWord b}\bigr)\\[-2pt]
    &\quad\rightarrow(\WordRumpf A w=u\land c=b)\Bigr)
  \end{aligned}}{\rRE{2,\rUE{9}}}
  \proofstepwidestar[1,2]{\begin{aligned}[t]
    &\forall b\in A\,\Bigl(\bigl(
      w=\WordConcat{\WordRumpf A w}{\LetterWord{\WordEnd A w}}\land{}\\[-2pt]
    &\quad w=\WordConcat u{\LetterWord b}\bigr)\\[-2pt]
    &\quad\rightarrow(\WordRumpf A w=u\land\WordEnd A w=b)\Bigr)
  \end{aligned}}{\rRE{6,\rUE{10}}}
  \proofstepwidestar[1,2,3]{\begin{aligned}[t]
    &\bigl(w=\WordConcat{\WordRumpf A w}{\LetterWord{\WordEnd A w}}\land{}\\[-2pt]
    &\quad w=\WordConcat u{\LetterWord a}\bigr)\\[-2pt]
    &\quad\rightarrow(\WordRumpf A w=u\land\WordEnd A w=a)
  \end{aligned}}{\rRE{3,\rUE{11}}}
  \proofstepwidestar[1,4]{\begin{aligned}[t]
    &w=\WordConcat{\WordRumpf A w}{\LetterWord{\WordEnd A w}}\land{}\\[-2pt]
    &\quad w=\WordConcat u{\LetterWord a}
  \end{aligned}}{\rAI{7,4}}
  \proofstepwidestar[1,2,3,4]{\WordRumpf A w=u\land\WordEnd A w=a}{\rRE{13,12}}
\end{tabproofwide}

\Needspace{10\baselineskip}
\FormulaThmDeltaK[Ein angehängter Buchstabe macht ein Wort nichtleer]{%
  u\in\FiniteWords A\dsep a\in A\vdash\WordConcat u{\LetterWord a}\in\NonemptyWords A
}{WordAppendNonempty}{\DeltaRow{Mengen}{A\dsep u\dsep a}}
\begin{tabproofwide}
  \proofstepwidestar[1]{u\in\FiniteWords A}{\rA}
  \proofstepwidestar[2]{a\in A}{\rA}
  \proofstepwidestar[2]{\LetterWord a\in\NonemptyWords A}{\FormulaRefAuto{LetterWordNonempty}[thm]{2}}
  \proofstepwidestar[2]{\LetterWord a\in\FiniteWords A}{\FormulaRefAuto{NonemptyWordFinite}[thm]{3}}
  \proofstepwidestar[1,2]{\WordConcat u{\LetterWord a}\in\FiniteWords A}{%
    \FormulaRefAuto{WordConcatenationClosure}[thm]{\rAI{1,4}}}
  \proofstepwidestar[1]{\WordLength u\in\mathbb N}{\FormulaRefAuto{FiniteWordLengthNatural}[thm]{1}}
  \proofstepwidestar[1]{\WordLength u+1\neq0}{\FormulaRefAuto{PeanoAddOneNonzero}[thm]{6}}
  \proofstepwidestar[1,2]{\WordLength{\WordConcat u{\LetterWord a}}=\WordLength u+1}{%
    \FormulaRefAuto{WordAppendLetterLength}[thm]{1,2}}
  \proofstepwidestar[1,2]{\WordLength u+1=\WordLength{\WordConcat u{\LetterWord a}}}{%
    \FormulaRefAuto{a=b\vdash b=a}[thm]{8}}
  \proofstepwidestar[1,2]{\WordLength{\WordConcat u{\LetterWord a}}\neq0}{\rIE{9,7}}
  \proofstepwidestar[1,2]{\WordConcat u{\LetterWord a}\in\NonemptyWords A}{%
    \FormulaRefAuto{NonemptyWordMembershipIntroduction}[thm]{\rAI{5,10}}}
\end{tabproofwide}

\FormulaThmDeltaK[Anhängen und Entfernen eines Endbuchstabens]{%
  u\in\FiniteWords A\dsep a\in A\vdash
    \WordRumpf A{\WordConcat u{\LetterWord a}}=u
}{WordRumpfAppend}{\DeltaRow{Mengen}{A\dsep u\dsep a}}
\begin{tabproofwide}
  \proofstepwidestar[1]{u\in\FiniteWords A}{\rA}
  \proofstepwidestar[2]{a\in A}{\rA}
  \proofstepwidestar[1,2]{\WordConcat u{\LetterWord a}\in\NonemptyWords A}{%
    \FormulaRefAuto{WordAppendNonempty}[thm]{1,2}}
  \proofstepwidestar[]{\WordConcat u{\LetterWord a}=\WordConcat u{\LetterWord a}}{\rII}
  \proofstepwidestar[1,2]{\WordRumpf A{\WordConcat u{\LetterWord a}}=u\land
    \WordEnd A{\WordConcat u{\LetterWord a}}=a}{%
    \FormulaRefAuto{WordRightDecompositionIdentification}[thm]{3,1,2,4}}
  \proofstepwidestar[1,2]{\WordRumpf A{\WordConcat u{\LetterWord a}}=u}{\rAEa{5}}
\end{tabproofwide}

\FormulaThmDeltaK[Der angehängte Buchstabe ist der Endbuchstabe]{%
  u\in\FiniteWords A\dsep a\in A\vdash
    \WordEnd A{\WordConcat u{\LetterWord a}}=a
}{WordEndAppend}{\DeltaRow{Mengen}{A\dsep u\dsep a}}
\begin{tabproofwide}
  \proofstepwidestar[1]{u\in\FiniteWords A}{\rA}
  \proofstepwidestar[2]{a\in A}{\rA}
  \proofstepwidestar[1,2]{\WordConcat u{\LetterWord a}\in\NonemptyWords A}{%
    \FormulaRefAuto{WordAppendNonempty}[thm]{1,2}}
  \proofstepwidestar[]{\WordConcat u{\LetterWord a}=\WordConcat u{\LetterWord a}}{\rII}
  \proofstepwidestar[1,2]{\WordRumpf A{\WordConcat u{\LetterWord a}}=u\land
    \WordEnd A{\WordConcat u{\LetterWord a}}=a}{%
    \FormulaRefAuto{WordRightDecompositionIdentification}[thm]{3,1,2,4}}
  \proofstepwidestar[1,2]{\WordEnd A{\WordConcat u{\LetterWord a}}=a}{\rAEb{5}}
\end{tabproofwide}

\FormulaThmDeltaK[Der Rumpf ist um einen Buchstaben kürzer]{%
  w\in\NonemptyWords A\vdash\WordLength{\WordRumpf A w}+1=\WordLength w
}{WordRumpfLength}{\DeltaRow{Mengen}{A\dsep w}}
\begin{tabproofwide}
  \proofstepwidestar[1]{w\in\NonemptyWords A}{\rA}
  \proofstepwidestar[1]{\WordRumpf A w\in\FiniteWords A}{\FormulaRefAuto{WordRumpfTyping}[thm]{1}}
  \proofstepwidestar[1]{\WordEnd A w\in A}{\FormulaRefAuto{WordEndTyping}[thm]{1}}
  \proofstepwidestar[1]{\begin{aligned}[t]
    &\WordLength{\WordConcat{\WordRumpf A w}{\LetterWord{\WordEnd A w}}}\\[-2pt]
    &\quad=\WordLength{\WordRumpf A w}+1
  \end{aligned}}{\FormulaRefAuto{WordAppendLetterLength}[thm]{2,3}}
  \proofstepwidestar[1]{w=\WordConcat{\WordRumpf A w}{\LetterWord{\WordEnd A w}}}{%
    \FormulaRefAuto{WordRightDecompositionEquation}[thm]{1}}
  \proofstepwidestar[1]{\WordConcat{\WordRumpf A w}{\LetterWord{\WordEnd A w}}=w}{%
    \FormulaRefAuto{a=b\vdash b=a}[thm]{5}}
  \proofstepwidestar[1]{\WordLength w=\WordLength{\WordRumpf A w}+1}{\rIE{6,4}}
  \proofstepwidestar[1]{\WordLength{\WordRumpf A w}+1=\WordLength w}{%
    \FormulaRefAuto{a=b\vdash b=a}[thm]{7}}
\end{tabproofwide}

\FormulaThmDeltaK[Gleiche angehängte Wörter haben dieselben Bestandteile]{%
  \begin{aligned}[t]
    &u,v\in\FiniteWords A\dsep a,b\in A\dsep
      \WordConcat u{\LetterWord a}=\WordConcat v{\LetterWord b}\vdash{}\\[-2pt]
    &u=v\land a=b
  \end{aligned}
}{WordAppendJointInjectivity}{\DeltaRow{Mengen}{A\dsep u\dsep v\dsep a\dsep b}}
\begin{tabproofwide}
  \proofstepwidestar[1]{u,v\in\FiniteWords A}{\rA}
  \proofstepwidestar[2]{a,b\in A}{\rA}
  \proofstepwidestar[3]{\WordConcat u{\LetterWord a}=\WordConcat v{\LetterWord b}}{\rA}
  \proofstepwidestar[1,2]{\WordRumpf A{\WordConcat u{\LetterWord a}}=u}{%
    \FormulaRefAuto{WordRumpfAppend}[thm]{\rAEa{1},\rAEa{2}}}
  \proofstepwidestar[1,2]{\WordRumpf A{\WordConcat v{\LetterWord b}}=v}{%
    \FormulaRefAuto{WordRumpfAppend}[thm]{\rAEb{1},\rAEb{2}}}
  \proofstepwidestar[1,2,3]{\WordRumpf A{\WordConcat v{\LetterWord b}}=u}{\rIE{3,4}}
  \proofstepwidestar[1,2,3]{u=v}{\FormulaRefAuto{a=b\dsep a=c\vdash b=c}[thm]{6,5}}
  \proofstepwidestar[1,2]{\WordEnd A{\WordConcat u{\LetterWord a}}=a}{%
    \FormulaRefAuto{WordEndAppend}[thm]{\rAEa{1},\rAEa{2}}}
  \proofstepwidestar[1,2]{\WordEnd A{\WordConcat v{\LetterWord b}}=b}{%
    \FormulaRefAuto{WordEndAppend}[thm]{\rAEb{1},\rAEb{2}}}
  \proofstepwidestar[1,2,3]{\WordEnd A{\WordConcat v{\LetterWord b}}=a}{\rIE{3,8}}
  \proofstepwidestar[1,2,3]{a=b}{\FormulaRefAuto{a=b\dsep a=c\vdash b=c}[thm]{10,9}}
  \proofstepwidestar[1,2,3]{u=v\land a=b}{\rAI{7,11}}
\end{tabproofwide}


\subsection{Wortinduktion in Konkatenationsschreibweise}

Die Minimalität des konkreten Wortmodells wurde vor Einführung der Axiome
aus seiner Konstruktion als kleinste abgeschlossene Menge bewiesen.
Jetzt wenden wir das
axiomatische Induktionsprinzip W3 auf das konkrete Wortmodell an
und schreiben den Anfügungsschritt als Konkatenation.

\noindent\begin{minipage}{\linewidth}
\FormulaThmDeltaK[Strukturinduktion über alle endlichen Wörter]{%
  \begin{aligned}[t]
    &P(\EmptyWord A)\dsep{}\\[-2pt]
    &\forall u\in\FiniteWords A\,\forall a\in A\,
      (P(u)\rightarrow P(\WordConcat u{\LetterWord a}))\\[-2pt]
    &\vdash\forall w\in\FiniteWords A\,P(w)
  \end{aligned}
}{FiniteWordStructuralInduction}{%
  \DeltaRow{Mengen}{A\dsep u\dsep w\dsep a\dsep n}
  \DeltaRow{Prädikat}{P}
}
\end{minipage}\par
Wir wenden \FormulaRefAuto{WordStructureInduction}[ax] auf das Modell
\FormulaRefAuto{WordStructureConcreteModel}[thm] an. Die Brücke
\FormulaRefAuto{WordModelSuccessorConcatenation}[thm] übersetzt den
primitiven Nachfolger in das Anhängen eines Buchstabenwortes.
\begin{tabproofwide}
  \proofstepwidestar[1]{P(\EmptyWord A)}{\rA}
  \proofstepwidestar[2]{\forall u\in\FiniteWords A\,\forall a\in A\;(P(u)\rightarrow P(\WordConcat u{\LetterWord a}))}{\rA}
  \proofstepwidestar[]{\mathsf{WortStr}_A(\FiniteWords A,\EmptyWord A,\sigma_A)}{\FormulaRefAuto{WordStructureConcreteModel}[thm]}
  \proofstepwidestar[4]{u\in\FiniteWords A}{\rA}
  \proofstepwidestar[5]{a\in A}{\rA}
  \proofstepwidestar[2,4,5]{P(u)\rightarrow P(\WordConcat u{\LetterWord a})}{\rRE{5,\rUE{\rRE{4,\rUE{2}}}}}
  \proofstepwidestar[4,5]{\sigma_A(u,a)=\WordConcat u{\LetterWord a}}{\FormulaRefAuto{WordModelSuccessorConcatenation}[thm]{4,5}}
  \proofstepwidestar[2,4,5]{P(u)\rightarrow P(\sigma_A(u,a))}{\rIE{7,6}}
  \proofstepwidestar[2]{\forall u\in\FiniteWords A\,\forall a\in A\;(P(u)\rightarrow P(\sigma_A(u,a)))}{\rUI{\rRI{4,\rUI{\rRI{5,8}}}}}
  \proofstepwidestar[1,2]{\forall w\in\FiniteWords A\;P(w)}{\FormulaRefAuto{WordStructureInduction}[thm]{3,1,9}}
\end{tabproofwide}


\noindent\begin{minipage}{\linewidth}
\FormulaThmDeltaK[Die Wortmenge ist die kleinste unter Anfügung abgeschlossene Menge]{%
  \begin{aligned}[t]
    &\EmptyWord A\in C\dsep{}\\[-2pt]
    &\forall u\in C\cap\FiniteWords A\,\forall a\in A\,
       (u\cup\{(\WordLength u,a)\}\in C)\\[-2pt]
    &\vdash\FiniteWords A\subseteq C
  \end{aligned}
}{FiniteWordSetLeastAppendClosed}{%
  \DeltaRow{Mengen}{A\dsep C\dsep u\dsep a\dsep w}
}
\end{minipage}\par
Die Schnittmenge in der Voraussetzung stellt lediglich sicher, dass die
Wortlänge im Anfügungsterm definiert ist. Man kann insbesondere
\(C\subseteq\FiniteWords A\) voraussetzen; dann folgt \(C=\FiniteWords A\).
Zusammen mit dem Leerwort und der abgeschlossenen Anfügung ist dies die
gewünschte Erzeugungsregel für \(A^*\).
\begin{tabproofwide}
  \proofstepwidestar[1]{\EmptyWord A\in C}{\rA}
  \proofstepwidestar[2]{\forall u\in C\cap\FiniteWords A\,\forall a\in A\,
    (u\cup\{(\WordLength u,a)\}\in C)}{\rA}
  \proofstepwidestar[3]{u\in\FiniteWords A}{\rA}
  \proofstepwidestar[4]{a\in A}{\rA}
  \proofstepwidestar[5]{u\in C}{\rA}
  \proofstepwidestar[3,5]{u\in C\cap\FiniteWords A}{%
    \FormulaRefAuto{x\in A\cap B\eqvdash x\in A\land x\in B}[thm]{\rAI{5,3}}}
  \proofstepwidestar[2,3,4,5]{u\cup\{(\WordLength u,a)\}\in C}{\rRE{4,\rUE{\rRE{6,\rUE{2}}}}}
  \proofstepwidestar[3,4]{\WordConcat u{\LetterWord a}=u\cup\{(\WordLength u,a)\}}{%
    \FormulaRefAuto{WordAppendOccurrenceEquation}[thm]{3,4}}
  \proofstepwidestar[2,3,4,5]{\WordConcat u{\LetterWord a}\in C}{\rIE{\FormulaRefAuto{a=b\vdash b=a}[thm]{8},7}}
  \proofstepwidestar[2]{\forall u\in\FiniteWords A\,\forall a\in A\,
    (u\in C\rightarrow\WordConcat u{\LetterWord a}\in C)}{\rUI{\rRI{3,\rUI{\rRI{4,\rRI{5,9}}}}}}
  \proofstepwidestar[1,2]{\forall w\in\FiniteWords A\,(w\in C)}{\FormulaRefAuto{FiniteWordStructuralInduction}[thm]{1,10}}
  \proofstepwidestar[1,2]{\FiniteWords A\subseteq C}{%
    \FormulaRefAuto{A\subseteq B\coloneq\forall x\,(x\in A\rightarrow x\in B)}[def]{11}}
\end{tabproofwide}




\subsection{Induktion und Rekursion über nichtleere Wörter}

Bei den natürlichen Zahlen kann eine Induktion auch bei \(1\) beginnen.
Für nichtleere Wörter treten an die Stelle dieses einen Anfangs alle
Buchstabenwörter \(\LetterWord a\). Deshalb ist der Induktionsanfang
für jedes \(a\in A\) zu zeigen. Der Schritt fügt wie zuvor einen
beliebigen Buchstaben an ein bereits behandeltes Wort an.

\par\noindent\begin{minipage}{\linewidth}
\FormulaThmDeltaK[Induktion über nichtleere Wörter]{%
  \begin{aligned}[t]
    &\forall a\in A\,P(\LetterWord{a})\dsep{}\\[-2pt]
    &\forall u\in\NonemptyWords{A}\,\forall a\in A\,
      \bigl(P(u)\rightarrow
        P(\WordConcat{u}{\LetterWord{a}})\bigr)\\[-2pt]
    &\vdash\forall w\in\NonemptyWords{A}\,P(w)
  \end{aligned}%
}{NonemptyWordInduction}{%
  \DeltaRow{Mengen}{A\dsep a\dsep u\dsep w}
  \DeltaRow{Prädikate}{P}
}
\end{minipage}\par
Die bereits bewiesene Strukturinduktion wird auf
\(R(w)\coloneqq(w\in\NonemptyWords A\rightarrow P(w))\) angewandt.
Beim Leerwort ist diese Aussage leer erfüllt; im Anfügungsschritt
unterscheiden wir, ob der Rumpf leer ist.
\begin{tabproofsplitwide}
  \proofpartwide{Leerwort}
    \proofstepwidestar[1]{\EmptyWord A\in\NonemptyWords A}{\rA}
    \proofstepwidestar[1]{\WordLength{\EmptyWord A}\neq0}{%
      \rAEb{\FormulaRefAuto{NonemptyWordMembership}[thm]{1}}}
    \proofstepwidestar[]{\WordLength{\EmptyWord A}=0}{\FormulaRefAuto{EmptyWordLengthZero}[thm]}
    \proofstepwidestar[1]{\bot}{\rBI{3,2}}
    \proofstepwidestar[1]{P(\EmptyWord A)}{\FormulaRefAuto{\bot\vdash P}[thm]{4}}
    \proofstepwidestar[]{R(\EmptyWord A)}{\rRI{1,5}}
  \closeproofpartwide
  \Needspace{5\baselineskip}
  \proofpartwide{Anfügung}
    \proofstepwidestar[1]{\forall a\in A\,P(\LetterWord a)}{\rA}
    \proofstepwidestar[2]{\forall u\in\NonemptyWords A\,\forall a\in A\,
      (P(u)\rightarrow P(\WordConcat u{\LetterWord a}))}{\rA}
    \proofstepwidestar[3]{u\in\FiniteWords A}{\rA}
    \proofstepwidestar[4]{a\in A}{\rA}
    \proofstepwidestar[5]{R(u)}{\rA}
    \proofstepwidestar[]{u=\EmptyWord A\lor u\neq\EmptyWord A}{\FormulaRefAuto{P\lor\neg P}[thm]}
    \proofstepwidestar[7]{u=\EmptyWord A}{\rA}
    \proofstepwidestar[1,4]{P(\LetterWord a)}{\rRE{4,\rUE{1}}}
    \proofstepwidestar[4]{\LetterWord a\in\NonemptyWords A}{\FormulaRefAuto{LetterWordNonempty}[thm]{4}}
    \proofstepwidestar[4]{\LetterWord a\in\FiniteWords A}{\FormulaRefAuto{NonemptyWordFinite}[thm]{9}}
    \proofstepwidestar[4]{\WordConcat{\EmptyWord A}{\LetterWord a}=\LetterWord a}{%
      \FormulaRefAuto{WordConcatenationLeftIdentity}[thm]{10}}
    \proofstepwidestar[1,4,7]{P(\WordConcat u{\LetterWord a})}{\rIE{7,11,8}}
    \proofstepwidestar[13]{u\neq\EmptyWord A}{\rA}
    \proofstepwidestar[3]{u\neq\EmptyWord A\leftrightarrow\WordLength u\neq0}{%
      \FormulaRefAuto{FiniteWordLengthNonzeroIffNonempty}[thm]{3}}
    \proofstepwidestar[3,13]{\WordLength u\neq0}{\rRE{13,\rLREa{14}}}
    \proofstepwidestar[3,13]{u\in\NonemptyWords A}{\FormulaRefAuto{NonemptyWordMembershipIntroduction}[thm]{\rAI{3,15}}}
    \proofstepwidestar[3,5,13]{P(u)}{\rRE{5,16}}
    \proofstepwidestar[2,3,4,13]{P(u)\rightarrow P(\WordConcat u{\LetterWord a})}{\rRE{4,\rUE{\rRE{16,\rUE{2}}}}}
    \proofstepwidestar[2,3,4,5,13]{P(\WordConcat u{\LetterWord a})}{\rRE{17,18}}
    \proofstepwidestar[1,2,3,4,5]{P(\WordConcat u{\LetterWord a})}{\rOE{6,7,12,13,19}}
    \proofstepwidestar[21]{\WordConcat u{\LetterWord a}\in\NonemptyWords A}{\rA}
    \proofstepwidestar[1,2,3,4,5,21]{\WordConcat u{\LetterWord a}\in\NonemptyWords A\land P(\WordConcat u{\LetterWord a})}{\rAI{21,20}}
    \proofstepwidestar[1,2,3,4,5,21]{P(\WordConcat u{\LetterWord a})}{\rAEb{22}}
    \proofstepwidestar[1,2,3,4,5]{R(\WordConcat u{\LetterWord a})}{\rRI{21,23}}
    \proofstepwidestar[1,2]{\forall u\in\FiniteWords A\,\forall a\in A\,
      (R(u)\rightarrow R(\WordConcat u{\LetterWord a}))}{\rUI{\rRI{3,\rUI{\rRI{4,\rRI{5,24}}}}}}
  \closeproofpartwide
  \Needspace{20\baselineskip}
  \proofpartwide{Schluss}
    \proofstepwidestar[1]{\forall a\in A\,P(\LetterWord a)}{\rA}
    \proofstepwidestar[2]{\forall u\in\NonemptyWords A\,\forall a\in A\,
      (P(u)\rightarrow P(\WordConcat u{\LetterWord a}))}{\rA}
    \proofstepwidestar[1,2]{\forall w\in\FiniteWords A\,R(w)}{%
      \FormulaRefAuto{FiniteWordStructuralInduction}[thm]{\text{(i),(ii)}}}
    \proofstepwidestar[4]{w\in\NonemptyWords A}{\rA}
    \proofstepwidestar[4]{w\in\FiniteWords A}{\FormulaRefAuto{NonemptyWordFinite}[thm]{4}}
    \proofstepwidestar[1,2,4]{R(w)}{\rRE{5,\rUE{3}}}
    \proofstepwidestar[1,2,4]{P(w)}{\rRE{6,4}}
    \proofstepwidestar[1,2]{\forall w\in\NonemptyWords A\,P(w)}{\rUI{\rRI{4,7}}}
  \closeproofpartwide
\end{tabproofsplitwide}


\FormulaThmDeltaK[Rekursion über nichtleere Wörter]{%
  \begin{aligned}[t]
    &f\colon A\to X\dsep s\colon X\times A\to X\vdash{}\\[-2pt]
    &\exists!F\colon\NonemptyWords{A}\to X\,\Bigl(
      (\forall a\in A\,
        F(\LetterWord{a})=f(a))\ \land{}\\[-2pt]
    &\qquad(\forall u\in\NonemptyWords{A}\,\forall a\in A\,
      F(\WordConcat{u}{\LetterWord{a}})=s(F(u),a))
    \Bigr)
  \end{aligned}%
}{NonemptyWordRecursion}{%
  \DeltaRow{Mengen}{A\dsep X\dsep a\dsep u}
  \DeltaRow{Funktionen}{f\dsep s\dsep F}
}
Die Hilfswerte liegen in \(\powerset(X)\). Wir verwenden den Term
\[
 \rho(S,a)\coloneqq\IfThenElse{S=\varnothing}{\{f(a)\}}
   {\{s(x,a)\mid x\in S\}}.
\]
Die Projektionen in den ersten Schritten gehören zu
\(\powerset(X)\times A\). Die folgenden Prädikate sind Abkürzungen:
\[
\begin{aligned}
 \mathsf D(r)&\coloneqq r\colon\powerset(X)\times A\to\powerset(X)\land{}\\[-2pt]
 &\qquad(\forall p\in\powerset(X)\times A\;r(p)=\rho(\pi_1(p),\pi_2(p))),\\
 P(w)&\coloneqq\exists x\in X\;H(w)=\{x\},\\
 Q(K)&\coloneqq K\colon\NonemptyWords A\to X\land{}\\[-2pt]
 &\quad(\forall a\in A\;K(\LetterWord a)=f(a))\land{}\\[-2pt]
 &\quad(\forall u\in\NonemptyWords A\,\forall a\in A\;
 K(\WordConcat u{\LetterWord a})=s(K(u),a)).
\end{aligned}
\]
Die Funktionen \(\widehat s,H,F\) werden jeweils erst nach dem
Nachweis ihrer eindeutigen Bestimmtheit durch \(\iota\) festgelegt.
\begin{tabproofwide}
  \proofstepwidestar[1]{f\colon A\to X}{\rA}
  \proofstepwidestar[2]{s\colon X\times A\to X}{\rA}
  \proofstepwidestar[3]{p\in\powerset(X)\times A}{\rA}
  \proofstepwidestar[3]{\pi_1(p)\in\powerset(X)}{\FormulaRefAuto{F\colon A\to B\dsep x\in A\vdash F(x)\in B}[thm]{\FormulaRefAuto{FirstProjectionDef}[def],3}}
  \proofstepwidestar[3]{\pi_2(p)\in A}{\FormulaRefAuto{F\colon A\to B\dsep x\in A\vdash F(x)\in B}[thm]{\FormulaRefAuto{SecondProjectionDef}[def],3}}
  \proofstepwidestar[3]{\pi_1(p)\subseteq X}{\FormulaRefAuto{x\in\powerset(A)\eqvdash x\subseteq A}[thm]{4}}
  \proofstepwidestar[7]{x\in\pi_1(p)}{\rA}
  \proofstepwidestar[3,7]{x\in X}{\FormulaRefAuto{A\subseteq B\dsep x\in A\vdash x\in B}[thm]{6,7}}
  \proofstepwidestar[3,7]{(x,\pi_2(p))\in X\times A}{\FormulaRefAuto{(a,b)\in A\times B\eqvdash a\in A\land b\in B}[thm]{\rAI{8,5}}}
  \proofstepwidestar[2,3,7]{s(x,\pi_2(p))\in X}{\FormulaRefAuto{F\colon A\to B\dsep x\in A\vdash F(x)\in B}[thm]{2,9}}
  \proofstepwidestar[2,3]{\forall x\in\pi_1(p)\;s(x,\pi_2(p))\in X}{\rUI{\rRI{7,10}}}
  \proofstepwidestar[2,3]{\{s(x,\pi_2(p))\mid x\in\pi_1(p)\}\subseteq X}{\FormulaRefAuto{B27TermImageSubset}[thm]{11}}
  \proofstepwidestar[2,3]{\{s(x,\pi_2(p))\mid x\in\pi_1(p)\}\in\powerset(X)}{\FormulaRefAuto{x\in\powerset(A)\eqvdash x\subseteq A}[thm]{12}}
  \proofstepwidestar[1,3]{f(\pi_2(p))\in X}{\FormulaRefAuto{F\colon A\to B\dsep x\in A\vdash F(x)\in B}[thm]{1,5}}
  \proofstepwidestar[1,3]{\{f(\pi_2(p))\}\subseteq X}{\FormulaRefAuto{a\in A\vdash\{a\}\subseteq A}[thm]{14}}
  \proofstepwidestar[1,3]{\{f(\pi_2(p))\}\in\powerset(X)}{\FormulaRefAuto{x\in\powerset(A)\eqvdash x\subseteq A}[thm]{15}}
  \proofstepwidestar[]{\pi_1(p)=\varnothing\lor\pi_1(p)\ne\varnothing}{\FormulaRefAuto{P\lor\neg P}[thm]}
  \proofstepwidestar[18]{\pi_1(p)=\varnothing}{\rA}
  \proofstepwidestar[18]{\rho(\pi_1(p),\pi_2(p))=\{f(\pi_2(p))\}}{\FormulaRefAuto{P\vdash\IfThenElse{P}{a}{b}=a}[thm]{18}}
  \proofstepwidestar[1,3,18]{\rho(\pi_1(p),\pi_2(p))\in\powerset(X)}{\rIE{19,16}}
  \proofstepwidestar[21]{\pi_1(p)\ne\varnothing}{\rA}
  \proofstepwidestar[21]{\rho(\pi_1(p),\pi_2(p))=\{s(x,\pi_2(p))\mid x\in\pi_1(p)\}}{\FormulaRefAuto{\neg P\vdash\IfThenElse{P}{a}{b}=b}[thm]{21}}
  \proofstepwidestar[2,3,21]{\rho(\pi_1(p),\pi_2(p))\in\powerset(X)}{\rIE{22,13}}
  \proofstepwidestar[1,2,3]{\rho(\pi_1(p),\pi_2(p))\in\powerset(X)}{\rOE{17,18,20,21,23}}
  \proofstepwidestar[1,2]{\forall p\in\powerset(X)\times A\;\rho(\pi_1(p),\pi_2(p))\in\powerset(X)}{\rUI{\rRI{3,24}}}
  \proofstepwidestar[1,2]{\begin{aligned}[t]&\exists!r\;\Bigl(r\colon\powerset(X)\times A\to\powerset(X)\land{}\\[-2pt]&\quad(\forall p\in\powerset(X)\times A\;r(p)=\rho(\pi_1(p),\pi_2(p)))\Bigr)\end{aligned}}{\FormulaRefAuto{TypedFunctionDefinitionPrinciple}[thm]{25}}
  \proofstepwidestar[1,2]{\widehat s\coloneqq\iota r\;\mathsf D(r)}{\FormulaRefAuto{IotaDefinitionDef}[def]{26}}
  \proofstepwidestar[1,2]{\mathsf D(\widehat s)}{\rIE{27,\FormulaRefAuto{IotaSymbolDef}[def]{26}}}
  \proofstepwidestar[1,2]{\widehat s\colon\powerset(X)\times A\to\powerset(X)}{\rAEa{28}}
  \proofstepwidestar[30]{S\in\powerset(X)}{\rA}
  \proofstepwidestar[31]{a\in A}{\rA}
  \proofstepwidestar[30,31]{(S,a)\in\powerset(X)\times A}{\FormulaRefAuto{(a,b)\in A\times B\eqvdash a\in A\land b\in B}[thm]{\rAI{30,31}}}
  \proofstepwidestar[1,2,30,31]{\widehat s(S,a)=\rho(\pi_1(S,a),\pi_2(S,a))}{\rRE{32,\rUE{\rAEb{28}}}}
  \proofstepwidestar[30,31]{\pi_1(S,a)=S}{\FormulaRefAuto{(x,y)\in A\times B\vdash\pi_1(x,y)=x}[thm]{32}}
  \proofstepwidestar[30,31]{\pi_2(S,a)=a}{\FormulaRefAuto{(x,y)\in A\times B\vdash\pi_2(x,y)=y}[thm]{32}}
  \proofstepwidestar[1,2,30,31]{\widehat s(S,a)=\rho(S,a)}{\rIE{34,35,33}}
  \proofstepwidestar[1,2]{\forall S\in\powerset(X)\,\forall a\in A\;\widehat s(S,a)=\rho(S,a)}{\rUI{\rRI{30,\rUI{\rRI{31,36}}}}}
  \proofstepwidestar[]{\varnothing\in\powerset(X)}{\FormulaRefAuto{x\in\powerset(A)\eqvdash x\subseteq A}[thm]{\FormulaRefAuto{\varnothing\subseteq A}[thm]}}
  \proofstepwidestar[1,2]{\exists!H\;\mathsf{WRec}_A(H;\FiniteWords A,\EmptyWord A,\sigma_A;\powerset(X),\varnothing,\widehat s)}{\FormulaRefAuto{FiniteWordRecursion}[thm]{38,29}}
  \proofstepwidestar[1,2]{\begin{aligned}[t]H\coloneqq\iota K\;&\mathsf{WRec}_A(K;\FiniteWords A,\EmptyWord A,\sigma_A;\\[-2pt]&\powerset(X),\varnothing,\widehat s)\end{aligned}}{\FormulaRefAuto{IotaDefinitionDef}[def]{39}}
  \proofstepwidestar[1,2]{\mathsf{WRec}_A(H;\FiniteWords A,\EmptyWord A,\sigma_A;\powerset(X),\varnothing,\widehat s)}{\rIE{40,\FormulaRefAuto{IotaSymbolDef}[def]{39}}}
  \proofstepwidestar[1,2]{\begin{aligned}[t]&H\colon\FiniteWords A\to\powerset(X)\land H(\EmptyWord A)=\varnothing\land{}\\[-2pt]&\forall u\in\FiniteWords A\,\forall a\in A\;H(\sigma_A(u,a))=\widehat s(H(u),a)\end{aligned}}{\FormulaRefAuto{WordRecursionSpecificationDef}[def]{41}}
  \proofstepwidestar[43]{a\in A}{\rA}
  \proofstepwidestar[]{\EmptyWord A\in\FiniteWords A}{\FormulaRefAuto{EmptyWordFinite}[thm]}
  \proofstepwidestar[43]{\sigma_A(\EmptyWord A,a)=\LetterWord a}{\rIE{\FormulaRefAuto{WordConcatenationLeftIdentity}[thm]{\FormulaRefAuto{NonemptyWordFinite}[thm]{\FormulaRefAuto{LetterWordNonempty}[thm]{43}}},\FormulaRefAuto{WordModelSuccessorConcatenation}[thm]{44,43}}}
  \proofstepwidestar[1,2,43]{H(\sigma_A(\EmptyWord A,a))=\widehat s(H(\EmptyWord A),a)}{\rRE{43,\rUE{\rRE{44,\rUE{\rAEb{\rAEb{42}}}}}}}
  \proofstepwidestar[1,2,43]{\widehat s(\varnothing,a)=\rho(\varnothing,a)}{\rRE{43,\rUE{\rRE{38,\rUE{37}}}}}
  \proofstepwidestar[]{\rho(\varnothing,a)=\{f(a)\}}{\FormulaRefAuto{P\vdash\IfThenElse{P}{a}{b}=a}[thm]{\rII}}
  \proofstepwidestar[1,2,43]{H(\LetterWord a)=\{f(a)\}}{\rIE{45,\rAEa{\rAEb{42}},47,48,46}}
  \proofstepwidestar[1,43]{f(a)\in X}{\FormulaRefAuto{F\colon A\to B\dsep x\in A\vdash F(x)\in B}[thm]{1,43}}
  \proofstepwidestar[1,2,43]{P(\LetterWord a)}{\rEI{\rAI{50,49}}}
  \proofstepwidestar[1,2]{\forall a\in A\;P(\LetterWord a)}{\rUI{\rRI{43,51}}}
  \proofstepwidestar[1,2]{\forall a\in A\;H(\LetterWord a)=\{f(a)\}}{\rUI{\rRI{43,49}}}
  \proofstepwidestar[54]{u\in\NonemptyWords A}{\rA}
  \proofstepwidestar[55]{a\in A}{\rA}
  \proofstepwidestar[56]{x\in X}{\rA}
  \proofstepwidestar[57]{H(u)=\{x\}}{\rA}
  \proofstepwidestar[54]{u\in\FiniteWords A}{\FormulaRefAuto{NonemptyWordFinite}[thm]{54}}
  \proofstepwidestar[56,55]{(x,a)\in X\times A}{\FormulaRefAuto{(a,b)\in A\times B\eqvdash a\in A\land b\in B}[thm]{\rAI{56,55}}}
  \proofstepwidestar[2,56,55]{s(x,a)\in X}{\FormulaRefAuto{F\colon A\to B\dsep x\in A\vdash F(x)\in B}[thm]{2,59}}
  \proofstepwidestar[]{x\in\{x\}}{\FormulaRefAuto{a\in\{a\}}[thm]}
  \proofstepwidestar[]{\{x\}\ne\varnothing}{\FormulaRefAuto{A\neq\varnothing \eqvdash \exists x\,(x \in A)}[thm]{\rEI{61}}}
  \proofstepwidestar[56]{\{x\}\in\powerset(X)}{\FormulaRefAuto{x\in\powerset(A)\eqvdash x\subseteq A}[thm]{\FormulaRefAuto{a\in A\vdash\{a\}\subseteq A}[thm]{56}}}
  \proofstepwidestar[1,2,56,55]{\widehat s(\{x\},a)=\rho(\{x\},a)}{\rRE{55,\rUE{\rRE{63,\rUE{37}}}}}
  \proofstepwidestar[]{\rho(\{x\},a)=\{s(y,a)\mid y\in\{x\}\}}{\FormulaRefAuto{\neg P\vdash\IfThenElse{P}{a}{b}=b}[thm]{62}}
  \proofstepwidestar[]{\rho(\{x\},a)=\{s(x,a)\}}{\rIE{\FormulaRefAuto{B27TermImageSingleton}[thm],65}}
  \proofstepwidestar[1,2,54,55]{H(\sigma_A(u,a))=\widehat s(H(u),a)}{\rRE{55,\rUE{\rRE{58,\rUE{\rAEb{\rAEb{42}}}}}}}
  \proofstepwidestar[1,2,54,55,56,57]{H(\WordConcat u{\LetterWord a})=\{s(x,a)\}}{\rIE{57,64,66,\FormulaRefAuto{WordModelSuccessorConcatenation}[thm]{58,55},67}}
  \proofstepwidestar[1,2]{\begin{aligned}[t]&\forall u\in\NonemptyWords A\,\forall a\in A\,\forall x\in X\\[-2pt]&\quad(H(u)=\{x\}\rightarrow H(\WordConcat u{\LetterWord a})=\{s(x,a)\})\end{aligned}}{\rUI{\rRI{54,\rUI{\rRI{55,\rUI{\rRI{56,\rRI{57,68}}}}}}}}
  \proofstepwidestar[70]{P(u)}{\rA}
  \proofstepwidestar[71]{x\in X\land H(u)=\{x\}}{\rA}
  \proofstepwidestar[1,2,54,55,71]{H(\WordConcat u{\LetterWord a})=\{s(x,a)\}}{\rRE{\rAEb{71},\rRE{\rAEa{71},\rUE{\rRE{55,\rUE{\rRE{54,\rUE{69}}}}}}}}
  \proofstepwidestar[2,55,71]{s(x,a)\in X}{\FormulaRefAuto{F\colon A\to B\dsep x\in A\vdash F(x)\in B}[thm]{2,\FormulaRefAuto{(a,b)\in A\times B\eqvdash a\in A\land b\in B}[thm]{\rAI{\rAEa{71},55}}}}
  \proofstepwidestar[1,2,54,55,71]{P(\WordConcat u{\LetterWord a})}{\rEI{\rAI{73,72}}}
  \proofstepwidestar[1,2,54,55,70]{P(\WordConcat u{\LetterWord a})}{\rEE{70,71,74}}
  \proofstepwidestar[1,2]{\forall u\in\NonemptyWords A\,\forall a\in A\;(P(u)\rightarrow P(\WordConcat u{\LetterWord a}))}{\rUI{\rRI{54,\rUI{\rRI{55,\rRI{70,75}}}}}}
  \proofstepwidestar[1,2]{\forall w\in\NonemptyWords A\;P(w)}{\FormulaRefAuto{NonemptyWordInduction}[thm]{52,76}}
  \proofstepwidestar[78]{w\in\NonemptyWords A}{\rA}
  \proofstepwidestar[1,2,78]{\exists x\in X\;H(w)=\{x\}}{\rRE{78,\rUE{77}}}
  \proofstepwidestar[80]{y\in X\land H(w)=\{y\}}{\rA}
  \proofstepwidestar[80]{\bigcup H(w)=y}{\rIE{\rAEb{80},\FormulaRefAuto{B27SingletonUnion}[thm]}}
  \proofstepwidestar[80]{\bigcup H(w)\in X}{\rIE{81,\rAEa{80}}}
  \proofstepwidestar[1,2,78]{\bigcup H(w)\in X}{\rEE{79,80,82}}
  \proofstepwidestar[1,2]{\forall w\in\NonemptyWords A\;\bigcup H(w)\in X}{\rUI{\rRI{78,83}}}
  \proofstepwidestar[1,2]{\exists!F\;(F\colon\NonemptyWords A\to X\land\forall w\in\NonemptyWords A\;F(w)=\bigcup H(w))}{\FormulaRefAuto{TypedFunctionDefinitionPrinciple}[thm]{84}}
  \proofstepwidestar[1,2]{F\coloneqq\iota K\;(K\colon\NonemptyWords A\to X\land\forall w\in\NonemptyWords A\;K(w)=\bigcup H(w))}{\FormulaRefAuto{IotaDefinitionDef}[def]{85}}
  \proofstepwidestar[1,2]{F\colon\NonemptyWords A\to X\land\forall w\in\NonemptyWords A\;F(w)=\bigcup H(w)}{\rIE{86,\FormulaRefAuto{IotaSymbolDef}[def]{85}}}
  \proofstepwidestar[1,2,78]{F(w)=\bigcup H(w)}{\rRE{78,\rUE{\rAEb{87}}}}
  \proofstepwidestar[1,2,78,80]{H(w)=\{F(w)\}}{\rIE{81,88,\rAEb{80}}}
  \proofstepwidestar[1,2,78]{H(w)=\{F(w)\}}{\rEE{79,80,89}}
  \proofstepwidestar[1,2]{\forall w\in\NonemptyWords A\;H(w)=\{F(w)\}}{\rUI{\rRI{78,90}}}
  \proofstepwidestar[92]{a\in A}{\rA}
  \proofstepwidestar[92]{\LetterWord a\in\NonemptyWords A}{\FormulaRefAuto{LetterWordNonempty}[thm]{92}}
  \proofstepwidestar[1,2,92]{F(\LetterWord a)=\bigcup H(\LetterWord a)}{\rRE{93,\rUE{\rAEb{87}}}}
  \proofstepwidestar[1,2,92]{F(\LetterWord a)=f(a)}{\rIE{\rRE{92,\rUE{53}},\FormulaRefAuto{B27SingletonUnion}[thm],94}}
  \proofstepwidestar[1,2]{\forall a\in A\;F(\LetterWord a)=f(a)}{\rUI{\rRI{92,95}}}
  \proofstepwidestar[97]{u\in\NonemptyWords A}{\rA}
  \proofstepwidestar[98]{a\in A}{\rA}
  \proofstepwidestar[1,2,97]{F(u)\in X}{\FormulaRefAuto{F\colon A\to B\dsep x\in A\vdash F(x)\in B}[thm]{\rAEa{87},97}}
  \proofstepwidestar[1,2,97]{H(u)=\{F(u)\}}{\rRE{97,\rUE{91}}}
  \proofstepwidestar[1,2,97,98]{H(\WordConcat u{\LetterWord a})=\{s(F(u),a)\}}{\rRE{100,\rRE{99,\rUE{\rRE{98,\rUE{\rRE{97,\rUE{69}}}}}}}}
  \proofstepwidestar[97,98]{\WordConcat u{\LetterWord a}\in\NonemptyWords A}{\FormulaRefAuto{WordAppendNonempty}[thm]{\FormulaRefAuto{NonemptyWordFinite}[thm]{97},98}}
  \proofstepwidestar[1,2,97,98]{F(\WordConcat u{\LetterWord a})=\bigcup H(\WordConcat u{\LetterWord a})}{\rRE{102,\rUE{\rAEb{87}}}}
  \proofstepwidestar[1,2,97,98]{F(\WordConcat u{\LetterWord a})=s(F(u),a)}{\rIE{101,\FormulaRefAuto{B27SingletonUnion}[thm],103}}
  \proofstepwidestar[1,2]{\forall u\in\NonemptyWords A\,\forall a\in A\;F(\WordConcat u{\LetterWord a})=s(F(u),a)}{\rUI{\rRI{97,\rUI{\rRI{98,104}}}}}
  \proofstepwidestar[1,2]{Q(F)}{\rAI{\rAEa{87},\rAI{96,105}}}
  \proofstepwidestar[107]{Q(G)}{\rA}
  \proofstepwidestar[108]{Q(K)}{\rA}
  \proofstepwidestar[109]{a\in A}{\rA}
  \proofstepwidestar[107,109]{G(\LetterWord a)=f(a)}{\rRE{109,\rUE{\rAEa{\rAEb{107}}}}}
  \proofstepwidestar[108,109]{K(\LetterWord a)=f(a)}{\rRE{109,\rUE{\rAEa{\rAEb{108}}}}}
  \proofstepwidestar[107,108,109]{G(\LetterWord a)=K(\LetterWord a)}{\FormulaRefAuto{a=b\dsep c=b\vdash a=c}[thm]{110,111}}
  \proofstepwidestar[107,108]{\forall a\in A\;G(\LetterWord a)=K(\LetterWord a)}{\rUI{\rRI{109,112}}}
  \proofstepwidestar[114]{u\in\NonemptyWords A}{\rA}
  \proofstepwidestar[115]{a\in A}{\rA}
  \proofstepwidestar[116]{G(u)=K(u)}{\rA}
  \proofstepwidestar[107,114,115]{G(\WordConcat u{\LetterWord a})=s(G(u),a)}{\rRE{115,\rUE{\rRE{114,\rUE{\rAEb{\rAEb{107}}}}}}}
  \proofstepwidestar[108,114,115]{K(\WordConcat u{\LetterWord a})=s(K(u),a)}{\rRE{115,\rUE{\rRE{114,\rUE{\rAEb{\rAEb{108}}}}}}}
  \proofstepwidestar[107,108,114,115,116]{G(\WordConcat u{\LetterWord a})=K(\WordConcat u{\LetterWord a})}{\FormulaRefAuto{a=b\dsep c=b\vdash a=c}[thm]{117,\rIE{116,118}}}
  \proofstepwidestar[107,108]{\forall u\in\NonemptyWords A\,\forall a\in A\;(G(u)=K(u)\rightarrow G(\WordConcat u{\LetterWord a})=K(\WordConcat u{\LetterWord a}))}{\rUI{\rRI{114,\rUI{\rRI{115,\rRI{116,119}}}}}}
  \proofstepwidestar[107,108]{\forall w\in\NonemptyWords A\;G(w)=K(w)}{\FormulaRefAuto{NonemptyWordInduction}[thm]{113,120}}
  \proofstepwidestar[107,108]{G=K}{\FormulaRefAuto{FunctionExtensionality}[thm]{\rAEa{107},\rAEa{108},121}}
  \proofstepwidestar[1,2]{\exists!F\;Q(F)}{\UEI{\rEI{106},107,108,122}}
\end{tabproofwide}


\section{Linksfaltung eines Wortes}

Die Linksfaltung ist eine Anwendung der soeben bewiesenen Rekursion.
Der erste Buchstabe ist der Anfangswert; jeder weitere Buchstabe wird
mit dem bisher berechneten Wert durch \(\star\) verknüpft. So wie
eine Zahlenrekursion ihren Wert von Schritt zu Schritt fortschreibt,
liest diese Wortrekursion die Buchstaben der Reihe nach. Für drei
Buchstaben entsteht \((a\star b)\star c\). Eine Assoziativitätsannahme
ist dafür nicht erforderlich: Die Rekursionsrichtung legt bereits fest,
welche Klammerung gemeint ist.

\Needspace{15\baselineskip}
\FormulaThmDeltaK[Existenz und Eindeutigkeit der Linksfaltung]{%
  \begin{aligned}[t]
    &\mathsf{BinOp}(A,\star)\vdash{}\\[-2pt]
    &\exists!F\colon\NonemptyWords{A}\to A\,\Bigl(
      \forall a\in A\,
        F(\LetterWord{a})=a\ \land{}\\[-2pt]
    &\qquad \forall u\in\NonemptyWords{A}\,\forall a\in A\,
      F(\WordConcat{u}{\LetterWord{a}})=F(u)\star a
      \Bigr)
  \end{aligned}%
}{LeftWordFoldExistence}{%
  \DeltaRow{Mengen}{A\dsep a\dsep u}
  \DeltaRow{Funktionen}{\star\dsep F}
}
\begin{tabproofsplitwide}
  \proofpartwide{Übersetzung der Rekursionsbedingungen}
  \proofstepwidestar[1]{\mathsf{BinOp}(A,\star)}{\rA}
  \proofstepwidestar{\Id_A\colon A\to A}{%
    \FormulaRefAuto{IdA}[thm]}
  \proofstepwidestar[1]{%
    \BinaryOpFunction{A}{\star}\colon A\times A\to A}{%
    \FormulaRefAuto{BinaryOperationFunctionDef}[def]{1}}
  \proofstepwidestar{%
    \begin{aligned}[t]
      &Q(F)\coloneqq{}\\[-2pt]
      &\quad F\colon\NonemptyWords{A}\to A\land{}\\[-2pt]
      &\quad \forall a\in A\,
        F(\LetterWord{a})=\Id_A(a)\land{}\\[-2pt]
      &\quad \forall u\in\NonemptyWords{A}\,\forall a\in A\,{}\\[-2pt]
      &\qquad F(\WordConcat{u}{\LetterWord{a}})\\[-2pt]
      &\qquad\quad =\BinaryOpFunction{A}{\star}(F(u),a)
    \end{aligned}}{\rII}
  \proofstepwidestar{%
    \begin{aligned}[t]
      &R(F)\coloneqq F\colon\NonemptyWords{A}\to A\land{}\\[-2pt]
      &\forall a\in A\,F(\LetterWord{a})=a\land{}\\[-2pt]
      &\forall u\in\NonemptyWords{A}\,\forall a\in A\,\\[-2pt]
      &\quad F(\WordConcat{u}{\LetterWord{a}})=F(u)\star a
    \end{aligned}}{\rII}
  \proofstepwidestar[1]{\exists!F\,Q(F)}{%
    \FormulaRefAuto{NonemptyWordRecursion}[thm]{2,3,4}}
  \proofstepwidestar[1]{%
    \forall a\in A\,\Id_A(a)=a}{%
    \FormulaRefAuto{x\in A\vdash\Id_A(x)=x}[thm]}
  \proofstepwidestar[8]{%
    F\colon\NonemptyWords{A}\to A}{\rA}
  \proofstepwidestar[9]{u\in\NonemptyWords{A}}{\rA}
  \proofstepwidestar[10]{a\in A}{\rA}
  \proofstepwidestar[8,9]{F(u)\in A}{%
    \FormulaRefAuto{F\colon A\to B\dsep x\in A\vdash F(x)\in B}[thm]{8,9}}
  \proofstepwidestar[8,9,10]{(F(u),a)\in A\times A}{%
    \FormulaRefAuto{(a,b)\in A\times B\eqvdash
      a\in A\land b\in B}[thm]{11,10}}
  \proofstepwidestar[1,8,9,10]{%
    \BinaryOpFunction{A}{\star}(F(u),a)=F(u)\star a}{%
    \FormulaRefAuto{BinaryOperationFunctionDef}[def]{1,11,10,12}}
  \proofstepwidestar[1]{%
    \begin{aligned}[t]
      &\forall F\colon\NonemptyWords{A}\to A\,
        \forall u\in\NonemptyWords{A}\,\forall a\in A\,{}\\[-2pt]
      &\qquad \BinaryOpFunction{A}{\star}(F(u),a)=F(u)\star a
    \end{aligned}}{%
    \rUI{\rRI{8,\rUI{\rRI{9,\rUI{\rRI{10,13}}}}}}}
  \proofstepwidestar[1]{%
    \forall F\,\bigl(Q(F)\leftrightarrow R(F)\bigr)}{%
    \FormulaRefAuto{a=b\dsep c=d\dsep P(b,d)\vdash P(a,c)}[thm]{%
      4,5,7,14}}
  \closeproofpartwide

  \proofpartwide{Existenz der Linksfaltung}
  \setcounter{proofstepnr}{15}
  \proofstepwidestar[1]{\exists F\,Q(F)}{%
    \FormulaRefAuto{\exists! x\,P(x)\vdash \exists x\,P(x)}[thm]{6}}
  \proofstepwidestar[17]{Q(F)}{\rA}
  \proofstepwidestar[1,17]{Q(F)\rightarrow R(F)}{\rUE{15}}
  \proofstepwidestar[1,17]{R(F)}{\rRE{18,17}}
  \proofstepwidestar[1,17]{\exists F\,R(F)}{\rEI{19}}
  \proofstepwidestar[1]{\exists F\,R(F)}{\rEE{16,17,20}}
  \closeproofpartwide

  \proofpartwide{Eindeutigkeit der Linksfaltung}
  \setcounter{proofstepnr}{21}
  \proofstepwidestar[1]{\existsleone F\,Q(F)}{%
    \FormulaRefAuto{\exists! x\,P(x)\vdash \existsleone x\,P(x)}[thm]{6}}
  \proofstepwidestar[1]{\existsleone F\,R(F)}{%
    \FormulaRefAuto{\existsleone x\,P(x),\;
      \forall x(P(x)\leftrightarrow Q(x))
      \vdash\existsleone x\,Q(x)}[thm]{22,15}}
  \proofstepwidestar[1]{\exists!F\,R(F)}{%
    \FormulaRefAuto{\exists x\,P(x),\; \existsleone x\,P(x)
      \vdash \exists! x\,P(x)}[thm]{21,23}}
  \proofstepwidestar[1]{%
    \exists!F\colon\NonemptyWords{A}\to A\,\Bigl(
      \forall a\in A\,F(\LetterWord{a})=a\land
      \forall u\in\NonemptyWords{A}\,\forall a\in A\,
      F(\WordConcat{u}{\LetterWord{a}})=F(u)\star a\Bigr)}{%
    \rIE{5,24}}
  \closeproofpartwide
\end{tabproofsplitwide}

\FormulaDefDeltaK[Linksfaltung eines nichtleeren Wortes]{%
  \begin{aligned}[t]
    &\mathsf{BinOp}(A,\star)\vdash
      \Pi_{\star}\coloneqq
      \iota F\colon\NonemptyWords{A}\to A\,\Bigl(\\[-2pt]
    &\quad \forall a\in A\,
      F(\LetterWord{a})=a\ \land{}\\[-2pt]
    &\quad \forall u\in\NonemptyWords{A}\,\forall a\in A\,
      F(\WordConcat{u}{\LetterWord{a}})=F(u)\star a
      \Bigr)
  \end{aligned}%
}{LeftWordFoldDef}{%
  \DeltaRow{Mengen}{A\dsep a\dsep u}
  \DeltaRow{Funktionen}{\star\dsep\Pi_{\star}}
  \DeltaRow{\textbf{Neues Symbol}}{}
  \DeltaPrem{Linksfaltung}{\Pi_{\star}}
}

\FormulaThmDeltaK[Kennzeichnung der Linksfaltung]{%
  \mathsf{BinOp}(A,\star)\vdash
  \begin{aligned}[t]
    &\text{(i)}\;\Pi_{\star}\colon\NonemptyWords{A}\to A\\[-2pt]
    &\text{(ii)}\;\forall a\in A\,
      \Pi_{\star}(\LetterWord{a})=a\\[-2pt]
    &\text{(iii)}\;\forall u\in\NonemptyWords{A}\,\forall a\in A\,
      \Pi_{\star}(\WordConcat{u}{\LetterWord{a}})
        =\Pi_{\star}(u)\star a
  \end{aligned}%
}{LeftWordFoldCharacterization}{%
  \DeltaRow{Mengen}{A\dsep a\dsep u}
  \DeltaRow{Funktionen}{\star\dsep F\dsep\Pi_{\star}}
}
\providecommand{\BXXVIIProofLeftFoldCharacterization}{%
  \proofstepwidestar[1]{\mathsf{BinOp}(A,\star)}{\rA}
  \proofstepwidestar[1]{%
    \exists!F\colon\NonemptyWords{A}\to A\,\Bigl(
      \forall a\in A\,F(\LetterWord{a})=a\land
      \forall u\in\NonemptyWords{A}\,\forall a\in A\,
      F(\WordConcat{u}{\LetterWord{a}})=F(u)\star a\Bigr)}{%
    \FormulaRefAuto{LeftWordFoldExistence}[thm]{1}}
    \proofstepwidestar[1]{%
    \Pi_\star\colon\NonemptyWords{A}\to A\land
    \forall a\in A\,\Pi_\star(\LetterWord{a})=a\land
    \forall u\in\NonemptyWords{A}\,\forall a\in A\,
    \Pi_\star(\WordConcat{u}{\LetterWord{a}})
      =\Pi_\star(u)\star a}{%
    \FormulaRefAuto{LeftWordFoldDef}[def]{1}}}
\begin{tabproofsplitwide}
  \proofpartwideindR[Typisierung]{%
    \mathsf{BinOp}(A,\star)\vdash
      \Pi_{\star}\colon\NonemptyWords{A}\to A
  }{\mathsf{BinOp}(A,\star)\vdash
      \Pi_{\star}\colon\NonemptyWords{A}\to A}%
    [LeftWordFoldCharacterizationTyping]
    \BXXVIIProofLeftFoldCharacterization
    \proofstepwidestar[1]{\Pi_\star\colon\NonemptyWords{A}\to A}{%
      \rAEa{3}}
  \closeproofpartwideindR

  \proofpartwideindR[Buchstabenwörter]{%
    \mathsf{BinOp}(A,\star)\vdash
      \forall a\in A\,\Pi_{\star}(\LetterWord{a})=a
  }{\mathsf{BinOp}(A,\star)\vdash
      \forall a\in A\,\Pi_{\star}(\LetterWord{a})=a}%
    [LeftWordFoldCharacterizationLetter]
    \BXXVIIProofLeftFoldCharacterization
    \proofstepwidestar[1]{%
      \forall a\in A\,\Pi_{\star}(\LetterWord{a})=a}{%
      \rAEa{\rAEb{3}}}
  \closeproofpartwideindR

  \proofpartwideindR[Rekursionsschritt]{%
    \begin{aligned}[t]
      &\mathsf{BinOp}(A,\star)\vdash{}\\[-2pt]
      &\forall u\in\NonemptyWords{A}\,\forall a\in A\,
        \Pi_{\star}(\WordConcat{u}{\LetterWord{a}})
          =\Pi_{\star}(u)\star a
    \end{aligned}
  }{\mathsf{BinOp}(A,\star)\vdash
      \forall u\in\NonemptyWords{A}\,\forall a\in A\,
      \Pi_{\star}(\WordConcat{u}{\LetterWord{a}})
        =\Pi_{\star}(u)\star a}%
    [LeftWordFoldCharacterizationStep]
    \BXXVIIProofLeftFoldCharacterization
    \proofstepwidestar[1]{%
      \forall u\in\NonemptyWords{A}\,\forall a\in A\,
      \Pi_{\star}(\WordConcat{u}{\LetterWord{a}})
        =\Pi_{\star}(u)\star a}{%
      \rAEb{\rAEb{3}}}
  \closeproofpartwideindR
\end{tabproofsplitwide}

\FormulaThmDeltaK[Typisierung der Linksfaltung]{%
  \mathsf{BinOp}(A,\star)\vdash
  \Pi_{\star}\colon\NonemptyWords{A}\to A%
}{LeftWordFoldFunction}{%
  \DeltaRow{Mengen}{A}
  \DeltaRow{Funktionen}{\star\dsep\Pi_{\star}}
}
\begin{tabproofwide}
  \proofstepwidestar[1]{\mathsf{BinOp}(A,\star)}{\rA}
  \proofstepwidestar[1]{\Pi_\star\colon\NonemptyWords{A}\to A}{%
    \FormulaRefAuto{LeftWordFoldCharacterizationTyping}[thm]{1}}
\end{tabproofwide}

\FormulaThmDeltaK[Rekursionsgleichungen der Linksfaltung]{%
  \begin{aligned}[t]
    &\text{(i)}\;\mathsf{BinOp}(A,\star)\dsep a\in A\vdash
      \WordFold{\star}{\LetterWord{a}}=a\\[-2pt]
    &\text{(ii)}\;\mathsf{BinOp}(A,\star)\dsep
      u\in\NonemptyWords{A}\dsep a\in A\vdash
      \WordFold{\star}{\WordConcat{u}{\LetterWord{a}}}
      =\WordFold{\star}{u}\star a
  \end{aligned}%
}{LeftWordFoldEquations}{%
  \DeltaRow{Mengen}{A\dsep a\dsep u}
  \DeltaRow{Funktionen}{\star\dsep\Pi_{\star}}
}
\begin{tabproofsplitwide}
  \proofpartwideindR[Buchstabenwort]{%
    \mathsf{BinOp}(A,\star)\dsep a\in A\vdash
      \WordFold{\star}{\LetterWord{a}}=a
  }{\mathsf{BinOp}(A,\star)\dsep a\in A\vdash
      \WordFold{\star}{\LetterWord{a}}=a}[LeftWordFoldLetter]
    \proofstepwidestar[1]{\mathsf{BinOp}(A,\star)}{\rA}
    \proofstepwidestar[2]{a\in A}{\rA}
    \proofstepwidestar[1]{%
      \forall a\in A\,\Pi_\star(\LetterWord{a})=a}{%
      \FormulaRefAuto{LeftWordFoldCharacterizationLetter}[thm]{1}}
    \proofstepwidestar[1,2]{\WordFold{\star}{\LetterWord{a}}=a}{%
      \rRE{\rUE{3},2}}
  \closeproofpartwideindR

  \proofpartwideindR[Rekursionsschritt]{%
    \begin{aligned}[t]
      &\mathsf{BinOp}(A,\star)\dsep
        u\in\NonemptyWords{A}\dsep a\in A\vdash{}\\[-2pt]
      &\WordFold{\star}{\WordConcat{u}{\LetterWord{a}}}
        =\WordFold{\star}{u}\star a
    \end{aligned}
  }{\mathsf{BinOp}(A,\star)\dsep
      u\in\NonemptyWords{A}\dsep a\in A\vdash
      \WordFold{\star}{\WordConcat{u}{\LetterWord{a}}}
      =\WordFold{\star}{u}\star a}[LeftWordFoldStep]
    \proofstepwidestar[1]{\mathsf{BinOp}(A,\star)}{\rA}
    \proofstepwidestar[2]{u\in\NonemptyWords{A}}{\rA}
    \proofstepwidestar[3]{a\in A}{\rA}
    \proofstepwidestar[1]{%
      \forall u\in\NonemptyWords{A}\,\forall a\in A\,
      \Pi_\star(\WordConcat{u}{\LetterWord{a}})
      =\Pi_\star(u)\star a}{%
      \FormulaRefAuto{LeftWordFoldCharacterizationStep}[thm]{1}}
    \proofstepwidestar[1,2,3]{%
      \WordFold{\star}{\WordConcat{u}{\LetterWord{a}}}
      =\WordFold{\star}{u}\star a}{%
      \rRE{\rUE{\rRE{\rUE{4},2}},3}}
  \closeproofpartwideindR
\end{tabproofsplitwide}

\chapter{Aufbau voller planarer Klammerungsbäume}

Ein Klammerungsbaum ist entweder ein beschriftetes Blatt oder entsteht
aus zwei bereits gebildeten Bäumen durch Einfügen einer neuen Wurzel.
Die Reihenfolge der beiden Teilbäume gehört zum Gegenstand. Schreiben wir
\(L_a=\LeafTree A a\), so entsprechen
\[
 \NodeTree A{\NodeTree A{L_a}{L_b}}{L_c}
 \quad\text{und}\quad
 \NodeTree A{L_a}{\NodeTree A{L_b}{L_c}}
\]
den Klammerungen \(((ab)c)\) und \((a(bc))\). Beide besitzen das
Blattwort \(\langle a,b,c\rangle\), aber verschiedene Baumformen.
Die Zeichen \(\operatorname{Blatt}\) und \(\operatorname{Knoten}\)
erhalten im folgenden Abschnitt eine konkrete mengentheoretische Bedeutung.

% Vektorzeichnung: zwei Baumformen zum selben Blattwort.
\begin{center}
\setlength{\unitlength}{1pt}
\begin{picture}(270,102)
  \qbezier(50,88)(37.5,72)(25,56)
  \qbezier(50,88)(67.5,56)(85,24)
  \qbezier(25,56)(17.5,40)(10,24)
  \qbezier(25,56)(35,40)(45,24)
  \put(50,88){\circle*{3}}
  \put(25,56){\circle*{3}}
  \put(10,24){\circle*{3}}
  \put(45,24){\circle*{3}}
  \put(85,24){\circle*{3}}
  \put(10,12){\makebox(0,0){\(a\)}}
  \put(45,12){\makebox(0,0){\(b\)}}
  \put(85,12){\makebox(0,0){\(c\)}}
  \put(47,-4){\makebox(0,0){\(((ab)c)\)}}
  \qbezier(215,88)(197.5,56)(180,24)
  \qbezier(215,88)(227.5,72)(240,56)
  \qbezier(240,56)(230,40)(220,24)
  \qbezier(240,56)(247.5,40)(255,24)
  \put(215,88){\circle*{3}}
  \put(240,56){\circle*{3}}
  \put(180,24){\circle*{3}}
  \put(220,24){\circle*{3}}
  \put(255,24){\circle*{3}}
  \put(180,12){\makebox(0,0){\(a\)}}
  \put(220,12){\makebox(0,0){\(b\)}}
  \put(255,12){\makebox(0,0){\(c\)}}
  \put(218,-4){\makebox(0,0){\((a(bc))\)}}
\end{picture}
\par\medskip
\small Dasselbe Blattwort, zwei verschiedene Klammerungsbäume.
\end{center}


Die Konstruktion beantwortet drei Fragen: Die Baumstufen sichern die
Existenz der Menge aller so erzeugten Gegenstände; getrennte, injektive
Konstruktoren sichern die eindeutige Zerlegung; Strukturinduktion und
Strukturrekursion rechtfertigen Beweise und Definitionen nach diesem
Aufbau. Anschließend werden die Knoten als endliche Positionsmenge
organisiert und mit der Baumtheorie aus Band~26 verbunden.

\section{Ein Mengenmodell der Blatt- und Knotenregel}

Die Baumcodierung verwendet zwei Sorten von Zeichen. Ein Blattzeichen trägt
ein Element von \(A\); ein Knotenzeichen speichert zusätzlich die Länge des
linken Teilbaums. Diese Längenangabe macht die Zerlegung eines Knotencodes
ohne syntaktische Metasprache eindeutig.

\FormulaDefDeltaK[Alphabet der Baumcodes]{%
  \TreeAlphabet{A}\coloneqq
  (\{0\}\times A)\cup(\{1\}\times\mathbb N)%
}{TreeTokenAlphabetDef}{%
  \DeltaRow{Mengen}{A\dsep\TreeAlphabet{A}}
  \DeltaRow{\textbf{Neues Symbol}}{}
  \DeltaPrem{Baumalphabet}{\TreeAlphabet{A}}
}

\FormulaDefDeltaK[Blattcode]{%
  a\in A\vdash
  \LeafTree{A}{a}\coloneqq\LetterWord{(0,a)}%
}{TreeLeafCodeDef}{%
  \DeltaRow{Mengen}{A\dsep a}
  \DeltaRow{\textbf{Neues Symbol}}{}
  \DeltaPrem{Blattcode}{\LeafTree{A}{a}}
}

\FormulaThmDeltaK[Wohldefiniertheit der Knotencodedaten]{%
  S,T\in\NonemptyWords{\TreeAlphabet{A}}\vdash
  \begin{aligned}[t]
    &\text{(i)}\;S\in\FiniteWords{\TreeAlphabet{A}}\\[-2pt]
    &\text{(ii)}\;T\in\FiniteWords{\TreeAlphabet{A}}\\[-2pt]
    &\text{(iii)}\;\WordLength{S}\in\mathbb N\\[-2pt]
    &\text{(iv)}\;\WordLength{T}\in\mathbb N\\[-2pt]
    &\text{(v)}\;(1,\WordLength{S})\in\TreeAlphabet{A}\\[-2pt]
    &\text{(vi)}\;\LetterWord{(1,\WordLength{S})}
      \in\NonemptyWords{\TreeAlphabet{A}}\\[-2pt]
    &\text{(vii)}\;\WordConcat{S}{T}
      \in\NonemptyWords{\TreeAlphabet{A}}\\[-2pt]
    &\text{(viii)}\;\LetterWord{(1,\WordLength{S})}
      \in\FiniteWords{\TreeAlphabet{A}}\\[-2pt]
    &\text{(ix)}\;\WordConcat{S}{T}
      \in\FiniteWords{\TreeAlphabet{A}}
  \end{aligned}%
}{TreeNodeCodeTyping}{%
  \DeltaRow{Mengen}{A\dsep S\dsep T}
}
\begin{tabproofsplitwide}
  \proofpartwideindR[Endlichkeit des linken Teilcodes]{%
    S,T\in\NonemptyWords{\TreeAlphabet{A}}\vdash
      S\in\FiniteWords{\TreeAlphabet{A}}
  }{S,T\in\NonemptyWords{\TreeAlphabet{A}}\vdash
      S\in\FiniteWords{\TreeAlphabet{A}}}%
    [TreeNodeCodeLeftFinite]
    \proofstepwidestar[1]{%
      S,T\in\NonemptyWords{\TreeAlphabet{A}}}{\rA}
    \proofstepwidestar[1]{S\in\FiniteWords{\TreeAlphabet{A}}}{%
      \FormulaRefAuto{NonemptyWordFinite}[thm]{1}}
  \closeproofpartwideindR

  \proofpartwideindR[Endlichkeit des rechten Teilcodes]{%
    S,T\in\NonemptyWords{\TreeAlphabet{A}}\vdash
      T\in\FiniteWords{\TreeAlphabet{A}}
  }{S,T\in\NonemptyWords{\TreeAlphabet{A}}\vdash
      T\in\FiniteWords{\TreeAlphabet{A}}}%
    [TreeNodeCodeRightFinite]
    \proofstepwidestar[1]{%
      S,T\in\NonemptyWords{\TreeAlphabet{A}}}{\rA}
    \proofstepwidestar[1]{T\in\FiniteWords{\TreeAlphabet{A}}}{%
      \FormulaRefAuto{NonemptyWordFinite}[thm]{1}}
  \closeproofpartwideindR

  \proofpartwideindR[Natürlichkeit der linken Teillänge]{%
    S,T\in\NonemptyWords{\TreeAlphabet{A}}\vdash
      \WordLength{S}\in\mathbb N
  }{S,T\in\NonemptyWords{\TreeAlphabet{A}}\vdash
      \WordLength{S}\in\mathbb N}%
    [TreeNodeCodeLeftLengthNatural]
    \proofstepwidestar[1]{%
      S,T\in\NonemptyWords{\TreeAlphabet{A}}}{\rA}
    \proofstepwidestar[1]{\WordLength{S}\in\mathbb N}{%
      \FormulaRefAuto{NonemptyWordLengthNatural}[thm]{1}}
  \closeproofpartwideindR

  \proofpartwideindR[Natürlichkeit der rechten Teillänge]{%
    S,T\in\NonemptyWords{\TreeAlphabet{A}}\vdash
      \WordLength{T}\in\mathbb N
  }{S,T\in\NonemptyWords{\TreeAlphabet{A}}\vdash
      \WordLength{T}\in\mathbb N}%
    [TreeNodeCodeRightLengthNatural]
    \proofstepwidestar[1]{%
      S,T\in\NonemptyWords{\TreeAlphabet{A}}}{\rA}
    \proofstepwidestar[1]{\WordLength{T}\in\mathbb N}{%
      \FormulaRefAuto{NonemptyWordLengthNatural}[thm]{1}}
  \closeproofpartwideindR

  \proofpartwideindR[Typisierung des Knotenzeichens]{%
    S,T\in\NonemptyWords{\TreeAlphabet{A}}\vdash
      (1,\WordLength{S})\in\TreeAlphabet{A}
  }{S,T\in\NonemptyWords{\TreeAlphabet{A}}\vdash
      (1,\WordLength{S})\in\TreeAlphabet{A}}%
    [TreeNodeCodeTagTyping]
    \proofstepwidestar[1]{%
      S,T\in\NonemptyWords{\TreeAlphabet{A}}}{\rA}
    \proofstepwidestar[1]{\WordLength{S}\in\mathbb N}{%
      \FormulaRefAuto{NonemptyWordLengthNatural}[thm]{1}}
    \proofstepwidestar[1]{%
      (1,\WordLength{S})\in\TreeAlphabet{A}}{%
      \FormulaRefAuto{TreeTokenAlphabetDef}[def]{2}}
  \closeproofpartwideindR

  \proofpartwideindR[Nichtleerheit des Knotenzeichenwortes]{%
    S,T\in\NonemptyWords{\TreeAlphabet{A}}\vdash
      \LetterWord{(1,\WordLength{S})}
        \in\NonemptyWords{\TreeAlphabet{A}}
  }{S,T\in\NonemptyWords{\TreeAlphabet{A}}\vdash
      \LetterWord{(1,\WordLength{S})}
        \in\NonemptyWords{\TreeAlphabet{A}}}%
    [TreeNodeCodeTagNonempty]
    \proofstepwidestar[1]{%
      S,T\in\NonemptyWords{\TreeAlphabet{A}}}{\rA}
    \proofstepwidestar[1]{%
      (1,\WordLength{S})\in\TreeAlphabet{A}}{%
      \FormulaRefAuto{TreeNodeCodeTagTyping}[thm]{1}}
    \proofstepwidestar[1]{%
      \LetterWord{(1,\WordLength{S})}
        \in\NonemptyWords{\TreeAlphabet{A}}}{%
      \FormulaRefAuto{LetterWordNonempty}[thm]{2}}
  \closeproofpartwideindR

  \proofpartwideindR[Nichtleerheit des Nutzwortes]{%
    S,T\in\NonemptyWords{\TreeAlphabet{A}}\vdash
      \WordConcat{S}{T}\in\NonemptyWords{\TreeAlphabet{A}}
  }{S,T\in\NonemptyWords{\TreeAlphabet{A}}\vdash
      \WordConcat{S}{T}\in\NonemptyWords{\TreeAlphabet{A}}}%
    [TreeNodeCodePayloadNonempty]
    \proofstepwidestar[1]{%
      S,T\in\NonemptyWords{\TreeAlphabet{A}}}{\rA}
    \proofstepwidestar[1]{%
      \WordConcat{S}{T}\in\NonemptyWords{\TreeAlphabet{A}}}{%
      \FormulaRefAuto{NonemptyWordConcatenationClosure}[thm]{1}}
  \closeproofpartwideindR

  \proofpartwideindR[Endlichkeit des Knotenzeichenwortes]{%
    S,T\in\NonemptyWords{\TreeAlphabet{A}}\vdash
      \LetterWord{(1,\WordLength{S})}
        \in\FiniteWords{\TreeAlphabet{A}}
  }{S,T\in\NonemptyWords{\TreeAlphabet{A}}\vdash
      \LetterWord{(1,\WordLength{S})}
        \in\FiniteWords{\TreeAlphabet{A}}}%
    [TreeNodeCodeTagFinite]
    \proofstepwidestar[1]{%
      S,T\in\NonemptyWords{\TreeAlphabet{A}}}{\rA}
    \proofstepwidestar[1]{%
      \LetterWord{(1,\WordLength{S})}
        \in\NonemptyWords{\TreeAlphabet{A}}}{%
      \FormulaRefAuto{TreeNodeCodeTagNonempty}[thm]{1}}
    \proofstepwidestar[1]{%
      \LetterWord{(1,\WordLength{S})}
        \in\FiniteWords{\TreeAlphabet{A}}}{%
      \FormulaRefAuto{NonemptyWordFinite}[thm]{2}}
  \closeproofpartwideindR

  \proofpartwideindR[Endlichkeit des Nutzwortes]{%
    S,T\in\NonemptyWords{\TreeAlphabet{A}}\vdash
      \WordConcat{S}{T}\in\FiniteWords{\TreeAlphabet{A}}
  }{S,T\in\NonemptyWords{\TreeAlphabet{A}}\vdash
      \WordConcat{S}{T}\in\FiniteWords{\TreeAlphabet{A}}}%
    [TreeNodeCodePayloadFinite]
    \proofstepwidestar[1]{%
      S,T\in\NonemptyWords{\TreeAlphabet{A}}}{\rA}
    \proofstepwidestar[1]{%
      \WordConcat{S}{T}\in\NonemptyWords{\TreeAlphabet{A}}}{%
      \FormulaRefAuto{TreeNodeCodePayloadNonempty}[thm]{1}}
    \proofstepwidestar[1]{%
      \WordConcat{S}{T}\in\FiniteWords{\TreeAlphabet{A}}}{%
      \FormulaRefAuto{NonemptyWordFinite}[thm]{2}}
  \closeproofpartwideindR
\end{tabproofsplitwide}

Nach \FormulaRefAuto{TreeNodeCodeTagFinite}[thm] und
\FormulaRefAuto{TreeNodeCodePayloadFinite}[thm] sind die beiden Argumente der
äußeren Wortkonkatenation endliche Wörter. Damit ist die folgende
Definition unter ihren angegebenen Voraussetzungen wohldefiniert.

\FormulaDefDeltaK[Knotencode]{%
  S,T\in\NonemptyWords{\TreeAlphabet{A}}\vdash
  \NodeTree{A}{S}{T}\coloneqq
  \WordConcat{\LetterWord{(1,\WordLength{S})}}
    {\WordConcat{S}{T}}%
}{TreeNodeCodeDef}{%
  \DeltaRow{Mengen}{A\dsep S\dsep T}
  \DeltaRow{\textbf{Neues Symbol}}{}
  \DeltaPrem{Knotencode}{\NodeTree{A}{S}{T}}
}

\FormulaThmDeltaK[Konstruktortrennung und -injektivität]{%
  \begin{aligned}[t]
    &\text{(i)}\;a,b\in A\vdash
      \bigl(\LeafTree{A}{a}=\LeafTree{A}{b}\leftrightarrow a=b\bigr)\\[-2pt]
    &\text{(ii)}\;a\in A\dsep
      S,T\in\NonemptyWords{\TreeAlphabet{A}}\vdash
      \LeafTree{A}{a}\neq\NodeTree{A}{S}{T}\\[-2pt]
    &\text{(iii)}\;S,T,S',T'\in
      \NonemptyWords{\TreeAlphabet{A}}\dsep
      \NodeTree{A}{S}{T}=\NodeTree{A}{S'}{T'}\\[-2pt]
    &\qquad\vdash S=S'\\[-2pt]
    &\text{(iv)}\;S,T,S',T'\in
      \NonemptyWords{\TreeAlphabet{A}}\dsep
      \NodeTree{A}{S}{T}=\NodeTree{A}{S'}{T'}\\[-2pt]
    &\qquad\vdash T=T'
  \end{aligned}%
}{TreeConstructorNoConfusion}{%
  \DeltaRow{Mengen}{A\dsep a\dsep b\dsep S\dsep T\dsep S'\dsep T'}
}
\newcommand{\BXXVIIProofTreeNodeConstructorCancellation}{%
  \proofstepwidestar[1]{%
    S,T,S',T'\in\NonemptyWords{\TreeAlphabet{A}}}{\rA}
  \proofstepwidestar[2]{%
    \NodeTree{A}{S}{T}=\NodeTree{A}{S'}{T'}}{\rA}
  \proofstepwidestar[1]{%
    (1,\WordLength{S})\in\TreeAlphabet{A}}{%
    \FormulaRefAuto{TreeNodeCodeTagTyping}[thm]{1}}
  \proofstepwidestar[1]{%
    (1,\WordLength{S'})\in\TreeAlphabet{A}}{%
    \FormulaRefAuto{TreeNodeCodeTagTyping}[thm]{1}}
  \proofstepwidestar[1]{%
    \LetterWord{(1,\WordLength{S})}
      \in\FiniteWords{\TreeAlphabet{A}}}{%
    \FormulaRefAuto{TreeNodeCodeTagFinite}[thm]{1}}
  \proofstepwidestar[1]{%
    \WordConcat{S}{T}\in\FiniteWords{\TreeAlphabet{A}}}{%
    \FormulaRefAuto{TreeNodeCodePayloadFinite}[thm]{1}}
  \proofstepwidestar[1]{%
    \LetterWord{(1,\WordLength{S'})}
      \in\FiniteWords{\TreeAlphabet{A}}}{%
    \FormulaRefAuto{TreeNodeCodeTagFinite}[thm]{1}}
  \proofstepwidestar[1]{%
    \WordConcat{S'}{T'}\in\FiniteWords{\TreeAlphabet{A}}}{%
    \FormulaRefAuto{TreeNodeCodePayloadFinite}[thm]{1}}
  \proofstepwidestar[1]{%
    \WordLength{\LetterWord{(1,\WordLength{S})}}=1}{%
    \FormulaRefAuto{LetterWordLengthOne}[thm]{3}}
  \proofstepwidestar[1]{%
    \WordLength{\LetterWord{(1,\WordLength{S'})}}=1}{%
    \FormulaRefAuto{LetterWordLengthOne}[thm]{4}}
  \proofstepwidestar[1]{%
    \WordLength{\LetterWord{(1,\WordLength{S})}}
      =\WordLength{\LetterWord{(1,\WordLength{S'})}}}{%
    \FormulaRefAuto{a=b\dsep c=b\vdash a=c}[thm]{9,10}}
  \proofstepwidestar[1,2]{%
    \begin{aligned}[t]
      &\WordConcat{\LetterWord{(1,\WordLength{S})}}
        {\WordConcat{S}{T}}\\[-2pt]
      &\quad=\WordConcat{\LetterWord{(1,\WordLength{S'})}}
        {\WordConcat{S'}{T'}}
    \end{aligned}}{%
    \FormulaRefAuto{TreeNodeCodeDef}[def]{2}}
  \proofstepwidestar[1,2]{%
    \LetterWord{(1,\WordLength{S})}
      =\LetterWord{(1,\WordLength{S'})}}{%
    \FormulaRefAuto{WordConcatenationFixedCutLeft}[thm]{%
      5,6,7,8,11,12}}
  \proofstepwidestar[1,2]{%
    \WordConcat{S}{T}=\WordConcat{S'}{T'}}{%
    \FormulaRefAuto{WordConcatenationFixedCutRight}[thm]{%
      5,6,7,8,11,12}}
  \proofstepwidestar[1,2]{%
    (1,\WordLength{S})=(1,\WordLength{S'})}{%
    \FormulaRefAuto{LetterWordInjectivity}[thm]{3,4,13}}
  \proofstepwidestar[1,2]{%
    1=1\land\WordLength{S}=\WordLength{S'}}{%
    \FormulaRefAuto{(a,b)=(c,d)\eqvdash
      a=c\land b=d}[thm]{15}}
  \proofstepwidestar[1,2]{%
    \WordLength{S}=\WordLength{S'}}{\rAEb{16}}
  \proofstepwidestar[1]{S\in\FiniteWords{\TreeAlphabet{A}}}{%
    \FormulaRefAuto{NonemptyWordFinite}[thm]{1}}
  \proofstepwidestar[1]{T\in\FiniteWords{\TreeAlphabet{A}}}{%
    \FormulaRefAuto{NonemptyWordFinite}[thm]{1}}
  \proofstepwidestar[1]{S'\in\FiniteWords{\TreeAlphabet{A}}}{%
    \FormulaRefAuto{NonemptyWordFinite}[thm]{1}}
  \proofstepwidestar[1]{T'\in\FiniteWords{\TreeAlphabet{A}}}{%
    \FormulaRefAuto{NonemptyWordFinite}[thm]{1}}%
}
\begin{tabproofsplitwide}
  \proofpartwideindR[Blattcodes]{%
    a,b\in A\vdash
      \bigl(\LeafTree{A}{a}=\LeafTree{A}{b}\leftrightarrow a=b\bigr)
  }{%
    a,b\in A\vdash
      \bigl(\LeafTree{A}{a}=\LeafTree{A}{b}\leftrightarrow a=b\bigr)
  }[TreeLeafConstructorInjective]
  \proofstepwidestar[1]{a\in A}{\rA}
  \proofstepwidestar[2]{b\in A}{\rA}
  \proofstepwidestar[3]{%
      \LeafTree{A}{a}=\LeafTree{A}{b}}{\rA}
  \proofstepwidestar[1]{(0,a)\in\TreeAlphabet{A}}{%
      \FormulaRefAuto{TreeTokenAlphabetDef}[def]{1}}
  \proofstepwidestar[2]{(0,b)\in\TreeAlphabet{A}}{%
      \FormulaRefAuto{TreeTokenAlphabetDef}[def]{2}}
  \proofstepwidestar[1]{%
      \LeafTree{A}{a}=\LetterWord{(0,a)}}{%
      \FormulaRefAuto{TreeLeafCodeDef}[def]{1}}
  \proofstepwidestar[2]{%
      \LeafTree{A}{b}=\LetterWord{(0,b)}}{%
      \FormulaRefAuto{TreeLeafCodeDef}[def]{2}}
  \proofstepwidestar[1,2,3]{(0,a)=(0,b)}{\FormulaRefAuto{LetterWordInjectivity}[thm]{4,5,\rIE{6,7,3}}}
  \proofstepwidestar[1,2,3]{0=0\land a=b}{%
      \FormulaRefAuto{(a,b)=(c,d)\eqvdash
        a=c\land b=d}[thm]{8}}
  \proofstepwidestar[1,2,3]{a=b}{\rAEb{9}}
  \proofstepwidestar[1,2]{%
      \LeafTree{A}{a}=\LeafTree{A}{b}\rightarrow a=b}{\rRI{3,10}}
  \proofstepwidestar[12]{a=b}{\rA}
  \proofstepwidestar[1,2,12]{%
      \LeafTree{A}{a}=\LeafTree{A}{b}}{\rIE{12}}
  \proofstepwidestar[1,2]{%
      a=b\rightarrow\LeafTree{A}{a}=\LeafTree{A}{b}}{\rRI{12,13}}
  \proofstepwidestar[1,2]{%
      \LeafTree{A}{a}=\LeafTree{A}{b}\leftrightarrow a=b}{%
      \rLRI{11,14}}
\closeproofpartwideindR

  \proofpartwideindR[Trennung der Konstruktoren]{%
    a\in A\dsep S,T\in\NonemptyWords{\TreeAlphabet{A}}
      \vdash\LeafTree{A}{a}\neq\NodeTree{A}{S}{T}
  }{%
    a\in A\dsep S,T\in\NonemptyWords{\TreeAlphabet{A}}
      \vdash\LeafTree{A}{a}\neq\NodeTree{A}{S}{T}
  }[TreeConstructorsDisjoint]
    \proofstepwidestar[1]{a\in A}{\rA}
    \proofstepwidestar[2]{%
      S,T\in\NonemptyWords{\TreeAlphabet{A}}}{\rA}
    \proofstepwidestar[1]{(0,a)\in\TreeAlphabet{A}}{%
      \FormulaRefAuto{TreeTokenAlphabetDef}[def]{1}}
    \proofstepwidestar[2]{%
      (1,\WordLength{S})\in\TreeAlphabet{A}}{%
      \FormulaRefAuto{TreeNodeCodeTagTyping}[thm]{2}}
    \proofstepwidestar[1]{%
      \LeafTree{A}{a}=\LetterWord{(0,a)}}{%
      \FormulaRefAuto{TreeLeafCodeDef}[def]{1}}
    \proofstepwidestar[1]{\LeafTree{A}{a}=\{(0,(0,a))\}}{%
      \rIE{5,\FormulaRefAuto{LetterWordDef}[def]{3}}}
    \proofstepwidestar[2]{%
      \LetterWord{(1,\WordLength{S})}
        \in\FiniteWords{\TreeAlphabet{A}}}{%
      \FormulaRefAuto{TreeNodeCodeTagFinite}[thm]{2}}
    \proofstepwidestar[2]{%
      \WordConcat{S}{T}\in\FiniteWords{\TreeAlphabet{A}}}{%
      \FormulaRefAuto{TreeNodeCodePayloadFinite}[thm]{2}}
    \proofstepwidestar[2]{%
      \begin{aligned}[t]
        &\LetterWord{(1,\WordLength{S})}\subseteq{}\\[-2pt]
        &\quad\WordConcat{\LetterWord{(1,\WordLength{S})}}
          {\WordConcat{S}{T}}
      \end{aligned}}{%
      \rIE{\FormulaRefAuto{WordConcatenationDef}[def]{\rAI{7,8}},
        \FormulaRefAuto{A\subseteq A\cup B}[thm]}}
    \proofstepwidestar[2]{%
      \begin{aligned}[t]
        \NodeTree{A}{S}{T}
        &=\WordConcat{\LetterWord{(1,\WordLength{S})}}
          {\WordConcat{S}{T}}
      \end{aligned}}{%
      \FormulaRefAuto{TreeNodeCodeDef}[def]{2}}
    \proofstepwidestar[2]{%
      \LetterWord{(1,\WordLength{S})}\subseteq\NodeTree{A}{S}{T}}{%
      \rIE{10,9}}
    \proofstepwidestar[2]{%
      (0,(1,\WordLength{S}))\in\LetterWord{(1,\WordLength{S})}}{%
      \rIE{\FormulaRefAuto{LetterWordDef}[def]{4},
        \FormulaRefAuto{x\in\{x\}}[thm]}}
    \proofstepwidestar[2]{%
      (0,(1,\WordLength{S}))\in\NodeTree{A}{S}{T}}{%
      \FormulaRefAuto{A\subseteq B\dsep x\in A\vdash x\in B}[thm]{11,12}}
    \proofstepwidestar[]{1\neq0}{\FormulaRefAuto{PeanoOneNeZero}[thm]}
    \proofstepwidestar[]{(1,\WordLength{S})\neq(0,a)}{%
      \FormulaRefAuto{p\neq q\vdash(p,x)\neq(q,y)}[thm]{14}}
    \proofstepwidestar[]{(0,(1,\WordLength{S}))\neq(0,(0,a))}{%
      \FormulaRefAuto{a\neq b\vdash(c,a)\neq(d,b)}[thm]{15}}
    \proofstepwidestar[]{(0,(1,\WordLength{S}))\notin\{(0,(0,a))\}}{%
      \FormulaRefAuto{x\notin\{a\}\eqvdash x\neq a}[thm]{16}}
    \proofstepwidestar[1]{(0,(1,\WordLength{S}))\notin\LeafTree{A}{a}}{%
      \rIE{6,17}}
    \proofstepwidestar[1,2]{%
      \LeafTree{A}{a}\neq\NodeTree{A}{S}{T}}{%
      \FormulaRefAuto{x\in B, x\not\in A\vdash A\neq B}[thm]{13,18}}
  \closeproofpartwideindR

  \proofpartwideindR[Linksinjektivität des Knotenkonstruktors]{%
    \begin{aligned}[t]
      &S,T,S',T'\in\NonemptyWords{\TreeAlphabet{A}}\dsep{}\\[-2pt]
      &\NodeTree{A}{S}{T}=\NodeTree{A}{S'}{T'}\vdash S=S'
    \end{aligned}
  }{S,T,S',T'\in\NonemptyWords{\TreeAlphabet{A}}\dsep
    \NodeTree{A}{S}{T}=\NodeTree{A}{S'}{T'}\vdash S=S'}%
    [TreeNodeConstructorLeftInjective]
    \BXXVIIProofTreeNodeConstructorCancellation
    \proofstepwidestar[1,2]{S=S'}{%
      \FormulaRefAuto{WordConcatenationFixedCutLeft}[thm]{%
        18,19,20,21,17,14}}
  \closeproofpartwideindR

  \proofpartwideindR[Rechtsinjektivität des Knotenkonstruktors]{%
    \begin{aligned}[t]
      &S,T,S',T'\in\NonemptyWords{\TreeAlphabet{A}}\dsep{}\\[-2pt]
      &\NodeTree{A}{S}{T}=\NodeTree{A}{S'}{T'}\vdash T=T'
    \end{aligned}
  }{S,T,S',T'\in\NonemptyWords{\TreeAlphabet{A}}\dsep
    \NodeTree{A}{S}{T}=\NodeTree{A}{S'}{T'}\vdash T=T'}%
    [TreeNodeConstructorRightInjective]
    \BXXVIIProofTreeNodeConstructorCancellation
    \proofstepwidestar[1,2]{T=T'}{%
      \FormulaRefAuto{WordConcatenationFixedCutRight}[thm]{%
        18,19,20,21,17,14}}
  \closeproofpartwideindR
\end{tabproofsplitwide}

\section{Baumstufen und Baumträger}

Die erste Stufe enthält noch keinen Baum. Im nächsten Schritt werden
Blätter sowie Knoten aus bereits vorhandenen Teilbäumen aufgenommen.
Eine Stufe begrenzt damit die Zahl der aufeinander aufbauenden Schritte,
nicht die Größe des Alphabets: Bei unendlichem \(A\) kann schon die
Menge aller Blätter unendlich sein. Jeder einzelne Baum besitzt dennoch
nur endlich viele Vorkommen. Diese beiden Endlichkeitsaussagen sind
auseinanderzuhalten.

\FormulaDefDeltaK[Erzeugungsschritt für Baumcodes]{%
  X\in\powerset(\NonemptyWords{\TreeAlphabet{A}})\vdash
  \begin{aligned}[t]
    \TreeClosure{A}{X}\coloneqq{}&
      \{\LeafTree{A}{a}\mid a\in A\}\\[-2pt]
      &\cup\{\NodeTree{A}{S}{T}\mid S,T\in X\}
  \end{aligned}%
}{TreeGenerationStepDef}{%
  \DeltaRow{Mengen}{A\dsep X\dsep a\dsep S\dsep T}
  \DeltaRow{\textbf{Neues Symbol}}{}
  \DeltaPrem{Erzeugungsschritt}{\mathcal C_A}
}

\FormulaThmDeltaK[Der Erzeugungsschritt ist eine Abbildung]{%
  \mathcal C_A\colon
  \powerset(\NonemptyWords{\TreeAlphabet{A}})
  \to\powerset(\NonemptyWords{\TreeAlphabet{A}})%
}{TreeGenerationStepFunction}{%
  \DeltaRow{Mengen}{A\dsep X}
  \DeltaRow{Funktionen}{\mathcal C_A}
}
\begin{tabproofsplitwide}
  \proofpartwide{Typisierung der Blattcodes}
  \proofstepwidestar[1]{%
    X\in\powerset(\NonemptyWords{\TreeAlphabet{A}})}{\rA}
  \proofstepwidestar[1]{%
    X\subseteq\NonemptyWords{\TreeAlphabet{A}}}{%
    \FormulaRefAuto{x\in\powerset(A)\eqvdash x\subseteq A}[thm]{1}}
  \proofstepwidestar[3]{a\in A}{\rA}
  \proofstepwidestar[3]{(0,a)\in\TreeAlphabet{A}}{%
    \FormulaRefAuto{TreeTokenAlphabetDef}[def]{3}}
  \proofstepwidestar[3]{%
    \LetterWord{(0,a)}\in\NonemptyWords{\TreeAlphabet{A}}}{%
    \FormulaRefAuto{LetterWordNonempty}[thm]{4}}
  \proofstepwidestar[3]{%
    \LeafTree{A}{a}=\LetterWord{(0,a)}}{%
    \FormulaRefAuto{TreeLeafCodeDef}[def]{3}}
  \proofstepwidestar[3]{%
    \LeafTree{A}{a}\in\NonemptyWords{\TreeAlphabet{A}}}{\rIE{6,5}}
  \proofstepwidestar[]{%
    \forall a\in A\,
      \LeafTree{A}{a}\in\NonemptyWords{\TreeAlphabet{A}}}{%
    \rUI{\rRI{3,7}}}
  \closeproofpartwide

  \proofpartwide{Typisierung der Knotencodes}
  \setcounter{proofstepnr}{8}
  \proofstepwidestar[9]{S,T\in X}{\rA}
  \proofstepwidestar[1,9]{%
    S,T\in\NonemptyWords{\TreeAlphabet{A}}}{%
    \FormulaRefAuto{A\subseteq B\dsep x\in A\vdash x\in B}[thm]{2,9}}
  \proofstepwidestar[1,9]{%
    \LetterWord{(1,\WordLength{S})}
      \in\NonemptyWords{\TreeAlphabet{A}}}{%
    \FormulaRefAuto{TreeNodeCodeTagNonempty}[thm]{10}}
  \proofstepwidestar[1,9]{%
    \WordConcat{S}{T}\in\NonemptyWords{\TreeAlphabet{A}}}{%
    \FormulaRefAuto{TreeNodeCodePayloadNonempty}[thm]{10}}
  \proofstepwidestar[1,9]{%
    \NodeTree{A}{S}{T}
      =\WordConcat{\LetterWord{(1,\WordLength{S})}}
        {\WordConcat{S}{T}}}{%
    \FormulaRefAuto{TreeNodeCodeDef}[def]{10}}
  \proofstepwidestar[1,9]{%
    \WordConcat{\LetterWord{(1,\WordLength{S})}}
      {\WordConcat{S}{T}}
      \in\NonemptyWords{\TreeAlphabet{A}}}{%
    \FormulaRefAuto{NonemptyWordConcatenationClosure}[thm]{11,12}}
  \proofstepwidestar[1,9]{%
    \NodeTree{A}{S}{T}\in
      \NonemptyWords{\TreeAlphabet{A}}}{%
    \rIE{13,14}}
  \proofstepwidestar[1]{%
    \forall S,T\in X\,
      \NodeTree{A}{S}{T}\in
        \NonemptyWords{\TreeAlphabet{A}}}{%
    \rUI{\rUI{\rRI{9,15}}}}
  \closeproofpartwide

  \proofpartwide{Abschluss und Abbildungseigenschaft}
  \setcounter{proofstepnr}{16}
  \proofstepwidestar[1]{%
    \TreeClosure{A}{X}\subseteq
      \NonemptyWords{\TreeAlphabet{A}}}{%
    \FormulaRefAuto{TreeGenerationStepDef}[def]{8,16}}
  \proofstepwidestar[1]{%
    \TreeClosure{A}{X}\in
      \powerset(\NonemptyWords{\TreeAlphabet{A}})}{%
    \FormulaRefAuto{x\in\powerset(A)\eqvdash x\subseteq A}[thm]{17}}
  \proofstepwidestar[]{%
    \forall X\in\powerset(\NonemptyWords{\TreeAlphabet{A}})\,
      \TreeClosure{A}{X}\in
        \powerset(\NonemptyWords{\TreeAlphabet{A}})}{%
    \rUI{\rRI{1,18}}}
  \proofstepwidestar[]{%
    \mathcal C_A\colon
      \powerset(\NonemptyWords{\TreeAlphabet{A}})
      \to\powerset(\NonemptyWords{\TreeAlphabet{A}})}{%
    \FormulaRefAuto{TypedFunctionDefinitionPrinciple}[thm]{19}}
  \closeproofpartwide
\end{tabproofsplitwide}

\FormulaThmDeltaK[Monotonie des Erzeugungsschritts]{%
  \begin{aligned}[t]
    &X,Y\in\powerset(\NonemptyWords{\TreeAlphabet{A}})\dsep
      X\subseteq Y\vdash{}\\[-2pt]
    &\TreeClosure{A}{X}\subseteq\TreeClosure{A}{Y}
  \end{aligned}%
}{TreeGenerationStepMonotonicity}{%
  \DeltaRow{Mengen}{A\dsep X\dsep Y\dsep R\dsep a\dsep S\dsep T}
}
\begin{tabproofsplitwide}
    \proofpartwideindR[Fall 1: Blattcode]{%
    \begin{aligned}[t]
      &X,Y\in\powerset(\NonemptyWords{\TreeAlphabet{A}})\dsep{}\\[-2pt]
      &\exists a\in A\,(R=\LeafTree{A}{a})\vdash{}\\[-2pt]
      &R\in\TreeClosure{A}{Y}
    \end{aligned}
  }{\begin{aligned}[t]
      &X,Y\in\powerset(\NonemptyWords{\TreeAlphabet{A}})\dsep{}\\[-2pt]
      &\exists a\in A\,(R=\LeafTree{A}{a})\vdash{}\\[-2pt]
      &R\in\TreeClosure{A}{Y}
    \end{aligned}}[TreeGenerationMonotonicityLeafCase]
    \proofstepwidestar[1]{%
      X,Y\in\powerset(\NonemptyWords{\TreeAlphabet{A}})}{%
      \rA}
    \proofstepwidestar[2]{%
      \exists a\in A\,(R=\LeafTree{A}{a})}{%
      \rA}
    \proofstepwidestar[3]{%
      a\in A\land R=\LeafTree{A}{a}}{%
      \rA}
    \proofstepwidestar[1,3]{%
      R\in\TreeClosure{A}{Y}}{%
      \FormulaRefAuto{TreeGenerationStepDef}[def]{1,\rAEa{3},\rAEb{3}}}
    \proofstepwidestar[1,2]{%
      R\in\TreeClosure{A}{Y}}{%
      \rEE{2,3,4}}
  \closeproofpartwideindR

  \proofpartwideindR[Fall 2: Knotencode]{%
    \begin{aligned}[t]
      &X,Y\in\powerset(\NonemptyWords{\TreeAlphabet{A}})\dsep{}\\[-2pt]
      &X\subseteq Y\dsep{}\\[-2pt]
      &\exists S,T\in X\,(R=\NodeTree{A}{S}{T})\vdash{}\\[-2pt]
      &R\in\TreeClosure{A}{Y}
    \end{aligned}
  }{\begin{aligned}[t]
      &X,Y\in\powerset(\NonemptyWords{\TreeAlphabet{A}})\dsep{}\\[-2pt]
      &X\subseteq Y\dsep{}\\[-2pt]
      &\exists S,T\in X\,(R=\NodeTree{A}{S}{T})\vdash{}\\[-2pt]
      &R\in\TreeClosure{A}{Y}
    \end{aligned}}[TreeGenerationMonotonicityNodeCase]
    \proofstepwidestar[1]{%
      X,Y\in\powerset(\NonemptyWords{\TreeAlphabet{A}})}{%
      \rA}
    \proofstepwidestar[2]{%
      X\subseteq Y}{%
      \rA}
    \proofstepwidestar[3]{%
      \exists S,T\in X\,(R=\NodeTree{A}{S}{T})}{%
      \rA}
    \proofstepwidestar[4]{%
      S,T\in X\land R=\NodeTree{A}{S}{T}}{%
      \rA}
    \proofstepwidestar[2,4]{%
      S,T\in Y}{%
      \FormulaRefAuto{A\subseteq B\dsep x\in A\vdash x\in B}[thm]{
      2,\rAEa{4}}}
    \proofstepwidestar[1,2,4]{%
      R\in\TreeClosure{A}{Y}}{%
      \FormulaRefAuto{TreeGenerationStepDef}[def]{1,5,\rAEb{4}}}
    \proofstepwidestar[1,2,3]{%
      R\in\TreeClosure{A}{Y}}{%
      \rEE{3,4,6}}
  \closeproofpartwideindR
  \proofpartwide{Zusammenführung der Fälle}
    \proofstepwidestar[1]{%
      X,Y\in\powerset(\NonemptyWords{\TreeAlphabet{A}})}{%
      \rA}
    \proofstepwidestar[2]{%
      X\subseteq Y}{%
      \rA}
    \proofstepwidestar[3]{%
      R\in\TreeClosure{A}{X}}{%
      \rA}
    \proofstepwidestar[1,3]{%
      \begin{aligned}[t]
      &\exists a\in A\,(R=\LeafTree{A}{a})\ \lor{}\\[-2pt]
      &\exists S,T\in X\,(R=\NodeTree{A}{S}{T})
    \end{aligned}}{%
      \FormulaRefAuto{TreeGenerationStepDef}[def]{1,3}}
    \proofstepwidestar[5]{%
      \exists a\in A\,(R=\LeafTree{A}{a})}{%
      \rA}
    \proofstepwidestar[1,5]{%
      R\in\TreeClosure{A}{Y}}{%
      \FormulaRefAuto{TreeGenerationMonotonicityLeafCase}[thm]{1,5}}
    \proofstepwidestar[7]{%
      \exists S,T\in X\,(R=\NodeTree{A}{S}{T})}{%
      \rA}
    \proofstepwidestar[1,2,7]{%
      R\in\TreeClosure{A}{Y}}{%
      \FormulaRefAuto{TreeGenerationMonotonicityNodeCase}[thm]{1,2,7}}
    \proofstepwidestar[1,2,3]{%
      R\in\TreeClosure{A}{Y}}{%
      \rOE{4,5,6,7,8}}
    \proofstepwidestar[1,2]{%
      \TreeClosure{A}{X}\subseteq\TreeClosure{A}{Y}}{%
      \FormulaRefAuto{SubsetIntroductionRule}[thm]{[3]\vdots 9}}
  \closeproofpartwide

\end{tabproofsplitwide}

\FormulaDefDeltaK[Stufen der Klammerungsbäume]{%
  n\in\mathbb N\vdash
  \TreeStage{A}{n}\coloneqq
  \RecFun{\varnothing}{\mathcal C_A}(n)%
}{TreeStageDef}{%
  \DeltaRow{Mengen}{A\dsep n\dsep\TreeStage{A}{n}}
  \DeltaRow{Funktionen}{\mathcal C_A}
  \DeltaRow{\textbf{Neues Symbol}}{}
  \DeltaPrem{Baumstufe}{\TreeStage{A}{n}}
}

\FormulaThmDeltaK[Typisierung der gesamten Baumstufenfolge]{%
  \RecFun{\varnothing}{\mathcal C_A}\colon
    \mathbb N\to\powerset(\NonemptyWords{\TreeAlphabet A})%
}{TreeStageSequenceTyping}{%
  \DeltaRow{Mengen}{A}
  \DeltaRow{Funktionen}{\mathcal C_A}
}
\begin{tabproofwide}
  \proofstepwidestar[]{%
    \mathcal C_A\colon\powerset(\NonemptyWords{\TreeAlphabet A})
      \to\powerset(\NonemptyWords{\TreeAlphabet A})}{%
    \FormulaRefAuto{TreeGenerationStepFunction}[thm]}
  \proofstepwidestar[]{\varnothing\subseteq\NonemptyWords{\TreeAlphabet A}}{%
    \FormulaRefAuto{\varnothing\subseteq A}[thm]}
  \proofstepwidestar[]{\varnothing\in\powerset(\NonemptyWords{\TreeAlphabet A})}{%
    \FormulaRefAuto{x\in\powerset(A)\eqvdash x\subseteq A}[thm]{2}}
  \proofstepwidestar[]{%
    \RecFun{\varnothing}{\mathcal C_A}\colon
      \mathbb N\to\powerset(\NonemptyWords{\TreeAlphabet A})}{%
    \FormulaRefAuto{m_0\in M\dsep f\colon M\to M\vdash
      \RecFun{m_0}{f}\colon\mathbb{N}\to M}[thm]{3,1}}
\end{tabproofwide}

\FormulaThmDeltaK[Nullgleichung der Baumstufen]{%
  \TreeStage{A}{0}=\varnothing%
}{TreeStageZeroEquation}{%
  \DeltaRow{Mengen}{A}
}
\begin{tabproofwide}
  \proofstepwidestar[]{%
    \mathcal C_A\colon
    \powerset(\NonemptyWords{\TreeAlphabet{A}})
    \to\powerset(\NonemptyWords{\TreeAlphabet{A}})}{%
    \FormulaRefAuto{TreeGenerationStepFunction}[thm]}
  \proofstepwidestar[]{%
    \varnothing\subseteq\NonemptyWords{\TreeAlphabet{A}}}{%
    \FormulaRefAuto{\varnothing\subseteq A}[thm]}
  \proofstepwidestar[]{\varnothing\in
    \powerset(\NonemptyWords{\TreeAlphabet{A}})}{%
    \FormulaRefAuto{x\in\powerset(A)\eqvdash x\subseteq A}[thm]{2}}
  \proofstepwidestar[]{0\in\mathbb N}{%
    \FormulaRefAuto{PeanoZero}[ax]}
  \proofstepwidestar[]{%
    \RecFun{\varnothing}{\mathcal C_A}(0)=\varnothing}{%
    \FormulaRefAuto{m_0\in M\dsep f\colon M\to M \vdash
      \RecFun{m_0}{f}(0)=m_0}[thm]{3,1}}
  \proofstepwidestar[]{\TreeStage{A}{0}=\varnothing}{%
    \FormulaRefAuto{TreeStageDef}[def]{4,5}}
\end{tabproofwide}

\FormulaThmDeltaK[Nachfolgergleichung der Baumstufen]{%
  n\in\mathbb N\vdash
  \TreeStage{A}{n+1}=\TreeClosure{A}{\TreeStage{A}{n}}%
}{TreeStageSuccessorEquation}{%
  \DeltaRow{Mengen}{A\dsep n}
}
\begin{tabproofwide}
  \proofstepwidestar[]{%
    \mathcal C_A\colon
    \powerset(\NonemptyWords{\TreeAlphabet{A}})
    \to\powerset(\NonemptyWords{\TreeAlphabet{A}})}{%
    \FormulaRefAuto{TreeGenerationStepFunction}[thm]}
  \proofstepwidestar[]{%
    \varnothing\subseteq\NonemptyWords{\TreeAlphabet{A}}}{%
    \FormulaRefAuto{\varnothing\subseteq A}[thm]}
  \proofstepwidestar[]{\varnothing\in
    \powerset(\NonemptyWords{\TreeAlphabet{A}})}{%
    \FormulaRefAuto{x\in\powerset(A)\eqvdash x\subseteq A}[thm]{2}}
  \proofstepwidestar[4]{n\in\mathbb N}{\rA}
  \proofstepwidestar[4]{n+1\in\mathbb N}{%
    \FormulaRefAuto{PeanoAddOneClosure}[thm]{4}}
  \proofstepwidestar[4]{%
    \RecFun{\varnothing}{\mathcal C_A}(n+1)
      =\mathcal C_A(\RecFun{\varnothing}{\mathcal C_A}(n))}{%
    \FormulaRefAuto{PeanoRecursionAddOne}[thm]{3,1,4}}
  \proofstepwidestar[4]{%
    \TreeStage{A}{n+1}=\TreeClosure{A}{\TreeStage{A}{n}}}{%
    \FormulaRefAuto{TreeStageDef}[def]{4,5,6}}
\end{tabproofwide}

\FormulaThmDeltaK[Typisierung der Baumstufen]{%
  n\in\mathbb N\vdash
  \begin{aligned}[t]
    &\text{(i)}\;\TreeStage{A}{n}\in
      \powerset(\NonemptyWords{\TreeAlphabet{A}})\\[-2pt]
    &\text{(ii)}\;\TreeStage{A}{n}\subseteq
      \NonemptyWords{\TreeAlphabet{A}}
  \end{aligned}%
}{TreeStageTyping}{%
  \DeltaRow{Mengen}{A\dsep n}
}
\begin{tabproofsplitwide}
  \proofpartwideindR[Potenzmengentypisierung]{%
    n\in\mathbb N\vdash\TreeStage{A}{n}\in
      \powerset(\NonemptyWords{\TreeAlphabet{A}})
  }{n\in\mathbb N\vdash\TreeStage{A}{n}\in
      \powerset(\NonemptyWords{\TreeAlphabet{A}})}%
    [TreeStagePowerSetTyping]
  \proofstepwidestar[1]{n\in\mathbb N}{\rA}
  \proofstepwidestar[]{%
      \RecFun{\varnothing}{\mathcal C_A}\colon
        \mathbb N\to\powerset(\NonemptyWords{\TreeAlphabet{A}})}{\FormulaRefAuto{TreeStageSequenceTyping}[thm]}
  \proofstepwidestar[1]{%
      \RecFun{\varnothing}{\mathcal C_A}(n)\in
        \powerset(\NonemptyWords{\TreeAlphabet{A}})}{%
      \FormulaRefAuto{F\colon A\to B\dsep x\in A
        \vdash F(x)\in B}[thm]{2,1}}
  \proofstepwidestar[1]{%
      \TreeStage{A}{n}\in
        \powerset(\NonemptyWords{\TreeAlphabet{A}})}{%
      \FormulaRefAuto{TreeStageDef}[def]{1,3}}
\closeproofpartwideindR

  \proofpartwideindR[Trägerinklusion]{%
    n\in\mathbb N\vdash\TreeStage{A}{n}\subseteq
      \NonemptyWords{\TreeAlphabet{A}}
  }{n\in\mathbb N\vdash\TreeStage{A}{n}\subseteq
      \NonemptyWords{\TreeAlphabet{A}}}%
    [TreeStageCarrierSubset]
    \proofstepwidestar[1]{n\in\mathbb N}{\rA}
    \proofstepwidestar[1]{%
      \TreeStage{A}{n}\in
        \powerset(\NonemptyWords{\TreeAlphabet{A}})}{%
      \FormulaRefAuto{TreeStagePowerSetTyping}[thm]{1}}
    \proofstepwidestar[1]{%
      \TreeStage{A}{n}\subseteq
        \NonemptyWords{\TreeAlphabet{A}}}{%
      \FormulaRefAuto{x\in\powerset(A)\eqvdash x\subseteq A}[thm]{2}}
  \closeproofpartwideindR
\end{tabproofsplitwide}

\FormulaThmDeltaK[Monotonie der Baumstufen]{%
  n\in\mathbb N\vdash
  \TreeStage{A}{n}\subseteq\TreeStage{A}{n+1}%
}{TreeStageMonotone}{%
  \DeltaRow{Mengen}{A\dsep n}
}
\begin{tabproofsplitwide}
  \proofpartwideindR[Induktionsanfang]{%
    \TreeStage{A}{0}\subseteq\TreeStage{A}{0+1}
  }{%
    \TreeStage{A}{0}\subseteq\TreeStage{A}{0+1}
  }[TreeStageMonotoneBase]
    \proofstepwidestar[]{\TreeStage{A}{0}=\varnothing}{%
      \FormulaRefAuto{TreeStageZeroEquation}[thm]}
    \proofstepwidestar[]{%
      \varnothing\subseteq\TreeStage{A}{0+1}}{%
      \FormulaRefAuto{\varnothing\subseteq A}[thm]}
    \proofstepwidestar[]{%
      \TreeStage{A}{0}\subseteq\TreeStage{A}{0+1}}{\rIE{1,2}}
  \closeproofpartwideindR

  \proofpartwideindR[Induktionsschritt]{%
    \begin{aligned}[t]
      &n\in\mathbb N\dsep
        \TreeStage{A}{n}\subseteq\TreeStage{A}{n+1}\vdash{}\\[-2pt]
      &\TreeStage{A}{n+1}
        \subseteq\TreeStage{A}{(n+1)+1}
    \end{aligned}
  }{%
    \begin{aligned}[t]
      &n\in\mathbb N\dsep
        \TreeStage{A}{n}\subseteq\TreeStage{A}{n+1}\vdash{}\\[-2pt]
      &\TreeStage{A}{n+1}
        \subseteq\TreeStage{A}{(n+1)+1}
    \end{aligned}
  }[TreeStageMonotoneStep]
    \proofstepwidestar[1]{n\in\mathbb N}{\rA}
    \proofstepwidestar[2]{%
      \TreeStage{A}{n}\subseteq\TreeStage{A}{n+1}}{\rA}
    \proofstepwidestar[1]{n+1\in\mathbb N}{%
      \FormulaRefAuto{PeanoAddOneClosure}[thm]{1}}
    \proofstepwidestar[1]{%
      \TreeStage{A}{n}
        \in\powerset(\NonemptyWords{\TreeAlphabet{A}})}{%
      \FormulaRefAuto{TreeStagePowerSetTyping}[thm]{1}}
    \proofstepwidestar[1]{%
      \TreeStage{A}{n+1}
        \in\powerset(\NonemptyWords{\TreeAlphabet{A}})}{%
      \FormulaRefAuto{TreeStagePowerSetTyping}[thm]{3}}
    \proofstepwidestar[1]{%
      \TreeStage{A}{n},\TreeStage{A}{n+1}
        \in\powerset(\NonemptyWords{\TreeAlphabet{A}})}{\rAI{4,5}}
    \proofstepwidestar[1,2]{%
      \TreeClosure{A}{\TreeStage{A}{n}}
      \subseteq
      \TreeClosure{A}{\TreeStage{A}{n+1}}}{%
      \FormulaRefAuto{TreeGenerationStepMonotonicity}[thm]{6,2}}
    \proofstepwidestar[1,2]{%
      \TreeStage{A}{n+1}
      \subseteq\TreeStage{A}{(n+1)+1}}{%
      \FormulaRefAuto{TreeStageSuccessorEquation}[thm]{1,3,7}}
  \closeproofpartwideindR

  \proofpartwide{Induktionsschluss}
    \proofstepwidestar[1]{n\in\mathbb N}{\rA}
    \proofstepwidestar[]{%
      \TreeStage{A}{0}\subseteq\TreeStage{A}{0+1}}{%
      \FormulaRefAuto{TreeStageMonotoneBase}[thm]}
    \proofstepwidestar[]{%
      \forall k\in\mathbb N\,\bigl(
        \TreeStage{A}{k}\subseteq\TreeStage{A}{k+1}
        \rightarrow
        \TreeStage{A}{k+1}
          \subseteq\TreeStage{A}{(k+1)+1}\bigr)}{%
      \FormulaRefAuto{TreeStageMonotoneStep}[thm]}
    \proofstepwidestar[1]{%
      \TreeStage{A}{n}\subseteq\TreeStage{A}{n+1}}{%
      \FormulaRefAuto{PeanoInductionRuleAddOne}[thm]{2,3,1}}
  \closeproofpartwide
\end{tabproofsplitwide}

\FormulaThmDeltaK[Additive Verschachtelung der Baumstufen]{%
  n,k\in\mathbb N\vdash
  \TreeStage{A}{n}\subseteq\TreeStage{A}{n+k}%
}{TreeStageAdditiveNested}{%
  \DeltaRow{Mengen}{A\dsep n\dsep k}
}
\begin{tabproofwide}
  \proofstepwidestar[1]{n,k\in\mathbb N}{\rA}
  \proofstepwidestar[]{%
    \RecFun{\varnothing}{\mathcal C_A}\colon
      \mathbb N\to\powerset(\NonemptyWords{\TreeAlphabet A})}{%
    \FormulaRefAuto{TreeStageSequenceTyping}[thm]}
  \proofstepwidestar[]{%
    \forall j\in\mathbb N\,
      \TreeStage A j\subseteq\TreeStage A{j+1}}{%
    \FormulaRefAuto{TreeStageMonotone}[thm]}
  \proofstepwidestar[]{%
    \forall j\in\mathbb N\,
      \RecFun{\varnothing}{\mathcal C_A}(j)
        \subseteq\RecFun{\varnothing}{\mathcal C_A}(j+1)}{%
    \FormulaRefAuto{TreeStageDef}[def]{3}}
  \proofstepwidestar[1]{%
    \RecFun{\varnothing}{\mathcal C_A}(n)
      \subseteq\RecFun{\varnothing}{\mathcal C_A}(n+k)}{%
    \FormulaRefAuto{SubsetSequenceAdditiveNesting}[thm]{2,4,1}}
  \proofstepwidestar[1]{n+k\in\mathbb N}{%
    \FormulaRefAuto{PeanoAddClosure}[thm]{1}}
  \proofstepwidestar[1]{\TreeStage A n\subseteq\TreeStage A{n+k}}{%
    \FormulaRefAuto{TreeStageDef}[def]{1,6,5}}
\end{tabproofwide}

\FormulaThmDeltaK[Verschachtelung der Baumstufen]{%
  m,n\in\mathbb N\dsep m\leq n\vdash
  \TreeStage{A}{m}\subseteq\TreeStage{A}{n}%
}{TreeStageNested}{%
  \DeltaRow{Mengen}{A\dsep m\dsep n\dsep k}
  \DeltaRow{Relationsoperatoren}{\leq_{\mathbb N}}
}
\begin{tabproofwide}
  \proofstepwidestar[1]{m,n\in\mathbb N}{\rA}
  \proofstepwidestar[2]{m\leq n}{\rA}
  \proofstepwidestar[1,2]{\exists k\in\mathbb N\,(m+k=n)}{%
    \FormulaRefAuto{m\in\mathbb{N}\dsep n\in\mathbb{N}\vdash
      m\leq n \leftrightarrow
      \exists k\in\mathbb{N}\,(m+k=n)}[thm]{1,2}}
  \proofstepwidestar[4]{k\in\mathbb N\land m+k=n}{\rA}
  \proofstepwidestar[4]{k\in\mathbb N}{\rAEa{4}}
  \proofstepwidestar[1,4]{\TreeStage A m\subseteq\TreeStage A{m+k}}{%
    \FormulaRefAuto{TreeStageAdditiveNested}[thm]{1,5}}
  \proofstepwidestar[1,4]{\TreeStage A m\subseteq\TreeStage A n}{\rIE{\rAEb{4},6}}
  \proofstepwidestar[1,2]{\TreeStage A m\subseteq\TreeStage A n}{\rEE{3,4,7}}
\end{tabproofwide}

\FormulaDefDeltaK[Menge der vollen planaren Klammerungsbäume]{%
  \BracketTreeSet{A}\coloneqq
  \bigl\{R\in\NonemptyWords{\TreeAlphabet{A}}\bigm|
    \exists n\in\mathbb N\,
      (R\in\TreeStage{A}{n})\bigr\}%
}{FullPlanarBinaryTreeSetDef}{%
  \DeltaRow{Mengen}{A\dsep R\dsep n\dsep\BracketTreeSet{A}}
  \DeltaRow{\textbf{Neues Symbol}}{}
  \DeltaPrem{Klammerungsbäume}{\BracketTreeSet{A}}
}

\FormulaThmDeltaK[Jede Baumstufe liegt im Baumträger]{%
  n\in\mathbb N\vdash
  \TreeStage{A}{n}\subseteq\BracketTreeSet{A}%
}{TreeStageSubsetTreeSet}{%
  \DeltaRow{Mengen}{A\dsep n\dsep R}
}
\begin{tabproofwide}
  \proofstepwidestar[1]{n\in\mathbb N}{\rA}
  \proofstepwidestar[2]{R\in\TreeStage{A}{n}}{\rA}
  \proofstepwidestar[1]{%
    \TreeStage{A}{n}\subseteq
      \NonemptyWords{\TreeAlphabet{A}}}{%
    \FormulaRefAuto{TreeStageCarrierSubset}[thm]{1}}
  \proofstepwidestar[1,2]{%
    R\in\NonemptyWords{\TreeAlphabet{A}}}{%
    \FormulaRefAuto{A\subseteq B\dsep x\in A\vdash x\in B}[thm]{3,2}}
  \proofstepwidestar[1,2]{%
    \exists k\in\mathbb N\,(R\in\TreeStage{A}{k})}{%
    \rEI{\rAI{1,2}}}
  \proofstepwidestar[1,2]{R\in\BracketTreeSet{A}}{%
    \FormulaRefAuto{FullPlanarBinaryTreeSetDef}[def]{4,5}}
  \proofstepwidestar[1]{%
    \TreeStage{A}{n}\subseteq\BracketTreeSet{A}}{%
    \FormulaRefAuto{SubsetIntroductionRule}[thm]{[2]\vdots 6}}
\end{tabproofwide}

\FormulaThmDeltaK[Gemeinsame Baumstufe]{%
  S,T\in\BracketTreeSet{A}\vdash
    \exists q\in\mathbb N\,
      (S,T\in\TreeStage{A}{q})%
}{TreeCommonStage}{%
  \DeltaRow{Mengen}{A\dsep S\dsep T\dsep m\dsep n\dsep q}
}
\begin{tabproofsplitwide}
    \proofpartwideindR[Fall 1: Die erste Stufe liegt vor der zweiten]{%
    \begin{aligned}[t]
      &m\in\mathbb N\land S\in\TreeStage{A}{m}\dsep{}\\[-2pt]
      &n\in\mathbb N\land T\in\TreeStage{A}{n}\dsep{}\\[-2pt]
      &m\leq n\vdash{}\\[-2pt]
      &\exists q\in\mathbb N\,
      (S,T\in\TreeStage{A}{q})
    \end{aligned}
  }{\begin{aligned}[t]
      &m\in\mathbb N\land S\in\TreeStage{A}{m}\dsep{}\\[-2pt]
      &n\in\mathbb N\land T\in\TreeStage{A}{n}\dsep{}\\[-2pt]
      &m\leq n\vdash{}\\[-2pt]
      &\exists q\in\mathbb N\,
      (S,T\in\TreeStage{A}{q})
    \end{aligned}}[TreeCommonStageFirstBeforeSecond]
    \proofstepwidestar[1]{%
      m\in\mathbb N\land S\in\TreeStage{A}{m}}{%
      \rA}
    \proofstepwidestar[2]{%
      n\in\mathbb N\land T\in\TreeStage{A}{n}}{%
      \rA}
    \proofstepwidestar[1,2]{%
      m,n\in\mathbb N}{%
      \rAI{\rAEa{1},\rAEa{2}}}
    \proofstepwidestar[1,2]{%
      S\in\TreeStage{A}{m}\land T\in\TreeStage{A}{n}}{%
      \rAI{\rAEb{1},\rAEb{2}}}
    \proofstepwidestar[5]{%
      m\leq n}{%
      \rA}
    \proofstepwidestar[1,2,5]{%
      \TreeStage{A}{m}\subseteq\TreeStage{A}{n}}{%
      \FormulaRefAuto{TreeStageNested}[thm]{3,5}}
    \proofstepwidestar[1,2,5]{%
      S\in\TreeStage{A}{n}}{%
      \FormulaRefAuto{A\subseteq B\dsep x\in A\vdash x\in B}[thm]{
      6,\rAEa{4}}}
    \proofstepwidestar[1,2,5]{%
      \exists q\in\mathbb N\,
      (S,T\in\TreeStage{A}{q})}{%
      \rEI{\rAI{\rAEa{2},\rAI{7,\rAEb{4}}}}}
  \closeproofpartwideindR

  \proofpartwideindR[Fall 2: Die zweite Stufe liegt vor der ersten]{%
    \begin{aligned}[t]
      &m\in\mathbb N\land S\in\TreeStage{A}{m}\dsep{}\\[-2pt]
      &n\in\mathbb N\land T\in\TreeStage{A}{n}\dsep{}\\[-2pt]
      &n\leq m\vdash{}\\[-2pt]
      &\exists q\in\mathbb N\,
      (S,T\in\TreeStage{A}{q})
    \end{aligned}
  }{\begin{aligned}[t]
      &m\in\mathbb N\land S\in\TreeStage{A}{m}\dsep{}\\[-2pt]
      &n\in\mathbb N\land T\in\TreeStage{A}{n}\dsep{}\\[-2pt]
      &n\leq m\vdash{}\\[-2pt]
      &\exists q\in\mathbb N\,
      (S,T\in\TreeStage{A}{q})
    \end{aligned}}[TreeCommonStageSecondBeforeFirst]
    \proofstepwidestar[1]{%
      m\in\mathbb N\land S\in\TreeStage{A}{m}}{%
      \rA}
    \proofstepwidestar[2]{%
      n\in\mathbb N\land T\in\TreeStage{A}{n}}{%
      \rA}
    \proofstepwidestar[1,2]{%
      m,n\in\mathbb N}{%
      \rAI{\rAEa{1},\rAEa{2}}}
    \proofstepwidestar[1,2]{%
      S\in\TreeStage{A}{m}\land T\in\TreeStage{A}{n}}{%
      \rAI{\rAEb{1},\rAEb{2}}}
    \proofstepwidestar[5]{%
      n\leq m}{%
      \rA}
    \proofstepwidestar[1,2,5]{%
      \TreeStage{A}{n}\subseteq\TreeStage{A}{m}}{%
      \FormulaRefAuto{TreeStageNested}[thm]{3,5}}
    \proofstepwidestar[1,2,5]{%
      T\in\TreeStage{A}{m}}{%
      \FormulaRefAuto{A\subseteq B\dsep x\in A\vdash x\in B}[thm]{
      6,\rAEb{4}}}
    \proofstepwidestar[1,2,5]{%
      \exists q\in\mathbb N\,
      (S,T\in\TreeStage{A}{q})}{%
      \rEI{\rAI{\rAEa{1},\rAI{\rAEa{4},7}}}}
  \closeproofpartwideindR
  \proofpartwide{Zusammenführung der Fälle}
    \proofstepwidestar[1]{%
      S,T\in\BracketTreeSet{A}}{%
      \rA}
    \proofstepwidestar[1]{%
      \exists m\in\mathbb N\,(S\in\TreeStage{A}{m})}{%
      \FormulaRefAuto{FullPlanarBinaryTreeSetDef}[def]{1}}
    \proofstepwidestar[1]{%
      \exists n\in\mathbb N\,(T\in\TreeStage{A}{n})}{%
      \FormulaRefAuto{FullPlanarBinaryTreeSetDef}[def]{1}}
    \proofstepwidestar[4]{%
      m\in\mathbb N\land S\in\TreeStage{A}{m}}{%
      \rA}
    \proofstepwidestar[5]{%
      n\in\mathbb N\land T\in\TreeStage{A}{n}}{%
      \rA}
    \proofstepwidestar[4,5]{%
      m,n\in\mathbb N}{%
      \rAI{\rAEa{4},\rAEa{5}}}
    \proofstepwidestar[4,5]{%
      S\in\TreeStage{A}{m}\land T\in\TreeStage{A}{n}}{%
      \rAI{\rAEb{4},\rAEb{5}}}
    \proofstepwidestar[4,5]{%
      m\leq n\lor n\leq m}{%
      \FormulaRefAuto{PeanoOrderComparison}[thm]{6}}
    \proofstepwidestar[9]{%
      m\leq n}{%
      \rA}
    \proofstepwidestar[4,5,9]{%
      \exists q\in\mathbb N\,
      (S,T\in\TreeStage{A}{q})}{%
      \FormulaRefAuto{TreeCommonStageFirstBeforeSecond}[thm]{4,5,9}}
    \proofstepwidestar[11]{%
      n\leq m}{%
      \rA}
    \proofstepwidestar[4,5,11]{%
      \exists q\in\mathbb N\,
      (S,T\in\TreeStage{A}{q})}{%
      \FormulaRefAuto{TreeCommonStageSecondBeforeFirst}[thm]{4,5,11}}
    \proofstepwidestar[4,5]{%
      \exists q\in\mathbb N\,
      (S,T\in\TreeStage{A}{q})}{%
      \rOE{8,9,10,11,12}}
    \proofstepwidestar[1,4]{%
      \exists q\in\mathbb N\,
      (S,T\in\TreeStage{A}{q})}{%
      \rEE{3,5,13}}
    \proofstepwidestar[1]{%
      \exists q\in\mathbb N\,
      (S,T\in\TreeStage{A}{q})}{%
      \rEE{2,4,14}}
  \closeproofpartwide

\end{tabproofsplitwide}

\FormulaThmDeltaK[Konstruktorabschluss der Klammerungsbäume]{%
  \begin{aligned}[t]
    &\text{(i)}\;a\in A\vdash
      \LeafTree{A}{a}\in\BracketTreeSet{A}\\[-2pt]
    &\text{(ii)}\;S,T\in\BracketTreeSet{A}\vdash
      \NodeTree{A}{S}{T}\in\BracketTreeSet{A}
  \end{aligned}%
}{TreeConstructorClosure}{%
  \DeltaRow{Mengen}{A\dsep a\dsep S\dsep T\dsep q}
}
\begin{tabproofsplitwide}
  \proofpartwideindR[Blattabschluss]{%
    a\in A\vdash\LeafTree{A}{a}\in\BracketTreeSet{A}
  }{%
    a\in A\vdash\LeafTree{A}{a}\in\BracketTreeSet{A}
  }[TreeLeafClosure]
    \proofstepwidestar[1]{a\in A}{\rA}
    \proofstepwidestar[]{0\in\mathbb N}{%
      \FormulaRefAuto{PeanoZero}[ax]}
    \proofstepwidestar[]{%
      \TreeStage{A}{0}\in
        \powerset(\NonemptyWords{\TreeAlphabet{A}})}{%
      \FormulaRefAuto{TreeStagePowerSetTyping}[thm]{2}}
    \proofstepwidestar[1]{%
      \LeafTree{A}{a}\in
        \TreeClosure{A}{\TreeStage{A}{0}}}{%
      \FormulaRefAuto{TreeGenerationStepDef}[def]{3,1}}
    \proofstepwidestar[1]{%
      \LeafTree{A}{a}\in\TreeStage{A}{0+1}}{%
      \FormulaRefAuto{TreeStageSuccessorEquation}[thm]{2,4}}
    \proofstepwidestar[]{0+1\in\mathbb N}{%
      \FormulaRefAuto{PeanoAddOneClosure}[thm]{2}}
    \proofstepwidestar[]{%
      \TreeStage{A}{0+1}\subseteq\BracketTreeSet{A}}{%
      \FormulaRefAuto{TreeStageSubsetTreeSet}[thm]{6}}
    \proofstepwidestar[1]{%
      \LeafTree{A}{a}\in\BracketTreeSet{A}}{%
      \FormulaRefAuto{A\subseteq B\dsep x\in A\vdash x\in B}[thm]{7,5}}
  \closeproofpartwideindR

  \proofpartwideindR[Knotenabschluss]{%
    S,T\in\BracketTreeSet{A}\vdash
      \NodeTree{A}{S}{T}\in\BracketTreeSet{A}
  }{%
    S,T\in\BracketTreeSet{A}\vdash
      \NodeTree{A}{S}{T}\in\BracketTreeSet{A}
  }[TreeNodeClosure]
    \proofstepwidestar[1]{S,T\in\BracketTreeSet{A}}{\rA}
    \proofstepwidestar[1]{%
      \exists q\in\mathbb N\,
        (S,T\in\TreeStage{A}{q})}{%
      \FormulaRefAuto{TreeCommonStage}[thm]{1}}
    \proofstepwidestar[3]{%
      q\in\mathbb N\land S,T\in\TreeStage{A}{q}}{\rA}
    \proofstepwidestar[3]{q\in\mathbb N}{\rAEa{3}}
    \proofstepwidestar[3]{S,T\in\TreeStage{A}{q}}{\rAEb{3}}
    \proofstepwidestar[3]{%
      \TreeStage{A}{q}\in
        \powerset(\NonemptyWords{\TreeAlphabet{A}})}{%
      \FormulaRefAuto{TreeStagePowerSetTyping}[thm]{4}}
    \proofstepwidestar[3]{%
      \NodeTree{A}{S}{T}\in
        \TreeClosure{A}{\TreeStage{A}{q}}}{%
      \FormulaRefAuto{TreeGenerationStepDef}[def]{6,5}}
    \proofstepwidestar[3]{%
      \NodeTree{A}{S}{T}\in\TreeStage{A}{q+1}}{%
      \FormulaRefAuto{TreeStageSuccessorEquation}[thm]{4,7}}
    \proofstepwidestar[3]{q+1\in\mathbb N}{%
      \FormulaRefAuto{PeanoAddOneClosure}[thm]{4}}
    \proofstepwidestar[3]{%
      \TreeStage{A}{q+1}\subseteq\BracketTreeSet{A}}{%
      \FormulaRefAuto{TreeStageSubsetTreeSet}[thm]{9}}
    \proofstepwidestar[3]{%
      \NodeTree{A}{S}{T}\in\BracketTreeSet{A}}{%
      \FormulaRefAuto{A\subseteq B\dsep x\in A\vdash x\in B}[thm]{10,8}}
    \proofstepwidestar[1]{%
      \NodeTree{A}{S}{T}\in\BracketTreeSet{A}}{\rEE{2,3,11}}
  \closeproofpartwideindR

\end{tabproofsplitwide}

% Nach TreeConstructorClosure und vor der abstrakten Baumstruktur einfügen.
\section{Kennzeichnende Eigenschaften der Baumkonstruktion}

Die Codekonstruktion hat nun ihren Zweck erfüllt: Der Baumträger
existiert, und beide Konstruktoren bleiben in ihm. Bevor daraus ein
allgemeiner Strukturbegriff entsteht, werden seine fünf entscheidenden
Eigenschaften am konkreten Modell bewiesen. Nur für diesen Nachweis
wird noch auf die Baumstufen zurückgegriffen.

\noindent\begin{minipage}{\linewidth}
\FormulaDefDeltaK[Konstruktorfunktionen des Baummodells]{%
  \begin{aligned}[t]
    \lambda_A&\coloneqq
      \{(a,z)\in A\times\BracketTreeSet A\mid z=\LeafTree A a\},\\[-2pt]
    \kappa_A&\coloneqq
      \{(p,z)\in(\BracketTreeSet A\times\BracketTreeSet A)\times\BracketTreeSet A\\[-2pt]
      &\hspace{35mm}\mid z=\NodeTree A{\pi_1(p)}{\pi_2(p)}\}
  \end{aligned}
}{BinaryTreeModelConstructorsDef}{%
  \DeltaRow{Mengen}{A\dsep a\dsep z\dsep p}
  \DeltaRow{Projektionen auf \(\BracketTreeSet A\times\BracketTreeSet A\)}{\pi_1\dsep\pi_2}
  \DeltaRow{Funktionen}{\lambda_A\dsep\kappa_A}
  \DeltaRow{\textbf{Neue Symbole}}{}
  \DeltaPrem{Blattfunktion}{\lambda_A}\DeltaPrem{Knotenfunktion}{\kappa_A}}
\end{minipage}\par

Die Graphen sind durch \(A\times\BracketTreeSet A\) beziehungsweise
\(((\BracketTreeSet A\times\BracketTreeSet A)\times\BracketTreeSet A)\)
gebunden. Aussonderung und die bereits bewiesene Abgeschlossenheit
rechtfertigen daher diese Funktionsdefinitionen.

\noindent\begin{minipage}{\linewidth}
\FormulaThmDeltaK[Typisierung der konkreten Baumkonstruktoren (B0)]{%
  \lambda_A\colon A\to\BracketTreeSet A\land
  \kappa_A\colon\BracketTreeSet A\times\BracketTreeSet A\to\BracketTreeSet A
}{TreeCodeConstructorTyping}{\DeltaRow{Mengen}{A\dsep a\dsep p}}
\end{minipage}\par
\begin{tabproofwide}
  \proofstepwidestar[1]{a\in A}{%
    \rA}
  \proofstepwidestar[1]{\LeafTree A a\in\BracketTreeSet A}{%
    \FormulaRefAuto{TreeLeafClosure}[thm]{1}}
  \proofstepwidestar[]{\forall a\in A\;\LeafTree A a\in\BracketTreeSet A}{%
    \rUI{\rRI{1,2}}}
  \proofstepwidestar[]{\lambda_A\colon A\to\BracketTreeSet A}{%
    \rIE{\FormulaRefAuto{a=b\vdash b=a}[thm]{\FormulaRefAuto{BinaryTreeModelConstructorsDef}[def]},\FormulaRefAuto{TypedTermGraphFunction}[thm]{3}}}
  \proofstepwidestar[5]{p\in\BracketTreeSet A\times\BracketTreeSet A}{%
    \rA}
  \proofstepwidestar[5]{\pi_1(p)\in\BracketTreeSet A}{%
    \FormulaRefAuto{z\in A\times B\vdash\pi_1(z)\in A}[thm]{5}}
  \proofstepwidestar[5]{\pi_2(p)\in\BracketTreeSet A}{%
    \FormulaRefAuto{z\in A\times B\vdash\pi_2(z)\in B}[thm]{5}}
  \proofstepwidestar[5]{\NodeTree A{\pi_1(p)}{\pi_2(p)}\in\BracketTreeSet A}{%
    \FormulaRefAuto{TreeNodeClosure}[thm]{\rAI{6,7}}}
  \proofstepwidestar[]{\forall p\in\BracketTreeSet A\times\BracketTreeSet A\;\NodeTree A{\pi_1(p)}{\pi_2(p)}\in\BracketTreeSet A}{%
    \rUI{\rRI{5,8}}}
  \proofstepwidestar[]{\kappa_A\colon\BracketTreeSet A\times\BracketTreeSet A\to\BracketTreeSet A}{%
    \rIE{\FormulaRefAuto{a=b\vdash b=a}[thm]{\FormulaRefAuto{BinaryTreeModelConstructorsDef}[def]},\FormulaRefAuto{TypedTermGraphFunction}[thm]{9}}}
  \proofstepwidestar[]{\begin{aligned}[t]&\lambda_A\colon A\to\BracketTreeSet A\land{}\\&\kappa_A\colon\BracketTreeSet A\times\BracketTreeSet A\to\BracketTreeSet A\end{aligned}}{%
    \rAI{4,10}}
\end{tabproofwide}

\noindent\begin{minipage}{\linewidth}
\FormulaThmDeltaK[Wert des konkreten Blattkonstruktors]{%
a\in A\vdash\lambda_A(a)=\LeafTree A a
}{TreeCodeLeafValue}{%
\DeltaRow{Mengen}{A\dsep a\dsep b}}
\end{minipage}\par
\begin{tabproofwide}
  \proofstepwidestar[1]{b\in A}{%
    \rA}
  \proofstepwidestar[1]{\LeafTree A b\in\BracketTreeSet A}{%
    \FormulaRefAuto{TreeLeafClosure}[thm]{1}}
  \proofstepwidestar[]{\forall b\in A\;\LeafTree A b\in\BracketTreeSet A}{%
    \rUI{\rRI{1,2}}}
  \proofstepwidestar[]{\forall b\in A\;\lambda_A(b)=\LeafTree A b}{%
    \rIE{\FormulaRefAuto{a=b\vdash b=a}[thm]{\FormulaRefAuto{BinaryTreeModelConstructorsDef}[def]},\FormulaRefAuto{TypedTermGraphValues}[thm]{3}}}
  \proofstepwidestar[5]{a\in A}{%
    \rA}
  \proofstepwidestar[5]{\lambda_A(a)=\LeafTree A a}{%
    \rRE{\rUE{4},5}}
\end{tabproofwide}

\noindent\begin{minipage}{\linewidth}
\FormulaThmDeltaK[Wert des konkreten Knotenkonstruktors]{%
S,R\in\BracketTreeSet A\vdash\kappa_A(S,R)=\NodeTree A S R
}{TreeCodeNodeValue}{%
\DeltaRow{Mengen}{A\dsep S\dsep R\dsep p}}
\end{minipage}\par
\begin{tabproofwide}
  \proofstepwidestar[1]{p\in\BracketTreeSet A\times\BracketTreeSet A}{%
    \rA}
  \proofstepwidestar[1]{\pi_1(p)\in\BracketTreeSet A}{%
    \FormulaRefAuto{z\in A\times B\vdash\pi_1(z)\in A}[thm]{1}}
  \proofstepwidestar[1]{\pi_2(p)\in\BracketTreeSet A}{%
    \FormulaRefAuto{z\in A\times B\vdash\pi_2(z)\in B}[thm]{1}}
  \proofstepwidestar[1]{\NodeTree A{\pi_1(p)}{\pi_2(p)}\in\BracketTreeSet A}{%
    \FormulaRefAuto{TreeNodeClosure}[thm]{\rAI{2,3}}}
  \proofstepwidestar[]{\forall p\in\BracketTreeSet A\times\BracketTreeSet A\;\NodeTree A{\pi_1(p)}{\pi_2(p)}\in\BracketTreeSet A}{%
    \rUI{\rRI{1,4}}}
  \proofstepwidestar[]{\forall p\in\BracketTreeSet A\times\BracketTreeSet A\;\kappa_A(p)=\NodeTree A{\pi_1(p)}{\pi_2(p)}}{%
    \rIE{\FormulaRefAuto{a=b\vdash b=a}[thm]{\FormulaRefAuto{BinaryTreeModelConstructorsDef}[def]},\FormulaRefAuto{TypedTermGraphValues}[thm]{5}}}
  \proofstepwidestar[7]{S,R\in\BracketTreeSet A}{%
    \rA}
  \proofstepwidestar[7]{(S,R)\in\BracketTreeSet A\times\BracketTreeSet A}{%
    \FormulaRefAuto{a\in A\dsep b\in B\vdash(a,b)\in A\times B}[thm]{\rAEa{7},\rAEb{7}}}
  \proofstepwidestar[7]{\kappa_A(S,R)=\NodeTree A{\pi_1(S,R)}{\pi_2(S,R)}}{%
    \rRE{\rUE{6},8}}
  \proofstepwidestar[7]{\pi_1(S,R)=S}{%
    \FormulaRefAuto{(x,y)\in A\times B\vdash\pi_1(x,y)=x}[thm]{8}}
  \proofstepwidestar[7]{\pi_2(S,R)=R}{%
    \FormulaRefAuto{(x,y)\in A\times B\vdash\pi_2(x,y)=y}[thm]{8}}
  \proofstepwidestar[7]{\kappa_A(S,R)=\NodeTree A S R}{%
    \rIE{10,\rIE{11,9}}}
\end{tabproofwide}

\noindent\begin{minipage}{\linewidth}
\FormulaThmDeltaK[Blatteindeutigkeit des konkreten Modells (B1)]{%
  a,b\in A\dsep\lambda_A(a)=\lambda_A(b)\vdash a=b
}{TreeCodeLeafInjectivity}{\DeltaRow{Mengen}{A\dsep a\dsep b}}
\end{minipage}\par
\begin{tabproofwide}
  \proofstepwidestar[1]{a,b\in A}{%
    \rA}
  \proofstepwidestar[2]{\lambda_A(a)=\lambda_A(b)}{%
    \rA}
  \proofstepwidestar[1]{\lambda_A(a)=\LeafTree A a}{%
    \FormulaRefAuto{TreeCodeLeafValue}[thm]{\rAEa{1}}}
  \proofstepwidestar[1]{\lambda_A(b)=\LeafTree A b}{%
    \FormulaRefAuto{TreeCodeLeafValue}[thm]{\rAEb{1}}}
  \proofstepwidestar[1,2]{\LeafTree A a=\LeafTree A b}{%
    \rIE{3,\rIE{4,2}}}
  \proofstepwidestar[1,2]{a=b}{%
    \FormulaRefAuto{TreeLeafConstructorInjective}[thm]{1,5}}
\end{tabproofwide}

\noindent\begin{minipage}{\linewidth}
\FormulaThmDeltaK[Knoteneindeutigkeit des konkreten Modells (B2)]{%
  \begin{aligned}[t]
    &S,R,S',R'\in\BracketTreeSet A\dsep{}\\[-2pt]
    &\kappa_A(S,R)=\kappa_A(S',R')\vdash S=S'\land R=R'
  \end{aligned}
}{TreeCodeNodeInjectivity}{\DeltaRow{Mengen}{A\dsep S\dsep R\dsep S'\dsep R'}}
\end{minipage}\par
\begin{tabproofwide}
  \proofstepwidestar[1]{S,R,S',R'\in\BracketTreeSet A}{%
    \rA}
  \proofstepwidestar[2]{\kappa_A(S,R)=\kappa_A(S',R')}{%
    \rA}
  \proofstepwidestar[1]{S,R,S',R'\in\NonemptyWords{\TreeAlphabet A}}{%
    \FormulaRefAuto{FullPlanarBinaryTreeSetDef}[def]{1}}
  \proofstepwidestar[1]{\kappa_A(S,R)=\NodeTree A S R}{%
    \FormulaRefAuto{TreeCodeNodeValue}[thm]{\rAEn{1}}}
  \proofstepwidestar[1]{\kappa_A(S',R')=\NodeTree A {S'}{R'}}{%
    \FormulaRefAuto{TreeCodeNodeValue}[thm]{\rAEn{1}}}
  \proofstepwidestar[1,2]{\NodeTree A S R=\NodeTree A {S'}{R'}}{%
    \rIE{4,\rIE{5,2}}}
  \proofstepwidestar[1,2]{S=S'}{%
    \FormulaRefAuto{TreeNodeConstructorLeftInjective}[thm]{3,6}}
  \proofstepwidestar[1,2]{R=R'}{%
    \FormulaRefAuto{TreeNodeConstructorRightInjective}[thm]{3,6}}
  \proofstepwidestar[1,2]{S=S'\land R=R'}{%
    \rAI{7,8}}
\end{tabproofwide}

\noindent\begin{minipage}{\linewidth}
\FormulaThmDeltaK[Trennung der konkreten Konstruktoren (B3)]{%
  a\in A\dsep S,R\in\BracketTreeSet A\vdash
    \lambda_A(a)\ne\kappa_A(S,R)
}{TreeCodeConstructorSeparation}{\DeltaRow{Mengen}{A\dsep a\dsep S\dsep R}}
\end{minipage}\par
\begin{tabproofwide}
  \proofstepwidestar[1]{a\in A}{%
    \rA}
  \proofstepwidestar[2]{S,R\in\BracketTreeSet A}{%
    \rA}
  \proofstepwidestar[2]{S,R\in\NonemptyWords{\TreeAlphabet A}}{%
    \FormulaRefAuto{FullPlanarBinaryTreeSetDef}[def]{2}}
  \proofstepwidestar[1,2]{\LeafTree A a\ne\NodeTree A S R}{%
    \FormulaRefAuto{TreeConstructorsDisjoint}[thm]{1,3}}
  \proofstepwidestar[1]{\lambda_A(a)=\LeafTree A a}{%
    \FormulaRefAuto{TreeCodeLeafValue}[thm]{1}}
  \proofstepwidestar[2]{\kappa_A(S,R)=\NodeTree A S R}{%
    \FormulaRefAuto{TreeCodeNodeValue}[thm]{2}}
  \proofstepwidestar[1,2]{\lambda_A(a)\ne\kappa_A(S,R)}{%
    \rIE{\FormulaRefAuto{a=b\vdash b=a}[thm]{5},\rIE{\FormulaRefAuto{a=b\vdash b=a}[thm]{6},4}}}
\end{tabproofwide}

Die letzte Eigenschaft besagt, dass keine echte Teilmenge des Trägers
alle Blätter enthält und unter der Knotenbildung abgeschlossen ist.
Ihr Beweis benutzt zunächst nur die Gleichung für die nächste
Baumstufe. Die Zerlegung eines beliebigen Baumes und die
Strukturinduktion werden dabei noch nicht vorausgesetzt.

\noindent\begin{minipage}{\linewidth}
\FormulaThmDeltaK[Abgeschlossene Teilmengen enthalten die nächste Baumstufe]{%
  \begin{aligned}[t]
    &U\subseteq\BracketTreeSet A\dsep\lambda_A[A]\subseteq U\dsep{}\\[-2pt]
    &\forall S,R\in U\;\kappa_A(S,R)\in U\dsep{}\\[-2pt]
    &n\in\mathbb N\dsep\TreeStage A n\subseteq U
      \vdash\TreeStage A {n+1}\subseteq U
  \end{aligned}
}{TreeCodeStageContainmentStep}{\DeltaRow{Mengen}{A\dsep U\dsep n\dsep a\dsep b\dsep S\dsep R\dsep Q}}
\end{minipage}\par
\begin{tabproofwide}
  \proofstepwidestar[1]{U\subseteq\BracketTreeSet A}{%
    \rA}
  \proofstepwidestar[2]{\lambda_A[A]\subseteq U}{%
    \rA}
  \proofstepwidestar[3]{\forall S,R\in U\;\kappa_A(S,R)\in U}{%
    \rA}
  \proofstepwidestar[4]{n\in\mathbb N}{%
    \rA}
  \proofstepwidestar[5]{\TreeStage A n\subseteq U}{%
    \rA}
  \proofstepwidestar[]{\lambda_A\colon A\to\BracketTreeSet A}{%
    \rAEa{\FormulaRefAuto{TreeCodeConstructorTyping}[thm]}}
  \proofstepwidestar[4]{\TreeStage A n\in\powerset(\NonemptyWords{\TreeAlphabet A})}{%
    \FormulaRefAuto{TreeStagePowerSetTyping}[thm]{4}}
  \proofstepwidestar[8]{Q\in\TreeStage A {n+1}}{%
    \rA}
  \proofstepwidestar[4]{\TreeStage A {n+1}=\TreeClosure A{\TreeStage A n}}{%
    \FormulaRefAuto{TreeStageSuccessorEquation}[thm]{4}}
  \proofstepwidestar[4,8]{Q\in\TreeClosure A{\TreeStage A n}}{%
    \rIE{9,8}}
  \proofstepwidestar[4]{\begin{aligned}[t]&\TreeClosure A{\TreeStage A n}\\&=\{\LeafTree A a\mid a\in A\}\\&\quad\cup\{\NodeTree A S R\mid S,R\in\TreeStage A n\}\end{aligned}}{%
    \FormulaRefAuto{TreeGenerationStepDef}[def]{7}}
  \proofstepwidestar[4,8]{\begin{aligned}[t]&Q\in\{\LeafTree A a\mid a\in A\}\\&\lor Q\in\{\NodeTree A S R\mid S,R\in\TreeStage A n\}\end{aligned}}{%
    \FormulaRefAuto{x\in A\cup B\eqvdash x\in A\lor x\in B}[thm]{\rIE{11,10}}}
  \proofstepwidestar[]{Q\in\{\LeafTree A a\mid a\in A\}\leftrightarrow\exists a\in A\;Q=\LeafTree A a}{%
    \FormulaRefAuto{B27TermImageMembership}[thm]}
  \proofstepwidestar[]{\begin{aligned}[t]&Q\in\{\NodeTree A S R\mid S,R\in\TreeStage A n\}\\&\leftrightarrow\exists S,R\in\TreeStage A n\;Q=\NodeTree A S R\end{aligned}}{%
    \FormulaRefAuto{TreeCodeBinaryTermImageMembership}[thm]}
  \proofstepwidestar[4,8]{\begin{aligned}[t]&(\exists a\in A\;Q=\LeafTree A a)\lor{}\\&(\exists S,R\in\TreeStage A n\;Q=\NodeTree A S R)\end{aligned}}{%
    \rLRS{14,\rLRS{13,12}}}
  \proofstepwidestar[16]{\exists a\in A\;Q=\LeafTree A a}{%
    \rA}
  \proofstepwidestar[17]{a\in A\land Q=\LeafTree A a}{%
    \rA}
  \proofstepwidestar[17]{\lambda_A(a)=\LeafTree A a}{%
    \FormulaRefAuto{TreeCodeLeafValue}[thm]{\rAEa{17}}}
  \proofstepwidestar[17]{\exists b\in A\;\lambda_A(a)=\lambda_A(b)}{%
    \rEI{\rAI{\rAEa{17},\rII}}}
  \proofstepwidestar[17]{\lambda_A(a)\in\lambda_A[A]}{%
    \FormulaRefAuto{F\colon A\to B\dsep\exists x\in A\;y=F(x)\vdash y\in F[A]}[thm]{6,19}}
  \proofstepwidestar[2,17]{\lambda_A(a)\in U}{%
    \FormulaRefAuto{A\subseteq B\dsep x\in A\vdash x\in B}[thm]{2,20}}
  \proofstepwidestar[2,17]{Q\in U}{%
    \rIE{\FormulaRefAuto{a=b\vdash b=a}[thm]{\rAEb{17}},\rIE{18,21}}}
  \proofstepwidestar[2,16]{Q\in U}{%
    \rEE{16,17,22}}
  \proofstepwidestar[24]{\exists S,R\in\TreeStage A n\;Q=\NodeTree A S R}{%
    \rA}
  \proofstepwidestar[25]{S\in\TreeStage A n\land\exists R\in\TreeStage A n\;Q=\NodeTree A S R}{%
    \rA}
  \proofstepwidestar[26]{R\in\TreeStage A n\land Q=\NodeTree A S R}{%
    \rA}
  \proofstepwidestar[5,26,25]{S\in U}{%
    \FormulaRefAuto{A\subseteq B\dsep x\in A\vdash x\in B}[thm]{5,\rAEn{\rAI{25,26}}}}
  \proofstepwidestar[5,26,25]{R\in U}{%
    \FormulaRefAuto{A\subseteq B\dsep x\in A\vdash x\in B}[thm]{5,\rAEn{\rAI{25,26}}}}
  \proofstepwidestar[1,5,26,25]{S\in\BracketTreeSet A}{%
    \FormulaRefAuto{A\subseteq B\dsep x\in A\vdash x\in B}[thm]{1,27}}
  \proofstepwidestar[1,5,26,25]{R\in\BracketTreeSet A}{%
    \FormulaRefAuto{A\subseteq B\dsep x\in A\vdash x\in B}[thm]{1,28}}
  \proofstepwidestar[1,5,26,25]{\kappa_A(S,R)=\NodeTree A S R}{%
    \FormulaRefAuto{TreeCodeNodeValue}[thm]{\rAI{29,30}}}
  \proofstepwidestar[3,5,26,25]{\kappa_A(S,R)\in U}{%
    \rRE{\rUE{\rRE{\rUE{3},27}},28}}
  \proofstepwidestar[1,3,5,26,25]{Q\in U}{%
    \rIE{\FormulaRefAuto{a=b\vdash b=a}[thm]{\rAEn{\rAI{25,26}}},\rIE{31,32}}}
  \proofstepwidestar[1,3,5,25]{Q\in U}{%
    \rEE{\rAEb{25},26,33}}
  \proofstepwidestar[1,3,5,24]{Q\in U}{%
    \rEE{24,25,34}}
  \proofstepwidestar[1,2,3,4,5,8]{Q\in U}{%
    \rOE{15,16,23,24,35}}
  \proofstepwidestar[1,2,3,4,5]{\TreeStage A {n+1}\subseteq U}{%
    \FormulaRefAuto{SubsetIntroductionRule}[thm]{[8]\vdots 36}}
\end{tabproofwide}

\noindent\begin{minipage}{\linewidth}
\FormulaThmDeltaK[Minimalität des konkreten Baumträgers (B4)]{%
  \begin{aligned}[t]
    &U\subseteq\BracketTreeSet A\dsep\lambda_A[A]\subseteq U\dsep{}\\[-2pt]
    &\forall S,R\in U\;\kappa_A(S,R)\in U\vdash U=\BracketTreeSet A
  \end{aligned}
}{TreeCodeMinimality}{\DeltaRow{Mengen}{A\dsep U\dsep n\dsep S\dsep R}}
\end{minipage}\par
\begin{tabproofwide}
  \proofstepwidestar[1]{U\subseteq\BracketTreeSet A}{%
    \rA}
  \proofstepwidestar[2]{\lambda_A[A]\subseteq U}{%
    \rA}
  \proofstepwidestar[3]{\forall S,R\in U\;\kappa_A(S,R)\in U}{%
    \rA}
  \proofstepwidestar[]{\TreeStage A0=\varnothing}{%
    \FormulaRefAuto{TreeStageZeroEquation}[thm]}
  \proofstepwidestar[]{\varnothing\subseteq U}{%
    \FormulaRefAuto{\varnothing\subseteq A}[thm]}
  \proofstepwidestar[]{\TreeStage A0\subseteq U}{%
    \rIE{\FormulaRefAuto{a=b\vdash b=a}[thm]{4},5}}
  \proofstepwidestar[7]{n\in\mathbb N}{%
    \rA}
  \proofstepwidestar[8]{\TreeStage A n\subseteq U}{%
    \rA}
  \proofstepwidestar[1,2,3,7,8]{\TreeStage A {n+1}\subseteq U}{%
    \FormulaRefAuto{TreeCodeStageContainmentStep}[thm]{1,2,3,7,8}}
  \proofstepwidestar[1,2,3]{\forall n\in\mathbb N\;(\TreeStage A n\subseteq U\rightarrow\TreeStage A {n+1}\subseteq U)}{%
    \rUI{\rRI{7,\rRI{8,9}}}}
  \proofstepwidestar[1,2,3]{\forall n\in\mathbb N\;\TreeStage A n\subseteq U}{%
    \FormulaRefAuto{PeanoInductionAddOne}[thm]{6,10}}
  \proofstepwidestar[12]{R\in\BracketTreeSet A}{%
    \rA}
  \proofstepwidestar[12]{\exists n\in\mathbb N\;R\in\TreeStage A n}{%
    \FormulaRefAuto{FullPlanarBinaryTreeSetDef}[def]{12}}
  \proofstepwidestar[14]{n\in\mathbb N\land R\in\TreeStage A n}{%
    \rA}
  \proofstepwidestar[1,2,3,14]{\TreeStage A n\subseteq U}{%
    \rRE{\rUE{11},\rAEa{14}}}
  \proofstepwidestar[1,2,3,14]{R\in U}{%
    \FormulaRefAuto{A\subseteq B\dsep x\in A\vdash x\in B}[thm]{15,\rAEb{14}}}
  \proofstepwidestar[1,2,3,12]{R\in U}{%
    \rEE{13,14,16}}
  \proofstepwidestar[1,2,3]{\BracketTreeSet A\subseteq U}{%
    \FormulaRefAuto{SubsetIntroductionRule}[thm]{[12]\vdots 17}}
  \proofstepwidestar[1,2,3]{U=\BracketTreeSet A}{%
    \FormulaRefAuto{A\subseteq B\dsep B\subseteq A\vdash A=B}[thm]{1,18}}
\end{tabproofwide}

Auch für \(A=\varnothing\) gelten alle fünf Eigenschaften. In diesem
Fall ist \(U=\varnothing\) unter den Konstruktoren abgeschlossen, und
die Minimalität ergibt \(\BracketTreeSet\varnothing=\varnothing\).
Ein Baum ohne beschriftetes Blatt wird durch die Konstruktion also
nicht erzeugt.

% Frühe Baumaxiome; nach b27-tree-early-foundations einfügen.
\section{Das Axiomsystem der Baumstrukturen}

Die soeben bewiesenen Eigenschaften motivieren den folgenden
Strukturbegriff. Darin besitzt jedes Blatt eine Beschriftung in
\(A\), und jeder innere Knoten hat genau zwei geordnete Kinder. Die
Reihenfolge der Argumente von \(k\) entspricht der Links-rechts-Ordnung
des planaren Baumes. Eine Vertauschung wird durch die Axiome nicht
identifiziert.

Die Symbole beschreiben ganze Bäume und ihre Konstruktion:
\begin{center}
\begin{tabularx}{\linewidth}{@{}p{0.20\linewidth}X@{}}
\(A;\ a,b\) & \(A\) ist das Alphabet der Blattbeschriftungen;
  \(a\) und \(b\) bezeichnen Elemente von \(A\).\\[3pt]
\(T;\ x,y,x',y'\) & \(T\) ist der Träger der Baumstruktur.
  Jedes seiner Elemente \(x,y,x',y'\) steht für einen ganzen Baum.\\[3pt]
\(\ell\) & Der Blattkonstruktor ordnet einer Beschriftung \(a\)
  den einzelnen Blattbaum \(\ell(a)\) zu.\\[3pt]
\(k\) & Der Knotenkonstruktor verbindet zwei ganze Bäume \(x,y\)
  unter einer neuen Wurzel zu \(k(x,y)\): \(x\) ist der linke,
  \(y\) der rechte Teilbaum.\\[3pt]
\(T\times T\) & Die Menge der geordneten Paare ganzer Bäume;
  sie ist der Definitionsbereich von \(k\).\\[3pt]
\(U;\ \ell[A]\) & \(U\subseteq T\) bezeichnet eine beliebige
  Teilmenge des Baumträgers. \(\ell[A]\) ist die Menge aller
  durch \(\ell\) erzeugten Blattbäume.
\end{tabularx}
\end{center}

Im zuvor konstruierten Modell ist die Zuordnung konkret
\[
  T=\BracketTreeSet A,\qquad \ell=\lambda_A,\qquad k=\kappa_A.
\]
Die folgenden Bedingungen B0 bis B4 werden nun für die Daten
\((T,\ell,k)\) gefordert. Die jeweils genannten Herkunftssätze
weisen sie für dieses konkrete Modell nach. Der Ausdruck
\(\mathsf{BaumStr}_A(T,\ell,k)\) besagt, dass die gewählten Daten
alle fünf Bedingungen erfüllen.

\noindent\begin{minipage}{\linewidth}
\FormulaDefDeltaK[Freie volle planare Binärbaumstruktur]{%
  \mathsf{BaumStr}_A(T,\ell,k)
}{BinaryTreeStructureDef}{%
  \DeltaRow{Mengen}{A\dsep T\dsep U\dsep a\dsep b\dsep x\dsep y\dsep x'\dsep y'}
  \DeltaRow{Funktionen}{\ell\dsep k}
  \DeltaRow{\textbf{Strukturaxiome}}{}
  \DeltaRow{B0: Typisierung}{\ell\colon A\to T\land k\colon T\times T\to T}
  \DeltaRow{B1: Blatteindeutigkeit}{%
    \forall a,b\in A\;(\ell(a)=\ell(b)\rightarrow a=b)}
  \DeltaRow{B2: Knoteneindeutigkeit}{%
    \begin{aligned}[t]
      &\forall x,y,x',y'\in T\\[-2pt]
      &\quad(k(x,y)=k(x',y')\rightarrow x=x'\land y=y')
    \end{aligned}}
  \DeltaRow{B3: Trennung}{\forall a\in A\,\forall x,y\in T\;\ell(a)\ne k(x,y)}
  \DeltaRow{B4: Minimalität}{%
    \begin{aligned}[t]
      &\forall U\subseteq T\,\Bigl(\bigl(\ell[A]\subseteq U\land{}\\[-2pt]
      &\quad(\forall x,y\in U\;k(x,y)\in U)\bigr)\rightarrow U=T\Bigr)
    \end{aligned}}
  \DeltaRow{\textbf{Neues Symbol}}{}
  \DeltaPrem{Freie Binärbaumstruktur}{\mathsf{BaumStr}_A(T,\ell,k)}}
\end{minipage}\par

\Needspace{14\baselineskip}
\noindent\begin{minipage}{\linewidth}
\FormulaAxiomDeltaK[Typisierung einer Baumstruktur]{%
  \mathsf{BaumStr}_A(T,\ell,k)\vdash
    \ell\colon A\to T\land k\colon T\times T\to T
}{BinaryTreeStructureTypingAxiom}{%
  \DeltaRow{Mengen}{A\dsep T}\DeltaRow{Funktionen}{\ell\dsep k}}
\par\smallskip\noindent
B0 verlangt, dass beide Konstruktoren wieder Bäume in \(T\)
liefern. Für \(\lambda_A\) und \(\kappa_A\) auf dem konkreten
Baumträger wurde dies in
\FormulaRefAuto{TreeCodeConstructorTyping}[thm] bewiesen.
\end{minipage}\par

\Needspace{14\baselineskip}
\noindent\begin{minipage}{\linewidth}
\FormulaAxiomDeltaK[Injektivität des Blattkonstruktors]{%
  \mathsf{BaumStr}_A(T,\ell,k)\dsep a,b\in A\dsep\ell(a)=\ell(b)\vdash a=b
}{BinaryTreeStructureLeafAxiom}{%
  \DeltaRow{Mengen}{A\dsep T\dsep a\dsep b}\DeltaRow{Funktionen}{\ell\dsep k}}
\par\smallskip\noindent
B1 besagt: Gleiche Blattbäume haben dieselbe Beschriftung.
Für den konkreten Blattkonstruktor \(\lambda_A\) ist dies
\FormulaRefAuto{TreeCodeLeafInjectivity}[thm].
\end{minipage}\par

\Needspace{15\baselineskip}
\noindent\begin{minipage}{\linewidth}
\FormulaAxiomDeltaK[Injektivität des Knotenkonstruktors]{%
  \begin{aligned}[t]
    &\mathsf{BaumStr}_A(T,\ell,k)\dsep x,y,x',y'\in T\dsep{}\\[-2pt]
    &k(x,y)=k(x',y')\vdash x=x'\land y=y'
  \end{aligned}
}{BinaryTreeStructureNodeAxiom}{%
  \DeltaRow{Mengen}{A\dsep T\dsep x\dsep y\dsep x'\dsep y'}
  \DeltaRow{Funktionen}{\ell\dsep k}}
\par\smallskip\noindent
B2 besagt: Gleiche zusammengesetzte Bäume haben denselben linken
und denselben rechten Teilbaum. Für \(\kappa_A\) wurde dies in
\FormulaRefAuto{TreeCodeNodeInjectivity}[thm] bewiesen.
\end{minipage}\par

\Needspace{14\baselineskip}
\noindent\begin{minipage}{\linewidth}
\FormulaAxiomDeltaK[Trennung von Blättern und Knoten]{%
  \mathsf{BaumStr}_A(T,\ell,k)\dsep a\in A\dsep x,y\in T
    \vdash\ell(a)\ne k(x,y)
}{BinaryTreeStructureDisjointnessAxiom}{%
  \DeltaRow{Mengen}{A\dsep T\dsep a\dsep x\dsep y}
  \DeltaRow{Funktionen}{\ell\dsep k}}
\par\smallskip\noindent
B3 schließt aus, dass ein Blattbaum zugleich aus zwei Teilbäumen
zusammengesetzt ist. Die entsprechende Trennung von
\(\lambda_A\) und \(\kappa_A\) liefert
\FormulaRefAuto{TreeCodeConstructorSeparation}[thm].
\end{minipage}\par

\Needspace{16\baselineskip}
\noindent\begin{minipage}{\linewidth}
\FormulaAxiomDeltaK[Minimalität einer Baumstruktur]{%
  \begin{aligned}[t]
    &\mathsf{BaumStr}_A(T,\ell,k)\dsep U\subseteq T\dsep\ell[A]\subseteq U\dsep{}\\[-2pt]
    &\forall x,y\in U\;k(x,y)\in U\vdash U=T
  \end{aligned}
}{BinaryTreeStructureMinimalityAxiom}{%
  \DeltaRow{Mengen}{A\dsep T\dsep U\dsep x\dsep y}
  \DeltaRow{Funktionen}{\ell\dsep k}}
\par\smallskip\noindent
B4 besagt: Enthält \(U\) alle Blattbäume \(\ell[A]\) und mit
\(x,y\) jeweils auch \(k(x,y)\), so enthält \(U\) bereits alle
Bäume von \(T\). Diese Minimalität des konkreten Baumträgers
ist \FormulaRefAuto{TreeCodeMinimality}[thm].
\end{minipage}\par

Die zuvor bewiesenen Eigenschaften werden nun in einem Strukturbegriff
gebündelt. Die registrierten Axiome sind Zugriffsregeln auf die
Bedingungen B0 bis B4; sie fordern keine neuen Mengenexistenzaxiome.

\noindent\begin{minipage}{\linewidth}
\FormulaThmDeltaK[Existenz einer freien Klammerungsbaumstruktur]{%
  \mathsf{BaumStr}_A(\BracketTreeSet A,\lambda_A,\kappa_A)
}{BinaryTreeStructureConcreteModel}{%
  \DeltaRow{Mengen}{A\dsep U\dsep S\dsep R\dsep S'\dsep R'\dsep a\dsep b}}
\end{minipage}\par
\begin{tabproofwide}
  \proofstepwidestar[]{\begin{aligned}[t]&\lambda_A\colon A\to\BracketTreeSet A\land{}\\&\kappa_A\colon\BracketTreeSet A\times\BracketTreeSet A\to\BracketTreeSet A\end{aligned}}{%
    \FormulaRefAuto{TreeCodeConstructorTyping}[thm]}
  \proofstepwidestar[2]{a\in A}{%
    \rA}
  \proofstepwidestar[3]{b\in A}{%
    \rA}
  \proofstepwidestar[4]{\lambda_A(a)=\lambda_A(b)}{%
    \rA}
  \proofstepwidestar[2,3,4]{a=b}{%
    \FormulaRefAuto{TreeCodeLeafInjectivity}[thm]{\rAI{2,3},4}}
  \proofstepwidestar[]{\forall a,b\in A\;(\lambda_A(a)=\lambda_A(b)\rightarrow a=b)}{%
    \rUI{\rRI{2,\rUI{\rRI{3,\rRI{4,5}}}}}}
  \proofstepwidestar[7]{S\in\BracketTreeSet A}{%
    \rA}
  \proofstepwidestar[8]{R\in\BracketTreeSet A}{%
    \rA}
  \proofstepwidestar[9]{S'\in\BracketTreeSet A}{%
    \rA}
  \proofstepwidestar[10]{R'\in\BracketTreeSet A}{%
    \rA}
  \proofstepwidestar[11]{\kappa_A(S,R)=\kappa_A(S',R')}{%
    \rA}
  \proofstepwidestar[7,8,9,10,11]{S=S'\land R=R'}{%
    \FormulaRefAuto{TreeCodeNodeInjectivity}[thm]{\rAI{7,\rAI{8,\rAI{9,10}}},11}}
  \proofstepwidestar[]{\begin{aligned}[t]&\forall S,R,S',R'\in\BracketTreeSet A\\&(\kappa_A(S,R)=\kappa_A(S',R')\rightarrow S=S'\land R=R')\end{aligned}}{%
    \rUI{\rRI{7,\rUI{\rRI{8,\rUI{\rRI{9,\rUI{\rRI{10,\rRI{11,12}}}}}}}}}}
  \proofstepwidestar[14]{a\in A}{%
    \rA}
  \proofstepwidestar[15]{S\in\BracketTreeSet A}{%
    \rA}
  \proofstepwidestar[16]{R\in\BracketTreeSet A}{%
    \rA}
  \proofstepwidestar[14,15,16]{\lambda_A(a)\ne\kappa_A(S,R)}{%
    \FormulaRefAuto{TreeCodeConstructorSeparation}[thm]{14,\rAI{15,16}}}
  \proofstepwidestar[]{\forall a\in A\,\forall S,R\in\BracketTreeSet A\;\lambda_A(a)\ne\kappa_A(S,R)}{%
    \rUI{\rRI{14,\rUI{\rRI{15,\rUI{\rRI{16,17}}}}}}}
  \proofstepwidestar[19]{U\subseteq\BracketTreeSet A}{%
    \rA}
  \proofstepwidestar[20]{\lambda_A[A]\subseteq U\land(\forall S,R\in U\;\kappa_A(S,R)\in U)}{%
    \rA}
  \proofstepwidestar[19,20]{U=\BracketTreeSet A}{%
    \FormulaRefAuto{TreeCodeMinimality}[thm]{19,\rAEa{20},\rAEb{20}}}
  \proofstepwidestar[]{\begin{aligned}[t]&\forall U\subseteq\BracketTreeSet A\;\Bigl((\lambda_A[A]\subseteq U\land{}\\&\quad(\forall S,R\in U\;\kappa_A(S,R)\in U))\rightarrow U=\BracketTreeSet A\Bigr)\end{aligned}}{%
    \rUI{\rRI{19,\rRI{20,21}}}}
  \proofstepwidestar[]{\mathsf{BaumStr}_A(\BracketTreeSet A,\lambda_A,\kappa_A)}{%
    \FormulaRefAuto{BinaryTreeStructureDef}[def]{1,6,13,18,22}}
\end{tabproofwide}

\subsection{Erste Folgerungen aus den Baumaxiomen}

Von hier an werden Induktion und Zerlegung aus den Strukturaxiomen
hergeleitet. Die einzelnen Codepositionen und Baumstufen gehen in diese
Beweise nicht mehr ein. Die gewonnenen Sätze gelten deshalb für jede
Menge mit Konstruktoren, welche B0 bis B4 erfüllt.

\noindent\begin{minipage}{\linewidth}
\FormulaThmDeltaK[Blätter liegen im Strukturträger]{%
\mathsf{BaumStr}_A(T,\ell,k)\dsep a\in A\vdash\ell(a)\in T
}{BinaryTreeStructureLeafTyping}{%
\DeltaRow{Mengen}{A\dsep T\dsep a}\DeltaRow{Funktionen}{\ell\dsep k}}
\end{minipage}\par
\begin{tabproofwide}
  \proofstepwidestar[1]{\mathsf{BaumStr}_A(T,\ell,k)}{%
    \rA}
  \proofstepwidestar[2]{a\in A}{%
    \rA}
  \proofstepwidestar[1]{\ell\colon A\to T}{%
    \rAEa{\FormulaRefAuto{BinaryTreeStructureTypingAxiom}[ax]{1}}}
  \proofstepwidestar[1,2]{\ell(a)\in T}{%
    \FormulaRefAuto{F\colon A\to B\dsep x\in A\vdash F(x)\in B}[thm]{3,2}}
\end{tabproofwide}

\noindent\begin{minipage}{\linewidth}
\FormulaThmDeltaK[Knoten liegen im Strukturträger]{%
\mathsf{BaumStr}_A(T,\ell,k)\dsep x,y\in T\vdash k(x,y)\in T
}{BinaryTreeStructureNodeTyping}{%
\DeltaRow{Mengen}{A\dsep T\dsep x\dsep y}\DeltaRow{Funktionen}{\ell\dsep k}}
\end{minipage}\par
\begin{tabproofwide}
  \proofstepwidestar[1]{\mathsf{BaumStr}_A(T,\ell,k)}{%
    \rA}
  \proofstepwidestar[2]{x,y\in T}{%
    \rA}
  \proofstepwidestar[1]{k\colon T\times T\to T}{%
    \rAEb{\FormulaRefAuto{BinaryTreeStructureTypingAxiom}[ax]{1}}}
  \proofstepwidestar[2]{(x,y)\in T\times T}{%
    \FormulaRefAuto{a\in A\dsep b\in B\vdash(a,b)\in A\times B}[thm]{\rAEa{2},\rAEb{2}}}
  \proofstepwidestar[1,2]{k(x,y)\in T}{%
    \FormulaRefAuto{F\colon A\to B\dsep x\in A\vdash F(x)\in B}[thm]{3,4}}
\end{tabproofwide}

\noindent\begin{minipage}{\linewidth}
\FormulaThmDeltaK[Eindeutigkeit des geordneten Kinderpaares]{%
\mathsf{BaumStr}_A(T,\ell,k)\dsep p,q\in T\times T\dsep k(p)=k(q)\vdash p=q
}{BinaryTreeStructurePairInjective}{%
\DeltaRow{Mengen}{A\dsep T\dsep p\dsep q\dsep x\dsep y\dsep x'\dsep y'}\DeltaRow{Funktionen}{\ell\dsep k}}
\end{minipage}\par
\begin{tabproofwide}
  \proofstepwidestar[1]{\mathsf{BaumStr}_A(T,\ell,k)}{%
    \rA}
  \proofstepwidestar[2]{p,q\in T\times T}{%
    \rA}
  \proofstepwidestar[3]{k(p)=k(q)}{%
    \rA}
  \proofstepwidestar[2]{\pi_1(p)\in T}{%
    \FormulaRefAuto{z\in A\times B\vdash\pi_1(z)\in A}[thm]{\rAEa{2}}}
  \proofstepwidestar[2]{\pi_2(p)\in T}{%
    \FormulaRefAuto{z\in A\times B\vdash\pi_2(z)\in B}[thm]{\rAEa{2}}}
  \proofstepwidestar[2]{\pi_1(q)\in T}{%
    \FormulaRefAuto{z\in A\times B\vdash\pi_1(z)\in A}[thm]{\rAEb{2}}}
  \proofstepwidestar[2]{\pi_2(q)\in T}{%
    \FormulaRefAuto{z\in A\times B\vdash\pi_2(z)\in B}[thm]{\rAEb{2}}}
  \proofstepwidestar[2]{p=(\pi_1(p),\pi_2(p))}{%
    \FormulaRefAuto{z\in A\times B\vdash z=\bigl(\pi_1(z),\pi_2(z)\bigr)}[thm]{\rAEa{2}}}
  \proofstepwidestar[2]{q=(\pi_1(q),\pi_2(q))}{%
    \FormulaRefAuto{z\in A\times B\vdash z=\bigl(\pi_1(z),\pi_2(z)\bigr)}[thm]{\rAEb{2}}}
  \proofstepwidestar[2,3]{k(\pi_1(p),\pi_2(p))=k(\pi_1(q),\pi_2(q))}{%
    \rIE{8,\rIE{9,3}}}
  \proofstepwidestar[1,2,3]{\pi_1(p)=\pi_1(q)\land\pi_2(p)=\pi_2(q)}{%
    \FormulaRefAuto{BinaryTreeStructureNodeAxiom}[ax]{1,\rAI{4,\rAI{5,\rAI{6,7}}},10}}
  \proofstepwidestar[1,2,3]{(\pi_1(p),\pi_2(p))=(\pi_1(q),\pi_2(q))}{%
    \FormulaRefAuto{(a,b)=(c,d)\eqvdash a=c\land b=d}[thm]{11}}
  \proofstepwidestar[1,2,3]{p=q}{%
    \rIE{\FormulaRefAuto{a=b\vdash b=a}[thm]{8},\rIE{\FormulaRefAuto{a=b\vdash b=a}[thm]{9},12}}}
\end{tabproofwide}

\noindent\begin{minipage}{\linewidth}
\FormulaThmDeltaK[Induktion in einer Baumstruktur]{%
  \begin{aligned}[t]
    &\mathsf{BaumStr}_A(T,\ell,k)\dsep\forall a\in A\;P(\ell(a))\dsep{}\\[-2pt]
    &\forall x,y\in T\;(P(x)\land P(y)\rightarrow P(k(x,y)))
      \vdash\forall t\in T\;P(t)
  \end{aligned}
}{BinaryTreeStructureInduction}{%
  \DeltaRow{Mengen}{A\dsep T\dsep U\dsep a\dsep x\dsep y\dsep t\dsep z}
  \DeltaRow{Funktionen}{\ell\dsep k}\DeltaRow{Prädikat}{P}
  \DeltaRow{Abkürzung}{U\coloneqq\{t\in T\mid P(t)\}}}
\end{minipage}\par
\begin{tabproofwide}
  \proofstepwidestar[1]{\mathsf{BaumStr}_A(T,\ell,k)}{%
    \rA}
  \proofstepwidestar[2]{\forall a\in A\;P(\ell(a))}{%
    \rA}
  \proofstepwidestar[3]{\forall x,y\in T\;(P(x)\land P(y)\rightarrow P(k(x,y)))}{%
    \rA}
  \proofstepwidestar[]{U\subseteq T}{%
    \FormulaRefAuto{\{x\in A\mid P(x)\}\subseteq A}[thm]}
  \proofstepwidestar[1]{\ell\colon A\to T}{%
    \rAEa{\FormulaRefAuto{BinaryTreeStructureTypingAxiom}[ax]{1}}}
  \proofstepwidestar[6]{z\in\ell[A]}{%
    \rA}
  \proofstepwidestar[1,6]{\exists a\in A\;z=\ell(a)}{%
    \FormulaRefAuto{F\colon A\to B\dsep y\in F[A]\vdash\exists x\in A\;y=F(x)}[thm]{5,6}}
  \proofstepwidestar[8]{a\in A\land z=\ell(a)}{%
    \rA}
  \proofstepwidestar[1,8]{\ell(a)\in T}{%
    \FormulaRefAuto{BinaryTreeStructureLeafTyping}[thm]{1,\rAEa{8}}}
  \proofstepwidestar[2,8]{P(\ell(a))}{%
    \rRE{\rUE{2},\rAEa{8}}}
  \proofstepwidestar[1,2,8]{\ell(a)\in U}{%
    \FormulaRefAuto{x\in A\dsep P(x)\vdash x\in\{x\in A\mid P(x)\}}[thm]{9,10}}
  \proofstepwidestar[1,2,8]{z\in U}{%
    \rIE{\FormulaRefAuto{a=b\vdash b=a}[thm]{\rAEb{8}},11}}
  \proofstepwidestar[1,2,6]{z\in U}{%
    \rEE{7,8,12}}
  \proofstepwidestar[1,2]{\ell[A]\subseteq U}{%
    \FormulaRefAuto{SubsetIntroductionRule}[thm]{[6]\vdots 13}}
  \proofstepwidestar[15]{x\in U}{%
    \rA}
  \proofstepwidestar[16]{y\in U}{%
    \rA}
  \proofstepwidestar[15,16]{x\in T\land P(x)}{%
    \FormulaRefAuto{x\in\{u\in A\mid P(u)\}\eqvdash x\in A\land P(x)}[thm]{15}}
  \proofstepwidestar[15,16]{y\in T\land P(y)}{%
    \FormulaRefAuto{x\in\{u\in A\mid P(u)\}\eqvdash x\in A\land P(x)}[thm]{16}}
  \proofstepwidestar[1,15,16]{k(x,y)\in T}{%
    \FormulaRefAuto{BinaryTreeStructureNodeTyping}[thm]{1,\rAI{\rAEa{17},\rAEa{18}}}}
  \proofstepwidestar[3,15,16]{P(k(x,y))}{%
    \rRE{\rRE{\rUE{\rRE{\rUE{3},\rAEa{17}}},\rAEa{18}},\rAI{\rAEb{17},\rAEb{18}}}}
  \proofstepwidestar[1,3,15,16]{k(x,y)\in U}{%
    \FormulaRefAuto{x\in A\dsep P(x)\vdash x\in\{x\in A\mid P(x)\}}[thm]{19,20}}
  \proofstepwidestar[1,3]{\forall x,y\in U\;k(x,y)\in U}{%
    \rUI{\rRI{15,\rUI{\rRI{16,21}}}}}
  \proofstepwidestar[1,2,3]{U=T}{%
    \FormulaRefAuto{BinaryTreeStructureMinimalityAxiom}[ax]{1,4,14,22}}
  \proofstepwidestar[24]{t\in T}{%
    \rA}
  \proofstepwidestar[1,2,3,24]{t\in U}{%
    \rIE{\FormulaRefAuto{a=b\vdash b=a}[thm]{23},24}}
  \proofstepwidestar[1,2,3,24]{P(t)}{%
    \rAEb{\FormulaRefAuto{x\in\{u\in A\mid P(u)\}\eqvdash x\in A\land P(x)}[thm]{25}}}
  \proofstepwidestar[1,2,3]{\forall t\in T\;P(t)}{%
    \rUI{\rRI{24,26}}}
\end{tabproofwide}

Auch hier folgt die Minimalität umgekehrt durch Einsetzen von
\(P(t)\coloneqq t\in U\). Insbesondere schließt das Induktionsaxiom
zusätzliche Elemente ohne endliche Konstruktion aus.

\noindent\begin{minipage}{\linewidth}
\FormulaThmDeltaK[Eindeutige Zerlegung in einer Baumstruktur]{%
  \begin{aligned}[t]
    &\mathsf{BaumStr}_A(T,\ell,k)\dsep t\in T\vdash{}\\[-2pt]
    &\bigl(\exists!a\in A\;t=\ell(a)\bigr)
      \lor\bigl(\exists!p\in T\times T\;t=k(p)\bigr)
  \end{aligned}
}{BinaryTreeStructureDecomposition}{%
  \DeltaRow{Mengen}{A\dsep T\dsep t\dsep a\dsep b\dsep c\dsep p\dsep q\dsep x\dsep y}
  \DeltaRow{Funktionen}{\ell\dsep k}
  \DeltaRow{Prädikatsabkürzung}{\begin{aligned}[t]
    Q(t)\coloneqq{}&(\exists!a\in A\;t=\ell(a))\lor{}\\[-2pt]
    &(\exists!p\in T\times T\;t=k(p))
  \end{aligned}}}
\end{minipage}\par
\begin{tabproofwide}
  \proofstepwidestar[1]{\mathsf{BaumStr}_A(T,\ell,k)}{%
    \rA}
  \proofstepwidestar[2]{t\in T}{%
    \rA}
  \proofstepwidestar[3]{a\in A}{%
    \rA}
  \proofstepwidestar[3]{\exists b\in A\;\ell(a)=\ell(b)}{%
    \rEI{\rAI{3,\rII}}}
  \proofstepwidestar[5]{b\in A\land\ell(a)=\ell(b)}{%
    \rA}
  \proofstepwidestar[6]{c\in A\land\ell(a)=\ell(c)}{%
    \rA}
  \proofstepwidestar[5,6]{\ell(b)=\ell(c)}{%
    \rIE{\rAEb{5},\rAEb{6}}}
  \proofstepwidestar[1,5,6]{b=c}{%
    \FormulaRefAuto{BinaryTreeStructureLeafAxiom}[ax]{1,\rAEa{5},\rAEa{6},7}}
  \proofstepwidestar[1,3]{\exists!b\in A\;\ell(a)=\ell(b)}{%
    \UEI{4,5,6,8}}
  \proofstepwidestar[1,3]{Q(\ell(a))}{%
    \rOIa{9}}
  \proofstepwidestar[1]{\forall a\in A\;Q(\ell(a))}{%
    \rUI{\rRI{3,10}}}
  \proofstepwidestar[12]{x\in T}{%
    \rA}
  \proofstepwidestar[13]{y\in T}{%
    \rA}
  \proofstepwidestar[12,13]{(x,y)\in T\times T}{%
    \FormulaRefAuto{a\in A\dsep b\in B\vdash(a,b)\in A\times B}[thm]{12,13}}
  \proofstepwidestar[12,13]{\exists p\in T\times T\;k(x,y)=k(p)}{%
    \rEI{\rAI{14,\rII}}}
  \proofstepwidestar[16]{p\in T\times T\land k(x,y)=k(p)}{%
    \rA}
  \proofstepwidestar[17]{q\in T\times T\land k(x,y)=k(q)}{%
    \rA}
  \proofstepwidestar[16,17]{k(p)=k(q)}{%
    \rIE{\rAEb{16},\rAEb{17}}}
  \proofstepwidestar[1,16,17]{p=q}{%
    \FormulaRefAuto{BinaryTreeStructurePairInjective}[thm]{1,\rAI{\rAEa{16},\rAEa{17}},18}}
  \proofstepwidestar[1,12,13]{\exists!p\in T\times T\;k(x,y)=k(p)}{%
    \UEI{15,16,17,19}}
  \proofstepwidestar[1,12,13]{Q(k(x,y))}{%
    \rOIb{20}}
  \proofstepwidestar[22]{Q(x)\land Q(y)}{%
    \rA}
  \proofstepwidestar[1,12,13]{Q(x)\land Q(y)\rightarrow Q(k(x,y))}{%
    \rRI{22,21}}
  \proofstepwidestar[1]{\forall x,y\in T\;(Q(x)\land Q(y)\rightarrow Q(k(x,y)))}{%
    \rUI{\rRI{12,\rUI{\rRI{13,23}}}}}
  \proofstepwidestar[1]{\forall t\in T\;Q(t)}{%
    \FormulaRefAuto{BinaryTreeStructureInduction}[thm]{1,11,24}}
  \proofstepwidestar[1,2]{(\exists!a\in A\;t=\ell(a))\lor(\exists!p\in T\times T\;t=k(p))}{%
    \rRE{\rUE{25},2}}
\end{tabproofwide}

Die beiden Fälle sind nach
\FormulaRefAuto{BinaryTreeStructureDisjointnessAxiom}[ax] disjunkt.


\subsection{Rekursion aus den Baumaxiomen}

Die rekursiven Vorschriften werden zunächst als Abschlussbedingungen
für Teilmengen von \(T\times X\) formuliert. Die Existenz ihres kleinsten
abgeschlossenen Graphen folgt aus Aussonderung. Die nachfolgenden
Filterbeweise zeigen getrennt, dass jeder Baum genau einen Wert besitzt.

\noindent\begin{minipage}{\linewidth}
\FormulaDefDeltaK[Rekursionsgleichungen einer Baumstruktur]{%
  \begin{aligned}[t]
    &\mathsf{TRec}_A(F;T,\ell,k;X,f,g)\coloneqq{}\\[-2pt]
    &\quad F\colon T\to X\land
       \bigl(\forall a\in A\;F(\ell(a))=f(a)\bigr)\land{}\\[-2pt]
    &\quad\forall x,y\in T\;F(k(x,y))=g(F(x),F(y))
  \end{aligned}
}{BinaryTreeRecursionSpecificationDef}{%
  \DeltaRow{Mengen}{A\dsep T\dsep X\dsep a\dsep x\dsep y}
  \DeltaRow{Funktionen}{F\dsep\ell\dsep k\dsep f\dsep g}
  \DeltaRow{\textbf{Neues Symbol}}{}
  \DeltaPrem{Rekursionsgleichungen}{\mathsf{TRec}_A(F;T,\ell,k;X,f,g)}}
\end{minipage}\par


\noindent\begin{minipage}{\linewidth}
\FormulaDefDeltaK[Abgeschlossene Relationsgraphen für die Baumrekursion]{%
 \begin{aligned}[t]
 \mathsf{BCl}_A(H;T,\ell,k;X,f,g)\coloneqq{}&
 H\subseteq T\times X\land{}\\[-2pt]
 &(\forall a\in A\;(\ell(a),f(a))\in H)\land{}\\[-2pt]
 &\forall u,v\in T\,\forall x,y\in X\\[-2pt]
 &\quad((u,x)\in H\land(v,y)\in H\\[-2pt]
 &\qquad\rightarrow(k(u,v),g(x,y))\in H).
 \end{aligned}
}{BinaryTreeClosedGraphDef}{%
 \DeltaRow{Mengen}{A\dsep T\dsep X\dsep H\dsep a\dsep u\dsep v\dsep x\dsep y}
 \DeltaRow{Funktionen}{\ell\dsep k\dsep f\dsep g}
 \DeltaRow{\textbf{Neues Symbol}}{}
 \DeltaPrem{Abgeschlossener Graph}{\mathsf{BCl}_A(H;T,\ell,k;X,f,g)}}
\end{minipage}\par

In den folgenden Hilfssätzen stehen \(C\) und \(M\) für die jeweils
im Kontext angegebenen Prädikate. Die Mengenabkürzungen stehen ebenfalls
im Kontext; ihre Aussonderungsterme werden durch
\FormulaRefAuto{\{x\in A\mid P(x)\}\coloneq
 \iota B(\forall x(x\in B\leftrightarrow(x\in A\land P(x))))}[def]
aufgrund des Aussonderungsaxioms definiert.

\noindent\begin{minipage}{\linewidth}
\FormulaThmDeltaK[Träger eines abgeschlossenen Baumgraphen]{%
C(H)\vdash H\subseteq T\times X
}{BinaryTreeClosedGraphSubset}{%
\DeltaRow{Mengen}{A\dsep T\dsep X\dsep G\dsep H\dsep t\dsep u\dsep v\dsep x\dsep y\dsep a}
\DeltaRow{Funktionen}{\ell\dsep k\dsep f\dsep g}
\DeltaRow{Prädikatsabkürzungen}{\begin{aligned}[t]
 C(H)&\coloneqq\mathsf{BCl}_A(H;T,\ell,k;X,f,g),\\[-2pt]
 M(G)&\coloneqq C(G)\land\forall H(C(H)\rightarrow G\subseteq H)
\end{aligned}}}
\end{minipage}\par
\begin{tabproofwide}
  \proofstepwidestar[1]{C(H)}{%
    \rA}
  \proofstepwidestar[1]{H\subseteq T\times X}{%
    \rAEa{\FormulaRefAuto{BinaryTreeClosedGraphDef}[def]{1}}}
\end{tabproofwide}
\noindent\begin{minipage}{\linewidth}
\FormulaThmDeltaK[Blattregel eines abgeschlossenen Baumgraphen]{%
C(H)\dsep a\in A\vdash(\ell(a),f(a))\in H
}{BinaryTreeClosedGraphLeaf}{%
\DeltaRow{Mengen}{A\dsep T\dsep X\dsep G\dsep H\dsep t\dsep u\dsep v\dsep x\dsep y\dsep a}
\DeltaRow{Funktionen}{\ell\dsep k\dsep f\dsep g}
\DeltaRow{Prädikatsabkürzungen}{\begin{aligned}[t]
 C(H)&\coloneqq\mathsf{BCl}_A(H;T,\ell,k;X,f,g),\\[-2pt]
 M(G)&\coloneqq C(G)\land\forall H(C(H)\rightarrow G\subseteq H)
\end{aligned}}}
\end{minipage}\par
\begin{tabproofwide}
  \proofstepwidestar[1]{C(H)}{%
    \rA}
  \proofstepwidestar[2]{a\in A}{%
    \rA}
  \proofstepwidestar[1]{\forall a\in A\;(\ell(a),f(a))\in H}{%
    \rAEn{\FormulaRefAuto{BinaryTreeClosedGraphDef}[def]{1}}}
  \proofstepwidestar[1,2]{(\ell(a),f(a))\in H}{%
    \rRE{\rUE{3},2}}
\end{tabproofwide}
\noindent\begin{minipage}{\linewidth}
\FormulaThmDeltaK[Knotenregel eines abgeschlossenen Baumgraphen]{%
\begin{aligned}[t]&C(H)\dsep u,v\in T\dsep x,y\in X\dsep{}\\&(u,x)\in H\dsep(v,y)\in H\vdash(k(u,v),g(x,y))\in H\end{aligned}
}{BinaryTreeClosedGraphNode}{%
\DeltaRow{Mengen}{A\dsep T\dsep X\dsep G\dsep H\dsep t\dsep u\dsep v\dsep x\dsep y\dsep a}
\DeltaRow{Funktionen}{\ell\dsep k\dsep f\dsep g}
\DeltaRow{Prädikatsabkürzungen}{\begin{aligned}[t]
 C(H)&\coloneqq\mathsf{BCl}_A(H;T,\ell,k;X,f,g),\\[-2pt]
 M(G)&\coloneqq C(G)\land\forall H(C(H)\rightarrow G\subseteq H)
\end{aligned}}}
\end{minipage}\par
\begin{tabproofwide}
  \proofstepwidestar[1]{C(H)}{%
    \rA}
  \proofstepwidestar[2]{u,v\in T}{%
    \rA}
  \proofstepwidestar[3]{x,y\in X}{%
    \rA}
  \proofstepwidestar[4]{(u,x)\in H}{%
    \rA}
  \proofstepwidestar[5]{(v,y)\in H}{%
    \rA}
  \proofstepwidestar[1]{\begin{aligned}[t]&\forall u,v\in T\,\forall x,y\in X\\&((u,x)\in H\land(v,y)\in H\\&\quad\rightarrow(k(u,v),g(x,y))\in H)\end{aligned}}{%
    \rAEn{\FormulaRefAuto{BinaryTreeClosedGraphDef}[def]{1}}}
  \proofstepwidestar[1,2,3,4,5]{(k(u,v),g(x,y))\in H}{%
    \rRE{\rRE{\rUE{\rRE{\rUE{\rRE{\rUE{\rRE{\rUE{6},\rAEa{2}}},\rAEb{2}}},\rAEa{3}}},\rAEb{3}},\rAI{4,5}}}
\end{tabproofwide}
\noindent\begin{minipage}{\linewidth}
\FormulaThmDeltaK[Minimalität des abgeschlossenen Baumgraphen]{%
M(G)\dsep C(H)\vdash G\subseteq H
}{BinaryTreeMinimalGraphSubset}{%
\DeltaRow{Mengen}{A\dsep T\dsep X\dsep G\dsep H\dsep t\dsep u\dsep v\dsep x\dsep y\dsep a}
\DeltaRow{Funktionen}{\ell\dsep k\dsep f\dsep g}
\DeltaRow{Prädikatsabkürzungen}{\begin{aligned}[t]
 C(H)&\coloneqq\mathsf{BCl}_A(H;T,\ell,k;X,f,g),\\[-2pt]
 M(G)&\coloneqq C(G)\land\forall H(C(H)\rightarrow G\subseteq H)
\end{aligned}}}
\end{minipage}\par
\begin{tabproofwide}
  \proofstepwidestar[1]{M(G)}{%
    \rA}
  \proofstepwidestar[2]{C(H)}{%
    \rA}
  \proofstepwidestar[1]{\forall H\;(C(H)\rightarrow G\subseteq H)}{%
    \rAEb{1}}
  \proofstepwidestar[1,2]{G\subseteq H}{%
    \rRE{\rUE{3},2}}
\end{tabproofwide}

\noindent\begin{minipage}{\linewidth}
\FormulaThmDeltaK[Existenz des kleinsten abgeschlossenen Baumgraphen]{%
\begin{aligned}[t]&\mathsf{BaumStr}_A(T,\ell,k)\dsep f\colon A\to X\dsep g\colon X\times X\to X\\&\vdash\exists G\;M(G)\end{aligned}
}{BinaryTreeLeastClosedGraphExists}{%
\DeltaRow{Mengen}{A\dsep T\dsep X\dsep G_0\dsep H\dsep t\dsep u\dsep v\dsep x\dsep y\dsep a}
\DeltaRow{Funktionen}{\ell\dsep k\dsep f\dsep g}
\DeltaRow{Prädikatsabkürzungen}{\begin{aligned}[t]
 C(H)&\coloneqq\mathsf{BCl}_A(H;T,\ell,k;X,f,g),\\[-2pt]
 M(G)&\coloneqq C(G)\land\forall H(C(H)\rightarrow G\subseteq H)
\end{aligned}}
\DeltaRow{Weitere Mengen}{B\dsep p\dsep K\dsep\mathcal C}
\DeltaRow{Mengenabkürzungen}{\begin{aligned}[t]
 B&\coloneqq T\times X,\\
 \mathcal C&\coloneqq\{H\in\powerset B\mid C(H)\},\\
 G_0&\coloneqq\{p\in B\mid\forall H\in\mathcal C\;p\in H\}
\end{aligned}}}
\end{minipage}\par
\begin{tabproofwide}
  \proofstepwidestar[1]{\mathsf{BaumStr}_A(T,\ell,k)}{%
    \rA}
  \proofstepwidestar[2]{f\colon A\to X}{%
    \rA}
  \proofstepwidestar[3]{g\colon X\times X\to X}{%
    \rA}
  \proofstepwidestar[]{B\subseteq B}{%
    \FormulaRefAuto{A\subseteq A}[thm]}
  \proofstepwidestar[]{B\in\powerset B}{%
    \FormulaRefAuto{x\in\powerset(A)\eqvdash x\subseteq A}[thm]{4}}
  \proofstepwidestar[6]{a\in A}{%
    \rA}
  \proofstepwidestar[1,6]{\ell(a)\in T}{%
    \FormulaRefAuto{BinaryTreeStructureLeafTyping}[thm]{1,6}}
  \proofstepwidestar[2,6]{f(a)\in X}{%
    \FormulaRefAuto{F\colon A\to B\dsep x\in A\vdash F(x)\in B}[thm]{2,6}}
  \proofstepwidestar[1,2,6]{(\ell(a),f(a))\in B}{%
    \FormulaRefAuto{a\in A\dsep b\in B\vdash(a,b)\in A\times B}[thm]{7,8}}
  \proofstepwidestar[1,2]{\forall a\in A\;(\ell(a),f(a))\in B}{%
    \rUI{\rRI{6,9}}}
  \proofstepwidestar[11]{u\in T}{%
    \rA}
  \proofstepwidestar[12]{v\in T}{%
    \rA}
  \proofstepwidestar[13]{x\in X}{%
    \rA}
  \proofstepwidestar[14]{y\in X}{%
    \rA}
  \proofstepwidestar[15]{(u,x)\in B\land(v,y)\in B}{%
    \rA}
  \proofstepwidestar[11,12,13,14,15]{\begin{aligned}[t]&u,v\in T\land x,y\in X\land{}\\&(u,x)\in B\land(v,y)\in B\end{aligned}}{%
    \rAI{\rAI{11,12},\rAI{\rAI{13,14},15}}}
  \proofstepwidestar[1,11,12,13,14,15]{k(u,v)\in T}{%
    \FormulaRefAuto{BinaryTreeStructureNodeTyping}[thm]{1,\rAEn{16}}}
  \proofstepwidestar[11,12,13,14,15]{(x,y)\in X\times X}{%
    \FormulaRefAuto{a\in A\dsep b\in B\vdash(a,b)\in A\times B}[thm]{\rAEn{16},\rAEn{16}}}
  \proofstepwidestar[3,11,12,13,14,15]{g(x,y)\in X}{%
    \FormulaRefAuto{F\colon A\to B\dsep x\in A\vdash F(x)\in B}[thm]{3,18}}
  \proofstepwidestar[1,3,11,12,13,14,15]{(k(u,v),g(x,y))\in B}{%
    \FormulaRefAuto{a\in A\dsep b\in B\vdash(a,b)\in A\times B}[thm]{17,19}}
  \proofstepwidestar[1,3]{\begin{aligned}[t]&\forall u,v\in T\,\forall x,y\in X\\&((u,x)\in B\land(v,y)\in B\rightarrow(k(u,v),g(x,y))\in B)\end{aligned}}{%
    \rUI{\rRI{11,\rUI{\rRI{12,\rUI{\rRI{13,\rUI{\rRI{14,\rRI{15,20}}}}}}}}}}
  \proofstepwidestar[1,2,3]{C(B)}{%
    \FormulaRefAuto{BinaryTreeClosedGraphDef}[def]{\rAI{4,\rAI{10,21}}}}
  \proofstepwidestar[1,2,3]{B\in\mathcal C}{%
    \FormulaRefAuto{x\in A\dsep P(x)\vdash x\in\{x\in A\mid P(x)\}}[thm]{5,22}}
  \proofstepwidestar[]{G_0\subseteq B}{%
    \FormulaRefAuto{\{x\in A\mid P(x)\}\subseteq A}[thm]}
  \proofstepwidestar[25]{H\in\mathcal C}{%
    \rA}
  \proofstepwidestar[26]{p\in G_0}{%
    \rA}
  \proofstepwidestar[26]{\forall H\in\mathcal C\;p\in H}{%
    \rAEb{\FormulaRefAuto{x\in\{u\in A\mid P(u)\}\eqvdash x\in A\land P(x)}[thm]{26}}}
  \proofstepwidestar[25,26]{p\in H}{%
    \rRE{\rUE{27},25}}
  \proofstepwidestar[25]{G_0\subseteq H}{%
    \FormulaRefAuto{SubsetIntroductionRule}[thm]{[26]\vdots 28}}
  \proofstepwidestar[]{\forall H\in\mathcal C\;G_0\subseteq H}{%
    \rUI{\rRI{25,29}}}
  \proofstepwidestar[31]{a\in A}{%
    \rA}
  \proofstepwidestar[32]{H\in\mathcal C}{%
    \rA}
  \proofstepwidestar[32]{C(H)}{%
    \rAEb{\FormulaRefAuto{x\in\{u\in A\mid P(u)\}\eqvdash x\in A\land P(x)}[thm]{32}}}
  \proofstepwidestar[32]{\forall a\in A\;(\ell(a),f(a))\in H}{%
    \rAEn{\FormulaRefAuto{BinaryTreeClosedGraphDef}[def]{33}}}
  \proofstepwidestar[31,32]{(\ell(a),f(a))\in H}{%
    \rRE{\rUE{34},31}}
  \proofstepwidestar[31]{\forall H\in\mathcal C\;(\ell(a),f(a))\in H}{%
    \rUI{\rRI{32,35}}}
  \proofstepwidestar[1,2,31]{(\ell(a),f(a))\in B}{%
    \rRE{\rUE{10},31}}
  \proofstepwidestar[1,2,31]{(\ell(a),f(a))\in G_0}{%
    \FormulaRefAuto{x\in A\dsep P(x)\vdash x\in\{x\in A\mid P(x)\}}[thm]{37,36}}
  \proofstepwidestar[1,2]{\forall a\in A\;(\ell(a),f(a))\in G_0}{%
    \rUI{\rRI{31,38}}}
  \proofstepwidestar[40]{u\in T}{%
    \rA}
  \proofstepwidestar[41]{v\in T}{%
    \rA}
  \proofstepwidestar[42]{x\in X}{%
    \rA}
  \proofstepwidestar[43]{y\in X}{%
    \rA}
  \proofstepwidestar[44]{(u,x)\in G_0\land(v,y)\in G_0}{%
    \rA}
  \proofstepwidestar[40,41,42,43,44]{\begin{aligned}[t]&u,v\in T\land x,y\in X\land{}\\&(u,x)\in G_0\land(v,y)\in G_0\end{aligned}}{%
    \rAI{\rAI{40,41},\rAI{\rAI{42,43},44}}}
  \proofstepwidestar[46]{H\in\mathcal C}{%
    \rA}
  \proofstepwidestar[46]{C(H)}{%
    \rAEb{\FormulaRefAuto{x\in\{u\in A\mid P(u)\}\eqvdash x\in A\land P(x)}[thm]{46}}}
  \proofstepwidestar[46]{G_0\subseteq H}{%
    \rRE{\rUE{30},46}}
  \proofstepwidestar[40,46,41,42,43,44]{(u,x)\in H}{%
    \FormulaRefAuto{A\subseteq B\dsep x\in A\vdash x\in B}[thm]{48,\rAEn{45}}}
  \proofstepwidestar[40,46,41,42,43,44]{(v,y)\in H}{%
    \FormulaRefAuto{A\subseteq B\dsep x\in A\vdash x\in B}[thm]{48,\rAEn{45}}}
  \proofstepwidestar[46]{\begin{aligned}[t]&\forall u,v\in T\,\forall x,y\in X\\&((u,x)\in H\land(v,y)\in H\\&\quad\rightarrow(k(u,v),g(x,y))\in H)\end{aligned}}{%
    \rAEn{\FormulaRefAuto{BinaryTreeClosedGraphDef}[def]{47}}}
  \proofstepwidestar[40,46,41,42,43,44]{(k(u,v),g(x,y))\in H}{%
    \rRE{\rRE{\rUE{\rRE{\rUE{\rRE{\rUE{\rRE{\rUE{51},\rAEn{45}}},\rAEn{45}}},\rAEn{45}}},\rAEn{45}},\rAI{49,50}}}
  \proofstepwidestar[40,41,42,43,44]{\forall H\in\mathcal C\;(k(u,v),g(x,y))\in H}{%
    \rUI{\rRI{46,52}}}
  \proofstepwidestar[1,3,40,41,42,43,44]{(k(u,v),g(x,y))\in B}{%
    \FormulaRefAuto{a\in A\dsep b\in B\vdash(a,b)\in A\times B}[thm]{\FormulaRefAuto{BinaryTreeStructureNodeTyping}[thm]{1,\rAEn{45}},\FormulaRefAuto{F\colon A\to B\dsep x\in A\vdash F(x)\in B}[thm]{3,\FormulaRefAuto{a\in A\dsep b\in B\vdash(a,b)\in A\times B}[thm]{\rAEn{45},\rAEn{45}}}}}
  \proofstepwidestar[1,3,40,41,42,43,44]{(k(u,v),g(x,y))\in G_0}{%
    \FormulaRefAuto{x\in A\dsep P(x)\vdash x\in\{x\in A\mid P(x)\}}[thm]{54,53}}
  \proofstepwidestar[1,3]{\begin{aligned}[t]&\forall u,v\in T\,\forall x,y\in X\\&((u,x)\in G_0\land(v,y)\in G_0\\&\quad\rightarrow(k(u,v),g(x,y))\in G_0)\end{aligned}}{%
    \rUI{\rRI{40,\rUI{\rRI{41,\rUI{\rRI{42,\rUI{\rRI{43,\rRI{44,55}}}}}}}}}}
  \proofstepwidestar[1,2,3]{C(G_0)}{%
    \FormulaRefAuto{BinaryTreeClosedGraphDef}[def]{\rAI{24,\rAI{39,56}}}}
  \proofstepwidestar[58]{C(K)}{%
    \rA}
  \proofstepwidestar[58]{K\subseteq B}{%
    \rAEa{\FormulaRefAuto{BinaryTreeClosedGraphDef}[def]{58}}}
  \proofstepwidestar[58]{K\in\powerset B}{%
    \FormulaRefAuto{x\in\powerset(A)\eqvdash x\subseteq A}[thm]{59}}
  \proofstepwidestar[58]{K\in\mathcal C}{%
    \FormulaRefAuto{x\in A\dsep P(x)\vdash x\in\{x\in A\mid P(x)\}}[thm]{60,58}}
  \proofstepwidestar[58]{G_0\subseteq K}{%
    \rRE{\rUE{30},61}}
  \proofstepwidestar[]{\forall K\;(C(K)\rightarrow G_0\subseteq K)}{%
    \rUI{\rRI{58,62}}}
  \proofstepwidestar[1,2,3]{\exists G\;M(G)}{%
    \rEI{\rAI{57,63}}}
\end{tabproofwide}


\noindent\begin{minipage}{\linewidth}
\FormulaThmDeltaK[Ein zulässiger Faserfilter erzwingt einen einzigen Wert]{%
\begin{aligned}[t]&M(G)\dsep C(\Phi^{T,X}_{t,c}(G))\dsep c\in X\dsep(t,c)\in G\\&\vdash\exists!x\in X\;(t,x)\in G\end{aligned}
}{BinaryTreeRecursionFilteredFiberUnique}{%
\DeltaRow{Mengen}{A\dsep T\dsep X\dsep G\dsep H\dsep t\dsep u\dsep v\dsep x\dsep y\dsep a}
\DeltaRow{Funktionen}{\ell\dsep k\dsep f\dsep g}
\DeltaRow{Prädikatsabkürzungen}{\begin{aligned}[t]
 C(H)&\coloneqq\mathsf{BCl}_A(H;T,\ell,k;X,f,g),\\[-2pt]
 M(G)&\coloneqq C(G)\land\forall H(C(H)\rightarrow G\subseteq H)
\end{aligned}}\DeltaRow{Weitere Mengen}{c\dsep z}}
\end{minipage}\par
\begin{tabproofwide}
  \proofstepwidestar[1]{M(G)}{%
    \rA}
  \proofstepwidestar[2]{C(\Phi^{T,X}_{t,c}(G))}{%
    \rA}
  \proofstepwidestar[3]{c\in X}{%
    \rA}
  \proofstepwidestar[4]{(t,c)\in G}{%
    \rA}
  \proofstepwidestar[1]{G\subseteq T\times X}{%
    \FormulaRefAuto{BinaryTreeClosedGraphSubset}[thm]{\rAEa{1}}}
  \proofstepwidestar[1,2]{G\subseteq\Phi^{T,X}_{t,c}(G)}{%
    \FormulaRefAuto{BinaryTreeMinimalGraphSubset}[thm]{1,2}}
  \proofstepwidestar[7]{(t,z)\in G}{%
    \rA}
  \proofstepwidestar[1,2,7]{(t,z)\in\Phi^{T,X}_{t,c}(G)}{%
    \FormulaRefAuto{A\subseteq B\dsep x\in A\vdash x\in B}[thm]{6,7}}
  \proofstepwidestar[1]{\begin{aligned}[t]&(t,z)\in\Phi^{T,X}_{t,c}(G)\leftrightarrow{}\\&((t,z)\in G\land(t\ne t\lor z=c))\end{aligned}}{%
    \FormulaRefAuto{WordRecursionFilterMembership}[thm]{5}}
  \proofstepwidestar[1,2,7]{t\ne t\lor z=c}{%
    \rAEb{\rRE{\rLREa{9},8}}}
  \proofstepwidestar[]{\neg(t\ne t)}{%
    \FormulaRefAuto{A\eqvdash\neg\neg A}[thm]{\rII}}
  \proofstepwidestar[1,2,7]{z=c}{%
    \FormulaRefAuto{A\lor B\dsep\neg A\vdash B}[thm]{10,11}}
  \proofstepwidestar[13]{z=c}{%
    \rA}
  \proofstepwidestar[4,13]{(t,z)\in G}{%
    \rIE{\FormulaRefAuto{a=b\vdash b=a}[thm]{13},4}}
  \proofstepwidestar[1,2,4]{(t,z)\in G\leftrightarrow z=c}{%
    \rLRI{\rRI{7,12},\rRI{13,14}}}
  \proofstepwidestar[1,2,4]{\forall z\;((t,z)\in G\leftrightarrow z=c)}{%
    \rUI{15}}
  \proofstepwidestar[1,2,3,4]{\exists!x\in X\;(t,x)\in G}{%
    \FormulaRefAuto{B27UniqueFromSingletonPredicate}[thm]{3,16}}
\end{tabproofwide}

\noindent\begin{minipage}{\linewidth}
\FormulaThmDeltaK[Eindeutiger Wert am Blatt]{%
\begin{aligned}[t]&\mathsf{BaumStr}_A(T,\ell,k)\dsep f\colon A\to X\dsep M(G)\dsep a\in A\\&\vdash\exists!x\in X\;(\ell(a),x)\in G\end{aligned}
}{BinaryTreeRecursionLeafFiberUnique}{%
\DeltaRow{Mengen}{A\dsep T\dsep X\dsep G\dsep H\dsep t\dsep u\dsep v\dsep x\dsep y\dsep a}
\DeltaRow{Funktionen}{\ell\dsep k\dsep f\dsep g}
\DeltaRow{Prädikatsabkürzungen}{\begin{aligned}[t]
 C(H)&\coloneqq\mathsf{BCl}_A(H;T,\ell,k;X,f,g),\\[-2pt]
 M(G)&\coloneqq C(G)\land\forall H(C(H)\rightarrow G\subseteq H)
\end{aligned}}\DeltaRow{Weitere Menge}{b}\DeltaRow{Mengenabkürzung}{H\coloneqq\Phi^{T,X}_{\ell(a),f(a)}(G)}}
\end{minipage}\par
\begin{tabproofwide}
  \proofstepwidestar[1]{\mathsf{BaumStr}_A(T,\ell,k)}{%
    \rA}
  \proofstepwidestar[2]{f\colon A\to X}{%
    \rA}
  \proofstepwidestar[3]{M(G)}{%
    \rA}
  \proofstepwidestar[4]{a\in A}{%
    \rA}
  \proofstepwidestar[3]{C(G)}{%
    \rAEa{3}}
  \proofstepwidestar[3]{G\subseteq T\times X}{%
    \FormulaRefAuto{BinaryTreeClosedGraphSubset}[thm]{5}}
  \proofstepwidestar[]{H\subseteq G}{%
    \FormulaRefAuto{WordRecursionFilterSubset}[thm]}
  \proofstepwidestar[3]{H\subseteq T\times X}{%
    \FormulaRefAuto{A\subseteq B\dsep B\subseteq C\vdash A\subseteq C}[thm]{7,6}}
  \proofstepwidestar[9]{b\in A}{%
    \rA}
  \proofstepwidestar[3,9]{(\ell(b),f(b))\in G}{%
    \FormulaRefAuto{BinaryTreeClosedGraphLeaf}[thm]{5,9}}
  \proofstepwidestar[]{\ell(b)=\ell(a)\lor\ell(b)\ne\ell(a)}{%
    \FormulaRefAuto{A\lor\neg A}[thm]}
  \proofstepwidestar[12]{\ell(b)=\ell(a)}{%
    \rA}
  \proofstepwidestar[1,4,9,12]{b=a}{%
    \FormulaRefAuto{BinaryTreeStructureLeafAxiom}[ax]{1,9,4,12}}
  \proofstepwidestar[1,4,9,12]{f(b)=f(a)}{%
    \rIE{13,\rII}}
  \proofstepwidestar[1,4,9,12]{\ell(b)\ne\ell(a)\lor f(b)=f(a)}{%
    \rOIb{14}}
  \proofstepwidestar[16]{\ell(b)\ne\ell(a)}{%
    \rA}
  \proofstepwidestar[16]{\ell(b)\ne\ell(a)\lor f(b)=f(a)}{%
    \rOIa{16}}
  \proofstepwidestar[1,4,9]{\ell(b)\ne\ell(a)\lor f(b)=f(a)}{%
    \rOE{11,12,15,16,17}}
  \proofstepwidestar[1,3,4,9]{(\ell(b),f(b))\in H}{%
    \rRE{\rLREb{\FormulaRefAuto{WordRecursionFilterMembership}[thm]{6}},\rAI{10,18}}}
  \proofstepwidestar[1,3,4]{\forall b\in A\;(\ell(b),f(b))\in H}{%
    \rUI{\rRI{9,19}}}
  \proofstepwidestar[21]{u\in T}{%
    \rA}
  \proofstepwidestar[22]{v\in T}{%
    \rA}
  \proofstepwidestar[23]{x\in X}{%
    \rA}
  \proofstepwidestar[24]{y\in X}{%
    \rA}
  \proofstepwidestar[25]{(u,x)\in H\land(v,y)\in H}{%
    \rA}
  \proofstepwidestar[25]{(u,x)\in G}{%
    \FormulaRefAuto{A\subseteq B\dsep x\in A\vdash x\in B}[thm]{7,\rAEa{25}}}
  \proofstepwidestar[25]{(v,y)\in G}{%
    \FormulaRefAuto{A\subseteq B\dsep x\in A\vdash x\in B}[thm]{7,\rAEb{25}}}
  \proofstepwidestar[3,21,22,23,24,25]{(k(u,v),g(x,y))\in G}{%
    \FormulaRefAuto{BinaryTreeClosedGraphNode}[thm]{5,\rAI{21,22},\rAI{23,24},26,27}}
  \proofstepwidestar[1,4,21,22]{\ell(a)\ne k(u,v)}{%
    \FormulaRefAuto{BinaryTreeStructureDisjointnessAxiom}[ax]{1,4,\rAI{21,22}}}
  \proofstepwidestar[30]{k(u,v)=\ell(a)}{%
    \rA}
  \proofstepwidestar[30]{\ell(a)=k(u,v)}{%
    \FormulaRefAuto{a=b\vdash b=a}[thm]{30}}
  \proofstepwidestar[1,4,21,22,30]{\bot}{%
    \rBI{31,29}}
  \proofstepwidestar[1,4,21,22]{k(u,v)\ne\ell(a)}{%
    \rCI{30,32}}
  \proofstepwidestar[1,4,21,22]{k(u,v)\ne\ell(a)\lor g(x,y)=f(a)}{%
    \rOIa{33}}
  \proofstepwidestar[1,3,4,21,22,23,24,25]{(k(u,v),g(x,y))\in H}{%
    \rRE{\rLREb{\FormulaRefAuto{WordRecursionFilterMembership}[thm]{6}},\rAI{28,34}}}
  \proofstepwidestar[1,3,4]{\begin{aligned}[t]&\forall u,v\in T\,\forall x,y\in X\\&((u,x)\in H\land(v,y)\in H\\&\quad\rightarrow(k(u,v),g(x,y))\in H)\end{aligned}}{%
    \rUI{\rRI{21,\rUI{\rRI{22,\rUI{\rRI{23,\rUI{\rRI{24,\rRI{25,35}}}}}}}}}}
  \proofstepwidestar[1,3,4]{C(H)}{%
    \FormulaRefAuto{BinaryTreeClosedGraphDef}[def]{\rAI{8,\rAI{20,36}}}}
  \proofstepwidestar[2,4]{f(a)\in X}{%
    \FormulaRefAuto{F\colon A\to B\dsep x\in A\vdash F(x)\in B}[thm]{2,4}}
  \proofstepwidestar[3,4]{(\ell(a),f(a))\in G}{%
    \FormulaRefAuto{BinaryTreeClosedGraphLeaf}[thm]{5,4}}
  \proofstepwidestar[1,2,3,4]{\exists!x\in X\;(\ell(a),x)\in G}{%
    \FormulaRefAuto{BinaryTreeRecursionFilteredFiberUnique}[thm]{3,37,38,39}}
\end{tabproofwide}

\noindent\begin{minipage}{\linewidth}
\FormulaThmDeltaK[Eindeutige Teilwerte bestimmen einen eindeutigen Knotenwert]{%
\begin{aligned}[t]&\mathsf{BaumStr}_A(T,\ell,k)\dsep g\colon X\times X\to X\dsep M(G)\dsep u,v\in T\\&\dsep\exists!z\in X\;(u,z)\in G\dsep\exists!z\in X\;(v,z)\in G\\&\vdash\exists!z\in X\;(k(u,v),z)\in G\end{aligned}
}{BinaryTreeRecursionNodeFiberUnique}{%
\DeltaRow{Mengen}{A\dsep T\dsep X\dsep G\dsep H\dsep t\dsep u\dsep v\dsep x\dsep y\dsep a}
\DeltaRow{Funktionen}{\ell\dsep k\dsep f\dsep g}
\DeltaRow{Prädikatsabkürzungen}{\begin{aligned}[t]
 C(H)&\coloneqq\mathsf{BCl}_A(H;T,\ell,k;X,f,g),\\[-2pt]
 M(G)&\coloneqq C(G)\land\forall H(C(H)\rightarrow G\subseteq H)
\end{aligned}}\DeltaRow{Weitere Mengen}{p\dsep q\dsep z\dsep w}
\DeltaRow{Iota-Abkürzungen}{\begin{aligned}[t]
 x&\coloneqq\iota z\,(z\in X\land(u,z)\in G),\\[-2pt]
 y&\coloneqq\iota z\,(z\in X\land(v,z)\in G)
\end{aligned}}
\DeltaRow{Mengenabkürzung}{H\coloneqq\Phi^{T,X}_{k(u,v),g(x,y)}(G)}}
\end{minipage}\par
\begin{tabproofwide}
  \proofstepwidestar[1]{\mathsf{BaumStr}_A(T,\ell,k)}{%
    \rA}
  \proofstepwidestar[2]{g\colon X\times X\to X}{%
    \rA}
  \proofstepwidestar[3]{M(G)}{%
    \rA}
  \proofstepwidestar[4]{u\in T}{%
    \rA}
  \proofstepwidestar[5]{v\in T}{%
    \rA}
  \proofstepwidestar[6]{\exists!z\in X\;(u,z)\in G}{%
    \rA}
  \proofstepwidestar[7]{\exists!z\in X\;(v,z)\in G}{%
    \rA}
  \proofstepwidestar[6]{x\in X\land(u,x)\in G}{%
    \FormulaRefAuto{IotaSymbolDef}[def]{6}}
  \proofstepwidestar[7]{y\in X\land(v,y)\in G}{%
    \FormulaRefAuto{IotaSymbolDef}[def]{7}}
  \proofstepwidestar[3]{C(G)}{%
    \rAEa{3}}
  \proofstepwidestar[3]{G\subseteq T\times X}{%
    \FormulaRefAuto{BinaryTreeClosedGraphSubset}[thm]{10}}
  \proofstepwidestar[]{H\subseteq G}{%
    \FormulaRefAuto{WordRecursionFilterSubset}[thm]}
  \proofstepwidestar[3]{H\subseteq T\times X}{%
    \FormulaRefAuto{A\subseteq B\dsep B\subseteq C\vdash A\subseteq C}[thm]{12,11}}
  \proofstepwidestar[14]{a\in A}{%
    \rA}
  \proofstepwidestar[3,14]{(\ell(a),f(a))\in G}{%
    \FormulaRefAuto{BinaryTreeClosedGraphLeaf}[thm]{10,14}}
  \proofstepwidestar[1,4,5,14]{\ell(a)\ne k(u,v)}{%
    \FormulaRefAuto{BinaryTreeStructureDisjointnessAxiom}[ax]{1,14,\rAI{4,5}}}
  \proofstepwidestar[1,3,4,5,14]{(\ell(a),f(a))\in H}{%
    \rRE{\rLREb{\FormulaRefAuto{WordRecursionFilterMembership}[thm]{11}},\rAI{15,\rOIa{16}}}}
  \proofstepwidestar[1,3,4,5]{\forall a\in A\;(\ell(a),f(a))\in H}{%
    \rUI{\rRI{14,17}}}
  \proofstepwidestar[19]{p\in T}{%
    \rA}
  \proofstepwidestar[20]{q\in T}{%
    \rA}
  \proofstepwidestar[21]{z\in X}{%
    \rA}
  \proofstepwidestar[22]{w\in X}{%
    \rA}
  \proofstepwidestar[23]{(p,z)\in H\land(q,w)\in H}{%
    \rA}
  \proofstepwidestar[23]{(p,z)\in G}{%
    \FormulaRefAuto{A\subseteq B\dsep x\in A\vdash x\in B}[thm]{12,\rAEa{23}}}
  \proofstepwidestar[23]{(q,w)\in G}{%
    \FormulaRefAuto{A\subseteq B\dsep x\in A\vdash x\in B}[thm]{12,\rAEb{23}}}
  \proofstepwidestar[3,19,20,21,22,23]{(k(p,q),g(z,w))\in G}{%
    \FormulaRefAuto{BinaryTreeClosedGraphNode}[thm]{10,\rAI{19,20},\rAI{21,22},24,25}}
  \proofstepwidestar[]{k(p,q)=k(u,v)\lor k(p,q)\ne k(u,v)}{%
    \FormulaRefAuto{A\lor\neg A}[thm]}
  \proofstepwidestar[28]{k(p,q)=k(u,v)}{%
    \rA}
  \proofstepwidestar[1,4,5,19,20,28]{p=u\land q=v}{%
    \FormulaRefAuto{BinaryTreeStructureNodeAxiom}[ax]{1,\rAI{19,\rAI{20,\rAI{4,5}}},28}}
  \proofstepwidestar[1,4,5,19,20,28,23]{(u,z)\in G}{%
    \rIE{\rAEa{29},24}}
  \proofstepwidestar[1,4,5,19,20,28,23]{(v,w)\in G}{%
    \rIE{\rAEb{29},25}}
  \proofstepwidestar[1,4,5,19,20,28,23,6,21]{z=x}{%
    \FormulaRefAuto{\exists! x\,P(x),\;P(a),\;P(b)\vdash a=b}[thm]{6,\rAI{21,30},8}}
  \proofstepwidestar[1,4,5,19,20,28,23,7,22]{w=y}{%
    \FormulaRefAuto{\exists! x\,P(x),\;P(a),\;P(b)\vdash a=b}[thm]{7,\rAI{22,31},9}}
  \proofstepwidestar[1,4,5,19,20,28,23,6,7,21,22]{g(z,w)=g(x,y)}{%
    \rIE{32,\rIE{33,\rII}}}
  \proofstepwidestar[1,4,5,19,20,28,23,6,7,21,22]{k(p,q)\ne k(u,v)\lor g(z,w)=g(x,y)}{%
    \rOIb{34}}
  \proofstepwidestar[36]{k(p,q)\ne k(u,v)}{%
    \rA}
  \proofstepwidestar[36]{k(p,q)\ne k(u,v)\lor g(z,w)=g(x,y)}{%
    \rOIa{36}}
  \proofstepwidestar[1,4,5,19,20,23,6,7,21,22]{k(p,q)\ne k(u,v)\lor g(z,w)=g(x,y)}{%
    \rOE{27,28,35,36,37}}
  \proofstepwidestar[1,3,4,5,19,20,23,6,7,21,22]{(k(p,q),g(z,w))\in H}{%
    \rRE{\rLREb{\FormulaRefAuto{WordRecursionFilterMembership}[thm]{11}},\rAI{26,38}}}
  \proofstepwidestar[1,3,4,5,6,7]{\begin{aligned}[t]&\forall p,q\in T\,\forall z,w\in X\\&((p,z)\in H\land(q,w)\in H\\&\quad\rightarrow(k(p,q),g(z,w))\in H)\end{aligned}}{%
    \rUI{\rRI{19,\rUI{\rRI{20,\rUI{\rRI{21,\rUI{\rRI{22,\rRI{23,39}}}}}}}}}}
  \proofstepwidestar[1,3,4,5,6,7]{C(H)}{%
    \FormulaRefAuto{BinaryTreeClosedGraphDef}[def]{\rAI{13,\rAI{18,40}}}}
  \proofstepwidestar[6,7]{(x,y)\in X\times X}{%
    \FormulaRefAuto{a\in A\dsep b\in B\vdash(a,b)\in A\times B}[thm]{\rAEa{8},\rAEa{9}}}
  \proofstepwidestar[2,6,7]{g(x,y)\in X}{%
    \FormulaRefAuto{F\colon A\to B\dsep x\in A\vdash F(x)\in B}[thm]{2,42}}
  \proofstepwidestar[3,4,5,6,7]{(k(u,v),g(x,y))\in G}{%
    \FormulaRefAuto{BinaryTreeClosedGraphNode}[thm]{10,\rAI{4,5},\rAI{\rAEa{8},\rAEa{9}},\rAEb{8},\rAEb{9}}}
  \proofstepwidestar[1,2,3,4,5,6,7]{\exists!z\in X\;(k(u,v),z)\in G}{%
    \FormulaRefAuto{BinaryTreeRecursionFilteredFiberUnique}[thm]{3,41,43,44}}
\end{tabproofwide}

\noindent\begin{minipage}{\linewidth}
\FormulaThmDeltaK[Der kleinste Baumgraph ist eine Funktion]{%
\begin{aligned}[t]&\mathsf{BaumStr}_A(T,\ell,k)\dsep f\colon A\to X\dsep g\colon X\times X\to X\\&\dsep M(G)\vdash G\colon T\to X\end{aligned}
}{BinaryTreeRecursionGraphFunction}{%
\DeltaRow{Mengen}{A\dsep T\dsep X\dsep G\dsep H\dsep t\dsep u\dsep v\dsep x\dsep y\dsep a}
\DeltaRow{Funktionen}{\ell\dsep k\dsep f\dsep g}
\DeltaRow{Prädikatsabkürzungen}{\begin{aligned}[t]
 C(H)&\coloneqq\mathsf{BCl}_A(H;T,\ell,k;X,f,g),\\[-2pt]
 M(G)&\coloneqq C(G)\land\forall H(C(H)\rightarrow G\subseteq H)
\end{aligned}}\DeltaRow{Prädikatsabkürzung}{P(t)\coloneqq\exists!x\in X\;(t,x)\in G}}
\end{minipage}\par
\begin{tabproofwide}
  \proofstepwidestar[1]{\mathsf{BaumStr}_A(T,\ell,k)}{%
    \rA}
  \proofstepwidestar[2]{f\colon A\to X}{%
    \rA}
  \proofstepwidestar[3]{g\colon X\times X\to X}{%
    \rA}
  \proofstepwidestar[4]{M(G)}{%
    \rA}
  \proofstepwidestar[5]{a\in A}{%
    \rA}
  \proofstepwidestar[1,2,4,5]{P(\ell(a))}{%
    \FormulaRefAuto{BinaryTreeRecursionLeafFiberUnique}[thm]{1,2,4,5}}
  \proofstepwidestar[1,2,4]{\forall a\in A\;P(\ell(a))}{%
    \rUI{\rRI{5,6}}}
  \proofstepwidestar[8]{u\in T}{%
    \rA}
  \proofstepwidestar[9]{v\in T}{%
    \rA}
  \proofstepwidestar[10]{P(u)\land P(v)}{%
    \rA}
  \proofstepwidestar[1,3,4,8,9,10]{P(k(u,v))}{%
    \FormulaRefAuto{BinaryTreeRecursionNodeFiberUnique}[thm]{1,3,4,\rAI{8,9},\rAEa{10},\rAEb{10}}}
  \proofstepwidestar[1,3,4]{\forall u,v\in T\;(P(u)\land P(v)\rightarrow P(k(u,v)))}{%
    \rUI{\rRI{8,\rUI{\rRI{9,\rRI{10,11}}}}}}
  \proofstepwidestar[1,2,3,4]{\forall t\in T\;\exists!x\in X\;(t,x)\in G}{%
    \FormulaRefAuto{BinaryTreeStructureInduction}[thm]{1,7,12}}
  \proofstepwidestar[4]{G\subseteq T\times X}{%
    \FormulaRefAuto{BinaryTreeClosedGraphSubset}[thm]{\rAEa{4}}}
  \proofstepwidestar[1,2,3,4]{G\colon T\to X}{%
    \FormulaRefAuto{UniqueValuedGraphFunction}[thm]{14,13}}
\end{tabproofwide}

\noindent\begin{minipage}{\linewidth}
\FormulaThmDeltaK[Ein abgeschlossener Funktionsgraph erfüllt die Rekursionsgleichungen]{%
\mathsf{BaumStr}_A(T,\ell,k)\dsep G\colon T\to X\dsep C(G)\vdash\mathsf{TRec}_A(G;T,\ell,k;X,f,g)
}{BinaryTreeRecursionGraphEquations}{%
\DeltaRow{Mengen}{A\dsep T\dsep X\dsep G\dsep H\dsep t\dsep u\dsep v\dsep x\dsep y\dsep a}
\DeltaRow{Funktionen}{\ell\dsep k\dsep f\dsep g}
\DeltaRow{Prädikatsabkürzungen}{\begin{aligned}[t]
 C(H)&\coloneqq\mathsf{BCl}_A(H;T,\ell,k;X,f,g),\\[-2pt]
 M(G)&\coloneqq C(G)\land\forall H(C(H)\rightarrow G\subseteq H)
\end{aligned}}}
\end{minipage}\par
\begin{tabproofwide}
  \proofstepwidestar[1]{\mathsf{BaumStr}_A(T,\ell,k)}{%
    \rA}
  \proofstepwidestar[2]{G\colon T\to X}{%
    \rA}
  \proofstepwidestar[3]{C(G)}{%
    \rA}
  \proofstepwidestar[4]{a\in A}{%
    \rA}
  \proofstepwidestar[3,4]{(\ell(a),f(a))\in G}{%
    \FormulaRefAuto{BinaryTreeClosedGraphLeaf}[thm]{3,4}}
  \proofstepwidestar[2,3,4]{G(\ell(a))=f(a)}{%
    \FormulaRefAuto{F\colon A\to B\dsep(x,y)\in F\vdash F(x)=y}[thm]{2,5}}
  \proofstepwidestar[2,3]{\forall a\in A\;G(\ell(a))=f(a)}{%
    \rUI{\rRI{4,6}}}
  \proofstepwidestar[8]{u\in T}{%
    \rA}
  \proofstepwidestar[9]{v\in T}{%
    \rA}
  \proofstepwidestar[2,8]{G(u)\in X}{%
    \FormulaRefAuto{F\colon A\to B\dsep x\in A\vdash F(x)\in B}[thm]{2,8}}
  \proofstepwidestar[2,9]{G(v)\in X}{%
    \FormulaRefAuto{F\colon A\to B\dsep x\in A\vdash F(x)\in B}[thm]{2,9}}
  \proofstepwidestar[2,8]{(u,G(u))\in G}{%
    \FormulaRefAuto{F\colon A\to B\dsep x\in A\vdash(x,F(x))\in F}[thm]{2,8}}
  \proofstepwidestar[2,9]{(v,G(v))\in G}{%
    \FormulaRefAuto{F\colon A\to B\dsep x\in A\vdash(x,F(x))\in F}[thm]{2,9}}
  \proofstepwidestar[2,3,8,9]{(k(u,v),g(G(u),G(v)))\in G}{%
    \FormulaRefAuto{BinaryTreeClosedGraphNode}[thm]{3,\rAI{8,9},\rAI{10,11},12,13}}
  \proofstepwidestar[2,3,8,9]{G(k(u,v))=g(G(u),G(v))}{%
    \FormulaRefAuto{F\colon A\to B\dsep(x,y)\in F\vdash F(x)=y}[thm]{2,14}}
  \proofstepwidestar[2,3]{\forall u,v\in T\;G(k(u,v))=g(G(u),G(v))}{%
    \rUI{\rRI{8,\rUI{\rRI{9,15}}}}}
  \proofstepwidestar[2,3]{\mathsf{TRec}_A(G;T,\ell,k;X,f,g)}{%
    \FormulaRefAuto{BinaryTreeRecursionSpecificationDef}[def]{\rAI{2,\rAI{7,16}}}}
\end{tabproofwide}

\noindent\begin{minipage}{\linewidth}
\FormulaThmDeltaK[Zwei Lösungen derselben Baumrekursion stimmen überein]{%
\begin{aligned}[t]&\mathsf{BaumStr}_A(T,\ell,k)\dsep\mathsf{TRec}_A(F;T,\ell,k;X,f,g)\\&\dsep\mathsf{TRec}_A(F';T,\ell,k;X,f,g)\vdash F=F'\end{aligned}
}{BinaryTreeRecursionTwoSolutionsEqual}{%
\DeltaRow{Mengen}{A\dsep T\dsep X\dsep G\dsep H\dsep t\dsep u\dsep v\dsep x\dsep y\dsep a}
\DeltaRow{Funktionen}{\ell\dsep k\dsep f\dsep g}
\DeltaRow{Prädikatsabkürzungen}{\begin{aligned}[t]
 C(H)&\coloneqq\mathsf{BCl}_A(H;T,\ell,k;X,f,g),\\[-2pt]
 M(G)&\coloneqq C(G)\land\forall H(C(H)\rightarrow G\subseteq H)
\end{aligned}}\DeltaRow{Funktionen}{F\dsep F'}\DeltaRow{Prädikatsabkürzung}{Q(t)\coloneqq F(t)=F'(t)}}
\end{minipage}\par
\begin{tabproofwide}
  \proofstepwidestar[1]{\mathsf{BaumStr}_A(T,\ell,k)}{%
    \rA}
  \proofstepwidestar[2]{\mathsf{TRec}_A(F;T,\ell,k;X,f,g)}{%
    \rA}
  \proofstepwidestar[3]{\mathsf{TRec}_A(F';T,\ell,k;X,f,g)}{%
    \rA}
  \proofstepwidestar[2]{F\colon T\to X}{%
    \rAEa{\FormulaRefAuto{BinaryTreeRecursionSpecificationDef}[def]{2}}}
  \proofstepwidestar[3]{F'\colon T\to X}{%
    \rAEa{\FormulaRefAuto{BinaryTreeRecursionSpecificationDef}[def]{3}}}
  \proofstepwidestar[2]{\forall a\in A\;F(\ell(a))=f(a)}{%
    \rAEn{\FormulaRefAuto{BinaryTreeRecursionSpecificationDef}[def]{2}}}
  \proofstepwidestar[3]{\forall a\in A\;F'(\ell(a))=f(a)}{%
    \rAEn{\FormulaRefAuto{BinaryTreeRecursionSpecificationDef}[def]{3}}}
  \proofstepwidestar[2]{\forall u,v\in T\;F(k(u,v))=g(F(u),F(v))}{%
    \rAEn{\FormulaRefAuto{BinaryTreeRecursionSpecificationDef}[def]{2}}}
  \proofstepwidestar[3]{\forall u,v\in T\;F'(k(u,v))=g(F'(u),F'(v))}{%
    \rAEn{\FormulaRefAuto{BinaryTreeRecursionSpecificationDef}[def]{3}}}
  \proofstepwidestar[10]{a\in A}{%
    \rA}
  \proofstepwidestar[2,10]{F(\ell(a))=f(a)}{%
    \rRE{\rUE{6},10}}
  \proofstepwidestar[3,10]{F'(\ell(a))=f(a)}{%
    \rRE{\rUE{7},10}}
  \proofstepwidestar[2,3,10]{Q(\ell(a))}{%
    \FormulaRefAuto{a=b\dsep c=b\vdash a=c}[thm]{11,12}}
  \proofstepwidestar[2,3]{\forall a\in A\;Q(\ell(a))}{%
    \rUI{\rRI{10,13}}}
  \proofstepwidestar[15]{u\in T}{%
    \rA}
  \proofstepwidestar[16]{v\in T}{%
    \rA}
  \proofstepwidestar[17]{Q(u)\land Q(v)}{%
    \rA}
  \proofstepwidestar[2,15,16]{F(k(u,v))=g(F(u),F(v))}{%
    \rRE{\rUE{\rRE{\rUE{8},15}},16}}
  \proofstepwidestar[3,15,16]{F'(k(u,v))=g(F'(u),F'(v))}{%
    \rRE{\rUE{\rRE{\rUE{9},15}},16}}
  \proofstepwidestar[3,15,16,17]{F'(k(u,v))=g(F(u),F(v))}{%
    \rIE{\FormulaRefAuto{a=b\vdash b=a}[thm]{\rAEa{17}},\rIE{\FormulaRefAuto{a=b\vdash b=a}[thm]{\rAEb{17}},19}}}
  \proofstepwidestar[2,3,15,16,17]{Q(k(u,v))}{%
    \FormulaRefAuto{a=b\dsep c=b\vdash a=c}[thm]{18,20}}
  \proofstepwidestar[2,3]{\forall u,v\in T\;(Q(u)\land Q(v)\rightarrow Q(k(u,v)))}{%
    \rUI{\rRI{15,\rUI{\rRI{16,\rRI{17,21}}}}}}
  \proofstepwidestar[1,2,3]{\forall t\in T\;F(t)=F'(t)}{%
    \FormulaRefAuto{BinaryTreeStructureInduction}[thm]{1,14,22}}
  \proofstepwidestar[1,2,3]{F=F'}{%
    \FormulaRefAuto{FunctionExtensionality}[thm]{4,5,23}}
\end{tabproofwide}

\noindent\begin{minipage}{\linewidth}
\FormulaThmDeltaK[Rekursion in jeder freien Baumstruktur]{%
\begin{aligned}[t]&\mathsf{BaumStr}_A(T,\ell,k)\dsep f\colon A\to X\dsep g\colon X\times X\to X\\&\vdash\exists!F\;\mathsf{TRec}_A(F;T,\ell,k;X,f,g)\end{aligned}
}{BinaryTreeStructureRecursion}{%
\DeltaRow{Mengen}{A\dsep T\dsep X\dsep G\dsep H\dsep t\dsep u\dsep v\dsep x\dsep y\dsep a}
\DeltaRow{Funktionen}{\ell\dsep k\dsep f\dsep g}
\DeltaRow{Prädikatsabkürzungen}{\begin{aligned}[t]
 C(H)&\coloneqq\mathsf{BCl}_A(H;T,\ell,k;X,f,g),\\[-2pt]
 M(G)&\coloneqq C(G)\land\forall H(C(H)\rightarrow G\subseteq H)
\end{aligned}}\DeltaRow{Funktionen}{F\dsep F'}}
\end{minipage}\par
\begin{tabproofwide}
  \proofstepwidestar[1]{\mathsf{BaumStr}_A(T,\ell,k)}{%
    \rA}
  \proofstepwidestar[2]{f\colon A\to X}{%
    \rA}
  \proofstepwidestar[3]{g\colon X\times X\to X}{%
    \rA}
  \proofstepwidestar[1,2,3]{\exists G\;M(G)}{%
    \FormulaRefAuto{BinaryTreeLeastClosedGraphExists}[thm]{1,2,3}}
  \proofstepwidestar[5]{M(G)}{%
    \rA}
  \proofstepwidestar[1,2,3,5]{G\colon T\to X}{%
    \FormulaRefAuto{BinaryTreeRecursionGraphFunction}[thm]{1,2,3,5}}
  \proofstepwidestar[1,2,3,5]{\mathsf{TRec}_A(G;T,\ell,k;X,f,g)}{%
    \FormulaRefAuto{BinaryTreeRecursionGraphEquations}[thm]{1,6,\rAEa{5}}}
  \proofstepwidestar[1,2,3,5]{\exists F\;\mathsf{TRec}_A(F;T,\ell,k;X,f,g)}{%
    \rEI{7}}
  \proofstepwidestar[1,2,3]{\exists F\;\mathsf{TRec}_A(F;T,\ell,k;X,f,g)}{%
    \rEE{4,5,8}}
  \proofstepwidestar[10]{\mathsf{TRec}_A(F;T,\ell,k;X,f,g)}{%
    \rA}
  \proofstepwidestar[11]{\mathsf{TRec}_A(F';T,\ell,k;X,f,g)}{%
    \rA}
  \proofstepwidestar[1,10,11]{F=F'}{%
    \FormulaRefAuto{BinaryTreeRecursionTwoSolutionsEqual}[thm]{1,10,11}}
  \proofstepwidestar[1,2,3]{\exists!F\;\mathsf{TRec}_A(F;T,\ell,k;X,f,g)}{%
    \UEI{9,10,11,12}}
\end{tabproofwide}

Die gesamte Konstruktion gilt auch bei leerem Alphabet und leerem
Zielträger, sofern die angegebenen Abbildungen existieren. Insbesondere
wird im Existenzbeweis kein zusätzlicher Wert aus \(X\) ausgewählt.


\noindent\begin{minipage}{\linewidth}
\FormulaThmDeltaK[Konstruktorzerlegung]{%
  \begin{aligned}[t]
    R\in\BracketTreeSet{A}\eqvdash{}&
      \exists a\in A\,(R=\LeafTree{A}{a})\ \lor{}\\[-2pt]
      &\exists S,T\in\BracketTreeSet{A}\,
        (R=\NodeTree{A}{S}{T})
  \end{aligned}%
}{TreeConstructorDecomposition}{%
  \DeltaRow{Mengen}{A\dsep R\dsep a\dsep S\dsep T\dsep n\dsep k}
}
\end{minipage}\par

% Konkrete Instanz der abstrakten Zerlegung mit expliziten Zeugen.
\begin{tabproofwide}
  \proofstepwidestar[]{\mathsf{BaumStr}_A(\BracketTreeSet A,\lambda_A,\kappa_A)}{%
    \FormulaRefAuto{BinaryTreeStructureConcreteModel}[thm]}
  \proofstepwidestar[2]{R\in\BracketTreeSet A}{%
    \rA}
  \proofstepwidestar[2]{\begin{aligned}[t]&(\exists!a\in A\;R=\lambda_A(a))\lor{}\\[-2pt]&(\exists!p\in\BracketTreeSet A\times\BracketTreeSet A\;R=\kappa_A(p))\end{aligned}}{%
    \FormulaRefAuto{BinaryTreeStructureDecomposition}[thm]{1,2}}
  \proofstepwidestar[4]{\exists!a\in A\;R=\lambda_A(a)}{%
    \rA}
  \proofstepwidestar[4]{\exists a\in A\;R=\lambda_A(a)}{%
    \FormulaRefAuto{\exists! x\,P(x)\vdash \exists x\,P(x)}[thm]{4}}
  \proofstepwidestar[6]{a\in A\land R=\lambda_A(a)}{%
    \rA}
  \proofstepwidestar[6]{\lambda_A(a)=\LeafTree A a}{%
    \FormulaRefAuto{TreeCodeLeafValue}[thm]{\rAEa{6}}}
  \proofstepwidestar[6]{R=\LeafTree A a}{%
    \rIE{7,\rAEb{6}}}
  \proofstepwidestar[6]{\exists a\in A\;R=\LeafTree A a}{%
    \rEI{\rAI{\rAEa{6},8}}}
  \proofstepwidestar[4]{\exists a\in A\;R=\LeafTree A a}{%
    \rEE{5,6,9}}
  \proofstepwidestar[4]{\begin{aligned}[t]&(\exists a\in A\;R=\LeafTree A a)\lor{}\\[-2pt]&(\exists S,T\in\BracketTreeSet A\;R=\NodeTree A S T)\end{aligned}}{%
    \rOIa{10}}
  \proofstepwidestar[12]{\exists!p\in\BracketTreeSet A\times\BracketTreeSet A\;R=\kappa_A(p)}{%
    \rA}
  \proofstepwidestar[12]{\exists p\in\BracketTreeSet A\times\BracketTreeSet A\;R=\kappa_A(p)}{%
    \FormulaRefAuto{\exists! x\,P(x)\vdash \exists x\,P(x)}[thm]{12}}
  \proofstepwidestar[14]{p\in\BracketTreeSet A\times\BracketTreeSet A\land R=\kappa_A(p)}{%
    \rA}
  \proofstepwidestar[14]{\exists S\in\BracketTreeSet A\,\exists T\in\BracketTreeSet A\;p=(S,T)}{%
    \FormulaRefAuto{z\in A\times B\eqvdash\exists a\in A\exists b\in B\;z=(a,b)}[thm]{\rAEa{14}}}
  \proofstepwidestar[16]{S\in\BracketTreeSet A\land\exists T\in\BracketTreeSet A\;p=(S,T)}{%
    \rA}
  \proofstepwidestar[17]{T\in\BracketTreeSet A\land p=(S,T)}{%
    \rA}
  \proofstepwidestar[16,17]{\kappa_A(S,T)=\NodeTree A S T}{%
    \FormulaRefAuto{TreeCodeNodeValue}[thm]{\rAEa{16},\rAEa{17}}}
  \proofstepwidestar[14,16,17]{R=\NodeTree A S T}{%
    \rIE{\rAEb{17},18,\rAEb{14}}}
  \proofstepwidestar[14,16,17]{\exists S,T\in\BracketTreeSet A\;R=\NodeTree A S T}{%
    \rEI{\rAI{\rAEa{16},\rEI{\rAI{\rAEa{17},19}}}}}
  \proofstepwidestar[14,16]{\exists S,T\in\BracketTreeSet A\;R=\NodeTree A S T}{%
    \rEE{\rAEb{16},17,20}}
  \proofstepwidestar[14]{\exists S,T\in\BracketTreeSet A\;R=\NodeTree A S T}{%
    \rEE{15,16,21}}
  \proofstepwidestar[12]{\exists S,T\in\BracketTreeSet A\;R=\NodeTree A S T}{%
    \rEE{13,14,22}}
  \proofstepwidestar[12]{\begin{aligned}[t]&(\exists a\in A\;R=\LeafTree A a)\lor{}\\[-2pt]&(\exists S,T\in\BracketTreeSet A\;R=\NodeTree A S T)\end{aligned}}{%
    \rOIb{23}}
  \proofstepwidestar[2]{\begin{aligned}[t]&(\exists a\in A\;R=\LeafTree A a)\lor{}\\[-2pt]&(\exists S,T\in\BracketTreeSet A\;R=\NodeTree A S T)\end{aligned}}{%
    \rOE{3,4,11,12,24}}
  \proofstepwidestar[]{R\in\BracketTreeSet A\rightarrow\bigl(\begin{aligned}[t]&(\exists a\in A\;R=\LeafTree A a)\lor{}\\[-2pt]&(\exists S,T\in\BracketTreeSet A\;R=\NodeTree A S T)\end{aligned}\bigr)}{%
    \rRI{2,25}}
  \proofstepwidestar[27]{\begin{aligned}[t]&(\exists a\in A\;R=\LeafTree A a)\lor{}\\[-2pt]&(\exists S,T\in\BracketTreeSet A\;R=\NodeTree A S T)\end{aligned}}{%
    \rA}
  \proofstepwidestar[28]{\exists a\in A\;R=\LeafTree A a}{%
    \rA}
  \proofstepwidestar[29]{a\in A\land R=\LeafTree A a}{%
    \rA}
  \proofstepwidestar[29]{\LeafTree A a\in\BracketTreeSet A}{%
    \FormulaRefAuto{TreeLeafClosure}[thm]{\rAEa{29}}}
  \proofstepwidestar[29]{R\in\BracketTreeSet A}{%
    \rIE{\FormulaRefAuto{x=y\vdash y=x}[thm]{\rAEb{29}},30}}
  \proofstepwidestar[28]{R\in\BracketTreeSet A}{%
    \rEE{28,29,31}}
  \proofstepwidestar[33]{\exists S,T\in\BracketTreeSet A\;R=\NodeTree A S T}{%
    \rA}
  \proofstepwidestar[34]{S\in\BracketTreeSet A\land\exists T\in\BracketTreeSet A\;R=\NodeTree A S T}{%
    \rA}
  \proofstepwidestar[35]{T\in\BracketTreeSet A\land R=\NodeTree A S T}{%
    \rA}
  \proofstepwidestar[34,35]{\NodeTree A S T\in\BracketTreeSet A}{%
    \FormulaRefAuto{TreeNodeClosure}[thm]{\rAEa{34},\rAEa{35}}}
  \proofstepwidestar[34,35]{R\in\BracketTreeSet A}{%
    \rIE{\FormulaRefAuto{x=y\vdash y=x}[thm]{\rAEb{35}},36}}
  \proofstepwidestar[34]{R\in\BracketTreeSet A}{%
    \rEE{\rAEb{34},35,37}}
  \proofstepwidestar[33]{R\in\BracketTreeSet A}{%
    \rEE{33,34,38}}
  \proofstepwidestar[27]{R\in\BracketTreeSet A}{%
    \rOE{27,28,32,33,39}}
  \proofstepwidestar[]{\bigl(\begin{aligned}[t]&(\exists a\in A\;R=\LeafTree A a)\lor{}\\[-2pt]&(\exists S,T\in\BracketTreeSet A\;R=\NodeTree A S T)\end{aligned}\bigr)\rightarrow R\in\BracketTreeSet A}{%
    \rRI{27,40}}
  \proofstepwidestar[]{R\in\BracketTreeSet A\leftrightarrow\bigl(\begin{aligned}[t]&(\exists a\in A\;R=\LeafTree A a)\lor{}\\[-2pt]&(\exists S,T\in\BracketTreeSet A\;R=\NodeTree A S T)\end{aligned}\bigr)}{%
    \rLRI{26,41}}
\end{tabproofwide}


\section{Die Strukturgesetze im konkreten Baummodell}

Induktion und Rekursion sind bereits aus den Baumaxiomen bewiesen.
Nun werden diese Sätze auf den konstruierten Baumträger angewendet.
Dabei ersetzen wir die Funktionen \(\lambda_A,\kappa_A\) durch
die Schreibweisen für Blätter und Knoten. Alle anschließenden
Positions-, Blattwort- und Auswertungsfunktionen beruhen auf diesen
Struktursätzen.

\noindent\begin{minipage}{\linewidth}
\FormulaThmDeltaK[Strukturinduktion auf Klammerungsbäumen]{%
  \begin{aligned}[t]
    &\forall a\in A\,P(\LeafTree{A}{a})\dsep{}\\[-2pt]
    &\forall S,T\in\BracketTreeSet{A}\,
      \bigl(P(S)\land P(T)\rightarrow
        P(\NodeTree{A}{S}{T})\bigr)\\[-2pt]
    &\vdash\forall R\in\BracketTreeSet{A}\,P(R)
  \end{aligned}%
}{FullPlanarBinaryTreeStructuralInduction}{%
  \DeltaRow{Mengen}{A\dsep a\dsep R\dsep S\dsep T\dsep n}
  \DeltaRow{Prädikate}{P}
}
\end{minipage}\par

% Ersatz nur für den Beweis von FullPlanarBinaryTreeStructuralInduction.
Die Strukturinduktion ist jetzt die Spezialisierung der bereits aus B4
gewonnenen allgemeinen Induktionsregel. Die Baumstufen werden hierfür
nicht erneut benötigt.

\begin{tabproofwide}
  \proofstepwidestar[1]{\forall a\in A\;P(\LeafTree A a)}{\rA}
  \proofstepwidestar[2]{\forall S,T\in\BracketTreeSet A\;
    (P(S)\land P(T)\rightarrow P(\NodeTree A S T))}{\rA}
  \proofstepwidestar[]{\mathsf{BaumStr}_A(\BracketTreeSet A,\lambda_A,\kappa_A)}{%
    \FormulaRefAuto{BinaryTreeStructureConcreteModel}[thm]}
  \proofstepwidestar[1]{\forall a\in A\;P(\lambda_A(a))}{%
    \FormulaRefAuto{BinaryTreeModelConstructorsDef}[def]{1}}
  \proofstepwidestar[2]{\forall S,T\in\BracketTreeSet A\;
    (P(S)\land P(T)\rightarrow P(\kappa_A(S,T)))}{%
    \FormulaRefAuto{BinaryTreeModelConstructorsDef}[def]{2}}
  \proofstepwidestar[1,2]{\forall R\in\BracketTreeSet A\;P(R)}{%
    \FormulaRefAuto{BinaryTreeStructureInduction}[thm]{3,4,5}}
\end{tabproofwide}


\noindent\begin{minipage}{\linewidth}
\FormulaThmDeltaK[Strukturrekursion auf Klammerungsbäumen]{%
  \begin{aligned}[t]
    &f\colon A\to X\dsep g\colon X\times X\to X\vdash{}\\[-2pt]
    &\exists!E\colon\BracketTreeSet{A}\to X\,\Bigl(%
      (\forall a\in A\,
        E(\LeafTree{A}{a})=f(a))\ \land{}\\[-2pt]
    &\qquad(\forall S,T\in\BracketTreeSet{A}\,
      E(\NodeTree{A}{S}{T})=g(E(S),E(T)))
    \Bigr)
  \end{aligned}%
}{FullPlanarBinaryTreeRecursion}{\DeltaRow{Mengen}{A\dsep X\dsep a\dsep S\dsep T}
  \DeltaRow{Funktionen}{f\dsep g\dsep E}
}
\end{minipage}\par

\begin{tabproofwide}
  \proofstepwidestar[1]{f\colon A\to X}{\rA}
  \proofstepwidestar[2]{g\colon X\times X\to X}{\rA}
  \proofstepwidestar[]{\mathsf{BaumStr}_A
    (\BracketTreeSet A,\lambda_A,\kappa_A)}{%
    \FormulaRefAuto{BinaryTreeStructureConcreteModel}[thm]}
  \proofstepwidestar[1,2]{\exists!E\;
    \mathsf{TRec}_A(E;\BracketTreeSet A,\lambda_A,\kappa_A;X,f,g)}{%
    \FormulaRefAuto{BinaryTreeStructureRecursion}[thm]{3,1,2}}
  \proofstepwidestar[]{\begin{aligned}[t]
    &\forall a\in A\;\lambda_A(a)=\LeafTree A a,\\[-2pt]
    &\forall S,T\in\BracketTreeSet A\;
      \kappa_A(S,T)=\NodeTree A S T
  \end{aligned}}{\FormulaRefAuto{BinaryTreeModelConstructorsDef}[def]}
  \proofstepwidestar[1,2]{\begin{aligned}[t]
    &\exists!E\colon\BracketTreeSet A\to X\;\Bigl(\\[-2pt]
    &\quad(\forall a\in A\;E(\LeafTree A a)=f(a))\land{}\\[-2pt]
    &\quad\forall S,T\in\BracketTreeSet A\\[-2pt]
    &\qquad E(\NodeTree A S T)=g(E(S),E(T))\Bigr)
  \end{aligned}}{\parbox[t]{\linewidth}{\raggedright
    \(\FormulaRefAuto{BinaryTreeRecursionSpecificationDef}[def]\), (4,5)}}
\end{tabproofwide}


\subsection{Kanonische Darstellung der Baumstruktur}

Die Rekursionsgleichungen werden zunächst für Identitäten und
Kompositionen nachgewiesen. Dann erhält jede strukturtreue Abbildung
durch eine Gegenrekursion eine beidseitige Umkehrung. Die Bijektivität
folgt damit aus den Funktionssätzen von Band~8.

\noindent\begin{minipage}{\linewidth}
\FormulaThmDeltaK[Die Identität erfüllt die Baumrekursionsgleichungen]{%
\mathsf{BaumStr}_A(T,\ell,k)\vdash\mathsf{TRec}_A(\Id_T;T,\ell,k;T,\ell,k)
}{BinaryTreeRecursionIdentity}{%
\DeltaRow{Mengen}{A\dsep T\dsep V\dsep Z\dsep a\dsep u\dsep v}
\DeltaRow{Funktionen}{\ell\dsep k\dsep m\dsep n\dsep b\dsep c\dsep h\dsep j\dsep q}}
\end{minipage}\par
\begin{tabproofwide}
  \proofstepwidestar[1]{\mathsf{BaumStr}_A(T,\ell,k)}{%
    \rA}
  \proofstepwidestar[]{\Id_T\colon T\to T}{%
    \FormulaRefAuto{IdA}[thm]}
  \proofstepwidestar[3]{a\in A}{%
    \rA}
  \proofstepwidestar[1,3]{\ell(a)\in T}{%
    \FormulaRefAuto{BinaryTreeStructureLeafTyping}[thm]{1,3}}
  \proofstepwidestar[1,3]{\Id_T(\ell(a))=\ell(a)}{%
    \FormulaRefAuto{x\in A\vdash\Id_A(x)=x}[thm]{4}}
  \proofstepwidestar[1]{\forall a\in A\;\Id_T(\ell(a))=\ell(a)}{%
    \rUI{\rRI{3,5}}}
  \proofstepwidestar[7]{u\in T}{%
    \rA}
  \proofstepwidestar[8]{v\in T}{%
    \rA}
  \proofstepwidestar[1,7,8]{k(u,v)\in T}{%
    \FormulaRefAuto{BinaryTreeStructureNodeTyping}[thm]{1,\rAI{7,8}}}
  \proofstepwidestar[1,7,8]{\Id_T(k(u,v))=k(u,v)}{%
    \FormulaRefAuto{x\in A\vdash\Id_A(x)=x}[thm]{9}}
  \proofstepwidestar[7]{\Id_T(u)=u}{%
    \FormulaRefAuto{x\in A\vdash\Id_A(x)=x}[thm]{7}}
  \proofstepwidestar[8]{\Id_T(v)=v}{%
    \FormulaRefAuto{x\in A\vdash\Id_A(x)=x}[thm]{8}}
  \proofstepwidestar[1,7,8]{\Id_T(k(u,v))=k(\Id_T(u),\Id_T(v))}{%
    \rIE{\FormulaRefAuto{a=b\vdash b=a}[thm]{11},\rIE{\FormulaRefAuto{a=b\vdash b=a}[thm]{12},10}}}
  \proofstepwidestar[1]{\forall u,v\in T\;\Id_T(k(u,v))=k(\Id_T(u),\Id_T(v))}{%
    \rUI{\rRI{7,\rUI{\rRI{8,13}}}}}
  \proofstepwidestar[1]{\mathsf{TRec}_A(\Id_T;T,\ell,k;T,\ell,k)}{%
    \FormulaRefAuto{BinaryTreeRecursionSpecificationDef}[def]{\rAI{2,\rAI{6,14}}}}
\end{tabproofwide}

\Needspace{19\baselineskip}
\noindent\begin{minipage}{\linewidth}
\FormulaThmDeltaK[Komposition strukturtreuer Baumabbildungen]{%
\begin{aligned}[t]&\mathsf{BaumStr}_A(T,\ell,k)\dsep\mathsf{BaumStr}_A(V,m,n)\\&\dsep\mathsf{TRec}_A(h;T,\ell,k;V,m,n)\\&\dsep\mathsf{TRec}_A(j;V,m,n;Z,b,c)\\&\vdash\mathsf{TRec}_A(j\circ h;T,\ell,k;Z,b,c)\end{aligned}
}{BinaryTreeRecursionComposition}{%
\DeltaRow{Mengen}{A\dsep T\dsep V\dsep Z\dsep a\dsep u\dsep v}
\DeltaRow{Funktionen}{\ell\dsep k\dsep m\dsep n\dsep b\dsep c\dsep h\dsep j\dsep q}}
\end{minipage}\par
\begin{tabproofwide}
  \proofstepwidestar[1]{\mathsf{BaumStr}_A(T,\ell,k)}{%
    \rA}
  \proofstepwidestar[2]{\mathsf{BaumStr}_A(V,m,n)}{%
    \rA}
  \proofstepwidestar[3]{\mathsf{TRec}_A(h;T,\ell,k;V,m,n)}{%
    \rA}
  \proofstepwidestar[4]{\mathsf{TRec}_A(j;V,m,n;Z,b,c)}{%
    \rA}
  \proofstepwidestar[3]{h\colon T\to V}{%
    \rAEa{\FormulaRefAuto{BinaryTreeRecursionSpecificationDef}[def]{3}}}
  \proofstepwidestar[4]{j\colon V\to Z}{%
    \rAEa{\FormulaRefAuto{BinaryTreeRecursionSpecificationDef}[def]{4}}}
  \proofstepwidestar[3,4]{j\circ h\colon T\to Z}{%
    \FormulaRefAuto{CompositionDef}[def]{5,6}}
  \proofstepwidestar[3]{\forall a\in A\;h(\ell(a))=m(a)}{%
    \rAEn{\FormulaRefAuto{BinaryTreeRecursionSpecificationDef}[def]{3}}}
  \proofstepwidestar[4]{\forall a\in A\;j(m(a))=b(a)}{%
    \rAEn{\FormulaRefAuto{BinaryTreeRecursionSpecificationDef}[def]{4}}}
  \proofstepwidestar[3]{\forall u,v\in T\;h(k(u,v))=n(h(u),h(v))}{%
    \rAEn{\FormulaRefAuto{BinaryTreeRecursionSpecificationDef}[def]{3}}}
  \proofstepwidestar[4]{\forall u,v\in V\;j(n(u,v))=c(j(u),j(v))}{%
    \rAEn{\FormulaRefAuto{BinaryTreeRecursionSpecificationDef}[def]{4}}}
  \proofstepwidestar[12]{a\in A}{%
    \rA}
  \proofstepwidestar[1,12]{\ell(a)\in T}{%
    \FormulaRefAuto{BinaryTreeStructureLeafTyping}[thm]{1,12}}
  \proofstepwidestar[3,4,1,12]{(j\circ h)(\ell(a))=j(h(\ell(a)))}{%
    \FormulaRefAuto{x\in A\vdash(F\circ G)(x)=F(G(x))}[thm]{5,6,13}}
  \proofstepwidestar[3,12]{h(\ell(a))=m(a)}{%
    \rRE{\rUE{8},12}}
  \proofstepwidestar[4,12]{j(m(a))=b(a)}{%
    \rRE{\rUE{9},12}}
  \proofstepwidestar[3,4,1,12]{(j\circ h)(\ell(a))=j(m(a))}{%
    \rIE{15,14}}
  \proofstepwidestar[3,4,1,12]{(j\circ h)(\ell(a))=b(a)}{%
    \FormulaRefAuto{a=b\dsep b=c\vdash a=c}[thm]{17,16}}
  \proofstepwidestar[3,4,1]{\forall a\in A\;(j\circ h)(\ell(a))=b(a)}{%
    \rUI{\rRI{12,18}}}
  \proofstepwidestar[20]{u\in T}{%
    \rA}
  \proofstepwidestar[21]{v\in T}{%
    \rA}
  \proofstepwidestar[3,20]{h(u)\in V}{%
    \FormulaRefAuto{F\colon A\to B\dsep x\in A\vdash F(x)\in B}[thm]{5,20}}
  \proofstepwidestar[3,21]{h(v)\in V}{%
    \FormulaRefAuto{F\colon A\to B\dsep x\in A\vdash F(x)\in B}[thm]{5,21}}
  \proofstepwidestar[1,20,21]{k(u,v)\in T}{%
    \FormulaRefAuto{BinaryTreeStructureNodeTyping}[thm]{1,\rAI{20,21}}}
  \proofstepwidestar[1,3,4,20,21]{(j\circ h)(k(u,v))=j(h(k(u,v)))}{%
    \FormulaRefAuto{x\in A\vdash(F\circ G)(x)=F(G(x))}[thm]{5,6,24}}
  \proofstepwidestar[3,20,21]{h(k(u,v))=n(h(u),h(v))}{%
    \rRE{\rUE{\rRE{\rUE{10},20}},21}}
  \proofstepwidestar[3,4,20,21]{j(n(h(u),h(v)))=c(j(h(u)),j(h(v)))}{%
    \rRE{\rUE{\rRE{\rUE{11},22}},23}}
  \proofstepwidestar[1,3,4,20,21]{(j\circ h)(k(u,v))=j(n(h(u),h(v)))}{%
    \rIE{26,25}}
  \proofstepwidestar[1,3,4,20,21]{(j\circ h)(k(u,v))=c(j(h(u)),j(h(v)))}{%
    \FormulaRefAuto{a=b\dsep b=c\vdash a=c}[thm]{28,27}}
  \proofstepwidestar[3,4,20]{(j\circ h)(u)=j(h(u))}{%
    \FormulaRefAuto{x\in A\vdash(F\circ G)(x)=F(G(x))}[thm]{5,6,20}}
  \proofstepwidestar[3,4,21]{(j\circ h)(v)=j(h(v))}{%
    \FormulaRefAuto{x\in A\vdash(F\circ G)(x)=F(G(x))}[thm]{5,6,21}}
  \proofstepwidestar[1,3,4,20,21]{\begin{aligned}[t]&(j\circ h)(k(u,v))\\&\quad=c((j\circ h)(u),(j\circ h)(v))\end{aligned}}{%
    \rIE{\FormulaRefAuto{a=b\vdash b=a}[thm]{30},\rIE{\FormulaRefAuto{a=b\vdash b=a}[thm]{31},29}}}
  \proofstepwidestar[1,3,4]{\forall u,v\in T\;(j\circ h)(k(u,v))=c((j\circ h)(u),(j\circ h)(v))}{%
    \rUI{\rRI{20,\rUI{\rRI{21,32}}}}}
  \proofstepwidestar[1,3,4]{\mathsf{TRec}_A(j\circ h;T,\ell,k;Z,b,c)}{%
    \FormulaRefAuto{BinaryTreeRecursionSpecificationDef}[def]{\rAI{7,\rAI{19,33}}}}
\end{tabproofwide}

\noindent\begin{minipage}{\linewidth}
\FormulaThmDeltaK[Strukturtreue Rekursoren zwischen freien Baumstrukturen sind bijektiv]{%
\begin{aligned}[t]&\mathsf{BaumStr}_A(T,\ell,k)\dsep\mathsf{BaumStr}_A(V,m,n)\\&\dsep\mathsf{TRec}_A(h;T,\ell,k;V,m,n)\vdash h\colon T\bij V\end{aligned}
}{BinaryTreeStructureRecursorIsomorphism}{%
\DeltaRow{Mengen}{A\dsep T\dsep V\dsep Z\dsep a\dsep u\dsep v}
\DeltaRow{Funktionen}{\ell\dsep k\dsep m\dsep n\dsep b\dsep c\dsep h\dsep j\dsep q}}
\end{minipage}\par
\begin{tabproofwide}
  \proofstepwidestar[1]{\mathsf{BaumStr}_A(T,\ell,k)}{%
    \rA}
  \proofstepwidestar[2]{\mathsf{BaumStr}_A(V,m,n)}{%
    \rA}
  \proofstepwidestar[3]{\mathsf{TRec}_A(h;T,\ell,k;V,m,n)}{%
    \rA}
  \proofstepwidestar[1]{\ell\colon A\to T\land k\colon T\times T\to T}{%
    \FormulaRefAuto{BinaryTreeStructureTypingAxiom}[ax]{1}}
  \proofstepwidestar[1,2]{\exists!q\;\mathsf{TRec}_A(q;V,m,n;T,\ell,k)}{%
    \FormulaRefAuto{BinaryTreeStructureRecursion}[thm]{2,\rAEa{4},\rAEb{4}}}
  \proofstepwidestar[1,2]{\exists q\;\mathsf{TRec}_A(q;V,m,n;T,\ell,k)}{%
    \FormulaRefAuto{\exists! x\,P(x)\vdash\exists x\,P(x)}[thm]{5}}
  \proofstepwidestar[7]{\mathsf{TRec}_A(q;V,m,n;T,\ell,k)}{%
    \rA}
  \proofstepwidestar[1,2,3,7]{\mathsf{TRec}_A(q\circ h;T,\ell,k;T,\ell,k)}{%
    \FormulaRefAuto{BinaryTreeRecursionComposition}[thm]{1,2,3,7}}
  \proofstepwidestar[1]{\mathsf{TRec}_A(\Id_T;T,\ell,k;T,\ell,k)}{%
    \FormulaRefAuto{BinaryTreeRecursionIdentity}[thm]{1}}
  \proofstepwidestar[1,2,3,7]{q\circ h=\Id_T}{%
    \FormulaRefAuto{BinaryTreeRecursionTwoSolutionsEqual}[thm]{1,8,9}}
  \proofstepwidestar[1,2,3,7]{\mathsf{TRec}_A(h\circ q;V,m,n;V,m,n)}{%
    \FormulaRefAuto{BinaryTreeRecursionComposition}[thm]{2,1,7,3}}
  \proofstepwidestar[2]{\mathsf{TRec}_A(\Id_V;V,m,n;V,m,n)}{%
    \FormulaRefAuto{BinaryTreeRecursionIdentity}[thm]{2}}
  \proofstepwidestar[1,2,3,7]{h\circ q=\Id_V}{%
    \FormulaRefAuto{BinaryTreeRecursionTwoSolutionsEqual}[thm]{2,11,12}}
  \proofstepwidestar[3]{h\colon T\to V}{%
    \rAEa{\FormulaRefAuto{BinaryTreeRecursionSpecificationDef}[def]{3}}}
  \proofstepwidestar[7]{q\colon V\to T}{%
    \rAEa{\FormulaRefAuto{BinaryTreeRecursionSpecificationDef}[def]{7}}}
  \proofstepwidestar[1,2,3,7]{h\colon T\bij V}{%
    \FormulaRefAuto{MutuallyInverseCompanionBijection}[thm]{14,15,10,13}}
  \proofstepwidestar[1,2,3]{h\colon T\bij V}{%
    \rEE{6,7,16}}
\end{tabproofwide}

\noindent\begin{minipage}{\linewidth}
\FormulaThmDeltaK[Kanonische Darstellung jeder Baumstruktur]{%
\begin{aligned}[t]&\mathsf{BaumStr}_A(T,\ell,k)\vdash{}\\&\exists!h\;\Bigl(h\colon\BracketTreeSet A\bij T\land{}\\&\qquad\mathsf{TRec}_A(h;\BracketTreeSet A,\lambda_A,\kappa_A;T,\ell,k)\Bigr)\end{aligned}
}{BinaryTreeStructureCanonicalIsomorphism}{%
\DeltaRow{Mengen}{A\dsep T\dsep V\dsep Z\dsep a\dsep u\dsep v}
\DeltaRow{Funktionen}{\ell\dsep k\dsep m\dsep n\dsep b\dsep c\dsep h\dsep j\dsep q}}
\end{minipage}\par
\begin{tabproofwide}
  \proofstepwidestar[1]{\mathsf{BaumStr}_A(T,\ell,k)}{%
    \rA}
  \proofstepwidestar[]{\mathsf{BaumStr}_A(\BracketTreeSet A,\lambda_A,\kappa_A)}{%
    \FormulaRefAuto{BinaryTreeStructureConcreteModel}[thm]}
  \proofstepwidestar[1]{\ell\colon A\to T\land k\colon T\times T\to T}{%
    \FormulaRefAuto{BinaryTreeStructureTypingAxiom}[ax]{1}}
  \proofstepwidestar[1]{\exists!h\;\mathsf{TRec}_A(h;\BracketTreeSet A,\lambda_A,\kappa_A;T,\ell,k)}{%
    \FormulaRefAuto{BinaryTreeStructureRecursion}[thm]{2,\rAEa{3},\rAEb{3}}}
  \proofstepwidestar[1]{\exists h\;\mathsf{TRec}_A(h;\BracketTreeSet A,\lambda_A,\kappa_A;T,\ell,k)}{%
    \FormulaRefAuto{\exists! x\,P(x)\vdash\exists x\,P(x)}[thm]{4}}
  \proofstepwidestar[6]{\mathsf{TRec}_A(h;\BracketTreeSet A,\lambda_A,\kappa_A;T,\ell,k)}{%
    \rA}
  \proofstepwidestar[1,6]{h\colon\BracketTreeSet A\bij T}{%
    \FormulaRefAuto{BinaryTreeStructureRecursorIsomorphism}[thm]{2,1,6}}
  \proofstepwidestar[1,6]{\exists h\;(h\colon\BracketTreeSet A\bij T\land\mathsf{TRec}_A(h;\BracketTreeSet A,\lambda_A,\kappa_A;T,\ell,k))}{%
    \rEI{\rAI{7,6}}}
  \proofstepwidestar[1]{\exists h\;(h\colon\BracketTreeSet A\bij T\land\mathsf{TRec}_A(h;\BracketTreeSet A,\lambda_A,\kappa_A;T,\ell,k))}{%
    \rEE{5,6,8}}
  \proofstepwidestar[10]{h\colon\BracketTreeSet A\bij T\land\mathsf{TRec}_A(h;\BracketTreeSet A,\lambda_A,\kappa_A;T,\ell,k)}{%
    \rA}
  \proofstepwidestar[11]{j\colon\BracketTreeSet A\bij T\land\mathsf{TRec}_A(j;\BracketTreeSet A,\lambda_A,\kappa_A;T,\ell,k)}{%
    \rA}
  \proofstepwidestar[10,11]{h=j}{%
    \FormulaRefAuto{BinaryTreeRecursionTwoSolutionsEqual}[thm]{2,\rAEb{10},\rAEb{11}}}
  \proofstepwidestar[1]{\exists!h\;(h\colon\BracketTreeSet A\bij T\land\mathsf{TRec}_A(h;\BracketTreeSet A,\lambda_A,\kappa_A;T,\ell,k))}{%
    \UEI{9,10,11,12}}
\end{tabproofwide}


% Einbindung nach dem vollständig bewiesenen Strukturrekursionssatz.
\subsection{Der Baumrekursor: Blätter ersetzen, Ergebnisse verknüpfen}

Der Baumrekursor führt zwei festgelegte Schritte aus: Er ersetzt ein Blatt
mit Beschriftung \(a\) durch \(f(a)\), und er verknüpft die Ergebnisse der
beiden Teilbäume mit \(g\). Der Strukturrekursionssatz liefert genau eine
Funktion mit diesen Eigenschaften.

\Needspace{16\baselineskip}
\FormulaDefDeltaK[Baumrekursor]{%
  \begin{aligned}[t]
    &f\colon A\to X\dsep g\colon X\times X\to X\vdash{}\\[-2pt]
    &\TreeRecursor{A}{X}{f}{g}\coloneqq
      \iota E\colon\BracketTreeSet A\to X\,\Bigl(\\[-2pt]
    &\quad\forall a\in A\,E(\LeafTree A a)=f(a)\land{}\\[-2pt]
    &\quad\forall S,T\in\BracketTreeSet A\,
      E(\NodeTree A S T)=g(E(S),E(T))\Bigr)
  \end{aligned}%
}{TreeRecursorDef}{%
  \DeltaRow{Mengen}{A\dsep X\dsep a\dsep S\dsep T}
  \DeltaRow{Funktionen}{f\dsep g\dsep E}
  \DeltaRow{\textbf{Neues Symbol}}{}
  \DeltaPrem{Baumrekursor}{\TreeRecursor{A}{X}{f}{g}}
}

Für die folgenden Aussagen und Beweise wird
\(R\coloneqq\TreeRecursor{A}{X}{f}{g}\) abgekürzt.

\FormulaThmDeltaK[Charakterisierung des Baumrekursors]{%
  \begin{aligned}[t]
    &f\colon A\to X\dsep g\colon X\times X\to X\vdash{}\\[-2pt]
    &R\colon\BracketTreeSet A\to X\land\Bigl(\\[-2pt]
    &\quad\forall a\in A\,R(\LeafTree A a)=f(a)\land{}\\[-2pt]
    &\quad\forall S,T\in\BracketTreeSet A\,
      R(\NodeTree A S T)=g(R(S),R(T))\Bigr)
  \end{aligned}%
}{TreeRecursorSpecification}{%
  \DeltaRow{Mengen}{A\dsep X\dsep a\dsep S\dsep T}
  \DeltaRow{Funktionen}{f\dsep g\dsep E\dsep R}
  \DeltaRow{Abkürzung}{R\coloneqq\TreeRecursor{A}{X}{f}{g}}
}
\begin{tabproofwide}
  \proofstepwidestar[1]{f\colon A\to X}{\rA}
  \proofstepwidestar[2]{g\colon X\times X\to X}{\rA}
  \proofstepwidestar[1,2]{%
    \begin{aligned}[t]
      &\exists!E\colon\BracketTreeSet A\to X\,\Bigl(\\[-2pt]
      &\quad\forall a\in A\,E(\LeafTree A a)=f(a)\land{}\\[-2pt]
      &\quad\forall S,T\in\BracketTreeSet A\,
        E(\NodeTree A S T)\\[-2pt]
      &\qquad=g(E(S),E(T))\Bigr)
    \end{aligned}}{%
    \FormulaRefAuto{FullPlanarBinaryTreeRecursion}[thm]{1,2}}
  \proofstepwidestar[1,2]{%
    \begin{aligned}[t]
      &R\colon\BracketTreeSet A\to X\land\Bigl(\\[-2pt]
      &\quad\forall a\in A\,R(\LeafTree A a)=f(a)\land{}\\[-2pt]
      &\quad\forall S,T\in\BracketTreeSet A\,
        R(\NodeTree A S T)\\[-2pt]
      &\qquad=g(R(S),R(T))\Bigr)
    \end{aligned}}{%
    \FormulaRefAuto{TreeRecursorDef}[def]{1,2,3}}
\end{tabproofwide}

\Needspace{20\baselineskip}
\FormulaThmDeltaK[Typisierung, Blattregel und Knotenregel]{%
  \begin{aligned}[t]
    &\text{(i)}\;f\colon A\to X\dsep g\colon X\times X\to X\vdash
      R\colon\BracketTreeSet A\to X,\\[-2pt]
    &\text{(ii)}\;f\colon A\to X\dsep g\colon X\times X\to X\dsep
      a\in A\vdash R(\LeafTree A a)=f(a),\\[-2pt]
    &\text{(iii)}\;f\colon A\to X\dsep g\colon X\times X\to X\dsep
      S,T\in\BracketTreeSet A\vdash{}\\[-2pt]
    &\qquad R(\NodeTree A S T)=g(R(S),R(T)).
  \end{aligned}%
}{TreeRecursorEquations}{%
  \DeltaRow{Mengen}{A\dsep X\dsep a\dsep S\dsep T}
  \DeltaRow{Funktionen}{f\dsep g\dsep R}
  \DeltaRow{Abkürzung}{R\coloneqq\TreeRecursor{A}{X}{f}{g}}
}
\begin{tabproofsplitwide}
  \Needspace{10\baselineskip}
  \proofpartwideindR[Typisierung]{%
    f\colon A\to X\dsep g\colon X\times X\to X\vdash
      R\colon\BracketTreeSet A\to X
  }{f\colon A\to X\dsep g\colon X\times X\to X\vdash
      R\colon\BracketTreeSet A\to X}[TreeRecursorTyping]
    \proofstepwidestar[1]{f\colon A\to X}{\rA}
    \proofstepwidestar[2]{g\colon X\times X\to X}{\rA}
    \proofstepwidestar[1,2]{%
      \begin{aligned}[t]
        &R\colon\BracketTreeSet A\to X\land\Bigl(\\[-2pt]
        &\quad\forall a\in A\,R(\LeafTree A a)=f(a)\land{}\\[-2pt]
        &\quad\forall S,T\in\BracketTreeSet A\,R(\NodeTree A S T)\\[-2pt]
        &\qquad=g(R(S),R(T))\Bigr)
      \end{aligned}}{\FormulaRefAuto{TreeRecursorSpecification}[thm]{1,2}}
    \proofstepwidestar[1,2]{%
      R\colon\BracketTreeSet A\to X}{%
      \rAEa{3}}
  \closeproofpartwideindR
  \Needspace{10\baselineskip}
  \proofpartwideindR[Blattregel]{%
    f\colon A\to X\dsep g\colon X\times X\to X\dsep a\in A\vdash
      R(\LeafTree A a)=f(a)
  }{f\colon A\to X\dsep g\colon X\times X\to X\dsep a\in A\vdash
      R(\LeafTree A a)=f(a)}[TreeRecursorLeafEquation]
    \proofstepwidestar[1]{f\colon A\to X}{\rA}
    \proofstepwidestar[2]{g\colon X\times X\to X}{\rA}
    \proofstepwidestar[3]{a\in A}{\rA}
    \proofstepwidestar[1,2]{%
      \begin{aligned}[t]
        &R\colon\BracketTreeSet A\to X\land\Bigl(\\[-2pt]
        &\quad\forall a\in A\,R(\LeafTree A a)=f(a)\land{}\\[-2pt]
        &\quad\forall S,T\in\BracketTreeSet A\,R(\NodeTree A S T)\\[-2pt]
        &\qquad=g(R(S),R(T))\Bigr)
      \end{aligned}}{\FormulaRefAuto{TreeRecursorSpecification}[thm]{1,2}}
    \proofstepwidestar[1,2]{%
      \forall a\in A\,R(\LeafTree A a)=f(a)}{%
      \rAEa{\rAEb{4}}}
    \proofstepwidestar[1,2]{a\in A\rightarrow
      R(\LeafTree A a)=f(a)}{\rUE{5}}
    \proofstepwidestar[1,2,3]{R(\LeafTree A a)=f(a)}{\rRE{3,6}}
  \closeproofpartwideindR
  \Needspace{10\baselineskip}
  \proofpartwideindR[Knotenregel]{%
    \begin{aligned}[t]
      &f\colon A\to X\dsep g\colon X\times X\to X\dsep
        S,T\in\BracketTreeSet A\vdash{}\\[-2pt]
      &R(\NodeTree A S T)=g(R(S),R(T))
    \end{aligned}
  }{\begin{aligned}[t]
      &f\colon A\to X\dsep g\colon X\times X\to X\dsep
        S,T\in\BracketTreeSet A\vdash{}\\[-2pt]
      &R(\NodeTree A S T)=g(R(S),R(T))
    \end{aligned}}[TreeRecursorNodeEquation]
    \proofstepwidestar[1]{f\colon A\to X}{\rA}
    \proofstepwidestar[2]{g\colon X\times X\to X}{\rA}
    \proofstepwidestar[3]{S\in\BracketTreeSet A}{\rA}
    \proofstepwidestar[4]{T\in\BracketTreeSet A}{\rA}
    \proofstepwidestar[1,2]{%
      \begin{aligned}[t]
        &R\colon\BracketTreeSet A\to X\land\Bigl(\\[-2pt]
        &\quad\forall a\in A\,R(\LeafTree A a)=f(a)\land{}\\[-2pt]
        &\quad\forall S,T\in\BracketTreeSet A\,R(\NodeTree A S T)\\[-2pt]
        &\qquad=g(R(S),R(T))\Bigr)
      \end{aligned}}{\FormulaRefAuto{TreeRecursorSpecification}[thm]{1,2}}
    \proofstepwidestar[1,2]{%
      \forall S,T\in\BracketTreeSet A\,
        R(\NodeTree A S T)=g(R(S),R(T))}{%
      \rAEb{\rAEb{5}}}
    \proofstepwidestar[1,2]{%
      S\in\BracketTreeSet A\rightarrow
      \forall T\in\BracketTreeSet A\,
        R(\NodeTree A S T)=g(R(S),R(T))}{\rUE{6}}
    \proofstepwidestar[1,2,3]{%
      \forall T\in\BracketTreeSet A\,
        R(\NodeTree A S T)=g(R(S),R(T))}{\rRE{3,7}}
    \proofstepwidestar[1,2,3]{%
      T\in\BracketTreeSet A\rightarrow
        R(\NodeTree A S T)=g(R(S),R(T))}{\rUE{8}}
    \proofstepwidestar[1,2,3,4]{%
      R(\NodeTree A S T)=g(R(S),R(T))}{\rRE{4,9}}
  \closeproofpartwideindR
\end{tabproofsplitwide}

Die nächste Regel identifiziert eine bereits gegebene Funktion: Es genügt,
ihre Typisierung sowie ihre Blatt- und Knotenregel nachzuweisen.

\FormulaThmDeltaK[Identifikationsregel für den Baumrekursor]{%
  \begin{aligned}[t]
    &f\colon A\to X\dsep g\colon X\times X\to X\dsep
      E\colon\BracketTreeSet A\to X\dsep{}\\[-2pt]
    &\forall a\in A\,E(\LeafTree A a)=f(a)\dsep{}\\[-2pt]
    &\forall S,T\in\BracketTreeSet A\,
      E(\NodeTree A S T)=g(E(S),E(T))\vdash E=R
  \end{aligned}%
}{TreeRecursorIdentification}{%
  \DeltaRow{Mengen}{A\dsep X\dsep a\dsep S\dsep T}
  \DeltaRow{Funktionen}{f\dsep g\dsep E\dsep R}
  \DeltaRow{Abkürzung}{R\coloneqq\TreeRecursor{A}{X}{f}{g}}
}
\begin{tabproofwide}
  \proofstepwidestar[1]{f\colon A\to X}{\rA}
  \proofstepwidestar[2]{g\colon X\times X\to X}{\rA}
  \proofstepwidestar[3]{E\colon\BracketTreeSet A\to X}{\rA}
  \proofstepwidestar[4]{\forall a\in A\,E(\LeafTree A a)=f(a)}{\rA}
  \proofstepwidestar[5]{%
    \forall S,T\in\BracketTreeSet A\,
      E(\NodeTree A S T)=g(E(S),E(T))}{\rA}
  \proofstepwidestar[1,2]{%
    \begin{aligned}[t]
      &R\colon\BracketTreeSet A\to X\land\Bigl(\\[-2pt]
      &\quad\forall a\in A\,R(\LeafTree A a)=f(a)\land{}\\[-2pt]
      &\quad\forall S,T\in\BracketTreeSet A\,
        R(\NodeTree A S T)\\[-2pt]
      &\qquad=g(R(S),R(T))\Bigr)
    \end{aligned}}{\FormulaRefAuto{TreeRecursorSpecification}[thm]{1,2}}
  \proofstepwidestar[1,2]{R\colon\BracketTreeSet A\to X}{\rAEa{6}}
  \proofstepwidestar[1,2]{%
    \forall a\in A\,R(\LeafTree A a)=f(a)}{\rAEa{\rAEb{6}}}
  \proofstepwidestar[1,2]{%
    \forall S,T\in\BracketTreeSet A\,
      R(\NodeTree A S T)=g(R(S),R(T))}{\rAEb{\rAEb{6}}}
  \proofstepwidestar[1,2,3]{E,R\colon\BracketTreeSet A\to X}{\rAI{3,7}}
  \proofstepwidestar[1,2,3,4,5]{E=R}{%
    \FormulaRefAuto{\exists! x\,P(x),\; P(a),\; P(b) \vdash a = b}[thm]{%
      \FormulaRefAuto{FullPlanarBinaryTreeRecursion}[thm]{1,2},
      \rAI{3,\rAI{4,5}},6}}
\end{tabproofwide}

\subsection{Rekursion mit einer binären Operation}

Bei einer binären Operation \(\star\) wird ihre bereits definierte
zugehörige Funktion eingesetzt. Wir schreiben hier
\(R_\star\coloneqq
\TreeRecursor{A}{X}{f}{\BinaryOpFunction{X}{\star}}\).
Die Knotenregel verwendet dann unmittelbar das Operationszeichen.

\FormulaThmDeltaK[Baumrekursor mit binärer Operation]{%
  \begin{aligned}[t]
    &\text{(i)}\;f\colon A\to X\dsep\mathsf{BinOp}(X,\star)\vdash
      R_\star\colon\BracketTreeSet A\to X,\\[-2pt]
    &\text{(ii)}\;f\colon A\to X\dsep\mathsf{BinOp}(X,\star)\dsep
      a\in A\vdash R_\star(\LeafTree A a)=f(a),\\[-2pt]
    &\text{(iii)}\;f\colon A\to X\dsep\mathsf{BinOp}(X,\star)\dsep
      S,T\in\BracketTreeSet A\vdash{}\\[-2pt]
    &\qquad R_\star(\NodeTree A S T)=R_\star(S)\star R_\star(T).
  \end{aligned}%
}{TreeRecursorBinaryOperationEquations}{%
  \DeltaRow{Mengen}{A\dsep X\dsep a\dsep S\dsep T}
  \DeltaRow{Funktionen}{f\dsep\star\dsep R_\star}
  \DeltaRow{Abkürzung}{R_\star\coloneqq
    \TreeRecursor{A}{X}{f}{\BinaryOpFunction{X}{\star}}}
}
\begin{tabproofsplitwide}
  \Needspace{10\baselineskip}
  \proofpartwideindR[Typisierung]{%
    f\colon A\to X\dsep\mathsf{BinOp}(X,\star)\vdash
      R_\star\colon\BracketTreeSet A\to X
  }{f\colon A\to X\dsep\mathsf{BinOp}(X,\star)\vdash
      R_\star\colon\BracketTreeSet A\to X}[TreeRecursorBinaryOperationTyping]
    \proofstepwidestar[1]{f\colon A\to X}{\rA}
    \proofstepwidestar[2]{\mathsf{BinOp}(X,\star)}{\rA}
    \proofstepwidestar[2]{%
      \BinaryOpFunction{X}{\star}\colon X\times X\to X}{%
      \FormulaRefAuto{BinaryOperationFunctionDef}[def]{2}}
    \proofstepwidestar[1,2]{R_\star\colon\BracketTreeSet A\to X}{%
      \FormulaRefAuto{TreeRecursorTyping}[thm]{1,3}}
  \closeproofpartwideindR
  \Needspace{10\baselineskip}
  \proofpartwideindR[Blattregel]{%
    f\colon A\to X\dsep\mathsf{BinOp}(X,\star)\dsep a\in A\vdash
      R_\star(\LeafTree A a)=f(a)
  }{f\colon A\to X\dsep\mathsf{BinOp}(X,\star)\dsep a\in A\vdash
      R_\star(\LeafTree A a)=f(a)}[TreeRecursorBinaryOperationLeafEquation]
    \proofstepwidestar[1]{f\colon A\to X}{\rA}
    \proofstepwidestar[2]{\mathsf{BinOp}(X,\star)}{\rA}
    \proofstepwidestar[3]{a\in A}{\rA}
    \proofstepwidestar[2]{%
      \BinaryOpFunction{X}{\star}\colon X\times X\to X}{%
      \FormulaRefAuto{BinaryOperationFunctionDef}[def]{2}}
    \proofstepwidestar[1,2,3]{R_\star(\LeafTree A a)=f(a)}{%
      \FormulaRefAuto{TreeRecursorLeafEquation}[thm]{1,4,3}}
  \closeproofpartwideindR
  \Needspace{10\baselineskip}
  \proofpartwideindR[Knotenregel]{%
    \begin{aligned}[t]
      &f\colon A\to X\dsep\mathsf{BinOp}(X,\star)\dsep
        S,T\in\BracketTreeSet A\vdash{}\\[-2pt]
      &R_\star(\NodeTree A S T)=R_\star(S)\star R_\star(T)
    \end{aligned}
  }{\begin{aligned}[t]
      &f\colon A\to X\dsep\mathsf{BinOp}(X,\star)\dsep
        S,T\in\BracketTreeSet A\vdash{}\\[-2pt]
      &R_\star(\NodeTree A S T)=R_\star(S)\star R_\star(T)
    \end{aligned}}[TreeRecursorBinaryOperationNodeEquation]
    \proofstepwidestar[1]{f\colon A\to X}{\rA}
    \proofstepwidestar[2]{\mathsf{BinOp}(X,\star)}{\rA}
    \proofstepwidestar[3]{S\in\BracketTreeSet A}{\rA}
    \proofstepwidestar[4]{T\in\BracketTreeSet A}{\rA}
    \proofstepwidestar[2]{%
      \BinaryOpFunction{X}{\star}\colon X\times X\to X}{%
      \FormulaRefAuto{BinaryOperationFunctionDef}[def]{2}}
    \proofstepwidestar[1,2]{R_\star\colon\BracketTreeSet A\to X}{%
      \FormulaRefAuto{TreeRecursorTyping}[thm]{1,5}}
    \proofstepwidestar[1,2,3]{R_\star(S)\in X}{%
      \FormulaRefAuto{F\colon A\to B\dsep x\in A\vdash F(x)\in B}[thm]{6,3}}
    \proofstepwidestar[1,2,4]{R_\star(T)\in X}{%
      \FormulaRefAuto{F\colon A\to B\dsep x\in A\vdash F(x)\in B}[thm]{6,4}}
    \proofstepwidestar[1,2,3,4]{%
      R_\star(\NodeTree A S T)=
        \BinaryOpFunction{X}{\star}(R_\star(S),R_\star(T))}{%
      \FormulaRefAuto{TreeRecursorNodeEquation}[thm]{1,5,\rAI{3,4}}}
    \proofstepwidestar[1,2,3,4]{%
      \BinaryOpFunction{X}{\star}(R_\star(S),R_\star(T))=
        R_\star(S)\star R_\star(T)}{%
      \FormulaRefAuto{BinaryOperationFunctionDef}[def]{2,7,8}}
    \proofstepwidestar[1,2,3,4]{%
      R_\star(\NodeTree A S T)=R_\star(S)\star R_\star(T)}{%
      \FormulaRefAuto{a=b\dsep b=c\vdash a=c}[thm]{9,10}}
  \closeproofpartwideindR
\end{tabproofsplitwide}

\FormulaThmDeltaK[Identifikation bei einer binären Operation]{%
  \begin{aligned}[t]
    &f\colon A\to X\dsep\mathsf{BinOp}(X,\star)\dsep
      E\colon\BracketTreeSet A\to X\dsep{}\\[-2pt]
    &\forall a\in A\,E(\LeafTree A a)=f(a)\dsep{}\\[-2pt]
    &\forall S,T\in\BracketTreeSet A\,
      E(\NodeTree A S T)=E(S)\star E(T)\vdash E=R_\star
  \end{aligned}%
}{TreeRecursorBinaryOperationIdentification}{%
  \DeltaRow{Mengen}{A\dsep X\dsep a\dsep S\dsep T}
  \DeltaRow{Funktionen}{f\dsep\star\dsep E\dsep R_\star}
  \DeltaRow{Abkürzung}{R_\star\coloneqq
    \TreeRecursor{A}{X}{f}{\BinaryOpFunction{X}{\star}}}
}
\begin{tabproofwide}
  \proofstepwidestar[1]{f\colon A\to X}{\rA}
  \proofstepwidestar[2]{\mathsf{BinOp}(X,\star)}{\rA}
  \proofstepwidestar[3]{E\colon\BracketTreeSet A\to X}{\rA}
  \proofstepwidestar[4]{\forall a\in A\,E(\LeafTree A a)=f(a)}{\rA}
  \proofstepwidestar[5]{%
    \forall S,T\in\BracketTreeSet A\,
      E(\NodeTree A S T)=E(S)\star E(T)}{\rA}
  \proofstepwidestar[2]{%
    \BinaryOpFunction{X}{\star}\colon X\times X\to X}{%
    \FormulaRefAuto{BinaryOperationFunctionDef}[def]{2}}
  \proofstepwidestar[7]{S\in\BracketTreeSet A}{\rA}
  \proofstepwidestar[8]{T\in\BracketTreeSet A}{\rA}
  \proofstepwidestar[3,7]{E(S)\in X}{%
    \FormulaRefAuto{F\colon A\to B\dsep x\in A\vdash F(x)\in B}[thm]{3,7}}
  \proofstepwidestar[3,8]{E(T)\in X}{%
    \FormulaRefAuto{F\colon A\to B\dsep x\in A\vdash F(x)\in B}[thm]{3,8}}
  \proofstepwidestar[5,7,8]{%
    E(\NodeTree A S T)=E(S)\star E(T)}{%
    \rRE{8,\rUE{\rRE{7,\rUE{5}}}}}
  \proofstepwidestar[2,3,7,8]{%
    \BinaryOpFunction{X}{\star}(E(S),E(T))=E(S)\star E(T)}{%
    \FormulaRefAuto{BinaryOperationFunctionDef}[def]{2,9,10}}
  \proofstepwidestar[2,3,5,7,8]{%
    E(\NodeTree A S T)=\BinaryOpFunction{X}{\star}(E(S),E(T))}{%
    \FormulaRefAuto{a=b\dsep c=b\vdash a=c}[thm]{11,12}}
  \proofstepwidestar[2,3,5]{%
    \forall S,T\in\BracketTreeSet A\,
      E(\NodeTree A S T)=\BinaryOpFunction{X}{\star}(E(S),E(T))}{%
    \rUI{\rRI{7,\rUI{\rRI{8,13}}}}}
  \proofstepwidestar[1,2,3,4,5]{E=R_\star}{%
    \FormulaRefAuto{TreeRecursorIdentification}[thm]{1,6,3,4,14}}
\end{tabproofwide}


\chapter{Positionsmengen und Graphen der Bäume}

\section{Positionsadressen und der zugehörige Baum}

Ein Baumcode ist ein Wort und kann daher gleiche Teilbaumcodes an mehreren
Stellen enthalten. Solche Vorkommen dürfen nicht als ein Knoten
identifiziert werden. Wir verwenden deshalb \emph{Positionsadressen}: Das
Leerwort bezeichnet die Wurzel, und das Voranstellen von (0) bzw. (1)
bezeichnet den Schritt zum linken bzw. rechten Teilbaum. Die Knoten des
folgenden Graphen sind somit Vorkommen, nicht Teilbaumwerte.

\FormulaDefDeltaK[Binäre Positionsadressen]{%
  \BinaryAddresses\coloneqq\FiniteWords{\{0,1\}}%
}{BinaryAddressWordSetDef}{%
  \DeltaRow{Mengen}{\BinaryAddresses\dsep p\dsep q\dsep i}
  \DeltaRow{\textbf{Neues Symbol}}{}
  \DeltaPrem{Binäre Adressen}{\BinaryAddresses}
}

\FormulaDefDeltaK[Voranstellen eines Adressbits]{%
  i\in\{0,1\}\vdash
  \begin{aligned}[t]
    &\AddressPrefixMap{i}\colon\BinaryAddresses\to\BinaryAddresses,\\[-2pt]
    &\forall p\in\BinaryAddresses\,
      \bigl(\AddressPrefix{i}{p}\coloneqq
        \WordConcat{\LetterWord{i}}{p}\bigr)
  \end{aligned}%
}{BinaryAddressPrefixDef}{%
  \DeltaRow{Mengen}{i\dsep p\dsep\BinaryAddresses}
  \DeltaRow{Funktionen}{\AddressPrefixMap{i}}
  \DeltaRow{\textbf{Neues Symbol}}{}
  \DeltaPrem{Adresspräfix}{\AddressPrefixMap{i}}
}
\begin{tabproofwide}
  \proofstepwidestar[1]{i\in\{0,1\}}{\rA}
  \proofstepwidestar[2]{p\in\BinaryAddresses}{\rA}
  \proofstepwidestar[1]{\LetterWord{i}\in\FiniteWords{\{0,1\}}}{%
    \FormulaRefAuto{LetterWordFacts}[thm]{1}}
  \proofstepwidestar[2]{p\in\FiniteWords{\{0,1\}}}{%
    \FormulaRefAuto{BinaryAddressWordSetDef}[def]{2}}
  \proofstepwidestar[1,2]{\WordConcat{\LetterWord i}p\in\FiniteWords{\{0,1\}}}{\FormulaRefAuto{WordConcatenationClosure}[thm]{3,4}}
  \proofstepwidestar[1]{%
    \forall p\in\BinaryAddresses\,
      \bigl(\WordConcat{\LetterWord{i}}{p}\in\BinaryAddresses\bigr)}{%
    \rUI{\rRI{2,5}}}
  \proofstepwidestar[1]{%
    \AddressPrefixMap{i}\colon\BinaryAddresses\to\BinaryAddresses}{%
    \FormulaRefAuto{TypedFunctionDefinitionPrinciple}[thm]{6}}
\end{tabproofwide}

\FormulaThmDeltaK[Trennung und Injektivität der Adresspräfixe]{%
  \begin{aligned}[t]
    &\text{(i)}\;i,j\in\{0,1\}\dsep p,q\in\BinaryAddresses\dsep{}\\[-2pt]
    &\qquad\AddressPrefix{i}{p}=\AddressPrefix{j}{q}\vdash i=j\\[-2pt]
    &\text{(ii)}\;i,j\in\{0,1\}\dsep p,q\in\BinaryAddresses\dsep{}\\[-2pt]
    &\qquad\AddressPrefix{i}{p}=\AddressPrefix{j}{q}\vdash p=q\\[-2pt]
    &\text{(iii)}\;
      \AddressPrefixMap{0}[X]\cap\AddressPrefixMap{1}[Y]=\varnothing
  \end{aligned}%
}{BinaryAddressPrefixSeparation}{%
  \DeltaRow{Mengen}{i\dsep j\dsep p\dsep q\dsep X\dsep Y}
  \DeltaRow{Funktionen}{\AddressPrefixMap{0}\dsep\AddressPrefixMap{1}}
}
\begin{tabproofsplitwide}
  \proofpartwideindR[Gleichheit der Präfixzeichen]{%
    \begin{aligned}[t]
      &i,j\in\{0,1\}\dsep p,q\in\BinaryAddresses\dsep{}\\[-2pt]
      &\AddressPrefix{i}{p}=\AddressPrefix{j}{q}\vdash i=j
    \end{aligned}
  }{\begin{aligned}[t]
      &i,j\in\{0,1\}\dsep p,q\in\BinaryAddresses\dsep{}\\[-2pt]
      &\AddressPrefix{i}{p}=\AddressPrefix{j}{q}\vdash i=j
    \end{aligned}}[BinaryAddressPrefixLetterInjective]
    \proofstepwidestar[1]{%
      i,j\in\{0,1\}}{%
      \rA}
    \proofstepwidestar[2]{%
      p,q\in\BinaryAddresses}{%
      \rA}
    \proofstepwidestar[3]{%
      \AddressPrefix{i}{p}=\AddressPrefix{j}{q}}{%
      \rA}
    \proofstepwidestar[1]{%
      \WordLength{\LetterWord{i}}=1=
      \WordLength{\LetterWord{j}}}{%
      \FormulaRefAuto{LetterWordFacts}[thm]{1}}
    \proofstepwidestar[1,2,3]{\LetterWord i,\LetterWord j\in\NonemptyWords{\{0,1\}}}{%
      \FormulaRefAuto{LetterWordNonempty}[thm]{1}}
    \proofstepwidestar[1,2,3]{\LetterWord i,\LetterWord j\in\FiniteWords{\{0,1\}}}{%
      \FormulaRefAuto{NonemptyWordFinite}[thm]{5}}
    \proofstepwidestar[1,2,3]{%
      \LetterWord{i}=\LetterWord{j}}{%
      \FormulaRefAuto{WordConcatenationFixedCutLeft}[thm]{
        6,2,4,3}}
    \proofstepwidestar[1,2,3]{%
      i=j}{%
      \FormulaRefAuto{LetterWordInjectivity}[thm]{1,7}}
  \closeproofpartwideindR

  \proofpartwideindR[Gleichheit der Restadressen]{%
    \begin{aligned}[t]
      &i,j\in\{0,1\}\dsep p,q\in\BinaryAddresses\dsep{}\\[-2pt]
      &\AddressPrefix{i}{p}=\AddressPrefix{j}{q}\vdash p=q
    \end{aligned}
  }{\begin{aligned}[t]
      &i,j\in\{0,1\}\dsep p,q\in\BinaryAddresses\dsep{}\\[-2pt]
      &\AddressPrefix{i}{p}=\AddressPrefix{j}{q}\vdash p=q
    \end{aligned}}[BinaryAddressPrefixWordInjective]
    \proofstepwidestar[1]{%
      i,j\in\{0,1\}}{%
      \rA}
    \proofstepwidestar[2]{%
      p,q\in\BinaryAddresses}{%
      \rA}
    \proofstepwidestar[3]{%
      \AddressPrefix{i}{p}=\AddressPrefix{j}{q}}{%
      \rA}
    \proofstepwidestar[1]{%
      \WordLength{\LetterWord{i}}=1=
      \WordLength{\LetterWord{j}}}{%
      \FormulaRefAuto{LetterWordFacts}[thm]{1}}
    \proofstepwidestar[1,2,3]{\LetterWord i,\LetterWord j\in\NonemptyWords{\{0,1\}}}{%
      \FormulaRefAuto{LetterWordNonempty}[thm]{1}}
    \proofstepwidestar[1,2,3]{\LetterWord i,\LetterWord j\in\FiniteWords{\{0,1\}}}{%
      \FormulaRefAuto{NonemptyWordFinite}[thm]{5}}
    \proofstepwidestar[1,2,3]{%
      p=q}{%
      \FormulaRefAuto{WordConcatenationFixedCutRight}[thm]{
        6,2,4,3}}
  \closeproofpartwideindR

  \proofpartwideindR[Disjunktheit der beiden Präfixbilder]{%
    \AddressPrefixMap{0}[X]\cap\AddressPrefixMap{1}[Y]=\varnothing
  }{\AddressPrefixMap{0}[X]\cap\AddressPrefixMap{1}[Y]=\varnothing}%
    [BinaryAddressPrefixImagesDisjoint]
    \proofstepwidestar[1]{%
      z\in\AddressPrefixMap{0}[X]\cap\AddressPrefixMap{1}[Y]}{\rA}
    \proofstepwidestar[1]{%
      \exists p\in X\,(z=\AddressPrefix{0}{p})\land
      \exists q\in Y\,(z=\AddressPrefix{1}{q})}{%
      \FormulaRefAuto{x\in A\cap B\eqvdash x\in A\land x\in B}[thm]{1}}
    \proofstepwidestar[3]{%
      p\in X\land q\in Y\land
      z=\AddressPrefix{0}{p}\land z=\AddressPrefix{1}{q}}{\rA}
    \proofstepwidestar[3]{%
      \AddressPrefix{0}{p}=\AddressPrefix{1}{q}}{%
      \FormulaRefAuto{a=b\dsep a=c\vdash b=c}[thm]{\rAEn{3},\rAEn{3}}}
    \proofstepwidestar[3]{0=1}{%
      \FormulaRefAuto{BinaryAddressPrefixLetterInjective}[thm]{4}}
    \proofstepwidestar[]{0\neq1}{%
      \FormulaRefAuto{PeanoZeroNeOne}[thm]}
    \proofstepwidestar[3]{\bot}{\rBI{5,6}}
    \proofstepwidestar[1]{\bot}{\rEE{2,3,7}}
    \proofstepwidestar[]{%
      \AddressPrefixMap{0}[X]\cap\AddressPrefixMap{1}[Y]=\varnothing}{%
      \FormulaRefAuto{\forall x\,x\notin A\vdash A=\varnothing}[thm]{%
        \rUI{\rRI{1,8}}}}
  \closeproofpartwideindR
\end{tabproofsplitwide}

\FormulaThmDeltaK[Die beiden Adresskinder sind verschieden]{%
  p\in\BinaryAddresses\vdash
    \WordConcat p{\LetterWord0}\neq\WordConcat p{\LetterWord1}%
}{BinaryAddressChildrenDistinct}{%
  \DeltaRow{Mengen}{p\dsep\BinaryAddresses}
}
\begin{tabproofwide}
  \proofstepwidestar[1]{p\in\BinaryAddresses}{\rA}
  \proofstepwidestar[1]{p\in\FiniteWords{\{0,1\}}}{%
    \FormulaRefAuto{BinaryAddressWordSetDef}[def]{1}}
  \proofstepwidestar[]{0,1\in\{0,1\}}{%
    \FormulaRefAuto{a\in\{a,b\}}[thm],
    \FormulaRefAuto{b\in\{a,b\}}[thm]}
  \proofstepwidestar[1]{\WordConcat p{\LetterWord0}=p\cup\{(\WordLength p,0)\}}{%
    \FormulaRefAuto{WordAppendOccurrenceEquation}[thm]{2,3}}
  \proofstepwidestar[1]{\WordConcat p{\LetterWord1}=p\cup\{(\WordLength p,1)\}}{%
    \FormulaRefAuto{WordAppendOccurrenceEquation}[thm]{2,3}}
  \proofstepwidestar[1]{(\WordLength p,1)\in\WordConcat p{\LetterWord1}}{%
    \rIE{5,\FormulaRefAuto{a\in A\cup\{a\}}[thm]}}
  \proofstepwidestar[1]{(\WordLength p,1)\notin p}{%
    \FormulaRefAuto{WordAppendOccurrenceFresh}[thm]{2}}
  \proofstepwidestar[]{1\neq0}{\FormulaRefAuto{PeanoOneNeZero}[thm]}
  \proofstepwidestar[]{(\WordLength p,1)\notin\{(\WordLength p,0)\}}{%
    \FormulaRefAuto{x\notin\{a\}\eqvdash x\neq a}[thm]{%
      \FormulaRefAuto{a\neq b\vdash(c,a)\neq(d,b)}[thm]{8}}}
  \noalign{\penalty10000}% Die drei abschließenden Schritte zusammenhalten.
  \proofstepwidestar[1]{(\WordLength p,1)\notin\WordConcat p{\LetterWord0}}{%
    \rIE{4,\FormulaRefAuto{NotInUnion}[thm]{7,9}}}
  \noalign{\penalty10000}%
  \proofstepwidestar[1]{\WordConcat p{\LetterWord0}\neq\WordConcat p{\LetterWord1}}{%
    \FormulaRefAuto{x\in B, x\not\in A\vdash A\neq B}[thm]{6,10}}
\end{tabproofwide}

\FormulaDefDeltaK[Pfropfen zweier Adressmengen]{%
  X,Y\in\powerset(\BinaryAddresses)\vdash
  \AddressGraft{X}{Y}\coloneqq
  \{\EmptyWord{\{0,1\}}\}\cup
  \AddressPrefixMap{0}[X]\cup\AddressPrefixMap{1}[Y]%
}{BinaryAddressGraftDef}{%
  \DeltaRow{Mengen}{X\dsep Y\dsep\BinaryAddresses}
  \DeltaRow{Funktionen}{\Gamma}
  \DeltaRow{\textbf{Neues Symbol}}{}
  \DeltaPrem{Adresspfropfung}{\Gamma}
}

\FormulaThmDeltaK[Die Adresspfropfung ist eine binäre Operation]{%
  \mathsf{BinOp}\bigl(\powerset(\BinaryAddresses),\Gamma\bigr)%
}{BinaryAddressGraftBinaryOperation}{%
  \DeltaRow{Mengen}{X\dsep Y\dsep\BinaryAddresses}
  \DeltaRow{Funktionen}{\Gamma}
}
\begin{tabproofwide}
  \proofstepwidestar[1]{X,Y\in\powerset(\BinaryAddresses)}{\rA}
  \proofstepwidestar[]{0,1\in\{0,1\}}{%
    \FormulaRefAuto{a\in\{a,b\}}[thm],
    \FormulaRefAuto{b\in\{a,b\}}[thm]}
  \proofstepwidestar[]{%
    \AddressPrefixMap{0},\AddressPrefixMap{1}\colon
      \BinaryAddresses\to\BinaryAddresses}{%
    \FormulaRefAuto{BinaryAddressPrefixDef}[def]{2}}
  \proofstepwidestar[1]{X\subseteq\BinaryAddresses\land Y\subseteq\BinaryAddresses}{%
      \FormulaRefAuto{x\in\powerset(A)\eqvdash x\subseteq A}[thm]{1}}
  \proofstepwidestar[1]{%
    \AddressPrefixMap{0}[X],\AddressPrefixMap{1}[Y]
      \subseteq\BinaryAddresses}{%
    \FormulaRefAuto{C\subseteq A\dsep F\colon A\to B
      \vdash F[C]\subseteq B}[thm]{%
        4,3}}
  \proofstepwidestar[]{%
    \EmptyWord{\{0,1\}}\in\BinaryAddresses}{%
    \FormulaRefAuto{EmptyWordFinite}[thm]}
  \proofstepwidestar[]{%
    \{\EmptyWord{\{0,1\}}\}\subseteq\BinaryAddresses}{%
    \FormulaRefAuto{a\in A\vdash\{a\}\subseteq A}[thm]{6}}
  \proofstepwidestar[1]{%
    \AddressPrefixMap{0}[X]\cup\AddressPrefixMap{1}[Y]
      \subseteq\BinaryAddresses}{%
    \FormulaRefAuto{A\subseteq C\dsep B\subseteq C
      \vdash A\cup B\subseteq C}[thm]{\rAEa{5},\rAEb{5}}}
  \proofstepwidestar[1]{\{\EmptyWord{\{0,1\}}\}\cup(\AddressPrefixMap{0}[X]\cup\AddressPrefixMap{1}[Y])\subseteq\BinaryAddresses}{%
      \FormulaRefAuto{A\subseteq C\dsep B\subseteq C
        \vdash A\cup B\subseteq C}[thm]{7,8}}
  \proofstepwidestar[1]{%
    \AddressGraft{X}{Y}\subseteq\BinaryAddresses}{%
    \FormulaRefAuto{BinaryAddressGraftDef}[def]{1,
      9}}
  \proofstepwidestar[1]{%
    \AddressGraft{X}{Y}\in\powerset(\BinaryAddresses)}{%
    \FormulaRefAuto{x\in\powerset(A)\eqvdash x\subseteq A}[thm]{10}}
  \proofstepwidestar[]{%
    \mathsf{BinOp}\bigl(\powerset(\BinaryAddresses),\Gamma\bigr)}{%
    \FormulaRefAuto{BinaryOperationDef}[def]{11}}
\end{tabproofwide}

Eine volle endliche Präfixmenge wird durch fünf getrennte Axiome
beschrieben.

\par\noindent\begin{minipage}{\linewidth}
\FormulaDefDeltaK[Volle endliche Präfixmenge]{\GoodAddressSet{P}}{GoodBinaryAddressSetDef}{%
  \DeltaRow{Mengen}{P\dsep p\dsep q\dsep i\dsep\BinaryAddresses}
  \DeltaRow{\textbf{Axiome}}{}
  \DeltaRow{Träger}{P\in\powerset(\BinaryAddresses)}
  \DeltaRow{Endlichkeit}{\Finite P}
  \DeltaRow{Wurzel}{\EmptyWord{\{0,1\}}\in P}
  \DeltaRow{Präfixabschluss}{\begin{aligned}[t]&q\in P\setminus\{\EmptyWord{\{0,1\}}\}\vdash{}\\[-2pt]&\exists p\in P\,\exists i\in\{0,1\}\,q=\WordConcat p{\LetterWord i}\end{aligned}}
  \DeltaRow{Volle Verzweigung}{\begin{aligned}[t]&p\in P\vdash{}\\[-2pt]&\WordConcat p{\LetterWord0}\in P\leftrightarrow\WordConcat p{\LetterWord1}\in P\end{aligned}}
  \DeltaRow{\textbf{Neues Symbol}}{}
  \DeltaPrem{Volle Präfixmenge}{\GoodAddressSet{P}}
}
\end{minipage}\par

\par\noindent\begin{minipage}{\linewidth}
\FormulaAxiomDeltaK[Träger einer vollen Präfixmenge]{\GoodAddressSet P\vdash P\in\powerset(\BinaryAddresses)}{GoodBinaryAddressSetCarrierAxiom}{%
  \DeltaRow{Mengen}{P\dsep p\dsep q\dsep i\dsep\BinaryAddresses}
}
\end{minipage}\par

\par\noindent\begin{minipage}{\linewidth}
\FormulaAxiomDeltaK[Endlichkeit einer vollen Präfixmenge]{\GoodAddressSet P\vdash \Finite P}{GoodBinaryAddressSetFiniteAxiom}{%
  \DeltaRow{Mengen}{P\dsep p\dsep q\dsep i\dsep\BinaryAddresses}
}
\end{minipage}\par

\par\noindent\begin{minipage}{\linewidth}
\FormulaAxiomDeltaK[Wurzel einer vollen Präfixmenge]{\GoodAddressSet P\vdash \EmptyWord{\{0,1\}}\in P}{GoodBinaryAddressSetRootAxiom}{%
  \DeltaRow{Mengen}{P\dsep p\dsep q\dsep i\dsep\BinaryAddresses}
}
\end{minipage}\par

\par\noindent\begin{minipage}{\linewidth}
\FormulaAxiomDeltaK[Präfixabschluss einer vollen Präfixmenge]{\begin{aligned}[t]&\GoodAddressSet P\dsep q\in P\setminus\{\EmptyWord{\{0,1\}}\}\vdash{}\\[-2pt]&\exists p\in P\,\exists i\in\{0,1\}\,q=\WordConcat p{\LetterWord i}\end{aligned}}{GoodBinaryAddressSetPrefixAxiom}{%
  \DeltaRow{Mengen}{P\dsep p\dsep q\dsep i\dsep\BinaryAddresses}
}
\end{minipage}\par

\par\noindent\begin{minipage}{\linewidth}
\FormulaAxiomDeltaK[Volle Verzweigung einer Präfixmenge]{\begin{aligned}[t]&\GoodAddressSet P\dsep p\in P\vdash{}\\[-2pt]&\WordConcat p{\LetterWord0}\in P\leftrightarrow\WordConcat p{\LetterWord1}\in P\end{aligned}}{GoodBinaryAddressSetFullnessAxiom}{%
  \DeltaRow{Mengen}{P\dsep p\dsep q\dsep i\dsep\BinaryAddresses}
}
\end{minipage}\par

\par\noindent\begin{minipage}{\linewidth}
\FormulaThmDeltaK[Kriterium für volle endliche Präfixmengen]{%
  \GoodAddressSet P\eqvdash \begin{aligned}[t]
      &P\in\powerset(\BinaryAddresses)\land\Finite P\land
        \EmptyWord{\{0,1\}}\in P\land{}\\[-2pt]
      &\Bigl(\forall q\in P\setminus\{\EmptyWord{\{0,1\}}\}\,
        \exists p\in P\,\exists i\in\{0,1\}\,{}\\[-2pt]
      &\qquad q=\WordConcat p{\LetterWord i}\Bigr)\land{}\\[-2pt]
      &\forall p\in P\,\Bigl(\WordConcat p{\LetterWord0}\in P
        \leftrightarrow{}\\[-2pt]
      &\qquad\WordConcat p{\LetterWord1}\in P\Bigr)
    \end{aligned}%
}{GoodBinaryAddressSetCriterion}{%
  \DeltaRow{Mengen}{P\dsep p\dsep q\dsep i\dsep\BinaryAddresses}
}
\end{minipage}\par
\begingroup
\ProofWideReasonsNearFormula
\begin{tabproofsplitwide}
  \proofpartwide{$\vdash$}
    \proofstepwidestarbreakkeep[1]{\GoodAddressSet P}{\rA}
    \proofstepwidestarbreakkeep[1]{P\in\powerset(\BinaryAddresses)}{\FormulaRefAuto{GoodBinaryAddressSetCarrierAxiom}[ax]{1}}
    \proofstepwidestarbreakkeep[1]{\Finite P}{\FormulaRefAuto{GoodBinaryAddressSetFiniteAxiom}[ax]{1}}
    \proofstepwidestarbreakkeep[1]{\EmptyWord{\{0,1\}}\in P}{\FormulaRefAuto{GoodBinaryAddressSetRootAxiom}[ax]{1}}
    \proofstepwidestarbreakkeep[5]{q\in P\setminus\{\EmptyWord{\{0,1\}}\}}{\rA}
    \proofstepwidestarbreakkeep[1,5]{\exists p\in P\,\exists i\in\{0,1\}\,q=\WordConcat p{\LetterWord i}}{\FormulaRefAuto{GoodBinaryAddressSetPrefixAxiom}[ax]{1,5}}
    \proofstepwidestarbreakkeep[1]{\begin{aligned}[t]
      &\forall q\in P\setminus\{\EmptyWord{\{0,1\}}\}\,
        \exists p\in P\,\exists i\in\{0,1\}\,{}\\[-2pt]
      &\qquad q=\WordConcat p{\LetterWord i}
    \end{aligned}}{\rUI{\rRI{5,6}}}
    \proofstepwidestarbreakkeep[8]{p\in P}{\rA}
    \proofstepwidestarbreakkeep[1,8]{\WordConcat p{\LetterWord0}\in P\leftrightarrow\WordConcat p{\LetterWord1}\in P}{\FormulaRefAuto{GoodBinaryAddressSetFullnessAxiom}[ax]{1,8}}
    \proofstepwidestarbreakkeep[1]{\forall p\in P\,\Bigl(\WordConcat p{\LetterWord0}\in P\leftrightarrow\WordConcat p{\LetterWord1}\in P\Bigr)}{\rUI{\rRI{8,9}}}
    \proofstepwidestarbreakkeep[1]{\begin{aligned}[t]
      &P\in\powerset(\BinaryAddresses)\land\Finite P\land
        \EmptyWord{\{0,1\}}\in P\land{}\\[-2pt]
      &\Bigl(\forall q\in P\setminus\{\EmptyWord{\{0,1\}}\}\,
        \exists p\in P\,\exists i\in\{0,1\}\,{}\\[-2pt]
      &\qquad q=\WordConcat p{\LetterWord i}\Bigr)\land{}\\[-2pt]
      &\forall p\in P\,\Bigl(\WordConcat p{\LetterWord0}\in P
        \leftrightarrow{}\\[-2pt]
      &\qquad\WordConcat p{\LetterWord1}\in P\Bigr)
    \end{aligned}}{\rAI{2,\rAI{3,\rAI{4,\rAI{7,10}}}}}
  \closeproofpartwide
  \proofpartwide{$\dashv$}
    \proofstepwidestarbreakkeep[1]{\begin{aligned}[t]
      &P\in\powerset(\BinaryAddresses)\land\Finite P\land
        \EmptyWord{\{0,1\}}\in P\land{}\\[-2pt]
      &\Bigl(\forall q\in P\setminus\{\EmptyWord{\{0,1\}}\}\,
        \exists p\in P\,\exists i\in\{0,1\}\,{}\\[-2pt]
      &\qquad q=\WordConcat p{\LetterWord i}\Bigr)\land{}\\[-2pt]
      &\forall p\in P\,\Bigl(\WordConcat p{\LetterWord0}\in P
        \leftrightarrow{}\\[-2pt]
      &\qquad\WordConcat p{\LetterWord1}\in P\Bigr)
    \end{aligned}}{\rA}
    \proofstepwidestarbreakkeep[1]{P\in\powerset(\BinaryAddresses)}{\rAEa{1}}
    \proofstepwidestarbreakkeep[1]{\Finite P}{\rAEa{\rAEb{1}}}
    \proofstepwidestarbreakkeep[1]{\EmptyWord{\{0,1\}}\in P}{\rAEa{\rAEb{\rAEb{1}}}}
    \proofstepwidestarbreakkeep[1]{\begin{aligned}[t]
      &\forall q\in P\setminus\{\EmptyWord{\{0,1\}}\}\,
        \exists p\in P\,\exists i\in\{0,1\}\,{}\\[-2pt]
      &\qquad q=\WordConcat p{\LetterWord i}
    \end{aligned}}{\rAEa{\rAEb{\rAEb{\rAEb{1}}}}}
    \proofstepwidestarbreakkeep[1]{\forall p\in P\,\Bigl(\WordConcat p{\LetterWord0}\in P\leftrightarrow\WordConcat p{\LetterWord1}\in P\Bigr)}{\rAEb{\rAEb{\rAEb{\rAEb{1}}}}}
    \proofstepwidestarbreakkeep[1]{\GoodAddressSet P}{\FormulaRefAuto{GoodBinaryAddressSetDef}[def]{2,3,4,5,6}}
  \closeproofpartwide
\end{tabproofsplitwide}
\endgroup

% Ausgeschriebene Hilfsbeweise für die beiden Pfropfungsbedingungen.
\FormulaThmDeltaK[Ein angehängtes Adressbit erzeugt nicht die Wurzel]{%
  p\in\BinaryAddresses\dsep i\in\{0,1\}\vdash
    \WordConcat p{\LetterWord i}\notin\{\EmptyWord{\{0,1\}}\}
}{B27BinaryAppendNotRoot}{\DeltaRow{Mengen}{p\dsep i}}
\begin{tabproofwide}
  \proofstepwidestar[1]{p\in\BinaryAddresses}{\rA}
  \proofstepwidestar[2]{i\in\{0,1\}}{\rA}
  \proofstepwidestar[1]{p\in\FiniteWords{\{0,1\}}}{\FormulaRefAuto{BinaryAddressWordSetDef}[def]{1}}
  \proofstepwidestar[1,2]{\sigma_{\{0,1\}}(p,i)\ne\EmptyWord{\{0,1\}}}{\FormulaRefAuto{EarlyWordSuccessorNonempty}[thm]{3,2}}
  \proofstepwidestar[1,2]{\sigma_{\{0,1\}}(p,i)=\WordConcat p{\LetterWord i}}{\FormulaRefAuto{WordModelSuccessorConcatenation}[thm]{3,2}}
  \proofstepwidestar[1,2]{\WordConcat p{\LetterWord i}\ne\EmptyWord{\{0,1\}}}{\rIE{5,4}}
  \proofstepwidestar[1,2]{\WordConcat p{\LetterWord i}\notin\{\EmptyWord{\{0,1\}}\}}{\FormulaRefAuto{x \notin \{a\} \eqvdash x \neq a}[thm]{6}}
\end{tabproofwide}

\FormulaThmDeltaK[Ein vorangestelltes Adressbit erzeugt nicht die Wurzel]{%
  p\in\BinaryAddresses\dsep i\in\{0,1\}\vdash
    \AddressPrefix i p\ne\EmptyWord{\{0,1\}}
}{B27BinaryPrefixNotRoot}{\DeltaRow{Mengen}{p\dsep i}}
\begin{tabproofwide}
  \proofstepwidestar[1]{p\in\BinaryAddresses}{\rA}
  \proofstepwidestar[2]{i\in\{0,1\}}{\rA}
  \proofstepwidestar[1]{p\in\FiniteWords{\{0,1\}}}{\FormulaRefAuto{BinaryAddressWordSetDef}[def]{1}}
  \proofstepwidestar[2]{\LetterWord i\in\NonemptyWords{\{0,1\}}}{\FormulaRefAuto{LetterWordNonempty}[thm]{2}}
  \proofstepwidestar[2]{\LetterWord i\in\FiniteWords{\{0,1\}}}{\FormulaRefAuto{NonemptyWordFinite}[thm]{4}}
  \proofstepwidestar[2]{\WordLength{\LetterWord i}\in\mathbb N}{\FormulaRefAuto{FiniteWordLengthNatural}[thm]{5}}
  \proofstepwidestar[1]{\WordLength p\in\mathbb N}{\FormulaRefAuto{FiniteWordLengthNatural}[thm]{3}}
  \proofstepwidestar[2]{\WordLength{\LetterWord i}\ne0}{\rAEb{\FormulaRefAuto{NonemptyWordMembership}[thm]{4}}}
  \proofstepwidestar[1,2]{\WordLength{\LetterWord i}+\WordLength p\ne0}{\FormulaRefAuto{PeanoAddNonzeroLeft}[thm]{6,7,8}}
  \proofstepwidestar[1,2]{\WordConcat{\LetterWord i}p\in\FiniteWords{\{0,1\}}}{\FormulaRefAuto{WordConcatenationClosure}[thm]{5,3}}
  \proofstepwidestar[1,2]{\WordLength{\WordConcat{\LetterWord i}p}=\WordLength{\LetterWord i}+\WordLength p}{\FormulaRefAuto{WordConcatenationLength}[thm]{5,3}}
  \proofstepwidestar[1,2]{\WordLength{\WordConcat{\LetterWord i}p}\ne0}{\rIE{\FormulaRefAuto{x=y\vdash y=x}[thm]{11},9}}
  \proofstepwidestar[1,2]{\WordConcat{\LetterWord i}p\ne\EmptyWord{\{0,1\}}}{\rRE{12,\rLREb{\FormulaRefAuto{FiniteWordLengthNonzeroIffNonempty}[thm]{10}}}}
  \proofstepwidestar[1,2]{\AddressPrefix i p=\WordConcat{\LetterWord i}p}{\FormulaRefAuto{BinaryAddressPrefixDef}[def]{2,1}}
  \proofstepwidestar[1,2]{\AddressPrefix i p\ne\EmptyWord{\{0,1\}}}{\rIE{\FormulaRefAuto{x=y\vdash y=x}[thm]{14},13}}
\end{tabproofwide}

\FormulaThmDeltaK[Die beiden Präfixbilder in einer Pfropfung]{%
  \begin{aligned}[t]
    &X,Y\in\powerset(\BinaryAddresses)\dsep p\in\BinaryAddresses\vdash{}\\[-2pt]
    &(\AddressPrefix0p\in\AddressGraft X Y\leftrightarrow p\in X)\land{}\\[-2pt]
    &(\AddressPrefix1p\in\AddressGraft X Y\leftrightarrow p\in Y)
  \end{aligned}
}{B27GraftPrefixMembership}{\DeltaRow{Mengen}{X\dsep Y\dsep p\dsep q}}
\begin{proof}
\emergencystretch=2em
Nach \FormulaRefAuto{x\in\powerset(A)\eqvdash x\subseteq A}[thm]
liegen \(X,Y\) im Träger der Präfixabbildungen; deren Typisierung
liefert \FormulaRefAuto{BinaryAddressPrefixDef}[def].
Sei \(\AddressPrefix0p\in\AddressGraft X Y\). Die Definition
\FormulaRefAuto{BinaryAddressGraftDef}[def] und das Vereinigungskriterium
\FormulaRefAuto{x\in A\cup B\eqvdash x\in A\lor x\in B}[thm]
geben drei Fälle. Die Zugehörigkeit zur Einermenge der Wurzel ist
wegen \FormulaRefAuto{B27BinaryPrefixNotRoot}[thm] und
\FormulaRefAuto{x\in\{a\}\eqvdash x=a}[thm] ausgeschlossen.
Im zweiten Fall liefert
\FormulaRefAuto{C\subseteq A\dsep F\colon A\to B\dsep y\in F[C]\vdash\exists x\in C\;y=F(x)}[thm]
einen Zeugen \(q\in X\) mit \(\AddressPrefix0p=\AddressPrefix0q\).
\FormulaRefAuto{BinaryAddressPrefixWordInjective}[thm] gibt \(p=q\),
also \(p\in X\). Im dritten Fall gäbe es \(q\in Y\) mit
\(\AddressPrefix0p=\AddressPrefix1q\).
\FormulaRefAuto{BinaryAddressPrefixLetterInjective}[thm] ergäbe
\(0=1\), entgegen \FormulaRefAuto{PeanoZeroNeOne}[thm].

Ist umgekehrt \(p\in X\), liefert derselbe Bildmengensatz in
Einführungsrichtung
\FormulaRefAuto{C\subseteq A\dsep F\colon A\to B\dsep\exists x\in C\;y=F(x)\vdash y\in F[C]}[thm]
mit dem Zeugen \(p\) die Zugehörigkeit
\(\AddressPrefix0p\in\AddressPrefixMap0[X]\).
Zweimalige Vereinigungseinführung und
\FormulaRefAuto{BinaryAddressGraftDef}[def] geben
\(\AddressPrefix0p\in\AddressGraft X Y\).
Für das zweite Präfixbild werden \(0,X\) und \(1,Y\) vertauscht:
Der Wurzelfall bleibt ausgeschlossen, der gleich markierte Bildfall
liefert \(p\in Y\), und der fremd markierte Fall ergäbe \(1=0\),
das durch Symmetrie der Gleichheit wieder \(0=1\) liefert.
Die beiden bewiesenen Richtungen werden jeweils durch \rLRI{} und
die beiden Äquivalenzen durch \rAI{} verbunden.
\end{proof}

\FormulaThmDeltaK[Präfixabschluss der gepfropften Adressmenge]{%
  \begin{aligned}[t]
    &\GoodAddressSet X\dsep\GoodAddressSet Y\vdash{}\\[-2pt]
    &\forall q\in\AddressGraft X Y\setminus\{\EmptyWord{\{0,1\}}\}\,
      \exists p\in\AddressGraft X Y\,\exists i\in\{0,1\}\;q=\WordConcat p{\LetterWord i}
  \end{aligned}
}{B27GraftPrefixClosure}{\DeltaRow{Mengen}{X\dsep Y\dsep q\dsep p\dsep i\dsep u\dsep v\dsep j}}
\begin{proof}
\emergencystretch=2em
\FormulaRefAuto{GoodBinaryAddressSetCarrierAxiom}[ax] liefert die
Trägerzugehörigkeiten von \(X,Y\).
Sei \(q\in\AddressGraft X Y\setminus\{\EmptyWord{\{0,1\}}\}\).
\FormulaRefAuto{BinaryAddressGraftDef}[def], das Vereinigungskriterium
\FormulaRefAuto{x\in A\cup B\eqvdash x\in A\lor x\in B}[thm]
und das Differenzkriterium
\FormulaRefAuto{x\in A\setminus B\vdash x\notin B}[thm]
geben \(q\in\AddressPrefixMap0[X]\) oder
\(q\in\AddressPrefixMap1[Y]\).

Im ersten Fall liefert das Bildmengenkriterium
\FormulaRefAuto{C\subseteq A\dsep F\colon A\to B\dsep y\in F[C]\vdash\exists x\in C\;y=F(x)}[thm]
einen Zeugen \(u\in X\) mit \(q=\AddressPrefix0u\).
Falls \(u=\EmptyWord{\{0,1\}}\), ist
\(q=\LetterWord0=\WordConcat{\EmptyWord{\{0,1\}}}{\LetterWord0}\)
nach \FormulaRefAuto{BinaryAddressPrefixDef}[def] und den beiden
Neutralitätsgesetzen \FormulaRefAuto{WordConcatenationIdentity}[thm].
Die Wurzel liegt in der Pfropfung, da sie in deren erster
Vereinigungskomponente liegt. Die gesuchten Zeugen sind also
\(p=\EmptyWord{\{0,1\}}\), \(i=0\).

Falls \(u\ne\EmptyWord{\{0,1\}}\), liegt \(u\) in
\(X\setminus\{\EmptyWord{\{0,1\}}\}\), nach
\FormulaRefAuto{x \notin \{a\} \eqvdash x \neq a}[thm] und
\FormulaRefAuto{x\in A\setminus B\eqvdash x\in A\land x\notin B}[thm].
\FormulaRefAuto{GoodBinaryAddressSetPrefixAxiom}[ax] liefert
\(v\in X\) und \(j\in\{0,1\}\) mit
\(u=\WordConcat v{\LetterWord j}\). Setze
\(p=\AddressPrefix0v\); nach
\FormulaRefAuto{B27GraftPrefixMembership}[thm] liegt \(p\) in der
Pfropfung. \FormulaRefAuto{BinaryAddressPrefixDef}[def] und
\FormulaRefAuto{WordConcatenationAssociative}[thm] ergeben
\[
 q=\WordConcat{\LetterWord0}{\WordConcat v{\LetterWord j}}
  =\WordConcat{\WordConcat{\LetterWord0}v}{\LetterWord j}
  =\WordConcat p{\LetterWord j}.
\]
Die Typisierungen folgen aus der Trägerzugehörigkeit von \(v\),
\FormulaRefAuto{LetterWordTyping}[thm] sowie
\FormulaRefAuto{WordConcatenationClosure}[thm].
Somit sind \(p\) und \(j\) die gesuchten Zeugen.
Der zweite Bildfall verwendet denselben Nachweis mit \(Y,1\):
Für \(u=\EmptyWord{\{0,1\}}\) sind die Zeugen Wurzel und \(1\),
sonst \(p=\AddressPrefix1v\), \(j\).
\rEE{} entlässt die gewählten Bild- und Vorgängerzeugen;
\rOE{} verbindet die Fälle, und \rUI{} schließt
über das beliebige \(q\).
\end{proof}

\FormulaThmDeltaK[Volle Verzweigung der gepfropften Adressmenge]{%
  \begin{aligned}[t]
    &\GoodAddressSet X\dsep\GoodAddressSet Y\vdash{}\\[-2pt]
    &\forall p\in\AddressGraft X Y\,
      (\WordConcat p{\LetterWord0}\in\AddressGraft X Y
        \leftrightarrow\WordConcat p{\LetterWord1}\in\AddressGraft X Y)
  \end{aligned}
}{B27GraftFullness}{\DeltaRow{Mengen}{X\dsep Y\dsep p\dsep u\dsep j}}
\begin{proof}
\emergencystretch=2em
Sei \(p\in\AddressGraft X Y\). Wie im vorigen Beweis liefern
\FormulaRefAuto{BinaryAddressGraftDef}[def], das Vereinigungs- und
das Einermengenkriterium die drei Fälle Wurzel, linkes Präfixbild und
rechtes Präfixbild. Die nötigen Trägerzugehörigkeiten stammen aus
\FormulaRefAuto{GoodBinaryAddressSetCarrierAxiom}[ax].

Für \(p=\EmptyWord{\{0,1\}}\) liegen nach
\FormulaRefAuto{GoodBinaryAddressSetRootAxiom}[ax] die Wurzeln in
\(X\) und \(Y\). \FormulaRefAuto{B27GraftPrefixMembership}[thm]
zeigt daher \(\AddressPrefix0{\EmptyWord{\{0,1\}}}\) und
\(\AddressPrefix1{\EmptyWord{\{0,1\}}}\) in der Pfropfung.
\FormulaRefAuto{WordConcatenationIdentity}[thm] identifiziert diese
Wörter mit \(\WordConcat p{\LetterWord0}\) bzw.
\(\WordConcat p{\LetterWord1}\). Beide Aussagen sind wahr, also
folgen beide Implikationen und ihre Äquivalenz.

Im linken Bildfall liefert der Bildmengensatz
\FormulaRefAuto{C\subseteq A\dsep F\colon A\to B\dsep y\in F[C]\vdash\exists x\in C\;y=F(x)}[thm]
ein \(u\in X\) mit \(p=\AddressPrefix0u\). Für jedes
\(j\in\{0,1\}\) liefern \FormulaRefAuto{BinaryAddressPrefixDef}[def]
und \FormulaRefAuto{WordConcatenationAssociative}[thm]
\[
 \WordConcat p{\LetterWord j}
 =\AddressPrefix0{\WordConcat u{\LetterWord j}}.
\]
Die zugehörigen Worttypisierungen folgen aus
\FormulaRefAuto{LetterWordTyping}[thm] und
\FormulaRefAuto{WordConcatenationClosure}[thm]. Daher gilt nach
\FormulaRefAuto{B27GraftPrefixMembership}[thm]
\[
 \WordConcat p{\LetterWord j}\in\AddressGraft X Y
 \leftrightarrow\WordConcat u{\LetterWord j}\in X.
\]
Für \(j=0\) und \(j=1\) verbindet
\FormulaRefAuto{GoodBinaryAddressSetFullnessAxiom}[ax] die rechten
Seiten. Äquivalenzersetzung \rLRS{} liefert die verlangte
Äquivalenz der linken Seiten. Im rechten Bildfall wird der Bildzeuge
\(u\in Y\) mit \(p=\AddressPrefix1u\) gewählt; dieselbe Rechnung
führt auf das Verzweigungsaxiom von \(Y\).
Die Bildzeugen werden durch \rEE{} entlassen, die drei Fälle durch
\rOE{} verbunden und das beliebige \(p\) durch \rUI{} gebunden.
\end{proof}


\FormulaThmDeltaK[Grund- und Pfropfschritt voller Präfixmengen]{%
  \begin{aligned}[t]
  &\text{(i)}\;\GoodAddressSet{\{\EmptyWord{\{0,1\}}\}},\\[-2pt]
  &\text{(ii)}\;\GoodAddressSet X\dsep\GoodAddressSet Y
    \vdash\GoodAddressSet{\AddressGraft X Y}.
  \end{aligned}%
}{GoodBinaryAddressSetConstructors}{%
  \DeltaRow{Mengen}{X\dsep Y\dsep P\dsep p\dsep q\dsep i}
  \DeltaRow{Prädikate}{\GoodAddressSet{\,\cdot\,}}
}
\begin{tabproofsplitwide}
  \proofpartwideindR[Grundmenge]{%
    \GoodAddressSet{\{\EmptyWord{\{0,1\}}\}}
  }{\GoodAddressSet{\{\EmptyWord{\{0,1\}}\}}}%
    [GoodBinaryAddressSetLeaf]
    \proofstepwidestar[]{\EmptyWord{\{0,1\}}\in\BinaryAddresses}{%
      \FormulaRefAuto{EmptyWordFinite}[thm]}
    \proofstepwidestar[]{\{\EmptyWord{\{0,1\}}\}\subseteq\BinaryAddresses}{%
      \FormulaRefAuto{a\in A\vdash\{a\}\subseteq A}[thm]{%
          1}}
    \proofstepwidestar[]{%
      \{\EmptyWord{\{0,1\}}\}\in\powerset(\BinaryAddresses)}{%
      \FormulaRefAuto{x\in\powerset(A)\eqvdash x\subseteq A}[thm]{%
        2}}
    \proofstepwidestar[]{\Finite{\{\EmptyWord{\{0,1\}}\}}}{%
      \FormulaRefAuto{FiniteSingleton}[thm]}
    \proofstepwidestar[]{%
      \EmptyWord{\{0,1\}}\in\{\EmptyWord{\{0,1\}}\}}{%
      \FormulaRefAuto{a\in\{a\}}[thm]}
    \proofstepwidestar[6]{%
      q\in\{\EmptyWord{\{0,1\}}\}\setminus
        \{\EmptyWord{\{0,1\}}\}}{\rA}
    \proofstepwidestar[6]{q\in\{\EmptyWord{\{0,1\}}\}}{%
      \FormulaRefAuto{x\in A\setminus B\vdash x\in A}[thm]{6}}
    \proofstepwidestar[6]{q\notin\{\EmptyWord{\{0,1\}}\}}{%
      \FormulaRefAuto{x\in A\setminus B\vdash x\notin B}[thm]{6}}
    \proofstepwidestar[6]{\bot}{\rBI{7,8}}
    \proofstepwidestar[6]{%
      \begin{aligned}[t]
      &\exists p\in\{\EmptyWord{\{0,1\}}\}\,
        \exists i\in\{0,1\}\,{}\\[-2pt]
      &\qquad q=\WordConcat p{\LetterWord i}
      \end{aligned}}{%
      \FormulaRefAuto{\bot\vdash P}[thm]{9}}
    \proofstepwidestar[]{%
      \begin{aligned}[t]
      &\forall q\in
        \{\EmptyWord{\{0,1\}}\}\setminus
          \{\EmptyWord{\{0,1\}}\}\,{}\\[-2pt]
      &\qquad\exists p\in\{\EmptyWord{\{0,1\}}\}\,
        \exists i\in\{0,1\}\,
          q=\WordConcat p{\LetterWord i}
      \end{aligned}}{\rUI{\rRI{6,10}}}
    \proofstepwidestar[]{%
      \WordConcat{\EmptyWord{\{0,1\}}}{\LetterWord0}
      \notin\{\EmptyWord{\{0,1\}}\}}{%
      \FormulaRefAuto{B27BinaryAppendNotRoot}[thm]{1,\FormulaRefAuto{a\in\{a,b\}}[thm]}}
    \proofstepwidestar[]{%
      \WordConcat{\EmptyWord{\{0,1\}}}{\LetterWord1}
      \notin\{\EmptyWord{\{0,1\}}\}}{%
      \FormulaRefAuto{B27BinaryAppendNotRoot}[thm]{1,\FormulaRefAuto{b\in\{a,b\}}[thm]}}
    \proofstepwidestar[14]{%
      p\in\{\EmptyWord{\{0,1\}}\}}{\rA}
    \proofstepwidestar[14]{p=\EmptyWord{\{0,1\}}}{%
      \FormulaRefAuto{x\in\{a\}\eqvdash x=a}[thm]{14}}
    \proofstepwidestar[14]{%
      \WordConcat p{\LetterWord0}
        \notin\{\EmptyWord{\{0,1\}}\}}{\rIE{15,12}}
    \proofstepwidestar[14]{%
      \WordConcat p{\LetterWord1}
        \notin\{\EmptyWord{\{0,1\}}\}}{\rIE{15,13}}
    \proofstepwidestar[14]{%
      \begin{aligned}[t]
      &\neg\bigl(\WordConcat p{\LetterWord0}
        \in\{\EmptyWord{\{0,1\}}\}\bigr)\land{}\\[-2pt]
      &\neg\bigl(\WordConcat p{\LetterWord1}
        \in\{\EmptyWord{\{0,1\}}\}\bigr)
      \end{aligned}}{\rAI{16,17}}
    \proofstepwidestar[14]{%
      \begin{aligned}[t]
      &\Bigl(\WordConcat p{\LetterWord0}
        \in\{\EmptyWord{\{0,1\}}\}\land{}\\[-2pt]
      &\quad\WordConcat p{\LetterWord1}
        \in\{\EmptyWord{\{0,1\}}\}\Bigr)\lor{}\\[-2pt]
      &\Bigl(\neg\bigl(\WordConcat p{\LetterWord0}
        \in\{\EmptyWord{\{0,1\}}\}\bigr)\land{}\\[-2pt]
      &\quad\neg\bigl(\WordConcat p{\LetterWord1}
        \in\{\EmptyWord{\{0,1\}}\}\bigr)\Bigr)
      \end{aligned}}{\rOIb{18}}
    \proofstepwidestar[14]{%
      \WordConcat p{\LetterWord0}\in\{\EmptyWord{\{0,1\}}\}
      \leftrightarrow
      \WordConcat p{\LetterWord1}\in\{\EmptyWord{\{0,1\}}\}}{%
      \FormulaRefAuto{P \leftrightarrow Q \eqvdash
        (P \land Q) \lor (\neg P \land \neg Q)}[thm]{19}}
    \proofstepwidestar[]{%
      \forall p\in\{\EmptyWord{\{0,1\}}\}\,\Bigl(
      \WordConcat p{\LetterWord0}\in\{\EmptyWord{\{0,1\}}\}
      \leftrightarrow
      \WordConcat p{\LetterWord1}\in\{\EmptyWord{\{0,1\}}\}\Bigr)}{%
      \rUI{\rRI{14,20}}}
    \proofstepwidestar[]{%
      \GoodAddressSet{\{\EmptyWord{\{0,1\}}\}}}{%
      \FormulaRefAuto{GoodBinaryAddressSetCriterion}[thm]{3,4,5,11,21}}
  \closeproofpartwideindR

  \proofpartwideindR[Pfropfung]{%
    \GoodAddressSet X\dsep\GoodAddressSet Y
      \vdash\GoodAddressSet{\AddressGraft X Y}
  }{%
    \GoodAddressSet X\dsep\GoodAddressSet Y
      \vdash\GoodAddressSet{\AddressGraft X Y}
  }[GoodBinaryAddressSetGraft]
    \proofstepwidestar[1]{\GoodAddressSet X}{\rA}
    \proofstepwidestar[2]{\GoodAddressSet Y}{\rA}
    \proofstepwidestar[1,2]{X,Y\in\powerset(\BinaryAddresses)}{%
      \FormulaRefAuto{GoodBinaryAddressSetCriterion}[thm]{1,2}}
    \proofstepwidestar[1,2]{\Finite X\land\Finite Y}{%
      \FormulaRefAuto{GoodBinaryAddressSetCriterion}[thm]{1,2}}
    \proofstepwidestar[]{0\in\{0,1\}}{%
      \FormulaRefAuto{a\in\{a,b\}}[thm]}
    \proofstepwidestar[]{1\in\{0,1\}}{%
      \FormulaRefAuto{b\in\{a,b\}}[thm]}
    \proofstepwidestar[]{%
      \AddressPrefixMap{0},\AddressPrefixMap{1}\colon
      \BinaryAddresses\to\BinaryAddresses}{%
      \FormulaRefAuto{BinaryAddressPrefixDef}[def]{%
        5,%
        6}}
    \proofstepwidestar[1,2]{%
      \Finite{\AddressPrefixMap{0}[X]}\land
      \Finite{\AddressPrefixMap{1}[Y]}}{%
      \FormulaRefAuto{FiniteImage}[thm]{3,7,4}}
    \proofstepwidestar[]{\Finite{\{\EmptyWord{\{0,1\}}\}}}{%
      \FormulaRefAuto{FiniteSingleton}[thm]}
    \proofstepwidestar[1,2]{%
      \Finite{\bigl(\AddressPrefixMap{0}[X]\cup
        \AddressPrefixMap{1}[Y]\bigr)}}{%
      \FormulaRefAuto{FiniteUnion}[thm]{\rAEa{8},\rAEb{8}}}
    \proofstepwidestar[1,2]{\Finite{\AddressGraft X Y}}{%
      \FormulaRefAuto{FiniteUnion}[thm]{9,10}}
    \proofstepwidestar[1,2]{%
      \AddressGraft X Y\in\powerset(\BinaryAddresses)}{%
      \FormulaRefAuto{BinaryAddressGraftBinaryOperation}[thm]{3}}
    \proofstepwidestar[1,2]{\EmptyWord{\{0,1\}}\in\{\EmptyWord{\{0,1\}}\}}{%
      \FormulaRefAuto{a\in\{a\}}[thm]}
    \proofstepwidestar[1,2]{\EmptyWord{\{0,1\}}\in\{\EmptyWord{\{0,1\}}\}\cup\AddressPrefixMap{0}[X]}{%
      \FormulaRefAuto{x\in A\vdash x\in A\cup B}[thm]{%
            13}}
    \proofstepwidestar[1,2]{\EmptyWord{\{0,1\}}\in(\{\EmptyWord{\{0,1\}}\}\cup\AddressPrefixMap{0}[X])\cup\AddressPrefixMap{1}[Y]}{%
      \FormulaRefAuto{x\in A\vdash x\in A\cup B}[thm]{%
          14}}
    \proofstepwidestar[1,2]{%
      \EmptyWord{\{0,1\}}\in\AddressGraft X Y}{%
      \FormulaRefAuto{BinaryAddressGraftDef}[def]{3,
        15}}
    \proofstepwidestarforcebreak[1,2]{%
      \begin{aligned}[t]
      &\forall q\in\AddressGraft X Y\setminus
        \{\EmptyWord{\{0,1\}}\}\,\exists p\in\AddressGraft X Y\,{}\\[-2pt]
      &\qquad\exists i\in\{0,1\}\,
        q=\WordConcat p{\LetterWord i}
      \end{aligned}}{\FormulaRefAuto{B27GraftPrefixClosure}[thm]{1,2}}
    \proofstepwidestarforcebreak[1,2]{%
      \begin{aligned}[t]
      &\forall p\in\AddressGraft X Y\,\Bigl(
        \WordConcat p{\LetterWord0}\in\AddressGraft X Y\\[-2pt]
      &\qquad\leftrightarrow
        \WordConcat p{\LetterWord1}\in\AddressGraft X Y\Bigr)
      \end{aligned}}{\FormulaRefAuto{B27GraftFullness}[thm]{1,2}}
    \proofstepwidestar[1,2]{\GoodAddressSet{\AddressGraft X Y}}{%
      \FormulaRefAuto{GoodBinaryAddressSetCriterion}[thm]{12,11,16,17,18}}
  \closeproofpartwideindR
\end{tabproofsplitwide}

\FormulaDefDeltaK[Konstante Blattadressabbildung]{%
  \begin{aligned}[t]
  &c_A\colon A\to\powerset(\BinaryAddresses),\\[-2pt]
  &\forall a\in A\,
    \bigl(c_A(a)\coloneqq\{\EmptyWord{\{0,1\}}\}\bigr)
  \end{aligned}%
}{TreePositionLeafMapDef}{%
  \DeltaRow{Mengen}{A\dsep a\dsep\BinaryAddresses}
  \DeltaRow{Funktionen}{c_A\dsep F}
  \DeltaRow{\textbf{Neues Symbol}}{}
  \DeltaPrem{Blattadressabbildung}{c_A}
}
\begin{tabproofwide}
  \proofstepwidestar[]{\EmptyWord{\{0,1\}}\in\BinaryAddresses}{%
    \FormulaRefAuto{EmptyWordFinite}[thm]}
  \proofstepwidestar[]{\{\EmptyWord{\{0,1\}}\}\subseteq\BinaryAddresses}{%
    \FormulaRefAuto{a\in A\vdash\{a\}\subseteq A}[thm]{1}}
  \proofstepwidestar[]{\{\EmptyWord{\{0,1\}}\}\in\powerset(\BinaryAddresses)}{%
    \FormulaRefAuto{x\in\powerset(A)\eqvdash x\subseteq A}[thm]{2}}
  \proofstepwidestar[4]{a\in A}{\rA}
  \proofstepwidestar[4]{%
    a\in A\land\{\EmptyWord{\{0,1\}}\}\in\powerset(\BinaryAddresses)}{\rAI{4,3}}
  \proofstepwidestar[4]{\{\EmptyWord{\{0,1\}}\}\in\powerset(\BinaryAddresses)}{\rAEb{5}}
  \proofstepwidestar[]{%
    a\in A\rightarrow\{\EmptyWord{\{0,1\}}\}\in\powerset(\BinaryAddresses)}{\rRI{4,6}}
  \proofstepwidestar[]{%
    \forall a\in A\,\{\EmptyWord{\{0,1\}}\}\in\powerset(\BinaryAddresses)}{\rUI{7}}
  \proofstepwidestar[]{%
    \begin{aligned}[t]
      &\exists!F\,\Bigl(F\colon A\to\powerset(\BinaryAddresses)\land{}\\[-2pt]
      &\quad\forall a\in A\,F(a)=\{\EmptyWord{\{0,1\}}\}\Bigr)
    \end{aligned}}{\FormulaRefAuto{TypedFunctionDefinitionPrinciple}[thm]{8}}
\end{tabproofwide}

\Needspace{10\baselineskip}
\FormulaThmDeltaK[Existenz und Eindeutigkeit der Positionsabbildung]{%
  \begin{aligned}[t]
  &\exists!P\colon\BracketTreeSet A\to\powerset(\BinaryAddresses)\,\Bigl(\\[-2pt]
  &\quad\forall a\in A\,
      P(\LeafTree A a)=\{\EmptyWord{\{0,1\}}\}\land{}\\[-2pt]
  &\quad\forall S,T\in\BracketTreeSet A\,
      P(\NodeTree A S T)=\AddressGraft{P(S)}{P(T)}\Bigr).
  \end{aligned}%
}{TreePositionMapExistenceUnique}{%
  \DeltaRow{Mengen}{A\dsep a\dsep S\dsep T\dsep\BinaryAddresses}
  \DeltaRow{Funktionen}{P\dsep Q\dsep R\dsep c_A\dsep\Gamma}
  \DeltaRow{Prädikate}{\Phi}
}
% Beweis zu TreePositionMapExistenceUnique; der alte Satz bleibt unverändert.
Für diesen Beweis werden
\[
  R\coloneqq\TreeRecursor{A}{\powerset(\BinaryAddresses)}{c_A}
    {\BinaryOpFunction{\powerset(\BinaryAddresses)}{\Gamma}}
\]
und die folgende Eigenschaft abgekürzt:
\[
  \begin{aligned}
    \Phi(P)\coloneqq{}&
      P\colon\BracketTreeSet A\to\powerset(\BinaryAddresses)\land\Bigl(\\[-2pt]
    &\forall a\in A\,P(\LeafTree A a)=\{\EmptyWord{\{0,1\}}\}\land{}\\[-2pt]
    &\forall S,T\in\BracketTreeSet A\,
      P(\NodeTree A S T)=\AddressGraft{P(S)}{P(T)}\Bigr).
  \end{aligned}
\]
Der Rekursor liefert den Zeugen. Die Identifikationsregel zeigt anschließend,
dass jede Funktion mit diesen Blatt- und Knotenregeln mit ihm übereinstimmt.

\begin{tabproofsplitwide}
  \proofpartwideindR[Identifikation eines Kandidaten]{%
    \Phi(P)\vdash P=R
  }{\Phi(P)\vdash P=R}[TreePositionRecursorIdentification]
    \proofstepwidestar[1]{\Phi(P)}{\rA}
    \proofstepwidestar[1]{%
      P\colon\BracketTreeSet A\to\powerset(\BinaryAddresses)}{\rAEa{1}}
    \proofstepwidestar[1]{%
      \forall a\in A\,P(\LeafTree A a)=\{\EmptyWord{\{0,1\}}\}}{%
      \rAEa{\rAEb{1}}}
    \proofstepwidestar[1]{%
      \forall S,T\in\BracketTreeSet A\,
        P(\NodeTree A S T)=\AddressGraft{P(S)}{P(T)}}{\rAEb{\rAEb{1}}}
    \proofstepwidestar[]{c_A\colon A\to\powerset(\BinaryAddresses)}{%
      \FormulaRefAuto{TreePositionLeafMapDef}[def]}
    \proofstepwidestar[]{%
      \mathsf{BinOp}\bigl(\powerset(\BinaryAddresses),\Gamma\bigr)}{%
      \FormulaRefAuto{BinaryAddressGraftBinaryOperation}[thm]}
    \proofstepwidestar[7]{a\in A}{\rA}
    \proofstepwidestar[1,7]{%
      P(\LeafTree A a)=\{\EmptyWord{\{0,1\}}\}}{\rRE{7,\rUE{3}}}
    \proofstepwidestar[7]{c_A(a)=\{\EmptyWord{\{0,1\}}\}}{%
      \FormulaRefAuto{TreePositionLeafMapDef}[def]{7}}
    \proofstepwidestar[1,7]{P(\LeafTree A a)=c_A(a)}{%
      \FormulaRefAuto{a=b\dsep c=b\vdash a=c}[thm]{8,9}}
    \proofstepwidestar[1]{\forall a\in A\,P(\LeafTree A a)=c_A(a)}{%
      \rUI{\rRI{7,10}}}
    \proofstepwidestar[1]{P=R}{%
      \FormulaRefAuto{TreeRecursorBinaryOperationIdentification}[thm]{5,6,2,11,4}}
  \closeproofpartwideindR

  \proofpartwide{Existenz und Eindeutigkeit}
    \proofstepwidestar[]{c_A\colon A\to\powerset(\BinaryAddresses)}{%
      \FormulaRefAuto{TreePositionLeafMapDef}[def]}
    \proofstepwidestar[]{%
      \mathsf{BinOp}\bigl(\powerset(\BinaryAddresses),\Gamma\bigr)}{%
      \FormulaRefAuto{BinaryAddressGraftBinaryOperation}[thm]}
    \proofstepwidestar[]{%
      R\colon\BracketTreeSet A\to\powerset(\BinaryAddresses)}{%
      \FormulaRefAuto{TreeRecursorBinaryOperationTyping}[thm]{1,2}}
    \proofstepwidestar[4]{a\in A}{\rA}
    \proofstepwidestar[4]{R(\LeafTree A a)=c_A(a)}{%
      \FormulaRefAuto{TreeRecursorBinaryOperationLeafEquation}[thm]{1,2,4}}
    \proofstepwidestar[4]{c_A(a)=\{\EmptyWord{\{0,1\}}\}}{%
      \FormulaRefAuto{TreePositionLeafMapDef}[def]{4}}
    \proofstepwidestar[4]{R(\LeafTree A a)=\{\EmptyWord{\{0,1\}}\}}{%
      \FormulaRefAuto{a=b\dsep b=c\vdash a=c}[thm]{5,6}}
    \proofstepwidestar[]{%
      \forall a\in A\,R(\LeafTree A a)=\{\EmptyWord{\{0,1\}}\}}{%
      \rUI{\rRI{4,7}}}
    \proofstepwidestar[9]{S\in\BracketTreeSet A}{\rA}
    \proofstepwidestar[10]{T\in\BracketTreeSet A}{\rA}
    \proofstepwidestar[9,10]{%
      R(\NodeTree A S T)=\AddressGraft{R(S)}{R(T)}}{%
      \FormulaRefAuto{TreeRecursorBinaryOperationNodeEquation}[thm]{1,2,\rAI{9,10}}}
    \proofstepwidestar[]{%
      \forall S,T\in\BracketTreeSet A\,
        R(\NodeTree A S T)=\AddressGraft{R(S)}{R(T)}}{%
      \rUI{\rRI{9,\rUI{\rRI{10,11}}}}}
    \proofstepwidestar[]{\Phi(R)}{\rAI{3,\rAI{8,12}}}
    \proofstepwidestar[]{\exists P\,\Phi(P)}{\rEI{13}}
    \proofstepwidestar[15]{\Phi(P)}{\rA}
    \proofstepwidestar[16]{\Phi(Q)}{\rA}
    \proofstepwidestar[15]{P=R}{%
      \FormulaRefAuto{TreePositionRecursorIdentification}[thm]{15}}
    \proofstepwidestar[16]{Q=R}{%
      \FormulaRefAuto{TreePositionRecursorIdentification}[thm]{16}}
    \proofstepwidestar[15,16]{P=Q}{%
      \FormulaRefAuto{a=b\dsep c=b\vdash a=c}[thm]{17,18}}
    \proofstepwidestar[]{\exists!P\,\Phi(P)}{\UEI{14,15,16,19}}
  \closeproofpartwide
\end{tabproofsplitwide}


\FormulaDefDeltaK[Positionsmenge eines Baumcodes]{%
  \begin{aligned}[t]
    &\operatorname{Pos}_A\coloneqq{}\\[-2pt]
    &\qquad \TreeRecursor{A}{\powerset(\BinaryAddresses)}{c_A}{\BinaryOpFunction{\powerset(\BinaryAddresses)}{\Gamma}}
  \end{aligned}%
}{TreePositionSetDef}{%
  \DeltaRow{Mengen}{A\dsep R\dsep a\dsep S\dsep T\dsep\BinaryAddresses}
  \DeltaRow{Funktionen}{\operatorname{Pos}_A}
  \DeltaRow{\textbf{Neues Symbol}}{}
  \DeltaPrem{Positionsmenge}{\TreePositions A R}
}

\Needspace{12\baselineskip}
\FormulaThmDeltaK[Rekursionsgleichungen der Positionsmenge]{%
  \begin{aligned}[t]
  &\text{(i)}\;\operatorname{Pos}_A\colon
    \BracketTreeSet A\to\powerset(\BinaryAddresses),\\[-2pt]
  &\text{(ii)}\;a\in A\vdash
    \TreePositions A{\LeafTree A a}=\{\EmptyWord{\{0,1\}}\},\\[-2pt]
  &\text{(iii)}\;S,T\in\BracketTreeSet A\vdash
    \TreePositions A{\NodeTree A S T}=
      \AddressGraft{\TreePositions A S}{\TreePositions A T}.
  \end{aligned}%
}{TreePositionSetEquations}{%
  \DeltaRow{Mengen}{A\dsep a\dsep S\dsep T\dsep\BinaryAddresses}
  \DeltaRow{Funktionen}{\operatorname{Pos}_A\dsep R}
}
Im folgenden Beweis bezeichnet \(R\) die rechte Seite von
\FormulaRefAuto{TreePositionSetDef}[def]. Die Blattregel ersetzt die
Beschriftung; die Knotenregel verknüpft die beiden Teilbaumergebnisse.
\begin{tabproofsplitwide}
  \proofpartwideindR[Typisierung]{%
    \operatorname{Pos}_A\colon\BracketTreeSet A\to\powerset(\BinaryAddresses)
  }{\operatorname{Pos}_A\colon\BracketTreeSet A\to\powerset(\BinaryAddresses)}%
    [TreePositionSetTyping]
    \proofstepwidestar[]{%
      c_A\colon A\to \powerset(\BinaryAddresses)}{%
      \FormulaRefAuto{TreePositionLeafMapDef}[def]}
    \proofstepwidestar[]{%
      \mathsf{BinOp}\bigl(\powerset(\BinaryAddresses),\Gamma\bigr)}{%
      \FormulaRefAuto{BinaryAddressGraftBinaryOperation}[thm]}
    \proofstepwidestar[]{%
      \operatorname{Pos}_A=R}{%
      \FormulaRefAuto{TreePositionSetDef}[def]}
    \proofstepwidestar[]{%
      R=\operatorname{Pos}_A}{%
      \FormulaRefAuto{a=b\vdash b=a}[thm]{3}}
    \proofstepwidestar[]{%
      R\colon\BracketTreeSet A\to \powerset(\BinaryAddresses)}{%
      \FormulaRefAuto{TreeRecursorBinaryOperationTyping}[thm]{1,2}}
    \proofstepwidestar[]{%
      \operatorname{Pos}_A\colon\BracketTreeSet A\to \powerset(\BinaryAddresses)}{%
      \rIE{4,5}}
  \closeproofpartwideindR
  \proofpartwideindR[Blattgleichung]{%
    a\in A\vdash
      \TreePositions A{\LeafTree A a}=\{\EmptyWord{\{0,1\}}\}
  }{a\in A\vdash
      \TreePositions A{\LeafTree A a}=\{\EmptyWord{\{0,1\}}\}}%
    [TreePositionSetLeafEquation]
    \proofstepwidestar[1]{%
      a\in A}{%
      \rA}
    \proofstepwidestar[]{%
      c_A\colon A\to \powerset(\BinaryAddresses)}{%
      \FormulaRefAuto{TreePositionLeafMapDef}[def]}
    \proofstepwidestar[]{%
      \mathsf{BinOp}\bigl(\powerset(\BinaryAddresses),\Gamma\bigr)}{%
      \FormulaRefAuto{BinaryAddressGraftBinaryOperation}[thm]}
    \proofstepwidestar[]{%
      \operatorname{Pos}_A=R}{%
      \FormulaRefAuto{TreePositionSetDef}[def]}
    \proofstepwidestar[]{%
      R=\operatorname{Pos}_A}{%
      \FormulaRefAuto{a=b\vdash b=a}[thm]{4}}
    \proofstepwidestar[1]{%
      R(\LeafTree A a)=c_A(a)}{%
      \FormulaRefAuto{TreeRecursorBinaryOperationLeafEquation}[thm]{2,3,1}}
    \proofstepwidestar[1]{%
      c_A(a)=\{\EmptyWord{\{0,1\}}\}}{%
      \FormulaRefAuto{TreePositionLeafMapDef}[def]{1}}
    \proofstepwidestar[1]{%
      \operatorname{Pos}_A(\LeafTree A a)=c_A(a)}{%
      \rIE{5,6}}
    \proofstepwidestar[1]{%
      \operatorname{Pos}_A(\LeafTree A a)=\{\EmptyWord{\{0,1\}}\}}{%
      \FormulaRefAuto{a=b\dsep b=c\vdash a=c}[thm]{8,7}}
  \closeproofpartwideindR
  \proofpartwideindR[Knotengleichung]{%
    S,T\in\BracketTreeSet A\vdash
      \TreePositions A{\NodeTree A S T}=
        \AddressGraft{\TreePositions A S}{\TreePositions A T}
  }{S,T\in\BracketTreeSet A\vdash
      \TreePositions A{\NodeTree A S T}=
        \AddressGraft{\TreePositions A S}{\TreePositions A T}}%
    [TreePositionSetNodeEquation]
    \proofstepwidestar[1]{%
      S\in\BracketTreeSet A}{%
      \rA}
    \proofstepwidestar[2]{%
      T\in\BracketTreeSet A}{%
      \rA}
    \proofstepwidestar[]{%
      c_A\colon A\to \powerset(\BinaryAddresses)}{%
      \FormulaRefAuto{TreePositionLeafMapDef}[def]}
    \proofstepwidestar[]{%
      \mathsf{BinOp}\bigl(\powerset(\BinaryAddresses),\Gamma\bigr)}{%
      \FormulaRefAuto{BinaryAddressGraftBinaryOperation}[thm]}
    \proofstepwidestar[]{%
      \operatorname{Pos}_A=R}{%
      \FormulaRefAuto{TreePositionSetDef}[def]}
    \proofstepwidestar[]{%
      R=\operatorname{Pos}_A}{%
      \FormulaRefAuto{a=b\vdash b=a}[thm]{5}}
    \proofstepwidestar[1,2]{%
      R(\NodeTree A S T)=\AddressGraft{R(S)}{R(T)}}{%
      \FormulaRefAuto{TreeRecursorBinaryOperationNodeEquation}[thm]{3,4,\rAI{1,2}}}
    \proofstepwidestar[1,2]{%
      \begin{aligned}[t]&\operatorname{Pos}_A(\NodeTree A S T)\\[-2pt]
      &\quad=\AddressGraft{\operatorname{Pos}_A(S)}{\operatorname{Pos}_A(T)}\end{aligned}}{%
      \rIE{6,7}}
  \closeproofpartwideindR
\end{tabproofsplitwide}

\FormulaThmDeltaK[Positionsmengen sind volle endliche Präfixmengen]{%
  R\in\BracketTreeSet A\vdash\GoodAddressSet{\TreePositions A R}%
}{TreePositionSetIsGood}{%
  \DeltaRow{Mengen}{A\dsep R\dsep a\dsep S\dsep T}
  \DeltaRow{Prädikate}{\GoodAddressSet{\,\cdot\,}}
}
\begin{tabproofsplitwide}
  \proofpartwideindR[Blatt]{%
    \forall a\in A\,\GoodAddressSet{\TreePositions A{\LeafTree A a}}
  }{\forall a\in A\,\GoodAddressSet{\TreePositions A{\LeafTree A a}}}%
    [TreePositionSetGoodLeaf]
    \proofstepwidestar[1]{a\in A}{\rA}
    \proofstepwidestar[1]{%
      \TreePositions A{\LeafTree A a}=\{\EmptyWord{\{0,1\}}\}}{%
      \FormulaRefAuto{TreePositionSetLeafEquation}[thm]{1}}
    \proofstepwidestar[]{\GoodAddressSet{\{\EmptyWord{\{0,1\}}\}}}{%
      \FormulaRefAuto{GoodBinaryAddressSetLeaf}[thm]}
    \proofstepwidestar[1]{%
      \GoodAddressSet{\TreePositions A{\LeafTree A a}}}{\rIE{2,3}}
    \proofstepwidestar[]{%
      \forall a\in A\,\GoodAddressSet{\TreePositions A{\LeafTree A a}}}{%
      \rUI{\rRI{1,4}}}
  \closeproofpartwideindR

  \proofpartwideindR[Knoten]{%
    \begin{aligned}[t]
    &\forall S,T\in\BracketTreeSet A\,\Bigl(
      \GoodAddressSet{\TreePositions A S}\land
      \GoodAddressSet{\TreePositions A T}\\[-2pt]
    &\qquad\rightarrow
      \GoodAddressSet{\TreePositions A{\NodeTree A S T}}\Bigr)
    \end{aligned}
  }{%
    \begin{aligned}[t]
    &\forall S,T\in\BracketTreeSet A\,\Bigl(
      \GoodAddressSet{\TreePositions A S}\land
      \GoodAddressSet{\TreePositions A T}\rightarrow
      \GoodAddressSet{\TreePositions A{\NodeTree A S T}}\Bigr)
    \end{aligned}
  }[TreePositionSetGoodNode]
    \proofstepwidestar[1]{S,T\in\BracketTreeSet A}{\rA}
    \proofstepwidestar[2]{%
      \GoodAddressSet{\TreePositions A S}\land
      \GoodAddressSet{\TreePositions A T}}{\rA}
    \proofstepwidestar[1]{%
      \TreePositions A{\NodeTree A S T}=
      \AddressGraft{\TreePositions A S}{\TreePositions A T}}{%
      \FormulaRefAuto{TreePositionSetNodeEquation}[thm]{1}}
    \proofstepwidestar[2]{%
      \GoodAddressSet{\AddressGraft{\TreePositions A S}{\TreePositions A T}}}{%
      \FormulaRefAuto{GoodBinaryAddressSetGraft}[thm]{\rAEa{2},\rAEb{2}}}
    \proofstepwidestar[1,2]{%
      \GoodAddressSet{\TreePositions A{\NodeTree A S T}}}{\rIE{3,4}}
    \proofstepwidestar[1]{%
      \GoodAddressSet{\TreePositions A S}\land
      \GoodAddressSet{\TreePositions A T}\rightarrow
      \GoodAddressSet{\TreePositions A{\NodeTree A S T}}}{\rRI{2,5}}
    \proofstepwidestar[]{\forall S,T\in\BracketTreeSet A\,(\cdots)}{%
      \rUI{\rUI{\rRI{1,6}}}}
  \closeproofpartwideindR

  \proofpartwide{Strukturinduktion}
    \proofstepwidestar[1]{R\in\BracketTreeSet A}{\rA}
    \proofstepwidestar[]{\forall a\in A\,\GoodAddressSet{\TreePositions A{\LeafTree A a}}}{%
      \FormulaRefAuto{TreePositionSetGoodLeaf}[thm]}
    \proofstepwidestar[]{\forall S,T\in\BracketTreeSet A\,\bigl(\GoodAddressSet{\TreePositions A S}\land\GoodAddressSet{\TreePositions A T}\rightarrow\GoodAddressSet{\TreePositions A{\NodeTree A S T}}\bigr)}{%
      \FormulaRefAuto{TreePositionSetGoodNode}[thm]}
    \proofstepwidestar[]{%
      \forall R\in\BracketTreeSet A\,
        \GoodAddressSet{\TreePositions A R}}{%
      \FormulaRefAuto{FullPlanarBinaryTreeStructuralInduction}[thm]{%
        2,
        3}}
    \proofstepwidestar[1]{\GoodAddressSet{\TreePositions A R}}{\rRE{1,4}}
  \closeproofpartwide
\end{tabproofsplitwide}


\FormulaDefDeltaK[Adressgraph und geordnete Kindoperatoren]{%
  \GoodAddressSet P\vdash
  \begin{aligned}[t]
  \AddressEdges P&\coloneqq
    \left\{(p,q)\in P\times P\;\middle|\;
      \substack{\exists i\in\{0,1\}:\\[-2pt]
        q=\WordConcat p{\LetterWord i}\lor
        p=\WordConcat q{\LetterWord i}}\right\},\\[-2pt]
  \AddressInner P&\coloneqq
    \bigl\{p\in P\bigm|\WordConcat p{\LetterWord0}\in P\bigr\},\\[-2pt]
  \AddressLeft P&\coloneqq
    \bigl\{(p,\WordConcat p{\LetterWord0})\bigm|
      p\in\AddressInner P\bigr\},\\[-2pt]
  \AddressRight P&\coloneqq
    \bigl\{(p,\WordConcat p{\LetterWord1})\bigm|
      p\in\AddressInner P\bigr\}.
  \end{aligned}%
}{BinaryAddressGraphDef}{%
  \DeltaRow{Mengen}{P\dsep p\dsep q\dsep i\dsep
    \AddressEdges P\dsep\AddressInner P}
  \DeltaRow{Funktionen}{\AddressLeft P\dsep\AddressRight P}
  \DeltaPrem{Volle Präfixmenge}{\GoodAddressSet P}
  \DeltaRow{\textbf{Neue Symbole}}{}
  \DeltaPrem{Adresskanten}{\AddressEdges P}
  \DeltaPrem{Innere Adressmenge}{\AddressInner P}
  \DeltaPrem{Linkes Adresskind}{\AddressLeft P}
  \DeltaPrem{Rechtes Adresskind}{\AddressRight P}
}

\FormulaDefDeltaK[Die Menge der beiden unmittelbaren Adressfortsetzungen]{%
  \{\WordConcat p{\LetterWord i}\mid i\in\{0,1\}\}\coloneqq
  \iota S\,\forall q\,\Bigl(q\in S\leftrightarrow
    \exists i\in\{0,1\}\,(q=\WordConcat p{\LetterWord i})\Bigr)%
}{B27BinaryChildTermSetDef}{%
  \DeltaRow{Mengen}{p\dsep q\dsep i\dsep S}
  \DeltaPrem{Adresse}{p\in\BinaryAddresses}
  \DeltaRow{Eindeutige Existenz}{%
    \begin{aligned}[t]
      &\exists!S\,\forall q\,\Bigl(q\in S\leftrightarrow{}\\[-2pt]
      &\quad\exists i\in\{0,1\}\,(q=\WordConcat p{\LetterWord i})\Bigr)
    \end{aligned}}
    [\FormulaRefAuto{TermImageSetUniqueExists}[thm]]
}

\FormulaThmDeltaK[Adressfortsetzungen bilden eine Paarmenge]{%
  p\in\BinaryAddresses\vdash
  \{\WordConcat p{\LetterWord i}\mid i\in\{0,1\}\}
    =\{\WordConcat p{\LetterWord0},\WordConcat p{\LetterWord1}\}%
}{B27BinaryChildTermSetPair}{%
  \DeltaRow{Mengen}{p\dsep q\dsep i}
}
\begin{tabproofwide}
  \proofstepwidestar[1]{p\in\BinaryAddresses}{\rA}
  \proofstepwidestar[1]{q\in\{\WordConcat p{\LetterWord i}\mid i\in\{0,1\}\}
    \leftrightarrow\exists i\in\{0,1\}\,(q=\WordConcat p{\LetterWord i})}{%
    \FormulaRefAuto{B27BinaryChildTermSetDef}[def]{1}}
  \proofstepwidestar[3]{q\in\{\WordConcat p{\LetterWord i}\mid i\in\{0,1\}\}}{\rA}
  \proofstepwidestar[1,3]{\exists i\in\{0,1\}\,(q=\WordConcat p{\LetterWord i})}{\rRE{\rLREa{2},3}}
  \proofstepwidestar[1,3]{q=\WordConcat p{\LetterWord0}\lor q=\WordConcat p{\LetterWord1}}{%
    \FormulaRefAuto{\exists x\in\{a,b\}P(x)\vdash P(a)\lor P(b)}[thm]{4}}
  \proofstepwidestar[1,3]{q\in\{\WordConcat p{\LetterWord0},\WordConcat p{\LetterWord1}\}}{%
    \FormulaRefAuto{x\in\{a,b\}\eqvdash(x=a\lor x=b)}[thm]{5}}
  \proofstepwidestar[7]{q\in\{\WordConcat p{\LetterWord0},\WordConcat p{\LetterWord1}\}}{\rA}
  \proofstepwidestar[7]{q=\WordConcat p{\LetterWord0}\lor q=\WordConcat p{\LetterWord1}}{%
    \FormulaRefAuto{x\in\{a,b\}\eqvdash(x=a\lor x=b)}[thm]{7}}
  \proofstepwidestar[9]{q=\WordConcat p{\LetterWord0}}{\rA}
  \proofstepwidestar[9]{\exists i\in\{0,1\}\,(q=\WordConcat p{\LetterWord i})}{%
    \rEI{\rAI{\FormulaRefAuto{a\in\{a,b\}}[thm],9}}}
  \proofstepwidestar[11]{q=\WordConcat p{\LetterWord1}}{\rA}
  \proofstepwidestar[11]{\exists i\in\{0,1\}\,(q=\WordConcat p{\LetterWord i})}{%
    \rEI{\rAI{\FormulaRefAuto{b\in\{a,b\}}[thm],11}}}
  \proofstepwidestar[7]{\exists i\in\{0,1\}\,(q=\WordConcat p{\LetterWord i})}{\rOE{8,9,10,11,12}}
  \proofstepwidestar[1,7]{q\in\{\WordConcat p{\LetterWord i}\mid i\in\{0,1\}\}}{\rRE{\rLREb{2},13}}
  \proofstepwidestar[1]{q\in\{\WordConcat p{\LetterWord i}\mid i\in\{0,1\}\}
    \leftrightarrow q\in\{\WordConcat p{\LetterWord0},\WordConcat p{\LetterWord1}\}}{%
    \rLRI{\rRI{3,6},\rRI{7,14}}}
  \proofstepwidestar[1]{\{\WordConcat p{\LetterWord i}\mid i\in\{0,1\}\}
    =\{\WordConcat p{\LetterWord0},\WordConcat p{\LetterWord1}\}}{%
    \FormulaRefAuto{\forall x\,(x\in A\leftrightarrow x\in B)\vdash A=B}[ax]{\rUI{15}}}
\end{tabproofwide}

\Needspace{12\baselineskip}
\subsection{Grundregeln für Adressen und Pfade}

Eine Adresskante fügt ein Bit an oder entfernt das letzte Bit. Ein einfacher
Pfad aus der Wurzel kann jedoch keinen Rückschritt machen: Er würde dabei
einen bereits besuchten Knoten wiederholen. Daher zählt die Wortlänge seine
Schritte, und der Rumpf liefert jeweils den vorherigen Knoten. Die allgemeinen Index- und Pfadregeln stehen in Band~10 und Band~26
bereit; hier werden sie auf die Wortadressen angewendet.

% Formale Hilfssätze für die Wurzelwege des Adressgraphen.
% Der Graph wird ausdrücklich vorausgesetzt; seine Schlichtheit wird erst
% im anschließenden Struktursatz aus der Adressdefinition gewonnen.
\newcommand{\BXXVIIAddressRoot}{\EmptyWord{\{0,1\}}}
\newcommand{\BXXVIIAddressPath}[2]{\DirectedPath{P}{\AddressEdges P}{#1}{#2}}
\newcommand{\BXXVIIAddressForward}[2]{%
  \forall i\in\NatSeg{#2}\,\exists a\in\{0,1\}\,
  #1(i+1)=\WordConcat{#1(i)}{\LetterWord a}}

\Needspace{10\baselineskip}
\FormulaThmDeltaK[Adresszugehörigkeit]{%
  \GoodAddressSet P\dsep q\in P\vdash q\in\FiniteWords{\{0,1\}}
}{B27AddressWordTyping}{\DeltaRow{Mengen}{P\dsep q}}
\begin{tabproofwide}
  \proofstepwidestar[1]{\GoodAddressSet P}{\rA}
  \proofstepwidestar[2]{q\in P}{\rA}
  \proofstepwidestar[1]{P\in\powerset(\BinaryAddresses)}{\FormulaRefAuto{GoodBinaryAddressSetCarrierAxiom}[ax]{1}}
  \proofstepwidestar[1]{P\subseteq\BinaryAddresses}{\FormulaRefAuto{x\in\powerset(A)\eqvdash x\subseteq A}[thm]{3}}
  \proofstepwidestar[1,2]{q\in\BinaryAddresses}{\FormulaRefAuto{A\subseteq B\dsep x\in A\vdash x\in B}[thm]{4,2}}
  \proofstepwidestar[1,2]{q\in\FiniteWords{\{0,1\}}}{\FormulaRefAuto{BinaryAddressWordSetDef}[def]{5}}
\end{tabproofwide}

\Needspace{10\baselineskip}
\FormulaThmDeltaK[Ein angehängter Buchstabe ergibt kein Leerwort]{%
  u\in\FiniteWords A\dsep a\in A\vdash
  \WordConcat u{\LetterWord a}\neq\EmptyWord A
}{B27WordAppendNotEmpty}{\DeltaRow{Mengen}{A\dsep u\dsep a}}
\begin{tabproofwide}
  \proofstepwidestar[1]{u\in\FiniteWords A}{\rA}
  \proofstepwidestar[2]{a\in A}{\rA}
  \proofstepwidestar[1]{\WordLength u\in\mathbb N}{\FormulaRefAuto{FiniteWordLengthNatural}[thm]{1}}
  \proofstepwidestar[1,2]{\WordLength{\WordConcat u{\LetterWord a}}=\WordLength u+1}{\FormulaRefAuto{WordAppendLetterLength}[thm]{1,2}}
  \proofstepwidestar[1]{\WordLength u+1\neq0}{\FormulaRefAuto{PeanoAddOneNonzero}[thm]{3}}
  \proofstepwidestar[]{\WordLength{\EmptyWord A}=0}{\FormulaRefAuto{EmptyWordLengthZero}[thm]}
  \proofstepwidestar[7]{\WordConcat u{\LetterWord a}=\EmptyWord A}{\rA}
  \proofstepwidestar[1,2,7]{\WordLength{\EmptyWord A}=\WordLength u+1}{\rIE{7,4}}
  \proofstepwidestar[1,2,7]{\WordLength u+1=0}{\FormulaRefAuto{a=b\dsep a=c\vdash b=c}[thm]{8,6}}
  \proofstepwidestar[1,2,7]{\bot}{\rBI{9,5}}
  \proofstepwidestar[1,2]{\WordConcat u{\LetterWord a}\neq\EmptyWord A}{\rCI{7,10}}
\end{tabproofwide}

\Needspace{10\baselineskip}
\FormulaThmDeltaK[Elementkriterium der Adresskanten]{\begin{aligned}[t]
    &\GoodAddressSet P\vdash\\[-2pt]
    &(p,q)\in\AddressEdges P\leftrightarrow\\[-2pt]
    &{}p\in P\land q\in P\land{}\\[-2pt]
    &\exists a\in\{0,1\}\,(q=\WordConcat p{\LetterWord a}\\[-2pt]
    &\lor p=\WordConcat q{\LetterWord a})
    \end{aligned}}{B27AddressEdgeCriterion}{\DeltaRow{Mengen}{P\dsep p\dsep q\dsep a}}
\begin{tabproofwide}
  \proofstepwidestar[1]{\GoodAddressSet P}{\rA}
  \proofstepwidestar[1]{\begin{aligned}[t]
    &\AddressEdges P=\{(p,q)\in P\times P\mid\\[-2pt]
    &\exists a\in\{0,1\}\,(q=\WordConcat p{\LetterWord a}\\[-2pt]
    &\lor p=\WordConcat q{\LetterWord a})\}
    \end{aligned}}{\FormulaRefAuto{BinaryAddressGraphDef}[def]{1}}
  \proofstepwidestar[]{\begin{aligned}[t]
    &(p,q)\in\{(u,v)\in P\times P\mid\\[-2pt]
    &{} \qquad\exists a\in\{0,1\}\,(v=\WordConcat u{\LetterWord a}\\[-2pt]
    &\lor u=\WordConcat v{\LetterWord a})\}\\[-2pt]
    &\leftrightarrow p\in P\land q\in P\land{} \qquad\\[-2pt]
    &\exists a\in\{0,1\}\,(q=\WordConcat p{\LetterWord a}\\[-2pt]
    &\lor p=\WordConcat q{\LetterWord a})
    \end{aligned}}{\FormulaRefAuto{ProductComprehensionMembership}[thm]}
  \proofstepwidestar[1]{\begin{aligned}[t]
    &\{(u,v)\in P\times P\mid\\[-2pt]
    &{} \quad\exists a\in\{0,1\}\,(v=\WordConcat u{\LetterWord a}\\[-2pt]
    &\lor u=\WordConcat v{\LetterWord a})\}\\[-2pt]
    &=\AddressEdges P
    \end{aligned}}{\FormulaRefAuto{a=b\vdash b=a}[thm]{2}}
  \proofstepwidestar[1]{\begin{aligned}[t]
    &(p,q)\in\AddressEdges P\leftrightarrow\\[-2pt]
    &{}p\in P\land q\in P\land{}\\[-2pt]
    &\exists a\in\{0,1\}\,(q=\WordConcat p{\LetterWord a}\\[-2pt]
    &\lor p=\WordConcat q{\LetterWord a})
    \end{aligned}}{\rIE{4,3}}
\end{tabproofwide}

\Needspace{10\baselineskip}
\FormulaThmDeltaK[Ein einfacher Adresspfad kehrt nicht sofort um]{\begin{aligned}[t]
    &\GoodAddressSet P\dsep\GraphStruct{P,\AddressEdges P}\dsep\\[-2pt]
    &\BXXVIIAddressPath pn\dsep i\in\NatSeg n\dsep i+1\in\NatSeg n\dsep\\[-2pt]
    &{} a,b\in\{0,1\}\dsep\\[-2pt]
    &p(i+1)=\WordConcat{p(i)}{\LetterWord a}\dsep\\[-2pt]
    &p(i+1)=\WordConcat{p((i+1)+1)}{\LetterWord b}\vdash\bot
    \end{aligned}}{B27AddressNoImmediateReturn}{\DeltaRow{Mengen}{P\dsep p\dsep n\dsep i\dsep a\dsep b}}
\begin{tabproofwide}
  \proofstepwidestar[1]{\GoodAddressSet P}{\rA}
  \proofstepwidestar[2]{\GraphStruct{P,\AddressEdges P}}{\rA}
  \proofstepwidestar[3]{\BXXVIIAddressPath pn}{\rA}
  \proofstepwidestar[4]{i\in\NatSeg n}{\rA}
  \proofstepwidestar[5]{i+1\in\NatSeg n}{\rA}
  \proofstepwidestar[6]{a,b\in\{0,1\}}{\rA}
  \proofstepwidestar[7]{p(i+1)=\WordConcat{p(i)}{\LetterWord a}}{\rA}
  \proofstepwidestar[8]{p(i+1)=\WordConcat{p((i+1)+1)}{\LetterWord b}}{\rA}
  \proofstepwidestar[2,3,4]{\begin{aligned}[t]
    &i\in\mathbb N\land i\in\NatSeg{n+1}\land i+1\in\NatSeg{n+1}\\[-2pt]
    &\land{}p(i)\in P\land p(i+1)\in P\\[-2pt]
    &\land(p(i),p(i+1))\in\AddressEdges P
    \end{aligned}}{\FormulaRefAuto{B27PathStepData}[thm]{2,3,4}}
  \proofstepwidestar[2,3,5]{\begin{aligned}[t]
    &i+1\in\mathbb N\land i+1\in\NatSeg{n+1}\\[-2pt]
    &\land(i+1)+1\in\NatSeg{n+1}\\[-2pt]
    &\land{}p(i+1)\in P\land p((i+1)+1)\in P\\[-2pt]
    &\land(p(i+1),p((i+1)+1))\in\AddressEdges P
    \end{aligned}}{\FormulaRefAuto{B27PathStepData}[thm]{2,3,5}}
  \proofstepwidestar[1,2,3,4]{p(i)\in\FiniteWords{\{0,1\}}}{\FormulaRefAuto{B27AddressWordTyping}[thm]{1,\rAEa{\rAEb{\rAEb{\rAEb{9}}}}}}
  \proofstepwidestar[1,2,3,5]{p((i+1)+1)\in\FiniteWords{\{0,1\}}}{\FormulaRefAuto{B27AddressWordTyping}[thm]{1,\rAEa{\rAEb{\rAEb{\rAEb{\rAEb{10}}}}}}}
  \proofstepwidestar[7,8]{\WordConcat{p(i)}{\LetterWord a}=\WordConcat{p((i+1)+1)}{\LetterWord b}}{\FormulaRefAuto{a=b\dsep a=c\vdash b=c}[thm]{7,8}}
  \proofstepwidestar[1,2,3,4,5,6,7,8]{p(i)=p((i+1)+1)\land a=b}{\FormulaRefAuto{WordAppendJointInjectivity}[thm]{\rAI{11,12},6,13}}
  \proofstepwidestar[1,2,3,4,5,6,7,8]{p(i)=p((i+1)+1)}{\rAEa{14}}
  \proofstepwidestar[2,3]{p\colon\NatSeg{n+1}\inj P}{\FormulaRefAuto{DirectedPathInjectiveAxiom}[ax]{2,3}}
  \proofstepwidestar[1,2,3,4,5,6,7,8]{i=(i+1)+1}{\FormulaRefAuto{Injektivität}[ax]{16,\rAEa{\rAEb{9}},\rAEa{\rAEb{\rAEb{10}}},15}}
  \proofstepwidestar[2,3,4]{i<i+1}{\FormulaRefAuto{PeanoLtAddOne}[thm]{\rAEa{9}}}
  \proofstepwidestar[2,3,5]{i+1<(i+1)+1}{\FormulaRefAuto{PeanoLtAddOne}[thm]{\rAEa{10}}}
  \proofstepwidestar[2,3,5]{(i+1)+1\in\mathbb N}{\FormulaRefAuto{PeanoAddOneClosure}[thm]{\rAEa{10}}}
  \proofstepwidestar[2,3,4,5]{i<(i+1)+1}{\FormulaRefAuto{PeanoLtTransitive}[thm]{\rAEa{9},\rAEa{10},20,18,19}}
  \proofstepwidestar[1,2,3,4,5,6,7,8]{\bot}{\FormulaRefAuto{PeanoStrictEqualityImpossible}[thm]{\rAEa{9},20,21,17}}
\end{tabproofwide}

\Needspace{10\baselineskip}
\FormulaThmDeltaK[Adressfortsetzung ist eine Kante]{%
  \GoodAddressSet P\dsep p\in P\dsep a\in\{0,1\}\dsep
  \WordConcat p{\LetterWord a}\in P\vdash
  (p,\WordConcat p{\LetterWord a})\in\AddressEdges P
}{B27AddressAppendEdge}{\DeltaRow{Mengen}{P\dsep p\dsep a}}
\begin{tabproofwide}
  \proofstepwidestar[1]{\GoodAddressSet P}{\rA}
  \proofstepwidestar[2]{p\in P}{\rA}
  \proofstepwidestar[3]{a\in\{0,1\}}{\rA}
  \proofstepwidestar[4]{\WordConcat p{\LetterWord a}\in P}{\rA}
  \proofstepwidestar[]{\WordConcat p{\LetterWord a}=\WordConcat p{\LetterWord a}}{\rII}
  \proofstepwidestar[3]{\begin{aligned}[t]
    &\exists i\in\{0,1\}\, (\WordConcat p{\LetterWord a}=\WordConcat p{\LetterWord i}\\[-2pt]
    &\lor p=\WordConcat{\WordConcat p{\LetterWord a}}{\LetterWord i})
    \end{aligned}}{\rEI{\rAI{3,\rOIa{5}}}}
  \proofstepwidestar[1]{\begin{aligned}[t]
    &(p,\WordConcat p{\LetterWord a})\in\AddressEdges P\leftrightarrow\\[-2pt]
    &{} p\in P\land\WordConcat p{\LetterWord a}\in P\land{}\\[-2pt]
    &\exists i\in\{0,1\}\, (\WordConcat p{\LetterWord a}=\WordConcat p{\LetterWord i}\\[-2pt]
    &\lor p=\WordConcat{\WordConcat p{\LetterWord a}}{\LetterWord i})
    \end{aligned}}{\FormulaRefAuto{B27AddressEdgeCriterion}[thm]{1}}
  \proofstepwidestar[2,3,4]{\begin{aligned}[t]
    &p\in P\land\WordConcat p{\LetterWord a}\in P\land{}\\[-2pt]
    &\exists i\in\{0,1\}\, (\WordConcat p{\LetterWord a}=\WordConcat p{\LetterWord i}\\[-2pt]
    &\lor p=\WordConcat{\WordConcat p{\LetterWord a}}{\LetterWord i})
    \end{aligned}}{\rAIStar{2,4,6}}
  \proofstepwidestar[1,2,3,4]{(p,\WordConcat p{\LetterWord a})\in\AddressEdges P}{\rRE{\rLREb{7},8}}
\end{tabproofwide}

\Needspace{10\baselineskip}
\FormulaThmDeltaK[Der Rumpf bleibt in der Präfixmenge]{%
  \GoodAddressSet P\dsep q\in P\dsep q\neq\BXXVIIAddressRoot\vdash
  q\in\NonemptyWords{\{0,1\}}\land\WordRumpf{\{0,1\}}q\in P
}{B27AddressRumpfClosure}{\DeltaRow{Mengen}{P\dsep q\dsep u\dsep a}}
\begin{tabproofwide}
  \proofstepwidestar[1]{\GoodAddressSet P}{\rA}
  \proofstepwidestar[2]{q\in P}{\rA}
  \proofstepwidestar[3]{q\neq\BXXVIIAddressRoot}{\rA}
  \proofstepwidestar[3]{q\notin\{\BXXVIIAddressRoot\}}{\FormulaRefAuto{x\notin\{a\}\eqvdash x\neq a}[thm]{3}}
  \proofstepwidestar[2,3]{q\in P\setminus\{\BXXVIIAddressRoot\}}{\FormulaRefAuto{x\in A\dsep x\notin B\vdash x\in A\setminus B}[thm]{2,4}}
  \proofstepwidestar[1,2,3]{\exists u\in P\,\exists a\in\{0,1\}\,q=\WordConcat u{\LetterWord a}}{\FormulaRefAuto{GoodBinaryAddressSetPrefixAxiom}[ax]{1,5}}
  \proofstepwidestar[7]{u\in P\land a\in\{0,1\}\land q=\WordConcat u{\LetterWord a}}{\rA}
  \proofstepwidestar[7]{u\in P}{\rAEa{7}}
  \proofstepwidestar[7]{a\in\{0,1\}}{\rAEa{\rAEb{7}}}
  \proofstepwidestar[7]{q=\WordConcat u{\LetterWord a}}{\rAEb{\rAEb{7}}}
  \proofstepwidestar[1,7]{u\in\FiniteWords{\{0,1\}}}{\FormulaRefAuto{B27AddressWordTyping}[thm]{1,8}}
  \proofstepwidestar[1,7]{\WordConcat u{\LetterWord a}\in\NonemptyWords{\{0,1\}}}{\FormulaRefAuto{WordAppendNonempty}[thm]{11,9}}
  \proofstepwidestar[7]{\WordConcat u{\LetterWord a}=q}{\FormulaRefAuto{a=b\vdash b=a}[thm]{10}}
  \proofstepwidestar[1,7]{q\in\NonemptyWords{\{0,1\}}}{\rIE{13,12}}
  \proofstepwidestar[1,7]{\WordRumpf{\{0,1\}}{\WordConcat u{\LetterWord a}}=u}{\FormulaRefAuto{WordRumpfAppend}[thm]{11,9}}
  \proofstepwidestar[1,7]{\WordRumpf{\{0,1\}}q=u}{\rIE{13,15}}
  \proofstepwidestar[1,7]{u=\WordRumpf{\{0,1\}}q}{\FormulaRefAuto{a=b\vdash b=a}[thm]{16}}
  \proofstepwidestar[1,7]{\WordRumpf{\{0,1\}}q\in P}{\rIE{17,8}}
  \proofstepwidestar[1,7]{q\in\NonemptyWords{\{0,1\}}\land\WordRumpf{\{0,1\}}q\in P}{\rAI{14,18}}
  \proofstepwidestar[1,2,3]{q\in\NonemptyWords{\{0,1\}}\land\WordRumpf{\{0,1\}}q\in P}{\rEEStar{6,7,19}}
\end{tabproofwide}

\Needspace{10\baselineskip}
\FormulaThmDeltaK[Die beiden Richtungen einer Adresskante]{\begin{aligned}[t]
    &\GoodAddressSet P\dsep(p,q)\in\AddressEdges P\vdash\\[-2pt]
    &\exists a\in\{0,1\}\,(q=\WordConcat p{\LetterWord a}\\[-2pt]
    &\lor p=\WordConcat q{\LetterWord a})
    \end{aligned}}{B27AddressEdgeDirections}{\DeltaRow{Mengen}{P\dsep p\dsep q\dsep a}}
\begin{tabproofwide}
  \proofstepwidestar[1]{\GoodAddressSet P}{\rA}
  \proofstepwidestar[2]{(p,q)\in\AddressEdges P}{\rA}
  \proofstepwidestar[1]{\begin{aligned}[t]
    &(p,q)\in\AddressEdges P\leftrightarrow\\[-2pt]
    &{}p\in P\land q\in P\land{}\\[-2pt]
    &\exists a\in\{0,1\}\,(q=\WordConcat p{\LetterWord a}\\[-2pt]
    &\lor p=\WordConcat q{\LetterWord a})
    \end{aligned}}{\FormulaRefAuto{B27AddressEdgeCriterion}[thm]{1}}
  \proofstepwidestar[1,2]{\begin{aligned}[t]
    &p\in P\land q\in P\land{}\\[-2pt]
    &\exists a\in\{0,1\}\,(q=\WordConcat p{\LetterWord a}\\[-2pt]
    &\lor p=\WordConcat q{\LetterWord a})
    \end{aligned}}{\rRE{\rLREa{3},2}}
  \proofstepwidestar[1,2]{\begin{aligned}[t]
    &\exists a\in\{0,1\}\,(q=\WordConcat p{\LetterWord a}\\[-2pt]
    &\lor p=\WordConcat q{\LetterWord a})
    \end{aligned}}{\rAEb{\rAEb{4}}}
\end{tabproofwide}

\Needspace{10\baselineskip}
\FormulaThmDeltaK[Eine Adresskante aus der Wurzel führt vorwärts]{%
  \GoodAddressSet P\dsep q\in P\dsep
  (\BXXVIIAddressRoot,q)\in\AddressEdges P\vdash
  \exists a\in\{0,1\}\,q=\WordConcat{\BXXVIIAddressRoot}{\LetterWord a}
}{B27AddressRootEdgeForward}{\DeltaRow{Mengen}{P\dsep q\dsep a}}
\begin{tabproofwide}
  \proofstepwidestar[1]{\GoodAddressSet P}{\rA}
  \proofstepwidestar[2]{q\in P}{\rA}
  \proofstepwidestar[3]{(\BXXVIIAddressRoot,q)\in\AddressEdges P}{\rA}
  \proofstepwidestar[1,3]{\begin{aligned}[t]
    &\exists a\in\{0,1\}\, (q=\WordConcat{\BXXVIIAddressRoot}{\LetterWord a}\\[-2pt]
    &\lor \BXXVIIAddressRoot=\WordConcat q{\LetterWord a})
    \end{aligned}}{\FormulaRefAuto{B27AddressEdgeDirections}[thm]{1,3}}
  \proofstepwidestar[5]{\begin{aligned}[t]
    &a\in\{0,1\}\land (q=\WordConcat{\BXXVIIAddressRoot}{\LetterWord a}\\[-2pt]
    &\lor \BXXVIIAddressRoot=\WordConcat q{\LetterWord a})
    \end{aligned}}{\rA}
  \proofstepwidestar[5]{a\in\{0,1\}}{\rAEa{5}}
  \proofstepwidestar[1,2]{q\in\FiniteWords{\{0,1\}}}{\FormulaRefAuto{B27AddressWordTyping}[thm]{1,2}}
  \proofstepwidestar[1,2,5]{\WordConcat q{\LetterWord a}\neq\BXXVIIAddressRoot}{\FormulaRefAuto{B27WordAppendNotEmpty}[thm]{7,6}}
  \proofstepwidestar[9]{\BXXVIIAddressRoot=\WordConcat q{\LetterWord a}}{\rA}
  \proofstepwidestar[9]{\WordConcat q{\LetterWord a}=\BXXVIIAddressRoot}{\FormulaRefAuto{a=b\vdash b=a}[thm]{9}}
  \proofstepwidestar[1,2,5,9]{\bot}{\rBI{10,8}}
  \proofstepwidestar[1,2,5]{\BXXVIIAddressRoot\neq\WordConcat q{\LetterWord a}}{\rCI{9,11}}
  \proofstepwidestar[1,2,5]{q=\WordConcat{\BXXVIIAddressRoot}{\LetterWord a}}{\FormulaRefAuto{P\lor Q,\neg Q\vdash P}[thm]{\rAEb{5},12}}
  \proofstepwidestar[1,2,5]{\exists a\in\{0,1\}\,q=\WordConcat{\BXXVIIAddressRoot}{\LetterWord a}}{\rEI{\rAI{6,13}}}
  \proofstepwidestar[1,2,3]{\exists a\in\{0,1\}\,q=\WordConcat{\BXXVIIAddressRoot}{\LetterWord a}}{\rEE{4,5,14}}
\end{tabproofwide}

\subsection{Wurzelwege wachsen mit ihren Adressen}

Wortinduktion liefert einen Weg zu jeder Adresse. Für die Eindeutigkeit
wird anschließend vom Endknoten rückwärts argumentiert: Gleiche Adressen
haben gleiche Rümpfe und damit dieselben vorherigen Wegknoten.

\Needspace{10\baselineskip}
\FormulaThmDeltaK[Jede Adresse besitzt einen Wurzelpfad]{%
  \GoodAddressSet P\dsep\GraphStruct{P,\AddressEdges P}\dsep q\in P\vdash
  \exists z\,(z\in\GraphPathsBetween P{\AddressEdges P}{\BXXVIIAddressRoot}q)
}{B27AddressRootPathExists}{\DeltaRow{Mengen}{P\dsep q\dsep u\dsep a\dsep z}}
Hier steht \(D(q)\) für
\(\exists z\,(z\in\GraphPathsBetween P{\AddressEdges P}{\BXXVIIAddressRoot}q)\)
und \(J(q)\) für \(q\in P\rightarrow D(q)\).
\begin{tabproofsplitwide}
  \proofpartwideindR[Buchstabenadressen]{%
  \GoodAddressSet P\dsep\GraphStruct{P,\AddressEdges P}\vdash
  \forall a\in\{0,1\}\,J(\LetterWord a)
  }{\GoodAddressSet P\dsep\GraphStruct{P,\AddressEdges P}\vdash
  \forall a\in\{0,1\}\,J(\LetterWord a)}[B27AddressRootPathExistsBase]
  \proofstepwidestar[1]{\GoodAddressSet P}{\rA}
  \proofstepwidestar[2]{\GraphStruct{P,\AddressEdges P}}{\rA}
  \proofstepwidestar[3]{a\in\{0,1\}}{\rA}
  \proofstepwidestar[4]{\LetterWord a\in P}{\rA}
  \proofstepwidestar[1]{\BXXVIIAddressRoot\in P}{\FormulaRefAuto{GoodBinaryAddressSetRootAxiom}[ax]{1}}
  \proofstepwidestar[1,2]{D(\BXXVIIAddressRoot)}{\FormulaRefAuto{B27RootPathSingleton}[thm]{2,5}}
  \proofstepwidestar[3]{\LetterWord a\in\NonemptyWords{\{0,1\}}}{\FormulaRefAuto{LetterWordNonempty}[thm]{3}}
  \proofstepwidestar[3]{\LetterWord a\in\FiniteWords{\{0,1\}}}{\FormulaRefAuto{NonemptyWordFinite}[thm]{7}}
  \proofstepwidestar[3]{\WordConcat{\BXXVIIAddressRoot}{\LetterWord a}=\LetterWord a}{\FormulaRefAuto{WordConcatenationLeftIdentity}[thm]{8}}
  \proofstepwidestar[3]{\LetterWord a=\WordConcat{\BXXVIIAddressRoot}{\LetterWord a}}{\FormulaRefAuto{a=b\vdash b=a}[thm]{9}}
  \proofstepwidestar[3,4]{\WordConcat{\BXXVIIAddressRoot}{\LetterWord a}\in P}{\rIE{10,4}}
  \proofstepwidestar[1,3,4]{(\BXXVIIAddressRoot,\WordConcat{\BXXVIIAddressRoot}{\LetterWord a})\in\AddressEdges P}{\FormulaRefAuto{B27AddressAppendEdge}[thm]{1,5,3,11}}
  \proofstepwidestar[1,3,4]{(\BXXVIIAddressRoot,\LetterWord a)\in\AddressEdges P}{\rIE{9,12}}
  \proofstepwidestar[1,2,3,4]{D(\LetterWord a)}{\FormulaRefAuto{B27RootPathExtendExistence}[thm]{2,6,13}}
  \proofstepwidestar[1,2,3]{J(\LetterWord a)}{\rRI{4,14}}
  \proofstepwidestar[1,2]{\forall a\in\{0,1\}\,J(\LetterWord a)}{\rUI{\rRI{3,15}}}
  \closeproofpartwideindR
  \Needspace{15\baselineskip}
  \proofpartwideindR[Eine weitere Stelle]{%
  \begin{aligned}[t]&\GoodAddressSet P\dsep\GraphStruct{P,\AddressEdges P}\vdash{}\\[-2pt]
  &\forall u\in\NonemptyWords{\{0,1\}}\,\forall a\in\{0,1\}\,
  (J(u)\rightarrow J(\WordConcat u{\LetterWord a}))\end{aligned}
  }{\GoodAddressSet P\dsep\GraphStruct{P,\AddressEdges P}\vdash
  \forall u\in\NonemptyWords{\{0,1\}}\,\forall a\in\{0,1\}\,
  (J(u)\rightarrow J(\WordConcat u{\LetterWord a}))}[B27AddressRootPathExistsStep]
  \proofstepwidestar[1]{\GoodAddressSet P}{\rA}
  \proofstepwidestar[2]{\GraphStruct{P,\AddressEdges P}}{\rA}
  \proofstepwidestar[3]{u\in\NonemptyWords{\{0,1\}}}{\rA}
  \proofstepwidestar[4]{a\in\{0,1\}}{\rA}
  \proofstepwidestar[5]{J(u)}{\rA}
  \proofstepwidestar[6]{\WordConcat u{\LetterWord a}\in P}{\rA}
  \proofstepwidestar[3]{u\in\FiniteWords{\{0,1\}}}{\FormulaRefAuto{NonemptyWordFinite}[thm]{3}}
  \proofstepwidestar[3,4]{\WordConcat u{\LetterWord a}\neq\BXXVIIAddressRoot}{\FormulaRefAuto{B27WordAppendNotEmpty}[thm]{7,4}}
  \proofstepwidestar[1,3,4,6]{\begin{aligned}[t]&\WordConcat u{\LetterWord a}\in\NonemptyWords{\{0,1\}}\land{}\\[-2pt]&\WordRumpf{\{0,1\}}{\WordConcat u{\LetterWord a}}\in P\end{aligned}}{\FormulaRefAuto{B27AddressRumpfClosure}[thm]{1,6,8}}
  \proofstepwidestar[3,4]{\WordRumpf{\{0,1\}}{\WordConcat u{\LetterWord a}}=u}{\FormulaRefAuto{WordRumpfAppend}[thm]{7,4}}
  \proofstepwidestar[1,3,4,6]{u\in P}{\rIE{10,\rAEb{9}}}
  \proofstepwidestar[1,3,4,5,6]{D(u)}{\rRE{5,11}}
  \proofstepwidestar[1,3,4,6]{(u,\WordConcat u{\LetterWord a})\in\AddressEdges P}{\FormulaRefAuto{B27AddressAppendEdge}[thm]{1,11,4,6}}
  \proofstepwidestar[1,2,3,4,5,6]{D(\WordConcat u{\LetterWord a})}{\FormulaRefAuto{B27RootPathExtendExistence}[thm]{2,12,13}}
  \proofstepwidestar[1,2,3,4,5]{J(\WordConcat u{\LetterWord a})}{\rRI{6,14}}
  \proofstepwidestar[1,2,3,4]{J(u)\rightarrow J(\WordConcat u{\LetterWord a})}{\rRI{5,15}}
  \proofstepwidestar[1,2]{\begin{aligned}[t]&\forall u\in\NonemptyWords{\{0,1\}}\,\forall a\in\{0,1\}\,\\[-2pt]&(J(u)\rightarrow J(\WordConcat u{\LetterWord a}))\end{aligned}}{\rUI{\rRI{3,\rUI{\rRI{4,16}}}}}
  \closeproofpartwideindR
  \Needspace{28\baselineskip}
  \proofpartwide{Leerwort und Wortinduktion}
  \proofstepwidestar[1]{\GoodAddressSet P}{\rA}
  \proofstepwidestar[2]{\GraphStruct{P,\AddressEdges P}}{\rA}
  \proofstepwidestar[3]{q\in P}{\rA}
  \proofstepwidestar[1,2]{\forall a\in\{0,1\}\,J(\LetterWord a)}{\FormulaRefAuto{B27AddressRootPathExistsBase}[thm]{1,2}}
  \proofstepwidestar[1,2]{\begin{aligned}[t]&\forall u\in\NonemptyWords{\{0,1\}}\,\forall a\in\{0,1\}\,\\[-2pt]&(J(u)\rightarrow J(\WordConcat u{\LetterWord a}))\end{aligned}}{\FormulaRefAuto{B27AddressRootPathExistsStep}[thm]{1,2}}
  \proofstepwidestar[1,2]{\forall q\in\NonemptyWords{\{0,1\}}\,J(q)}{\FormulaRefAuto{NonemptyWordInduction}[thm]{4,5}}
  \proofstepwidestar[]{q=\BXXVIIAddressRoot\lor q\neq\BXXVIIAddressRoot}{\FormulaRefAuto{P\lor\neg P}[thm]}
  \proofstepwidestar[8]{q=\BXXVIIAddressRoot}{\rA}
  \proofstepwidestar[1]{\BXXVIIAddressRoot\in P}{\FormulaRefAuto{GoodBinaryAddressSetRootAxiom}[ax]{1}}
  \proofstepwidestar[1,2]{D(\BXXVIIAddressRoot)}{\FormulaRefAuto{B27RootPathSingleton}[thm]{2,9}}
  \proofstepwidestar[8]{\BXXVIIAddressRoot=q}{\FormulaRefAuto{a=b\vdash b=a}[thm]{8}}
  \proofstepwidestar[1,2,8]{D(q)}{\rIE{11,10}}
  \proofstepwidestar[13]{q\neq\BXXVIIAddressRoot}{\rA}
  \proofstepwidestar[1,3,13]{q\in\NonemptyWords{\{0,1\}}\land\WordRumpf{\{0,1\}}q\in P}{\FormulaRefAuto{B27AddressRumpfClosure}[thm]{1,3,13}}
  \proofstepwidestar[1,2,3,13]{J(q)}{\rRE{\rUE{6},\rAEa{14}}}
  \proofstepwidestar[1,2,3,13]{D(q)}{\rRE{15,3}}
  \proofstepwidestar[1,2,3]{D(q)}{\rOE{7,8,12,13,16}}
  \closeproofpartwide
\end{tabproofsplitwide}

\Needspace{10\baselineskip}
\FormulaThmDeltaK[Jeder Wurzelpfadschritt fügt ein Bit an]{%
  \begin{aligned}[t]&\GoodAddressSet P\dsep\GraphStruct{P,\AddressEdges P}\dsep
  \BXXVIIAddressPath pn\dsep p(0)=\BXXVIIAddressRoot\vdash{}\\[-2pt]
  &\BXXVIIAddressForward pn\end{aligned}
}{B27AddressRootPathForward}{\DeltaRow{Mengen}{P\dsep p\dsep n\dsep i\dsep a\dsep b}}
Für diesen Beweis sei
\[
  Q(i)\coloneqq i\in\NatSeg n\rightarrow
  \exists a\in\{0,1\}\,p(i+1)=\WordConcat{p(i)}{\LetterWord a}.
\]
\Needspace{12\baselineskip}
\begin{tabproofsplitwide}
  \proofpartwideindR[Erster Schritt]{%
    \GoodAddressSet P\dsep\GraphStruct{P,\AddressEdges P}\dsep
    \BXXVIIAddressPath pn\dsep p(0)=\BXXVIIAddressRoot\vdash Q(0)
  }{\GoodAddressSet P\dsep\GraphStruct{P,\AddressEdges P}\dsep
    \BXXVIIAddressPath pn\dsep p(0)=\BXXVIIAddressRoot\vdash Q(0)}[B27AddressRootPathForwardBase]
  \proofstepwidestar[1]{\GoodAddressSet P}{\rA}
  \proofstepwidestar[2]{\GraphStruct{P,\AddressEdges P}}{\rA}
  \proofstepwidestar[3]{\BXXVIIAddressPath pn}{\rA}
  \proofstepwidestar[4]{p(0)=\BXXVIIAddressRoot}{\rA}
  \proofstepwidestar[5]{0\in\NatSeg n}{\rA}
  \proofstepwidestar[2,3]{\DirectedWalk{P}{\AddressEdges P}{p}{n}}{\FormulaRefAuto{DirectedPathWalkAxiom}[ax]{2,3}}
  \proofstepwidestar[2,3,5]{(p(0),p(0+1))\in\AddressEdges P}{\FormulaRefAuto{DirectedWalkStepEdge}[thm]{2,6,5}}
  \proofstepwidestar[2,3,5]{p(0)\in P\land p(0+1)\in P}{\FormulaRefAuto{GraphEdgeEndpointTyping}[thm]{2,7}}
  \proofstepwidestar[2,3,5]{p(0+1)\in P}{\rAEb{8}}
  \proofstepwidestar[2,3,4,5]{(\BXXVIIAddressRoot,p(0+1))\in\AddressEdges P}{\rIE{4,7}}
  \proofstepwidestar[1,2,3,4,5]{\exists a\in\{0,1\}\,p(0+1)=\WordConcat{\BXXVIIAddressRoot}{\LetterWord a}}{\FormulaRefAuto{B27AddressRootEdgeForward}[thm]{1,9,10}}
  \proofstepwidestar[4]{\BXXVIIAddressRoot=p(0)}{\FormulaRefAuto{a=b\vdash b=a}[thm]{4}}
  \proofstepwidestar[1,2,3,4,5]{\exists a\in\{0,1\}\,p(0+1)=\WordConcat{p(0)}{\LetterWord a}}{\rIE{12,11}}
  \proofstepwidestar[1,2,3,4]{Q(0)}{\rRI{5,13}}
  \closeproofpartwideindR

  \Needspace{15\baselineskip}
  \proofpartwideindR[Ein Rückschritt würde einen Knoten wiederholen]{%
    \begin{aligned}[t]&\GoodAddressSet P\dsep\GraphStruct{P,\AddressEdges P}\dsep
    \BXXVIIAddressPath pn\dsep{}\\[-2pt]&i\in\mathbb N\dsep Q(i)\vdash Q(i+1)\end{aligned}
  }{\GoodAddressSet P\dsep\GraphStruct{P,\AddressEdges P}\dsep
    \BXXVIIAddressPath pn\dsep i\in\mathbb N\dsep Q(i)\vdash Q(i+1)}[B27AddressRootPathForwardStep]
  \proofstepwidestar[1]{\GoodAddressSet P}{\rA}
  \proofstepwidestar[2]{\GraphStruct{P,\AddressEdges P}}{\rA}
  \proofstepwidestar[3]{\BXXVIIAddressPath pn}{\rA}
  \proofstepwidestar[4]{i\in\mathbb N}{\rA}
  \proofstepwidestar[5]{Q(i)}{\rA}
  \proofstepwidestar[6]{i+1\in\NatSeg n}{\rA}
  \proofstepwidestar[2,3]{n\in\mathbb N\land p\colon\NatSeg{n+1}\to P\land p\colon\NatSeg{n+1}\inj P}{\FormulaRefAuto{FinitePathTyping}[thm]{2,3}}
  \proofstepwidestar[2,3]{n\in\mathbb N}{\rAEa{7}}
  \proofstepwidestar[4]{i+1\in\mathbb N}{\FormulaRefAuto{PeanoAddOneClosure}[thm]{4}}
  \proofstepwidestar[4]{i<i+1}{\FormulaRefAuto{PeanoLtAddOne}[thm]{4}}
  \proofstepwidestar[2,3,4,6]{i\in\NatSeg n}{\FormulaRefAuto{B27SmallerIndexInSegment}[thm]{\rAIStar{8,4,9},10,6}}
  \proofstepwidestar[2,3,4,5,6]{\exists a\in\{0,1\}\,p(i+1)=\WordConcat{p(i)}{\LetterWord a}}{\rRE{11,5}}
  \proofstepwidestar[13]{a\in\{0,1\}\land p(i+1)=\WordConcat{p(i)}{\LetterWord a}}{\rA}
  \proofstepwidestar[2,3]{\DirectedWalk{P}{\AddressEdges P}{p}{n}}{\FormulaRefAuto{DirectedPathWalkAxiom}[ax]{2,3}}
  \proofstepwidestar[2,3,6]{(p(i+1),p((i+1)+1))\in\AddressEdges P}{\FormulaRefAuto{DirectedWalkStepEdge}[thm]{2,14,6}}
  \proofstepwidestar[1,2,3,6]{\begin{aligned}[t]&\exists b\in\{0,1\}\,(\\[-2pt]&\quad p((i+1)+1)=\WordConcat{p(i+1)}{\LetterWord b}\lor{}\\[-2pt]&\quad p(i+1)=\WordConcat{p((i+1)+1)}{\LetterWord b})\end{aligned}}{\FormulaRefAuto{B27AddressEdgeDirections}[thm]{1,15}}
  \proofstepwidestar[17]{\begin{aligned}[t]&b\in\{0,1\}\land(\\[-2pt]&\quad p((i+1)+1)=\WordConcat{p(i+1)}{\LetterWord b}\lor{}\\[-2pt]&\quad p(i+1)=\WordConcat{p((i+1)+1)}{\LetterWord b})\end{aligned}}{\rA}
  \proofstepwidestar[18]{p(i+1)=\WordConcat{p((i+1)+1)}{\LetterWord b}}{\rA}
  \proofstepwidestar[1,2,3,4,6,13,17,18]{\bot}{\FormulaRefAuto{B27AddressNoImmediateReturn}[thm]{1,2,3,11,6,\rAI{\rAEa{13},\rAEa{17}},\rAEb{13},18}}
  \proofstepwidestar[1,2,3,4,6,13,17]{p(i+1)\neq\WordConcat{p((i+1)+1)}{\LetterWord b}}{\rCI{18,19}}
  \proofstepwidestar[1,2,3,4,6,13,17]{p((i+1)+1)=\WordConcat{p(i+1)}{\LetterWord b}}{\FormulaRefAuto{P\lor Q,\neg Q\vdash P}[thm]{\rAEb{17},20}}
  \proofstepwidestar[1,2,3,4,6,13,17]{\begin{aligned}[t]&\exists b\in\{0,1\}\,\\[-2pt]&p((i+1)+1)=\WordConcat{p(i+1)}{\LetterWord b}\end{aligned}}{\rEI{\rAI{\rAEa{17},21}}}
  \proofstepwidestar[1,2,3,4,6,13]{\begin{aligned}[t]&\exists b\in\{0,1\}\,\\[-2pt]&p((i+1)+1)=\WordConcat{p(i+1)}{\LetterWord b}\end{aligned}}{\rEE{16,17,22}}
  \proofstepwidestar[1,2,3,4,5,6]{\begin{aligned}[t]&\exists b\in\{0,1\}\,\\[-2pt]&p((i+1)+1)=\WordConcat{p(i+1)}{\LetterWord b}\end{aligned}}{\rEE{12,13,23}}
  \proofstepwidestar[1,2,3,4,5]{Q(i+1)}{\rRI{6,24}}
  \closeproofpartwideindR

  \Needspace{28\baselineskip}
  \proofpartwide{Induktionsschluss}
  \proofstepwidestar[1]{\GoodAddressSet P}{\rA}
  \proofstepwidestar[2]{\GraphStruct{P,\AddressEdges P}}{\rA}
  \proofstepwidestar[3]{\BXXVIIAddressPath pn}{\rA}
  \proofstepwidestar[4]{p(0)=\BXXVIIAddressRoot}{\rA}
  \proofstepwidestar[1,2,3,4]{Q(0)}{\FormulaRefAuto{B27AddressRootPathForwardBase}[thm]{1,2,3,4}}
  \proofstepwidestar[6]{i\in\mathbb N}{\rA}
  \proofstepwidestar[7]{Q(i)}{\rA}
  \proofstepwidestar[1,2,3,6,7]{Q(i+1)}{\FormulaRefAuto{B27AddressRootPathForwardStep}[thm]{1,2,3,6,7}}
  \proofstepwidestar[1,2,3,6]{Q(i)\rightarrow Q(i+1)}{\rRI{7,8}}
  \proofstepwidestar[1,2,3]{i\in\mathbb N\rightarrow(Q(i)\rightarrow Q(i+1))}{\rRI{6,9}}
  \proofstepwidestar[1,2,3]{\forall i\in\mathbb N\,(Q(i)\rightarrow Q(i+1))}{\rUI{10}}
  \proofstepwidestar[1,2,3,4]{\forall i\in\mathbb N\,Q(i)}{\FormulaRefAuto{PeanoInductionAddOne}[thm]{5,11}}
  \proofstepwidestar[13]{i\in\NatSeg n}{\rA}
  \proofstepwidestar[2,3]{n\in\mathbb N\land p\colon\NatSeg{n+1}\to P\land p\colon\NatSeg{n+1}\inj P}{\FormulaRefAuto{FinitePathTyping}[thm]{2,3}}
  \proofstepwidestar[2,3]{n\in\mathbb N}{\rAEa{14}}
  \proofstepwidestar[2,3,13]{i\in\mathbb N}{\FormulaRefAuto{PeanoNatSegElementNatural}[thm]{15,13}}
  \proofstepwidestar[1,2,3,4]{i\in\mathbb N\rightarrow Q(i)}{\rUE{12}}
  \proofstepwidestar[1,2,3,4,13]{Q(i)}{\rRE{16,17}}
  \proofstepwidestar[1,2,3,4,13]{\exists a\in\{0,1\}\,p(i+1)=\WordConcat{p(i)}{\LetterWord a}}{\rRE{13,18}}
  \proofstepwidestar[1,2,3,4]{\BXXVIIAddressForward pn}{\rUI{\rRI{13,19}}}
  \closeproofpartwide
\end{tabproofsplitwide}

\Needspace{32\baselineskip}
\FormulaThmDeltaK[Die Adresslänge zählt die Pfadschritte]{%
  \begin{aligned}[t]&\GoodAddressSet P\dsep\GraphStruct{P,\AddressEdges P}\dsep
  \BXXVIIAddressPath pn\dsep p(0)=\BXXVIIAddressRoot\vdash{}\\[-2pt]
  &\forall i\in\NatSeg{n+1}\,\WordLength{p(i)}=i\end{aligned}
}{B27AddressRootPathRanks}{\DeltaRow{Mengen}{P\dsep p\dsep n\dsep i\dsep a}}
Hier bezeichnet \(R(i)\) die Formel
\(i\in\NatSeg{n+1}\rightarrow\WordLength{p(i)}=i\).
\begin{tabproofsplitwide}
  \proofpartwideindR[Länge an der Wurzel]{%
    p(0)=\BXXVIIAddressRoot\vdash R(0)
  }{p(0)=\BXXVIIAddressRoot\vdash R(0)}[B27AddressRootPathRanksBase]
\proofstepwidestar[1]{p(0)=\BXXVIIAddressRoot}{\rA}
  \proofstepwidestar[2]{0\in\NatSeg{n+1}}{\rA}
  \proofstepwidestar[]{\WordLength{\BXXVIIAddressRoot}=0}{\FormulaRefAuto{EmptyWordLengthZero}[thm]}
  \proofstepwidestar[1]{\BXXVIIAddressRoot=p(0)}{\FormulaRefAuto{a=b\vdash b=a}[thm]{1}}
  \proofstepwidestar[1]{\WordLength{p(0)}=0}{\rIE{4,3}}
  \proofstepwidestar[1,2]{0\in\NatSeg{n+1}\land\WordLength{p(0)}=0}{\rAI{2,5}}
  \proofstepwidestar[1,2]{\WordLength{p(0)}=0}{\rAEb{6}}
  \proofstepwidestar[1]{R(0)}{\rRI{2,7}}
  \closeproofpartwideindR
  \Needspace{16\baselineskip}
  \proofpartwideindR[Längenzuwachs um eins]{%
    \begin{aligned}[t]&\GoodAddressSet P\dsep\GraphStruct{P,\AddressEdges P}\dsep
    \BXXVIIAddressPath pn\dsep{}\\[-2pt]
    &\BXXVIIAddressForward pn\dsep{}\\[-2pt]
    &i\in\mathbb N\dsep R(i)\vdash R(i+1)\end{aligned}
  }{\GoodAddressSet P\dsep\GraphStruct{P,\AddressEdges P}\dsep
    \BXXVIIAddressPath pn\dsep\BXXVIIAddressForward pn\dsep
    i\in\mathbb N\dsep R(i)\vdash R(i+1)}[B27AddressRootPathRanksStep]
\proofstepwidestar[1]{\GoodAddressSet P}{\rA}
  \proofstepwidestar[2]{\GraphStruct{P,\AddressEdges P}}{\rA}
  \proofstepwidestar[3]{\BXXVIIAddressPath pn}{\rA}
  \proofstepwidestar[4]{\BXXVIIAddressForward pn}{\rA}
  \proofstepwidestar[5]{i\in\mathbb N}{\rA}
  \proofstepwidestar[6]{R(i)}{\rA}
  \proofstepwidestar[7]{i+1\in\NatSeg{n+1}}{\rA}
  \proofstepwidestar[2,3]{n\in\mathbb N\land p\colon\NatSeg{n+1}\to P\land p\colon\NatSeg{n+1}\inj P}{\FormulaRefAuto{FinitePathTyping}[thm]{2,3}}
  \proofstepwidestar[2,3]{n\in\mathbb N}{\rAEa{8}}
  \proofstepwidestar[2,3,5]{n,i\in\mathbb N}{\rAI{9,5}}
  \proofstepwidestar[2,3,5,7]{i\in\NatSeg n}{\FormulaRefAuto{B27SuccessorIndexBackward}[thm]{10,7}}
  \proofstepwidestar[2,3,5,7]{\begin{aligned}[t]&i\in\mathbb N\land i\in\NatSeg{n+1}\land i+1\in\NatSeg{n+1}\land{}\\[-2pt]&p(i)\in P\land p(i+1)\in P\land{}\\[-2pt]&(p(i),p(i+1))\in\AddressEdges P\end{aligned}}{\FormulaRefAuto{B27PathStepData}[thm]{2,3,11}}
  \proofstepwidestar[2,3,5,7]{i\in\NatSeg{n+1}}{\rAEa{\rAEb{12}}}
  \proofstepwidestar[2,3,5,6,7]{\WordLength{p(i)}=i}{\rRE{6,13}}
  \proofstepwidestar[2,3,4,5,7]{\exists a\in\{0,1\}\,p(i+1)=\WordConcat{p(i)}{\LetterWord a}}{\rRE{\rUE{4},11}}
  \proofstepwidestar[16]{a\in\{0,1\}\land p(i+1)=\WordConcat{p(i)}{\LetterWord a}}{\rA}
  \proofstepwidestar[2,3,5,7]{p(i)\in P}{\rAEa{\rAEb{\rAEb{\rAEb{12}}}}}
  \proofstepwidestar[1,2,3,5,7]{p(i)\in\FiniteWords{\{0,1\}}}{\FormulaRefAuto{B27AddressWordTyping}[thm]{1,17}}
  \proofstepwidestar[1,2,3,5,7,16]{\WordLength{\WordConcat{p(i)}{\LetterWord a}}=\WordLength{p(i)}+1}{\FormulaRefAuto{WordAppendLetterLength}[thm]{18,\rAEa{16}}}
  \proofstepwidestar[16]{\WordConcat{p(i)}{\LetterWord a}=p(i+1)}{\FormulaRefAuto{a=b\vdash b=a}[thm]{\rAEb{16}}}
  \proofstepwidestar[1,2,3,5,7,16]{\WordLength{p(i+1)}=\WordLength{p(i)}+1}{\rIE{20,19}}
  \proofstepwidestar[1,2,3,5,6,7,16]{\WordLength{p(i+1)}=i+1}{\rIE{14,21}}
  \proofstepwidestar[1,2,3,4,5,6,7]{\WordLength{p(i+1)}=i+1}{\rEE{15,16,22}}
  \proofstepwidestar[1,2,3,4,5,6]{R(i+1)}{\rRI{7,23}}
  \closeproofpartwideindR
  \Needspace{10\baselineskip}
  \proofpartwide{Induktionsschluss}
  \proofstepwidestar[1]{\GoodAddressSet P}{\rA}
  \proofstepwidestar[2]{\GraphStruct{P,\AddressEdges P}}{\rA}
  \proofstepwidestar[3]{\BXXVIIAddressPath pn}{\rA}
  \proofstepwidestar[4]{p(0)=\BXXVIIAddressRoot}{\rA}
  \proofstepwidestar[1,2,3,4]{\BXXVIIAddressForward pn}{\FormulaRefAuto{B27AddressRootPathForward}[thm]{1,2,3,4}}
  \proofstepwidestar[4]{R(0)}{\FormulaRefAuto{B27AddressRootPathRanksBase}[thm]{4}}
  \proofstepwidestar[7]{j\in\mathbb N}{\rA}
  \proofstepwidestar[8]{R(j)}{\rA}
  \proofstepwidestar[1,2,3,4,7,8]{R(j+1)}{\FormulaRefAuto{B27AddressRootPathRanksStep}[thm]{1,2,3,5,7,8}}
  \proofstepwidestar[1,2,3,4,7]{R(j)\rightarrow R(j+1)}{\rRI{8,9}}
  \proofstepwidestar[1,2,3,4]{\forall j\in\mathbb N\,(R(j)\rightarrow R(j+1))}{\rUI{\rRI{7,10}}}
  \proofstepwidestar[1,2,3,4]{\forall j\in\mathbb N\,R(j)}{\FormulaRefAuto{PeanoInductionAddOne}[thm]{6,11}}
  \proofstepwidestar[13]{i\in\NatSeg{n+1}}{\rA}
  \proofstepwidestar[2,3]{n\in\mathbb N\land p\colon\NatSeg{n+1}\to P\land p\colon\NatSeg{n+1}\inj P}{\FormulaRefAuto{FinitePathTyping}[thm]{2,3}}
  \proofstepwidestar[2,3]{n\in\mathbb N}{\rAEa{14}}
  \proofstepwidestar[2,3]{n+1\in\mathbb N}{\FormulaRefAuto{PeanoAddOneClosure}[thm]{15}}
  \proofstepwidestar[2,3,13]{i\in\mathbb N}{\FormulaRefAuto{PeanoNatSegElementNatural}[thm]{16,13}}
  \proofstepwidestar[1,2,3,4,13]{R(i)}{\rRE{\rUE{12},17}}
  \proofstepwidestar[1,2,3,4,13]{\WordLength{p(i)}=i}{\rRE{18,13}}
  \proofstepwidestar[1,2,3,4]{\forall i\in\NatSeg{n+1}\,\WordLength{p(i)}=i}{\rUI{\rRI{13,19}}}
  \closeproofpartwide
\end{tabproofsplitwide}

\Needspace{10\baselineskip}
\FormulaThmDeltaK[Der vorherige Pfadknoten ist der Rumpf]{%
  \begin{aligned}[t]&\GoodAddressSet P\dsep\GraphStruct{P,\AddressEdges P}\dsep
  \BXXVIIAddressPath pn\dsep p(0)=\BXXVIIAddressRoot\dsep
  i\in\NatSeg n\vdash{}\\[-2pt]
  &p(i+1)\in\NonemptyWords{\{0,1\}}\land
  p(i)=\WordRumpf{\{0,1\}}{p(i+1)}\end{aligned}
}{B27AddressRootPathRumpf}{\DeltaRow{Mengen}{P\dsep p\dsep n\dsep i\dsep a}}
\begin{tabproofwide}
  \proofstepwidestar[1]{\GoodAddressSet P}{\rA}
  \proofstepwidestar[2]{\GraphStruct{P,\AddressEdges P}}{\rA}
  \proofstepwidestar[3]{\BXXVIIAddressPath pn}{\rA}
  \proofstepwidestar[4]{p(0)=\BXXVIIAddressRoot}{\rA}
  \proofstepwidestar[5]{i\in\NatSeg n}{\rA}
  \proofstepwidestar[1,2,3,4]{\BXXVIIAddressForward pn}{\FormulaRefAuto{B27AddressRootPathForward}[thm]{1,2,3,4}}
  \proofstepwidestar[1,2,3,4,5]{\exists a\in\{0,1\}\,p(i+1)=\WordConcat{p(i)}{\LetterWord a}}{\rRE{\rUE{6},5}}
  \proofstepwidestar[8]{a\in\{0,1\}\land p(i+1)=\WordConcat{p(i)}{\LetterWord a}}{\rA}
  \proofstepwidestar[2,3,5]{\begin{aligned}[t]&i\in\mathbb N\land i\in\NatSeg{n+1}\land i+1\in\NatSeg{n+1}\land{}\\[-2pt]&p(i)\in P\land p(i+1)\in P\land{}\\[-2pt]&(p(i),p(i+1))\in\AddressEdges P\end{aligned}}{\FormulaRefAuto{B27PathStepData}[thm]{2,3,5}}
  \proofstepwidestar[2,3,5]{p(i)\in P}{\rAEa{\rAEb{\rAEb{\rAEb{9}}}}}
  \proofstepwidestar[1,2,3,5]{p(i)\in\FiniteWords{\{0,1\}}}{\FormulaRefAuto{B27AddressWordTyping}[thm]{1,10}}
  \proofstepwidestar[1,2,3,5,8]{\WordConcat{p(i)}{\LetterWord a}\in\NonemptyWords{\{0,1\}}}{\FormulaRefAuto{WordAppendNonempty}[thm]{11,\rAEa{8}}}
  \proofstepwidestar[1,2,3,5,8]{\WordRumpf{\{0,1\}}{\WordConcat{p(i)}{\LetterWord a}}=p(i)}{\FormulaRefAuto{WordRumpfAppend}[thm]{11,\rAEa{8}}}
  \proofstepwidestar[8]{\WordConcat{p(i)}{\LetterWord a}=p(i+1)}{\FormulaRefAuto{a=b\vdash b=a}[thm]{\rAEb{8}}}
  \proofstepwidestar[1,2,3,5,8]{p(i+1)\in\NonemptyWords{\{0,1\}}}{\rIE{14,12}}
  \proofstepwidestar[1,2,3,5,8]{\WordRumpf{\{0,1\}}{p(i+1)}=p(i)}{\rIE{14,13}}
  \proofstepwidestar[1,2,3,5,8]{p(i)=\WordRumpf{\{0,1\}}{p(i+1)}}{\FormulaRefAuto{a=b\vdash b=a}[thm]{16}}
  \proofstepwidestar[1,2,3,5,8]{\begin{aligned}[t]&p(i+1)\in\NonemptyWords{\{0,1\}}\land{}\\[-2pt]&p(i)=\WordRumpf{\{0,1\}}{p(i+1)}\end{aligned}}{\rAI{15,17}}
  \proofstepwidestar[1,2,3,4,5]{\begin{aligned}[t]&p(i+1)\in\NonemptyWords{\{0,1\}}\land{}\\[-2pt]&p(i)=\WordRumpf{\{0,1\}}{p(i+1)}\end{aligned}}{\rEE{7,8,18}}
\end{tabproofwide}

\Needspace{10\baselineskip}
\FormulaThmDeltaK[Wurzelpfade mit gleichem Endknoten stimmen überein]{%
  \begin{aligned}[t]&\GoodAddressSet P\dsep\GraphStruct{P,\AddressEdges P}\dsep
  \BXXVIIAddressPath pn\dsep\BXXVIIAddressPath qm\dsep{}\\[-2pt]
  &p(0)=\BXXVIIAddressRoot\dsep q(0)=\BXXVIIAddressRoot\dsep
  p(n)=q(m)\vdash(n,p)=(m,q)\end{aligned}
}{B27AddressRootPathEquality}{\DeltaRow{Mengen}{P\dsep p\dsep q\dsep n\dsep m\dsep i}}
\begin{tabproofsplitwide}
  \proofpartwide{Gleiche Länge}
  \proofstepwidestar[1]{\GoodAddressSet P}{\rA}
  \proofstepwidestar[2]{\GraphStruct{P,\AddressEdges P}}{\rA}
  \proofstepwidestar[3]{\BXXVIIAddressPath pn}{\rA}
  \proofstepwidestar[4]{\BXXVIIAddressPath qm}{\rA}
  \proofstepwidestar[5]{p(0)=\BXXVIIAddressRoot}{\rA}
  \proofstepwidestar[6]{q(0)=\BXXVIIAddressRoot}{\rA}
  \proofstepwidestar[7]{p(n)=q(m)}{\rA}
  \proofstepwidestar[2,3]{n\in\mathbb N\land p\colon\NatSeg{n+1}\to P\land p\colon\NatSeg{n+1}\inj P}{\FormulaRefAuto{FinitePathTyping}[thm]{2,3}}
  \proofstepwidestar[2,3]{n\in\mathbb N}{\rAEa{8}}
  \proofstepwidestar[2,4]{m\in\mathbb N\land q\colon\NatSeg{m+1}\to P\land q\colon\NatSeg{m+1}\inj P}{\FormulaRefAuto{FinitePathTyping}[thm]{2,4}}
  \proofstepwidestar[2,4]{m\in\mathbb N}{\rAEa{10}}
  \proofstepwidestar[2,3]{n\in\NatSeg{n+1}}{\FormulaRefAuto{PeanoNatSegEndpointInAddOne}[thm]{9}}
  \proofstepwidestar[2,4]{m\in\NatSeg{m+1}}{\FormulaRefAuto{PeanoNatSegEndpointInAddOne}[thm]{11}}
  \proofstepwidestar[1,2,3,5]{\forall i\in\NatSeg{n+1}\,\WordLength{p(i)}=i}{\FormulaRefAuto{B27AddressRootPathRanks}[thm]{1,2,3,5}}
  \proofstepwidestar[1,2,4,6]{\forall i\in\NatSeg{m+1}\,\WordLength{q(i)}=i}{\FormulaRefAuto{B27AddressRootPathRanks}[thm]{1,2,4,6}}
  \proofstepwidestar[1,2,3,5]{\WordLength{p(n)}=n}{\rRE{\rUE{14},12}}
  \proofstepwidestar[1,2,4,6]{\WordLength{q(m)}=m}{\rRE{\rUE{15},13}}
  \proofstepwidestar[1,2,3,5,7]{\WordLength{q(m)}=n}{\rIE{7,16}}
  \proofstepwidestar[1,2,3,4,5,6,7]{n=m}{\FormulaRefAuto{a=b\dsep a=c\vdash b=c}[thm]{18,17}}
  \proofstepwidestar[1,2,3,4,5,6,7]{m=n}{\FormulaRefAuto{a=b\vdash b=a}[thm]{19}}
  \proofstepwidestar[1,2,3,4,5,6,7]{\BXXVIIAddressPath qn}{\rIE{20,4}}
  \proofstepwidestar[1,2,3,4,5,6,7]{p(n)=q(n)}{\rIE{20,7}}
  \closeproofpartwide
  \proofpartwide{Vom Endknoten rückwärts}
  \setcounter{proofstepnr}{22}
  \proofstepwidestar[23]{i\in\NatSeg n}{\rA}
  \proofstepwidestar[24]{p(i+1)=q(i+1)}{\rA}
  \proofstepwidestar[1,2,3,5,23]{p(i+1)\in\NonemptyWords{\{0,1\}}\land p(i)=\WordRumpf{\{0,1\}}{p(i+1)}}{\FormulaRefAuto{B27AddressRootPathRumpf}[thm]{1,2,3,5,23}}
  \proofstepwidestar[1,2,3,4,5,6,7,23]{q(i+1)\in\NonemptyWords{\{0,1\}}\land q(i)=\WordRumpf{\{0,1\}}{q(i+1)}}{\FormulaRefAuto{B27AddressRootPathRumpf}[thm]{1,2,21,6,23}}
  \proofstepwidestar[1,2,3,5,23,24]{p(i)=\WordRumpf{\{0,1\}}{q(i+1)}}{\rIE{24,\rAEb{25}}}
  \proofstepwidestar[1,2,3,4,5,6,7,23,24]{p(i)=q(i)}{\FormulaRefAuto{a=b\dsep c=b\vdash a=c}[thm]{27,\rAEb{26}}}
  \proofstepwidestar[1,2,3,4,5,6,7,23]{p(i+1)=q(i+1)\rightarrow p(i)=q(i)}{\rRI{24,28}}
  \proofstepwidestar[1,2,3,4,5,6,7]{\forall i\in\NatSeg n\,(p(i+1)=q(i+1)\rightarrow p(i)=q(i))}{\rUI{\rRI{23,29}}}
  \proofstepwidestar[1,2,3,4,5,6,7]{\forall i\in\NatSeg{n+1}\,p(i)=q(i)}{\FormulaRefAuto{B27FiniteBackwardInduction}[thm]{9,22,30}}
  \proofstepwidestar[2,3]{p\colon\NatSeg{n+1}\to P}{\FormulaRefAuto{DirectedPathFunction}[thm]{2,3}}
  \proofstepwidestar[1,2,3,4,5,6,7]{q\colon\NatSeg{n+1}\to P}{\FormulaRefAuto{DirectedPathFunction}[thm]{2,21}}
  \proofstepwidestar[1,2,3,4,5,6,7]{p=q}{\FormulaRefAuto{FunctionExtensionality}[thm]{32,33,31}}
  \proofstepwidestar[1,2,3,4,5,6,7]{(n,p)=(m,q)}{\FormulaRefAuto{a=c,\, b=d\vdash (a,b) = (c,d)}[thm]{19,34}}
  \closeproofpartwide
\end{tabproofsplitwide}

\Needspace{10\baselineskip}
\FormulaThmDeltaK[Eindeutigkeit gespeicherter Wurzelpfade]{%
  \begin{aligned}[t]&\GoodAddressSet P\dsep\GraphStruct{P,\AddressEdges P}\dsep
  z_1\in\GraphPathsBetween P{\AddressEdges P}{\BXXVIIAddressRoot}q\dsep{}\\[-2pt]
  &z_2\in\GraphPathsBetween P{\AddressEdges P}{\BXXVIIAddressRoot}q\vdash z_1=z_2\end{aligned}
}{B27AddressStoredRootPathsEqual}{\DeltaRow{Mengen}{P\dsep q\dsep z_1\dsep z_2\dsep n\dsep m\dsep p\dsep s}}
\begin{tabproofwide}
  \proofstepwidestar[1]{\GoodAddressSet P}{\rA}
  \proofstepwidestar[2]{\GraphStruct{P,\AddressEdges P}}{\rA}
  \proofstepwidestar[3]{z_1\in\GraphPathsBetween P{\AddressEdges P}{\BXXVIIAddressRoot}q}{\rA}
  \proofstepwidestar[4]{z_2\in\GraphPathsBetween P{\AddressEdges P}{\BXXVIIAddressRoot}q}{\rA}
  \proofstepwidestar[2]{z_1\in\GraphPathsBetween P{\AddressEdges P}{\BXXVIIAddressRoot}q\leftrightarrow\exists n\in\mathbb N\,\exists p\in\powerset(\mathbb N\times P)\,(z_1=(n,p)\land\BXXVIIAddressPath pn\land p(0)=\BXXVIIAddressRoot\land p(n)=q)}{\FormulaRefAuto{FiniteSimplePathFamilyMembership}[thm]{2}}
  \proofstepwidestar[2,3]{\exists n\in\mathbb N\,\exists p\in\powerset(\mathbb N\times P)\,(z_1=(n,p)\land\BXXVIIAddressPath pn\land p(0)=\BXXVIIAddressRoot\land p(n)=q)}{\rRE{3,\rLREa{5}}}
  \proofstepwidestar[7]{n\in\mathbb N\land p\in\powerset(\mathbb N\times P)\land z_1=(n,p)\land\BXXVIIAddressPath pn\land p(0)=\BXXVIIAddressRoot\land p(n)=q}{\rA}
  \proofstepwidestar[2]{z_2\in\GraphPathsBetween P{\AddressEdges P}{\BXXVIIAddressRoot}q\leftrightarrow\exists m\in\mathbb N\,\exists s\in\powerset(\mathbb N\times P)\,(z_2=(m,s)\land\BXXVIIAddressPath sm\land s(0)=\BXXVIIAddressRoot\land s(m)=q)}{\FormulaRefAuto{FiniteSimplePathFamilyMembership}[thm]{2}}
  \proofstepwidestar[2,4]{\exists m\in\mathbb N\,\exists s\in\powerset(\mathbb N\times P)\,(z_2=(m,s)\land\BXXVIIAddressPath sm\land s(0)=\BXXVIIAddressRoot\land s(m)=q)}{\rRE{4,\rLREa{8}}}
  \proofstepwidestar[10]{m\in\mathbb N\land s\in\powerset(\mathbb N\times P)\land z_2=(m,s)\land\BXXVIIAddressPath sm\land s(0)=\BXXVIIAddressRoot\land s(m)=q}{\rA}
  \proofstepwidestar[7,10]{p(n)=s(m)}{\FormulaRefAuto{a=b\dsep c=b\vdash a=c}[thm]{\rAEb{\rAEb{\rAEb{\rAEb{\rAEb{7}}}}},\rAEb{\rAEb{\rAEb{\rAEb{\rAEb{10}}}}}}}
  \proofstepwidestar[1,2,7,10]{(n,p)=(m,s)}{\FormulaRefAuto{B27AddressRootPathEquality}[thm]{1,2,\rAEa{\rAEb{\rAEb{\rAEb{7}}}},\rAEa{\rAEb{\rAEb{\rAEb{10}}}},\rAEa{\rAEb{\rAEb{\rAEb{\rAEb{7}}}}},\rAEa{\rAEb{\rAEb{\rAEb{\rAEb{10}}}}},11}}
  \proofstepwidestar[7]{z_1=(n,p)}{\rAEa{\rAEb{\rAEb{7}}}}
  \proofstepwidestar[10]{z_2=(m,s)}{\rAEa{\rAEb{\rAEb{10}}}}
  \proofstepwidestar[1,2,7,10]{z_1=(m,s)}{\rChain{13,12}}
  \proofstepwidestar[1,2,7,10]{z_1=z_2}{\FormulaRefAuto{a=b\dsep c=b\vdash a=c}[thm]{15,14}}
  \proofstepwidestar[1,2,4,7]{z_1=z_2}{\rEEStar{9,10,16}}
  \proofstepwidestar[1,2,3,4]{z_1=z_2}{\rEEStar{6,7,17}}
\end{tabproofwide}

\Needspace{10\baselineskip}
\FormulaThmDeltaK[Jede Adresse besitzt genau einen Wurzelpfad]{%
  \GoodAddressSet P\dsep\GraphStruct{P,\AddressEdges P}\dsep q\in P\vdash
    \exists!z\,(z\in\GraphPathsBetween P{\AddressEdges P}{\BXXVIIAddressRoot}q)
}{B27AddressRootPathUnique}{%
  \DeltaRow{Mengen}{P\dsep q\dsep z\dsep z_1\dsep z_2}
}
\begin{tabproofwide}
  \proofstepwidestar[1]{\GoodAddressSet P}{\rA}
  \proofstepwidestar[2]{\GraphStruct{P,\AddressEdges P}}{\rA}
  \proofstepwidestar[3]{q\in P}{\rA}
  \proofstepwidestar[1,2,3]{%
    \exists z\,(z\in\GraphPathsBetween P{\AddressEdges P}{\BXXVIIAddressRoot}q)}{%
    \FormulaRefAuto{B27AddressRootPathExists}[thm]{1,2,3}}
  \proofstepwidestar[5]{z_1\in\GraphPathsBetween P{\AddressEdges P}{\BXXVIIAddressRoot}q}{\rA}
  \proofstepwidestar[6]{z_2\in\GraphPathsBetween P{\AddressEdges P}{\BXXVIIAddressRoot}q}{\rA}
  \proofstepwidestar[1,2,5,6]{z_1=z_2}{%
    \FormulaRefAuto{B27AddressStoredRootPathsEqual}[thm]{1,2,5,6}}
  \proofstepwidestar[1,2,3]{%
    \exists!z\,(z\in\GraphPathsBetween P{\AddressEdges P}{\BXXVIIAddressRoot}q)}{%
    \UEI{4,5,6,7}}
\end{tabproofwide}

\subsection{Eltern und Kinder als Wortoperationen}

Sobald der Adressgraph als verwurzelter Baum vorliegt, lassen sich seine
Eltern und Kinder mit den Wortoperationen beschreiben. Ein Elternschritt
entfernt das letzte Bit; ein Kinderschritt fügt ein Bit an, sofern die verlängerte Adresse vorhanden ist.

\Needspace{10\baselineskip}
\FormulaThmDeltaK[Der Baumelter einer Adresse ist ihr Rumpf]{%
  \begin{aligned}[t]&\GoodAddressSet P\dsep
  \RootedTreeStruct{P,\AddressEdges P,\BXXVIIAddressRoot}\dsep
  \TreeParentRel P{\AddressEdges P}{\BXXVIIAddressRoot}p q\vdash{}\\[-2pt]
  &q\in\NonemptyWords{\{0,1\}}\land p=\WordRumpf{\{0,1\}}q\end{aligned}
}{B27AddressParentIsRumpf}{\DeltaRow{Mengen}{P\dsep p\dsep q\dsep k\dsep s}}
\begin{tabproofwide}
  \proofstepwidestar[1]{\GoodAddressSet P}{\rA}
  \proofstepwidestar[2]{\RootedTreeStruct{P,\AddressEdges P,\BXXVIIAddressRoot}}{\rA}
  \proofstepwidestar[3]{\TreeParentRel P{\AddressEdges P}{\BXXVIIAddressRoot}p q}{\rA}
  \proofstepwidestar[2]{\TreeStruct{P,\AddressEdges P}}{\FormulaRefAuto{RootedTreeTreeAxiom}[ax]{2}}
  \proofstepwidestar[2]{\SimpleUndirectedGraphStruct{P,\AddressEdges P}}{\FormulaRefAuto{TreeGraphAxiom}[ax]{4}}
  \proofstepwidestar[2]{\UndirectedGraphStruct{P,\AddressEdges P}\land\GraphStruct{P,\AddressEdges P}}{\FormulaRefAuto{SimpleUndirectedGraphAccess}[thm]{5}}
  \proofstepwidestar[2]{\GraphStruct{P,\AddressEdges P}}{\rAEb{6}}
  \proofstepwidestar[2,3]{\exists k\in\mathbb N\,\exists s\in\powerset(\mathbb N\times P)\,((k+1,s)\in\GraphPathsBetween P{\AddressEdges P}{\BXXVIIAddressRoot}q\land s(k)=p)}{\FormulaRefAuto{TreeParentAccess}[thm]{2,3}}
  \proofstepwidestar[9]{k\in\mathbb N\land s\in\powerset(\mathbb N\times P)\land(k+1,s)\in\GraphPathsBetween P{\AddressEdges P}{\BXXVIIAddressRoot}q\land s(k)=p}{\rA}
  \proofstepwidestar[9]{k\in\mathbb N}{\rAEa{9}}
  \proofstepwidestar[9]{(k+1,s)\in\GraphPathsBetween P{\AddressEdges P}{\BXXVIIAddressRoot}q}{\rAEa{\rAEb{\rAEb{9}}}}
  \proofstepwidestar[9]{s(k)=p}{\rAEb{\rAEb{\rAEb{9}}}}
  \proofstepwidestar[2,9]{k+1\in\mathbb N\land s\in\powerset(\mathbb N\times P)\land\BXXVIIAddressPath s{k+1}\land s(0)=\BXXVIIAddressRoot\land s(k+1)=q}{\FormulaRefAuto{B27SimplePathFamilyPairElimination}[thm]{7,11}}
  \proofstepwidestar[2,9]{\begin{aligned}[t]&\BXXVIIAddressPath s{k+1}\\[-2pt]&\quad\land s(0)=\BXXVIIAddressRoot\land s(k+1)=q\end{aligned}}{\rAEb{\rAEb{13}}}
  \proofstepwidestar[9]{k\in\NatSeg{k+1}}{\FormulaRefAuto{PeanoNatSegEndpointInAddOne}[thm]{10}}
  \proofstepwidestar[1,2,9]{\begin{aligned}[t]&s(k+1)\in\NonemptyWords{\{0,1\}}\\[-2pt]&\quad\land s(k)=\WordRumpf{\{0,1\}}{s(k+1)}\end{aligned}}{\FormulaRefAuto{B27AddressRootPathRumpf}[thm]{1,7,\rAEa{14},\rAEa{\rAEb{14}},15}}
  \proofstepwidestar[1,2,9]{q\in\NonemptyWords{\{0,1\}}\land p=\WordRumpf{\{0,1\}}q}{\rIE{12,\rAEb{\rAEb{14}},16}}
  \proofstepwidestar[1,2,3]{q\in\NonemptyWords{\{0,1\}}\land p=\WordRumpf{\{0,1\}}q}{\rEEStar{8,9,17}}
\end{tabproofwide}

\Needspace{10\baselineskip}
\FormulaThmDeltaK[Der Rumpf ist ein Baumelter]{%
  \begin{aligned}[t]&\GoodAddressSet P\dsep
  \RootedTreeStruct{P,\AddressEdges P,\BXXVIIAddressRoot}\dsep{}\\[-2pt]
  &q\in P\dsep q\neq\BXXVIIAddressRoot\dsep p=\WordRumpf{\{0,1\}}q\vdash{}\\[-2pt]
  &\TreeParentRel P{\AddressEdges P}{\BXXVIIAddressRoot}p q\end{aligned}
}{B27AddressRumpfIsParent}{\DeltaRow{Mengen}{P\dsep p\dsep q\dsep b}}
\begin{tabproofwide}
  \proofstepwidestar[1]{\GoodAddressSet P}{\rA}
  \proofstepwidestar[2]{\RootedTreeStruct{P,\AddressEdges P,\BXXVIIAddressRoot}}{\rA}
  \proofstepwidestar[3]{q\in P}{\rA}
  \proofstepwidestar[4]{q\neq\BXXVIIAddressRoot}{\rA}
  \proofstepwidestar[5]{p=\WordRumpf{\{0,1\}}q}{\rA}
  \proofstepwidestar[]{q\notin\{\BXXVIIAddressRoot\}\leftrightarrow q\neq\BXXVIIAddressRoot}{\FormulaRefAuto{x\notin\{a\}\eqvdash x\neq a}[thm]}
  \proofstepwidestar[4]{q\notin\{\BXXVIIAddressRoot\}}{\rRE{4,\rLREb{6}}}
  \proofstepwidestar[3,4]{q\in P\setminus\{\BXXVIIAddressRoot\}}{\FormulaRefAuto{x\in A\dsep x\notin B\vdash x\in A\setminus B}[thm]{3,7}}
  \proofstepwidestar[2,3,4]{\exists b\in P\,\TreeParentRel P{\AddressEdges P}{\BXXVIIAddressRoot}b q}{\FormulaRefAuto{TreeParentExists}[thm]{2,8}}
  \proofstepwidestar[10]{b\in P\land\TreeParentRel P{\AddressEdges P}{\BXXVIIAddressRoot}b q}{\rA}
  \proofstepwidestar[1,2,10]{q\in\NonemptyWords{\{0,1\}}\land b=\WordRumpf{\{0,1\}}q}{\FormulaRefAuto{B27AddressParentIsRumpf}[thm]{1,2,\rAEb{10}}}
  \proofstepwidestar[1,2,5,10]{b=p}{\FormulaRefAuto{a=b\dsep c=b\vdash a=c}[thm]{\rAEb{11},5}}
  \proofstepwidestar[1,2,5,10]{\TreeParentRel P{\AddressEdges P}{\BXXVIIAddressRoot}p q}{\rIE{12,\rAEb{10}}}
  \proofstepwidestar[1,2,3,4,5]{\TreeParentRel P{\AddressEdges P}{\BXXVIIAddressRoot}p q}{\rEE{9,10,13}}
\end{tabproofwide}

\Needspace{10\baselineskip}
\FormulaThmDeltaK[Adresskinder sind die vorhandenen Fortsetzungen]{%
  \begin{aligned}[t]&\GoodAddressSet P\dsep
  \RootedTreeStruct{P,\AddressEdges P,\BXXVIIAddressRoot}\dsep p\in P\vdash{}\\[-2pt]
  &\TreeChildren P{\AddressEdges P}{\BXXVIIAddressRoot}p
  =\{\WordConcat p{\LetterWord i}\mid i\in\{0,1\}\}\cap P\end{aligned}
}{B27AddressChildrenSet}{\DeltaRow{Mengen}{P\dsep p\dsep q\dsep i}}
\begin{tabproofsplitwide}
  \proofpartwide{Vom Kind zur Fortsetzung}
  \proofstepwidestar[1]{\GoodAddressSet P}{\rA}
  \proofstepwidestar[2]{\RootedTreeStruct{P,\AddressEdges P,\BXXVIIAddressRoot}}{\rA}
  \proofstepwidestar[3]{p\in P}{\rA}
  \proofstepwidestar[1,3]{p\in\FiniteWords{\{0,1\}}}{\FormulaRefAuto{B27AddressWordTyping}[thm]{1,3}}
  \proofstepwidestar[1,3]{p\in\BinaryAddresses}{\FormulaRefAuto{BinaryAddressWordSetDef}[def]{4}}
  \proofstepwidestar[6]{q\in\TreeChildren P{\AddressEdges P}{\BXXVIIAddressRoot}p}{\rA}
  \proofstepwidestar[2,3]{q\in\TreeChildren P{\AddressEdges P}{\BXXVIIAddressRoot}p\leftrightarrow\TreeParentRel P{\AddressEdges P}{\BXXVIIAddressRoot}p q}{\FormulaRefAuto{TreeChildrenMembership}[thm]{2,3}}
  \proofstepwidestar[2,3,6]{\TreeParentRel P{\AddressEdges P}{\BXXVIIAddressRoot}p q}{\rRE{\rLREa{7},6}}
  \proofstepwidestar[2,3,6]{q\in P\setminus\{\BXXVIIAddressRoot\}}{\FormulaRefAuto{TreeParentAccess}[thm]{2,8}}
  \proofstepwidestar[2,3,6]{q\in P}{\FormulaRefAuto{x\in A\setminus B\vdash x\in A}[thm]{9}}
  \proofstepwidestar[1,2,3,6]{q\in\NonemptyWords{\{0,1\}}\land p=\WordRumpf{\{0,1\}}q}{\FormulaRefAuto{B27AddressParentIsRumpf}[thm]{1,2,8}}
  \proofstepwidestar[1,2,3,6]{q=\WordConcat{\WordRumpf{\{0,1\}}q}{\LetterWord{\WordEnd{\{0,1\}}q}}}{\FormulaRefAuto{WordRightDecompositionEquation}[thm]{\rAEa{11}}}
  \proofstepwidestar[1,2,3,6]{\WordRumpf{\{0,1\}}q=p}{\FormulaRefAuto{a=b\vdash b=a}[thm]{\rAEb{11}}}
  \proofstepwidestar[1,2,3,6]{q=\WordConcat p{\LetterWord{\WordEnd{\{0,1\}}q}}}{\rIE{13,12}}
  \proofstepwidestar[1,2,3,6]{\WordEnd{\{0,1\}}q\in\{0,1\}}{\FormulaRefAuto{WordEndTyping}[thm]{\rAEa{11}}}
  \proofstepwidestar[1,2,3,6]{\exists i\in\{0,1\}\,q=\WordConcat p{\LetterWord i}}{\rEI{\rAI{15,14}}}
  \proofstepwidestar[1,3]{\begin{aligned}[t]&\forall q\,\Bigl(q\in\{\WordConcat p{\LetterWord i}\mid i\in\{0,1\}\}\\[-2pt]&\quad\leftrightarrow\exists i\in\{0,1\}\,q=\WordConcat p{\LetterWord i}\Bigr)\end{aligned}}{\FormulaRefAuto{B27BinaryChildTermSetDef}[def]{5}}
  \proofstepwidestar[1,3]{q\in\{\WordConcat p{\LetterWord i}\mid i\in\{0,1\}\}\leftrightarrow\exists i\in\{0,1\}\,q=\WordConcat p{\LetterWord i}}{\rUE{17}}
  \proofstepwidestar[1,2,3,6]{q\in\{\WordConcat p{\LetterWord i}\mid i\in\{0,1\}\}}{\rRE{\rLREb{18},16}}
  \proofstepwidestar[]{q\in\{\WordConcat p{\LetterWord i}\mid i\in\{0,1\}\}\cap P\leftrightarrow(q\in\{\WordConcat p{\LetterWord i}\mid i\in\{0,1\}\}\land q\in P)}{\FormulaRefAuto{x\in A\cap B\eqvdash x\in A\land x\in B}[thm]}
  \proofstepwidestar[1,2,3,6]{q\in\{\WordConcat p{\LetterWord i}\mid i\in\{0,1\}\}\cap P}{\rRE{\rAI{19,10},\rLREb{20}}}
  \closeproofpartwide
  \proofpartwide{Von der Fortsetzung zum Kind}
  \setcounter{proofstepnr}{21}
  \proofstepwidestar[22]{q\in\{\WordConcat p{\LetterWord i}\mid i\in\{0,1\}\}\cap P}{\rA}
  \proofstepwidestar[22]{q\in\{\WordConcat p{\LetterWord i}\mid i\in\{0,1\}\}\land q\in P}{\rRE{22,\rLREa{20}}}
  \proofstepwidestar[1,3,22]{\exists i\in\{0,1\}\,q=\WordConcat p{\LetterWord i}}{\rRE{\rLREa{18},\rAEa{23}}}
  \proofstepwidestar[25]{i\in\{0,1\}\land q=\WordConcat p{\LetterWord i}}{\rA}
  \proofstepwidestar[1,3,25]{\WordConcat p{\LetterWord i}\neq\BXXVIIAddressRoot}{\FormulaRefAuto{B27WordAppendNotEmpty}[thm]{4,\rAEa{25}}}
  \proofstepwidestar[25]{\WordConcat p{\LetterWord i}=q}{\FormulaRefAuto{a=b\vdash b=a}[thm]{\rAEb{25}}}
  \proofstepwidestar[1,3,25]{q\neq\BXXVIIAddressRoot}{\rIE{27,26}}
  \proofstepwidestar[1,3,25]{\WordRumpf{\{0,1\}}{\WordConcat p{\LetterWord i}}=p}{\FormulaRefAuto{WordRumpfAppend}[thm]{4,\rAEa{25}}}
  \proofstepwidestar[1,3,25]{\WordRumpf{\{0,1\}}q=p}{\rIE{27,29}}
  \proofstepwidestar[1,3,25]{p=\WordRumpf{\{0,1\}}q}{\FormulaRefAuto{a=b\vdash b=a}[thm]{30}}
  \proofstepwidestar[1,2,3,22,25]{\TreeParentRel P{\AddressEdges P}{\BXXVIIAddressRoot}p q}{\FormulaRefAuto{B27AddressRumpfIsParent}[thm]{1,2,\rAEb{23},28,31}}
  \proofstepwidestar[1,2,3,22,25]{q\in\TreeChildren P{\AddressEdges P}{\BXXVIIAddressRoot}p}{\rRE{\rLREb{7},32}}
  \proofstepwidestar[1,2,3,22]{q\in\TreeChildren P{\AddressEdges P}{\BXXVIIAddressRoot}p}{\rEE{24,25,33}}
  \proofstepwidestar[1,2,3]{q\in\TreeChildren P{\AddressEdges P}{\BXXVIIAddressRoot}p\leftrightarrow q\in\{\WordConcat p{\LetterWord i}\mid i\in\{0,1\}\}\cap P}{\rLRI{\rRI{6,21},\rRI{22,34}}}
  \proofstepwidestar[1,2,3]{\begin{aligned}[t]&\TreeChildren P{\AddressEdges P}{\BXXVIIAddressRoot}p\\[-2pt]&\quad=\{\WordConcat p{\LetterWord i}\mid i\in\{0,1\}\}\cap P\end{aligned}}{\FormulaRefAuto{\forall x\,(x\in A\leftrightarrow x\in B)\vdash A=B}[ax]{\rUI{35}}}
  \closeproofpartwide
\end{tabproofsplitwide}



\subsection{Der Positionsgraph als voller geordneter Binärbaum}

\FormulaThmDeltaK[Volle Präfixmengen tragen endliche geordnete volle
Binärbäume]{%
  \GoodAddressSet P\vdash
  \begin{aligned}[t]
    &\text{(i)}\;\Finite P\\[-2pt]
    &\text{(ii)}\;\OrderedFullBinaryTreeStruct{P,\AddressEdges P,
      \EmptyWord{\{0,1\}},\AddressLeft P,\AddressRight P}
  \end{aligned}%
}{GoodBinaryAddressSetIsFiniteOrderedFullBinaryTree}{%
  \DeltaRow{Mengen}{P\dsep p\dsep q\dsep i\dsep z}
  \DeltaRow{Funktionen}{\AddressLeft P\dsep\AddressRight P}
}
\Needspace{11\baselineskip}
\begin{tabproofsplitwide}
  \proofpartwideindR[Endlichkeit der Positionsmenge]{%
    \GoodAddressSet P\vdash\Finite P
  }{\GoodAddressSet P\vdash\Finite P}[GoodBinaryAddressSetFinite]
    \proofstepwidestar[1]{%
      \GoodAddressSet P}{%
      \rA}
    \proofstepwidestar[1]{%
      \Finite P}{%
      \FormulaRefAuto{GoodBinaryAddressSetCriterion}[thm]{1}}
  \closeproofpartwideindR

  \begin{BandTwentySevenProofCases}

  \proofpartwideindR[Schlichter Adressgraph]{%
    \GoodAddressSet P\vdash
      \SimpleUndirectedGraphStruct{P,\AddressEdges P}
  }{\GoodAddressSet P\vdash
      \SimpleUndirectedGraphStruct{P,\AddressEdges P}}%
    [GoodBinaryAddressGraphSimple]
    \proofstepwidestar[1]{\GoodAddressSet P}{\rA}
    \proofstepwidestar[1]{\AddressEdges P\subseteq P\times P}{%
      \FormulaRefAuto{BinaryAddressGraphDef}[def]{1}}
    \proofstepwidestar[3]{p,q\in P}{\rA}
    \proofstepwidestar[4]{(p,q)\in\AddressEdges P}{\rA}
    \proofstepwidestar[1,3,4]{(q,p)\in\AddressEdges P}{%
      \FormulaRefAuto{BinaryAddressGraphDef}[def]{1,3,4}}
    \proofstepwidestar[1]{%
      \forall p,q\in P\,
      ((p,q)\in\AddressEdges P\rightarrow(q,p)\in\AddressEdges P)}{%
      \rUI{\rRI{3,\rUI{\rRI{4,5}}}}}
    \proofstepwidestar[7]{p\in P}{\rA}
    \proofstepwidestar[8]{(p,p)\in\AddressEdges P}{\rA}
    \proofstepwidestar[1,8]{%
      \exists i\in\{0,1\}\,p=\WordConcat p{\LetterWord i}}{%
      \FormulaRefAuto{BinaryAddressGraphDef}[def]{1,8}}
    \proofstepwidestar[10]{%
      i\in\{0,1\}\land p=\WordConcat p{\LetterWord i}}{\rA}
    \proofstepwidestarforcebreak[1,7,10]{%
      \WordLength p=\WordLength p+1}{\rIE{\FormulaRefAuto{a=b\vdash b=a}[thm]{\rAEb{10}},\FormulaRefAuto{LetterWordLengthOne}[thm]{\rAEa{10}},\FormulaRefAuto{WordConcatenationLength}[thm]{\FormulaRefAuto{GoodBinaryAddressSetCriterion}[thm]{1,7},\FormulaRefAuto{NonemptyWordFinite}[thm]{\FormulaRefAuto{LetterWordNonempty}[thm]{\rAEa{10}}}}}}
    \proofstepwidestar[]{0\in\mathbb N}{%
      \FormulaRefAuto{PeanoZero}[ax]}
    \proofstepwidestar[1,7]{p\in\FiniteWords{\{0,1\}}}{%
      \FormulaRefAuto{GoodBinaryAddressSetCriterion}[thm]{1,7}}
    \proofstepwidestar[1,7]{\WordLength p\in\mathbb N}{%
      \FormulaRefAuto{FiniteWordLengthNatural}[thm]{%
          13}}
    \proofstepwidestar[1,7]{%
      \WordLength p+1\neq\WordLength p}{%
      \FormulaRefAuto{PeanoAddPositiveDistance}[thm]{%
        14,
        12},
      \FormulaRefAuto{PeanoOne}[def]}
    \proofstepwidestar[1,7,10]{%
      \WordLength p+1=\WordLength p}{%
      \FormulaRefAuto{a=b\vdash b=a}[thm]{11}}
    \proofstepwidestar[1,7,10]{\bot}{\rBI{16,15}}
    \proofstepwidestar[1,7,8]{\bot}{\rEE{9,10,17}}
    \proofstepwidestar[1,7]{(p,p)\notin\AddressEdges P}{\rCI{8,18}}
    \proofstepwidestar[1]{%
      \forall p\in P\,((p,p)\notin\AddressEdges P)}{%
      \rUI{\rRI{7,19}}}
    \proofstepwidestar[1]{%
      \SimpleUndirectedGraphStruct{P,\AddressEdges P}}{%
      \FormulaRefAuto{SimpleUndirectedGraphIntroduction}[thm]{2,6,20}}
  \closeproofpartwideindR

  \Needspace{11\baselineskip}
  \proofpartwideindR[Eindeutige Wurzelwege]{%
    \GoodAddressSet P\vdash
    \forall q\in P\,\exists!z\,
      \bigl(z\in\GraphPathsBetween{P}{\AddressEdges P}
        {\EmptyWord{\{0,1\}}}{q}\bigr)
  }{%
    \GoodAddressSet P\vdash
    \forall q\in P\,\exists!z\,
      \bigl(z\in\GraphPathsBetween{P}{\AddressEdges P}
        {\EmptyWord{\{0,1\}}}{q}\bigr)
  }[GoodBinaryAddressRootPathsUnique]
\proofstepwidestar[1]{\GoodAddressSet P}{\rA}
    \proofstepwidestar[1]{\SimpleUndirectedGraphStruct{P,\AddressEdges P}}{\FormulaRefAuto{GoodBinaryAddressGraphSimple}[thm]{1}}
    \proofstepwidestar[1]{\UndirectedGraphStruct{P,\AddressEdges P}\land\GraphStruct{P,\AddressEdges P}}{\FormulaRefAuto{SimpleUndirectedGraphAccess}[thm]{2}}
    \proofstepwidestar[1]{\GraphStruct{P,\AddressEdges P}}{\rAEb{3}}
    \proofstepwidestar[5]{q\in P}{\rA}
    \proofstepwidestar[1,5]{\exists!z\,(z\in\GraphPathsBetween{P}{\AddressEdges P}{\EmptyWord{\{0,1\}}}{q})}{\FormulaRefAuto{B27AddressRootPathUnique}[thm]{1,4,5}}
    \proofstepwidestar[1]{\forall q\in P\,\exists!z\,(z\in\GraphPathsBetween{P}{\AddressEdges P}{\EmptyWord{\{0,1\}}}{q})}{\rUI{\rRI{5,6}}}
  \closeproofpartwideindR

  \Needspace{12\baselineskip}
  \proofpartwideindR[Adresskinder]{%
    \begin{aligned}[t]
    &\GoodAddressSet P\dsep p\in P\vdash{}\\[-2pt]
    &\quad\TreeChildren{P}{\AddressEdges P}{\EmptyWord{\{0,1\}}}{p}
      =\{\WordConcat p{\LetterWord i}\mid i\in\{0,1\}\}\cap P
    \end{aligned}
  }{%
    \begin{aligned}[t]
    &\GoodAddressSet P\dsep p\in P\vdash{}\\[-2pt]
    &\quad\TreeChildren{P}{\AddressEdges P}{\EmptyWord{\{0,1\}}}{p}
      =\{\WordConcat p{\LetterWord i}\mid i\in\{0,1\}\}\cap P
    \end{aligned}
  }[GoodBinaryAddressChildren]
\proofstepwidestar[1]{\GoodAddressSet P}{\rA}
    \proofstepwidestar[2]{p\in P}{\rA}
    \proofstepwidestar[1]{\SimpleUndirectedGraphStruct{P,\AddressEdges P}}{\FormulaRefAuto{GoodBinaryAddressGraphSimple}[thm]{1}}
    \proofstepwidestar[1]{\EmptyWord{\{0,1\}}\in P}{\FormulaRefAuto{GoodBinaryAddressSetRootAxiom}[ax]{1}}
    \proofstepwidestar[1]{\forall q\in P\,\exists!z\,(z\in\GraphPathsBetween{P}{\AddressEdges P}{\EmptyWord{\{0,1\}}}{q})}{\FormulaRefAuto{GoodBinaryAddressRootPathsUnique}[thm]{1}}
    \proofstepwidestar[1]{\RootedTreeStruct{P,\AddressEdges P,\EmptyWord{\{0,1\}}}}{\FormulaRefAuto{RootedTreeFromUniqueRootPaths}[thm]{3,4,5}}
    \proofstepwidestar[1,2]{\begin{aligned}[t]
      &\TreeChildren{P}{\AddressEdges P}{\EmptyWord{\{0,1\}}}{p}\\[-2pt]
      &\quad=\{\WordConcat p{\LetterWord i}\mid i\in\{0,1\}\}\cap P
    \end{aligned}}{\FormulaRefAuto{B27AddressChildrenSet}[thm]{1,6,2}}
  \closeproofpartwideindR

  \proofpartwideindR[Kinder und Fülle]{%
    \GoodAddressSet P\vdash
      \FullBinaryTreeStruct{P,\AddressEdges P,\EmptyWord{\{0,1\}}}
  }{\GoodAddressSet P\vdash
      \FullBinaryTreeStruct{P,\AddressEdges P,\EmptyWord{\{0,1\}}}}%
    [GoodBinaryAddressGraphFullBinary]
    \proofstepwidestar[1]{\GoodAddressSet P}{\rA}
    \proofstepwidestar[1]{%
      \SimpleUndirectedGraphStruct{P,\AddressEdges P}}{%
      \FormulaRefAuto{GoodBinaryAddressGraphSimple}[thm]{1}}
    \proofstepwidestar[1]{\EmptyWord{\{0,1\}}\in P}{%
      \FormulaRefAuto{GoodBinaryAddressSetCriterion}[thm]{1}}
    \proofstepwidestar[1]{%
      \forall q\in P\,\exists!z\,
      \bigl(z\in\GraphPathsBetween{P}{\AddressEdges P}
        {\EmptyWord{\{0,1\}}}{q}\bigr)}{%
      \FormulaRefAuto{GoodBinaryAddressRootPathsUnique}[thm]{1}}
    \proofstepwidestar[1]{%
      \RootedTreeStruct{P,\AddressEdges P,\EmptyWord{\{0,1\}}}}{%
      \FormulaRefAuto{RootedTreeFromUniqueRootPaths}[thm]{2,3,4}}
    \proofstepwidestar[6]{p\in P}{\rA}
    \proofstepwidestar[1,6]{%
      \TreeChildren{P}{\AddressEdges P}{\EmptyWord{\{0,1\}}}{p}
      =\{\WordConcat p{\LetterWord i}\mid i\in\{0,1\}\}\cap P}{%
      \FormulaRefAuto{GoodBinaryAddressChildren}[thm]{1,6}}
    \proofstepwidestarforcebreak[1,6]{%
      \TreeChildren{P}{\AddressEdges P}{\EmptyWord{\{0,1\}}}{p}
      \subseteq\{\WordConcat p{\LetterWord0},
                   \WordConcat p{\LetterWord1}\}}{\rIE{\FormulaRefAuto{a=b\vdash b=a}[thm]{\rIE{\FormulaRefAuto{B27BinaryChildTermSetPair}[thm]{\FormulaRefAuto{GoodBinaryAddressSetCriterion}[thm]{1,6}},7}},\FormulaRefAuto{A\cap B\subseteq A}[thm]}}
    \proofstepwidestarforcebreak[1,6]{%
      \begin{aligned}[t]
      &\forall a,b,c\in
        \TreeChildren{P}{\AddressEdges P}{\EmptyWord{\{0,1\}}}{p}\,\\[-2pt]
      &\quad(a=b\lor a=c\lor b=c)
      \end{aligned}}{%
      \FormulaRefAuto{ThreeElementsInPair}[thm]{8}}
    \proofstepwidestar[1]{%
      \forall p\in P\,\forall a,b,c\in
        \TreeChildren{P}{\AddressEdges P}{\EmptyWord{\{0,1\}}}{p}\,
        (a=b\lor a=c\lor b=c)}{%
      \rUI{\rRI{6,9}}}
    \proofstepwidestar[1]{%
      \BinaryTreeStruct{P,\AddressEdges P,\EmptyWord{\{0,1\}}}}{%
      \FormulaRefAuto{BinaryTreeDef}[def]{5,10}}
    \proofstepwidestarforcebreak[1,6]{%
      \begin{aligned}[t]
      &\TreeChildren{P}{\AddressEdges P}{\EmptyWord{\{0,1\}}}{p}
        =\varnothing\ \lor{}\\[-2pt]
      &\TreeChildren{P}{\AddressEdges P}{\EmptyWord{\{0,1\}}}{p}
        =\{\WordConcat p{\LetterWord0},
            \WordConcat p{\LetterWord1}\}
      \end{aligned}}{\rIE{\FormulaRefAuto{a=b\vdash b=a}[thm]{\rIE{\FormulaRefAuto{B27BinaryChildTermSetPair}[thm]{\FormulaRefAuto{GoodBinaryAddressSetCriterion}[thm]{1,6}},7}},\FormulaRefAuto{B27PairIntersectionFullness}[thm]{\FormulaRefAuto{GoodBinaryAddressSetFullnessAxiom}[ax]{1,6}}}}
    \proofstepwidestar[1,6]{p\in\BinaryAddresses}{%
      \FormulaRefAuto{GoodBinaryAddressSetCriterion}[thm]{1,6}}
    \proofstepwidestarforcebreak[1,6]{%
      \WordConcat p{\LetterWord0}\neq\WordConcat p{\LetterWord1}}{\FormulaRefAuto{BinaryAddressChildrenDistinct}[thm]{13}}
    \proofstepwidestarforcebreak[1,6]{%
      \begin{aligned}[t]
      &\TreeChildren{P}{\AddressEdges P}{\EmptyWord{\{0,1\}}}{p}
        =\varnothing\ \lor{}\\[-2pt]
      &\exists a,b\in
        \TreeChildren{P}{\AddressEdges P}{\EmptyWord{\{0,1\}}}{p}\,\\[-2pt]
      &\quad\Bigl(a\neq b\land{}\\[-2pt]
      &\quad\TreeChildren{P}{\AddressEdges P}{\EmptyWord{\{0,1\}}}{p}
          =\{a,b\}\Bigr)
      \end{aligned}}{\FormulaRefAuto{B27EmptyOrDistinctPairWitnesses}[thm]{12,14}}
    \proofstepwidestarbreak[1]{%
      \begin{aligned}[t]
      &\forall p\in P\,\Bigl(\\[-2pt]
      &\quad\TreeChildren{P}{\AddressEdges P}{\EmptyWord{\{0,1\}}}{p}
          =\varnothing\ \lor{}\\[-2pt]
      &\quad\exists a,b\in
        \TreeChildren{P}{\AddressEdges P}{\EmptyWord{\{0,1\}}}{p}\,\\[-2pt]
      &\qquad\bigl(a\neq b\land{}\\[-2pt]
      &\qquad\TreeChildren{P}{\AddressEdges P}{\EmptyWord{\{0,1\}}}{p}
          =\{a,b\}\bigr)\Bigr)
      \end{aligned}}{\rUI{\rRI{6,15}}}
    \proofstepwidestar[1]{%
      \FullBinaryTreeStruct{P,\AddressEdges P,\EmptyWord{\{0,1\}}}}{%
      \FormulaRefAuto{FullBinaryTreeDef}[def]{11,16}}
  \closeproofpartwideindR

    \end{BandTwentySevenProofCases}


  \proofpartwideindR[Die beiden Kinder eines inneren Adressknotens]{%
    \GoodAddressSet P\dsep p\in\TreeInnerVertices{P}{\AddressEdges P}{\EmptyWord{\{0,1\}}}\vdash\begin{aligned}[t]&\WordConcat p{\LetterWord0}\in P\land \WordConcat p{\LetterWord1}\in P\land{}\\[-2pt]&\TreeChildren{P}{\AddressEdges P}{\EmptyWord{\{0,1\}}}{p}=\{\WordConcat p{\LetterWord0},\WordConcat p{\LetterWord1}\}\end{aligned}
  }{%
    \GoodAddressSet P\dsep p\in\TreeInnerVertices{P}{\AddressEdges P}{\EmptyWord{\{0,1\}}}\vdash\begin{aligned}[t]&\WordConcat p{\LetterWord0}\in P\land \WordConcat p{\LetterWord1}\in P\land{}\\[-2pt]&\TreeChildren{P}{\AddressEdges P}{\EmptyWord{\{0,1\}}}{p}=\{\WordConcat p{\LetterWord0},\WordConcat p{\LetterWord1}\}\end{aligned}
  }[B27InnerAddressChildren]
    \proofstepwidestarbreak[1]{%
      \GoodAddressSet P}{%
      \rA}
    \proofstepwidestarbreak[2]{%
      p\in\TreeInnerVertices{P}{\AddressEdges P}{\EmptyWord{\{0,1\}}}}{%
      \rA}
    \proofstepwidestarbreak[1]{%
      \RootedTreeStruct{P,\AddressEdges P,\EmptyWord{\{0,1\}}}}{%
      \FormulaRefAuto{BinaryTreeRootedAxiom}[ax]{\FormulaRefAuto{FullBinaryTreeBinaryAxiom}[ax]{\FormulaRefAuto{GoodBinaryAddressGraphFullBinary}[thm]{1}}}}
    \proofstepwidestarbreak[1,2]{%
      p\in P\land \TreeChildren{P}{\AddressEdges P}{\EmptyWord{\{0,1\}}}{p}\neq\varnothing}{%
      \FormulaRefAuto{TreeLeafInnerDef}[def]{\FormulaRefAuto{TreeInnerVertexSetDef}[def]{3,2}}}
    \proofstepwidestarbreak[1,2]{%
      p\in\BinaryAddresses}{%
      \FormulaRefAuto{GoodBinaryAddressSetCriterion}[thm]{1,\rAEa{4}}}
    \proofstepwidestarbreak[1,2]{%
      \{\WordConcat p{\LetterWord i}\mid i\in\{0,1\}\}=\{\WordConcat p{\LetterWord0},\WordConcat p{\LetterWord1}\}}{%
      \FormulaRefAuto{B27BinaryChildTermSetPair}[thm]{5}}
    \proofstepwidestarbreak[1,2]{%
      \begin{aligned}[t]
        &\TreeChildren{P}{\AddressEdges P}{\EmptyWord{\{0,1\}}}{p}\\[-2pt]
        &\quad=\{\WordConcat p{\LetterWord i}\mid i\in\{0,1\}\}\cap P
      \end{aligned}}{%
      \FormulaRefAuto{GoodBinaryAddressChildren}[thm]{1,\rAEa{4}}}
    \proofstepwidestarbreak[1,2]{%
      \TreeChildren{P}{\AddressEdges P}{\EmptyWord{\{0,1\}}}{p}=\{\WordConcat p{\LetterWord0},\WordConcat p{\LetterWord1}\}\cap P}{%
      \rIE{6,7}}
    \proofstepwidestarbreak[1,2]{%
      \WordConcat p{\LetterWord0}\in P\leftrightarrow \WordConcat p{\LetterWord1}\in P}{%
      \FormulaRefAuto{GoodBinaryAddressSetFullnessAxiom}[ax]{1,\rAEa{4}}}
    \proofstepwidestarbreak[1,2]{%
      \{\WordConcat p{\LetterWord0},\WordConcat p{\LetterWord1}\}\cap P=\varnothing\lor \{\WordConcat p{\LetterWord0},\WordConcat p{\LetterWord1}\}\cap P=\{\WordConcat p{\LetterWord0},\WordConcat p{\LetterWord1}\}}{%
      \FormulaRefAuto{B27PairIntersectionFullness}[thm]{9}}
    \proofstepwidestarbreak[1,2]{%
      \TreeChildren{P}{\AddressEdges P}{\EmptyWord{\{0,1\}}}{p}=\varnothing\lor \TreeChildren{P}{\AddressEdges P}{\EmptyWord{\{0,1\}}}{p}=\{\WordConcat p{\LetterWord0},\WordConcat p{\LetterWord1}\}}{%
      \rIE{\FormulaRefAuto{a=b\vdash b=a}[thm]{8},10}}
    \proofstepwidestarbreak[1,2]{%
      \TreeChildren{P}{\AddressEdges P}{\EmptyWord{\{0,1\}}}{p}=\{\WordConcat p{\LetterWord0},\WordConcat p{\LetterWord1}\}}{%
      \FormulaRefAuto{P\lor Q,\neg P\vdash Q}[thm]{11,\rAEb{4}}}
    \proofstepwidestarbreak[1,2]{%
      \WordConcat p{\LetterWord0}\in\TreeChildren{P}{\AddressEdges P}{\EmptyWord{\{0,1\}}}{p}}{%
      \rIE{\FormulaRefAuto{a=b\vdash b=a}[thm]{12},\FormulaRefAuto{a\in\{a,b\}}[thm]}}
    \proofstepwidestarbreak[1,2]{%
      \WordConcat p{\LetterWord1}\in\TreeChildren{P}{\AddressEdges P}{\EmptyWord{\{0,1\}}}{p}}{%
      \rIE{\FormulaRefAuto{a=b\vdash b=a}[thm]{12},\FormulaRefAuto{b\in\{a,b\}}[thm]}}
    \proofstepwidestarbreak[1,2]{%
      \TreeChildren{P}{\AddressEdges P}{\EmptyWord{\{0,1\}}}{p}\subseteq P}{%
      \rIE{\FormulaRefAuto{a=b\vdash b=a}[thm]{8},\FormulaRefAuto{A\cap B\subseteq B}[thm]}}
    \proofstepwidestarbreak[1,2]{%
      \WordConcat p{\LetterWord0}\in P}{%
      \FormulaRefAuto{A\subseteq B\dsep x\in A\vdash x\in B}[thm]{15,13}}
    \proofstepwidestarbreak[1,2]{%
      \WordConcat p{\LetterWord1}\in P}{%
      \FormulaRefAuto{A\subseteq B\dsep x\in A\vdash x\in B}[thm]{15,14}}
    \proofstepwidestarbreak[1,2]{%
      \begin{aligned}[t]&\WordConcat p{\LetterWord0}\in P\land \WordConcat p{\LetterWord1}\in P\land{}\\[-2pt]&\TreeChildren{P}{\AddressEdges P}{\EmptyWord{\{0,1\}}}{p}=\{\WordConcat p{\LetterWord0},\WordConcat p{\LetterWord1}\}\end{aligned}}{%
      \rAI{16,\rAI{17,12}}}
  \closeproofpartwideindR

  \proofpartwideindR[Elementkriterium der inneren Adressknoten]{%
    \GoodAddressSet P\vdash p\in\TreeInnerVertices{P}{\AddressEdges P}{\EmptyWord{\{0,1\}}}\leftrightarrow p\in P\land \WordConcat p{\LetterWord0}\in P
  }{%
    \GoodAddressSet P\vdash p\in\TreeInnerVertices{P}{\AddressEdges P}{\EmptyWord{\{0,1\}}}\leftrightarrow p\in P\land \WordConcat p{\LetterWord0}\in P
  }[B27InnerAddressMembership]
    \proofstepwidestarbreak[1]{%
      \GoodAddressSet P}{%
      \rA}
    \proofstepwidestarbreak[1]{%
      \RootedTreeStruct{P,\AddressEdges P,\EmptyWord{\{0,1\}}}}{%
      \FormulaRefAuto{BinaryTreeRootedAxiom}[ax]{\FormulaRefAuto{FullBinaryTreeBinaryAxiom}[ax]{\FormulaRefAuto{GoodBinaryAddressGraphFullBinary}[thm]{1}}}}
    \proofstepwidestarbreak[3]{%
      p\in\TreeInnerVertices{P}{\AddressEdges P}{\EmptyWord{\{0,1\}}}}{%
      \rA}
    \proofstepwidestarbreak[1,3]{%
      p\in P}{%
      \FormulaRefAuto{TreeInnerVertexSetDef}[def]{2,3}}
    \proofstepwidestarbreak[1,3]{%
      \WordConcat p{\LetterWord0}\in P}{%
      \rAEa{\FormulaRefAuto{B27InnerAddressChildren}[thm]{1,3}}}
    \proofstepwidestarbreak[1,3]{%
      p\in P\land \WordConcat p{\LetterWord0}\in P}{%
      \rAI{4,5}}
    \proofstepwidestarbreak[7]{%
      p\in P\land \WordConcat p{\LetterWord0}\in P}{%
      \rA}
    \proofstepwidestarbreak[1,7]{%
      p\in\BinaryAddresses}{%
      \FormulaRefAuto{GoodBinaryAddressSetCriterion}[thm]{1,\rAEa{7}}}
    \proofstepwidestarbreak[1,7]{%
      \{\WordConcat p{\LetterWord i}\mid i\in\{0,1\}\}=\{\WordConcat p{\LetterWord0},\WordConcat p{\LetterWord1}\}}{%
      \FormulaRefAuto{B27BinaryChildTermSetPair}[thm]{8}}
    \proofstepwidestarbreak[1,7]{%
      \TreeChildren{P}{\AddressEdges P}{\EmptyWord{\{0,1\}}}{p}=\{\WordConcat p{\LetterWord0},\WordConcat p{\LetterWord1}\}\cap P}{%
      \rIE{9,\FormulaRefAuto{GoodBinaryAddressChildren}[thm]{1,\rAEa{7}}}}
    \proofstepwidestarbreak[]{%
      \WordConcat p{\LetterWord0}\in\{\WordConcat p{\LetterWord0},\WordConcat p{\LetterWord1}\}}{%
      \FormulaRefAuto{a\in\{a,b\}}[thm]}
    \proofstepwidestarbreak[7]{%
      \WordConcat p{\LetterWord0}\in\{\WordConcat p{\LetterWord0},\WordConcat p{\LetterWord1}\}\cap P}{%
      \FormulaRefAuto{x\in A\dsep x\in B\vdash x\in A\cap B}[thm]{11,\rAEb{7}}}
    \proofstepwidestarbreak[1,7]{%
      \WordConcat p{\LetterWord0}\in\TreeChildren{P}{\AddressEdges P}{\EmptyWord{\{0,1\}}}{p}}{%
      \rIE{\FormulaRefAuto{a=b\vdash b=a}[thm]{10},12}}
    \proofstepwidestarbreak[1,7]{%
      \TreeChildren{P}{\AddressEdges P}{\EmptyWord{\{0,1\}}}{p}\neq\varnothing}{%
      \FormulaRefAuto{a\in A\vdash A\neq\varnothing}[thm]{13}}
    \proofstepwidestarbreak[1,7]{%
      \TreeInnerPred{P}{\AddressEdges P}{\EmptyWord{\{0,1\}}}{p}}{%
      \FormulaRefAuto{TreeLeafInnerDef}[def]{2,\rAI{\rAEa{7},14}}}
    \proofstepwidestarbreak[1,7]{%
      p\in\TreeInnerVertices{P}{\AddressEdges P}{\EmptyWord{\{0,1\}}}}{%
      \FormulaRefAuto{TreeInnerVertexSetDef}[def]{2,\rAEa{7},15}}
    \proofstepwidestarbreak[1]{%
      p\in\TreeInnerVertices{P}{\AddressEdges P}{\EmptyWord{\{0,1\}}}\leftrightarrow p\in P\land \WordConcat p{\LetterWord0}\in P}{%
      \rLRI{\rRI{3,6},\rRI{7,16}}}
  \closeproofpartwideindR

  \proofpartwideindR[Geordnete volle Binärbaumstruktur]{%
    \GoodAddressSet P\vdash \OrderedFullBinaryTreeStruct{P,\AddressEdges P,\EmptyWord{\{0,1\}},\AddressLeft P,\AddressRight P}
  }{%
    \GoodAddressSet P\vdash \OrderedFullBinaryTreeStruct{P,\AddressEdges P,\EmptyWord{\{0,1\}},\AddressLeft P,\AddressRight P}
  }[GoodBinaryAddressSetOrderedFullBinaryTree]
    \proofstepwidestarbreak[1]{%
      \GoodAddressSet P}{%
      \rA}
    \proofstepwidestarbreak[1]{%
      \FullBinaryTreeStruct{P,\AddressEdges P,\EmptyWord{\{0,1\}}}}{%
      \FormulaRefAuto{GoodBinaryAddressGraphFullBinary}[thm]{1}}
    \proofstepwidestarbreak[1]{%
      \AddressInner P=\TreeInnerVertices{P}{\AddressEdges P}{\EmptyWord{\{0,1\}}}}{%
      \FormulaRefAuto{\forall x\,(x\in A\leftrightarrow x\in B)\vdash A=B}[ax]{\rUI{\FormulaRefAuto{P\leftrightarrow Q,R\leftrightarrow Q\vdash P\leftrightarrow R}[thm]{\FormulaRefAuto{BinaryAddressGraphDef}[def]{1},\FormulaRefAuto{B27InnerAddressMembership}[thm]{1}}}}}
    \proofstepwidestarbreak[4]{%
      p\in\TreeInnerVertices{P}{\AddressEdges P}{\EmptyWord{\{0,1\}}}}{%
      \rA}
    \proofstepwidestarbreak[1,4]{%
      \begin{aligned}[t]&\WordConcat p{\LetterWord0}\in P\land \WordConcat p{\LetterWord1}\in P\land{}\\[-2pt]&\TreeChildren{P}{\AddressEdges P}{\EmptyWord{\{0,1\}}}{p}=\{\WordConcat p{\LetterWord0},\WordConcat p{\LetterWord1}\}\end{aligned}}{%
      \FormulaRefAuto{B27InnerAddressChildren}[thm]{1,4}}
    \proofstepwidestarbreak[1,4]{%
      \WordConcat p{\LetterWord0}\in P\land \WordConcat p{\LetterWord1}\in P}{%
      \rAI{\rAEa{5},\rAEa{\rAEb{5}}}}
    \proofstepwidestarbreak[1]{%
      \forall p\in\TreeInnerVertices{P}{\AddressEdges P}{\EmptyWord{\{0,1\}}}\,(\WordConcat p{\LetterWord0}\in P\land \WordConcat p{\LetterWord1}\in P)}{%
      \rUI{\rRI{4,6}}}
    \proofstepwidestarbreak[1]{%
      \forall p\in\AddressInner P\,(\WordConcat p{\LetterWord0}\in P\land \WordConcat p{\LetterWord1}\in P)}{%
      \rIE{\FormulaRefAuto{a=b\vdash b=a}[thm]{3},7}}
    \proofstepwidestarbreak[9]{%
      p\in\AddressInner P}{%
      \rA}
    \proofstepwidestarbreak[1,9]{%
      \WordConcat p{\LetterWord0}\in P\land \WordConcat p{\LetterWord1}\in P}{%
      \rUE{8,9}}
    \proofstepwidestarbreak[1]{%
      \forall p\in\AddressInner P\,(\WordConcat p{\LetterWord0}\in P)}{%
      \rUI{\rRI{9,\rAEa{10}}}}
    \proofstepwidestarbreak[1]{%
      \forall p\in\AddressInner P\,(\WordConcat p{\LetterWord1}\in P)}{%
      \rUI{\rRI{9,\rAEb{10}}}}
    \proofstepwidestarbreak[1]{%
      \AddressLeft P\colon\AddressInner P\to P}{%
      \FormulaRefAuto{BinaryAddressGraphDef}[def]{1,\FormulaRefAuto{TypedFunctionDefinitionPrinciple}[thm]{11}}}
    \proofstepwidestarbreak[1]{%
      \AddressRight P\colon\AddressInner P\to P}{%
      \FormulaRefAuto{BinaryAddressGraphDef}[def]{1,\FormulaRefAuto{TypedFunctionDefinitionPrinciple}[thm]{12}}}
    \proofstepwidestarbreak[1]{%
      \AddressLeft P,\AddressRight P\colon\AddressInner P\to P}{%
      \rAI{13,14}}
    \proofstepwidestarbreak[1]{%
      \AddressLeft P,\AddressRight P\colon\TreeInnerVertices{P}{\AddressEdges P}{\EmptyWord{\{0,1\}}}\to P}{%
      \rIE{3,15}}
    \proofstepwidestarbreak[1,4]{%
      p\in P}{%
      \FormulaRefAuto{TreeInnerVertexSetDef}[def]{\FormulaRefAuto{BinaryTreeRootedAxiom}[ax]{\FormulaRefAuto{FullBinaryTreeBinaryAxiom}[ax]{\FormulaRefAuto{GoodBinaryAddressGraphFullBinary}[thm]{1}}},4}}
    \proofstepwidestarbreak[1,4]{%
      p\in\BinaryAddresses}{%
      \FormulaRefAuto{GoodBinaryAddressSetCriterion}[thm]{1,17}}
    \proofstepwidestarbreak[1,4]{%
      \WordConcat p{\LetterWord0}\neq \WordConcat p{\LetterWord1}}{%
      \FormulaRefAuto{BinaryAddressChildrenDistinct}[thm]{18}}
    \proofstepwidestarbreak[1,4]{%
      \AddressLeft P(p)=\WordConcat p{\LetterWord0}\land\AddressRight P(p)=\WordConcat p{\LetterWord1}}{%
      \FormulaRefAuto{BinaryAddressGraphDef}[def]{1,\rIE{\FormulaRefAuto{a=b\vdash b=a}[thm]{3},4}}}
    \proofstepwidestarbreak[1,4]{%
      \begin{aligned}[t]&\AddressLeft P(p)\neq\AddressRight P(p)\land{}\\[-2pt]&\TreeChildren{P}{\AddressEdges P}{\EmptyWord{\{0,1\}}}{p}=\{\AddressLeft P(p),\AddressRight P(p)\}\end{aligned}}{%
      \rIE{\FormulaRefAuto{a=b\vdash b=a}[thm]{\rAEa{20}},\FormulaRefAuto{a=b\vdash b=a}[thm]{\rAEb{20}},\rAI{19,\rAEb{\rAEb{5}}}}}
    \proofstepwidestarbreak[1]{%
      \begin{aligned}[t]&\forall p\in\TreeInnerVertices{P}{\AddressEdges P}{\EmptyWord{\{0,1\}}}\,\Bigl(\AddressLeft P(p)\neq\AddressRight P(p)\land{}\\[-2pt]&\quad\TreeChildren{P}{\AddressEdges P}{\EmptyWord{\{0,1\}}}{p}=\{\AddressLeft P(p),\AddressRight P(p)\}\Bigr)\end{aligned}}{%
      \rUI{\rRI{4,21}}}
    \proofstepwidestarbreak[1]{%
      \OrderedFullBinaryTreeStruct{P,\AddressEdges P,\EmptyWord{\{0,1\}},\AddressLeft P,\AddressRight P}}{%
      \FormulaRefAuto{OrderedFullBinaryTreeDef}[def]{2,16,22}}
  \closeproofpartwideindR

\end{tabproofsplitwide}

\enlargethispage{3\baselineskip}

\FormulaDefDeltaK[Positionsgraph eines Klammerungsbaumcodes]{%
  R\in\BracketTreeSet A\vdash
  \TreePositionGraph A R\coloneqq
  \bigl(\TreePositions A R,\AddressEdges{\TreePositions A R},
    \EmptyWord{\{0,1\}},\AddressLeft{\TreePositions A R},
    \AddressRight{\TreePositions A R}\bigr)%
}{TreePositionGraphDef}{%
  \DeltaRow{Mengen}{A\dsep R\dsep\TreePositions A R}
  \DeltaRow{\textbf{Neues Symbol}}{}
  \DeltaPrem{Positionsgraph}{\TreePositionGraph A R}
}

\vspace{-\baselineskip}
\FormulaThmDeltaK[Klammerungsbaumcodes tragen endliche geordnete volle
Binärbäume]{%
  \begin{aligned}[t]
    &R\in\BracketTreeSet A\vdash{}\\[-2pt]
    &\text{(i)}\;\Finite{\TreePositions A R}\\[-2pt]
    &\text{(ii)}\;\mathsf{GVollBinBaum}\!\bigl(
      \TreePositions A R,\AddressEdges{\TreePositions A R},\\[-2pt]
    &\qquad\EmptyWord{\{0,1\}},\AddressLeft{\TreePositions A R},
      \AddressRight{\TreePositions A R}\bigr)\\[-2pt]
    &\text{(iii)}\;\TreeStruct{\TreePositions A R,
      \AddressEdges{\TreePositions A R}}
  \end{aligned}%
}{BracketCodePositionGraphIsFiniteOrderedFullBinaryTree}{%
  \DeltaRow{Mengen}{A\dsep R\dsep\TreePositions A R}
}
\begin{tabproofsplitwide}
  \proofpartwideindR[Endlichkeit]{%
    R\in\BracketTreeSet A\vdash \Finite{\TreePositions A R}
  }{R\in\BracketTreeSet A\vdash \Finite{\TreePositions A R}}[BracketCodePositionSetFinite]
    \proofstepwidestar[1]{%
      R\in\BracketTreeSet A}{%
      \rA}
    \proofstepwidestar[1]{%
      \GoodAddressSet{\TreePositions A R}}{%
      \FormulaRefAuto{TreePositionSetIsGood}[thm]{1}}
    \proofstepwidestar[1]{%
      \Finite{\TreePositions A R}}{%
      \FormulaRefAuto{GoodBinaryAddressSetFinite}[thm]{2}}
  \closeproofpartwideindR

  \proofpartwideindR[Geordnete volle Binärbaumstruktur]{%
    \begin{aligned}[t]
      &R\in\BracketTreeSet A\vdash{}\\[-2pt]
      &\OrderedFullBinaryTreeStruct{
        \TreePositions A R,\AddressEdges{\TreePositions A R},
        \EmptyWord{\{0,1\}},\AddressLeft{\TreePositions A R},
        \AddressRight{\TreePositions A R}}
    \end{aligned}
  }{\begin{aligned}[t]
      &R\in\BracketTreeSet A\vdash{}\\[-2pt]
      &\OrderedFullBinaryTreeStruct{
        \TreePositions A R,\AddressEdges{\TreePositions A R},
        \EmptyWord{\{0,1\}},\AddressLeft{\TreePositions A R},
        \AddressRight{\TreePositions A R}}
    \end{aligned}}[BracketCodePositionGraphOrderedFullBinaryTree]
    \proofstepwidestar[1]{%
      R\in\BracketTreeSet A}{%
      \rA}
    \proofstepwidestar[1]{%
      \GoodAddressSet{\TreePositions A R}}{%
      \FormulaRefAuto{TreePositionSetIsGood}[thm]{1}}
    \proofstepwidestarbreak[1]{%
      \begin{aligned}[t]
      &\mathsf{GVollBinBaum}\!\bigl(\TreePositions A R,\\[-2pt]
      &\quad\AddressEdges{\TreePositions A R},\EmptyWord{\{0,1\}},\\[-2pt]
      &\quad\AddressLeft{\TreePositions A R},\AddressRight{\TreePositions A R}\bigr)
      \end{aligned}}{%
      \FormulaRefAuto{GoodBinaryAddressSetOrderedFullBinaryTree}[thm]{2}}
  \closeproofpartwideindR

  \proofpartwideindR[Zugrunde liegender Baum]{%
    R\in\BracketTreeSet A\vdash \TreeStruct{\TreePositions A R,\AddressEdges{\TreePositions A R}}
  }{R\in\BracketTreeSet A\vdash \TreeStruct{\TreePositions A R,\AddressEdges{\TreePositions A R}}}[BracketCodePositionGraphTree]
    \proofstepwidestar[1]{%
      R\in\BracketTreeSet A}{%
      \rA}
    \proofstepwidestarbreak[1]{%
      \begin{aligned}[t]
      &\mathsf{GVollBinBaum}\!\bigl(\TreePositions A R,\\[-2pt]
      &\quad\AddressEdges{\TreePositions A R},\EmptyWord{\{0,1\}},\\[-2pt]
      &\quad\AddressLeft{\TreePositions A R},\AddressRight{\TreePositions A R}\bigr)
      \end{aligned}}{%
      \FormulaRefAuto{BracketCodePositionGraphOrderedFullBinaryTree}[thm]{1}}
    \proofstepwidestar[1]{%
      \begin{aligned}[t]
        &\FullBinaryTreeStruct{\TreePositions A R,
          \AddressEdges{\TreePositions A R},\EmptyWord{\{0,1\}}}\land{}\\[-2pt]
        &\BinaryTreeStruct{\TreePositions A R,
          \AddressEdges{\TreePositions A R},\EmptyWord{\{0,1\}}}\land{}\\[-2pt]
        &\RootedTreeStruct{\TreePositions A R,
          \AddressEdges{\TreePositions A R},\EmptyWord{\{0,1\}}}
      \end{aligned}}{%
      \FormulaRefAuto{OrderedFullBinaryTreeForgetfulChain}[thm]{2}}
    \proofstepwidestar[1]{%
      \RootedTreeStruct{\TreePositions A R,
        \AddressEdges{\TreePositions A R},\EmptyWord{\{0,1\}}}}{%
      \rAEn{3}}
    \proofstepwidestar[1]{\TreeStruct{\TreePositions A R,\AddressEdges{\TreePositions A R}}\land\EmptyWord{\{0,1\}}\in\TreePositions A R}{%
      \FormulaRefAuto{RootedTreeAccess}[thm]{4}}
    \proofstepwidestar[1]{%
      \TreeStruct{\TreePositions A R,\AddressEdges{\TreePositions A R}}}{%
      \rAEa{5}}
  \closeproofpartwideindR
\end{tabproofsplitwide}


% Anwendungen des allgemeinen Elternsystem-Kriteriums aus Band 26.
\providecommand{\BXXVIIFiniteRooted}[3]{%
  \mathsf{EndWBaum}(#1,#2,#3)}
\providecommand{\BXXVIIParentSystem}[4]{%
  \mathsf{ElterSys}(#1,#2,#3,#4)}
\providecommand{\BXXVIISymmetricEdges}[2]{%
  \operatorname{sym}_{#1}(#2)}
\providecommand{\BXXVIIForwardEdges}[1]{%
  \operatorname{vor}(#1)}


\section{Die gerichtete Darstellung der Positionsbäume}

Beim Positionsbaum ist die Höhenmarke bereits vorhanden: Es ist die
Wortlänge der Adresse. Das Anhängen eines Bits ist eine gerichtete
Elternkante; erst durch das Vergessen ihrer Richtung erhält man die
bisher verwendete symmetrische Adressrelation.

\noindent\begin{minipage}{\linewidth}
\FormulaDefDeltaK[Vorwärtskanten einer Positionsmenge]{%
  \BXXVIIForwardEdges P\coloneqq
  \{(p,q)\in P\times P\mid
      \exists i\in\{0,1\}\,q=\WordConcat p{\LetterWord i}\}
}{B27ForwardAddressEdgesDef}{%
  \DeltaRow{Mengen}{P\dsep p\dsep q\dsep i}
  \DeltaPrem{Positionsmenge}{\GoodAddressSet P}
  \DeltaRow{Neue Menge}{\BXXVIIForwardEdges P}
}
\end{minipage}\par

\noindent\begin{minipage}{\linewidth}
\FormulaThmDeltaK[Gerichtete und ungerichtete Positionskanten]{%
  \begin{aligned}[t]
    \GoodAddressSet P\vdash{}&
      \BXXVIISymmetricEdges P{\BXXVIIForwardEdges P}=\AddressEdges P
  \end{aligned}
}{B27ForwardAddressSymmetrization}{%
  \DeltaRow{Mengen}{P\dsep p\dsep q\dsep i}
}
\end{minipage}\par
\noindent\textit{Beweis.}
Für \(p,q\in P\) bedeutet die Mitgliedschaft links, dass entweder
\(q=p\frown\LetterWord i\) für ein Bit \(i\) gilt oder
\(p=q\frown\LetterWord i\) für ein Bit \(i\) gilt.
Durch Fallunterscheidung ist dies gleichbedeutend mit der Existenz
eines Bits, für das eine dieser beiden Gleichungen gilt. Nach
\FormulaRefAuto{B27AddressEdgeCriterion}[thm] ist dies genau die
Mitgliedschaft in \(\AddressEdges P\). Außerhalb von \(P\times P\)
enthält keine der beiden Relationen ein Paar. Die Mengenextensionalität
liefert die Gleichheit.

\noindent\begin{minipage}{\linewidth}
\FormulaThmDeltaK[Positionsmengen erfüllen das allgemeine Elternsystem]{%
  \begin{aligned}[t]
    &\GoodAddressSet P\dsep h\colon P\to\mathbb N\dsep
      \forall p\in P\,(h(p)=\WordLength p)\vdash\\[-2pt]
    &\BXXVIIParentSystem P{\BXXVIIForwardEdges P}
      {\EmptyWord{\{0,1\}}}h
  \end{aligned}
}{B27AddressRankedParentSystem}{%
  \DeltaRow{Mengen}{P\dsep h\dsep p\dsep q\dsep i}
}
\end{minipage}\par
\noindent\textit{Beweis.}
\FormulaRefAuto{GoodBinaryAddressSetCriterion}[thm] liefert die
Endlichkeit von \(P\), die Wurzel und die Existenz eines unmittelbaren
Präfixes jeder Nichtwurzeladresse. Die Vorwärtskanten liegen nach ihrer
Definition in \(P\times P\). Aus
\FormulaRefAuto{EmptyWordLengthZero}[thm] folgt die Höhenmarke null
an der Wurzel. Jeder Vorgänger einer Nichtwurzeladresse \(q\) ist
\(\WordRumpf{\{0,1\}}q\), denn
\FormulaRefAuto{WordRumpfAppend}[thm] entfernt das angehängte Bit.
Der Vorgänger ist somit eindeutig. Schließlich liefert
\FormulaRefAuto{WordAppendLetterLength}[thm] für jede Vorwärtskante
\(h(q)=h(p)+1\), also das verlangte strenge Wachstum. Dies sind
sämtliche Bedingungen von
\FormulaRefAuto{B27RankedParentSystemDef}[def].

Eine solche Funktion \(h\) existiert: Jedem \(p\in P\) wird die
eindeutig bestimmte natürliche Zahl \(\WordLength p\) zugeordnet;
\FormulaRefAuto{B27AddressWordTyping}[thm] und
\FormulaRefAuto{FiniteWordLengthNatural}[thm] rechtfertigen die
Typisierung. Das Funktionsdefinitionsprinzip bildet daraus ihren Graphen.
Damit ist das allgemeine Einführungskriterium auf jede der zuvor
konstruierten Positionsmengen anwendbar.

Die volle binäre Verzweigung ist eine zusätzliche Bedingung. Im
allgemeinen Elternsystem sind beliebig viele Kinder möglich. Bei
\(\GoodAddressSet P\) kommen als Kinder einer Adresse nur ihre beiden
Erweiterungen um \(0\) und \(1\) infrage. Der Vollheitsabschluss
fordert, dass beide vorhanden sind oder beide fehlen. Die Funktionen
\(\AddressLeft P\) und \(\AddressRight P\) unterscheiden diese Kinder.
Genau dies hat
\FormulaRefAuto{GoodBinaryAddressSetOrderedFullBinaryTree}[thm] als
Struktur aus Band~26 nachgewiesen. Eine spätere Syntax darf daher
den allgemeinen endlichen Wurzelbaum verwenden und ihre jeweiligen
Konstruktoren durch zusätzliche Beschriftungs- und Verzweigungsregeln
kennzeichnen.


\chapter{Blattwörter, Klammerungen und Auswertung}

Die folgenden Konstruktionen benutzen die zuvor bewiesenen
Induktions- und Rekursionssätze als gemeinsame Grundlage. Das Blattwort
ersetzt ein Blatt durch seinen Buchstaben und einen Knoten durch
Konkatenation. Der Linksbaum übersetzt das Anfügen eines Buchstabens
in das Anfügen eines Blattes. Eine Auswertung ersetzt schließlich den
Knotenkonstruktor durch eine vorgegebene binäre Operation.
Damit haben die verschiedenen Beweise denselben Aufbau: Anfangswerte,
Fortsetzung und ein Induktions- oder Rekursionsschluss.

\section{Blattwort}

\FormulaDefDeltaK[Einbuchstabenabbildung]{%
  \begin{aligned}[t]
    &\operatorname{ein}_A\colon A\to\NonemptyWords{A},\\[-2pt]
    &\forall a\in A\,
      \bigl(\operatorname{ein}_A(a)\coloneqq\LetterWord{a}\bigr)
  \end{aligned}%
}{LetterWordFunctionDef}{%
  \DeltaRow{Mengen}{A\dsep a}
  \DeltaRow{Funktionen}{\operatorname{ein}_A}
  \DeltaRow{\textbf{Neues Symbol}}{}
  \DeltaPrem{Einbuchstabenabbildung}{\operatorname{ein}_A}
}
\begin{tabproofwide}
  \proofstepwidestar[1]{a\in A}{\rA}
  \proofstepwidestar[1]{\LetterWord{a}\in\NonemptyWords{A}}{%
    \FormulaRefAuto{LetterWordNonempty}[thm]{1}}
  \proofstepwidestar[]{%
    \forall a\in A\,
      \bigl(\LetterWord{a}\in\NonemptyWords{A}\bigr)}{%
    \rUI{\rRI{1,2}}}
  \proofstepwidestar[]{%
    \exists!F\,\Bigl(
      F\colon A\to\NonemptyWords{A}\land
      \forall a\in A\,F(a)=\LetterWord{a}
    \Bigr)}{%
    \FormulaRefAuto{TypedFunctionDefinitionPrinciple}[thm]{3}}
\end{tabproofwide}

\FormulaThmDeltaK[Die Konkatenation ist ein binärer Operator auf
nichtleeren Wörtern]{%
  \mathsf{BinOp}(\NonemptyWords{A},\frown)%
}{NonemptyWordConcatenationBinaryOperation}{%
  \DeltaRow{Mengen}{A\dsep u\dsep v}
  \DeltaRow{Funktionen}{\frown}
}
\begin{tabproofwide}
  \proofstepwidestar[1]{u,v\in\NonemptyWords{A}}{\rA}
  \proofstepwidestar[1]{%
    \WordConcat{u}{v}\in\NonemptyWords{A}}{%
    \FormulaRefAuto{NonemptyWordConcatenationClosure}[thm]{1}}
  \proofstepwidestar[]{%
    \mathsf{BinOp}(\NonemptyWords{A},\frown)}{%
    \FormulaRefAuto{BinaryOperationDef}[def]{2}}
\end{tabproofwide}

\FormulaThmDeltaK[Existenz und Eindeutigkeit der Blattwortabbildung]{%
  \begin{aligned}[t]
    &\exists!F\colon
      \BracketTreeSet{A}\to\NonemptyWords{A}\,\Bigl(\\[-2pt]
    &\quad \forall a\in A\,
      F(\LeafTree{A}{a})=\operatorname{ein}_A(a)\ \land{}\\[-2pt]
    &\quad \forall S,T\in\BracketTreeSet{A}\,
      F(\NodeTree{A}{S}{T})=
        \BinaryOpFunction{\NonemptyWords{A}}{\frown}(F(S),F(T))
      \Bigr)
  \end{aligned}%
}{TreeLeafWordExistenceUnique}{%
  \DeltaRow{Mengen}{A\dsep a\dsep S\dsep T}
  \DeltaRow{Funktionen}{F\dsep\operatorname{ein}_A\dsep\frown}
}
\begin{tabproofwide}
  \proofstepwidestar[]{%
    \operatorname{ein}_A\colon A\to\NonemptyWords{A}}{%
    \FormulaRefAuto{LetterWordFunctionDef}[def]}
  \proofstepwidestar[]{%
    \mathsf{BinOp}(\NonemptyWords{A},\frown)}{%
    \FormulaRefAuto{NonemptyWordConcatenationBinaryOperation}[thm]}
  \proofstepwidestar[]{%
    \BinaryOpFunction{\NonemptyWords{A}}{\frown}\colon
      \NonemptyWords{A}\times\NonemptyWords{A}
      \to\NonemptyWords{A}}{%
    \FormulaRefAuto{BinaryOperationFunctionDef}[def]{2}}
  \proofstepwidestar[]{%
    \begin{aligned}[t]
      &\exists!F\colon\BracketTreeSet{A}\\[-2pt]
      &\qquad\to\NonemptyWords{A}\,\Bigl(\\[-2pt]
      &\quad \forall a\in A\,
        F(\LeafTree{A}{a})=\operatorname{ein}_A(a)\ \land{}\\[-2pt]
      &\quad \forall S,T\in\BracketTreeSet{A}\,
        F(\NodeTree{A}{S}{T})\\[-2pt]
      &\qquad=\BinaryOpFunction{\NonemptyWords{A}}{\frown}\\[-2pt]
      &\qquad\quad(F(S),F(T))
        \Bigr)
    \end{aligned}}{%
    \FormulaRefAuto{FullPlanarBinaryTreeRecursion}[thm]{1,3}}
\end{tabproofwide}

\FormulaDefDeltaK[Blattwort eines Klammerungsbaums]{%
  \begin{aligned}[t]
    &\operatorname{wort}_A\coloneqq{}\\[-2pt]
    &\qquad \TreeRecursor{A}{\NonemptyWords{A}}{\operatorname{ein}_A}{\BinaryOpFunction{\NonemptyWords{A}}{\frown}}
  \end{aligned}%
}{TreeLeafWordDef}{%
  \DeltaRow{Mengen}{A\dsep a\dsep S\dsep T}
  \DeltaRow{Funktionen}{\operatorname{wort}_A}
  \DeltaRow{\textbf{Neues Symbol}}{}
  \DeltaPrem{Blattwort}{\operatorname{wort}_A}
}

\FormulaThmDeltaK[Typisierung und Rekursionsgleichungen des Blattwortes]{%
  \begin{aligned}[t]
    &\text{(i)}\;\operatorname{wort}_A\colon
      \BracketTreeSet{A}\to\NonemptyWords{A}\\[-2pt]
    &\text{(ii)}\;a\in A\vdash
      \TreeLeafWord{A}{\LeafTree{A}{a}}=\LetterWord{a}\\[-2pt]
    &\text{(iii)}\;S,T\in\BracketTreeSet{A}\vdash
      \TreeLeafWord{A}{\NodeTree{A}{S}{T}}
      =\WordConcat{\TreeLeafWord{A}{S}}{\TreeLeafWord{A}{T}}
  \end{aligned}%
}{TreeLeafWordEquations}{%
  \DeltaRow{Mengen}{A\dsep a\dsep S\dsep T}
  \DeltaRow{Funktionen}{\operatorname{wort}_A\dsep R}
}
Im folgenden Beweis bezeichnet \(R\) die rechte Seite von
\FormulaRefAuto{TreeLeafWordDef}[def]. Die Blattregel ersetzt die
Beschriftung; die Knotenregel verknüpft die beiden Teilbaumergebnisse.
\begin{tabproofsplitwide}
  \proofpartwideindR[Typisierung des Blattwortes]{%
    \operatorname{wort}_A\colon
      \BracketTreeSet{A}\to\NonemptyWords{A}
  }{\operatorname{wort}_A\colon
      \BracketTreeSet{A}\to\NonemptyWords{A}}%
    [TreeLeafWordTyping]
    \proofstepwidestar[]{%
      \operatorname{ein}_A\colon A\to \NonemptyWords{A}}{%
      \FormulaRefAuto{LetterWordFunctionDef}[def]}
    \proofstepwidestar[]{%
      \mathsf{BinOp}\bigl(\NonemptyWords{A},\frown\bigr)}{%
      \FormulaRefAuto{NonemptyWordConcatenationBinaryOperation}[thm]}
    \proofstepwidestar[]{%
      \operatorname{wort}_A=R}{%
      \FormulaRefAuto{TreeLeafWordDef}[def]}
    \proofstepwidestar[]{%
      R=\operatorname{wort}_A}{%
      \FormulaRefAuto{a=b\vdash b=a}[thm]{3}}
    \proofstepwidestar[]{%
      R\colon\BracketTreeSet A\to \NonemptyWords{A}}{%
      \FormulaRefAuto{TreeRecursorBinaryOperationTyping}[thm]{1,2}}
    \proofstepwidestar[]{%
      \operatorname{wort}_A\colon\BracketTreeSet A\to \NonemptyWords{A}}{%
      \rIE{4,5}}
  \closeproofpartwideindR

  \proofpartwideindR[Blattgleichung]{%
    a\in A\vdash
      \TreeLeafWord{A}{\LeafTree{A}{a}}=\LetterWord{a}
  }{a\in A\vdash
      \TreeLeafWord{A}{\LeafTree{A}{a}}=\LetterWord{a}}%
    [TreeLeafWordLeafEquation]
    \proofstepwidestar[1]{%
      a\in A}{%
      \rA}
    \proofstepwidestar[]{%
      \operatorname{ein}_A\colon A\to \NonemptyWords{A}}{%
      \FormulaRefAuto{LetterWordFunctionDef}[def]}
    \proofstepwidestar[]{%
      \mathsf{BinOp}\bigl(\NonemptyWords{A},\frown\bigr)}{%
      \FormulaRefAuto{NonemptyWordConcatenationBinaryOperation}[thm]}
    \proofstepwidestar[]{%
      \operatorname{wort}_A=R}{%
      \FormulaRefAuto{TreeLeafWordDef}[def]}
    \proofstepwidestar[]{%
      R=\operatorname{wort}_A}{%
      \FormulaRefAuto{a=b\vdash b=a}[thm]{4}}
    \proofstepwidestar[1]{%
      R(\LeafTree A a)=\operatorname{ein}_A(a)}{%
      \FormulaRefAuto{TreeRecursorBinaryOperationLeafEquation}[thm]{2,3,1}}
    \proofstepwidestar[1]{%
      \operatorname{ein}_A(a)=\LetterWord{a}}{%
      \FormulaRefAuto{LetterWordFunctionDef}[def]{1}}
    \proofstepwidestar[1]{%
      \operatorname{wort}_A(\LeafTree A a)=\operatorname{ein}_A(a)}{%
      \rIE{5,6}}
    \proofstepwidestar[1]{%
      \operatorname{wort}_A(\LeafTree A a)=\LetterWord{a}}{%
      \FormulaRefAuto{a=b\dsep b=c\vdash a=c}[thm]{8,7}}
  \closeproofpartwideindR

  \proofpartwideindR[Knotengleichung]{%
    S,T\in\BracketTreeSet{A}\vdash
      \TreeLeafWord{A}{\NodeTree{A}{S}{T}}
      =\WordConcat{\TreeLeafWord{A}{S}}{\TreeLeafWord{A}{T}}
  }{S,T\in\BracketTreeSet{A}\vdash
      \TreeLeafWord{A}{\NodeTree{A}{S}{T}}
      =\WordConcat{\TreeLeafWord{A}{S}}{\TreeLeafWord{A}{T}}}%
    [TreeLeafWordNodeEquation]
    \proofstepwidestar[1]{%
      S\in\BracketTreeSet A}{%
      \rA}
    \proofstepwidestar[2]{%
      T\in\BracketTreeSet A}{%
      \rA}
    \proofstepwidestar[]{%
      \operatorname{ein}_A\colon A\to \NonemptyWords{A}}{%
      \FormulaRefAuto{LetterWordFunctionDef}[def]}
    \proofstepwidestar[]{%
      \mathsf{BinOp}\bigl(\NonemptyWords{A},\frown\bigr)}{%
      \FormulaRefAuto{NonemptyWordConcatenationBinaryOperation}[thm]}
    \proofstepwidestar[]{%
      \operatorname{wort}_A=R}{%
      \FormulaRefAuto{TreeLeafWordDef}[def]}
    \proofstepwidestar[]{%
      R=\operatorname{wort}_A}{%
      \FormulaRefAuto{a=b\vdash b=a}[thm]{5}}
    \proofstepwidestar[1,2]{%
      R(\NodeTree A S T)=\WordConcat{R(S)}{R(T)}}{%
      \FormulaRefAuto{TreeRecursorBinaryOperationNodeEquation}[thm]{3,4,\rAI{1,2}}}
    \proofstepwidestar[1,2]{%
      \begin{aligned}[t]&\operatorname{wort}_A(\NodeTree A S T)\\[-2pt]
      &\quad=\WordConcat{\operatorname{wort}_A(S)}{\operatorname{wort}_A(T)}\end{aligned}}{%
      \rIE{6,7}}
  \closeproofpartwideindR
\end{tabproofsplitwide}

% Definierte Konstruktionsnotation auf Grundlage der Blattwortabbildung.
% Einbindung nach TreeLeafWordEquations und vor den Klammerungsmengen.

\subsection{Ein Baum klammert ein Wort}

Die Schreibweise \(\TreeBrackets{A}{T}{w}\) wird als
\glqq\(T\) klammert \(w\)\grqq\ gelesen. Sie bezeichnet die Zugehörigkeit
zum bereits konstruierten Funktionsgraphen des Blattwortes. Die folgenden
Sätze stellen dessen Typisierung und Konstruktionsgesetze für kurze
formale Ableitungen bereit.

\FormulaDefDeltaK[Ein Baum klammert ein Wort]{%
  \TreeBrackets{A}{T}{w}\coloneqq
    (T,w)\in\operatorname{wort}_A%
}{TreeBracketingRelationDef}{%
  \DeltaRow{Mengen}{A\dsep T\dsep w}
  \DeltaRow{Funktionen}{\operatorname{wort}_A}
  \DeltaRow{\textbf{Neues Symbol}}{}
  \DeltaPrem{Klammerungsrelation}{\TreeBrackets{A}{T}{w}}
}

\FormulaThmDeltaK[Baumtypisierung der Klammerungsrelation]{%
  \TreeBrackets{A}{T}{w}\vdash T\in\BracketTreeSet{A}%
}{TreeBracketingRelationTreeTyping}{%
  \DeltaRow{Mengen}{A\dsep T\dsep w}
}
\begin{tabproofwide}
  \proofstepwidestar[1]{\TreeBrackets{A}{T}{w}}{\rA}
  \proofstepwidestar[1]{(T,w)\in\operatorname{wort}_A}{%
    \FormulaRefAuto{TreeBracketingRelationDef}[def]{1}}
  \proofstepwidestar[]{%
    \operatorname{wort}_A\colon
      \BracketTreeSet{A}\to\NonemptyWords{A}}{%
    \FormulaRefAuto{TreeLeafWordTyping}[thm]}
  \proofstepwidestar[1]{T\in\BracketTreeSet{A}}{%
    \FormulaRefAuto{F\colon A\to B\dsep (x,y)\in F\vdash
      x\in A}[thm]{3,2}}
\end{tabproofwide}

\FormulaThmDeltaK[Worttypisierung der Klammerungsrelation]{%
  \TreeBrackets{A}{T}{w}\vdash w\in\NonemptyWords{A}%
}{TreeBracketingRelationWordTyping}{%
  \DeltaRow{Mengen}{A\dsep T\dsep w}
}
\begin{tabproofwide}
  \proofstepwidestar[1]{\TreeBrackets{A}{T}{w}}{\rA}
  \proofstepwidestar[1]{(T,w)\in\operatorname{wort}_A}{%
    \FormulaRefAuto{TreeBracketingRelationDef}[def]{1}}
  \proofstepwidestar[]{%
    \operatorname{wort}_A\colon
      \BracketTreeSet{A}\to\NonemptyWords{A}}{%
    \FormulaRefAuto{TreeLeafWordTyping}[thm]}
  \proofstepwidestar[1]{w\in\NonemptyWords{A}}{%
    \FormulaRefAuto{F\colon A\to B\dsep (x,y)\in F\vdash
      y\in B}[thm]{3,2}}
\end{tabproofwide}

\FormulaThmDeltaK[Auswertung der Klammerungsrelation]{%
  \TreeBrackets{A}{T}{w}\vdash\TreeLeafWord{A}{T}=w%
}{TreeBracketingRelationElimination}{%
  \DeltaRow{Mengen}{A\dsep T\dsep w}
  \DeltaRow{Funktionen}{\operatorname{wort}_A}
}
\begin{tabproofwide}
  \proofstepwidestar[1]{\TreeBrackets{A}{T}{w}}{\rA}
  \proofstepwidestar[1]{(T,w)\in\operatorname{wort}_A}{%
    \FormulaRefAuto{TreeBracketingRelationDef}[def]{1}}
  \proofstepwidestar[]{%
    \operatorname{wort}_A\colon
      \BracketTreeSet{A}\to\NonemptyWords{A}}{%
    \FormulaRefAuto{TreeLeafWordTyping}[thm]}
  \proofstepwidestar[1]{\TreeLeafWord{A}{T}=w}{%
    \FormulaRefAuto{F\colon A\to B\dsep (x,y)\in F\vdash
      F(x)=y}[thm]{3,2}}
\end{tabproofwide}

\FormulaThmDeltaK[Einführung der Klammerungsrelation]{%
  T\in\BracketTreeSet{A}\dsep\TreeLeafWord{A}{T}=w\vdash
    \TreeBrackets{A}{T}{w}%
}{TreeBracketingRelationIntroduction}{%
  \DeltaRow{Mengen}{A\dsep T\dsep w}
  \DeltaRow{Funktionen}{\operatorname{wort}_A}
}
\begin{tabproofwide}
  \proofstepwidestar[1]{T\in\BracketTreeSet{A}}{\rA}
  \proofstepwidestar[2]{\TreeLeafWord{A}{T}=w}{\rA}
  \proofstepwidestar[]{%
    \operatorname{wort}_A\colon
      \BracketTreeSet{A}\to\NonemptyWords{A}}{%
    \FormulaRefAuto{TreeLeafWordTyping}[thm]}
  \proofstepwidestar[1,2]{(T,w)\in\operatorname{wort}_A}{%
    \FormulaRefAuto{F\colon A\to B\dsep x\in A\dsep
      F(x)=y\vdash (x,y)\in F}[thm]{3,1,2}}
  \proofstepwidestar[1,2]{\TreeBrackets{A}{T}{w}}{%
    \FormulaRefAuto{TreeBracketingRelationDef}[def]{4}}
\end{tabproofwide}

\subsection{Blattregel und Knotenregel}

Ein Blatt klammert das zugehörige Buchstabenwort. Ein Knoten klammert
die Konkatenation der Wörter seiner beiden Teilbäume. Die Regeln
enthalten damit zugleich die benötigten Aussagen über den Baum und
sein Blattwort.

\FormulaThmDeltaK[Blattregel der Klammerungsrelation]{%
  a\in A\vdash
    \TreeBrackets{A}{\LeafTree{A}{a}}{\LetterWord{a}}%
}{TreeBracketingLeafIntroduction}{%
  \DeltaRow{Mengen}{A\dsep a}
}
\begin{tabproofwide}
  \proofstepwidestar[1]{a\in A}{\rA}
  \proofstepwidestar[1]{\LeafTree{A}{a}\in\BracketTreeSet{A}}{%
    \FormulaRefAuto{TreeLeafClosure}[thm]{1}}
  \proofstepwidestar[1]{%
    \TreeLeafWord{A}{\LeafTree{A}{a}}=\LetterWord{a}}{%
    \FormulaRefAuto{TreeLeafWordLeafEquation}[thm]{1}}
  \proofstepwidestar[1]{%
    \TreeBrackets{A}{\LeafTree{A}{a}}{\LetterWord{a}}}{%
    \FormulaRefAuto{TreeBracketingRelationIntroduction}[thm]{2,3}}
\end{tabproofwide}

\FormulaThmDeltaK[Knotenregel der Klammerungsrelation]{%
  \begin{aligned}[t]
    &\TreeBrackets{A}{S}{u}\dsep\TreeBrackets{A}{T}{v}\vdash{}\\[-2pt]
    &\TreeBrackets{A}{\NodeTree{A}{S}{T}}{\WordConcat{u}{v}}
  \end{aligned}%
}{TreeBracketingNodeIntroduction}{%
  \DeltaRow{Mengen}{A\dsep S\dsep T\dsep u\dsep v}
}
\begin{tabproofwide}
  \proofstepwidestar[1]{\TreeBrackets{A}{S}{u}}{\rA}
  \proofstepwidestar[2]{\TreeBrackets{A}{T}{v}}{\rA}
  \proofstepwidestar[1]{S\in\BracketTreeSet{A}}{%
    \FormulaRefAuto{TreeBracketingRelationTreeTyping}[thm]{1}}
  \proofstepwidestar[2]{T\in\BracketTreeSet{A}}{%
    \FormulaRefAuto{TreeBracketingRelationTreeTyping}[thm]{2}}
  \proofstepwidestar[1]{u\in\NonemptyWords{A}}{%
    \FormulaRefAuto{TreeBracketingRelationWordTyping}[thm]{1}}
  \proofstepwidestar[2]{v\in\NonemptyWords{A}}{%
    \FormulaRefAuto{TreeBracketingRelationWordTyping}[thm]{2}}
  \proofstepwidestar[1]{u\in\FiniteWords{A}}{%
    \FormulaRefAuto{NonemptyWordFinite}[thm]{5}}
  \proofstepwidestar[2]{v\in\FiniteWords{A}}{%
    \FormulaRefAuto{NonemptyWordFinite}[thm]{6}}
  \proofstepwidestar[1]{\TreeLeafWord{A}{S}=u}{%
    \FormulaRefAuto{TreeBracketingRelationElimination}[thm]{1}}
  \proofstepwidestar[2]{\TreeLeafWord{A}{T}=v}{%
    \FormulaRefAuto{TreeBracketingRelationElimination}[thm]{2}}
  \proofstepwidestar[1,2]{S,T\in\BracketTreeSet{A}}{\rAI{3,4}}
  \proofstepwidestar[1,2]{%
    \NodeTree{A}{S}{T}\in\BracketTreeSet{A}}{%
    \FormulaRefAuto{TreeNodeClosure}[thm]{11}}
  \proofstepwidestar[1,2]{%
    \begin{aligned}[t]
      &\TreeLeafWord{A}{\NodeTree{A}{S}{T}}\\[-2pt]
      &\quad=\WordConcat{\TreeLeafWord{A}{S}}{\TreeLeafWord{A}{T}}
    \end{aligned}}{%
    \FormulaRefAuto{TreeLeafWordNodeEquation}[thm]{11}}
  \proofstepwidestar[1,2]{%
    \begin{aligned}[t]
      &\TreeLeafWord{A}{\NodeTree{A}{S}{T}}\\[-2pt]
      &\quad=\WordConcat{u}{\TreeLeafWord{A}{T}}
    \end{aligned}}{\rIE{9,13}}
  \proofstepwidestar[1,2]{%
    \TreeLeafWord{A}{\NodeTree{A}{S}{T}}=\WordConcat{u}{v}}{%
    \rIE{10,14}}
  \proofstepwidestar[1,2]{%
    \TreeBrackets{A}{\NodeTree{A}{S}{T}}{\WordConcat{u}{v}}}{%
    \FormulaRefAuto{TreeBracketingRelationIntroduction}[thm]{12,15}}
\end{tabproofwide}

\subsection{Ein Blatt rechts anfügen}

Der folgende Operator bildet einen neuen Knoten mit dem bisherigen Baum
als linkem Teilbaum und einem neuen Blatt als rechtem Teilbaum. In der
Anfügeregel werden Baum und Wort dadurch in derselben Richtung erweitert.

\FormulaDefDeltaK[Rechtsanfügen eines Blattes]{%
  \begin{aligned}[t]
    &T\in\BracketTreeSet{A}\dsep a\in A\vdash{}\\[-2pt]
    &\TreeAppendLeaf{A}{T}{a}\coloneqq
      \NodeTree{A}{T}{\LeafTree{A}{a}}
  \end{aligned}%
}{TreeAppendLeafDef}{%
  \DeltaRow{Mengen}{A\dsep T\dsep a}
  \DeltaRow{\textbf{Neues Symbol}}{}
  \DeltaPrem{Blattanfügung}{\TreeAppendLeaf{A}{T}{a}}
}

\FormulaThmDeltaK[Die Blattanfügung ist ein Klammerungsbaum]{%
  T\in\BracketTreeSet{A}\dsep a\in A\vdash
    \TreeAppendLeaf{A}{T}{a}\in\BracketTreeSet{A}%
}{TreeAppendLeafTyping}{%
  \DeltaRow{Mengen}{A\dsep T\dsep a}
}
\begin{tabproofwide}
  \proofstepwidestar[1]{T\in\BracketTreeSet{A}}{\rA}
  \proofstepwidestar[2]{a\in A}{\rA}
  \proofstepwidestar[2]{\LeafTree{A}{a}\in\BracketTreeSet{A}}{%
    \FormulaRefAuto{TreeLeafClosure}[thm]{2}}
  \proofstepwidestar[1,2]{T,\LeafTree{A}{a}\in\BracketTreeSet{A}}{%
    \rAI{1,3}}
  \proofstepwidestar[1,2]{%
    \NodeTree{A}{T}{\LeafTree{A}{a}}\in\BracketTreeSet{A}}{%
    \FormulaRefAuto{TreeNodeClosure}[thm]{4}}
  \proofstepwidestar[1,2]{%
    \TreeAppendLeaf{A}{T}{a}=\NodeTree{A}{T}{\LeafTree{A}{a}}}{%
    \FormulaRefAuto{TreeAppendLeafDef}[def]{1,2}}
  \proofstepwidestar[1,2]{%
    \NodeTree{A}{T}{\LeafTree{A}{a}}=\TreeAppendLeaf{A}{T}{a}}{%
    \FormulaRefAuto{a=b\vdash b=a}[thm]{6}}
  \proofstepwidestar[1,2]{%
    \TreeAppendLeaf{A}{T}{a}\in\BracketTreeSet{A}}{\rIE{7,5}}
\end{tabproofwide}

\FormulaThmDeltaK[Anfügeregel der Klammerungsrelation]{%
  \begin{aligned}[t]
    &\TreeBrackets{A}{T}{u}\dsep a\in A\vdash{}\\[-2pt]
    &\TreeBrackets{A}{\TreeAppendLeaf{A}{T}{a}}
      {\WordConcat{u}{\LetterWord{a}}}
  \end{aligned}%
}{TreeBracketingAppendLeafIntroduction}{%
  \DeltaRow{Mengen}{A\dsep T\dsep u\dsep a}
}
\begin{tabproofwide}
  \proofstepwidestar[1]{\TreeBrackets{A}{T}{u}}{\rA}
  \proofstepwidestar[2]{a\in A}{\rA}
  \proofstepwidestar[1]{T\in\BracketTreeSet{A}}{%
    \FormulaRefAuto{TreeBracketingRelationTreeTyping}[thm]{1}}
  \proofstepwidestar[2]{%
    \TreeBrackets{A}{\LeafTree{A}{a}}{\LetterWord{a}}}{%
    \FormulaRefAuto{TreeBracketingLeafIntroduction}[thm]{2}}
  \proofstepwidestar[1,2]{%
    \TreeBrackets{A}{\NodeTree{A}{T}{\LeafTree{A}{a}}}
      {\WordConcat{u}{\LetterWord{a}}}}{%
    \FormulaRefAuto{TreeBracketingNodeIntroduction}[thm]{1,4}}
  \proofstepwidestar[1,2]{%
    \TreeAppendLeaf{A}{T}{a}=\NodeTree{A}{T}{\LeafTree{A}{a}}}{%
    \FormulaRefAuto{TreeAppendLeafDef}[def]{3,2}}
  \proofstepwidestar[1,2]{%
    \NodeTree{A}{T}{\LeafTree{A}{a}}=\TreeAppendLeaf{A}{T}{a}}{%
    \FormulaRefAuto{a=b\vdash b=a}[thm]{6}}
  \proofstepwidestar[1,2]{%
    \TreeBrackets{A}{\TreeAppendLeaf{A}{T}{a}}
      {\WordConcat{u}{\LetterWord{a}}}}{\rIE{7,5}}
\end{tabproofwide}

% Kanonische Linksklammerung durch die bereits bewiesene Wortrekursion.
% Einbindung nach b27-construction-rules und vor den Klammerungsmengen.

\subsection{Der Linksbaum eines Wortes}

Der Linksbaum beginnt mit einem Blatt und fügt jeden weiteren Buchstaben
rechts an. Die Wortrekursion liefert hierfür eine eindeutig bestimmte
Funktion. Zunächst werden ihre Anfangs- und Schrittfunktion als
Funktionsgraphen bereitgestellt.

\FormulaThmDeltaK[Existenz der Blattfunktion]{%
  \exists!f\colon A\to\BracketTreeSet A\,
    \forall a\in A\,f(a)=\LeafTree A a
}{LeftBracketingLeafFunctionExistenceUnique}{%
  \DeltaRow{Mengen}{A\dsep a}
  \DeltaRow{Funktionen}{f}
}
\begin{tabproofwide}
  \proofstepwidestar[1]{a\in A}{\rA}
  \proofstepwidestar[1]{\LeafTree A a\in\BracketTreeSet A}{%
    \FormulaRefAuto{TreeLeafClosure}[thm]{1}}
  \proofstepwidestar[]{\forall a\in A\,\LeafTree A a\in\BracketTreeSet A}{%
    \rUI{\rRI{1,2}}}
  \proofstepwidestar[]{%
    \exists!f\colon A\to\BracketTreeSet A\,
      \forall a\in A\,f(a)=\LeafTree A a}{%
    \FormulaRefAuto{TypedFunctionDefinitionPrinciple}[thm]{3}}
\end{tabproofwide}

\FormulaDefDeltaK[Blattfunktion für die Wortrekursion]{%
  \operatorname{blatt}_A\coloneqq
    \iota f\colon A\to\BracketTreeSet A\,
      \forall a\in A\,f(a)=\LeafTree A a
}{LeftBracketingLeafFunctionDef}{%
  \DeltaRow{Mengen}{A\dsep a}
  \DeltaRow{Funktionen}{f\dsep\operatorname{blatt}_A}
  \DeltaRow{\textbf{Neues Symbol}}{}
  \DeltaPrem{Blattfunktion}{\operatorname{blatt}_A}
}

\FormulaThmDeltaK[Kennzeichnung der Blattfunktion]{%
  \operatorname{blatt}_A\colon A\to\BracketTreeSet A\land
    \forall a\in A\,\operatorname{blatt}_A(a)=\LeafTree A a
}{LeftBracketingLeafFunctionCharacterization}{%
  \DeltaRow{Mengen}{A\dsep a}
  \DeltaRow{Funktionen}{f\dsep\operatorname{blatt}_A}
}
\begin{tabproofwide}
  \proofstepwidestar[]{%
    \exists!f\colon A\to\BracketTreeSet A\,
      \forall a\in A\,f(a)=\LeafTree A a}{%
    \FormulaRefAuto{LeftBracketingLeafFunctionExistenceUnique}[thm]}
  \proofstepwidestar[]{%
    \operatorname{blatt}_A\colon A\to\BracketTreeSet A\land
      \forall a\in A\,\operatorname{blatt}_A(a)=\LeafTree A a}{%
    \FormulaRefAuto{LeftBracketingLeafFunctionDef}[def]{1}}
\end{tabproofwide}

\FormulaThmDeltaK[Existenz der Anfügefunktion]{%
  \begin{aligned}[t]
    &\exists!s\colon\BracketTreeSet A\times A\to\BracketTreeSet A\,\\[-2pt]
    &\quad\forall p\in\BracketTreeSet A\times A\,
      s(p)=\TreeAppendLeaf A{\pi_1(p)}{\pi_2(p)}
  \end{aligned}
}{LeftBracketingAppendFunctionExistenceUnique}{%
  \DeltaRow{Mengen}{A\dsep p}
  \DeltaRow{Funktionen}{s}
}
\begin{tabproofwide}
  \proofstepwidestar[1]{p\in\BracketTreeSet A\times A}{\rA}
  \proofstepwidestar[1]{\pi_1(p)\in\BracketTreeSet A}{%
    \FormulaRefAuto{z\in A\times B\vdash\pi_1(z)\in A}[thm]{1}}
  \proofstepwidestar[1]{\pi_2(p)\in A}{%
    \FormulaRefAuto{z\in A\times B\vdash\pi_2(z)\in B}[thm]{1}}
  \proofstepwidestar[1]{%
    \TreeAppendLeaf A{\pi_1(p)}{\pi_2(p)}\in\BracketTreeSet A}{%
    \FormulaRefAuto{TreeAppendLeafTyping}[thm]{2,3}}
  \proofstepwidestar[]{%
    \forall p\in\BracketTreeSet A\times A\,
      \TreeAppendLeaf A{\pi_1(p)}{\pi_2(p)}\in\BracketTreeSet A}{%
    \rUI{\rRI{1,4}}}
  \proofstepwidestar[]{%
    \begin{aligned}[t]
      &\exists!s\colon\BracketTreeSet A\times A\to\BracketTreeSet A\,\\[-2pt]
      &\quad\forall p\in\BracketTreeSet A\times A\,
        s(p)=\TreeAppendLeaf A{\pi_1(p)}{\pi_2(p)}
    \end{aligned}}{%
    \FormulaRefAuto{TypedFunctionDefinitionPrinciple}[thm]{5}}
\end{tabproofwide}

\FormulaDefDeltaK[Anfügefunktion für die Wortrekursion]{%
  \begin{aligned}[t]
    &\operatorname{anf}_A\coloneqq
      \iota s\colon\BracketTreeSet A\times A\to\BracketTreeSet A\,\\[-2pt]
    &\quad\forall p\in\BracketTreeSet A\times A\,
      s(p)=\TreeAppendLeaf A{\pi_1(p)}{\pi_2(p)}
  \end{aligned}
}{LeftBracketingAppendFunctionDef}{%
  \DeltaRow{Mengen}{A\dsep p}
  \DeltaRow{Funktionen}{s\dsep\operatorname{anf}_A}
  \DeltaRow{\textbf{Neues Symbol}}{}
  \DeltaPrem{Anfügefunktion}{\operatorname{anf}_A}
}

\FormulaThmDeltaK[Kennzeichnung der Anfügefunktion]{%
  \begin{aligned}[t]
    &\operatorname{anf}_A\colon
      \BracketTreeSet A\times A\to\BracketTreeSet A\land{}\\[-2pt]
    &\forall p\in\BracketTreeSet A\times A\,
      \operatorname{anf}_A(p)=\TreeAppendLeaf A{\pi_1(p)}{\pi_2(p)}
  \end{aligned}
}{LeftBracketingAppendFunctionCharacterization}{%
  \DeltaRow{Mengen}{A\dsep p}
  \DeltaRow{Funktionen}{s\dsep\operatorname{anf}_A}
}
\begin{tabproofwide}
  \proofstepwidestar[]{%
    \begin{aligned}[t]
      &\exists!s\colon\BracketTreeSet A\times A\to\BracketTreeSet A\,\\[-2pt]
      &\quad\forall p\in\BracketTreeSet A\times A\,
        s(p)=\TreeAppendLeaf A{\pi_1(p)}{\pi_2(p)}
    \end{aligned}}{%
    \FormulaRefAuto{LeftBracketingAppendFunctionExistenceUnique}[thm]}
  \proofstepwidestar[]{%
    \begin{aligned}[t]
      &\operatorname{anf}_A\colon
        \BracketTreeSet A\times A\to\BracketTreeSet A\land{}\\[-2pt]
      &\forall p\in\BracketTreeSet A\times A\,
        \operatorname{anf}_A(p)=\TreeAppendLeaf A{\pi_1(p)}{\pi_2(p)}
    \end{aligned}}{%
    \FormulaRefAuto{LeftBracketingAppendFunctionDef}[def]{1}}
\end{tabproofwide}

\FormulaThmDeltaK[Auswertung der Anfügefunktion an einem Paar]{%
  T\in\BracketTreeSet A\dsep a\in A\vdash
    \operatorname{anf}_A(T,a)=\TreeAppendLeaf A T a
}{LeftBracketingAppendFunctionPairEquation}{%
  \DeltaRow{Mengen}{A\dsep T\dsep a\dsep p}
  \DeltaRow{Funktionen}{\operatorname{anf}_A}
}
\begin{tabproofwide}
  \proofstepwidestar[1]{T\in\BracketTreeSet A}{\rA}
  \proofstepwidestar[2]{a\in A}{\rA}
  \proofstepwidestar[]{%
    \begin{aligned}[t]
      &\operatorname{anf}_A\colon
        \BracketTreeSet A\times A\to\BracketTreeSet A\land{}\\[-2pt]
      &\forall p\in\BracketTreeSet A\times A\,
        \operatorname{anf}_A(p)=\TreeAppendLeaf A{\pi_1(p)}{\pi_2(p)}
    \end{aligned}}{%
    \FormulaRefAuto{LeftBracketingAppendFunctionCharacterization}[thm]}
  \proofstepwidestar[]{%
    \forall p\in\BracketTreeSet A\times A\,
      \operatorname{anf}_A(p)=\TreeAppendLeaf A{\pi_1(p)}{\pi_2(p)}}{\rAEb{3}}
  \proofstepwidestar[1,2]{(T,a)\in\BracketTreeSet A\times A}{%
    \FormulaRefAuto{a\in A\dsep b\in B\vdash(a,b)\in A\times B}[thm]{1,2}}
  \proofstepwidestar[1,2]{\pi_1(T,a)=T}{%
    \FormulaRefAuto{(x,y)\in A\times B\vdash\pi_1(x,y)=x}[thm]{5}}
  \proofstepwidestar[1,2]{\pi_2(T,a)=a}{%
    \FormulaRefAuto{(x,y)\in A\times B\vdash\pi_2(x,y)=y}[thm]{5}}
  \proofstepwidestar[1,2]{%
    \operatorname{anf}_A(T,a)=\TreeAppendLeaf A{\pi_1(T,a)}{\pi_2(T,a)}}{%
    \rRE{5,\rUE{4}}}
  \proofstepwidestar[1,2]{%
    \operatorname{anf}_A(T,a)=\TreeAppendLeaf A T{\pi_2(T,a)}}{\rIE{6,8}}
  \proofstepwidestar[1,2]{%
    \operatorname{anf}_A(T,a)=\TreeAppendLeaf A T a}{\rIE{7,9}}
\end{tabproofwide}

% Reine Satzhilfe: Jeder Aufruf druckt die vollständige Rekursionsformel.
\newcommand{\BXXVIILeftBracketingRecursionProperty}[1]{%
  \begin{aligned}[t]
    &#1\colon\NonemptyWords A\to\BracketTreeSet A\land{}\\[-2pt]
    &\Bigl(\forall a\in A\,
      #1(\LetterWord a)=\operatorname{blatt}_A(a)\land{}\\[-2pt]
    &\quad\forall u\in\NonemptyWords A\,\forall a\in A\,\\[-2pt]
    &\qquad #1(\WordConcat u{\LetterWord a})
      =\operatorname{anf}_A(#1(u),a)\Bigr)
  \end{aligned}}

\newcommand{\BXXVIILeftBracketingQuantifiedProperty}[1]{%
  \begin{aligned}[t]
    &#1 F\,\Bigl(F\colon\NonemptyWords A\to\BracketTreeSet A\land{}\\[-2pt]
    &\quad\Bigl(\forall a\in A\,
      F(\LetterWord a)=\operatorname{blatt}_A(a)\land{}\\[-2pt]
    &\qquad\forall u\in\NonemptyWords A\,\forall a\in A\,\\[-2pt]
    &\qquad\quad F(\WordConcat u{\LetterWord a})
      =\operatorname{anf}_A(F(u),a)\Bigr)\Bigr)
  \end{aligned}}

\FormulaThmDeltaK[Existenz und Eindeutigkeit des Linksbaums]{%
  \BXXVIILeftBracketingQuantifiedProperty{\exists!}
}{LeftBracketingExistenceUnique}{%
  \DeltaRow{Mengen}{A\dsep a\dsep u\dsep p}
  \DeltaRow{Funktionen}{F\dsep\operatorname{blatt}_A\dsep\operatorname{anf}_A}
}
\begin{tabproofwide}
  \proofstepwidestar[]{%
    \operatorname{blatt}_A\colon A\to\BracketTreeSet A\land
      \forall a\in A\,\operatorname{blatt}_A(a)=\LeafTree A a}{%
    \FormulaRefAuto{LeftBracketingLeafFunctionCharacterization}[thm]}
  \proofstepwidestar[]{\operatorname{blatt}_A\colon A\to\BracketTreeSet A}{\rAEa{1}}
  \proofstepwidestar[]{%
    \begin{aligned}[t]
      &\operatorname{anf}_A\colon
        \BracketTreeSet A\times A\to\BracketTreeSet A\land{}\\[-2pt]
      &\forall p\in\BracketTreeSet A\times A\,
        \operatorname{anf}_A(p)=\TreeAppendLeaf A{\pi_1(p)}{\pi_2(p)}
    \end{aligned}}{%
    \FormulaRefAuto{LeftBracketingAppendFunctionCharacterization}[thm]}
  \proofstepwidestar[]{%
    \operatorname{anf}_A\colon\BracketTreeSet A\times A\to\BracketTreeSet A}{\rAEa{3}}
  \proofstepwidestar[]{%
    \BXXVIILeftBracketingQuantifiedProperty{\exists!}}{%
    \FormulaRefAuto{NonemptyWordRecursion}[thm]{2,4}}
\end{tabproofwide}

\FormulaDefDeltaK[Linksbaum eines nichtleeren Wortes]{%
  \LeftBracketingMap A\coloneqq
    \BXXVIILeftBracketingQuantifiedProperty{\iota}
}{LeftBracketingDef}{%
  \DeltaRow{Mengen}{A\dsep a\dsep u}
  \DeltaRow{Funktionen}{F\dsep\LeftBracketingMap A\dsep
    \operatorname{blatt}_A\dsep\operatorname{anf}_A}
  \DeltaRow{\textbf{Neues Symbol}}{}
  \DeltaPrem{Linksbaumabbildung}{\LeftBracketingMap A}
}

\FormulaThmDeltaK[Rekursionskennzeichnung der Linksbaumabbildung]{%
  \BXXVIILeftBracketingRecursionProperty{\LeftBracketingMap A}
}{LeftBracketingRecursionCharacterization}{%
  \DeltaRow{Mengen}{A\dsep a\dsep u}
  \DeltaRow{Funktionen}{F\dsep\LeftBracketingMap A\dsep
    \operatorname{blatt}_A\dsep\operatorname{anf}_A}
}
\begin{tabproofwide}
  \proofstepwidestar[]{%
    \BXXVIILeftBracketingQuantifiedProperty{\exists!}}{%
    \FormulaRefAuto{LeftBracketingExistenceUnique}[thm]}
  \proofstepwidestar[]{%
    \BXXVIILeftBracketingRecursionProperty{\LeftBracketingMap A}}{%
    \FormulaRefAuto{LeftBracketingDef}[def]{1}}
\end{tabproofwide}

\FormulaThmDeltaK[Typisierung der Linksbaumabbildung]{%
  \LeftBracketingMap A\colon\NonemptyWords A\to\BracketTreeSet A
}{LeftBracketingTyping}{%
  \DeltaRow{Mengen}{A\dsep a\dsep u}
  \DeltaRow{Funktionen}{\LeftBracketingMap A\dsep
    \operatorname{blatt}_A\dsep\operatorname{anf}_A}
}
\begin{tabproofwide}
  \proofstepwidestar[]{%
    \BXXVIILeftBracketingRecursionProperty{\LeftBracketingMap A}}{%
    \FormulaRefAuto{LeftBracketingRecursionCharacterization}[thm]}
  \proofstepwidestar[]{%
    \LeftBracketingMap A\colon\NonemptyWords A\to\BracketTreeSet A}{\rAEa{1}}
\end{tabproofwide}

\FormulaThmDeltaK[Der Wert der Linksbaumabbildung ist ein Baum]{%
  w\in\NonemptyWords A\vdash\LeftBracketing A w\in\BracketTreeSet A
}{LeftBracketingValueTyping}{%
  \DeltaRow{Mengen}{A\dsep w}
  \DeltaRow{Funktionen}{\LeftBracketingMap A}
}
\begin{tabproofwide}
  \proofstepwidestar[1]{w\in\NonemptyWords A}{\rA}
  \proofstepwidestar[]{%
    \LeftBracketingMap A\colon\NonemptyWords A\to\BracketTreeSet A}{%
    \FormulaRefAuto{LeftBracketingTyping}[thm]}
  \proofstepwidestar[1]{\LeftBracketing A w\in\BracketTreeSet A}{%
    \FormulaRefAuto{F\colon A\to B\dsep x\in A\vdash F(x)\in B}[thm]{2,1}}
\end{tabproofwide}

\FormulaThmDeltaK[Der Linksbaum eines Buchstabenwortes ist ein Blatt]{%
  a\in A\vdash\LeftBracketing A{\LetterWord a}=\LeafTree A a
}{LeftBracketingLetterEquation}{%
  \DeltaRow{Mengen}{A\dsep a\dsep u}
  \DeltaRow{Funktionen}{\LeftBracketingMap A\dsep
    \operatorname{blatt}_A\dsep\operatorname{anf}_A}
}
\begin{tabproofwide}
  \proofstepwidestar[1]{a\in A}{\rA}
  \proofstepwidestar[]{%
    \BXXVIILeftBracketingRecursionProperty{\LeftBracketingMap A}}{%
    \FormulaRefAuto{LeftBracketingRecursionCharacterization}[thm]}
  \proofstepwidestar[]{%
    \forall a\in A\,\LeftBracketing A{\LetterWord a}=\operatorname{blatt}_A(a)}{%
    \rAEa{\rAEb{2}}}
  \proofstepwidestar[1]{%
    \LeftBracketing A{\LetterWord a}=\operatorname{blatt}_A(a)}{\rRE{1,\rUE{3}}}
  \proofstepwidestar[]{%
    \operatorname{blatt}_A\colon A\to\BracketTreeSet A\land
      \forall a\in A\,\operatorname{blatt}_A(a)=\LeafTree A a}{%
    \FormulaRefAuto{LeftBracketingLeafFunctionCharacterization}[thm]}
  \proofstepwidestar[1]{\operatorname{blatt}_A(a)=\LeafTree A a}{%
    \rRE{1,\rUE{\rAEb{5}}}}
  \proofstepwidestar[1]{\LeftBracketing A{\LetterWord a}=\LeafTree A a}{%
    \rChain{4,6}}
\end{tabproofwide}

\FormulaThmDeltaK[Der Linksbaum wächst durch Rechtsanfügen eines Blattes]{%
  \begin{aligned}[t]
    &u\in\NonemptyWords A\dsep a\in A\vdash{}\\[-2pt]
    &\LeftBracketing A{\WordConcat u{\LetterWord a}}
      =\TreeAppendLeaf A{\LeftBracketing A u}a
  \end{aligned}
}{LeftBracketingAppendEquation}{%
  \DeltaRow{Mengen}{A\dsep u\dsep a}
  \DeltaRow{Funktionen}{\LeftBracketingMap A\dsep
    \operatorname{blatt}_A\dsep\operatorname{anf}_A}
}
\begin{tabproofwide}
  \proofstepwidestar[1]{u\in\NonemptyWords A}{\rA}
  \proofstepwidestar[2]{a\in A}{\rA}
  \proofstepwidestar[]{%
    \BXXVIILeftBracketingRecursionProperty{\LeftBracketingMap A}}{%
    \FormulaRefAuto{LeftBracketingRecursionCharacterization}[thm]}
  \proofstepwidestar[]{%
    \LeftBracketingMap A\colon\NonemptyWords A\to\BracketTreeSet A}{\rAEa{3}}
  \proofstepwidestar[]{%
    \forall u\in\NonemptyWords A\,\forall a\in A\,
      \LeftBracketing A{\WordConcat u{\LetterWord a}}
        =\operatorname{anf}_A(\LeftBracketing A u,a)}{\rAEb{\rAEb{3}}}
  \proofstepwidestar[1]{\LeftBracketing A u\in\BracketTreeSet A}{%
    \FormulaRefAuto{F\colon A\to B\dsep x\in A\vdash F(x)\in B}[thm]{4,1}}
  \proofstepwidestar[1,2]{%
    \LeftBracketing A{\WordConcat u{\LetterWord a}}
      =\operatorname{anf}_A(\LeftBracketing A u,a)}{%
    \rRE{2,\rUE{\rRE{1,\rUE{5}}}}}
  \proofstepwidestar[1,2]{%
    \operatorname{anf}_A(\LeftBracketing A u,a)
      =\TreeAppendLeaf A{\LeftBracketing A u}a}{%
    \FormulaRefAuto{LeftBracketingAppendFunctionPairEquation}[thm]{6,2}}
  \proofstepwidestar[1,2]{%
    \LeftBracketing A{\WordConcat u{\LetterWord a}}
      =\TreeAppendLeaf A{\LeftBracketing A u}a}{\rChain{7,8}}
\end{tabproofwide}

Der Linksbaum ist damit für jedes nichtleere Wort festgelegt. Die nächste
Induktion zeigt, dass sein Blattwort tatsächlich das Ausgangswort ist.

\FormulaThmDeltaK[Der Linksbaum klammert sein Wort]{%
  w\in\NonemptyWords A\vdash
    \TreeBrackets A{\LeftBracketing A w}w
}{LeftBracketingProperty}{%
  \DeltaRow{Mengen}{A\dsep w\dsep u\dsep a}
  \DeltaRow{Funktionen}{\LeftBracketingMap A}
}
\begin{tabproofsplitwide}
  \proofpartwideindR[Ein Blatt klammert das Buchstabenwort]{%
    \forall a\in A\,
      \TreeBrackets A{\LeftBracketing A{\LetterWord a}}{\LetterWord a}
  }{%
    \forall a\in A\,
      \TreeBrackets A{\LeftBracketing A{\LetterWord a}}{\LetterWord a}
  }[LeftBracketingPropertyBase]
    \proofstepwidestar[1]{a\in A}{\rA}
    \proofstepwidestar[1]{\TreeBrackets A{\LeafTree A a}{\LetterWord a}}{%
      \FormulaRefAuto{TreeBracketingLeafIntroduction}[thm]{1}}
    \proofstepwidestar[1]{\LeftBracketing A{\LetterWord a}=\LeafTree A a}{%
      \FormulaRefAuto{LeftBracketingLetterEquation}[thm]{1}}
    \proofstepwidestar[1]{\LeafTree A a=\LeftBracketing A{\LetterWord a}}{%
      \FormulaRefAuto{a=b\vdash b=a}[thm]{3}}
    \proofstepwidestar[1]{%
      \TreeBrackets A{\LeftBracketing A{\LetterWord a}}{\LetterWord a}}{\rIE{4,2}}
    \proofstepwidestar[]{%
      \forall a\in A\,
        \TreeBrackets A{\LeftBracketing A{\LetterWord a}}{\LetterWord a}}{%
      \rUI{\rRI{1,5}}}
  \closeproofpartwideindR

  \proofpartwideindR[Baum und Wort wachsen gemeinsam]{%
    \begin{aligned}[t]
      &\forall u\in\NonemptyWords A\,\forall a\in A\,\Bigl(\\[-2pt]
      &\quad\TreeBrackets A{\LeftBracketing A u}u\rightarrow{}\\[-2pt]
      &\quad\TreeBrackets A{\LeftBracketing A{\WordConcat u{\LetterWord a}}}
        {\WordConcat u{\LetterWord a}}\Bigr)
    \end{aligned}
  }{%
    \begin{aligned}[t]
      &\forall u\in\NonemptyWords A\,\forall a\in A\,\Bigl(\\[-2pt]
      &\quad\TreeBrackets A{\LeftBracketing A u}u\rightarrow{}\\[-2pt]
      &\quad\TreeBrackets A{\LeftBracketing A{\WordConcat u{\LetterWord a}}}
        {\WordConcat u{\LetterWord a}}\Bigr)
    \end{aligned}
  }[LeftBracketingPropertyStep]
    \proofstepwidestar[1]{u\in\NonemptyWords A}{\rA}
    \proofstepwidestar[2]{a\in A}{\rA}
    \proofstepwidestar[3]{\TreeBrackets A{\LeftBracketing A u}u}{\rA}
    \proofstepwidestar[2,3]{%
      \TreeBrackets A{\TreeAppendLeaf A{\LeftBracketing A u}a}
        {\WordConcat u{\LetterWord a}}}{%
      \FormulaRefAuto{TreeBracketingAppendLeafIntroduction}[thm]{3,2}}
    \proofstepwidestar[1,2]{%
      \LeftBracketing A{\WordConcat u{\LetterWord a}}
        =\TreeAppendLeaf A{\LeftBracketing A u}a}{%
      \FormulaRefAuto{LeftBracketingAppendEquation}[thm]{1,2}}
    \proofstepwidestar[1,2]{%
      \TreeAppendLeaf A{\LeftBracketing A u}a
        =\LeftBracketing A{\WordConcat u{\LetterWord a}}}{%
      \FormulaRefAuto{a=b\vdash b=a}[thm]{5}}
    \proofstepwidestar[1,2,3]{%
      \TreeBrackets A{\LeftBracketing A{\WordConcat u{\LetterWord a}}}
        {\WordConcat u{\LetterWord a}}}{\rIE{6,4}}
    \proofstepwidestar[1,2]{%
      \begin{aligned}[t]
        &\TreeBrackets A{\LeftBracketing A u}u\rightarrow{}\\[-2pt]
        &\TreeBrackets A{\LeftBracketing A{\WordConcat u{\LetterWord a}}}
          {\WordConcat u{\LetterWord a}}
      \end{aligned}}{\rRI{3,7}}
    \proofstepwidestar[]{%
      \begin{aligned}[t]
        &\forall u\in\NonemptyWords A\,\forall a\in A\,\Bigl(\\[-2pt]
        &\quad\TreeBrackets A{\LeftBracketing A u}u\rightarrow{}\\[-2pt]
        &\quad\TreeBrackets A{\LeftBracketing A{\WordConcat u{\LetterWord a}}}
          {\WordConcat u{\LetterWord a}}\Bigr)
      \end{aligned}}{\rUI{\rRI{1,\rUI{\rRI{2,8}}}}}
  \closeproofpartwideindR

  \Needspace{6\baselineskip}
  \proofpartwide{Wortinduktion}
    \proofstepwidestar[]{%
      \forall a\in A\,
        \TreeBrackets A{\LeftBracketing A{\LetterWord a}}{\LetterWord a}}{%
      \FormulaRefAuto{LeftBracketingPropertyBase}[thm]}
    \proofstepwidestar[]{%
      \begin{aligned}[t]
        &\forall u\in\NonemptyWords A\,\forall a\in A\,\Bigl(\\[-2pt]
        &\quad\TreeBrackets A{\LeftBracketing A u}u\rightarrow{}\\[-2pt]
        &\quad\TreeBrackets A{\LeftBracketing A{\WordConcat u{\LetterWord a}}}
          {\WordConcat u{\LetterWord a}}\Bigr)
      \end{aligned}}{%
      \FormulaRefAuto{LeftBracketingPropertyStep}[thm]}
    \proofstepwidestar[]{%
      \forall w\in\NonemptyWords A\,
        \TreeBrackets A{\LeftBracketing A w}w}{%
      \FormulaRefAuto{NonemptyWordInduction}[thm]{1,2}}
    \proofstepwidestar[4]{w\in\NonemptyWords A}{\rA}
    \proofstepwidestar[4]{\TreeBrackets A{\LeftBracketing A w}w}{%
      \rRE{4,\rUE{3}}}
  \closeproofpartwide
\end{tabproofsplitwide}


\section{Klammerungen eines Wortes}

\FormulaDefDeltaK[Klammerungsmenge eines Nichtleerwortes]{%
  w\in\NonemptyWords{A}\vdash
  \WordBracketings{A}{w}\coloneqq
  \bigl\{T\in\BracketTreeSet{A}\bigm|
    \TreeLeafWord{A}{T}=w\bigr\}%
}{WordBracketingSetDef}{%
  \DeltaRow{Mengen}{A\dsep w\dsep T}
  \DeltaRow{\textbf{Neues Symbol}}{}
  \DeltaPrem{Klammerungsmenge}{\WordBracketings{A}{w}}
}

\FormulaThmDeltaK[Elementkriterium der Klammerungsmenge]{%
  w\in\NonemptyWords{A}\vdash
  \Bigl(T\in\WordBracketings{A}{w}\leftrightarrow
  \bigl(T\in\BracketTreeSet{A}\land\TreeLeafWord{A}{T}=w\bigr)\Bigr)%
}{WordBracketingMembership}{%
  \DeltaRow{Mengen}{A\dsep w\dsep T}
}
\begin{tabproofwide}
  \proofstepwidestar[1]{w\in\NonemptyWords{A}}{\rA}
  \proofstepwidestar[1]{%
    \WordBracketings{A}{w}
      =\{S\in\BracketTreeSet{A}\mid\TreeLeafWord{A}{S}=w\}}{%
    \FormulaRefAuto{WordBracketingSetDef}[def]{1}}
  \proofstepwidestar[]{%
    T\in\{S\in\BracketTreeSet{A}\mid\TreeLeafWord{A}{S}=w\}
      \leftrightarrow
    \bigl(T\in\BracketTreeSet{A}\land\TreeLeafWord{A}{T}=w\bigr)}{%
    \FormulaRefAuto{x\in\{u\in A\mid P(u)\}\eqvdash
      x\in A\land P(x)}[thm]}
  \proofstepwidestar[1]{%
    \begin{aligned}[t]
      &T\in\WordBracketings{A}{w}\leftrightarrow\\[-2pt]
      &\quad\bigl(T\in\BracketTreeSet{A}\land\TreeLeafWord{A}{T}=w\bigr)
    \end{aligned}}{%
    \rIE{2,3}}
\end{tabproofwide}

% Die Konstruktionsrelation verbindet die formalen Konstruktionsregeln mit
% der bereits definierten Klammerungsmenge. Der Linksbaum liefert
% anschließend die Zeugen für die beiden Existenzaussagen.

\FormulaThmDeltaK[Relation und Zugehörigkeit zur Klammerungsmenge]{%
  w\in\NonemptyWords A\vdash
  \bigl(\TreeBrackets A T w\leftrightarrow
    T\in\WordBracketings A w\bigr)%
}{TreeBracketingRelationMembership}{%
  \DeltaRow{Mengen}{A\dsep T\dsep w}
}
\begin{tabproofsplitwide}
  \proofpartwideindR[Von der Relation zur Klammerungsmenge]{%
    w\in\NonemptyWords A\dsep\TreeBrackets A T w\vdash
      T\in\WordBracketings A w
  }{%
    w\in\NonemptyWords A\dsep\TreeBrackets A T w\vdash
      T\in\WordBracketings A w
  }[TreeBracketingRelationToMembership]
    \proofstepwidestar[1]{w\in\NonemptyWords A}{\rA}
    \proofstepwidestar[2]{\TreeBrackets A T w}{\rA}
    \proofstepwidestar[2]{T\in\BracketTreeSet A}{%
      \FormulaRefAuto{TreeBracketingRelationTreeTyping}[thm]{2}}
    \proofstepwidestar[2]{\TreeLeafWord A T=w}{%
      \FormulaRefAuto{TreeBracketingRelationElimination}[thm]{2}}
    \proofstepwidestar[2]{%
      T\in\BracketTreeSet A\land\TreeLeafWord A T=w}{\rAI{3,4}}
    \proofstepwidestar[1]{%
      \begin{aligned}[t]
        &T\in\WordBracketings A w\leftrightarrow\\[-2pt]
        &\quad\bigl(T\in\BracketTreeSet A\land\TreeLeafWord A T=w\bigr)
      \end{aligned}}{%
      \FormulaRefAuto{WordBracketingMembership}[thm]{1}}
    \proofstepwidestar[1,2]{T\in\WordBracketings A w}{%
      \rRE{5,\rLREb{6}}}
  \closeproofpartwideindR

  \proofpartwideindR[Von der Klammerungsmenge zur Relation]{%
    w\in\NonemptyWords A\dsep T\in\WordBracketings A w\vdash
      \TreeBrackets A T w
  }{%
    w\in\NonemptyWords A\dsep T\in\WordBracketings A w\vdash
      \TreeBrackets A T w
  }[TreeBracketingMembershipToRelation]
    \proofstepwidestar[1]{w\in\NonemptyWords A}{\rA}
    \proofstepwidestar[2]{T\in\WordBracketings A w}{\rA}
    \proofstepwidestar[1]{%
      \begin{aligned}[t]
        &T\in\WordBracketings A w\leftrightarrow\\[-2pt]
        &\quad\bigl(T\in\BracketTreeSet A\land\TreeLeafWord A T=w\bigr)
      \end{aligned}}{%
      \FormulaRefAuto{WordBracketingMembership}[thm]{1}}
    \proofstepwidestar[1,2]{%
      T\in\BracketTreeSet A\land\TreeLeafWord A T=w}{%
      \rRE{2,\rLREa{3}}}
    \proofstepwidestar[1,2]{T\in\BracketTreeSet A}{\rAEa{4}}
    \proofstepwidestar[1,2]{\TreeLeafWord A T=w}{\rAEb{4}}
    \proofstepwidestar[1,2]{\TreeBrackets A T w}{%
      \FormulaRefAuto{TreeBracketingRelationIntroduction}[thm]{5,6}}
  \closeproofpartwideindR

  \Needspace{6\baselineskip}
  \proofpartwide{Äquivalenz}
    \proofstepwidestar[1]{w\in\NonemptyWords A}{\rA}
    \proofstepwidestar[2]{\TreeBrackets A T w}{\rA}
    \proofstepwidestar[1,2]{T\in\WordBracketings A w}{%
      \FormulaRefAuto{TreeBracketingRelationToMembership}[thm]{1,2}}
    \proofstepwidestar[1]{%
      \TreeBrackets A T w\rightarrow T\in\WordBracketings A w}{\rRI{2,3}}
    \proofstepwidestar[5]{T\in\WordBracketings A w}{\rA}
    \proofstepwidestar[1,5]{\TreeBrackets A T w}{%
      \FormulaRefAuto{TreeBracketingMembershipToRelation}[thm]{1,5}}
    \proofstepwidestar[1]{%
      T\in\WordBracketings A w\rightarrow\TreeBrackets A T w}{\rRI{5,6}}
    \proofstepwidestar[1]{%
      \TreeBrackets A T w\leftrightarrow T\in\WordBracketings A w}{\rLRI{4,7}}
  \closeproofpartwide
\end{tabproofsplitwide}

\FormulaThmDeltaK[Ein Klammerungsbaum bezeugt die Nichtleerheit]{%
  \TreeBrackets A T w\vdash\WordBracketings A w\neq\varnothing%
}{TreeBracketingRelationNonempty}{%
  \DeltaRow{Mengen}{A\dsep T\dsep w}
}
\begin{tabproofwide}
  \proofstepwidestar[1]{\TreeBrackets A T w}{\rA}
  \proofstepwidestar[1]{w\in\NonemptyWords A}{%
    \FormulaRefAuto{TreeBracketingRelationWordTyping}[thm]{1}}
  \proofstepwidestar[1]{T\in\WordBracketings A w}{%
    \FormulaRefAuto{TreeBracketingRelationToMembership}[thm]{2,1}}
  \proofstepwidestar[1]{\WordBracketings A w\neq\varnothing}{%
    \FormulaRefAuto{a\in A\vdash A\neq\varnothing}[thm]{3}}
\end{tabproofwide}

\Needspace{16\baselineskip}
\FormulaThmDeltaK[Ein Zeuge aus einer nichtleeren Klammerungsmenge]{%
  w\in\NonemptyWords A\dsep\WordBracketings A w\neq\varnothing
  \vdash\exists T\,(\TreeBrackets A T w)%
}{WordBracketingNonemptyRelationWitness}{%
  \DeltaRow{Mengen}{A\dsep T\dsep R\dsep w}
}
\begin{tabproofwide}
  \proofstepwidestar[1]{w\in\NonemptyWords A}{\rA}
  \proofstepwidestar[2]{\WordBracketings A w\neq\varnothing}{\rA}
  \proofstepwidestar[2]{\exists R\,(R\in\WordBracketings A w)}{%
    \FormulaRefAuto{A\neq\varnothing\eqvdash\exists x\,(x\in A)}[thm]{2}}
  \proofstepwidestar[4]{R\in\WordBracketings A w}{\rA}
  \proofstepwidestar[1,4]{\TreeBrackets A R w}{%
    \FormulaRefAuto{TreeBracketingMembershipToRelation}[thm]{1,4}}
  \proofstepwidestar[1,4]{\exists T\,(\TreeBrackets A T w)}{\rEI{5}}
  \proofstepwidestar[1,2]{\exists T\,(\TreeBrackets A T w)}{\rEE{3,4,6}}
\end{tabproofwide}

Der bereits konstruierte Linksbaum klammert sein Wort. Er liefert damit
für jedes nichtleere Wort unmittelbar einen Klammerungszeugen.

\Needspace{12\baselineskip}
\FormulaThmDeltaK[Jedes nichtleere Wort wird von einem Baum geklammert]{%
  w\in\NonemptyWords A\vdash\exists T\,(\TreeBrackets A T w)%
}{TreeBracketingExistence}{%
  \DeltaRow{Mengen}{A\dsep w\dsep T}
  \DeltaRow{Funktionen}{\LeftBracketingMap A}
}
\begin{tabproofwide}
  \proofstepwidestar[1]{w\in\NonemptyWords A}{\rA}
  \proofstepwidestar[1]{\TreeBrackets A{\LeftBracketing A w}w}{%
    \FormulaRefAuto{LeftBracketingProperty}[thm]{1}}
  \proofstepwidestar[1]{\exists T\,(\TreeBrackets A T w)}{\rEI{2}}
\end{tabproofwide}

Derselbe Linksbaum bezeugt auch die Nichtleerheit der Klammerungsmenge.

\Needspace{12\baselineskip}
\FormulaThmDeltaK[Existenz einer Klammerung]{%
  w\in\NonemptyWords A\vdash\WordBracketings A w\neq\varnothing%
}{WordBracketingExistence}{%
  \DeltaRow{Mengen}{A\dsep w}
  \DeltaRow{Funktionen}{\LeftBracketingMap A}
}
\begin{tabproofwide}
  \proofstepwidestar[1]{w\in\NonemptyWords A}{\rA}
  \proofstepwidestar[1]{\TreeBrackets A{\LeftBracketing A w}w}{%
    \FormulaRefAuto{LeftBracketingProperty}[thm]{1}}
  \proofstepwidestar[1]{\WordBracketings A w\neq\varnothing}{%
    \FormulaRefAuto{TreeBracketingRelationNonempty}[thm]{2}}
\end{tabproofwide}


\section{Auswertung eines Klammerungsbaums}

\FormulaThmDeltaK[Existenz und Eindeutigkeit der Baumauswertung]{%
  \begin{aligned}[t]
    &\mathsf{BinOp}(A,\star)\vdash
      \exists!E\colon\BracketTreeSet{A}\to A\,\Bigl(\\[-2pt]
    &\quad \forall a\in A\,
      E(\LeafTree{A}{a})=\Id_A(a)\ \land{}\\[-2pt]
    &\quad \forall S,T\in\BracketTreeSet{A}\,
      E(\NodeTree{A}{S}{T})=
        \BinaryOpFunction{A}{\star}(E(S),E(T))
      \Bigr)
  \end{aligned}%
}{TreeEvaluationExistenceUnique}{%
  \DeltaRow{Mengen}{A\dsep a\dsep S\dsep T}
  \DeltaRow{Funktionen}{\star\dsep E\dsep\Id_A}
}
\begin{tabproofwide}
  \proofstepwidestar[1]{\mathsf{BinOp}(A,\star)}{\rA}
  \proofstepwidestar[]{\Id_A\colon A\to A}{%
    \FormulaRefAuto{IdA}[thm]}
  \proofstepwidestar[1]{%
    \BinaryOpFunction{A}{\star}\colon A\times A\to A}{%
    \FormulaRefAuto{BinaryOperationFunctionDef}[def]{1}}
  \proofstepwidestar[1]{%
    \begin{aligned}[t]
      &\exists!E\colon\BracketTreeSet{A}\to A\,\Bigl(\\[-2pt]
      &\quad \forall a\in A\,
        E(\LeafTree{A}{a})=\Id_A(a)\ \land{}\\[-2pt]
      &\quad \forall S,T\in\BracketTreeSet{A}\,
        E(\NodeTree{A}{S}{T})\\[-2pt]
      &\qquad=\BinaryOpFunction{A}{\star}\\[-2pt]
      &\qquad\quad(E(S),E(T))
        \Bigr)
    \end{aligned}}{%
    \FormulaRefAuto{FullPlanarBinaryTreeRecursion}[thm]{2,3}}
\end{tabproofwide}

\FormulaDefDeltaK[Auswertung eines Klammerungsbaums]{%
  \begin{aligned}[t]
    &\mathsf{BinOp}(A,\star)\vdash \operatorname{ev}_{\star}\coloneqq{}\\[-2pt]
    &\qquad \TreeRecursor{A}{A}{\Id_A}{\BinaryOpFunction{A}{\star}}
  \end{aligned}%
}{TreeEvaluationDef}{%
  \DeltaRow{Mengen}{A\dsep a\dsep S\dsep T}
  \DeltaRow{Funktionen}{\star\dsep\operatorname{ev}_{\star}}
  \DeltaRow{\textbf{Neues Symbol}}{}
  \DeltaPrem{Baumauswertung}{\operatorname{ev}_{\star}}
}

\FormulaThmDeltaK[Typisierung und Rekursionsgleichungen der Baumauswertung]{%
  \begin{aligned}[t]
    &\text{(i)}\;\mathsf{BinOp}(A,\star)\vdash
      \operatorname{ev}_{\star}\colon\BracketTreeSet{A}\to A\\[-2pt]
    &\text{(ii)}\;\mathsf{BinOp}(A,\star)\dsep a\in A\vdash
      \TreeEvaluation{\star}{\LeafTree{A}{a}}=a\\[-2pt]
    &\text{(iii)}\;\mathsf{BinOp}(A,\star)\dsep
      S,T\in\BracketTreeSet{A}\vdash
      \TreeEvaluation{\star}{\NodeTree{A}{S}{T}}
      =\TreeEvaluation{\star}{S}\star\TreeEvaluation{\star}{T}
  \end{aligned}%
}{TreeEvaluationEquations}{%
  \DeltaRow{Mengen}{A\dsep a\dsep S\dsep T}
  \DeltaRow{Funktionen}{\star\dsep\operatorname{ev}_{\star}\dsep R}
}
Im folgenden Beweis bezeichnet \(R\) die rechte Seite von
\FormulaRefAuto{TreeEvaluationDef}[def]. Die Blattregel ersetzt die
Beschriftung; die Knotenregel verknüpft die beiden Teilbaumergebnisse.
\begin{tabproofsplitwide}
  \proofpartwideindR[Typisierung der Baumauswertung]{%
    \mathsf{BinOp}(A,\star)\vdash
      \operatorname{ev}_{\star}\colon\BracketTreeSet{A}\to A
  }{\mathsf{BinOp}(A,\star)\vdash
      \operatorname{ev}_{\star}\colon\BracketTreeSet{A}\to A}%
    [TreeEvaluationTyping]
    \proofstepwidestar[]{%
      \Id_A\colon A\to A}{%
      \FormulaRefAuto{IdA}[thm]}
    \proofstepwidestar[2]{%
      \mathsf{BinOp}\bigl(A,\star\bigr)}{%
      \rA}
    \proofstepwidestar[2]{%
      \operatorname{ev}_{\star}=R}{%
      \FormulaRefAuto{TreeEvaluationDef}[def]{2}}
    \proofstepwidestar[2]{%
      R=\operatorname{ev}_{\star}}{%
      \FormulaRefAuto{a=b\vdash b=a}[thm]{3}}
    \proofstepwidestar[2]{%
      R\colon\BracketTreeSet A\to A}{%
      \FormulaRefAuto{TreeRecursorBinaryOperationTyping}[thm]{1,2}}
    \proofstepwidestar[2]{%
      \operatorname{ev}_{\star}\colon\BracketTreeSet A\to A}{%
      \rIE{4,5}}
  \closeproofpartwideindR

  \proofpartwideindR[Blattgleichung der Baumauswertung]{%
    \mathsf{BinOp}(A,\star)\dsep a\in A\vdash
      \TreeEvaluation{\star}{\LeafTree{A}{a}}=a
  }{\mathsf{BinOp}(A,\star)\dsep a\in A\vdash
      \TreeEvaluation{\star}{\LeafTree{A}{a}}=a}%
    [TreeEvaluationLeafEquation]
    \proofstepwidestar[1]{%
      a\in A}{%
      \rA}
    \proofstepwidestar[]{%
      \Id_A\colon A\to A}{%
      \FormulaRefAuto{IdA}[thm]}
    \proofstepwidestar[3]{%
      \mathsf{BinOp}\bigl(A,\star\bigr)}{%
      \rA}
    \proofstepwidestar[3]{%
      \operatorname{ev}_{\star}=R}{%
      \FormulaRefAuto{TreeEvaluationDef}[def]{3}}
    \proofstepwidestar[3]{%
      R=\operatorname{ev}_{\star}}{%
      \FormulaRefAuto{a=b\vdash b=a}[thm]{4}}
    \proofstepwidestar[1,3]{%
      R(\LeafTree A a)=\Id_A(a)}{%
      \FormulaRefAuto{TreeRecursorBinaryOperationLeafEquation}[thm]{2,3,1}}
    \proofstepwidestar[1]{%
      \Id_A(a)=a}{%
      \FormulaRefAuto{x\in A\vdash\Id_A(x)=x}[thm]{1}}
    \proofstepwidestar[1,3]{%
      \operatorname{ev}_{\star}(\LeafTree A a)=\Id_A(a)}{%
      \rIE{5,6}}
    \proofstepwidestar[1,3]{%
      \operatorname{ev}_{\star}(\LeafTree A a)=a}{%
      \FormulaRefAuto{a=b\dsep b=c\vdash a=c}[thm]{8,7}}
  \closeproofpartwideindR

  \proofpartwideindR[Knotengleichung der Baumauswertung]{%
    \mathsf{BinOp}(A,\star)\dsep S,T\in\BracketTreeSet{A}\vdash
      \TreeEvaluation{\star}{\NodeTree{A}{S}{T}}
      =\TreeEvaluation{\star}{S}\star\TreeEvaluation{\star}{T}
  }{\mathsf{BinOp}(A,\star)\dsep S,T\in\BracketTreeSet{A}\vdash
      \TreeEvaluation{\star}{\NodeTree{A}{S}{T}}
      =\TreeEvaluation{\star}{S}\star\TreeEvaluation{\star}{T}}%
    [TreeEvaluationNodeEquation]
    \proofstepwidestar[1]{%
      S\in\BracketTreeSet A}{%
      \rA}
    \proofstepwidestar[2]{%
      T\in\BracketTreeSet A}{%
      \rA}
    \proofstepwidestar[]{%
      \Id_A\colon A\to A}{%
      \FormulaRefAuto{IdA}[thm]}
    \proofstepwidestar[4]{%
      \mathsf{BinOp}\bigl(A,\star\bigr)}{%
      \rA}
    \proofstepwidestar[4]{%
      \operatorname{ev}_{\star}=R}{%
      \FormulaRefAuto{TreeEvaluationDef}[def]{4}}
    \proofstepwidestar[4]{%
      R=\operatorname{ev}_{\star}}{%
      \FormulaRefAuto{a=b\vdash b=a}[thm]{5}}
    \proofstepwidestar[1,2,4]{%
      R(\NodeTree A S T)=R(S)\star R(T)}{%
      \FormulaRefAuto{TreeRecursorBinaryOperationNodeEquation}[thm]{3,4,\rAI{1,2}}}
    \proofstepwidestar[1,2,4]{%
      \begin{aligned}[t]&\operatorname{ev}_{\star}(\NodeTree A S T)\\[-2pt]
      &\quad=\operatorname{ev}_{\star}(S)\star \operatorname{ev}_{\star}(T)\end{aligned}}{%
      \rIE{6,7}}
  \closeproofpartwideindR
\end{tabproofsplitwide}

% Einbindung nach TreeEvaluationEquations.

\subsection{Auswertung des Linksbaums}

Beim Rechtsanfügen eines Blattes wird der bisherige Wert mit dessen
Buchstaben verknüpft. Deshalb stimmt die Auswertung des Linksbaums mit
der bereits definierten Linksfaltung überein. Die Aussage setzt nur
einen binären Operator voraus.

\FormulaThmDeltaK[Auswertung einer Blattanfügung]{%
  \begin{aligned}[t]
    &\mathsf{BinOp}(A,\star)\dsep T\in\BracketTreeSet A\dsep a\in A
      \vdash{}\\[-2pt]
    &\TreeEvaluation\star{\TreeAppendLeaf A T a}
      =\TreeEvaluation\star T\star a
  \end{aligned}
}{TreeAppendLeafEvaluation}{%
  \DeltaRow{Mengen}{A\dsep T\dsep a}
  \DeltaRow{Funktionen}{\star\dsep\operatorname{ev}_\star}
}
\begin{tabproofwide}
  \proofstepwidestar[1]{\mathsf{BinOp}(A,\star)}{\rA}
  \proofstepwidestar[2]{T\in\BracketTreeSet A}{\rA}
  \proofstepwidestar[3]{a\in A}{\rA}
  \proofstepwidestar[1]{\operatorname{ev}_\star\colon\BracketTreeSet A\to A}{%
    \FormulaRefAuto{TreeEvaluationTyping}[thm]{1}}
  \proofstepwidestar[1,2]{\TreeEvaluation\star T\in A}{%
    \FormulaRefAuto{F\colon A\to B\dsep x\in A\vdash F(x)\in B}[thm]{4,2}}
  \proofstepwidestar[3]{\LeafTree A a\in\BracketTreeSet A}{%
    \FormulaRefAuto{TreeLeafClosure}[thm]{3}}
  \proofstepwidestar[2,3]{T,\LeafTree A a\in\BracketTreeSet A}{\rAI{2,6}}
  \proofstepwidestar[1,2,3]{%
    \TreeEvaluation\star{\NodeTree A T{\LeafTree A a}}
      =\TreeEvaluation\star T\star\TreeEvaluation\star{\LeafTree A a}}{%
    \FormulaRefAuto{TreeEvaluationNodeEquation}[thm]{1,7}}
  \proofstepwidestar[1,3]{\TreeEvaluation\star{\LeafTree A a}=a}{%
    \FormulaRefAuto{TreeEvaluationLeafEquation}[thm]{1,3}}
  \proofstepwidestar[1,2,3]{%
    \TreeEvaluation\star{\NodeTree A T{\LeafTree A a}}
      =\TreeEvaluation\star T\star a}{\rIE{9,8}}
  \proofstepwidestar[2,3]{%
    \TreeAppendLeaf A T a=\NodeTree A T{\LeafTree A a}}{%
    \FormulaRefAuto{TreeAppendLeafDef}[def]{2,3}}
  \proofstepwidestar[2,3]{%
    \NodeTree A T{\LeafTree A a}=\TreeAppendLeaf A T a}{%
    \FormulaRefAuto{a=b\vdash b=a}[thm]{11}}
  \proofstepwidestar[1,2,3]{%
    \TreeEvaluation\star{\TreeAppendLeaf A T a}
      =\TreeEvaluation\star T\star a}{\rIE{12,10}}
\end{tabproofwide}

\FormulaThmDeltaK[Die Auswertung des Linksbaums ist die Linksfaltung]{%
  \mathsf{BinOp}(A,\star)\dsep w\in\NonemptyWords A\vdash
    \TreeEvaluation\star{\LeftBracketing A w}=\WordFold\star w
}{LeftBracketingEvaluation}{%
  \DeltaRow{Mengen}{A\dsep w\dsep u\dsep a}
  \DeltaRow{Funktionen}{\star\dsep\operatorname{ev}_\star\dsep
    \Pi_\star\dsep\LeftBracketingMap A}
}
\begin{tabproofsplitwide}
  \proofpartwideindR[Blatt und Buchstabenwort haben denselben Wert]{%
    \mathsf{BinOp}(A,\star)\vdash
      \forall a\in A\,
        \TreeEvaluation\star{\LeftBracketing A{\LetterWord a}}
          =\WordFold\star{\LetterWord a}
  }{%
    \mathsf{BinOp}(A,\star)\vdash
      \forall a\in A\,
        \TreeEvaluation\star{\LeftBracketing A{\LetterWord a}}
          =\WordFold\star{\LetterWord a}
  }[LeftBracketingEvaluationBase]
    \proofstepwidestar[1]{\mathsf{BinOp}(A,\star)}{\rA}
    \proofstepwidestar[2]{a\in A}{\rA}
    \proofstepwidestar[2]{\LetterWord a\in\NonemptyWords A}{%
      \FormulaRefAuto{LetterWordNonempty}[thm]{2}}
    \proofstepwidestar[2]{%
      \LeftBracketing A{\LetterWord a}\in\BracketTreeSet A}{%
      \FormulaRefAuto{LeftBracketingValueTyping}[thm]{3}}
    \proofstepwidestar[2]{\LeftBracketing A{\LetterWord a}=\LeafTree A a}{%
      \FormulaRefAuto{LeftBracketingLetterEquation}[thm]{2}}
    \proofstepwidestar[2]{\LeafTree A a=\LeftBracketing A{\LetterWord a}}{%
      \FormulaRefAuto{a=b\vdash b=a}[thm]{5}}
    \proofstepwidestar[1,2]{\TreeEvaluation\star{\LeafTree A a}=a}{%
      \FormulaRefAuto{TreeEvaluationLeafEquation}[thm]{1,2}}
    \proofstepwidestar[1,2]{%
      \TreeEvaluation\star{\LeftBracketing A{\LetterWord a}}=a}{\rIE{6,7}}
    \proofstepwidestar[1,2]{\WordFold\star{\LetterWord a}=a}{%
      \FormulaRefAuto{LeftWordFoldLetter}[thm]{1,2}}
    \proofstepwidestar[1,2]{a=\WordFold\star{\LetterWord a}}{%
      \FormulaRefAuto{a=b\vdash b=a}[thm]{9}}
    \proofstepwidestar[1,2]{%
      \TreeEvaluation\star{\LeftBracketing A{\LetterWord a}}
        =\WordFold\star{\LetterWord a}}{\rChain{8,10}}
    \proofstepwidestar[1]{%
      \forall a\in A\,
        \TreeEvaluation\star{\LeftBracketing A{\LetterWord a}}
          =\WordFold\star{\LetterWord a}}{\rUI{\rRI{2,11}}}
  \closeproofpartwideindR

  \proofpartwideindR[Denselben Buchstaben verknüpfen]{%
    \begin{aligned}[t]
      &\mathsf{BinOp}(A,\star)\vdash{}\\[-2pt]
      &\forall u\in\NonemptyWords A\,\forall a\in A\,\Bigl(\\[-2pt]
      &\quad\TreeEvaluation\star{\LeftBracketing A u}=\WordFold\star u
        \rightarrow{}\\[-2pt]
      &\quad\TreeEvaluation\star{\LeftBracketing A{\WordConcat u{\LetterWord a}}}
        =\WordFold\star{\WordConcat u{\LetterWord a}}\Bigr)
    \end{aligned}
  }{%
    \begin{aligned}[t]
      &\mathsf{BinOp}(A,\star)\vdash{}\\[-2pt]
      &\forall u\in\NonemptyWords A\,\forall a\in A\,\Bigl(\\[-2pt]
      &\quad\TreeEvaluation\star{\LeftBracketing A u}=\WordFold\star u
        \rightarrow{}\\[-2pt]
      &\quad\TreeEvaluation\star{\LeftBracketing A{\WordConcat u{\LetterWord a}}}
        =\WordFold\star{\WordConcat u{\LetterWord a}}\Bigr)
    \end{aligned}
  }[LeftBracketingEvaluationStep]
    \proofstepwidestar[1]{\mathsf{BinOp}(A,\star)}{\rA}
    \proofstepwidestar[2]{u\in\NonemptyWords A}{\rA}
    \proofstepwidestar[3]{a\in A}{\rA}
    \proofstepwidestar[4]{%
      \TreeEvaluation\star{\LeftBracketing A u}=\WordFold\star u}{\rA}
    \proofstepwidestar[2]{\LeftBracketing A u\in\BracketTreeSet A}{%
      \FormulaRefAuto{LeftBracketingValueTyping}[thm]{2}}
    \proofstepwidestar[1]{\Pi_\star\colon\NonemptyWords A\to A}{%
      \FormulaRefAuto{LeftWordFoldFunction}[thm]{1}}
    \proofstepwidestar[1,2]{\WordFold\star u\in A}{%
      \FormulaRefAuto{F\colon A\to B\dsep x\in A\vdash F(x)\in B}[thm]{6,2}}
    \proofstepwidestar[1,2,3]{%
      \TreeEvaluation\star{\TreeAppendLeaf A{\LeftBracketing A u}a}
        =\TreeEvaluation\star{\LeftBracketing A u}\star a}{%
      \FormulaRefAuto{TreeAppendLeafEvaluation}[thm]{1,5,3}}
    \proofstepwidestar[1,2,3,4]{%
      \TreeEvaluation\star{\TreeAppendLeaf A{\LeftBracketing A u}a}
        =\WordFold\star u\star a}{\rIE{4,8}}
    \proofstepwidestar[2,3]{%
      \LeftBracketing A{\WordConcat u{\LetterWord a}}
        =\TreeAppendLeaf A{\LeftBracketing A u}a}{%
      \FormulaRefAuto{LeftBracketingAppendEquation}[thm]{2,3}}
    \proofstepwidestar[2,3]{%
      \TreeAppendLeaf A{\LeftBracketing A u}a
        =\LeftBracketing A{\WordConcat u{\LetterWord a}}}{%
      \FormulaRefAuto{a=b\vdash b=a}[thm]{10}}
    \proofstepwidestar[1,2,3,4]{%
      \TreeEvaluation\star{\LeftBracketing A{\WordConcat u{\LetterWord a}}}
        =\WordFold\star u\star a}{\rIE{11,9}}
    \proofstepwidestar[1,2,3]{%
      \WordFold\star{\WordConcat u{\LetterWord a}}=\WordFold\star u\star a}{%
      \FormulaRefAuto{LeftWordFoldStep}[thm]{1,2,3}}
    \proofstepwidestar[1,2,3]{%
      \WordFold\star u\star a=\WordFold\star{\WordConcat u{\LetterWord a}}}{%
      \FormulaRefAuto{a=b\vdash b=a}[thm]{13}}
    \proofstepwidestar[1,2,3,4]{%
      \TreeEvaluation\star{\LeftBracketing A{\WordConcat u{\LetterWord a}}}
        =\WordFold\star{\WordConcat u{\LetterWord a}}}{\rChain{12,14}}
    \proofstepwidestar[1,2,3]{%
      \begin{aligned}[t]
        &\TreeEvaluation\star{\LeftBracketing A u}=\WordFold\star u
          \rightarrow{}\\[-2pt]
        &\TreeEvaluation\star{\LeftBracketing A{\WordConcat u{\LetterWord a}}}
          =\WordFold\star{\WordConcat u{\LetterWord a}}
      \end{aligned}}{\rRI{4,15}}
    \proofstepwidestar[1]{%
      \begin{aligned}[t]
        &\forall u\in\NonemptyWords A\,\forall a\in A\,\Bigl(\\[-2pt]
        &\quad\TreeEvaluation\star{\LeftBracketing A u}=\WordFold\star u
          \rightarrow{}\\[-2pt]
        &\quad\TreeEvaluation\star{\LeftBracketing A{\WordConcat u{\LetterWord a}}}
          =\WordFold\star{\WordConcat u{\LetterWord a}}\Bigr)
      \end{aligned}}{\rUI{\rRI{2,\rUI{\rRI{3,16}}}}}
  \closeproofpartwideindR

  \Needspace{6\baselineskip}
  \proofpartwide{Wortinduktion}
    \proofstepwidestar[1]{\mathsf{BinOp}(A,\star)}{\rA}
    \proofstepwidestar[2]{w\in\NonemptyWords A}{\rA}
    \proofstepwidestar[1]{%
      \forall a\in A\,
        \TreeEvaluation\star{\LeftBracketing A{\LetterWord a}}
          =\WordFold\star{\LetterWord a}}{%
      \FormulaRefAuto{LeftBracketingEvaluationBase}[thm]{1}}
    \proofstepwidestar[1]{%
      \begin{aligned}[t]
        &\forall u\in\NonemptyWords A\,\forall a\in A\,\Bigl(\\[-2pt]
        &\quad\TreeEvaluation\star{\LeftBracketing A u}=\WordFold\star u
          \rightarrow{}\\[-2pt]
        &\quad\TreeEvaluation\star{\LeftBracketing A{\WordConcat u{\LetterWord a}}}
          =\WordFold\star{\WordConcat u{\LetterWord a}}\Bigr)
      \end{aligned}}{%
      \FormulaRefAuto{LeftBracketingEvaluationStep}[thm]{1}}
    \proofstepwidestar[1]{%
      \forall w\in\NonemptyWords A\,
        \TreeEvaluation\star{\LeftBracketing A w}=\WordFold\star w}{%
      \FormulaRefAuto{NonemptyWordInduction}[thm]{3,4}}
    \proofstepwidestar[1,2]{%
      \TreeEvaluation\star{\LeftBracketing A w}=\WordFold\star w}{%
      \rRE{2,\rUE{5}}}
  \closeproofpartwide
\end{tabproofsplitwide}


% Abschließende Folgerungen aus den früh eingeführten Strukturaxiomen.
\chapter{Eindeutigkeit und Wiederverwendung der Strukturen}

Die Wortaxiome wurden unmittelbar nach dem Aufbau der elementaren
Anfügung eingeführt, die Baumaxiome unmittelbar nach dem Nachweis der
Konstruktor- und Trägereigenschaften. Die Wortinduktion steht durch W3
axiomatisch bereit; die Bauminduktion und die Rekursionssätze sind aus
den jeweiligen Strukturaxiomen bewiesen. Abschließend
werden ihre Konsequenzen für den Darstellungswechsel und die
Wiederverwendung zusammengeführt.

\section{Eindeutigkeit der Strukturen und das leere Alphabet}

Ein \emph{strukturtreuer Isomorphismus} ist hier eine Bijektion, welche
den Anfang und die Anfügung beziehungsweise sämtliche Blattbeschriftungen
und das geordnete Zusammenfügen erhält. Das Alphabet wird dabei
festgehalten; seine Buchstaben werden nicht umbenannt.

\FormulaThmDeltaK[Eindeutiger Isomorphismus freier Wortstrukturen]{%
  \begin{aligned}[t]
    &\mathsf{WortStr}_A(W,e,s)\dsep
      \mathsf{WortStr}_A(V,d,r)\vdash{}\\[-2pt]
    &\exists!j\;\bigl(j\colon W\bij V\land
      \mathsf{WRec}_A(j;W,e,s;V,d,r)\bigr)
  \end{aligned}
}{WordStructureUniqueIsomorphism}{%
  \DeltaRow{Mengen}{A\dsep W\dsep V\dsep e\dsep d}
  \DeltaRow{Funktionen}{s\dsep r\dsep h\dsep h'\dsep j}}
\begin{tabproofwide}
  \proofstepwidestar[1]{\mathsf{WortStr}_A(W,e,s)}{%
    \rA}
  \proofstepwidestar[2]{\mathsf{WortStr}_A(V,d,r)}{%
    \rA}
  \proofstepwidestar[2]{d\in V\land r\colon V\times A\to V}{%
    \FormulaRefAuto{WordStructureTypingAxiom}[ax]{2}}
  \proofstepwidestar[1,2]{\exists!j\;\mathsf{WRec}_A(j;W,e,s;V,d,r)}{%
    \FormulaRefAuto{WordStructureRecursion}[thm]{1,\rAEa{3},\rAEb{3}}}
  \proofstepwidestar[1,2]{\exists j\;\mathsf{WRec}_A(j;W,e,s;V,d,r)}{%
    \FormulaRefAuto{\exists! x\,P(x)\vdash\exists x\,P(x)}[thm]{4}}
  \proofstepwidestar[6]{\mathsf{WRec}_A(j;W,e,s;V,d,r)}{%
    \rA}
  \proofstepwidestar[1,2,6]{j\colon W\bij V}{%
    \FormulaRefAuto{WordStructureRecursorIsomorphism}[thm]{1,2,6}}
  \proofstepwidestar[1,2,6]{\exists j\;(j\colon W\bij V\land\mathsf{WRec}_A(j;W,e,s;V,d,r))}{%
    \rEI{\rAI{7,6}}}
  \proofstepwidestar[1,2]{\exists j\;(j\colon W\bij V\land\mathsf{WRec}_A(j;W,e,s;V,d,r))}{%
    \rEE{5,6,8}}
  \proofstepwidestar[10]{j\colon W\bij V\land\mathsf{WRec}_A(j;W,e,s;V,d,r)}{%
    \rA}
  \proofstepwidestar[11]{h\colon W\bij V\land\mathsf{WRec}_A(h;W,e,s;V,d,r)}{%
    \rA}
  \proofstepwidestar[1,2,10,11]{j=h}{%
    \FormulaRefAuto{WordRecursionTwoSolutionsEqual}[thm]{1,\rAEa{3},\rAEb{3},\rAEb{10},\rAEb{11}}}
  \proofstepwidestar[1,2]{\exists!j\;(j\colon W\bij V\land\mathsf{WRec}_A(j;W,e,s;V,d,r))}{%
    \UEI{9,10,11,12}}
\end{tabproofwide}

\FormulaThmDeltaK[Eindeutiger Isomorphismus freier Baumstrukturen]{%
  \begin{aligned}[t]
    &\mathsf{BaumStr}_A(T,\ell,k)\dsep
      \mathsf{BaumStr}_A(V,m,n)\vdash{}\\[-2pt]
    &\exists!j\;\bigl(j\colon T\bij V\land
      \mathsf{TRec}_A(j;T,\ell,k;V,m,n)\bigr)
  \end{aligned}
}{BinaryTreeStructureUniqueIsomorphism}{%
  \DeltaRow{Mengen}{A\dsep T\dsep V}
  \DeltaRow{Funktionen}{\ell\dsep k\dsep m\dsep n\dsep h\dsep h'\dsep j}}
\begin{tabproofwide}
  \proofstepwidestar[1]{\mathsf{BaumStr}_A(T,\ell,k)}{%
    \rA}
  \proofstepwidestar[2]{\mathsf{BaumStr}_A(V,m,n)}{%
    \rA}
  \proofstepwidestar[2]{m\colon A\to V\land n\colon V\times V\to V}{%
    \FormulaRefAuto{BinaryTreeStructureTypingAxiom}[ax]{2}}
  \proofstepwidestar[1,2]{\exists!j\;\mathsf{TRec}_A(j;T,\ell,k;V,m,n)}{%
    \FormulaRefAuto{BinaryTreeStructureRecursion}[thm]{1,\rAEa{3},\rAEb{3}}}
  \proofstepwidestar[1,2]{\exists j\;\mathsf{TRec}_A(j;T,\ell,k;V,m,n)}{%
    \FormulaRefAuto{\exists! x\,P(x)\vdash\exists x\,P(x)}[thm]{4}}
  \proofstepwidestar[6]{\mathsf{TRec}_A(j;T,\ell,k;V,m,n)}{%
    \rA}
  \proofstepwidestar[1,2,6]{j\colon T\bij V}{%
    \FormulaRefAuto{BinaryTreeStructureRecursorIsomorphism}[thm]{1,2,6}}
  \proofstepwidestar[1,2,6]{\exists j\;(j\colon T\bij V\land\mathsf{TRec}_A(j;T,\ell,k;V,m,n))}{%
    \rEI{\rAI{7,6}}}
  \proofstepwidestar[1,2]{\exists j\;(j\colon T\bij V\land\mathsf{TRec}_A(j;T,\ell,k;V,m,n))}{%
    \rEE{5,6,8}}
  \proofstepwidestar[10]{j\colon T\bij V\land\mathsf{TRec}_A(j;T,\ell,k;V,m,n)}{%
    \rA}
  \proofstepwidestar[11]{h\colon T\bij V\land\mathsf{TRec}_A(h;T,\ell,k;V,m,n)}{%
    \rA}
  \proofstepwidestar[1,10,11]{j=h}{%
    \FormulaRefAuto{BinaryTreeRecursionTwoSolutionsEqual}[thm]{1,\rAEb{10},\rAEb{11}}}
  \proofstepwidestar[1,2]{\exists!j\;(j\colon T\bij V\land\mathsf{TRec}_A(j;T,\ell,k;V,m,n))}{%
    \UEI{9,10,11,12}}
\end{tabproofwide}

\Needspace{12\baselineskip}
\FormulaThmDeltaK[Freie Strukturen über dem leeren Alphabet]{%
  \begin{aligned}[t]
    &\text{(i)}\quad\mathsf{WortStr}_{\varnothing}(W,e,s)\vdash W=\{e\},\\[-2pt]
    &\text{(ii)}\quad\mathsf{BaumStr}_{\varnothing}(T,\ell,k)\vdash T=\varnothing.
  \end{aligned}
}{FreeStructuresEmptyAlphabet}{%
  \DeltaRow{Mengen}{W\dsep T\dsep e\dsep u\dsep a\dsep x\dsep y}\DeltaRow{Funktionen}{s\dsep\ell\dsep k}}
\begin{tabproofsplitwide}
 \proofpartwideindR[Wortstruktur]{\mathsf{WortStr}_{\varnothing}(W,e,s)\vdash W=\{e\}}
 {\mathsf{WortStr}_{\varnothing}(W,e,s)\vdash W=\{e\}}[WordStructureEmptyAlphabet]
  \proofstepwidestar[1]{\mathsf{WortStr}_{\varnothing}(W,e,s)}{%
    \rA}
  \proofstepwidestar[1]{e\in W}{%
    \rAEa{\FormulaRefAuto{WordStructureTypingAxiom}[ax]{1}}}
  \proofstepwidestar[1]{\{e\}\subseteq W}{%
    \FormulaRefAuto{x\in A\vdash\{x\}\subseteq A}[thm]{2}}
  \proofstepwidestar[]{e\in\{e\}}{%
    \FormulaRefAuto{x\in\{x\}}[thm]}
  \proofstepwidestar[5]{u\in\{e\}}{%
    \rA}
  \proofstepwidestar[6]{a\in\varnothing}{%
    \rA}
  \proofstepwidestar[]{a\notin\varnothing}{%
    \FormulaRefAuto{x\not\in\varnothing}[thm]}
  \proofstepwidestar[6]{s(u,a)\in\{e\}}{%
    \FormulaRefAuto{A\dsep\neg A\vdash B}[thm]{6,7}}
  \proofstepwidestar[]{\forall u\in\{e\}\,\forall a\in\varnothing\;s(u,a)\in\{e\}}{%
    \rUI{\rRI{5,\rUI{\rRI{6,8}}}}}
  \proofstepwidestar[1]{\{e\}=W}{%
    \FormulaRefAuto{WordStructureMinimalityAxiom}[ax]{1,3,4,9}}
  \proofstepwidestar[1]{W=\{e\}}{%
    \FormulaRefAuto{a=b\vdash b=a}[thm]{10}}
 \closeproofpartwideindR
 \Needspace{14\baselineskip}
 \proofpartwideindR[Baumstruktur]{\mathsf{BaumStr}_{\varnothing}(T,\ell,k)\vdash T=\varnothing}
 {\mathsf{BaumStr}_{\varnothing}(T,\ell,k)\vdash T=\varnothing}[BinaryTreeStructureEmptyAlphabet]
  \proofstepwidestar[1]{\mathsf{BaumStr}_{\varnothing}(T,\ell,k)}{%
    \rA}
  \proofstepwidestar[1]{\ell\colon\varnothing\to T}{%
    \rAEa{\FormulaRefAuto{BinaryTreeStructureTypingAxiom}[ax]{1}}}
  \proofstepwidestar[]{\varnothing\subseteq T}{%
    \FormulaRefAuto{\varnothing\subseteq A}[thm]}
  \proofstepwidestar[1]{\ell[\varnothing]=\varnothing}{%
    \FormulaRefAuto{F\colon A\to B\vdash F[\varnothing]=\varnothing}[thm]{2}}
  \proofstepwidestar[]{\varnothing\subseteq\varnothing}{%
    \FormulaRefAuto{A\subseteq A}[thm]}
  \proofstepwidestar[1]{\ell[\varnothing]\subseteq\varnothing}{%
    \rIE{\FormulaRefAuto{a=b\vdash b=a}[thm]{4},5}}
  \proofstepwidestar[7]{x\in\varnothing}{%
    \rA}
  \proofstepwidestar[8]{y\in\varnothing}{%
    \rA}
  \proofstepwidestar[]{x\notin\varnothing}{%
    \FormulaRefAuto{x\not\in\varnothing}[thm]}
  \proofstepwidestar[7]{k(x,y)\in\varnothing}{%
    \FormulaRefAuto{A\dsep\neg A\vdash B}[thm]{7,9}}
  \proofstepwidestar[]{\forall x,y\in\varnothing\;k(x,y)\in\varnothing}{%
    \rUI{\rRI{7,\rUI{\rRI{8,10}}}}}
  \proofstepwidestar[1]{\varnothing=T}{%
    \FormulaRefAuto{BinaryTreeStructureMinimalityAxiom}[ax]{1,3,6,11}}
  \proofstepwidestar[1]{T=\varnothing}{%
    \FormulaRefAuto{a=b\vdash b=a}[thm]{12}}
 \closeproofpartwideindR
\end{tabproofsplitwide}

Der Unterschied ist beabsichtigt: Ein Wort darf leer sein; ein endlicher
voller Klammerungsbaum besitzt ein Blatt. Ohne Blattbeschriftungen gibt es
daher keinen solchen Baum. Die Rekursionssätze gelten auch in diesen
Grenzfällen. Bei Wörtern bleibt der vorgegebene Anfangswert erforderlich;
bei der leeren Baumstruktur ist die gesuchte Funktion die leere Funktion.

\section{Wiederverwendung in Algebra und formaler Syntax}

Die Axiome fassen den Band zu einer wiederverwendbaren Grundlage zusammen.
Die axiomatische Wortinduktion liefert die Existenz der Zerlegung;
Trennung und Injektivität der Konstruktoren sichern die Eindeutigkeit.
Bei Bäumen liefert die Minimalität das entsprechende Induktionsprinzip.
Die Rekursionssätze wurden direkt
aus den Strukturaxiomen bewiesen; auf ihnen beruht der eindeutige Transport. Ein späterer Beweis darf
deshalb die Strukturgesetze benutzen, ohne die Positionen eines Wortes
oder den Wortcode eines Baumes erneut auszuwerten.

Beispielsweise erhält die Baumrekursion mit
\(f(a)=\LetterWord a\) und \(g(u,v)=\WordConcat u v\) das Blattwort in
jeder freien Baumstruktur. Mit einer beliebigen binären Operation auf
\(X\) erhält man ihre Auswertung. Die Eindeutigkeit der Rekursion zeigt,
dass diese Funktionen unter dem eindeutigen Isomorphismus mit den früher
definierten Blattwort- und Auswertungsfunktionen übereinstimmen. Die
bereits bewiesenen Aussagen zu Klammerungen bleiben damit in jeder
Darstellung gültig.

Für Band~48, \emph{Axiomatische Mengenlehre II}, ist dieselbe Trennung
zwischen Konstruktion und Strukturgesetzen nützlich. Eine Formel besitzt
einen äußeren Konstruktor und dessen unmittelbare Teilformeln. Die
entsprechenden Regeln für Syntax müssen allerdings die tatsächlichen
Stelligkeiten berücksichtigen: Negation hat eine Teilformel, eine binäre
Verknüpfung zwei, und ein Quantor zusätzlich einen Variablenparameter.
Atomare Formeln tragen ihre eigenen Parameter. Die hier behandelte
Binärbaumstruktur beschreibt unmittelbar die volle binäre
Klammerungsstruktur; sie ist noch kein vollständiges Modell dieser
verschiedenen Formelkonstruktoren.

Für eine solche Syntax wird das vorliegende Muster auf die betreffende
Familie von Konstruktoren übertragen: Jeder Konstruktor ist auf seinen
Parametern und Teilobjekten injektiv, die Bilder verschiedener
Konstruktoren sind disjunkt, und jede unter allen Konstruktoren
abgeschlossene Teilmenge ist die gesamte Syntaxmenge. Nach Konstruktion
eines passenden Modells sind Strukturinduktion und Strukturrekursion
mit je einem Fall pro Konstruktor zu beweisen. Dabei dürfen Identifikationen
wie Gleichheit bis zur Umbenennung gebundener Variablen erst in einem
gesonderten Schritt eingeführt werden; sie verändern die Aussage der
eindeutigen Zerlegung.

Die Verbindung mit den Graphen- und Baumbänden liefert dazu die
geometrische Sicht: Vorkommen von Teilobjekten sind Knoten, die Beziehung
zum unmittelbaren Teilobjekt gibt gerichtete Kanten, und das Vergessen
ihrer Richtung gibt den zugrunde liegenden ungerichteten Baum.
Verschiedene Vorkommen gleicher Teilobjekte bleiben dabei verschiedene
Knoten. Die hier bewiesenen Strukturgesetze erklären, welche Information
die geordnete und beschriftete Baumdarstellung für spätere rekursive
Definitionen bewahren muss.


% Die bandlokalen Tabellenanpassungen dürfen beim Gesamtsatz nicht in die
% nachfolgenden Bände hineinwirken.
\let\RuleRefWithArg\BandTwentySevenOriginalRuleRefWithArg
\ExplSyntaxOn
\cs_set_eq:NN
  \mx_formula_ref_suffix:n
  \mx_band_twentyseven_original_formula_ref_suffix:n
\ExplSyntaxOff
\makeatletter
\newcol@{G}[0]{>{\scriptsize}r}
\makeatother

\end{document}
